% AUTO-GENERATED FROM MATH/Math.md BY build_math_from_md.py.
% Edit Math.md or the builder, then regenerate.
% Options for packages loaded elsewhere
\PassOptionsToPackage{unicode}{hyperref}
\PassOptionsToPackage{hyphens}{url}
\documentclass[
  10pt,
]{article}
\usepackage{xcolor}
\usepackage[margin=0.8in]{geometry}
\usepackage{amsmath,amssymb}
\setcounter{secnumdepth}{-\maxdimen} % remove section numbering
\usepackage{iftex}
\ifPDFTeX
  \usepackage[T1]{fontenc}
  \usepackage[utf8]{inputenc}
  \usepackage{textcomp} % provide euro and other symbols
\else % if luatex or xetex
  \usepackage{unicode-math} % this also loads fontspec
  \defaultfontfeatures{Scale=MatchLowercase}
  \defaultfontfeatures[\rmfamily]{Ligatures=TeX,Scale=1}
\fi
\usepackage{lmodern}
\ifPDFTeX\else
  % xetex/luatex font selection
\fi
% Use upquote if available, for straight quotes in verbatim environments
\IfFileExists{upquote.sty}{\usepackage{upquote}}{}
\IfFileExists{microtype.sty}{% use microtype if available
  \usepackage[]{microtype}
  \UseMicrotypeSet[protrusion]{basicmath} % disable protrusion for tt fonts
}{}
\makeatletter
\@ifundefined{KOMAClassName}{% if non-KOMA class
  \IfFileExists{parskip.sty}{%
    \usepackage{parskip}
  }{% else
    \setlength{\parindent}{0pt}
    \setlength{\parskip}{6pt plus 2pt minus 1pt}}
}{% if KOMA class
  \KOMAoptions{parskip=half}}
\makeatother
\usepackage{color}
\usepackage{fancyvrb}
\newcommand{\VerbBar}{|}
\newcommand{\VERB}{\Verb[commandchars=\\\{\}]}
\DefineVerbatimEnvironment{Highlighting}{Verbatim}{commandchars=\\\{\}}
% Add ',fontsize=\small' for more characters per line
\newenvironment{Shaded}{}{}
\newcommand{\AlertTok}[1]{\textcolor[rgb]{1.00,0.00,0.00}{\textbf{#1}}}
\newcommand{\AnnotationTok}[1]{\textcolor[rgb]{0.38,0.63,0.69}{\textbf{\textit{#1}}}}
\newcommand{\AttributeTok}[1]{\textcolor[rgb]{0.49,0.56,0.16}{#1}}
\newcommand{\BaseNTok}[1]{\textcolor[rgb]{0.25,0.63,0.44}{#1}}
\newcommand{\BuiltInTok}[1]{\textcolor[rgb]{0.00,0.50,0.00}{#1}}
\newcommand{\CharTok}[1]{\textcolor[rgb]{0.25,0.44,0.63}{#1}}
\newcommand{\CommentTok}[1]{\textcolor[rgb]{0.38,0.63,0.69}{\textit{#1}}}
\newcommand{\CommentVarTok}[1]{\textcolor[rgb]{0.38,0.63,0.69}{\textbf{\textit{#1}}}}
\newcommand{\ConstantTok}[1]{\textcolor[rgb]{0.53,0.00,0.00}{#1}}
\newcommand{\ControlFlowTok}[1]{\textcolor[rgb]{0.00,0.44,0.13}{\textbf{#1}}}
\newcommand{\DataTypeTok}[1]{\textcolor[rgb]{0.56,0.13,0.00}{#1}}
\newcommand{\DecValTok}[1]{\textcolor[rgb]{0.25,0.63,0.44}{#1}}
\newcommand{\DocumentationTok}[1]{\textcolor[rgb]{0.73,0.13,0.13}{\textit{#1}}}
\newcommand{\ErrorTok}[1]{\textcolor[rgb]{1.00,0.00,0.00}{\textbf{#1}}}
\newcommand{\ExtensionTok}[1]{#1}
\newcommand{\FloatTok}[1]{\textcolor[rgb]{0.25,0.63,0.44}{#1}}
\newcommand{\FunctionTok}[1]{\textcolor[rgb]{0.02,0.16,0.49}{#1}}
\newcommand{\ImportTok}[1]{\textcolor[rgb]{0.00,0.50,0.00}{\textbf{#1}}}
\newcommand{\InformationTok}[1]{\textcolor[rgb]{0.38,0.63,0.69}{\textbf{\textit{#1}}}}
\newcommand{\KeywordTok}[1]{\textcolor[rgb]{0.00,0.44,0.13}{\textbf{#1}}}
\newcommand{\NormalTok}[1]{#1}
\newcommand{\OperatorTok}[1]{\textcolor[rgb]{0.40,0.40,0.40}{#1}}
\newcommand{\OtherTok}[1]{\textcolor[rgb]{0.00,0.44,0.13}{#1}}
\newcommand{\PreprocessorTok}[1]{\textcolor[rgb]{0.74,0.48,0.00}{#1}}
\newcommand{\RegionMarkerTok}[1]{#1}
\newcommand{\SpecialCharTok}[1]{\textcolor[rgb]{0.25,0.44,0.63}{#1}}
\newcommand{\SpecialStringTok}[1]{\textcolor[rgb]{0.73,0.40,0.53}{#1}}
\newcommand{\StringTok}[1]{\textcolor[rgb]{0.25,0.44,0.63}{#1}}
\newcommand{\VariableTok}[1]{\textcolor[rgb]{0.10,0.09,0.49}{#1}}
\newcommand{\VerbatimStringTok}[1]{\textcolor[rgb]{0.25,0.44,0.63}{#1}}
\newcommand{\WarningTok}[1]{\textcolor[rgb]{0.38,0.63,0.69}{\textbf{\textit{#1}}}}
\usepackage{longtable,booktabs,array}
\newcounter{none} % for unnumbered tables
\usepackage{calc} % for calculating minipage widths
% Correct order of tables after \paragraph or \subparagraph
\usepackage{etoolbox}
\makeatletter
\patchcmd\longtable{\par}{\if@noskipsec\mbox{}\fi\par}{}{}
\makeatother
% Allow footnotes in longtable head/foot
\IfFileExists{footnotehyper.sty}{\usepackage{footnotehyper}}{\usepackage{footnote}}
\makesavenoteenv{longtable}
\usepackage{graphicx}
\makeatletter
\newsavebox\pandoc@box
\newcommand*\pandocbounded[1]{% scales image to fit in text height/width
  \sbox\pandoc@box{#1}%
  \Gscale@div\@tempa{\textheight}{\dimexpr\ht\pandoc@box+\dp\pandoc@box\relax}%
  \Gscale@div\@tempb{\linewidth}{\wd\pandoc@box}%
  \ifdim\@tempb\p@<\@tempa\p@\let\@tempa\@tempb\fi% select the smaller of both
  \ifdim\@tempa\p@<\p@\scalebox{\@tempa}{\usebox\pandoc@box}%
  \else\usebox{\pandoc@box}%
  \fi%
}
% Set default figure placement to htbp
\def\fps@figure{htbp}
\makeatother
\setlength{\emergencystretch}{3em} % prevent overfull lines
\providecommand{\tightlist}{%
  \setlength{\itemsep}{0pt}\setlength{\parskip}{0pt}}
\usepackage{bookmark}
\IfFileExists{xurl.sty}{\usepackage{xurl}}{} % add URL line breaks if available
\urlstyle{same}
\hypersetup{
  pdftitle={Hubbard-Holstein Mathematical Notes (Symbolic-Only Edition)},
  pdfauthor={Jake Skyler Strobel},
  hidelinks,
  pdfcreator={LaTeX via pandoc}}

\title{Hubbard-Holstein Mathematical Notes (Symbolic-Only Edition)}
\author{Jake Skyler Strobel}
\date{March 20, 2026}

\begin{document}
\maketitle

{
\setcounter{tocdepth}{3}
\tableofcontents
}
\section{Current Math Source Map}\label{current-math-source-map}

The former front-page \texttt{L=2\ Noisy\ Surface\ Run\ Handoff} has
been moved to Appendix F so this source opens with the active
Paper-I/Paper-II/Paper-III math rather than an archival run recipe. For
the local/offline L=2 HH noisy-surface recipe, see
\href{notes/l2_noisy_surface_run.md}{MATH/notes/l2\_noisy\_surface\_run.md}
and Appendix F below.

\begin{center}\rule{0.5\linewidth}{0.5pt}\end{center}

\section{Molecular-Vibronic H2 Benchmark: Physical Origin and
Parameters}\label{molecular-vibronic-h2-benchmark-physical-origin-and-parameters}

This section records what the molecular-vibronic H2 benchmark means
before it becomes a qubit Hamiltonian. It is a local harmonic model of
the H2 bond-stretch coordinate \(R\), with the electronic part computed
from molecular integrals and the nuclear motion reduced to one quantized
stretch mode. The reference geometry \(R_0=0.7414\,{\rm Angstrom}\) is
the center bond length for the Taylor expansion; it is the
near-equilibrium H2 bond length at which the molecular-orbital frame is
fixed. The finite-difference step is \(\Delta R=0.01\,{\rm Angstrom}\),
the electronic basis is STO-3G, and the three electronic snapshots are
evaluated at \(R_0-\Delta R\), \(R_0\), and \(R_0+\Delta R\). In the
repo generator these snapshots are RHF/STO-3G Psi4 calculations for
neutral singlet H2 with the molecular axis written as two H atoms
separated along \(z\), with no reorientation and no center-of-mass
recentering. Thus the answer to ``was this computed using the Psi4
interface?'' is yes for fixture generation: Psi4 supplies the SCF
wavefunctions, molecular orbital coefficients, one-electron integrals,
two-electron integrals, and nuclear repulsion at the three geometries.
For each geometry, the electronic molecular Hamiltonian is
\(H_{\rm el}(R)=E_{\rm nuc}(R)I+\sum_{pq}h_{pq}(R)c_p^\dagger c_q+\frac12\sum_{pqrs}g_{pqrs}(R)c_p^\dagger c_q^\dagger c_s c_r\),
where \(p,q,r,s\) run over the spin orbitals induced by the active
STO-3G molecular orbitals, \(h_{pq}(R)\) is the MO one-electron integral
matrix, and \(g_{pqrs}(R)\) is the spin-resolved two-electron Coulomb
tensor in the repo's second-quantized convention. The nuclear term is
the clamped-nuclei Coulomb energy
\(E_{\rm nuc}(R)=\sum_{A<B}Z_AZ_B/|\mathbf R_A-\mathbf R_B|\) in atomic
units, multiplying the electronic identity; for neutral H2,
\(Z_A=Z_B=1\) and \(E_{\rm nuc}(R)=1/R_B\), with
\(R_B=1.8897259886\,R[{\rm Angstrom}]\) the bond length in Bohr. Thus
\(H_{\rm el}\) means the ab initio electronic molecular operator, not an
electric field or a classical electric term; numerically it is a Psi4
integral package converted to a fermionic operator and then to a Pauli
polynomial by the repository's Jordan-Wigner map. Ordinary benchmark
runs read the cached fixture and do not call Psi4.

The electronic active space is the minimal H2/STO-3G closed-shell space:
two spatial orbitals, four spin orbitals, and two electrons, so the
electronic Hamiltonian occupies four fermionic qubits in the repo's
blocked Jordan-Wigner ordering before the oscillator register is
appended. The center snapshot gives \(H_{\rm el}(R_0)\). The off-center
snapshots first have to be compared in the same molecular-orbital gauge,
because the three independent SCF calculations at \(R_0-\Delta R\),
\(R_0\), and \(R_0+\Delta R\) do not automatically label their molecular
orbitals in a common coordinate system. A displaced orbital can acquire
a sign, rotate within a nearly degenerate occupied or virtual subspace,
or mix numerically with the center-frame partner that a chemist would
call ``the same'' orbital. A finite difference taken before this
alignment would therefore mix physical vibronic response with arbitrary
SCF orbital-frame motion.

The gauge alignment is an orthogonal change of molecular-orbital
coordinates before subtraction. Let \(\{\chi_\mu^{(0)}\}\) and
\(\{\chi_\nu^{(\pm)}\}\) denote the center and displaced AO bases, and
write the center and displaced molecular orbitals as
\(\phi_p^{(0)}=\sum_\mu C_{0,\mu p}\chi_\mu^{(0)}\) and
\(\phi_q^{(\pm)}=\sum_\nu C_{\pm,\nu q}\chi_\nu^{(\pm)}\). The AO
cross-overlap \(S_{\rm AO}^{0,\pm}\) has entries
\((S_{\rm AO}^{0,\pm})_{\mu\nu}=\langle\chi_\mu^{(0)},\chi_\nu^{(\pm)}\rangle\),
so \(S_{\rm MO}=C_0^\top S_{\rm AO}^{0,\pm}C_\pm\) has entries
\((S_{\rm MO})_{pq}=\langle\phi_p^{(0)},\phi_q^{(\pm)}\rangle\). The
generator computes \(S_{\rm MO}=U\Sigma V^\top\) and uses the orthogonal
Procrustes rotation \(Q=VU^\top\), the closest orthogonal map from the
displaced MO frame into the center MO frame in Frobenius norm. The
displaced one-electron integrals are then transformed as
\(h_\pm\mapsto Q^\top h_\pm Q\), and the displaced two-electron tensor
as \(g_{\pm,pqrs}\mapsto Q_{ap}Q_{bq}Q_{cr}Q_{ds}g_{\pm,abcd}\), with
repeated orbital indices summed. After this alignment,
\(H_{\rm el}(R_0+\Delta R)-H_{\rm el}(R_0-\Delta R)\) is a derivative of
one electronic operator in one common orbital frame, not a subtraction
of two unrelated MO coordinate systems.

The vibronic coupling is the first derivative of the electronic
Hamiltonian with respect to the physical bond length. The code converts
the step to Bohr, \(\Delta R_B=1.8897259886\,\Delta R\), then forms
\(\left.dH_{\rm el}/dR\right|_{R_0}\approx[H_{\rm el}(R_0+\Delta R)-H_{\rm el}(R_0-\Delta R)]/(2\Delta R_B)\).
The same three electronic Hamiltonians give the harmonic force constant
through exact electronic ground-state energies in the two-electron
sector, not through the RHF energies: with
\(E_-=E_{\rm exact}(R_0-\Delta R)\), \(E_0=E_{\rm exact}(R_0)\), and
\(E_+=E_{\rm exact}(R_0+\Delta R)\), the local curvature is
\(\kappa=(E_+-2E_0+E_-)/\Delta R_B^2\). The cached values are
\(E_-=-1.1372957462351274\), \(E_0=-1.1372701746571048\),
\(E_+=-1.137081149509897\), hence
\(\kappa=0.4577166774522122\,E_h/a_0^2\). With the H2 reduced nuclear
mass \(\mu=m_p/2=918.076336715\) in electron-mass atomic units, the
stretch frequency is
\(\omega=\sqrt{\kappa/\mu}=0.022328470326434775\,E_h\), and the
zero-point displacement length is
\(x_{\rm zpf}=\sqrt{1/(2\mu\omega)}=0.15617665597649316\,a_0\).

The resulting full one-mode Hamiltonian used by SNAKE is
\(H_{\rm H2vib}=H_{\rm el}(R_0)+\omega(b^\dagger b+\tfrac12 I)+g_{\rm ep}x_{\rm zpf}H'_R(R_0)(b+b^\dagger)\),
with \(H'_R(R_0)\equiv\left.dH_{\rm el}/dR\right|_{R_0}\). Here
\(H_{\rm el}(R_0)\) acts on the four electronic spin-orbital qubits,
\(b^\dagger b\) counts stretch quanta in the appended oscillator
register, \(b+b^\dagger\) is the dimensionless oscillator displacement,
and \(x_{\rm zpf}(b+b^\dagger)\) is the physical bond-length
displacement around \(R_0\). In the actual Pauli polynomial,
\(H_{\rm el}(R_0)\) and \(H'_R(R_0)\) are lifted by identity on the
oscillator register, while \(b^\dagger b\) and \(b+b^\dagger\) are
lifted by identity on the electronic register. The dimensionless knob
\(g_{\rm ep}\) scales the linear vibronic channel: \(g_{\rm ep}=1\)
means the finite-difference molecular derivative is used at its fixture
scale, while smaller values such as \(g_{\rm ep}=0.25\) define weaker
coupling lanes without recomputing the underlying electronic surface.
The checked-in fixture stores \(H_{\rm el}(R_0)\), \(H'_R(R_0)\), the
assembled \(H_{\rm H2vib}\), the restricted-HF times oscillator-vacuum
reference state, the dense physical-sector exact ground energy, and a
13-generator molecular-vibronic pool containing electronic UCCSD
channels, a pure oscillator momentum channel, and electron-oscillator
coupled momentum channels. Runtime reconstruction uses those cached
Pauli polynomials to rebuild the requested oscillator register and
coupling scale; regenerating the fixture itself requires Psi4.

\begin{center}\rule{0.5\linewidth}{0.5pt}\end{center}

\section{{[}QPU{]} L=2 Winner Framing and Real QPU Settings (April 11,
2026)}\label{qpu-l2-winner-framing-and-real-qpu-settings-april-11-2026}

\begin{quote}
Front-page contract: distinguish the best raw HF-start ADAPT result from
the separately validated real-QPU deployment anchor. Do not conflate
either one with later prune/readapt JSON post-processing or with
fixed-scaffold runtime submissions.
\end{quote}

\subsection{Two Front-Page Objects}\label{two-front-page-objects}

\begin{itemize}
\tightlist
\item
  \textbf{Scientific raw HF-start winner.} This means direct ADAPT from
  the Hartree-Fock reference with \texttt{adapt\_ref\_json\ =\ null},
  with no later JSON prune/recovery pass substituted in. Under the
  canonical \texttt{L=2,\ U=4.0,\ g=0.5} gate
  \texttt{\textbar{}ΔE\textbar{}\ \textbackslash{}sim\ 10\^{}\{-4\}},
  the best substantiated raw HF-start line currently in hand is the
  frozen legacy replay
  \texttt{20260409\_hh\_l2\_hist81\_legacy\_current\_compare\_d16\_v3/legacy\_20260322},
  which reproduces the March 22 \texttt{full\_meta} winner at \[
  |\Delta E|=5.6178234645\times 10^{-5}
  \] at \texttt{81} transpiled two-qubit gates on
  \texttt{FakeMarrakesh}.
\item
  \textbf{Validated real-QPU deployment anchor.} This is the separately
  validated April 5 \texttt{pareto\_lean\ +\ phase3\_v1\ +\ SPSA} raw
  HF-start command that is currently trusted for rerunability on the
  live public CLI surface. Its scaffold burden is heavier, but it
  remains the validated deployment recipe: \[
  |\Delta E|=1.0822209459\times 10^{-4},
  \] with \texttt{218} transpiled two-qubit gates on the same
  \texttt{FakeMarrakesh} proxy.
\end{itemize}

\subsection{Current checkout state (cost vs
energy)}\label{current-checkout-state-cost-vs-energy}

The current checkout now has a clean cost-vs-energy ranking at the
canonical \texttt{L=2,\ U=4.0,\ g=0.5} point:

\textbf{Default documentation comparison.} For HH ADAPT docs, the
headline physical-accuracy comparison is no longer the artifact's
same-cutoff exact target alone. The variational/scaffold run may still
use \texttt{n\_ph\_max=1}, but documentation should also compare its
final energy to an independently diagonalized higher-cutoff HH sector,
defaulting at the canonical L=2 point to \texttt{n\_ph\_max=3}. Keep the
two errors separate:

\begin{itemize}
\tightlist
\item
  same-cutoff ADAPT convergence:
  \(|E_{\mathrm{ADAPT}}^{(1)}-E_{\mathrm{exact}}^{(1)}|\);
\item
  external phonon-truncation benchmark:
  \(|E_{\mathrm{ADAPT}}^{(1)}-E_{\mathrm{exact}}^{(3)}|\).
\end{itemize}

At this point, \(E_{\mathrm{exact}}^{(1)}=0.158667904125724\) and
\(E_{\mathrm{exact}}^{(3)}=0.158439453796385\). Thus the documented
physical error of the best raw winner with
\(E_{\mathrm{ADAPT}}=0.1587240823603705\) is
\(2.8462856399\times10^{-4}\), while its same-cutoff convergence error
remains \(5.6178234645\times10^{-5}\).

\begin{itemize}
\tightlist
\item
  Powell legacy incumbents: best \texttt{75} 2Q at
  \texttt{\textbar{}\textbackslash{}Delta\ E\textbar{}=5.6178234644e-05};
  runner-up \texttt{81} 2Q at
  \texttt{\textbar{}\textbackslash{}Delta\ E\textbar{}=5.6178234645e-05}.
\item
  Powell current-route incumbents: best \texttt{98} 2Q at
  \texttt{\textbar{}\textbackslash{}Delta\ E\textbar{}=5.6178234645e-05};
  runner-up \texttt{118} 2Q at
  \texttt{\textbar{}\textbackslash{}Delta\ E\textbar{}=5.6178234645e-05}.
\item
  SPSA legacy incumbents: no in-band feasible candidate yet; search
  still running.
\item
  SPSA current-route incumbents: best so far \texttt{0} 2Q at
  \texttt{\textbar{}\textbackslash{}Delta\ E\textbar{}=5.6178234645e-05};
  runner-up pending.
\item
  The validated public/deployment anchor remains the April 5 SPSA route
  at \texttt{218} 2Q and
  \texttt{\textbar{}\textbackslash{}Delta\ E\textbar{}=1.0822209459e-04}.
\end{itemize}

\subsection{\texorpdfstring{Raw HF-Start Direct ADAPT Table
(\texttt{\textbar{}\textbackslash{}Delta\ E\textbar{}\textbackslash{}sim\ 10\^{}\{-4\}})}{Raw HF-Start Direct ADAPT Table (\textbar\textbackslash Delta E\textbar\textbackslash sim 10\^{}\{-4\})}}\label{raw-hf-start-direct-adapt-table-delta-esim-10-4}

{\def\LTcaptype{none} % do not increment counter
\begin{longtable}[]{@{}
  >{\raggedright\arraybackslash}p{(\linewidth - 14\tabcolsep) * \real{0.1034}}
  >{\raggedright\arraybackslash}p{(\linewidth - 14\tabcolsep) * \real{0.1034}}
  >{\raggedright\arraybackslash}p{(\linewidth - 14\tabcolsep) * \real{0.1034}}
  >{\raggedleft\arraybackslash}p{(\linewidth - 14\tabcolsep) * \real{0.1379}}
  >{\raggedleft\arraybackslash}p{(\linewidth - 14\tabcolsep) * \real{0.1379}}
  >{\raggedleft\arraybackslash}p{(\linewidth - 14\tabcolsep) * \real{0.1379}}
  >{\raggedleft\arraybackslash}p{(\linewidth - 14\tabcolsep) * \real{0.1379}}
  >{\raggedleft\arraybackslash}p{(\linewidth - 14\tabcolsep) * \real{0.1379}}@{}}
\toprule\noalign{}
\begin{minipage}[b]{\linewidth}\raggedright
Role
\end{minipage} & \begin{minipage}[b]{\linewidth}\raggedright
Artifact / surface
\end{minipage} & \begin{minipage}[b]{\linewidth}\raggedright
Pool
\end{minipage} & \begin{minipage}[b]{\linewidth}\raggedleft
\texttt{\textbar{}\textbackslash{}Delta\ E\textbar{}}
\end{minipage} & \begin{minipage}[b]{\linewidth}\raggedleft
Logical ops
\end{minipage} & \begin{minipage}[b]{\linewidth}\raggedleft
Runtime params
\end{minipage} & \begin{minipage}[b]{\linewidth}\raggedleft
Transpiled 2Q
\end{minipage} & \begin{minipage}[b]{\linewidth}\raggedleft
Depth
\end{minipage} \\
\midrule\noalign{}
\endhead
\bottomrule\noalign{}
\endlastfoot
Best raw HF-start ADAPT winner &
\texttt{hh\_backend\_adapt\_fullhorse\_powell\_20260322T194155Z\_fakenighthawk}
& \texttt{full\_meta} & \texttt{5.6178234645e-05} & 16 & 27 & 81 &
151 \\
Frozen legacy replay in current checkout &
\texttt{20260409\_hh\_l2\_hist81\_legacy\_current\_compare\_d16\_v3/legacy\_20260322}
& \texttt{full\_meta} frozen legacy replay & \texttt{5.6178234645e-05} &
16 & 27 & 81 & 151 \\
Best reduced-pool direct line &
\texttt{adapt\_hh\_L2\_ecut1\_pareto\_lean\_l2\_phase3\_powell\_rerun\_currentcode\_20260321T171615Z}
& \texttt{pareto\_lean\_l2} & \texttt{5.6178234644e-05} & 14 & 25 & 97 &
208 \\
Wide-shortlist direct comparison &
\texttt{campaign\_A7c\_L2\_shortlist\_wide\_phase3\_v1} &
\texttt{full\_meta} & \texttt{5.6178234644e-05} & 14 & 30 & 116 & 310 \\
Best current-route comparable replay &
\texttt{20260409\_hh\_l2\_current\_motif\_from\_legacy81\_v1} &
\texttt{full\_meta} current route + legacy motif &
\texttt{5.7102691670e-05} & 14 & 36 & 168 & 475 \\
Best current fullhorse-style replay &
\texttt{20260410\_hh\_l2\_current\_fullhorse\_recovery\_v1/fullhorse\_spliton\_norepeats\_motif}
& \texttt{full\_meta} current fullhorse-style route &
\texttt{5.7102730821e-05} & 16 & 41 & 193 & 592 \\
Current-code reproduction of March 22 low-energy regime &
\texttt{20260409\_hh\_l2\_hist81\_repro\_powell\_histbare\_d16\_v1} &
\texttt{full\_meta} historicalized POWELL surface &
\texttt{5.7217492583e-05} & 17 & 54 & 265 & 786 \\
Best current native recovery branch on today's diagnostic surface &
\texttt{20260409\_hh\_l2\_children\_repeat\_bridge\_diag\_v1}
(\texttt{children\_off\ +\ repeats\_off\ +\ hist}) & \texttt{full\_meta}
diagnostic bridge surface & \texttt{5.6178241861e-05} & 13 & 36 & 160 &
408 \\
\end{longtable}
}

Best Powell legacy-focused compatibility-search line \textbar{}
\texttt{artifacts/agent\_runs/20260414\_hh\_l2\_legacy\_focus\_optuna\_v1/legacy/eps\_6.200em05/trial\_0016}
\textbar{} legacy Powell Optuna lane (\texttt{runtime\_split} on, no
batching, transpile-single burden, no prune) \textbar{}
\texttt{5.6178234644e-05} \textbar{} 16 \textbar{} 29 \textbar{} 75
\textbar{} 178 \textbar{}\\
Runner-up Powell legacy-focused compatibility-search line \textbar{}
\texttt{artifacts/agent\_runs/20260414\_hh\_l2\_legacy\_focus\_optuna\_v1/legacy/eps\_6.200em05/trial\_0001}
\textbar{} legacy Powell Optuna lane (\texttt{runtime\_split} on,
batching on, transpile-single burden, prune on) \textbar{}
\texttt{5.6178234645e-05} \textbar{} 16 \textbar{} 27 \textbar{} 81
\textbar{} 151 \textbar{}\\
Best Powell current-code focused-bridge search line \textbar{}
\texttt{artifacts/agent\_runs/20260414\_hh\_l2\_bridge\_diag\_focus\_optuna\_v1/global/eps\_6.200em05/trial\_0024}
\textbar{} focused \texttt{bridge\_diag} Powell Optuna global lane
(\texttt{proxy\_reduced}, repeats off, shortlist split on, no prune)
\textbar{} \texttt{5.6178234645e-05} \textbar{} 13 \textbar{} 28
\textbar{} 98 \textbar{} 267 \textbar{}\\
Runner-up Powell current-code focused-bridge search line \textbar{}
\texttt{artifacts/agent\_runs/20260414\_hh\_l2\_bridge\_diag\_focus\_optuna\_v1/global/eps\_6.200em05/trial\_0005}
\textbar{} focused \texttt{bridge\_diag} Powell Optuna global lane
(\texttt{proxy\_reduced}, repeats off, shortlist split on, no prune)
\textbar{} \texttt{5.6178234645e-05} \textbar{} 13 \textbar{} 30
\textbar{} 118 \textbar{} 351 \textbar{}\\
Best SPSA legacy-focused compatibility-search line \textbar{}
\texttt{pending} \textbar{} live SPSA legacy search in progress
\textbar{} \texttt{-\/-} \textbar{} -- \textbar{} -- \textbar{} --
\textbar{} -- \textbar{}\\
Runner-up SPSA legacy-focused compatibility-search line \textbar{}
\texttt{pending} \textbar{} live SPSA legacy search in progress
\textbar{} \texttt{-\/-} \textbar{} -- \textbar{} -- \textbar{} --
\textbar{} -- \textbar{}\\
Best SPSA current-code focused search line \textbar{}
\texttt{artifacts/agent\_runs/20260415\_hh\_l2\_powell\_manifold\_spsa\_refine\_v2/current/trial\_0000}
\textbar{} focused \texttt{current\_98} SPSA Optuna global lane
(\texttt{powell\_handoff\_phase3\_v1}, repeats on/base,
split=phase3\_v1\_replay, prune=inherited) \textbar{}
\texttt{5.6178234645e-05} \textbar{} 13 \textbar{} 56 \textbar{} 0
\textbar{} 0 \textbar{}\\
Runner-up SPSA current-code focused search line \textbar{}
\texttt{pending} \textbar{} live SPSA current-route search in progress
\textbar{} \texttt{-\/-} \textbar{} -- \textbar{} -- \textbar{} --
\textbar{} -- \textbar{}

Validated April 5 deployment anchor \textbar{}
\texttt{pareto\_lean\ +\ phase3\_v1\ +\ SPSA} \textbar{}
\texttt{pareto\_lean} \textbar{} \texttt{1.0822209459e-04} \textbar{} 20
\textbar{} 52 \textbar{} 218 \textbar{} 633 \textbar{}

The final \texttt{218}-2Q row is the same raw HF-start scaffold family
as the confirmed April 5 authority command; it is heavier than the
scientific raw winner, but it is the currently validated deployment
anchor on the public CLI surface.

The repo now tracks Powell and SPSA separately: Powell incumbents are
\texttt{75} 2Q (legacy) and \texttt{98} 2Q (current), while SPSA
incumbents are \texttt{pending} 2Q (legacy) and \texttt{0} 2Q (current).

The post-processed \texttt{31/37/57}-2Q prune/readapt artifacts are
\textbf{not} the object a fresh real-QPU ADAPT run from HF would
execute. They remain useful comparison context only.

\subsubsection{Current native recovery note (April 9,
2026)}\label{current-native-recovery-note-april-9-2026}

The new \texttt{children\_off\ +\ repeats\_off\ +\ hist} branch from
\texttt{20260409\_hh\_l2\_children\_repeat\_bridge\_diag\_v1} deserves
front-page visibility because it is the best \textbf{current-code native
recovery branch} on today's diagnostic/public surface:

\[
|\Delta E|=5.6178241861\times 10^{-5},\qquad
F=0.9999719078,
\]

with \texttt{13} logical operators, \texttt{36} runtime parameters, and
\texttt{160} transpiled two-qubit gates at depth \texttt{408} on
\texttt{FakeMarrakesh}.

This is \textbf{not} the absolute best scaffold now available; the April
14 legacy-compat Optuna line at \texttt{75} 2Q and the March 22 /
fresh-replayed \texttt{81}-2Q legacy lines are all lighter overall. What
makes the April 9 branch important is different: it is the first current
diagnostic line that re-enters the good native basin while
simultaneously showing that disabling children and literal repeats helps
native recovery.

An April 14, 2026 search update strengthens that conclusion rather than
replacing it: the focused bridge-only Optuna follow-up now has two fresh
current-code in-band winners by transpiled two-qubit cost.

The current best fresh current-code bridge line is

\[
|\Delta E|=5.6178234645\times 10^{-5},\qquad \text{2Q}=98,
\]

from
\texttt{artifacts/agent\_runs/20260414\_hh\_l2\_bridge\_diag\_focus\_optuna\_v1/global/eps\_6.200em05/trial\_0024/},
with selected prefix

\begin{itemize}
\tightlist
\item
  \texttt{uccsd\_ferm\_lifted::uccsd\_sing(alpha:0-\textgreater{}1)}
\item
  \texttt{uccsd\_ferm\_lifted::uccsd\_sing(beta:2-\textgreater{}3)}
\item
  \texttt{paop\_lf\_full:paop\_dbl\_p(site=0-\textgreater{}phonon=0)}
\item
  \texttt{paop\_lf\_full:paop\_dbl\_p(site=1-\textgreater{}phonon=1)}
\item
  \texttt{paop\_full:paop\_hopdrag(0,1)::child\_set{[}1,3{]}}
\item
  \texttt{paop\_lf\_full:paop\_dbl\_p(site=0-\textgreater{}phonon=1)}
\item
  \texttt{paop\_full:paop\_disp(site=1)::child\_set{[}1{]}}
\item
  \texttt{paop\_full:paop\_disp(site=0)::child\_set{[}1{]}}
\end{itemize}

The current fresh runner-up is

\[
|\Delta E|=5.6178234645\times 10^{-5},\qquad \text{2Q}=118,
\]

from
\texttt{artifacts/agent\_runs/20260414\_hh\_l2\_bridge\_diag\_focus\_optuna\_v1/global/eps\_6.200em05/trial\_0005/},
with selected prefix

\begin{itemize}
\tightlist
\item
  \texttt{uccsd\_ferm\_lifted::uccsd\_sing(alpha:0-\textgreater{}1)}
\item
  \texttt{uccsd\_ferm\_lifted::uccsd\_sing(beta:2-\textgreater{}3)}
\item
  \texttt{paop\_lf\_full:paop\_dbl\_p(site=0-\textgreater{}phonon=0)}
\item
  \texttt{paop\_full:paop\_hopdrag(0,1)::child\_set{[}0,2{]}}
\item
  \texttt{paop\_full:paop\_disp(site=0)}
\item
  \texttt{paop\_full:paop\_disp(site=1)}
\end{itemize}

In parallel, the compatibility-only legacy lane is currently tracking
the same top two in-band legacy candidates by two-qubit cost as
\texttt{81} 2Q (the reproduced oracle) and \texttt{114} 2Q (the present
runner-up fresh legacy trial).

It still does \textbf{not} recover the historical lean family exactly.
Relative to the historical targets, it reaches \texttt{9/10} bridge
overlap and \texttt{5/6} lean overlap, with the remaining extras
concentrated in

\begin{itemize}
\tightlist
\item
  \texttt{paop\_full:paop\_disp(site=1)}
\item
  \texttt{paop\_lf\_full:paop\_dbl\_p(site=0-\textgreater{}phonon=1)}
\item
  \texttt{paop\_lf\_full:paop\_dbl\_p(site=1-\textgreater{}phonon=0)}
\item
  \texttt{paop\_full:paop\_hopdrag(0,1)}
\end{itemize}

This diagnostic result matters because it isolates the current native
failure mode more sharply than the older heavy broad-pool lines: the
good basin is real, but some quadrature and inside-family extras are
still acting as \textbf{path operators} rather than merely removable
terminal clutter.

\subsubsection{Current-code reproduction of the March 22 low-energy
regime}\label{current-code-reproduction-of-the-march-22-low-energy-regime}

The March 22 \texttt{81}-2Q \texttt{full\_meta} line is still a real
historical raw HF-start winner, but current code does \textbf{not}
replay that exact burden from its historical settings block. What
current code \textbf{does} reproduce is the same low-energy regime on a
historicalized POWELL surface with a hard cap at the first in-run depth
that re-enters the target band.

The saved current-code replay is

\[
|\Delta E|=5.7217492583\times 10^{-5},
\]

from
\texttt{20260409\_hh\_l2\_hist81\_repro\_powell\_histbare\_d16\_v1},
with \texttt{17} logical blocks, \texttt{54} runtime parameters, and
\texttt{265} transpiled two-qubit gates at depth \texttt{786} on
\texttt{FakeMarrakesh}.

So this run should be interpreted as a \textbf{current-code low-energy
reproduction}, not as a faithful current reproduction of the historical
\texttt{81}-2Q scaffold burden.

\textbf{Exact current-code replay CLI}

\begin{Shaded}
\begin{Highlighting}[]
\ExtensionTok{python}\NormalTok{ pipelines/hardcoded/adapt\_pipeline.py }\DataTypeTok{\textbackslash{}}
  \AttributeTok{{-}{-}L}\NormalTok{ 2 }\AttributeTok{{-}{-}problem}\NormalTok{ hh }\AttributeTok{{-}{-}t}\NormalTok{ 1.0 }\AttributeTok{{-}{-}u}\NormalTok{ 4.0 }\AttributeTok{{-}{-}dv}\NormalTok{ 0.0 }\AttributeTok{{-}{-}omega0}\NormalTok{ 1.0 }\AttributeTok{{-}{-}g{-}ep}\NormalTok{ 0.5 }\AttributeTok{{-}{-}n{-}ph{-}max}\NormalTok{ 1 }\DataTypeTok{\textbackslash{}}
  \AttributeTok{{-}{-}boson{-}encoding}\NormalTok{ binary }\AttributeTok{{-}{-}ordering}\NormalTok{ blocked }\AttributeTok{{-}{-}boundary}\NormalTok{ open }\AttributeTok{{-}{-}term{-}order}\NormalTok{ sorted }\DataTypeTok{\textbackslash{}}
  \AttributeTok{{-}{-}adapt{-}pool}\NormalTok{ full\_meta }\AttributeTok{{-}{-}adapt{-}continuation{-}mode}\NormalTok{ phase3\_v1 }\DataTypeTok{\textbackslash{}}
  \AttributeTok{{-}{-}adapt{-}max{-}depth}\NormalTok{ 16 }\AttributeTok{{-}{-}adapt{-}eps{-}grad}\NormalTok{ 5e{-}7 }\AttributeTok{{-}{-}adapt{-}eps{-}energy}\NormalTok{ 1e{-}9 }\DataTypeTok{\textbackslash{}}
  \AttributeTok{{-}{-}adapt{-}inner{-}optimizer}\NormalTok{ POWELL }\AttributeTok{{-}{-}adapt{-}state{-}backend}\NormalTok{ compiled }\DataTypeTok{\textbackslash{}}
  \AttributeTok{{-}{-}adapt{-}reopt{-}policy}\NormalTok{ windowed }\AttributeTok{{-}{-}adapt{-}window{-}size}\NormalTok{ 999999 }\AttributeTok{{-}{-}adapt{-}window{-}topk}\NormalTok{ 999999 }\DataTypeTok{\textbackslash{}}
  \AttributeTok{{-}{-}adapt{-}full{-}refit{-}every}\NormalTok{ 8 }\AttributeTok{{-}{-}adapt{-}final{-}full{-}refit}\NormalTok{ true }\DataTypeTok{\textbackslash{}}
  \AttributeTok{{-}{-}adapt{-}beam{-}live{-}branches}\NormalTok{ 1 }\AttributeTok{{-}{-}adapt{-}beam{-}children{-}per{-}parent}\NormalTok{ 1 }\AttributeTok{{-}{-}adapt{-}beam{-}terminated{-}keep}\NormalTok{ 1 }\DataTypeTok{\textbackslash{}}
  \AttributeTok{{-}{-}phase1{-}lambda{-}F}\NormalTok{ 1.0 }\AttributeTok{{-}{-}phase1{-}lambda{-}compile}\NormalTok{ 0.05 }\AttributeTok{{-}{-}phase1{-}lambda{-}measure}\NormalTok{ 0.02 }\AttributeTok{{-}{-}phase1{-}lambda{-}leak}\NormalTok{ 0.0 }\DataTypeTok{\textbackslash{}}
  \AttributeTok{{-}{-}phase1{-}score{-}z{-}alpha}\NormalTok{ 0.0 }\AttributeTok{{-}{-}phase1{-}shortlist{-}size}\NormalTok{ 256 }\AttributeTok{{-}{-}phase1{-}probe{-}max{-}positions}\NormalTok{ 999999 }\DataTypeTok{\textbackslash{}}
  \AttributeTok{{-}{-}phase1{-}plateau{-}patience}\NormalTok{ 2 }\AttributeTok{{-}{-}phase1{-}trough{-}margin{-}ratio}\NormalTok{ 1.0 }\DataTypeTok{\textbackslash{}}
  \AttributeTok{{-}{-}phase1{-}prune{-}enabled} \AttributeTok{{-}{-}phase1{-}prune{-}mode}\NormalTok{ live }\AttributeTok{{-}{-}phase1{-}prune{-}fraction}\NormalTok{ 0.25 }\AttributeTok{{-}{-}phase1{-}prune{-}max{-}candidates}\NormalTok{ 6 }\AttributeTok{{-}{-}phase1{-}prune{-}max{-}regression}\NormalTok{ 1e{-}8 }\DataTypeTok{\textbackslash{}}
  \AttributeTok{{-}{-}phase2{-}shortlist{-}fraction}\NormalTok{ 1.0 }\AttributeTok{{-}{-}phase2{-}shortlist{-}size}\NormalTok{ 128 }\DataTypeTok{\textbackslash{}}
  \AttributeTok{{-}{-}phase2{-}lambda{-}H}\NormalTok{ 1e{-}6 }\AttributeTok{{-}{-}phase2{-}rho}\NormalTok{ 0.25 }\AttributeTok{{-}{-}phase2{-}gamma{-}N}\NormalTok{ 1.0 }\DataTypeTok{\textbackslash{}}
  \AttributeTok{{-}{-}phase2{-}enable{-}batching} \AttributeTok{{-}{-}phase2{-}batch{-}target{-}size}\NormalTok{ 8 }\AttributeTok{{-}{-}phase2{-}batch{-}size{-}cap}\NormalTok{ 16 }\AttributeTok{{-}{-}phase2{-}batch{-}near{-}degenerate{-}ratio}\NormalTok{ 0.98 }\DataTypeTok{\textbackslash{}}
  \AttributeTok{{-}{-}phase2{-}w{-}depth}\NormalTok{ 0.0 }\AttributeTok{{-}{-}phase2{-}w{-}group}\NormalTok{ 0.0 }\AttributeTok{{-}{-}phase2{-}w{-}shot}\NormalTok{ 0.0 }\AttributeTok{{-}{-}phase2{-}w{-}optdim}\NormalTok{ 0.0 }\AttributeTok{{-}{-}phase2{-}w{-}reuse}\NormalTok{ 0.0 }\AttributeTok{{-}{-}phase2{-}w{-}lifetime}\NormalTok{ 0.0 }\DataTypeTok{\textbackslash{}}
  \AttributeTok{{-}{-}phase2{-}eta{-}L}\NormalTok{ 0.0 }\AttributeTok{{-}{-}phase2{-}motif{-}bonus{-}weight}\NormalTok{ 0.0 }\AttributeTok{{-}{-}phase2{-}duplicate{-}penalty{-}weight}\NormalTok{ 0.0 }\DataTypeTok{\textbackslash{}}
  \AttributeTok{{-}{-}phase2{-}compat{-}overlap{-}weight}\NormalTok{ 0.0 }\AttributeTok{{-}{-}phase2{-}compat{-}comm{-}weight}\NormalTok{ 0.0 }\AttributeTok{{-}{-}phase2{-}compat{-}curv{-}weight}\NormalTok{ 0.0 }\AttributeTok{{-}{-}phase2{-}compat{-}sched{-}weight}\NormalTok{ 0.0 }\AttributeTok{{-}{-}phase2{-}compat{-}measure{-}weight}\NormalTok{ 0.0 }\DataTypeTok{\textbackslash{}}
  \AttributeTok{{-}{-}phase2{-}frontier{-}ratio}\NormalTok{ 1.0 }\AttributeTok{{-}{-}phase3{-}frontier{-}ratio}\NormalTok{ 1.0 }\AttributeTok{{-}{-}phase2{-}remaining{-}evaluations{-}proxy{-}mode}\NormalTok{ none }\DataTypeTok{\textbackslash{}}
  \AttributeTok{{-}{-}phase3{-}tie{-}beam{-}max{-}branches}\NormalTok{ 1 }\DataTypeTok{\textbackslash{}}
  \AttributeTok{{-}{-}phase3{-}symmetry{-}mitigation{-}mode}\NormalTok{ verify\_only }\AttributeTok{{-}{-}phase3{-}enable{-}rescue} \DataTypeTok{\textbackslash{}}
  \AttributeTok{{-}{-}phase3{-}lifetime{-}cost{-}mode}\NormalTok{ phase3\_v1 }\AttributeTok{{-}{-}phase3{-}runtime{-}split{-}mode}\NormalTok{ shortlist\_pauli\_children\_v1 }\DataTypeTok{\textbackslash{}}
  \AttributeTok{{-}{-}phase3{-}backend{-}cost{-}mode}\NormalTok{ transpile\_single\_v1 }\AttributeTok{{-}{-}phase3{-}backend{-}name}\NormalTok{ FakeNighthawk }\DataTypeTok{\textbackslash{}}
  \AttributeTok{{-}{-}phase3{-}backend{-}transpile{-}seed}\NormalTok{ 7 }\AttributeTok{{-}{-}phase3{-}backend{-}optimization{-}level}\NormalTok{ 1 }\DataTypeTok{\textbackslash{}}
  \AttributeTok{{-}{-}adapt{-}seed}\NormalTok{ 7 }\AttributeTok{{-}{-}adapt{-}finite{-}angle{-}fallback} \AttributeTok{{-}{-}adapt{-}finite{-}angle}\NormalTok{ 0.1 }\AttributeTok{{-}{-}adapt{-}finite{-}angle{-}min{-}improvement}\NormalTok{ 1e{-}12 }\DataTypeTok{\textbackslash{}}
  \AttributeTok{{-}{-}adapt{-}drop{-}floor}\NormalTok{ 1e{-}11 }\AttributeTok{{-}{-}adapt{-}drop{-}patience}\NormalTok{ 5 }\AttributeTok{{-}{-}adapt{-}drop{-}min{-}depth}\NormalTok{ 12 }\AttributeTok{{-}{-}adapt{-}grad{-}floor} \AttributeTok{{-}1.0} \DataTypeTok{\textbackslash{}}
  \AttributeTok{{-}{-}paop{-}r}\NormalTok{ 1 }\AttributeTok{{-}{-}paop{-}normalization}\NormalTok{ none }\DataTypeTok{\textbackslash{}}
  \AttributeTok{{-}{-}initial{-}state{-}source}\NormalTok{ adapt\_vqe }\DataTypeTok{\textbackslash{}}
  \AttributeTok{{-}{-}skip{-}pdf}
\end{Highlighting}
\end{Shaded}

\subsection{Real QPU Deployment Anchor (April 5, 2026
Validated)}\label{real-qpu-deployment-anchor-april-5-2026-validated}

\begin{quote}
\textbf{This is the recommended configuration for real QPU deployment.}
Validated via narrow 36-case local sweep confirming rerunnable
deployment optimality in the \texttt{pareto\_lean} + \texttt{phase3\_v1}
regime. This is the deployment anchor, not the absolute lowest raw
HF-start two-qubit line.
\end{quote}

\subsection{Anchor Configuration}\label{anchor-configuration}

\textbf{Use this exact command for real QPU L=2 ADAPT runs:}

\begin{Shaded}
\begin{Highlighting}[]
\ExtensionTok{python} \AttributeTok{{-}u} \AttributeTok{{-}m}\NormalTok{ pipelines.hardcoded.adapt\_pipeline }\DataTypeTok{\textbackslash{}}
  \AttributeTok{{-}{-}L}\NormalTok{ 2 }\AttributeTok{{-}{-}problem}\NormalTok{ hh }\AttributeTok{{-}{-}t}\NormalTok{ 1.0 }\AttributeTok{{-}{-}u}\NormalTok{ 4.0 }\AttributeTok{{-}{-}dv}\NormalTok{ 0.0 }\DataTypeTok{\textbackslash{}}
  \AttributeTok{{-}{-}omega0}\NormalTok{ 1.0 }\AttributeTok{{-}{-}g{-}ep}\NormalTok{ 0.5 }\AttributeTok{{-}{-}n{-}ph{-}max}\NormalTok{ 1 }\DataTypeTok{\textbackslash{}}
  \AttributeTok{{-}{-}boson{-}encoding}\NormalTok{ binary }\AttributeTok{{-}{-}ordering}\NormalTok{ blocked }\AttributeTok{{-}{-}boundary}\NormalTok{ open }\AttributeTok{{-}{-}term{-}order}\NormalTok{ sorted }\DataTypeTok{\textbackslash{}}
  \AttributeTok{{-}{-}adapt{-}pool}\NormalTok{ pareto\_lean }\AttributeTok{{-}{-}adapt{-}continuation{-}mode}\NormalTok{ phase3\_v1 }\DataTypeTok{\textbackslash{}}
  \AttributeTok{{-}{-}adapt{-}max{-}depth}\NormalTok{ 160 }\AttributeTok{{-}{-}adapt{-}eps{-}grad}\NormalTok{ 5e{-}7 }\AttributeTok{{-}{-}adapt{-}eps{-}energy}\NormalTok{ 1e{-}9 }\AttributeTok{{-}{-}adapt{-}seed}\NormalTok{ 5 }\DataTypeTok{\textbackslash{}}
  \AttributeTok{{-}{-}adapt{-}inner{-}optimizer}\NormalTok{ SPSA }\DataTypeTok{\textbackslash{}}
  \AttributeTok{{-}{-}adapt{-}spsa{-}a}\NormalTok{ 0.1 }\AttributeTok{{-}{-}adapt{-}spsa{-}c}\NormalTok{ 0.02 }\AttributeTok{{-}{-}adapt{-}spsa{-}A}\NormalTok{ 5.0 }\DataTypeTok{\textbackslash{}}
  \AttributeTok{{-}{-}adapt{-}spsa{-}callback{-}every}\NormalTok{ 5 }\AttributeTok{{-}{-}adapt{-}spsa{-}progress{-}every{-}s}\NormalTok{ 30 }\DataTypeTok{\textbackslash{}}
  \AttributeTok{{-}{-}adapt{-}maxiter}\NormalTok{ 3200 }\DataTypeTok{\textbackslash{}}
  \AttributeTok{{-}{-}adapt{-}state{-}backend}\NormalTok{ compiled }\DataTypeTok{\textbackslash{}}
  \AttributeTok{{-}{-}adapt{-}reopt{-}policy}\NormalTok{ windowed }\AttributeTok{{-}{-}adapt{-}window{-}size}\NormalTok{ 999999 }\AttributeTok{{-}{-}adapt{-}window{-}topk}\NormalTok{ 999999 }\DataTypeTok{\textbackslash{}}
  \AttributeTok{{-}{-}adapt{-}full{-}refit{-}every}\NormalTok{ 8 }\AttributeTok{{-}{-}adapt{-}final{-}full{-}refit}\NormalTok{ true }\DataTypeTok{\textbackslash{}}
  \AttributeTok{{-}{-}adapt{-}beam{-}live{-}branches}\NormalTok{ 2 }\AttributeTok{{-}{-}adapt{-}beam{-}children{-}per{-}parent}\NormalTok{ 2 }\AttributeTok{{-}{-}adapt{-}beam{-}terminated{-}keep}\NormalTok{ 2 }\DataTypeTok{\textbackslash{}}
  \AttributeTok{{-}{-}phase1{-}prune{-}enabled} \AttributeTok{{-}{-}phase1{-}prune{-}fraction}\NormalTok{ 0.25 }\AttributeTok{{-}{-}phase1{-}prune{-}max{-}candidates}\NormalTok{ 6 }\AttributeTok{{-}{-}phase1{-}prune{-}max{-}regression}\NormalTok{ 1e{-}8 }\DataTypeTok{\textbackslash{}}
  \AttributeTok{{-}{-}phase1{-}probe{-}max{-}positions}\NormalTok{ 999999 }\AttributeTok{{-}{-}phase1{-}trough{-}margin{-}ratio}\NormalTok{ 1.0 }\AttributeTok{{-}{-}phase1{-}shortlist{-}size}\NormalTok{ 64 }\DataTypeTok{\textbackslash{}}
  \AttributeTok{{-}{-}phase2{-}shortlist{-}fraction}\NormalTok{ 1.0 }\AttributeTok{{-}{-}phase2{-}shortlist{-}size}\NormalTok{ 64 }\DataTypeTok{\textbackslash{}}
  \AttributeTok{{-}{-}phase2{-}frontier{-}ratio}\NormalTok{ 0.85 }\DataTypeTok{\textbackslash{}}
  \AttributeTok{{-}{-}phase2{-}lambda{-}H}\NormalTok{ 1e{-}06 }\AttributeTok{{-}{-}phase2{-}rho}\NormalTok{ 0.25 }\AttributeTok{{-}{-}phase2{-}gamma{-}N}\NormalTok{ 1.0 }\DataTypeTok{\textbackslash{}}
  \AttributeTok{{-}{-}phase2{-}w{-}depth}\NormalTok{ 0.1 }\AttributeTok{{-}{-}phase2{-}w{-}group}\NormalTok{ 0.075 }\AttributeTok{{-}{-}phase2{-}w{-}shot}\NormalTok{ 0.075 }\AttributeTok{{-}{-}phase2{-}w{-}optdim}\NormalTok{ 0.05 }\AttributeTok{{-}{-}phase2{-}w{-}reuse}\NormalTok{ 0.05 }\AttributeTok{{-}{-}phase2{-}w{-}lifetime}\NormalTok{ 0.05 }\DataTypeTok{\textbackslash{}}
  \AttributeTok{{-}{-}phase2{-}eta{-}L}\NormalTok{ 0.0 }\AttributeTok{{-}{-}phase2{-}motif{-}bonus{-}weight}\NormalTok{ 0.05 }\AttributeTok{{-}{-}phase2{-}duplicate{-}penalty{-}weight}\NormalTok{ 0.0 }\DataTypeTok{\textbackslash{}}
  \AttributeTok{{-}{-}phase2{-}compat{-}overlap{-}weight}\NormalTok{ 0.4 }\AttributeTok{{-}{-}phase2{-}compat{-}comm{-}weight}\NormalTok{ 0.2 }\AttributeTok{{-}{-}phase2{-}compat{-}curv{-}weight}\NormalTok{ 0.2 }\DataTypeTok{\textbackslash{}}
  \AttributeTok{{-}{-}phase2{-}compat{-}sched{-}weight}\NormalTok{ 0.2 }\AttributeTok{{-}{-}phase2{-}compat{-}measure{-}weight}\NormalTok{ 0.2 }\DataTypeTok{\textbackslash{}}
  \AttributeTok{{-}{-}phase2{-}remaining{-}evaluations{-}proxy{-}mode}\NormalTok{ auto }\DataTypeTok{\textbackslash{}}
  \AttributeTok{{-}{-}phase3{-}frontier{-}ratio}\NormalTok{ 0.85 }\DataTypeTok{\textbackslash{}}
  \AttributeTok{{-}{-}phase3{-}symmetry{-}mitigation{-}mode}\NormalTok{ verify\_only }\AttributeTok{{-}{-}phase3{-}enable{-}rescue} \DataTypeTok{\textbackslash{}}
  \AttributeTok{{-}{-}phase3{-}backend{-}cost{-}mode}\NormalTok{ proxy }\AttributeTok{{-}{-}phase3{-}runtime{-}split{-}mode}\NormalTok{ off }\AttributeTok{{-}{-}phase3{-}lifetime{-}cost{-}mode}\NormalTok{ off }\DataTypeTok{\textbackslash{}}
  \AttributeTok{{-}{-}adapt{-}drop{-}floor}\NormalTok{ 0.0005 }\AttributeTok{{-}{-}adapt{-}drop{-}patience}\NormalTok{ 3 }\AttributeTok{{-}{-}adapt{-}drop{-}min{-}depth}\NormalTok{ 12 }\DataTypeTok{\textbackslash{}}
  \AttributeTok{{-}{-}initial{-}state{-}source}\NormalTok{ adapt\_vqe }\AttributeTok{{-}{-}skip{-}pdf} \AttributeTok{{-}{-}phase2{-}no{-}batching}
\end{Highlighting}
\end{Shaded}

\subsection{Expected Performance}\label{expected-performance}

{\def\LTcaptype{none} % do not increment counter
\begin{longtable}[]{@{}
  >{\raggedright\arraybackslash}p{(\linewidth - 4\tabcolsep) * \real{0.3636}}
  >{\raggedright\arraybackslash}p{(\linewidth - 4\tabcolsep) * \real{0.3182}}
  >{\raggedright\arraybackslash}p{(\linewidth - 4\tabcolsep) * \real{0.3182}}@{}}
\toprule\noalign{}
\begin{minipage}[b]{\linewidth}\raggedright
Metric
\end{minipage} & \begin{minipage}[b]{\linewidth}\raggedright
Value
\end{minipage} & \begin{minipage}[b]{\linewidth}\raggedright
Notes
\end{minipage} \\
\midrule\noalign{}
\endhead
\bottomrule\noalign{}
\endlastfoot
Energy & 0.15877612622031317 & Exact ground state reference \\
\textbar ΔE\textbar{} & \textbf{1.082e-4} & Excellent accuracy for
NISQ \\
Logical scaffold size & 20 & Validated raw HF-start deployment
scaffold \\
Runtime parameters & 52 & Efficient parameterization \\
Marrakesh proxy burden & \texttt{218} 2Q, depth \texttt{633}, size
\texttt{939} & Same raw HF-start scaffold as the confirmed strong
template \\
Wall-clock & \textasciitilde43s & Fast convergence \\
\end{longtable}
}

\subsection{Validation Results (April 5
Sweep)}\label{validation-results-april-5-sweep}

A 36-case narrow local sweep confirmed this configuration is optimal:

\begin{itemize}
\tightlist
\item
  {[}OK{]} \textbf{SPSA maxiter=3200}: Goldilocks zone (2400
  undershoots, 4000 overshoots)
\item
  {[}OK{]} \textbf{Frontier ratio 0.85}: All values (0.80--0.90)
  converge to same energy; 0.85 balances quality
\item
  {[}OK{]} \textbf{Beam (2,2,2)}: Sufficient; wider beam (3,2,3) costs
  50\% more time, zero quality gain
\item
  {[}OK{]} \textbf{Drop policy (0.0005, 3, 12)}: Optimal shallowness;
  tighter drop \(\to\) shallower circuits for QPU
\end{itemize}

\subsubsection{Critical Finding}\label{critical-finding}

Increasing SPSA \texttt{maxiter} beyond 3200 \textbf{degrades}
performance (1.27--1.54e-4 vs 1.082e-4): - Suggests pool saturation:
\texttt{pareto\_lean} operators are exhausted at this budget - SPSA
noise amplification: longer optimization in finite ansatz space doesn't
help

\textbf{Do not deviate from maxiter=3200 in this regime.}

\subsection{What NOT to Change}\label{what-not-to-change}

Based on sweep validation: - {[}X{]} Pool: Do not switch from
\texttt{pareto\_lean} - {[}X{]} Continuation: Do not move away from
\texttt{phase3\_v1} - {[}X{]} Runtime split: Keep \texttt{off} for L=2 -
{[}X{]} SPSA budget: Do not increase \texttt{maxiter} beyond 3200 -
{[}X{]} Frontier: Do not vary frontier ratio (insensitive in 0.80--0.90
range) - {[}X{]} Beam: Do not use (3,2,3); keep (2,2,2)

\subsection{For Real QPU Deployment}\label{for-real-qpu-deployment}

\begin{enumerate}
\def\labelenumi{\arabic{enumi}.}
\tightlist
\item
  Use the command above exactly as written
\item
  Expect \textbar ΔE\textbar{} \(\approx\) 1.08e-4 (this is the
  best-known trustworthy result)
\item
  Ansatz depth = 20 is excellent for NISQ mitigation
\item
  No further local tuning will improve this configuration (sweep
  confirms)
\end{enumerate}

\subsection{Archival Driven Calibration
(non-default)}\label{archival-driven-calibration-non-default}

The historical heavier-target driven realtime-controller calibration for
the same L=2 HH scaffold family came from an archived April 2026
drive-heavy controller scan artifact.

This block is retained only as calibration evidence for the stronger
\texttt{A=0.55}, \texttt{tbar=4.0}, \texttt{t\_final=8.0} benchmark. It
is \textbf{not} the current route authority, must not set code defaults,
and must not be used for agent route selection. The canonical realtime
law is Chapter 17A AP-McLachlan append--prune support-patch control
below.

\subsubsection{Historical driven-controller headline
metrics}\label{historical-driven-controller-headline-metrics}

{\def\LTcaptype{none} % do not increment counter
\begin{longtable}[]{@{}
  >{\raggedright\arraybackslash}p{(\linewidth - 4\tabcolsep) * \real{0.3636}}
  >{\raggedright\arraybackslash}p{(\linewidth - 4\tabcolsep) * \real{0.3182}}
  >{\raggedright\arraybackslash}p{(\linewidth - 4\tabcolsep) * \real{0.3182}}@{}}
\toprule\noalign{}
\begin{minipage}[b]{\linewidth}\raggedright
Metric
\end{minipage} & \begin{minipage}[b]{\linewidth}\raggedright
Value
\end{minipage} & \begin{minipage}[b]{\linewidth}\raggedright
Notes
\end{minipage} \\
\midrule\noalign{}
\endhead
\bottomrule\noalign{}
\endlastfoot
Historical controller law & archived signed-energy-lead-tapered secant
blend & Calibration evidence only; non-default \\
Drive target & \texttt{A=0.55}, \texttt{omega=2.0}, \texttt{tbar=4.0},
\texttt{t0=0.0} & Heavier than the earlier fine-resolution regime \\
Mean \texttt{\textbar{}ΔE\textbar{}} & \textbf{1.633e-3} & Historical
whole-window value on that benchmark \\
Max \texttt{\textbar{}ΔE\textbar{}} & \textbf{6.501e-3} & Historical
improvement over prior signed-blend and raw secant anchors \\
Early-window MAE (\texttt{t≤3.35}) & \textbf{1.176e-3} & Historical
startup-dip recovery evidence \\
Late-window MAE (\texttt{t≥3.35}) & \textbf{1.964e-3} & Historical
late-time evidence \\
Energy span ratio & \texttt{1.032} & Slight overshoot, but close to
exact \\
Min fidelity & \texttt{0.9626} & Remains bounded with no append
events \\
\end{longtable}
}

\subsubsection{Current QPU-facing
anchors}\label{current-qpu-facing-anchors}

\begin{itemize}
\tightlist
\item
  \textbf{Static / ground-state anchor:} the validated
  \texttt{pareto\_lean\ +\ phase3\_v1} L=2 ADAPT command above.
\item
  \textbf{Realtime dynamics law anchor:} Chapter 17A AP-McLachlan
  append--prune support-patch control.
\item
  \textbf{Driven calibration evidence:} the archival heavier-target
  driven calibration result above, retained as historical data only.
\end{itemize}

Do not treat archived secant-lead calibration artifacts as leading
settings or defaults when optimizing toward a real QPU run.

\begin{center}\rule{0.5\linewidth}{0.5pt}\end{center}

\begin{quote}
Canonical naming note: \texttt{MATH/Math.md} is the markdown authoring
source for manuscript sync, and \texttt{MATH/Math.tex} is the generated
typesetting twin used to build \texttt{MATH/Math.pdf}. Regenerate the
synchronized TeX/PDF pair with
\texttt{python\ MATH/build/build\_math\_from\_md.py} or
\texttt{MATH/build/build\_math.sh}. The archival repo-oriented
manuscript is \texttt{MATH/archive/archaic\_repo\_math.md}.
\end{quote}

\section{1. Parameter Manifest and Reader
Contract}\label{parameter-manifest-and-reader-contract}

This duplicate is intentionally symbolic-first. It keeps the same broad
manuscript shape as the existing mathematical notes,

\begin{enumerate}
\def\labelenumi{\arabic{enumi}.}
\tightlist
\item
  primitives first,
\item
  composite operators second,
\item
  explicit substitutions third,
\item
  reduced master forms last,
\end{enumerate}

but it suppresses almost all implementation labels, file names, mode
strings, and executable-surface commentary. The goal is to make the
document readable as mathematics rather than as a repository audit.

\subsection{1.1 Required parameter
manifest}\label{required-parameter-manifest}

A complete symbolic run or derivation should specify at least the
following data.

\begin{itemize}
\tightlist
\item
  Model family: Hubbard or Hubbard-Holstein.
\item
  Lattice size: \(L\).
\item
  Ordering map: interleaved or blocked, i.e.~a choice of
  \(p(i,\sigma)\).
\item
  Boundary condition: open or periodic.
\item
  Fermionic parameters: hopping \(t\), onsite interaction \(U\), and
  static site potentials \(v_i\).
\item
  Bosonic parameters: phonon frequency \(\omega_0\), electron-phonon
  coupling \(g\), and cutoff \(n_{\mathrm{ph,max}}\).
\item
  Boson encoding choice: binary or unary.
\item
  Variational family, when relevant: layerwise, physical-termwise,
  excitation-based, or adaptive.
\item
  Control parameters, when relevant: optimizer family, propagation rule,
  drive waveform parameters, and benchmark grid.
\end{itemize}

\subsection{1.2 Reader contract}\label{reader-contract}

This manuscript favors explicit substitution over compressed
abstraction.

\begin{itemize}
\tightlist
\item
  If a primitive operator exists, it is written first.
\item
  If a composite object is built from primitives, the primitive form is
  shown.
\item
  If a later equation can be simplified by inserting an earlier
  primitive explicitly, that insertion is made.
\item
  If multiple mathematical surfaces exist, such as different index
  orderings or boson encodings, the surface is named locally rather than
  hidden behind an unnamed default.
\end{itemize}

\subsubsection{1.2.1 Quick notation glossary for adaptive and
continuation
sections}\label{quick-notation-glossary-for-adaptive-and-continuation-sections}

This subsection is a compact map for the symbols that recur in the
adaptive-selection, split-aware, replay, and beam-style sections.

\begin{itemize}
\tightlist
\item
  \(g,m\): generic candidate generators. The manuscript uses both
  letters in different sections; they play the same role.
\item
  \(p\): candidate insertion position in the ordered scaffold.
\item
  \(r=(m,p)\): a candidate-position record, i.e.~a generator together
  with a chosen insertion position.
\item
  \(\mathcal G\): the full master generator pool.
\item
  \(\mathcal G^{(1)},\mathcal G^{(2)},\mathcal G^{(3)}\): progressively
  more restricted candidate universes, with \[
  \mathcal G^{(3)}\subseteq\mathcal G^{(2)}\subseteq\mathcal G^{(1)}\subseteq\mathcal G.
  \]
\item
  \(\mathcal P_g\) or \(\mathcal P_m\): admissible insertion positions
  for candidate \(g\) or \(m\).
\item
  \(W(p)\): inherited refit window used if a candidate is inserted at
  position \(p\).
\item
  \(W_{\mathrm{new}}(p)\): the newest-parameter portion of the refit
  window.
\item
  \(W_{|\theta|}(p)\): the older large-amplitude carry subset of the
  refit window.
\item
  \(\mathcal P_{\mathrm{avail}}^{(3)}\): the full available set of
  Phase-3 candidate-position records.
\item
  \(\mathcal S_2,\mathcal S_3\): the Phase-2 and Phase-3 shortlists in
  the cumulative selector chain
  \(\mathcal S_1\to\mathcal S_2\to\mathcal S_3\); thus Phase 2 forms
  \(\mathcal S_2\) from \(S_2\), and Phase 3 forms \(\mathcal S_3\) from
  \(S_3\).
\item
  \(\mathcal C_{\mathrm{split}}(m)\): the family of split children
  generated from macro generator \(m\).
\item
  \(\Gamma_{\mathrm{prune}}^{\mathrm{perm}}\): prune-permissibility gate
  deciding whether a compression pass may be attempted after admission,
  plateau, scheduled checkpoint, or terminal scaffold.
\item
  \(\mathcal J_{\mathrm{exp}}\): legacy name for locally expendable
  scaffold coordinates or generators considered by older prune passes;
  the corrected rule uses the mature candidate set and recovery tests of
  §11.5.1.
\item
  \(\mathcal O_*\): the finally selected adaptive scaffold, i.e.~the
  ordered generator sequence admitted by the adaptive stage.
\item
  \(\mathcal M_{\mathrm{warm}},\mathcal M_{\mathrm{refine}},\mathcal M_{\mathrm{replay}}\):
  archival appendix-only warm-start, optional refine, and replay
  manifolds.
\item
  \(\mathcal B\): a portable handoff/state bundle carrying manifest,
  amplitudes, energies, and continuation metadata.
\item
  \(\mathcal C\): continuation payload / provenance payload attached to
  a handoff bundle.
\item
  \(\Gamma_{\mathrm{stage}}\): stage-admissibility gate.
\item
  \(\Gamma_{\mathrm{sym}}\): symmetry-admissibility gate.
\item
  \(\mathfrak b\): a beam branch state. In §11 it indexes alternative
  ADAPT scaffolds; in §17 it indexes projected-dynamics branches over
  checkpoint segments.
\item
  \(\mathcal H_b\): admitted-record history carried by scaffold branch
  \(b\).
\item
  \(\mathcal A_b\): branch-local admissible admission set retained for
  scaffold-beam expansion.
\item
  \(\mathcal Q(b)\): scaffold-beam proposal family, usually
  \(\{\mathrm{stop}\}\cup\mathcal A_b\).
\item
  \(\eta_b\): branch-local stage label for the scaffold beam,
  e.g.~core/seed versus residual.
\item
  \(W_{\mathrm{act}}(\mathfrak b)\): the currently active position-probe
  window on branch \(b\).
\item
  \(\mathfrak F_c\): live frontier at beam round \(c\).
\item
  \(\mathfrak T_c\): terminal-branch pool at beam round \(c\).
\item
  \(\mathfrak W_{\mathrm{final}}\): final union of surviving frontier
  and terminal branches.
\item
  \(J_b\): cumulative branch objective tuple used to rank or prune
  branch \(b\); its coordinates are section-specific.
\end{itemize}

\subsubsection{1.2.2 Quick parameter glossary for adaptive and
continuation
sections}\label{quick-parameter-glossary-for-adaptive-and-continuation-sections}

The symbols below are the main scalar controls that appear in the
adaptive-selection, split-aware, and beam-style formulas.

\begin{itemize}
\tightlist
\item
  \(\lambda_{\mathrm{2q}}\): weight on two-qubit proxy burden.
\item
  \(\lambda_d\): weight on depth-style proxy burden.
\item
  \(\lambda_\theta\): weight on marginal parameter-count burden.
\item
  \(\lambda_F\): cheap-stage normalization weight multiplying the
  tangent metric term.
\item
  \(\lambda_H\): regularization strength for inherited-window Hessian
  blocks.
\item
  \(\rho\): trust-region radius scale used in one-dimensional local
  improvement models.
\item
  \(N_{\mathrm{cheap}}^{\max}\): absolute cap on the horizon-sized
  Phase-3 coarse shortlist.
\item
  \(N_{\max}\): maximum shortlist size.
\item
  \(f_{\mathrm{short}}\): shortlist fraction used in Phase 3.
\item
  \(w_R\): weight on useful-horizon lifetime burden.
\item
  \(w_D\): weight on lifetime depth/compile burden.
\item
  \(w_G\): weight on newly introduced grouped-measurement burden.
\item
  \(w_C\): weight on newly introduced compile or circuit burden.
\item
  \(w_c\): weight on any extra cheap-path correction burden.
\item
  \(\gamma_N\): exponent applied to novelty factors in full rerank
  scores.
\item
  \(\varepsilon\): small positive stabilizer used in denominators and
  tie-safe quotients.
\item
  \(\varepsilon_{\mathrm{nov}}\): novelty regularizer used when
  inverting or stabilizing tangent-overlap matrices.
\item
  \(z_\alpha\): confidence multiplier used to convert a raw gradient
  estimate into a lower-confidence gradient.
\item
  \(\sigma(g,p)\) or \(\sigma_r\): uncertainty estimate attached to the
  candidate-position gradient signal.
\item
  \(\tau_1\): Phase-2 shortlist threshold.
\item
  \(\tau_{\mathrm{split}}\): split-replacement margin required before a
  split child or child-set beats its parent.
\item
  \(\tau_{\mathrm{drop}}\): minimum per-depth absolute-error improvement
  regarded as meaningful by the stage controller.
\item
  \(M\): patience window used when counting repeated small-drop depths.
\item
  \(D_{\mathrm{left}}(t)\): remaining controller depth budget at
  selector step \(t\).
\item
  \(S_t(k)\): survival factor for having a still-useful admission
  opportunity \(k\) selector steps ahead.
\item
  \(Q_t(k)\): conditional usefulness factor for the \(k\)-ahead
  admission opportunity.
\item
  \(\widehat N_{\mathrm{rem}}(t)\): heuristic forecast of useful
  remaining admissions at selector step \(t\).
\item
  \(C_t\): confidence attached to the current useful-runway forecast.
\item
  \(K_{\mathrm{coarse}}(t)\): dynamic coarse-cap induced by the current
  useful-runway forecast.
\item
  \(H_t\): effective useful-depth horizon, \[
  H_t=\min\bigl(D_{\mathrm{left}}(t),\widehat N_{\mathrm{rem}}(t)\bigr).
  \]
\item
  \(W_{\mathrm{sat}}(t)\): saturation pressure induced by the current
  useful-runway forecast.
\item
  \(\pi_t\): selector mode tag, e.g.~fast / mixed / high-fidelity.
\item
  \(\mathrm{regime}_{\mathrm{sat}}(t)\): hungry / watch / plateau
  saturation regime.
\item
  \(\eta_{\mathrm{ret}}\): optional retained-gain fraction used only
  when a just-admitted ADAPT gain exceeds the useful-gain floor.
\item
  \(\beta_{\mathrm{keep}}\): repo-facing retained-gain floor alias;
  after the compression-prune correction it is a conditional large-gain
  guard, not the primary prune safety threshold.
\item
  \(\tau_{\mathrm{miss}}\): projected-miss threshold used in
  manifold-growth logic.
\item
  \(B_{\mathrm{child}}\): beam child cap per parent branch-expansion
  round.
\item
  \(B_{\mathrm{target}}(t)\): active batch target size at selector step
  \(t\).
\item
  \(B_{\mathrm{live}}\): live-frontier beam width retained after
  pruning.
\item
  \(B_{\mathrm{term}}\): cap on stored terminal branches.
\item
  \(M_{\mathrm{probe}}\): cap on distinct probed insertion positions
  retained in order.
\item
  \(\operatorname{age}_{j,t}\): branch-local admission age of scaffold
  coordinate or generator \(j\) at selector step \(t\); in §11.5.1,
  \(a_j(q)\) denotes normalized coordinate-amplitude history.
\item
  \(\varrho_j(t)\): normalized admission-rank score of coordinate or
  generator \(j\) at selector step \(t\).
\item
  \(\mathcal P_{\mathrm{protect}}(t)\): hard-protected coordinate set
  excluded from the live prune pass.
\item
  \(c_j(t)\): local prune cooldown on coordinate \(j\).
\item
  \(R_j^{\mathrm{cheap}}(t)\): cheap recoverability-prior score used for
  ADAPT prune nomination; low values mean ``try delete/refit,'' not
  ``delete.''
\item
  \(S_\theta(j,t)\): legacy small-angle nomination feature; after the
  recoverability-ladder correction it is only an optional component of
  \(R_j^{\mathrm{cheap}}\).
\item
  \(\tau_{\mathrm{stale}}\): stagnation scale used in prune nomination
  features before recoverability and safety checks.
\item
  \(\Xi_t^{\mathrm{pr}}\): exact rollback snapshot taken before a local
  ADAPT prune/remove-refit trial.
\item
  \(\Delta k_{\mathrm{cp}}\): checkpoint advance used in the
  projected-dynamics beam of §17, not in the scaffold beam of §11.
\item
  \(\Delta k_{\mathrm{probe}}\): probe advance used in the
  projected-dynamics beam of §17, not in the scaffold beam of §11.
\item
  \(C\): final beam round index / number of retained beam-update rounds.
\end{itemize}

In §11, the index \(t\) denotes a selector/controller step rather than
the real-time variable used later in the dynamics sections.

\subsection{1.3 Non-negotiable representational
conventions}\label{non-negotiable-representational-conventions}

\subsubsection{1.3.1 Internal Pauli
alphabet}\label{internal-pauli-alphabet}

In formulas below we use the standard alphabet \[
\{I,X,Y,Z\}.
\] If one prefers the lower-case alphabet \(\{e,x,y,z\}\), the
identification is simply \[
e \leftrightarrow I,
\qquad
x \leftrightarrow X,
\qquad
y \leftrightarrow Y,
\qquad
z \leftrightarrow Z.
\]

\subsubsection{1.3.2 Pauli-word and qubit
ordering}\label{pauli-word-and-qubit-ordering}

Pauli words and computational-basis labels are written left-to-right as
\[
q_{N_q-1}\cdots q_1 q_0,
\] with qubit \(q_0\) the rightmost symbol and also the
least-significant bit in basis-index arithmetic.

\subsubsection{1.3.3 Canonical algebra
sources}\label{canonical-algebra-sources}

The canonical mathematical primitives are:

\begin{itemize}
\tightlist
\item
  Jordan-Wigner ladder operators \(\hat c_p^\dagger\) and \(\hat c_p\),
\item
  fermion number operators \(\hat n_p\),
\item
  boson ladder operators \(\hat b_i^\dagger\) and \(\hat b_i\),
\item
  Hermitian Pauli words and Pauli polynomials,
\item
  ordered exponentials generated from those primitives.
\end{itemize}

\subsubsection{1.3.4 Surface-specific
defaults}\label{surface-specific-defaults}

This symbolic version does not pretend there is a single universal
surface. The mathematics admits several equally valid choices:

\begin{itemize}
\tightlist
\item
  interleaved or blocked fermion ordering,
\item
  open or periodic boundary conditions,
\item
  binary or unary boson encoding,
\item
  static or driven Hamiltonians,
\item
  fixed-structure or adaptive variational manifolds.
\end{itemize}

Whenever a formula depends on the choice of surface, that dependence is
shown locally.

\subsection{1.4 Canonical mathematical
anchors}\label{canonical-mathematical-anchors}

The main objects carried through the manuscript are:

\begin{itemize}
\tightlist
\item
  the ordering map \(p(i,\sigma)\),
\item
  the ladder primitives \(\hat c_p^\dagger\), \(\hat c_p\),
  \(\hat b_i^\dagger\), \(\hat b_i\),
\item
  the density operators \(\hat n_{i\sigma}\), \(\hat n_i\), and
  \(\hat D_i\),
\item
  the Hubbard and Hubbard-Holstein Hamiltonians,
\item
  the reference state \(|\psi_{\mathrm{ref}}\rangle\),
\item
  the variational manifold \(\mathcal M\) and its ansatz map
  \(\theta \mapsto |\psi(\theta)\rangle\),
\item
  the propagation rule \(U(t)\),
\item
  the continuation or handoff bundle \(\mathcal B\).
\end{itemize}

\section{2. Ordering, Indexing, and Register
Layout}\label{ordering-indexing-and-register-layout}

\subsection{2.1 Site, spin, and mode
indices}\label{site-spin-and-mode-indices}

The site index is \[
i \in \{0,1,\dots,L-1\},
\] and the spin label is stored as \[
\sigma \in \{\uparrow,\downarrow\} \equiv \{0,1\}.
\] The fermion mode index is the Jordan-Wigner qubit index.

\subsubsection{2.1.1 Interleaved ordering}\label{interleaved-ordering}

The interleaved map is \[
p(i,\sigma)=2i+\sigma,
\] so that \[
p(i,\uparrow)=2i,
\qquad
p(i,\downarrow)=2i+1.
\]

\subsubsection{2.1.2 Blocked ordering}\label{blocked-ordering}

The blocked map is \[
p(i,\uparrow)=i,
\qquad
p(i,\downarrow)=L+i.
\]

\subsection{2.2 Pauli-word placement and basis-index
extraction}\label{pauli-word-placement-and-basis-index-extraction}

If a Pauli letter acts on qubit \(q\), then in a printed word of length
\(N_q\) it sits at string position \[
\operatorname{pos}(q)=N_q-1-q.
\]

If the computational-basis index is \(k\), then the occupation bit on
qubit \(q\) is \[
b_q(k)=\left\lfloor\frac{k}{2^q}\right\rfloor \bmod 2.
\] Thus the printed bitstring and integer basis index obey the same
rightmost-\(q_0\) convention.

\subsection{2.3 Full HH register layout}\label{full-hh-register-layout}

The fermion register uses \[
N_{\mathrm{ferm}}=2L
\] qubits.

If the local phonon cutoff is \(n_{\mathrm{ph,max}}\), then the local
phonon Hilbert dimension is \[
d=n_{\mathrm{ph,max}}+1.
\] The number of phonon qubits per site is \[
q_{\mathrm{pb}}=
\begin{cases}
\max\{1,\lceil \log_2 d\rceil\}, & \text{binary},\\
d, & \text{unary}.
\end{cases}
\] Therefore the total HH qubit count is \[
N_q=2L+Lq_{\mathrm{pb}}.
\]

In qubit-index order, the register is naturally partitioned as \[
[\text{fermions}\;|\;\text{site-0 phonons}\;|\;\cdots\;|\;\text{site-(L-1) phonons}].
\] In printed order \(q_{N_q-1}\cdots q_0\), the high-index phonon
blocks appear on the left.

\section{3. Fermionic Primitives and Direct
Substitution}\label{fermionic-primitives-and-direct-substitution}

\subsection{3.1 Jordan-Wigner ladder
primitives}\label{jordan-wigner-ladder-primitives}

For mode \(p\), the creation operator is \[
\hat c_p^\dagger
=
\frac{1}{2}(X_p-iY_p)\prod_{r=0}^{p-1}Z_r,
\] and the annihilation operator is \[
\hat c_p
=
\frac{1}{2}(X_p+iY_p)\prod_{r=0}^{p-1}Z_r.
\]

In fully expanded printed-word form, the operator acts on the word
ordered as \(q_{N_q-1}\cdots q_0\) with the Jordan-Wigner \(Z\)-string
extending through all modes below \(p\).

\subsection{3.2 Number primitive}\label{number-primitive}

The fermion number operator is \[
\hat n_p=\hat c_p^\dagger\hat c_p=\frac{I-Z_p}{2}.
\]

\subsection{3.3 Site densities and doublon
operator}\label{site-densities-and-doublon-operator}

At site \(i\), define \[
\hat n_i = \hat n_{i\uparrow}+\hat n_{i\downarrow},
\] and the doublon operator \[
\hat D_i = \hat n_{i\uparrow}\hat n_{i\downarrow}.
\] By direct substitution, \[
\hat D_i = \frac{1}{4}(I-Z_{p(i,\uparrow)})(I-Z_{p(i,\downarrow)}).
\]

\subsection{3.4 Worked ordering
substitutions}\label{worked-ordering-substitutions}

\subsubsection{3.4.1 Interleaved}\label{interleaved}

Under interleaved ordering, \[
\hat n_i = I - \frac{1}{2}(Z_{2i}+Z_{2i+1}),
\] and \[
\hat D_i = \frac{1}{4}(I-Z_{2i})(I-Z_{2i+1}).
\] For nearest-neighbor spin-up hopping in a one-dimensional chain, \[
\hat c_{i\uparrow}^\dagger \hat c_{i+1,\uparrow} + \hat c_{i+1,\uparrow}^\dagger \hat c_{i\uparrow}
=
\frac{1}{2}\Bigl(X_{2i} Z_{2i+1} X_{2i+2} + Y_{2i} Z_{2i+1} Y_{2i+2}\Bigr).
\]

\subsubsection{3.4.2 Blocked}\label{blocked}

Under blocked ordering, \[
\hat n_i = I - \frac{1}{2}(Z_i + Z_{L+i}),
\] and \[
\hat D_i = \frac{1}{4}(I-Z_i)(I-Z_{L+i}).
\] For spin-up nearest-neighbor hopping, the Jordan-Wigner string is
empty because the two blocked spin-up modes are adjacent, so \[
\hat c_{i\uparrow}^\dagger \hat c_{i+1,\uparrow} + \hat c_{i+1,\uparrow}^\dagger \hat c_{i\uparrow}
=
\frac{1}{2}\Bigl(X_i X_{i+1} + Y_i Y_{i+1}\Bigr).
\] The general blocked-ordering formula uses the same Jordan-Wigner
interval rule as before, with an empty product for adjacent modes.

\section{4. Boson Primitives and
Encodings}\label{boson-primitives-and-encodings}

\subsection{4.1 Local boson Hilbert
space}\label{local-boson-hilbert-space}

For each site \(i\), the truncated boson Hilbert space is \[
\mathcal H_{b,i}=\operatorname{span}\{|n_i\rangle : 0\le n_i\le n_{\mathrm{ph,max}}\}.
\] The ladder operators satisfy \[
\hat b_i^\dagger |n\rangle = \sqrt{n+1}\,|n+1\rangle,
\qquad
\hat b_i |n\rangle = \sqrt{n}\,|n-1\rangle,
\] with \[
\hat n_{b,i}=\hat b_i^\dagger \hat b_i,
\qquad
\hat x_i=\hat b_i+\hat b_i^\dagger,
\qquad
\hat P_i=i(\hat b_i^\dagger-\hat b_i).
\]

\subsection{4.2 Binary encoding}\label{binary-encoding}

In binary encoding, each local occupation number \(n\in\{0,\dots,d-1\}\)
is mapped to a binary word on \(q_{\mathrm{pb}}=\lceil \log_2 d\rceil\)
qubits: \[
|n\rangle \longmapsto |\operatorname{bin}(n)\rangle.
\] The truncated number operator is then \[
\hat n_{b,i} = \sum_{n=0}^{d-1} n\,|n\rangle\langle n|,
\] understood as an operator on the binary-encoded subspace.

\subsection{4.3 Unary encoding}\label{unary-encoding}

In unary encoding, each local occupation \(n\) is represented by a
one-hot basis state on \(d\) qubits: \[
|n\rangle \longmapsto |0\cdots 010\cdots 0\rangle,
\] with the single \(1\) marking the occupied unary slot. The number
operator remains \[
\hat n_{b,i} = \sum_{n=0}^{d-1} n\,|n\rangle\langle n|,
\] but now the basis states are embedded in a one-hot code space.

\subsection{4.4 Phonon vacuum}\label{phonon-vacuum}

\subsubsection{4.4.1 Binary vacuum}\label{binary-vacuum}

The local binary vacuum is \[
|0\rangle_{b,i} \longmapsto |00\cdots 0\rangle.
\]

\subsubsection{4.4.2 Unary vacuum}\label{unary-vacuum}

The local unary vacuum is the one-hot state carrying occupancy zero, \[
|0\rangle_{b,i} \longmapsto |10\cdots 0\rangle,
\] or whichever one-hot convention is chosen consistently for the lowest
occupation.

\subsection{4.5 Bose-Hubbard ADAPT
reference}\label{bose-hubbard-adapt-reference}

For the number-conserving Bose-Hubbard benchmark, \[
\hat H_{\mathrm{BH}}
=
-t\sum_{\langle i,j\rangle}
\left(\hat b_i^\dagger \hat b_j+\hat b_j^\dagger \hat b_i\right)
+\omega_0\sum_i \hat n_i
+\frac{U}{2}\sum_i \hat n_i(\hat n_i-I)
+\sum_i v_i \hat n_i,
\] the pure vacuum is a stationary sector by construction. If the exact
comparison is the unrestricted truncated boson register, the lowest
sector for the default small benchmarks can be a nonzero one-boson
sector. Therefore the static ADAPT reference used for the
problem-generic Bose-Hubbard route is a one-boson product state, \[
|\phi_{\mathrm{BH},0}\rangle
=
|0\rangle_0\otimes\cdots\otimes |1\rangle_{i_*}
\otimes\cdots\otimes |0\rangle_{L-1},
\qquad
i_*=\arg\min_i v_i,
\] rather than the bosonic vacuum. This avoids the false convergence
condition \[
g_A(0)=2\,\operatorname{Im}\langle H\phi_{\mathrm{vac}}|A\phi_{\mathrm{vac}}\rangle=0
\quad \forall A\in\mathcal P_{\mathrm{BH}},
\] which otherwise stops phase-3 ADAPT at depth zero even when the
unrestricted truncated-register exact energy lies below the vacuum
energy.

\subsection{4.6 Fixed-number binary boson transport and the BEMPA
pool}\label{fixed-number-binary-boson-transport-and-the-bempa-pool}

For number-conserving bosonic benchmarks, the binary-encoded multilevel
particle ansatz (BEMPA) of Bahrami and Sawaya gives a symmetry-native
candidate pool rather than a penalty-enforced sector rule. This
construction is canonical for fixed-\(N_B\) boson routes such as
Bose-Hubbard, and should not be treated as a hard phonon-number
constraint for the Hubbard-Holstein phonon sector.

With the repo convention that \(q_0\) is the least-significant and
rightmost printed bit, the binary boson number operator is \[
\widehat N_B
=
\sum_i\widehat n_i^{\mathrm{bin}}
=
\sum_i\sum_{k=0}^{q_{\mathrm{pb}}-1}
2^k\,\frac{I-Z_{i,k}}2 .
\] If the local cutoff dimension is not a power of two, the binary
register also contains invalid states. In that case a generator must be
checked against the valid code projector \(\Pi_{\mathrm{valid}}\) as
well as \(\widehat N_B\). The default first BEMPA benchmark should
therefore use power-of-two local dimensions.

The same-significance exchange generator between mode bits \((i,k)\) and
\((j,k)\) is \[
G_A^{(i,j;k)}
=
\frac12\left(X_{i,k}Y_{j,k}-Y_{i,k}X_{j,k}\right),
\] with sign fixed by the convention \(U_A(\theta)=\exp(-i\theta G_A)\).
It mixes only \(|01\rangle\) and \(|10\rangle\) within the same
significance block, so \[
[G_A^{(i,j;k)},\widehat N_B]=0.
\]

The binary carry-transfer generator trades two lower-significance
occupied bits for one higher-significance occupied bit. For lower bits
\(a=(i,k-1)\), \(b=(j,k-1)\), and higher bit \(c=(\ell,k)\), both sides
of the transition \[
|1_a1_b0_c\rangle \longleftrightarrow |0_a0_b1_c\rangle
\] carry weight \(2^{k-1}+2^{k-1}=2^k\). With \((a,b,c)\) ordered
consistently with this displayed transition, use \[
G_B^{(a,b,c)}
=
\frac14\left(
X_aX_bY_c
-X_aY_bX_c
-Y_aX_bX_c
-Y_aY_bY_c
\right),
\] and verify the basis-order sign on the two active basis states before
making the operator executable. Algebraically, \[
[G_B^{(a,b,c)},\widehat N_B]=0.
\] The four Pauli words in \(G_B\) commute pairwise, so the exponential
is an exact product of Pauli rotations: \[
e^{-i\alpha G_B}
=
\prod_{r=1}^{4} e^{-i\alpha c_rP_r},
\qquad
G_B=\sum_{r=1}^4 c_rP_r,
\qquad
[P_r,P_s]=0.
\]

The fixed-number BEMPA pool is \[
\mathcal G_{\mathrm{BEMPA}}
:=
\mathcal G_A\cup\mathcal G_B,
\] where \[
\mathcal G_A
:=
\{G_A^{(i,j;k)}:i<j,\ 0\le k<q_{\mathrm{pb}}\},
\] and \[
\mathcal G_B
:=
\{G_B^{(i,k-1;\,j,k-1;\,\ell,k)}:
0<k<q_{\mathrm{pb}},\ i,j,\ell\ \text{admissible}\}.
\] Every BEMPA generator satisfies the exact symmetry gate \[
\Gamma_{\mathrm{sym}}^{B}(m)
:=
\mathbf 1\!\left([G_m,\widehat N_B]=0\right)=1.
\] For an initialized fixed sector
\(\widehat N_B|\psi_0\rangle=N_0|\psi_0\rangle\), the leakage observable
\[
\mathcal L_N(\psi)
:=
\langle\psi|(\widehat N_B-N_0)^2|\psi\rangle
\] vanishes throughout a BEMPA-only variational manifold, up to
numerical and measurement error.

The \(G_A\) family alone is generally not a complete fixed-sector
transport pool, because it preserves the bit count inside each
significance block. For diagnostics, define the fixed-sector graph \[
\mathfrak G_{N_0}:=(V_{N_0},E_A\cup E_B),
\qquad
V_{N_0}:=\{\mathbf n:0\le n_i<d,\ \sum_i n_i=N_0\}.
\] The \(E_A\) edges come from same-significance exchanges and the
\(E_B\) edges come from carry transfers. The selector and prune layer
should treat \(G_B\) records as possible carry-bridge coordinates:
deleting all carry-transfer directions may fragment the reachable
fixed-\(N_B\) manifold even if a local energy refit looks temporarily
safe.

\section{5. Hubbard Hamiltonian by Explicit
Substitution}\label{hubbard-hamiltonian-by-explicit-substitution}

Write the static Hubbard Hamiltonian as \[
\hat H_{\mathrm{Hub}} = \hat H_t + \hat H_U + \hat H_v.
\]

\subsection{5.1 Kinetic term}\label{kinetic-term}

The kinetic term is \[
\hat H_t = -t\sum_{\langle i,j\rangle,\sigma}
\left(
\hat c_{i\sigma}^\dagger \hat c_{j\sigma}
+
\hat c_{j\sigma}^\dagger \hat c_{i\sigma}
\right).
\] For two fermion modes \(p<q\), direct Jordan-Wigner substitution
gives \[
\hat c_p^\dagger \hat c_q + \hat c_q^\dagger \hat c_p
=
\frac{1}{2}
\left(
X_p Z_{p+1}\cdots Z_{q-1} X_q
+
Y_p Z_{p+1}\cdots Z_{q-1} Y_q
\right).
\]

\subsection{5.2 Onsite interaction}\label{onsite-interaction}

The onsite interaction term is \[
\hat H_U = U\sum_i \hat n_{i\uparrow}\hat n_{i\downarrow}
= U\sum_i \hat D_i.
\] Substituting the number primitive, \[
\hat H_U = \frac{U}{4}\sum_i (I-Z_{p(i,\uparrow)})(I-Z_{p(i,\downarrow)}).
\]

\subsection{5.3 Static potential term}\label{static-potential-term}

The static potential term is \[
\hat H_v = \sum_{i,\sigma} v_i \hat n_{i\sigma}.
\] After substitution, \[
\hat H_v = \sum_i v_i\left(I-\frac{1}{2}(Z_{p(i,\uparrow)}+Z_{p(i,\downarrow)})\right).
\]

\subsection{5.4 Fully substituted Hubbard
Hamiltonian}\label{fully-substituted-hubbard-hamiltonian}

Combining the pieces, \[
\hat H_{\mathrm{Hub}}
=
-t\sum_{\langle i,j\rangle,\sigma}
\left(
\hat c_{i\sigma}^\dagger \hat c_{j\sigma} + \hat c_{j\sigma}^\dagger \hat c_{i\sigma}
\right)
+
\frac{U}{4}\sum_i (I-Z_{p(i,\uparrow)})(I-Z_{p(i,\downarrow)})
+
\sum_i v_i\left(I-\frac{1}{2}(Z_{p(i,\uparrow)}+Z_{p(i,\downarrow)})\right),
\] with the kinetic term expanded further into Pauli strings as needed.

\section{6. Hubbard-Holstein Hamiltonian by Explicit
Substitution}\label{hubbard-holstein-hamiltonian-by-explicit-substitution}

The driven Hubbard-Holstein Hamiltonian is \[
\hat H_{\mathrm{HH}}(t)
=
\hat H_t + \hat H_U + \hat H_v + \hat H_{\mathrm{ph}} + \hat H_g + \hat H_{\mathrm{drive}}(t).
\]

The driven onsite-density contribution is written in fixed-label Pauli
form as \[
\hat H_{\mathrm{drive}}(t)=\sum_{\lambda\in\mathcal L_{\mathrm{drive}}} c_\lambda(t) P_\lambda,
\] where the operator labels \(P_\lambda\) are static and only the
coefficients \(c_\lambda(t)\) vary in time. For the site-density drive
used in the active HH dynamics stack, \[
\Delta c[Z_{p(i,\uparrow)}](t)=\Delta c[Z_{p(i,\downarrow)}](t)=-\frac{1}{2}v_i(t),
\] and, if one retains the scalar offset explicitly, \[
\Delta c[I](t)=\sum_i v_i(t).
\]

\subsection{6.1 Phonon energy}\label{phonon-energy}

The phonon energy is \[
\hat H_{\mathrm{ph}} = \omega_0\sum_i \left(\hat n_{b,i}+\frac{1}{2}I\right).
\]

\subsubsection{6.1.1 Unary explicit form}\label{unary-explicit-form}

In unary form, \(\hat n_{b,i}\) is diagonal on the one-hot occupation
basis: \[
\hat n_{b,i} = \sum_{n=0}^{d-1} n\,|n\rangle_i\langle n|.
\] For the full qubit-level Pauli reduction in the spin-block JW + unary
setting, including the explicit \(I/Z\) form of
\(\hat H_{\mathrm{ph}}\), see Section 6.2.3.4(a).

\subsubsection{6.1.2 Binary explicit form}\label{binary-explicit-form}

In binary form, \[
\hat n_{b,i} = \sum_{n=0}^{d-1} n\,|\operatorname{bin}(n)\rangle_i\langle \operatorname{bin}(n)|,
\] understood as an operator on the truncated binary code space.

\subsection{6.2 Electron-phonon
coupling}\label{electron-phonon-coupling}

A natural density-shifted coupling is \[
\hat H_g = g\sum_i \hat x_i\bigl(\hat n_i-\bar n I\bigr),
\] where \(\bar n = N_e/L\) is the reference density, often \(\bar n=1\)
at half filling.

\subsubsection{6.2.1 Unary explicit form}\label{unary-explicit-form-1}

In unary encoding, \(\hat x_i\) acts as nearest-neighbor hopping on the
truncated local occupation chain: \[
\hat x_i = \sum_{n=0}^{d-2} \sqrt{n+1}\Bigl(|n+1\rangle_i\langle n| + |n\rangle_i\langle n+1|\Bigr).
\] Thus \[
\hat H_g = g\sum_i \hat x_i\bigl(\hat n_i-\bar n I\bigr)
\] with \(\hat x_i\) understood in that unary basis. For the detailed
spin-block JW + unary Pauli reduction of the Holstein sector, including
the explicit \(ZXX/ZYY\) expansion, see Sections 6.2.3.2--6.2.3.4.

\subsubsection{6.2.2 Binary explicit form}\label{binary-explicit-form-1}

In binary encoding, the same operator is written on the binary code
space: \[
\hat x_i = \sum_{n=0}^{d-2} \sqrt{n+1}\Bigl(|\operatorname{bin}(n+1)\rangle_i\langle \operatorname{bin}(n)| + |\operatorname{bin}(n)\rangle_i\langle \operatorname{bin}(n+1)|\Bigr),
\] with the same density factor \((\hat n_i-\bar n I)\).

The physical HH phonon sector is not a fixed-phonon-number sector. Since
\[
[\hat x_i,\widehat N_{\mathrm{ph}}]\ne0,
\qquad
\widehat N_{\mathrm{ph}}:=\sum_i \hat n_{b,i},
\] the electron-phonon coupling gives, in general, \[
[H_{\mathrm{HH}},\widehat N_{\mathrm{ph}}]\ne0.
\] Therefore the BEMPA fixed-\(N_B\) pool from Section 4.6 is not a hard
phonon-number-preserving production ansatz for HH. Its HH use is only
indirect: binary carry structure can refine phonon transition operators
when \(n_{\mathrm{ph,max}}\ge2\).

Define local binary transition quadratures \[
C_{i,n}^{(+)}
:=
|n+1\rangle_i\langle n|
+|n\rangle_i\langle n+1|,
\] and \[
C_{i,n}^{(-)}
:=
i\left(|n+1\rangle_i\langle n|-|n\rangle_i\langle n+1|\right).
\] Then \[
x_i
=
\sum_{n=0}^{d-2}\sqrt{n+1}\,C_{i,n}^{(+)},
\qquad
P_i
=
\sum_{n=0}^{d-2}\sqrt{n+1}\,C_{i,n}^{(-)}.
\] These transition-resolved operators are not BEMPA generators, because
they do not preserve phonon number. They are BEMPA-inspired binary
carry-resolved building blocks for HH phonon displacement channels. For
example, the existing polaronic families may be refined into child
records such as \[
G_{i,n}^{(\mathrm{cd})}:=\tilde n_i C_{i,n}^{(-)},
\qquad
G_{i,n}^{(\mathrm{dd})}:=D_i C_{i,n}^{(-)},
\] and \[
G_{ij,n}^{(\mathrm{hd})}
:=
T_{ij}^{(+)}\left(C_{i,n}^{(-)}-C_{j,n}^{(-)}\right),
\qquad
G_{ij,n}^{(\mathrm{od})}
:=
J_{ij}^{(-)}\left(C_{i,n}^{(+)}-C_{j,n}^{(+)}\right).
\] For \(d=2\), these reduce to the present one-qubit phonon
quadratures, \[
C_{i,0}^{(+)}=X_{i,0},
\qquad
C_{i,0}^{(-)}=Y_{i,0}.
\] Thus the transition-resolved refinement is mainly relevant for
\(n_{\mathrm{ph,max}}\ge2\), where binary carry strings can expose which
\(n\leftrightarrow n+1\) phonon channels are geometrically useful.

\subsection{6.2.3 Worked Holstein-only reduction after Jordan-Wigner
fermions and unary
phonons}\label{worked-holstein-only-reduction-after-jordan-wigner-fermions-and-unary-phonons}

Set \[
N_b = n_{\mathrm{ph,max}}.
\] In this worked subsection, isolate the Holstein part of the
Hubbard-Holstein Hamiltonian: \[
\hat H_{\mathrm{Hol}}
=
\omega_0\sum_{i=0}^{L-1}\hat b_i^\dagger \hat b_i
\;+
\;g\sum_{i=0}^{L-1}(\hat b_i^\dagger+\hat b_i)(\hat n_i-1),
\qquad
\hat n_i=\hat n_{i\uparrow}+\hat n_{i\downarrow}.
\]

\subsubsection{6.2.3.1 Qubit indexing used}\label{qubit-indexing-used}

For fermions, use spin-block Jordan-Wigner ordering: \[
q(i,\sigma)=i+\sigma L,
\qquad
\sigma=0\equiv\uparrow,
\quad
\sigma=1\equiv\downarrow.
\] Thus \[
q_{i\uparrow}=i,
\qquad
q_{i\downarrow}=i+L,
\qquad
i=0,\dots,L-1.
\]

For phonons, use a unary register with cutoff \(N_b\). Per site \(i\),
allocate \(N_b+1\) qubits labelled by \(n=0,\dots,N_b\), and place all
fermions first and phonons second: \[
q_{i,n}^{(b)} = 2L + i(N_b+1)+n,
\qquad
n=0,\dots,N_b.
\]

Let \(X_q\), \(Y_q\), and \(Z_q\) denote the Pauli operators on qubit
\(q\).

\subsubsection{6.2.3.2 Jordan-Wigner fermionic number
operators}\label{jordan-wigner-fermionic-number-operators}

Jordan-Wigner gives \[
\hat n_{i\sigma} = \frac{I-Z_{q_{i\sigma}}}{2},
\] so \[
\hat n_i-1
=
\left(\frac{I-Z_{q_{i\uparrow}}}{2}+\frac{I-Z_{q_{i\downarrow}}}{2}\right)-1
=
-\frac{Z_{q_{i\uparrow}}+Z_{q_{i\downarrow}}}{2}.
\]

\subsubsection{6.2.3.3 Unary phonon ladder operators in Pauli
form}\label{unary-phonon-ladder-operators-in-pauli-form}

Define the single-qubit combinations \[
(+)_{q}=\frac{X_q+iY_q}{2},
\qquad
(-)_{q}=\frac{X_q-iY_q}{2}.
\]

Then the unary truncated ladder operators are \[
\hat b_i^\dagger
=
\sum_{n=0}^{N_b-1}\sqrt{n+1}\,(+)_{q_{i,n}^{(b)}}\,(-)_{q_{i,n+1}^{(b)}},
\] \[
\hat b_i
=
\sum_{n=1}^{N_b}\sqrt{n}\,(-)_{q_{i,n-1}^{(b)}}\,(+)_{q_{i,n}^{(b)}}.
\]

A convenient real Pauli expansion is \[
\hat b_i^\dagger+\hat b_i
=
\frac12\sum_{n=0}^{N_b-1}\sqrt{n+1}
\Big(
X_{q_{i,n}^{(b)}}X_{q_{i,n+1}^{(b)}}
+
Y_{q_{i,n}^{(b)}}Y_{q_{i,n+1}^{(b)}}
\Big).
\]

\subsubsection{6.2.3.4 Explicit Holstein Hamiltonian in the qubit
basis}\label{explicit-holstein-hamiltonian-in-the-qubit-basis}

\paragraph{(a) Phonon energy term}\label{a-phonon-energy-term}

On the unary one-hot code space, the phonon number operator is diagonal
and may be written as \[
\hat b_i^\dagger\hat b_i \equiv \hat n_i^{(b)}
=
\sum_{n=0}^{N_b} n\,\frac{I-Z_{q_{i,n}^{(b)}}}{2}
\qquad
\text{(exact on the one-hot subspace)}.
\] Therefore \[
\boxed{
\hat H_{\mathrm{ph}}
=
\omega_0\sum_{i=0}^{L-1}\sum_{n=0}^{N_b} n\,\frac{I-Z_{q_{i,n}^{(b)}}}{2}
}.
\]

Equivalently, as constant plus Pauli-\(Z\), \[
\hat H_{\mathrm{ph}}
=
\underbrace{
\omega_0\frac12\sum_i\sum_{n=0}^{N_b} n
}_{\displaystyle \omega_0 L\frac{N_b(N_b+1)}{4}}
I
-
\frac{\omega_0}{2}\sum_{i=0}^{L-1}\sum_{n=0}^{N_b} n\,Z_{q_{i,n}^{(b)}}.
\] The global constant shift may be dropped if desired.

\paragraph{(b) Electron-phonon coupling
term}\label{b-electron-phonon-coupling-term}

Using \[
\hat n_i-1 = -\frac{Z_{q_{i\uparrow}}+Z_{q_{i\downarrow}}}{2}
\] and the unary expansion of \(\hat b_i^\dagger+\hat b_i\), one obtains
\[
\boxed{
\hat H_{e\text{-}ph}
=
-\frac{g}{4}\sum_{i=0}^{L-1}\sum_{n=0}^{N_b-1}\sqrt{n+1}
\bigl(Z_{q_{i\uparrow}}+Z_{q_{i\downarrow}}\bigr)
\Big(
X_{q_{i,n}^{(b)}}X_{q_{i,n+1}^{(b)}}
+
Y_{q_{i,n}^{(b)}}Y_{q_{i,n+1}^{(b)}}
\Big)
}.
\]

The fully expanded Pauli-string sum is \[
\hat H_{e\text{-}ph}
=
-\frac{g}{4}\sum_{i=0}^{L-1}\sum_{n=0}^{N_b-1}\sqrt{n+1}
\Big[
Z_{q_{i\uparrow}}X_{q_{i,n}^{(b)}}X_{q_{i,n+1}^{(b)}}
+
Z_{q_{i\uparrow}}Y_{q_{i,n}^{(b)}}Y_{q_{i,n+1}^{(b)}}
+
Z_{q_{i\downarrow}}X_{q_{i,n}^{(b)}}X_{q_{i,n+1}^{(b)}}
+
Z_{q_{i\downarrow}}Y_{q_{i,n}^{(b)}}Y_{q_{i,n+1}^{(b)}}
\Big].
\] Each summand is three-local.

\paragraph{Final Holstein Hamiltonian in JW + unary
form}\label{final-holstein-hamiltonian-in-jw-unary-form}

Collecting the phonon-energy and electron-phonon pieces, \[
\boxed{
\hat H_{\mathrm{Hol}}^{\mathrm{(JW+unary)}}
=
\omega_0\sum_{i=0}^{L-1}\sum_{n=0}^{N_b} n\,\frac{I-Z_{q_{i,n}^{(b)}}}{2}
-
\frac{g}{4}\sum_{i=0}^{L-1}\sum_{n=0}^{N_b-1}\sqrt{n+1}
\bigl(Z_{q_{i\uparrow}}+Z_{q_{i\downarrow}}\bigr)
\Big(
X_{q_{i,n}^{(b)}}X_{q_{i,n+1}^{(b)}}
+
Y_{q_{i,n}^{(b)}}Y_{q_{i,n+1}^{(b)}}
\Big)
}.
\]

\subsection{6.3 Time-dependent density
drive}\label{time-dependent-density-drive}

A general density drive may be written as \[
\hat H_{\mathrm{drive}}(t) = \sum_i v_i(t)\hat n_i.
\] A useful waveform class is \[
v_i(t)=A s_i \sin(\omega t + \phi)\exp\!\left[-\frac{(t-t_0)^2}{2\bar t^2}\right],
\] where \(A\) is the amplitude, \(s_i\) is a static spatial pattern,
\(\omega\) is the drive frequency, \(\phi\) is the phase, and \(\bar t\)
is the envelope width.

\subsubsection{6.3.1 Generic density-drive waveform
surface}\label{generic-density-drive-waveform-surface}

After inserting
\(\hat n_i=I-\tfrac12(Z_{p(i,\uparrow)}+Z_{p(i,\downarrow)})\), the
drive separates into an identity offset and a sum of time-dependent
\(Z\)-coefficients. Thus the time dependence resides in the coefficients
rather than in the operator labels.

\subsection{6.4 Two HH assembly
surfaces}\label{two-hh-assembly-surfaces}

\subsubsection{6.4.1 Exact Hamiltonian
surface}\label{exact-hamiltonian-surface}

The exact surface treats \(\hat H_{\mathrm{HH}}(t)\) as an operator to
be diagonalized, exponentiated, projected, or applied directly to
statevectors.

\subsubsection{6.4.2 Variational HH ansatz
surface}\label{variational-hh-ansatz-surface}

The variational surface treats the same physical pieces as generator
families from which one builds a structured manifold \[
|\psi(\theta)\rangle = U(\theta)|\psi_{\mathrm{ref}}\rangle.
\]

\subsection{6.5 Fully substituted HH master
form}\label{fully-substituted-hh-master-form}

The full symbolic master form is therefore \[
\hat H_{\mathrm{HH}}(t)
=
\hat H_{\mathrm{Hub}}
+
\omega_0\sum_i \left(\hat n_{b,i}+\frac{1}{2}I\right)
+
 g\sum_i \hat x_i(\hat n_i-\bar n I)
+
\sum_i v_i(t)\hat n_i,
\] with \(\hat n_i\), \(\hat n_{b,i}\), and \(\hat x_i\) expanded
according to the chosen fermion ordering and boson encoding.

\section{7. Reference States and Exact Sector
ED}\label{reference-states-and-exact-sector-ed}

\subsection{7.1 Fermionic Hartree-Fock
determinant}\label{fermionic-hartree-fock-determinant}

Let \(\mathcal O\) denote the occupied spin-orbital set in a Slater
determinant reference. Then \[
|\Phi_{\mathrm{HF}}\rangle = \prod_{p\in\mathcal O} \hat c_p^\dagger |\mathrm{vac}_f\rangle.
\]

\subsection{7.2 Hubbard-Holstein reference
state}\label{hubbard-holstein-reference-state}

The natural HH reference is the tensor product of the phonon vacuum and
fermionic Hartree-Fock state: \[
|\psi_{\mathrm{ref}}^{\mathrm{HH}}\rangle = |\mathrm{vac}_{\mathrm{ph}}\rangle \otimes |\Phi_{\mathrm{HF}}\rangle.
\]

\subsection{7.3 Exact HH sector basis}\label{exact-hh-sector-basis}

Fixing the fermion particle numbers \((N_\uparrow,N_\downarrow)\), the
exact sector basis is \[
\mathcal H_{\mathrm{sector}} = \operatorname{span}\{|f\rangle\otimes|\mathbf n_b\rangle : |f\rangle\in\mathcal H_{N_\uparrow,N_\downarrow},\; 0\le n_{b,i}\le n_{\mathrm{ph,max}}\}.
\]

\subsubsection{7.3.1 Binary index map}\label{binary-index-map}

For binary boson encoding, one may order basis states by the pair \[
(f,\mathbf n_b),
\] where \(f\) is the fermionic bit pattern and
\(\mathbf n_b=(n_{b,0},\dots,n_{b,L-1})\) is the vector of local phonon
occupations, each encoded in binary.

\subsubsection{7.3.2 Unary index map}\label{unary-index-map}

For unary boson encoding, the same sector is indexed by \[
(f,\mathbf e_{n_{b,0}},\dots,\mathbf e_{n_{b,L-1}}),
\] where \(\mathbf e_n\) denotes the unary one-hot representation of
occupation \(n\).

\subsection{7.4 Exact HH matrix
elements}\label{exact-hh-matrix-elements}

\subsubsection{7.4.1 Diagonal elements}\label{diagonal-elements}

For basis state \(|f\rangle\otimes|\mathbf n_b\rangle\), the diagonal
contribution is \[
U\sum_i n_{i\uparrow}(f)n_{i\downarrow}(f)
+
\sum_i v_i\,n_i(f)
+
\omega_0\sum_i \left(n_{b,i}+\frac{1}{2}\right).
\]

\subsubsection{7.4.2 Hopping matrix
elements}\label{hopping-matrix-elements}

If \(|f'\rangle\) differs from \(|f\rangle\) by a single allowed
nearest-neighbor hop with the same phonon occupations, then \[
\langle f',\mathbf n_b|\hat H_t|f,\mathbf n_b\rangle
=
-t\,(-1)^{\chi(f;p,q)},
\] where \(\chi(f;p,q)\) is the Jordan-Wigner parity count between the
two fermion modes \(p\) and \(q\).

\subsubsection{7.4.3 Electron-phonon matrix
elements}\label{electron-phonon-matrix-elements}

At fixed fermion configuration, the electron-phonon term changes the
local phonon number by one. For site \(i\), \[
\langle n_{b,i}+1|\hat x_i|n_{b,i}\rangle = \sqrt{n_{b,i}+1},
\qquad
\langle n_{b,i}-1|\hat x_i|n_{b,i}\rangle = \sqrt{n_{b,i}}.
\] Hence the off-diagonal electron-phonon matrix elements are
proportional to \[
g\,(n_i(f)-\bar n)\sqrt{n_{b,i}+1}
\quad\text{or}\quad
g\,(n_i(f)-\bar n)\sqrt{n_{b,i}},
\] according to whether the phonon occupation increases or decreases.

\section{8. Statevector Action, Expectation, and Exponential
Primitives}\label{statevector-action-expectation-and-exponential-primitives}

\subsection{8.1 Basis-state primitive}\label{basis-state-primitive}

A statevector on \(N_q\) qubits is written as \[
|\psi\rangle = \sum_{k=0}^{2^{N_q}-1} a_k |k\rangle
= \sum_{b\in\{0,1\}^{N_q}} a_b |b\rangle.
\]

\subsection{8.2 Explicit Pauli-word
action}\label{explicit-pauli-word-action}

Let \(P\) be a Pauli word. Then its action on a computational basis
state can be written in the form \[
P|b\rangle = \omega_P(b)\,|b\oplus \chi_P\rangle,
\] where \(\chi_P\) records which qubits are flipped by \(X\) or \(Y\),
and \(\omega_P(b)\) is the phase determined by the \(Y\) and \(Z\)
letters acting on the bitstring \(b\).

\subsection{8.3 Expectation values}\label{expectation-values}

If \[
\hat H = \sum_\alpha h_\alpha P_\alpha,
\] then the expectation value is \[
E(\psi)=\langle \psi|\hat H|\psi\rangle = \sum_\alpha h_\alpha \langle \psi|P_\alpha|\psi\rangle.
\]

\subsection{8.4 Pauli rotations}\label{pauli-rotations}

For any Hermitian Pauli word with \(P^2=I\), \[
e^{-i\theta P} = \cos\theta\,I - i\sin\theta\,P.
\] This identity is the primitive behind fast statevector application of
Pauli exponentials.

\subsection{8.5 Exponential of a Pauli
polynomial}\label{exponential-of-a-pauli-polynomial}

For a Pauli polynomial \[
A = \sum_\alpha a_\alpha P_\alpha,
\] one has \[
U_A(\theta)=e^{-i\theta A}.
\] If the \(P_\alpha\) commute pairwise, then \[
U_A(\theta)=\prod_\alpha e^{-i\theta a_\alpha P_\alpha}.
\] If they do not commute, one may use either exact exponentiation on
the working subspace or an ordered product formula.

\subsection{8.6 Exact sector energy
target}\label{exact-sector-energy-target}

The natural exact comparison target in a fixed fermion sector is \[
E_{\mathrm{exact,sector}} = \min_{\psi\in\mathcal H_{\mathrm{sector}}} \langle \psi|\hat H|\psi\rangle.
\] This is distinct from the unrestricted full-Hilbert minimum when a
sector constraint is imposed.

For HH ADAPT, distinguish the \emph{working} phonon cutoff used by the
variational Hilbert space from the \emph{external exact-evaluation}
cutoff used for documentation. If the ADAPT run uses
\(N_b^{\mathrm{work}}\) and the exact benchmark uses
\(N_b^{\mathrm{eval}}\), write \[
\Delta E_{\mathrm{same}}=
E_{\mathrm{ADAPT}}^{(N_b^{\mathrm{work}})}-E_{\mathrm{exact}}^{(N_b^{\mathrm{work}})},
\qquad
\Delta E_{\mathrm{ext}}=
E_{\mathrm{ADAPT}}^{(N_b^{\mathrm{work}})}-E_{\mathrm{exact}}^{(N_b^{\mathrm{eval}})}.
\] The same-cutoff error \(|\Delta E_{\mathrm{same}}|\) measures ADAPT
convergence inside the truncated working problem. The
documentation-default physical benchmark is
\(|\Delta E_{\mathrm{ext}}|\) with \(N_b^{\mathrm{eval}}=3\) for the
canonical L=2 HH point unless a report explicitly states a different
convergence cutoff.

\section{9. Fixed VQE Ansatz Families}\label{fixed-vqe-ansatz-families}

This section records the main symbolic manifold families without tying
them to any implementation names.

\subsection{9.1 Hubbard layerwise
ansatz}\label{hubbard-layerwise-ansatz}

For the static Hubbard model, write \[
\hat H_{\mathrm{Hub}} = \hat H_t + \hat H_U + \hat H_v.
\] A layerwise ansatz uses one parameter per physical group per layer:
\[
U_{\mathrm{Hub}}^{(\ell)}
=
\exp(-i\theta_t^{(\ell)}\hat H_t)
\exp(-i\theta_U^{(\ell)}\hat H_U)
\exp(-i\theta_v^{(\ell)}\hat H_v),
\] with the understanding that each grouped generator may itself be
split into Pauli exponentials internally while sharing the same scalar
parameter within the group.

\subsection{9.2 Hubbard excitation-based
surface}\label{hubbard-excitation-based-surface}

Relative to a Hartree-Fock occupied/virtual partition, single- and
double-excitation generators take the form \[
G_{i\to a}^{(\sigma)} = i\left(\hat c_{a\sigma}^\dagger \hat c_{i\sigma} - \hat c_{i\sigma}^\dagger \hat c_{a\sigma}\right),
\] and \[
G_{ij\to ab}^{(\sigma,\sigma')} = i\left(\hat c_{a\sigma}^\dagger \hat c_{b\sigma'}^\dagger \hat c_{j\sigma'}\hat c_{i\sigma} - \text{h.c.}\right).
\] A layered excitation manifold is then \[
U_{\mathrm{exc}}(\theta)=\prod_r e^{-i\theta_r G_r}.
\]

\subsection{9.3 HH layerwise ansatz}\label{hh-layerwise-ansatz}

For the HH model, define grouped generators \[
G_t,\;G_U,\;G_v,\;G_{\mathrm{ph}},\;G_g,\;G_{\mathrm{drive}}.
\] One layer is \[
U_{\mathrm{HH,lw}}^{(\ell)}
=
\exp(-i\theta_{t,\ell}G_t)
\exp(-i\theta_{U,\ell}G_U)
\exp(-i\theta_{v,\ell}G_v)
\exp(-i\theta_{\mathrm{ph},\ell}G_{\mathrm{ph}})
\exp(-i\theta_{g,\ell}G_g)
\exp(-i\theta_{\mathrm{drive},\ell}G_{\mathrm{drive}}),
\] with inactive groups omitted.

\subsection{9.4 HH physical-termwise
ansatz}\label{hh-physical-termwise-ansatz}

A physical-termwise layer sits between the grouped and fully split
extremes. Let \[
\mathcal A_{\mathrm{phys}} = \{A_m\}
\] be the collection of unsplit physical HH generators such as bond
hopping terms, onsite interactions, local phonon energies, local
couplings, and local drive terms. Then \[
U_{\mathrm{HH,phys}}^{(\ell)} = \prod_{A_m\in\mathcal A_{\mathrm{phys}}} e^{-i\phi_{m,\ell}A_m}.
\] Because the generators remain grouped at the physical-operator level,
this fixed-family surface stays aligned with the same conserved-sector
story as the HH layerwise surface. In this symbolic HH manuscript, the
canonical fixed-family surface is intentionally restricted to the
sector-preserving HH layerwise and HH physical-termwise families.

\subsection{9.5 Reference-state and ansatz pairing
principle}\label{reference-state-and-ansatz-pairing-principle}

The reference state and variational family should be paired by three
criteria:

\begin{enumerate}
\def\labelenumi{\arabic{enumi}.}
\tightlist
\item
  the conserved sector they preserve,
\item
  the expressive scale needed by the target Hamiltonian,
\item
  the conditioning and optimization burden induced by the chosen
  parameterization.
\end{enumerate}

\section{10. Polaronic Operator
Families}\label{polaronic-operator-families}

\subsection{10.1 Primitive polaronic
ingredients}\label{primitive-polaronic-ingredients}

\subsubsection{10.1.1 Shifted density}\label{shifted-density}

If the total electron number is \(N_e\), define the mean density \[
\bar n = \frac{N_e}{L},
\] and the shifted density \[
\tilde n_i = \hat n_i - \bar n I.
\] At half filling \(N_e=L\), this reduces to \[
\tilde n_i = -\frac{1}{2}(Z_{p(i,\uparrow)}+Z_{p(i,\downarrow)}).
\]

\subsubsection{\texorpdfstring{10.1.2 Phonon primitives \(P_i\) and
\(x_i\)}{10.1.2 Phonon primitives P\_i and x\_i}}\label{phonon-primitives-p_i-and-x_i}

Define \[
P_i=i(\hat b_i^\dagger-\hat b_i),
\qquad
x_i=\hat b_i+\hat b_i^\dagger.
\]

\subsubsection{10.1.3 Local squeeze
primitive}\label{local-squeeze-primitive}

A natural local squeeze generator is \[
S_i = i\left((\hat b_i^\dagger)^2-\hat b_i^2\right).
\]

\subsubsection{10.1.4 Doublon primitive}\label{doublon-primitive}

The doublon projector is \[
D_i = \hat n_{i\uparrow}\hat n_{i\downarrow}.
\]

\subsubsection{10.1.5 Even hopping channel}\label{even-hopping-channel}

A bond-even hopping operator is \[
T_{ij}^{(+)} = \sum_\sigma\left(\hat c_{i\sigma}^\dagger\hat c_{j\sigma}+\hat c_{j\sigma}^\dagger\hat c_{i\sigma}\right).
\]

\subsubsection{10.1.6 Odd current channel}\label{odd-current-channel}

A bond-odd current operator is \[
J_{ij}^{(-)} = i\sum_\sigma\left(\hat c_{i\sigma}^\dagger\hat c_{j\sigma}-\hat c_{j\sigma}^\dagger\hat c_{i\sigma}\right).
\]

\subsection{10.2 Representative dressed
channels}\label{representative-dressed-channels}

The exact family boundaries can vary, but the symbolic equivalents of
the common channels can be written as follows.

\subsubsection{10.2.1 Conditional
displacement}\label{conditional-displacement}

\[
G_i^{(\mathrm{cd})} = \tilde n_i P_i.
\]

\subsubsection{10.2.2 Doublon dressing}\label{doublon-dressing}

\[
G_i^{(\mathrm{dd})} = D_i P_i.
\]

\subsubsection{10.2.3 Dressed hopping drag}\label{dressed-hopping-drag}

\[
G_{ij}^{(\mathrm{hd})} = T_{ij}^{(+)}(P_i-P_j).
\]

\subsubsection{10.2.4 Odd-channel drag}\label{odd-channel-drag}

\[
G_{ij}^{(\mathrm{od})} = J_{ij}^{(-)}(x_i-x_j).
\]

\subsubsection{10.2.5 Second-order even
channel}\label{second-order-even-channel}

\[
G_{ij}^{(2)} = T_{ij}^{(+)}(x_i-x_j)^2.
\]

\subsubsection{10.2.6 Third-order odd current
channel}\label{third-order-odd-current-channel}

\[
G_{ij}^{(3)} = J_{ij}^{(-)}(x_i-x_j)^3.
\]

\subsubsection{10.2.7 Fourth-order even hopping
channel}\label{fourth-order-even-hopping-channel}

\[
G_{ij}^{(4)} = T_{ij}^{(+)}(x_i-x_j)^4.
\]

\subsubsection{10.2.8 Bond-conditioned displacement and squeeze
channels}\label{bond-conditioned-displacement-and-squeeze-channels}

With a bond-sensitive scalar observable \(B_{ij}\), one may write \[
G_{ij}^{(\mathrm{bd})}=B_{ij}P_i,
\qquad
G_{ij}^{(\mathrm{bs})}=B_{ij}S_i.
\] A simple choice is \(B_{ij}=\hat n_i-\hat n_j\).

\subsubsection{10.2.9 Local squeeze
channels}\label{local-squeeze-channels}

\[
G_i^{(\mathrm{sq})}=S_i,
\qquad
\tilde G_i^{(\mathrm{sq})}=\tilde n_i S_i.
\]

\subsubsection{10.2.10 Extended cloud
channels}\label{extended-cloud-channels}

A symbolic extended cloud family is \[
G_i^{(\mathrm{cloud})}=\tilde n_i\sum_{r\in\mathcal N(i)} \alpha_{ir} P_r,
\] with coefficients \(\alpha_{ir}\) controlling cloud shape.

\subsubsection{10.2.11 Doublon-translation and doublon-squeeze
channels}\label{doublon-translation-and-doublon-squeeze-channels}

\[
G_{ij}^{(\mathrm{DT})}=D_i T_{ij}^{(+)},
\qquad
G_i^{(\mathrm{DS})}=D_i S_i.
\]

\subsection{10.3 Pool-family map}\label{pool-family-map}

A useful abstract partition is \[
\mathcal G = \mathcal G_{\mathrm{local}} \cup \mathcal G_{\mathrm{bond}} \cup \mathcal G_{\mathrm{squeeze}} \cup \mathcal G_{\mathrm{cloud}} \cup \mathcal G_{\mathrm{doublon}}.
\] These families can be opened gradually rather than all at once.

\subsubsection{10.3.1 Compact displacement/squeeze macro
families}\label{compact-displacementsqueeze-macro-families}

Two especially compact subfamilies are \[
\mathcal G_{\mathrm{VLF}} = \{\tilde n_i P_i\}_i,
\qquad
\mathcal G_{\mathrm{SQ}} = \{S_i,\tilde n_i S_i\}_i.
\] They act as coarse polaronic and squeeze envelopes.

\subsection{10.4 Paper-I Hamiltonians and primitive SNAKE generator
menus}\label{paper-i-hamiltonians-and-primitive-snake-generator-menus}

This section is a Table-I dictionary for the fixed-accuracy benchmark
family. The Hamiltonians are written before Jordan--Wigner,
boson-register, or qubit-truncation encoding. SNAKE builds an adaptive
ordered product
\(U_{\rm SNAKE}(\theta)=\prod_{\ell=1}^{d}\exp(-i\theta_\ell G_\ell)\);
``available'' means the primitive generator records offered to the route
for that Hamiltonian, while ``selected'' means the records that the
current corrected SNAKE sidecars actually show, or the active primitive
menu when the row is still running or lacks a completed label-level
certificate.

Primitive generator vocabulary. Electronic single excitations are
\(G_{i\to a,\sigma}=i(c^\dagger_{a\sigma}c_{i\sigma}-c^\dagger_{i\sigma}c_{a\sigma})\);
electronic double excitations are
\(G_{ij\to ab,\sigma\sigma'}=i(c^\dagger_{a\sigma}c^\dagger_{b\sigma'}c_{j\sigma'}c_{i\sigma}-{\rm h.c.})\).
Spinful hopping and current primitives are
\(T_{ij}^{(\sigma)}=c^\dagger_{i\sigma}c_{j\sigma}+c^\dagger_{j\sigma}c_{i\sigma}\)
and
\(J_{ij}^{(\sigma)}=i(c^\dagger_{i\sigma}c_{j\sigma}-c^\dagger_{j\sigma}c_{i\sigma})\),
with density-assisted versions \(n_{k\bar\sigma}T_{ij}^{(\sigma)}\) and
\(n_{k\bar\sigma}J_{ij}^{(\sigma)}\). Pair primitives are
\(T_{ij}^{\rm pair}=c^\dagger_{i\uparrow}c^\dagger_{i\downarrow}c_{j\downarrow}c_{j\uparrow}+c^\dagger_{j\uparrow}c^\dagger_{j\downarrow}c_{i\downarrow}c_{i\uparrow}\)
and
\(J_{ij}^{\rm pair}=i(c^\dagger_{i\uparrow}c^\dagger_{i\downarrow}c_{j\downarrow}c_{j\uparrow}-c^\dagger_{j\uparrow}c^\dagger_{j\downarrow}c_{i\downarrow}c_{i\uparrow})\).
Exchange primitives are
\(T_{ij}^{\rm ex}=c^\dagger_{i\uparrow}c_{j\uparrow}c^\dagger_{j\downarrow}c_{i\downarrow}+{\rm h.c.}\)
and
\(J_{ij}^{\rm ex}=i(c^\dagger_{i\uparrow}c_{j\uparrow}c^\dagger_{j\downarrow}c_{i\downarrow}-{\rm h.c.})\),
with three-site bridge hop/current records where present. Spinless
primitives use the same \(T_{ij},J_{ij}\) grammar without spin, plus
\(n_i\), \(n_in_j\), and Hamiltonian-term rotations. Bosonic primitives
are \(X_i=b_i+b_i^\dagger\), \(P_i=i(b_i^\dagger-b_i)\),
\(n_i=b_i^\dagger b_i\), \(S_i=i((b_i^\dagger)^2-b_i^2)\), \(X_i^2\),
\(P_i^2\), \(n_i^2\), \(X_iP_i+P_iX_i\),
\(b_i^\dagger b_j+b_j^\dagger b_i\),
\(i(b_i^\dagger b_j-b_j^\dagger b_i)\), density-weighted hop/current
records, pair hop/current records, \(X_iX_j\), \(P_iP_j\), and Kerr
rotations \(n_i(n_i-I)\). Spin-boson primitives add emitter flip
\(T_{\rm em}\), imbalance \(Z_{\rm em}\),
\(Y_{\rm em}=i(\sigma^- - \sigma^+)\), oscillator \(X_b,P_b,n_b\) and
quadratics, and products such as \(X_bT_{\rm em}\), \(P_bT_{\rm em}\),
\(X_bZ_{\rm em}\), \(P_bZ_{\rm em}\), \(n_bT_{\rm em}\), and
\(n_bZ_{\rm em}\). HH polaronic primitives lift the electronic
primitives to electron-phonon space and add local \(X_i,P_i,n_i,S_i\),
Lang--Firsov density-momentum \((n_i-\bar n I)P_i\), density-squeeze
\((n_i-\bar n I)S_i\), cloud/displacement/doublon/bond-conditioned
phonon channels, and products of lifted electronic excitations with
these phonon channels. Molecular-vibronic primitives are electronic
singles/doubles, the stretch momentum \(P_b\), and coupled records
\(O_eP_b\) where \(O_e\) is an active-space occupation, bond, pair, or
mixed electronic channel. Hamiltonian-block primitives are rotations
generated by the full Hamiltonian, grouped physical blocks, or typed
terms after normalization.

\begin{itemize}
\tightlist
\item
  \textbf{Hubbard.} The clean two-site spinful Hubbard Hamiltonian is
  \(H_{\rm Hub}=-t\sum_{\langle i,j\rangle,\sigma}(c^\dagger_{i\sigma}c_{j\sigma}+c^\dagger_{j\sigma}c_{i\sigma})+U\sum_i n_{i\uparrow}n_{i\downarrow}+\sum_{i,\sigma}v_i n_{i\sigma}\),
  with \(n_{i\sigma}=c^\dagger_{i\sigma}c_{i\sigma}\) and \(v_i=0\) in
  the displayed clean row. SNAKE available: electronic singles/doubles,
  \(T_{ij}^{(\sigma)}\), \(J_{ij}^{(\sigma)}\), density-assisted
  hop/current, pair hop/current, exchange/current, and Hubbard
  Hamiltonian-block rotations. SNAKE selected: electronic
  singles/doubles in both weak \(U/t=2\) and strong \(U/t=8\) lanes.
\item
  \textbf{Ionic Hubbard.}
  \(H_{\rm ion}=-t\sum_{\langle i,j\rangle,\sigma}(c^\dagger_{i\sigma}c_{j\sigma}+{\rm h.c.})+U\sum_i n_{i\uparrow}n_{i\downarrow}-\Delta\sum_{i,\sigma}(-1)^i n_{i\sigma}\).
  SNAKE available: electronic singles/doubles, nearest-neighbor
  \(T_{ij}^{(\sigma)}\) and \(J_{ij}^{(\sigma)}\), density-assisted
  hop/current, ionic potential rotations \((-1)^in_{i\sigma}\), and
  ionic-Hubbard Hamiltonian blocks. SNAKE selected: electronic
  singles/doubles together with nearest-neighbor hopping-current
  records.
\item
  \textbf{\(t\)-\(t'\) Hubbard.}
  \(H_{tt'}=-t\sum_{\langle i,j\rangle,\sigma}(c^\dagger_{i\sigma}c_{j\sigma}+{\rm h.c.})-t'\sum_{\langle\langle i,j\rangle\rangle,\sigma}(c^\dagger_{i\sigma}c_{j\sigma}+{\rm h.c.})+U\sum_i n_{i\uparrow}n_{i\downarrow}+\sum_{i,\sigma}v_i n_{i\sigma}\).
  SNAKE available: electronic singles/doubles, nearest- and
  next-neighbor \(T_{ij}^{(\sigma)}\), \(J_{ij}^{(\sigma)}\),
  density-assisted currents, pair/exchange records, and \(t\)-\(t'\)
  Hamiltonian blocks. SNAKE selected: electronic singles/doubles in the
  weak lane; electronic singles/doubles plus hopping and
  assisted-current records in the strong lane.
\item
  \textbf{Extended Hubbard.}
  \(H_{\rm ext}=-t\sum_{\langle i,j\rangle,\sigma}(c^\dagger_{i\sigma}c_{j\sigma}+{\rm h.c.})+U\sum_i n_{i\uparrow}n_{i\downarrow}+V\sum_{\langle i,j\rangle}n_i n_j+\sum_{i,\sigma}v_i n_{i\sigma}\),
  with \(n_i=n_{i\uparrow}+n_{i\downarrow}\). SNAKE available:
  electronic singles/doubles, \(T_{ij}^{(\sigma)}\),
  \(J_{ij}^{(\sigma)}\), density-density rotations \(n_in_j\),
  density-assisted hop/current, exchange/current, pair hop/current, and
  extended-Hubbard Hamiltonian blocks. SNAKE selected: electronic
  singles/doubles plus hopping, exchange, and pair-current records.
\item
  \textbf{Spinless \(t\)-\(V\).}
  \(H_{tV}=-t\sum_{\langle i,j\rangle}(c^\dagger_i c_j+c^\dagger_j c_i)+V\sum_{\langle i,j\rangle}n_i n_j+\sum_i v_i n_i\).
  SNAKE available: spinless \(T_{ij}\), \(J_{ij}\), number rotations
  \(n_i\), interaction rotations \(n_in_j\), and spinless Hamiltonian
  blocks. SNAKE selected: nearest-neighbor hopping/current quadrature
  records in both lanes.
\item
  \textbf{Bose--Hubbard.}
  \(H_{\rm BH}=-J\sum_{\langle i,j\rangle}(b_i^\dagger b_j+b_j^\dagger b_i)+\omega_0\sum_i n_i+\frac{U_b}{2}\sum_i n_i(n_i-I)+\sum_i \epsilon_i n_i\),
  with local Fock cutoff \(n_i\le n_b^{\max}\) applied only after this
  symbolic definition. SNAKE available: local
  \(X_i,P_i,n_i,X_i^2,P_i^2,n_i^2,S_i,X_iP_i+P_iX_i\), boson
  hop/current, density-weighted hop/current, pair hop/current, Kerr, and
  Bose--Hubbard Hamiltonian blocks. SNAKE selected: boson-current and
  density-current records in both weak \(U_b/J=2\) and strong
  \(U_b/J=6\) lanes.
\item
  \textbf{Harmonic/Kerr chain.}
  \(H_{\rm HKC}=\omega_0\sum_i n_i+\frac{K}{2}\sum_i n_i(n_i-I)-t_x\sum_{\langle i,j\rangle}X_iX_j+d_x\sum_i X_i\),
  with \(X_i=b_i+b_i^\dagger\). SNAKE available: local
  \(X_i,P_i,n_i,X_i^2,P_i^2,n_i^2,S_i,X_iP_i+P_iX_i\), \(X_iX_j\),
  \(P_iP_j\), boson hop/current, Kerr rotations, and harmonic/Kerr
  Hamiltonian blocks. SNAKE selected: before the current target miss,
  observed records include local momentum/number/squeeze-like rotations,
  \(X_iX_j\), boson current, and Hamiltonian-term rotations.
\item
  \textbf{Spin--boson.}
  \(H_{\rm SB}=-t\,T_{\rm em}+d_v Z_{\rm em}+\omega_0(n_b+\frac12 I)+u\,X_bT_{\rm em}+g_{\rm ep}X_bZ_{\rm em}\),
  where \(T_{\rm em}\) is the emitter flip, \(Z_{\rm em}\) is the signed
  emitter imbalance, \(X_b=b+b^\dagger\), and \(n_b=b^\dagger b\). SNAKE
  available: \(T_{\rm em}\), \(Z_{\rm em}\), \(Y_{\rm em}\), oscillator
  \(X_b,P_b,n_b\) and quadratic/squeeze records, longitudinal products
  \(X_bZ_{\rm em},P_bZ_{\rm em},n_bZ_{\rm em}\), transverse products
  \(X_bT_{\rm em},P_bT_{\rm em},n_bT_{\rm em}\), and spin-boson
  Hamiltonian blocks. SNAKE selected: current status cells are present,
  but completed label-level payloads are failed or running, so cite the
  active spin-boson primitive menu until a later sidecar narrows it.
\item
  \textbf{Hubbard--Holstein.}
  \(H_{\rm HH}=H_{\rm Hub}+\omega_0\sum_i b_i^\dagger b_i+g\sum_i(n_i-\bar n I)(b_i+b_i^\dagger)\),
  where \(\bar n=N_e/L\), \(H_{\rm Hub}\) is the spinful Hubbard block
  above, and the clean static rows have no time-dependent drive term.
  SNAKE available: lifted electronic singles/doubles, lifted spinful
  \(T_{ij}^{(\sigma)}\), \(J_{ij}^{(\sigma)}\), density-assisted
  currents, exchange and pair-current records, local phonon
  \(X_i,P_i,n_i,S_i\), Lang--Firsov \((n_i-\bar n I)P_i\),
  density-squeeze \((n_i-\bar n I)S_i\),
  cloud/displacement/doublon/bond-conditioned phonon channels, products
  of lifted electronic excitations with polaronic channels, and HH
  Hamiltonian blocks. SNAKE selected: the corrected weak row selects a
  mixed electronic-polaronic-bosonic surface containing lifted
  hopping/current, polaronic channels, lifted electronic
  singles/doubles, and boson quadratures; the strong row is not a
  reached current resource row and has only broad mixed-surface evidence
  before target failure.
\item
  \textbf{Molecular-vibronic H2.}
  \(H_{\rm H2vib}=H_{\rm el}(R_0)+\omega(b^\dagger b+\frac12 I)+g_{\rm ep}x_{\rm zpf}H'_R(R_0)(b+b^\dagger)\),
  with
  \(H_{\rm el}(R_0)=E_{\rm nuc}(R_0)+\sum_{pq}h_{pq}(R_0)c_p^\dagger c_q+\frac12\sum_{pqrs}g_{pqrs}(R_0)c_p^\dagger c_q^\dagger c_s c_r\),
  \(H'_R(R_0)\approx[H_{\rm el}(R_0+\Delta R)-H_{\rm el}(R_0-\Delta R)]/(2\Delta R_B)\)
  after displaced molecular-orbital gauge alignment, and \(\Delta R_B\)
  the Bohr-valued finite-difference step. SNAKE available: electronic
  singles/doubles, pure stretch momentum \(P_b\), coupled
  electronic-stretch momentum records \(O_eP_b\) with active-space
  occupation/bond/pair/mixed \(O_e\), and molecular-vibronic Hamiltonian
  blocks. SNAKE selected: weak SNAKE selects electronic singles/doubles
  plus coupled electronic-stretch momentum records; the corrected strong
  \(n_b^{\rm work}=3\), \(n_b^{\rm ref}=6\) row is running without a
  completed selected-family certificate, and the older suppressed
  \(n_b^{\rm work}=1\), \(n_b^{\rm ref}=4\) strong sidecar selected the
  same primitive types but is not the current strong-cutoff row.
\end{itemize}

\subsection{10.5 Paper-I benchmark primitive-generator
ledger}\label{paper-i-benchmark-primitive-generator-ledger}

Exact diagonalization, when present, has no ansatz. The benchmark rows
below use these rules: HEA VQE is the fixed circuit with \(R_y\) and
\(R_z\) rotations on each qubit, a nearest-neighbor CNOT chain, and a
final \(R_y,R_z\) layer; family-informed fixed VQE is a non-adaptive
policy-selected subset of the problem-local primitive menu; append-only
ADAPT greedily appends the largest raw commutator-gradient primitive
from that menu and refits by BFGS; TETRIS-ADAPT uses the same menu but
greedily batches gradient-ranked primitives with disjoint qubit support;
Pos-Geo-ADAPT uses the same menu with projected-Fubini--Study scoring
and best-position local refit; Qubit/QEB-ADAPT uses qubit-excitation
singles/doubles expanded into Pauli words; SNAKE is the route-native
adaptive row summarized in Sec. 10.4.

\begin{itemize}
\tightlist
\item
  \textbf{Hubbard.}

  \begin{itemize}
  \tightlist
  \item
    \textbf{SNAKE.} Available: electronic singles/doubles, spinful
    hop/current and assisted-current primitives, pair/exchange records,
    and Hubbard Hamiltonian blocks. Selected: electronic
    singles/doubles.
  \item
    \textbf{HEA VQE.} Available: \(R_y,R_z\) layers plus
    nearest-neighbor CNOT chain. Selected: family selected equals family
    available.
  \item
    \textbf{Family-informed fixed VQE.} Available: electronic
    singles/doubles. Selected: fixed electronic singles/doubles.
  \item
    \textbf{Append-only ADAPT.} Available: the Hubbard problem-local
    primitive menu used by SNAKE. Selected: row-specific
    largest-gradient append sequence from that menu.
  \item
    \textbf{TETRIS-ADAPT.} Available: the same Hubbard primitive menu.
    Selected: row-specific gradient-ranked disjoint-support batches.
  \item
    \textbf{Pos-Geo-ADAPT.} Available: the same Hubbard primitive menu.
    Selected: row-specific projected-metric, best-position sequence.
  \item
    \textbf{Qubit/QEB-ADAPT.} Available: qubit-excitation
    singles/doubles expanded into Pauli words. Selected: row-specific
    qubit-excitation sequence.
  \end{itemize}
\item
  \textbf{Ionic Hubbard.}

  \begin{itemize}
  \tightlist
  \item
    \textbf{SNAKE.} Available: electronic singles/doubles, spinful
    hop/current and assisted-current primitives, ionic potential
    rotations, pair/exchange records, and ionic-Hubbard blocks.
    Selected: electronic singles/doubles plus nearest-neighbor
    hopping-current records.
  \item
    \textbf{HEA VQE.} Available: \(R_y,R_z\) layers plus
    nearest-neighbor CNOT chain. Selected: family selected equals family
    available.
  \item
    \textbf{Family-informed fixed VQE.} Available: electronic
    singles/doubles. Selected: fixed electronic singles/doubles.
  \item
    \textbf{Append-only ADAPT.} Available: the ionic-Hubbard
    problem-local primitive menu. Selected: row-specific
    largest-gradient append sequence from that menu.
  \item
    \textbf{TETRIS-ADAPT.} Available: the same ionic-Hubbard primitive
    menu. Selected: row-specific gradient-ranked disjoint-support
    batches.
  \item
    \textbf{Pos-Geo-ADAPT.} Available: the same ionic-Hubbard primitive
    menu. Selected: row-specific projected-metric, best-position
    sequence.
  \item
    \textbf{Qubit/QEB-ADAPT.} Available: qubit-excitation
    singles/doubles expanded into Pauli words. Selected: row-specific
    qubit-excitation sequence.
  \end{itemize}
\item
  \textbf{\(t\)-\(t'\) Hubbard.}

  \begin{itemize}
  \tightlist
  \item
    \textbf{SNAKE.} Available: electronic singles/doubles, nearest- and
    next-neighbor hop/current records, assisted currents, pair/exchange
    records, and \(t\)-\(t'\) Hamiltonian blocks. Selected: electronic
    singles/doubles in the weak lane; electronic singles/doubles plus
    hopping and assisted-current records in the strong lane.
  \item
    \textbf{HEA VQE.} Available: \(R_y,R_z\) layers plus
    nearest-neighbor CNOT chain. Selected: family selected equals family
    available.
  \item
    \textbf{Family-informed fixed VQE.} Available: electronic
    singles/doubles. Selected: fixed electronic singles/doubles.
  \item
    \textbf{Append-only ADAPT.} Available: the \(t\)-\(t'\)
    problem-local primitive menu. Selected: row-specific
    largest-gradient append sequence from that menu.
  \item
    \textbf{TETRIS-ADAPT.} Available: the same \(t\)-\(t'\) primitive
    menu. Selected: row-specific gradient-ranked disjoint-support
    batches.
  \item
    \textbf{Pos-Geo-ADAPT.} Available: the same \(t\)-\(t'\) primitive
    menu. Selected: row-specific projected-metric, best-position
    sequence.
  \item
    \textbf{Qubit/QEB-ADAPT.} Available: qubit-excitation
    singles/doubles expanded into Pauli words. Selected: row-specific
    qubit-excitation sequence.
  \end{itemize}
\item
  \textbf{Extended Hubbard.}

  \begin{itemize}
  \tightlist
  \item
    \textbf{SNAKE.} Available: electronic singles/doubles, hop/current
    and assisted-current primitives, density-density rotations,
    exchange/current, pair hop/current, and extended-Hubbard blocks.
    Selected: electronic singles/doubles plus hopping, exchange, and
    pair-current records.
  \item
    \textbf{HEA VQE.} Available: \(R_y,R_z\) layers plus
    nearest-neighbor CNOT chain. Selected: family selected equals family
    available.
  \item
    \textbf{Family-informed fixed VQE.} Available: electronic
    singles/doubles. Selected: fixed electronic singles/doubles.
  \item
    \textbf{Append-only ADAPT.} Available: the extended-Hubbard
    problem-local primitive menu. Selected: row-specific
    largest-gradient append sequence from that menu.
  \item
    \textbf{TETRIS-ADAPT.} Available: the same extended-Hubbard
    primitive menu. Selected: row-specific gradient-ranked
    disjoint-support batches.
  \item
    \textbf{Pos-Geo-ADAPT.} Available: the same extended-Hubbard
    primitive menu. Selected: row-specific projected-metric,
    best-position sequence.
  \item
    \textbf{Qubit/QEB-ADAPT.} Available: qubit-excitation
    singles/doubles expanded into Pauli words. Selected: row-specific
    qubit-excitation sequence.
  \end{itemize}
\item
  \textbf{Spinless \(t\)-\(V\).}

  \begin{itemize}
  \tightlist
  \item
    \textbf{SNAKE.} Available: spinless hop/current, number,
    density-density, and Hamiltonian-block rotations. Selected:
    nearest-neighbor hopping/current quadrature records.
  \item
    \textbf{HEA VQE.} Available: \(R_y,R_z\) layers plus
    nearest-neighbor CNOT chain. Selected: family selected equals family
    available.
  \item
    \textbf{Family-informed fixed VQE.} Available: spinless hop/current,
    density-density, and Hamiltonian-block rotations approved by policy.
    Selected: fixed number-conserving hopping-style subset.
  \item
    \textbf{Append-only ADAPT.} Available: the spinless \(t\)-\(V\)
    problem-local primitive menu. Selected: row-specific
    largest-gradient append sequence from that menu.
  \item
    \textbf{TETRIS-ADAPT.} Available: the same spinless primitive menu.
    Selected: row-specific gradient-ranked disjoint-support batches.
  \item
    \textbf{Pos-Geo-ADAPT.} Available: the same spinless primitive menu.
    Selected: row-specific projected-metric, best-position sequence.
  \item
    \textbf{Qubit/QEB-ADAPT.} Available: qubit-excitation
    singles/doubles expanded into Pauli words. Selected: row-specific
    qubit-excitation sequence.
  \end{itemize}
\item
  \textbf{Bose--Hubbard.}

  \begin{itemize}
  \tightlist
  \item
    \textbf{SNAKE.} Available: local oscillator
    quadratures/numbers/quadratics, boson hop/current, density-weighted
    hop/current, pair records, Kerr, and Bose--Hubbard blocks. Selected:
    boson-current and density-current records.
  \item
    \textbf{HEA VQE.} Available: \(R_y,R_z\) layers plus
    nearest-neighbor CNOT chain. Selected: family selected equals family
    available.
  \item
    \textbf{Family-informed fixed VQE.} Available: local oscillator
    quadratures/numbers/quadratics, boson hop/current, density-weighted
    records, Kerr, and Hamiltonian-block rotations approved by policy.
    Selected: fixed bosonic quadrature-or-Hamiltonian-block subset.
  \item
    \textbf{Append-only ADAPT.} Available: the Bose--Hubbard
    problem-local primitive menu. Selected: row-specific
    largest-gradient append sequence from that menu.
  \item
    \textbf{TETRIS-ADAPT.} Available: the same Bose--Hubbard primitive
    menu. Selected: row-specific gradient-ranked disjoint-support
    batches.
  \item
    \textbf{Pos-Geo-ADAPT.} Available: the same Bose--Hubbard primitive
    menu. Selected: row-specific projected-metric, best-position
    sequence.
  \item
    \textbf{Qubit/QEB-ADAPT.} Available: qubit-excitation
    singles/doubles expanded into Pauli words. Selected: row-specific
    qubit-excitation sequence.
  \end{itemize}
\item
  \textbf{Harmonic/Kerr chain.}

  \begin{itemize}
  \tightlist
  \item
    \textbf{SNAKE.} Available: local oscillator
    quadratures/numbers/quadratics, \(X_iX_j\), \(P_iP_j\), boson
    hop/current, Kerr, and harmonic/Kerr blocks. Selected: observed
    local momentum/number/squeeze-like rotations, \(X_iX_j\), boson
    current, and Hamiltonian-term rotations before target miss.
  \item
    \textbf{HEA VQE.} Available: \(R_y,R_z\) layers plus
    nearest-neighbor CNOT chain. Selected: family selected equals family
    available.
  \item
    \textbf{Family-informed fixed VQE.} Available: oscillator
    quadratures/numbers/quadratics, \(X_iX_j\), \(P_iP_j\), boson
    currents, Kerr, and Hamiltonian-block rotations approved by policy.
    Selected: fixed bosonic quadrature-or-Hamiltonian-block subset.
  \item
    \textbf{Append-only ADAPT.} Available: the harmonic/Kerr
    problem-local primitive menu. Selected: row-specific
    largest-gradient append sequence from that menu.
  \item
    \textbf{TETRIS-ADAPT.} Available: the same harmonic/Kerr primitive
    menu. Selected: row-specific gradient-ranked disjoint-support
    batches.
  \item
    \textbf{Pos-Geo-ADAPT.} Available: the same harmonic/Kerr primitive
    menu. Selected: row-specific projected-metric, best-position
    sequence.
  \item
    \textbf{Qubit/QEB-ADAPT.} Available: qubit-excitation
    singles/doubles expanded into Pauli words. Selected: row-specific
    qubit-excitation sequence.
  \end{itemize}
\item
  \textbf{Spin--boson.}

  \begin{itemize}
  \tightlist
  \item
    \textbf{SNAKE.} Available: emitter flip/imbalance/\(Y\), oscillator
    quadratures/numbers/quadratics, longitudinal and transverse
    emitter-oscillator products, and spin-boson blocks. Selected: active
    spin-boson primitive menu until a completed label-level sidecar
    narrows it.
  \item
    \textbf{HEA VQE.} Available: \(R_y,R_z\) layers plus
    nearest-neighbor CNOT chain. Selected: family selected equals family
    available.
  \item
    \textbf{Family-informed fixed VQE.} Available: emitter, oscillator,
    and emitter-oscillator product primitives approved by policy.
    Selected: fixed spin-boson displacement/product subset.
  \item
    \textbf{Append-only ADAPT.} Available: the spin-boson problem-local
    primitive menu. Selected: row-specific largest-gradient append
    sequence from that menu.
  \item
    \textbf{TETRIS-ADAPT.} Available: the same spin-boson primitive
    menu. Selected: row-specific gradient-ranked disjoint-support
    batches.
  \item
    \textbf{Pos-Geo-ADAPT.} Available: the same spin-boson primitive
    menu. Selected: row-specific projected-metric, best-position
    sequence.
  \item
    \textbf{Qubit/QEB-ADAPT.} Available: qubit-excitation
    singles/doubles expanded into Pauli words. Selected: row-specific
    qubit-excitation sequence.
  \end{itemize}
\item
  \textbf{Hubbard--Holstein.}

  \begin{itemize}
  \tightlist
  \item
    \textbf{SNAKE.} Available: lifted electronic singles/doubles, lifted
    spinful hop/current primitives, local phonon quadratures/numbers,
    Lang--Firsov and density-squeeze records,
    cloud/displacement/doublon/bond-conditioned phonon channels,
    electronic-polaronic products, and HH blocks. Selected: weak row
    uses a mixed lifted-electronic, polaronic, and bosonic surface;
    strong row has broad mixed-surface evidence before target failure.
  \item
    \textbf{HEA VQE.} Available: \(R_y,R_z\) layers plus
    nearest-neighbor CNOT chain. Selected: family selected equals family
    available.
  \item
    \textbf{Family-informed fixed VQE.} Available: lifted electronic
    singles/doubles, Lang--Firsov/density-squeeze, oscillator
    quadratures, polaronic channels, and Hamiltonian-block rotations
    approved by policy. Selected: fixed lifted-electronic plus
    polaronic/quadrature subset.
  \item
    \textbf{Append-only ADAPT.} Available: the HH problem-local
    primitive menu. Selected: row-specific largest-gradient append
    sequence from that menu.
  \item
    \textbf{TETRIS-ADAPT.} Available: the same HH primitive menu.
    Selected: row-specific gradient-ranked disjoint-support batches.
  \item
    \textbf{Pos-Geo-ADAPT.} Available: the same HH primitive menu.
    Selected: row-specific projected-metric, best-position sequence.
  \item
    \textbf{Qubit/QEB-ADAPT.} Available: qubit-excitation
    singles/doubles expanded into Pauli words. Selected: row-specific
    qubit-excitation sequence.
  \end{itemize}
\item
  \textbf{Molecular-vibronic H2.}

  \begin{itemize}
  \tightlist
  \item
    \textbf{SNAKE.} Available: electronic singles/doubles, pure stretch
    momentum, coupled electronic-stretch momentum \(O_eP_b\), and
    molecular-vibronic blocks. Selected: weak row uses electronic
    singles/doubles plus coupled electronic-stretch momentum records;
    corrected strong row is still running without a completed
    selected-family certificate.
  \item
    \textbf{HEA VQE.} Available: \(R_y,R_z\) layers plus
    nearest-neighbor CNOT chain. Selected: family selected equals family
    available.
  \item
    \textbf{Family-informed fixed VQE.} Available: electronic
    singles/doubles, pure stretch momentum, coupled \(O_eP_b\), and
    molecular-vibronic Hamiltonian blocks approved by policy. Selected:
    fixed molecular-vibronic electronic-plus-coupled-stretch subset.
  \item
    \textbf{Append-only ADAPT.} Available: the molecular-vibronic
    problem-local primitive menu. Selected: row-specific
    largest-gradient append sequence from that menu.
  \item
    \textbf{TETRIS-ADAPT.} Available: the same molecular-vibronic
    primitive menu. Selected: row-specific gradient-ranked
    disjoint-support batches.
  \item
    \textbf{Pos-Geo-ADAPT.} Available: the same molecular-vibronic
    primitive menu. Selected: row-specific projected-metric,
    best-position sequence.
  \item
    \textbf{Qubit/QEB-ADAPT.} Available: qubit-excitation
    singles/doubles expanded into Pauli words. Selected: row-specific
    qubit-excitation sequence.
  \end{itemize}
\end{itemize}

\section{11. Adaptive Selection and Staged
Continuation}\label{adaptive-selection-and-staged-continuation}

\subsection{11.1 Cumulative phase construction for adaptive
selection}\label{cumulative-phase-construction-for-adaptive-selection}

If a shortlisted macro generator \(m\) admits a split family
\(\mathcal C_{\mathrm{split}}(m)=\{m_1,\dots,m_{K_m}\}\), then each
split child inherits the same position \(p\) and therefore defines a
child record \(r_j=(m_j,p)\). For simplicity we drop the \(j\), and
assume each \(m\) below is a unique generator.

Record-first view: the primary object is a candidate-position record
\(r:=(m,p)\), with projections \(\pi_m(m,p):=m\) and \(\pi_p(m,p):=p\).
For each phase \(k \in \{1,2,3\}\), the one-way selector flow is \[
\mathcal R_k\longrightarrow\mathcal S_k\longrightarrow r_k^\star,
\] with \[
\mathcal R_k:=\{(m,p):m\in\mathcal G_k,\ p\in\mathcal P_m^{k}\},
\qquad
\mathcal S_k:=
\{(m,p)\in\mathcal R_k:S_k(m,p)\ge\tau_k\}
\quad\text{or}\quad
\operatorname{Top}_{N_k}\!\left(S_k;\mathcal R_k\right),
\qquad
r_k^\star\in\mathcal S_k.
\] where \(S_k\) denotes the phase-\(k\) score used for shortlist
construction. Generator image and position image:
\(\mathcal G^{(k)}:=\pi_m(\mathcal S_k)\) and
\(\mathcal P_m^{(k)}:=\{p:(m,p)\in\mathcal S_k\}\). Record inheritance:
when later phases inherit the same candidate-position records, the clean
handoff is \(\mathcal R_{k+1}:=\mathcal S_k\). Hence the cumulative
record chain is
\(\mathcal R_3 \subseteq \mathcal R_2 \subseteq \mathcal R_1\), the
cumulative generator chain is
\(\mathcal G^{(3)} \subseteq \mathcal G^{(2)} \subseteq \mathcal G^{(1)} \subseteq \mathcal G\),
the first chain follows from
\(\mathcal R_{k+1}:=\mathcal S_k\subseteq\mathcal R_k\), and the
retained position images inherit analogously via the corresponding
record shortlists \(\mathcal S_1\to\mathcal S_2\to\mathcal S_3\).
Post-shortlist admission, batching, pruning, and beam continuation are
then defined in terms of the final Phase-3 surface.

Seed broadening: in addition, we use a smooth broadening of the initial
candidate record \(\mathcal R\) as expected generator utility saturates,
anti-comensurate to scaffold depth, where the initial seeds are
low-order generators and late-stage seeds are higher-order excitation
and coupling terms.

Chronology: thus the presentation chronologically accords with the
algorithmic implementation: Phase 1 \(\to\) Phase 2 \(\to\) Phase 3,
with each to-arrow filtering its preceding phase.

\subsection{11.2 Phase 1}\label{phase-1}

\subsubsection{11.2.1 Core signal, append position, and one-coordinate
local
model}\label{core-signal-append-position-and-one-coordinate-local-model}

Core record, scaffold, and append slot: \(r=(m,p)\),
\(\mathcal O=(a_1,\dots,a_n)\), \(m\in\mathcal G\), and
\(p\in\mathcal P_m^{1}=\{0,\dots,n\}\). Here \(m\) is a candidate
generator already admitted by the generator gate (so
\(m\in\mathcal G_1\)), and \(p\) is the index of a slot in the
pre-insertion, ordered scaffold \(\mathcal O\). Window bias: note, the
insertion window near the terminal scaffold point is biased to be small
proportional to expected scaffold utility remaining.

Scaffold parameter space and energy map: WLOG consider \(\mathcal O\),
s.t. each \(a\) is assigned one scalar \(\theta (\mod 2 \pi)\), then
define \(\Theta_{\mathcal O}:=\mathbb R^{|\mathcal O|}\),
\(E(\mathcal O,\theta):=\langle\psi(\mathcal O,\theta)|H|\psi(\mathcal O,\theta)\rangle\),
and \(\theta\in\Theta_{\mathcal O}\). Scaffold insertion and parameter
insertion: define
\(\mathcal O\oplus_p m:=(a_1,\dots,a_p,m,a_{p+1},\dots,a_n)\) and
\(\theta\oplus_p\vartheta:=(\theta_1,\dots,\theta_p,\vartheta,\theta_{p+1},\dots,\theta_n)\in\Theta_{\mathcal O\oplus_p m}\).

Write \(|\psi\rangle:=U_{>p}U_{\le p}|\psi_0\rangle\),
\(\widetilde A_r:=U_{>p}A_mU_{>p}^\dagger\),
\(U_r(\alpha):=e^{-i\alpha \widetilde A_r}\), and
\(|\psi_r(\alpha)\rangle:=U_r(\alpha)|\psi\rangle\). Let
\(Q_\psi:=I-|\psi\rangle\langle\psi|\),
\(|d_r\rangle:=\left.\partial_\alpha|\psi_r(\alpha)\rangle\right|_{\alpha=0}=-i\widetilde A_r|\psi\rangle\),
\(\tau_r:=Q_\psi|d_r\rangle\), \(g_r:=2\Re\langle \psi|H|d_r\rangle\),
\(\sigma_r:=\operatorname{SE}\!\left[\widehat g(r)\right]\ge 0\),
\(g_{\mathrm{lcb}}(r):=\max\{|g_r|-z_\alpha\sigma_r,0\}\), and
\(F_{\mathrm{raw}}(r):=\langle \tau_r,\tau_r\rangle+\varepsilon_F\). The
live broad Phase-1 one-coordinate trust-region proxy is \[
\Delta E_{1,\mathrm{TR}}(r)
:=
\max_{|\alpha|\le \rho/\sqrt{F_{\mathrm{raw}}(r)}}
\left[
g_{\mathrm{lcb}}(r)|\alpha|
-\frac12 \lambda_F F_{\mathrm{raw}}(r)\alpha^2
\right].
\] The strict one-coordinate insertion drop is retained as the exact
auxiliary comparator \[
\delta E_1(m,p\mid\mathcal O,\theta)
=
E(\mathcal O,\theta)
-\inf_{\vartheta\in\mathbb R}E\bigl(\mathcal O\oplus_p m,\theta\oplus_p\vartheta\bigr).
\] The infimum notation may be replaced by \(\min\) or
\(\operatorname{Min}\) whenever the minimizer exists. Thus
\(\Delta E_{1,\mathrm{TR}}(r)\) is the live broad-pool second-order
local insertion proxy, whereas \(\delta E_1(m,p\mid\mathcal O,\theta)\)
is the exact one-coordinate insertion drop.
\footnote{User interest: How does window parameterization width scale with cost?}

\subsubsection{11.2.2 Gates, burdens, and feature
ingredients}\label{gates-burdens-and-feature-ingredients}

Base hardware ingredients: each candidate-position pair carries the
proxy triple \(C_{\mathrm{2q}}(m,p)\), \(D(m,p)\), and
\(\Delta K(m,p)\), with fixed nonnegative weights
\(\lambda_{\mathrm{2q}},\lambda_d,\lambda_\theta\).\footnote{Admissibility gate: For each selector stage $k$, let $\Gamma_k:\mathcal G_k^{\mathrm{in}}\to\{0,1\}$ be the generator-level admissibility gate and define $\mathcal G_k^{\mathrm{adm}}:=\{m\in\mathcal G_k^{\mathrm{in}}:\Gamma_k(m)=1\}$. Record domain: the working candidate-record domain is built only from this gated generator pool, $\mathcal R_k:=\{(m,p):m\in\mathcal G_k^{\mathrm{adm}},\ p\in\mathcal P_m^{\mathrm{in},k}\}$. Selector action: all selectors below act on $\mathcal R_k$ itself, so admissibility is enforced by pre-score domain construction rather than post hoc score masking. Tie order: fix once and for all a total order $\prec$ on records (for example, lexicographic order on $(m,p)$). Tie-broken maximizer: for any score $S:\mathcal R_k\to\mathbb R$, define $\arg\max^{\prec}_{r\in\mathcal R_k} S(r):=\min_{\prec}\{r\in\mathcal R_k:S(r)=\max_{u\in\mathcal R_k}S(u)\}$, and similarly for $\arg\min^{\prec}$. Shortlist surface: define $\operatorname{Top}_K^{\prec}(S;\mathcal R_k)$ as the ordered list of the first $K$ records obtained by sorting $\mathcal R_k$ in descending score, with ties broken by $\prec$. Unadorned $\arg\max$, $\arg\min$, and $\operatorname{Top}_K$ below follow this convention.}
Here \(C_{\mathrm{2q}}\) is a two-qubit proxy cost, \(D\) is an expected
depth-induced by \(m\) (at \(p\)) proxy, and \(\Delta K\) is the
marginal parameter count. Positivity assumptions: assume throughout that
\(\lambda_{\mathrm{2q}}\ge 0\), \(C_{\mathrm{2q}}(m,p)\ge 0\), and
\(\lambda_\theta,\lambda_d,\Delta K(m,p),D(m,p)>0\) on admissible
records.
\footnote{Depth and parameter count are constants for our L=2 pareto lean runs.}

\subsubsection{11.2.3 Anticipated utility left and controller
schedule}\label{anticipated-utility-left-and-controller-schedule}

Let \(t\) denote the current selector step, let \(D_{\mathrm{left}}(t)\)
be the remaining controller depth budget, and let
\(\widehat N_{\mathrm{rem}}(t)\) be a heuristic forecast (expected
count) of useful remaining admissions. For gain-process statistics,
define \(\Delta_{\mathrm{adm}}(s):=E(\theta^{(s-1)})-E(\theta^{(s)})\),
\(\Delta_{\mathrm{adm}}^{+}(s):=\operatorname{asinh}\!\left(\frac{\Delta_{\mathrm{adm}}(s)}{\tau_\Delta}\right)\),
and let \(\mathcal I_t^+\) be the index set of past selector steps with
valid admissions used for EWMA/MAD statistics.

Let \(D_{\mathrm{allow}}\) denote the total allowed depth budget, let
\(\ell_t^{\mathrm{plat}}:=\max\{1,\lceil P\,D_{\mathrm{left}}(t)/D_{\mathrm{allow}}\rceil\}\)
and
\(\ell_t^{\mathrm{low}}:=\max\{1,\lceil L\,D_{\mathrm{left}}(t)/D_{\mathrm{allow}}\rceil\}\),
and define the relative-flattening and absolute weak-gain thresholds by
\(\tau_{\mathrm{plat}}^{+}(t):=\bigl(\eta_0+\eta_1(1-D_{\mathrm{left}}(t)/D_{\mathrm{allow}})\bigr)m_t\)
and
\(\tau_{\mathrm{low}}^{+}(t):=\operatorname{asinh}(\tau_{\mathrm{use},t}/\tau_\Delta)\).
The corresponding depth-adaptive streak statistics are
\(\operatorname{plateau\_streak}(t):=\frac{P}{\ell_t^{\mathrm{plat}}}\max\{q\in\{0,\dots,\ell_t^{\mathrm{plat}}\}:\Delta_{\mathrm{adm}}^{+}(t-j+1)\le \tau_{\mathrm{plat}}^{+}(t)\ \forall\,1\le j\le q\}\)
and
\(\operatorname{low\_streak}(t):=\frac{L}{\ell_t^{\mathrm{low}}}\max\{q\in\{0,\dots,\ell_t^{\mathrm{low}}\}:\Delta_{\mathrm{adm}}^{+}(t-j+1)\le \tau_{\mathrm{low}}^{+}(t)\ \forall\,1\le j\le q\}\).

For stagnation, gain, decay, and scale summaries, define
\(u_{\mathrm{stag}}(t):=\operatorname{clip}\!\left(\max\!\left(\frac{\mathrm{plateau\_streak}(t)}{P},\frac{\mathrm{low\_streak}(t)}{L}\right),0,1\right)\),
\(m_t:=\operatorname{EWMA}_{s\in\mathcal I_t^+}\!\bigl(\Delta_{\mathrm{adm}}^{+}(s)\bigr)\),
\(\rho_t:=\operatorname{clip}\!\left(\frac{\operatorname{EWMA}_{\mathrm{recent}}(\Delta_{\mathrm{adm}})}{\operatorname{EWMA}_{\mathrm{older}}(\Delta_{\mathrm{adm}})+\varepsilon},\rho_{\min},1\right)\),
\(\gamma_t:=[-\log\rho_t]_+\), and
\(s_t:=\max\!\left(s_{\min},\operatorname{MAD}_{s\in\mathcal I_t^+}\!\bigl(\Delta_{\mathrm{adm}}^{+}(s)\bigr)\right)\).

For the frontier statistic, let
\(s_{\mathrm{raw},1}^{(1)}(t)\ge s_{\mathrm{raw},2}^{(1)}(t)\ge 0\)
denote the first two order statistics of the current broad Phase-1 score
multiset \(\{S_1(r;t):r\in\mathcal R_1\}\) before any
\(\tau_1(t)\)-thresholding or \(N_1(t)\)-truncation, with the
conventions
\(s_{\mathrm{raw},1}^{(1)}(t)=s_{\mathrm{raw},2}^{(1)}(t)=0\) if
\(\mathcal R_1=\varnothing\) and \(s_{\mathrm{raw},2}^{(1)}(t)=0\) if
\(|\mathcal R_1|=1\). Define \[
u_{\mathrm{front}}(t)
:=
\frac{s_{\mathrm{raw},2}^{(1)}(t)+\varepsilon}{s_{\mathrm{raw},1}^{(1)}(t)+\varepsilon}.
\] Thus \(u_{\mathrm{front}}(t)\in[0,1]\), with larger values
corresponding to a flatter, more crowded current frontier and smaller
values corresponding to a clearly separated leading
record.\footnote{With the empty-pool convention $s_{\mathrm{raw},1}^{(1)}(t)=s_{\mathrm{raw},2}^{(1)}(t)=0$, one gets $u_{\mathrm{front}}(t)=1$. This treats the absence of a usable current frontier as maximal compression rather than as a missing-data null signal. An explicit signal-to-noise controller statistic is a possible implementation here, but it may not be necessary if worsening signal-to-noise already manifests as frontier flattening, which the present expected-utility model scores directly through $u_{\mathrm{front}}(t)$. Here $h_t(j)$ is the conditional probability that usefulness terminates at future offset $j$, given survival through offset $j-1$; accordingly, $S_t(k)$ is the survival probability to the start of offset $k$.}

Then the survival hazard is
\(h_t(j):=\sigma\!\Bigl(\beta_0+\beta_{\mathrm{stag}}u_{\mathrm{stag}}(t)+\beta_{\mathrm{front}}u_{\mathrm{front}}(t)+\eta_h(j-1)\Bigr)\).

\[
S_t(k):=\prod_{j=1}^{k-1}\bigl(1-h_t(j)\bigr),
\qquad
Q_t(k):=\Phi\!\left(
\frac{m_t-\gamma_t(k-1)-\operatorname{asinh}\!\left(\frac{\tau_{\mathrm{use},t}}{\tau_\Delta}\right)}{s_t}
\right).
\] \[
\widehat N_{\mathrm{rem}}(t)
:=
\sum_{k=1}^{D_{\mathrm{left}}(t)}
S_t(k)\,Q_t(k).
\] For pessimistic/optimistic runway envelopes, set
\(h_t^{\mathrm{pess/opt}}(j):=\sigma\!\bigl(\operatorname{logit}h_t(j)\pm\Delta_h\bigr)\),
\(m_t^{\mathrm{pess/opt}}:=m_t\mp\Delta_m\),
\(s_t^{\mathrm{pess/opt}}:=s_t+\Delta_s\) (pess) or
\(s_t^{\mathrm{pess/opt}}:=\max\!\left(s_{\min},s_t-\Delta_s\right)\)
(opt),
\(S_t^{\mathrm{pess/opt}}(k):=\prod_{j=1}^{k-1}\bigl(1-h_t^{\mathrm{pess/opt}}(j)\bigr)\),
and \(Q_t^{\mathrm{pess/opt}}(k):=\Phi\!\left(
\frac{m_t^{\mathrm{pess/opt}}-\gamma_t(k-1)-\operatorname{asinh}\!\left(\frac{\tau_{\mathrm{use},t}}{\tau_\Delta}\right)}{s_t^{\mathrm{pess/opt}}}
\right)\). Then
\(\widehat N_{\mathrm{rem}}^{\mathrm{low/high}}(t):=\sum_{k=1}^{D_{\mathrm{left}}(t)} S_t^{\mathrm{pess/opt}}(k)\,Q_t^{\mathrm{pess/opt}}(k)\),
with
\((\mathrm{low},\mathrm{high})\leftrightarrow(\mathrm{pess},\mathrm{opt})\),
and the confidence ratio is
\(C_t:=\frac{\widehat N_{\mathrm{rem}}^{\mathrm{low}}(t)}{\widehat N_{\mathrm{rem}}^{\mathrm{high}}(t)+\varepsilon}\).

For the controller coordinate and directional transforms, define \[
R_t:=\frac{\widehat N_{\mathrm{rem}}(t)}{D_{\mathrm{left}}(t)+\varepsilon}\in[0,1],
\qquad
e_t:=R_t^{a},
\qquad
g_t:=(1-R_t)^{b},
\qquad
a,b>0.
\] Here \(R_t\) is remaining useful runway as a fraction of the
depth-budget permitted by hardware/run-time constraints, and \(g_t\) is
the late-stage activation coordinate that scales controller variables
depending on age, where age is a normalized quantity relative
to/determined by expected utility of more generators.

Then \(R_t\approx 1\) early and \(R_t\approx 0\) late. Use \(e_t\) for
early-dominant quantities and \(g_t\) for late-dominant quantities. In
particular, define \[
\tau_{\mathrm{pos}}(t):=\tau_{\mathrm{pos}}^{\min}+(\tau_{\mathrm{pos}}^{\max}-\tau_{\mathrm{pos}}^{\min})e_t,
\qquad
\lambda_F(t):=\lambda_F^{\min}+(\lambda_F^{\max}-\lambda_F^{\min})g_t,
\qquad
\nu_{\mathrm{sat}}(t):=1-R_t.
\] Here \(\tau_{\mathrm{pos}}(t)\) is a position-domain parameter: it
controls how broadly the candidate-position fibers \(\mathcal P_m^k(t)\)
or active probe window \(W_{\mathrm{act}}\) may open around append and
active-window positions before phase scores are evaluated. It is not a
score threshold and does not multiply \(S_1\). The separate
inherited-window retention \(\rho_W(t)\) below controls how many old
coordinates are allowed to relax when a surviving record is locally
refit; it changes the geometry used after a record is shortlisted, not
the prior eligibility of a position for scoring.

The quantity \(\widehat N_{\mathrm{rem}}(t)\) is controller telemetry
rather than a physical observable: it summarizes recent selector
progress and remaining useful runway, but it is not itself a hard
admissibility
certificate.\footnote{Interpretation: the model has three layers --- (i) gain process ($m_t,s_t$), (ii) horizon-decay process ($\gamma_t$), and (iii) survival process ($h_t(j)$); together these drive the runway forecast $\widehat N_{\mathrm{rem}}(t)$.}

For inherited refit windows, use the maturity-shrinking retention law \[
\rho_W(t):=\rho_W^{\min}+(\rho_W^{\max}-\rho_W^{\min})e_t,
\qquad
0\le \rho_W^{\min}\le \rho_W^{\max}\le 1,
\] so that the retained old-coordinate fraction is largest early and
contracts as useful runway is exhausted.

The controller doctrine is monotone in scaffold maturity: as
\(R_t\downarrow\) and \(g_t\uparrow\), shot count, prune pressure, and
shortlist pressure increase, while inherited refit breadth, shortlist
breadth, and live phase depth decrease. Phase live/dead decisions are
driven by the runway telemetry \(\widehat N_{\mathrm{rem}}(t)\), not by
raw chronological time. Set \(\Gamma_1^{\mathrm{live}}(t):=1\). For
\(p\in\{2,3\}\), use hysteretic switches \[
\Gamma_p^{\mathrm{live}}(t^+)
:=
\begin{cases}
0,
&
\widehat N_{\mathrm{rem}}^{\mathrm{high}}(t)\le \nu_p^-
\text{ for }h_p\text{ consecutive controller steps},
\\[1mm]
1,
&
\widehat N_{\mathrm{rem}}^{\mathrm{low}}(t)>\nu_p^+,
\\[1mm]
\Gamma_p^{\mathrm{live}}(t^-),
&
\text{otherwise},
\end{cases}
\qquad
0\le \nu_2^-\le \nu_2^+<\nu_3^-\le \nu_3^+.
\] Then enforce
\(\Gamma_3^{\mathrm{live}}(t)\leftarrow\Gamma_2^{\mathrm{live}}(t)\Gamma_3^{\mathrm{live}}(t)\),
so Phase 3 cannot remain live after Phase 2 is nulled. The terminal live
selector phase is \[
k_\star(t):=1+\Gamma_2^{\mathrm{live}}(t)+
\Gamma_2^{\mathrm{live}}(t)\Gamma_3^{\mathrm{live}}(t),
\qquad
r_t^\star\in\arg\max_{r\in\mathcal R_{k_\star}(t)}S_{k_\star}(r;t).
\] Thus early scaffold construction may use the full Phase-1 \(\to\)
Phase-2 \(\to\) Phase-3 chain; when useful runway is nearly exhausted,
Phase 3 is nulled first, then Phase 2 is nulled, leaving the cheaper
greedy Phase-1 selector.

For measurement effort, use \(g_t\) as the canonical maturity-controlled
scheduled-shot baseline. For \(k\in\{1,2,3\}\), define \[
N_{\mathrm{shot},k}^{g}(t)
:=
\left\lceil
N_{\mathrm{shot},k}^{\min}
+
\bigl(N_{\mathrm{shot},k}^{\max}-N_{\mathrm{shot},k}^{\min}\bigr)g_t
\right\rceil,
\qquad
N_{\mathrm{shot},k}^{\min}\le N_{\mathrm{shot},k}^{\max}.
\] Frontier flattening and phase-geometric uncertainty are allowed only
as clipped uplifts on the remaining shot headroom. For Phase 1 set \[
h_1(t)
:=
\operatorname{clip}_{[0,1]}
\!\left(
\beta_{\mathrm{front},1}u_{\mathrm{front}}(t)
\right),
\] while for \(k\in\{2,3\}\), once the phase-\(k\) uncertainty summary
\(u_{\sigma,k}(t)\) is available, set \[
h_k(t)
:=
\operatorname{clip}_{[0,1]}
\!\left(
\beta_{\mathrm{front},k}u_{\mathrm{front}}(t)
+
\beta_{\sigma,k}u_{\sigma,k}(t)
\right),
\qquad
k\in\{2,3\},
\] with nonnegative uplift weights \(\beta_{\mathrm{front},k}\) and
\(\beta_{\sigma,k}\). The bounded scheduled-shot fraction is \[
\Phi_k(t):=g_t+(1-g_t)h_k(t),
\qquad
0\le \Phi_k(t)\le 1,
\] and the scheduled shot law is \[
N_{\mathrm{shot},k}^{\mathrm{sched}}(t):=
\left\lceil
N_{\mathrm{shot},k}^{\min}
+
\bigl(N_{\mathrm{shot},k}^{\max}-N_{\mathrm{shot},k}^{\min}\bigr)\Phi_k(t)
\right\rceil,
\qquad
k\in\{1,2,3\}.
\] Thus
\(h_k(t)=0\Rightarrow N_{\mathrm{shot},k}^{\mathrm{sched}}(t)=N_{\mathrm{shot},k}^{g}(t)\),
\(g_t\le \Phi_k(t)\le 1\), and, holding telemetry fixed, \[
\frac{\partial \Phi_k}{\partial g_t}=1-h_k(t)\ge 0.
\] Therefore \(N_{\mathrm{shot},k}^{g}(t)\) is the doctrinal maturity
floor, while \(u_{\mathrm{front}}(t)\) and \(u_{\sigma,k}(t)\) can only
increase shots above that floor when the selector frontier is flat or
the phase geometry is uncertain. This invariant does not require
irreversible wall-clock monotonicity: if telemetry collapses,
\(N_{\mathrm{shot},k}^{\mathrm{sched}}(t)\) may fall from a
telemetry-inflated value back toward the maturity floor, but it remains
bounded below by \(N_{\mathrm{shot},k}^{g}(t)\).

The actually deployed phase-\(k\) shot count is the live, capped maximum
of this scheduled law and an empirical SNR floor: \[
N_{\mathrm{shot},k}^{\mathrm{snr}}(t)
:=
\left\lceil
\frac{\kappa_k^2\,\widehat\sigma_{t,k}^{\,2}}
{\max\{\widehat\delta_{t,k}^{\,2},\delta_{\mathrm{floor},k}^{\,2}\}}
\right\rceil,
\qquad
N_{\mathrm{shot},k}^{\mathrm{eff}}(t)
:=
\Gamma_k^{\mathrm{live}}(t)
\min\!\left\{
N_{\mathrm{shot},k}^{\mathrm{cap}},
\max\!\left[
N_{\mathrm{shot},k}^{\mathrm{sched}}(t),
N_{\mathrm{shot},k}^{\mathrm{snr}}(t)
\right]
\right\}.
\] Here \(\widehat\delta_{t,k}\) is the useful decision-scale signal,
\(\widehat\sigma_{t,k}\) is its estimator spread, \(\kappa_k\) is the
target SNR, and \(\delta_{\mathrm{floor},k}>0\) prevents pathological
shot explosion when the phase has no resolvable useful gain. If
\(\widehat\delta_{t,k}\) falls below the useful floor, the controller
reads this as no resolvable gain at that phase, not as permission to
spend unbounded shots. The total live measurement burden is \[
N_{\mathrm{tot}}^{\mathrm{eff}}(t)
:=
N_{\mathrm{shot},1}^{\mathrm{eff}}(t)
+
\Gamma_2^{\mathrm{live}}(t)N_{\mathrm{shot},2}^{\mathrm{eff}}(t)
+
\Gamma_2^{\mathrm{live}}(t)\Gamma_3^{\mathrm{live}}(t)N_{\mathrm{shot},3}^{\mathrm{eff}}(t).
\] For phase shortlist thresholds, use the late-stricter monotone law \[
\tau_k^{\mathrm{short}}(t):=\tau_k^{\min}+(\tau_k^{\max}-\tau_k^{\min})g_t,
\qquad
\tau_k^{\min}\le \tau_k^{\max},
\qquad
k\in\{1,2,3\},
\] so that shortlist admission pressure increases as the scaffold
matures.

For phase shortlist breadth, use the live bounded law \[
B_k^{\mathrm{short}}(t):=
\left\lceil
B_k^{\min}+(B_k^{\max}-B_k^{\min})e_t
\right\rceil,
\qquad
B_k^{\min}\le B_k^{\max},
\qquad
k\in\{1,2,3\}.
\] Thus \(R_t\approx1\) early gives low shortlist thresholds and broad
candidate breadth, while \(R_t\approx0\) late gives high shortlist
thresholds and narrow candidate breadth. The invariant is
\(R_t\downarrow\Rightarrow g_t\uparrow\),
\(N_{\mathrm{shot},k}^{\mathrm{sched}}\uparrow\),
\(\tau_k^{\mathrm{short}}\uparrow\), \(B_k^{\mathrm{short}}\downarrow\),
and \(k_\star\downarrow\).

\subsubsection{11.2.4 Phase-local cost and structural burden
definitions}\label{phase-local-cost-and-structural-burden-definitions}

For any current phase-input record family \(\mathcal R_k(t)\), with
\(\mathcal R_1(t):=\mathcal R_1\) and \(k\in\{1,2,3\}\), and any scalar
record feature \(x:\mathcal R_k(t)\to\mathbb R_{\ge 0}\), define the
empirical quantiles
\(q_\alpha^x(t):=\operatorname{Quantile}_\alpha\bigl(\{x(u):u\in\mathcal R_k(t)\}\bigr)\)
and \(m_x(t):=q_{0.5}^x(t)\), together with the robust scales
\(s_x^{\mathrm{MAD}}(t):=1.4826\,\operatorname{median}_{u\in\mathcal R_k(t)}\bigl|x(u)-m_x(t)\bigr|\),
\(s_x^{\mathrm{pct}}(t):=\frac{q_{0.9}^x(t)-q_{0.1}^x(t)}{2.563}\), and
\(s_x(t):=\max\!\bigl(\varepsilon_x,\;s_x^{\mathrm{MAD}}(t),\;s_x^{\mathrm{pct}}(t)\bigr)\).
Thus \(s_x(t)\) is a percentile/MAD hybrid reference scale, with
percentile fallback when the MAD degenerates on sparse discrete
features.

Define the one-sided robust normalized excess by
\(z_x(r;t):=\operatorname{asinh}\!\left(\frac{[x(r)-m_x(t)]_+}{s_x(t)}\right)\),
with \([y]_+:=\max(y,0)\). When a direct ratio form is preferred, use
the reference level \(x_{\mathrm{ref}}(t):=m_x(t)+s_x(t)\).

For a candidate-position record \(r=(m,p)\), let
\(\mathcal L_r:=\operatorname{supp}(m)=\{j:m_j\neq I\}\),
\(\mathrm{wt}(r):=|\mathcal L_r|\), and
\(\#_{X,Y}(r):=\left|\{j\in\mathcal L_r:m_j\in\{X,Y\}\}\right|\).

Define feature-specific robust excesses by
\(z_{\mathrm{wt}}(r;t):=z_x(r;t)\ \text{with }x=\mathrm{wt}\) and
\(z_{XY}(r;t):=z_x(r;t)\ \text{with }x=\#_{X,Y}\). The structural burden
term is \[
D_{\mathrm{struct}}(r;t):=
\beta_{\mathrm{wt}}\,z_{\mathrm{wt}}(r;t)
+\beta_{XY}\,z_{XY}(r;t),
\] with nonnegative coefficients
\(\beta_{\mathrm{wt}},\beta_{XY}\ge 0\).\footnote{Provenance note: generator novelty $G_{\mathrm{new}}(r)$, family novelty $\mathrm{groups}_{\mathrm{new}}(r)$ with single-valued $\mathfrak g:\mathcal G\to\mathcal C$, lifetime/staleness $\mathbf 1_{\mathrm{lifetime}}(r;t)$ with $t_{\mathrm{first}}(m)$ meaning first time ever seen in the run, and the normalized descendants $\bar D_{\mathrm{life}}^{\mathrm{struct}},\bar G_{\mathrm{new}},\bar C_{\mathrm{new}},\bar c,\bar P_{\mathrm{opt}}$ are omitted from the live hardware cost in the present draft and retained only as controller-side provenance/policy bookkeeping. Append-distance pressure is a position-domain construction policy for candidate-position records, not a live hardware-cost term and not a multiplicative Phase-1 score prior.}

For candidate-specific measurement burden, let
\(O_r:=i[H,\widetilde A_r]=\sum_{b\in\mathcal B_r} O_{r,b}\), with
\(O_{r,b}:=\sum_{\ell\in B_{r,b}} c_{r,b\ell} P_{r,b\ell}\), where
\(\mathcal B_r\) is the qubit-wise commuting partition of the
candidate's gradient observable and
\(\mathcal B_r^{\mathrm{new}}\subseteq\mathcal B_r\) is the subset not
already covered by measurement reuse. Here \(\widetilde A_r\) is the
dressed candidate generator at record \(r=(m,p)\), \(\mathcal B_r\) is
the partition of \(O_r=i[H,\widetilde A_r]\) into qubit-wise commuting
Pauli-string measurement groups, and \(\mathcal B_r^{\mathrm{new}}\) is
the subset that still requires new measurement effort. When exact
state-dependent variances are unavailable, use the commutator-based
proxy \[
\widehat v_{r,b}:=\sum_{\ell\in B_{r,b}} |c_{r,b\ell}|^2,
\qquad
\widehat{\bar C}_{\mathrm{shot}}(r):=
\frac{1}{\mathrm{shot}_{\mathrm{ref}}}
\cdot
\frac{
\left(
\sum_{b\in\mathcal B_r^{\mathrm{new}}}\sqrt{\widehat v_{r,b}}
\right)^2
}{
\sigma_{\star}^{2}
}.
\] Here \(\sigma_\star\) is the target standard error for the candidate
gradient estimate. The positive constant
\(\mathrm{shot}_{\mathrm{ref}}\) is a baseline measurement budget
against which candidate-specific shot cost is compared.

For the shared phase-\(k\) cost, define
\(z_{\mathrm{2q}}(r;t):=z_x(r;t)\ \text{with }x=C_{\mathrm{2q}}\),
\(z_d(r;t):=z_x(r;t)\ \text{with }x=D\),
\(z_{\theta}(r;t):=z_x(r;t)\ \text{with }x=\Delta K\),
\(z_{\mathrm{shot}}(r;t):=z_x(r;t)\ \text{with }x=\widehat{\bar C}_{\mathrm{shot}}\),
and \[
K_k(r;t):=
\lambda_{\mathrm{2q}}\,z_{\mathrm{2q}}(r;t)
+\lambda_d\,z_d(r;t)
+\lambda_\theta\,z_{\theta}(r;t)
+\lambda_{\mathrm{shot}}\,z_{\mathrm{shot}}(r;t)
+D_{\mathrm{struct}}(r;t).
\] Explicitly, \[
\begin{aligned}
K_k(r;t)
&=
\lambda_{\mathrm{2q}}\,z_{\mathrm{2q}}(r;t)
+\lambda_d\,z_d(r;t)
+\lambda_\theta\,z_{\theta}(r;t)
+\lambda_{\mathrm{shot}}\,z_{\mathrm{shot}}(r;t)\\
&\qquad
+\beta_{\mathrm{wt}}\,z_{\mathrm{wt}}(r;t)
+\beta_{XY}\,z_{XY}(r;t),
\qquad
k\in\{1,2,3\}.
\end{aligned}
\]

\subsubsection{Algebraic commutation metadata and flat-sector scaffold features}

Before the selector evaluates live energy geometry, the generator pool
may carry static algebraic metadata. This layer identifies support
overlap, bare-generator commutation, and local internally commuting
sectors in the pool or current scaffold. It is a structural prior only:
it is not a dressed runtime commutation claim, and it is not a
substitute for tangent overlap, Hessian coupling, reduced novelty, or
remove-refit recovery.

For bare generator labels define three static metadata entries: the
support-overlap bit, the pairwise commutation bit, and an exactness tag
that records whether the commutation entry was computed exactly or
retained only as approximate metadata: \[
\Omega_{mn}:=\mathbf 1[\operatorname{supp}(m)\cap\operatorname{supp}(n)\neq\varnothing],
\qquad
\mathsf C_{mn}:=\mathbf 1[[A_m,A_n]=0],
\qquad
\mathsf q_{mn}\in\{\mathrm{exact},\mathrm{approx}\}.
\] Here \(\Omega_{mn}\) records whether the two bare generators have
overlapping support, \(\mathsf C_{mn}\) records whether their bare
generators commute, and \(\mathsf q_{mn}\) records the reliability of
that commutation entry: \(\mathrm{exact}\) means the entry was computed
from an exact Pauli expansion or an exact algebraic rule, while
\(\mathrm{approx}\) means only weak metadata was available. The tag
\(\mathsf q_{mn}\) is an exactness label for the pair \((m,n)\), not a
qubit index. For Pauli-word generators
\(A_m=\bigotimes_\ell P_\ell^{(m)}\) and
\(A_n=\bigotimes_\ell P_\ell^{(n)}\), the exact rule is \[
\mathsf C_{mn}
=
\mathbf 1\!\left[
\sum_{\ell\in\operatorname{supp}(m)\cap\operatorname{supp}(n)}
\mathbf 1\!\big[P_\ell^{(m)}P_\ell^{(n)}=-P_\ell^{(n)}P_\ell^{(m)}\big]
\equiv 0\pmod 2
\right],
\qquad
\mathsf q_{mn}=\mathrm{exact}.
\] For non-Pauli composite generators the table is computed from the
expanded Pauli support whenever available; otherwise the corresponding
entries retain \(\mathsf q_{mn}=\mathrm{approx}\) and may be used only
as weak prior evidence.

The static metadata layer keeps three distinct relations: \[
\begin{aligned}
\mathcal E_{\mathrm{disj}}^{\mathrm{comm}}
&:=\{(m,n):\Omega_{mn}=0,\ \mathsf C_{mn}=1\},\qquad
&&\text{support-disjoint commutation},\\
\mathcal E_{\mathrm{flat}}^{\mathrm{comm}}
&:=\{(m,n):\Omega_{mn}=1,\ \mathsf C_{mn}=1\},
&&\text{overlapping flat-sector commutation},\\
\mathcal E_{\mathrm{curv}}^{\mathrm{nc}}
&:=\{(m,n):\Omega_{mn}=1,\ \mathsf C_{mn}=0\},
&&\text{overlapping noncommutation / curvature relation}.
\end{aligned}
\] Disjoint commutation is not redundancy evidence: two support-disjoint
commuting generators may be tangent-orthogonal, geometrically novel, and
simultaneously desirable. The overlapping flat-sector relation is the
relation relevant to low-risk local redundancy and cluster saturation.
The overlapping noncommutation relation is the relation relevant to
frame rotation, curvature-seeking selection, and bridge protection.

Let \(\mathcal N_{\mathrm{loc}}(m)\) be the deterministic locality
neighborhood used for the candidate pool, for example support overlap
plus one-hop hardware or lattice neighbors. This neighborhood is defined
by physical or algebraic locality, not by an interval in the ordered
generator list. Let \(\{\mathfrak C_a\}_a\) denote the internally
commuting clusters returned by a deterministic greedy pass over the
ordered pool, and let \(c(m)=a\) mean that generator \(m\) belongs to
cluster \(\mathfrak C_a\). Equivalently, the pass gives \[
\operatorname{comm\_cluster\_id}(m):=c(m),
\qquad
m,n\in\mathfrak C_a\Longrightarrow \mathsf C_{mn}=1,
\] with tie order inherited from the global record order \(\prec\). The
order \(\prec\) is used only for deterministic ties and reproducible
cluster assignment; it is not the definition of locality.

For Phase-1 stratification, let \(\mathcal C_1(r;t)\subseteq\mathcal G\)
be the local algebraic context of the record \(r=(m,p)\) at selector
step \(t\). This context is built from \(\mathcal N_{\mathrm{loc}}(m)\),
recent scaffold coordinates relevant to the selector, and any inherited
local window information used for the Phase-1 record; it is not an
ordered-list window. Define the exact local relation counts \[
\begin{aligned}
n_{\mathrm{flat}}(r;t)
&:=
\sum_{n\in\mathcal C_1(r;t)}
\mathbf 1[\mathsf q_{mn}=\mathrm{exact}]\,\Omega_{mn}\mathsf C_{mn},\\[1mm]
n_{\mathrm{curv}}(r;t)
&:=
\sum_{n\in\mathcal C_1(r;t)}
\mathbf 1[\mathsf q_{mn}=\mathrm{exact}]\,\Omega_{mn}(1-\mathsf C_{mn}),\\[1mm]
n_{\mathrm{disj}}(r;t)
&:=
\sum_{n\in\mathcal C_1(r;t)}
\mathbf 1[\mathsf q_{mn}=\mathrm{exact}]\,(1-\Omega_{mn})\mathsf C_{mn},\\[1mm]
n_{\mathrm{approx}}(r;t)
&:=
\sum_{n\in\mathcal C_1(r;t)}
\mathbf 1[\mathsf q_{mn}=\mathrm{approx}].
\end{aligned}
\] Let \[
\mathfrak L:=
\{\ell_{\mathrm{flat}},\ell_{\mathrm{curv}},\ell_{\mathrm{disj}},\ell_{\mathrm{mix}}\}
\] be the algebraic lane set, where \(\ell_{\mathrm{flat}}\) denotes
overlapping commuting records, \(\ell_{\mathrm{curv}}\) denotes
overlapping noncommuting records, \(\ell_{\mathrm{disj}}\) denotes
support-disjoint commuting records, and \(\ell_{\mathrm{mix}}\) denotes
mixed or ambiguous records. The support-disjoint lane is a coverage
lane, not a redundancy class. Define the conservative Phase-1 lane map
\[
\ell_1(r;t)
:=
\begin{cases}
\ell_{\mathrm{flat}},
&
n_{\mathrm{flat}}(r;t)>0,\quad
n_{\mathrm{curv}}(r;t)=0,\quad
n_{\mathrm{approx}}(r;t)=0,
\\[1mm]
\ell_{\mathrm{curv}},
&
n_{\mathrm{curv}}(r;t)>0,\quad
n_{\mathrm{flat}}(r;t)=0,\quad
n_{\mathrm{approx}}(r;t)=0,
\\[1mm]
\ell_{\mathrm{disj}},
&
n_{\mathrm{flat}}(r;t)=0,\quad
n_{\mathrm{curv}}(r;t)=0,\quad
n_{\mathrm{disj}}(r;t)>0,\quad
n_{\mathrm{approx}}(r;t)=0,
\\[1mm]
\ell_{\mathrm{mix}},
&
\text{otherwise}.
\end{cases}
\] Thus commutation determines a categorical lane; it does not multiply
the Phase-1 score.

Continuous commutation averages and bridge summaries are not part of the
canonical selector. They are defined only in Appendix 11-S for
telemetry, lane-budget diagnostics, prune-risk bookkeeping, and ablation
studies.

The selection doctrine is local: do not reward noncommutation globally.
Do not penalize consecutive commuting generators merely because they
commute. Use commutation to open Phase-1 lanes, use Phase-2 geometry to
keep or retire inherited lanes for the current selector step, and use
Phase-3 reduced geometry to decide final rank.

Safety convention: do not use commutation as a deletion certificate; do
not treat disjoint commuting generators as redundant; do not use global
commuting degree as a strong downweight; do not allow motif,
commutation, or flat-sector bonuses to rescue weak reduced geometry; and
do not require exact dressed commutation at runtime. Bare commutation is
the precomputed algebraic prior, while tangent overlap, Hessian
coupling, and reduced-window recovery are the effective dressed
geometric diagnostics.

The controller rule is: commutation opens Phase-1 lanes; Phase-2
geometry keeps or retires lanes; Phase-3 geometry decides the final
rank; prune recovery remains categorical through flat-sector versus
curvature-compensated recovery classes.

\subsubsection{11.2.5 Final Phase I Scoring}

For the current scaffold at selector step \(t\), let
\(p_{\mathrm{app}}(t):=|\mathcal O_t|\) be the append slot and define
the normalized append-distance \[
d_{\mathrm{app}}(r;t):=\frac{|p-p_{\mathrm{app}}(t)|}{\max\{1,|\mathcal O_t|\}},
\qquad
0\le d_{\mathrm{app}}(r;t)\le 1.
\] This append-distance is position-domain telemetry. It may be used to
construct or audit the admissible position fibers
\(\mathcal P_m^{1}(t)\), the active probe window
\(W_{\mathrm{act}}(t)\), or the cap \(M_{\mathrm{probe}}\), but it is
not a multiplicative Phase-1 score factor. Position preference is
imposed before scoring by choosing which insertion positions are
eligible for evaluation: \[
\mathcal P_m^{1}(t)\subseteq W_{\mathrm{act}}(t),
\qquad
|\mathcal P_m^{1}(t)|\le M_{\mathrm{probe}},
\qquad
\mathcal R_1(t):=\{(m,p):m\in\mathcal G_1,\ p\in\mathcal P_m^{1}(t)\}.
\] Thus append-like records may receive protected exposure by the
position-window policy, while far internal insertions may be omitted or
delayed by the probe policy. Once a record has entered
\(\mathcal R_1(t)\), the rank score is determined by its geometric gain
and cost burden, not by a separate append-distance multiplier.

For Phase 1, the set \(\mathcal R_1(t)\) is ranked inside algebraic
lanes by the geometric screen, \[
\begin{aligned}
S_{1,\mathrm{geom}}(r;t)
&:=
\frac{\Delta E_{1,\mathrm{TR}}(r)}{1+K_1(r;t)},\\[1mm]
S_{1,\mathrm{rank}}(r;t)
&:=
S_{1,\mathrm{geom}}(r;t),
\qquad r\in\mathcal R_1.
\end{aligned}
\] Here \(\Delta E_{1,\mathrm{TR}}(r)\) is the Phase-1 one-coordinate
trust-region proxy defined in \S 11.2.1, and the phase-1 cost is
evaluated on the current Phase-1 input family. The algebraic layer
assigns the lane \(\ell_1(r;t)\) and controls shortlist coverage; it
does not multiply the Phase-1 rank score. Position-window bias has
already acted by constructing \(\mathcal R_1(t)\); it does not reweight
records after their Phase-1 quantities have been computed. In
particular, Phase 1 does not penalize a candidate merely because it
commutes with the last generator, and support-disjoint commutation is
not a redundancy signal. This Phase-1 rank score is intentionally broad
as a screening surface: it privileges cheap local second-order insertion
estimates while the controller still has long useful runway. The
corresponding broad Phase-1 probing/evaluation budget is set by the
baseline controller shot law \(N_{\mathrm{shot},1}(t)\).

In the present draft, the live Phase-1 numerator is the one-coordinate
trust-region proxy \(\Delta E_{1,\mathrm{TR}}(r)\), so that the broad
screen remains in the same gradient-plus-second-order family later
refined in Phases 2 and 3. The strict one-coordinate insertion drop
\(\delta E_1(r\mid\mathcal O,\theta)\) is retained only as an auxiliary
exact comparator for selective frontier confirmation, ablation studies,
or future confirm-stage refinement. A broader early-stage local-refit
numerator remains only a non-live hook. To make such a refinement live,
one would first have to specify (i) the admissible refit-window family,
(ii) the corresponding constrained local-refit drop, and (iii) a
controller activation law independent of the current raw Phase-1 score
multiset, so that the live Phase-1 law is not self-referential.

\subsubsection{11.2.6 Position probing, lane-wise trough detection, and
effective
selector}\label{position-probing-lane-wise-trough-detection-and-effective-selector}

The Phase-1 scoring mechanism produces a stratified record shortlist for
further refined scoring. For each algebraic lane \(\ell\in\mathfrak L\),
define \[
\mathcal R_{1,\ell}(t):=\{r\in\mathcal R_1(t):\ell_1(r;t)=\ell\},
\qquad
\mathcal A_{1,\ell}(t):=\{r\in\mathcal R_{1,\ell}(t):S_{1,\mathrm{rank}}(r;t)\ge \tau_1(t)\}.
\] Let \(B_{1,\ell}(t)\in\mathbb N\) be the Phase-1 shortlist budget
assigned to lane \(\ell\). The default lane-budget doctrine is
asymmetric. The mixed lane receives protected exposure because mixed can
mean unresolved or approximate algebraic evidence rather than weak
evidence. For \(\mathcal R_{1,\mathrm{mix}}(t)\neq\varnothing\), define
\[
\bar f_{\mathrm{approx}}^{\mathrm{mix}}(t)
:=
\frac{
\sum_{r\in\mathcal R_{1,\mathrm{mix}}(t)} n_{\mathrm{approx}}(r;t)
}{
\sum_{r\in\mathcal R_{1,\mathrm{mix}}(t)}
\left[n_{\mathrm{flat}}(r;t)+n_{\mathrm{curv}}(r;t)+n_{\mathrm{disj}}(r;t)+n_{\mathrm{approx}}(r;t)\right]
\varepsilon
}.
\] A canonical protected mixed budget is \[
B_{1,\mathrm{mix}}(t)
:=
\min\!\left\{
B_{1,\mathrm{mix}}^{\max},
B_{1,\mathrm{mix}}^{\min}
+
\left\lceil
\kappa_{\mathrm{mix}}\,
\bar f_{\mathrm{approx}}^{\mathrm{mix}}(t)\,
B_1^{\mathrm{short}}(t)
\right\rceil
\right\},
\qquad
B_{1,\mathrm{mix}}^{\min}\ge 1.
\] Thus approximate metadata can increase mixed-lane exposure, but it
does not increase any record's score. The disjoint lane has the
complementary default: it receives modest Phase-1 exposure because
support-disjoint commutation is not redundancy evidence, but it is often
less locally corrective than overlapping flat or curvature evidence. A
typical default satisfies \[
B_{1,\mathrm{disj}}^{\min}
\le
B_{1,\mathrm{flat}}^{\min},
\qquad
B_{1,\mathrm{disj}}^{\min}
\le
B_{1,\mathrm{curv}}^{\min},
\] while still keeping \(B_{1,\mathrm{disj}}^{\min}>0\) whenever
disjoint exploration is enabled. These lane budgets control exposure
only; records inside every lane are still ordered by
\(S_{1,\mathrm{rank}}\).

Order the records in \(\mathcal A_{1,\ell}(t)\) as \[
r_{\ell,(1)}^{(1)},\dots,r_{\ell,(M_\ell)}^{(1)},
\qquad
s_{\ell,j}^{(1)}(t):=S_{1,\mathrm{rank}}\!\left(r_{\ell,(j)}^{(1)};t\right),
\] with
\(s_{\ell,1}^{(1)}(t)\ge s_{\ell,2}^{(1)}(t)\ge\cdots\ge s_{\ell,M_\ell}^{(1)}(t)\),
ties resolved by \(\prec\). For \(M_\ell\ge2\), define the lane-local
frontier ratio \[
u_{\mathrm{front},k}^{(1,\ell)}(t)
:=
\frac{s_{\ell,k+1}^{(1)}(t)+\varepsilon}{s_{\ell,k}^{(1)}(t)+\varepsilon},
\qquad
k=1,\dots,M_\ell-1.
\] With \(0<\eta_{1,\ell}<1\), set \[
\begin{aligned}
k_{1,\ell}^\star(t)&:=
\begin{cases}
0, & M_\ell=0,\\[4pt]
\min\!\left\{k\in\{1,\dots,\min(M_\ell-1,B_{1,\ell}(t)-1)\}:u_{\mathrm{front},k}^{(1,\ell)}(t)\le \eta_{1,\ell}\right\},
& \text{if this set is nonempty},\\[6pt]
\min\{M_\ell,B_{1,\ell}(t)\}, & \text{otherwise},
\end{cases}\\
\mathcal S_{1,\ell}(t)&:=\{r_{\ell,(j)}^{(1)}:1\le j\le k_{1,\ell}^\star(t)\}.
\end{aligned}
\] The retained Phase-1 shortlist is the lane union \[
\mathcal S_1(t):=\bigcup_{\ell\in\mathfrak L}\mathcal S_{1,\ell}(t).
\] The retained model set is
\(\mathcal G^{(1)}:=\pi_m(\mathcal S_1(t))\), the retained position
fibers are
\(\mathcal P_m^{(1)}:=\{p:(m,p)\in\mathcal S_1(t)\}\ \forall m\in\mathcal G^{(1)}\),
and the Phase-2 record pool is \(\mathcal R_2(t):=\mathcal S_1(t)\).
Thus Phase 1 gives each algebraic lane controlled exposure, while the
ranking within each lane remains geometric-cost based.

\subsection{11.3 Phase II}\label{phase-ii}

Phase II is the raw geometric selector stage. It acts on the Phase-1
survivors, restores the coarse geometric signal that distinguishes
directional usefulness beyond insertion-drop alone, and produces the
filtered shortlist \(\mathcal S_2\).

\subsubsection{11.3.1 Raw geometry, restored novelty, and auxiliary
overlap
window}\label{raw-geometry-restored-novelty-and-auxiliary-overlap-window}

The Phase-2 selector works over candidate-position records inherited
from Phase 1. In the record-inheritance mode used here, the inherited
and admissible domain is \(\mathcal R_2(t):=\mathcal S_1(t)\). Each
inherited record carries its Phase-1 lane label \[
\ell_2(r;t):=\ell_1(r;t),
\qquad
\mathcal R_{2,\ell}(t):=\{r\in\mathcal R_2(t):\ell_2(r;t)=\ell\},
\qquad
\ell\in\mathfrak L.
\] For \(r=(m,p)\), let \(W(r;t)\subseteq\{1,\dots,n+1\}\) denote the
inherited overlap window attached to \(r\). In the present three-phase
architecture this window enters the Phase-2 selector through the Fubini
tangent span of the current horizontal tangent family; the same
inherited window is then carried forward to Phase 3, where
reduced-window elimination and final reranking are
performed.\footnote{Here $p$ indexes the cut after $U_p$. Thus $U_{>p}:=U_K\cdots U_{p+1}$ and $U_{\le p}:=U_p\cdots U_1$, and insertion at cut $p$ means $ |\psi_r(\alpha)\rangle = U_{>p}e^{-i\alpha A_m}U_{\le p}|\psi_0\rangle = U_K\cdots U_{p+1}\,e^{-i\alpha A_m}\,U_p\cdots U_1|\psi_0\rangle $. So the new generator is inserted between $U_p$ and $U_{p+1}$.}

Reusing the record-local insertion objects already introduced in Phase
1, let \(|d_r^{(2)}\rangle:=-\widetilde A_r^2|\psi\rangle\). The
gradient estimator uncertainty is already accounted for through the
lower-confidence gradient
\(g_{\mathrm{lcb}}(r)=\max\{|g_r|-z_\alpha\sigma_r,0\}\), so the live
Phase-2 and Phase-3 scores do not include an additional multiplicative
confidence
factor.\footnote{Let $\widehat g(r)$ denote the empirical gradient estimator for record $r=(m,p)$ obtained from $N_r$ independent shot-level repetitions. If $X^{(r)}_1,\dots,X^{(r)}_{N_r}$ are the repetition-level gradient samples, then $\widehat g(r):=\frac{1}{N_r}\sum_{i=1}^{N_r} X^{(r)}_i$ and $\sigma_r:=\operatorname{SE}\!\big[\widehat g(r)\big]=\sqrt{\operatorname{Var}\!\big(\widehat g(r)\big)}$. Under independent repetitions, $\operatorname{Var}\!\big(\widehat g(r)\big)=\frac{\operatorname{Var}\!\big(X^{(r)}_1\big)}{N_r}$, so in practice one estimates $\sigma_r\approx \frac{s_r}{\sqrt{N_r}}$, where $s_r$ is the sample standard deviation of the repetition-level gradients. Thus $\sigma_r$ is the standard error of the mean gradient estimate, not the sample standard deviation itself.}
For phase-dependent shot allocation in later phases, define \[
u_{\sigma,k}(t):=
\operatorname{clip}\!\left(
\operatorname{median}_{r\in\mathcal R_k(t)}
\frac{z_\alpha\,\sigma_r}{|g_r|+\varepsilon_g},
\,0,\,1
\right),
\qquad
k\in\{2,3\},
\] where in Phase 2 the record-local gradient and uncertainty estimates
are evaluated using \(N_{\mathrm{shot},2}(t)\), and in Phase 3 the pair
\((g_r,\sigma_r)\) is inherited from the Phase-II raw geometry carried
by the surviving records.

The raw Phase-2 tangent-span quantities are \[
\begin{aligned}
F_{\mathrm{raw}}(r)&:=\langle \tau_r,\tau_r\rangle+\varepsilon_F,
\qquad \varepsilon_F>0,\\
(Q_r^{(2)})_{ij}
&:=
\Re\langle \tau_i,\tau_j\rangle,
\qquad i,j\in W(r;t),\\
(q_r^{(2)})_i
&:=
\Re\langle \tau_i,\tau_r\rangle,
\qquad i\in W(r;t),\\
z_r^{(2)}
&:=
\frac{(q_r^{(2)})^\top
\bigl(Q_r^{(2)}+\lambda_{2,\mathrm{nov}}I\bigr)^{-1}
q_r^{(2)}}
{F_{\mathrm{raw}}(r)},
\qquad \lambda_{2,\mathrm{nov}}>0,\\
\mathcal N_2(r;t)
&:=\operatorname{clip}_{[0,1]}\!\left(1-z_r^{(2)}\right),\\
h_r&:=2\Re\!\left(\langle d_r|H|d_r\rangle+\langle \psi|H|d_r^{(2)}\rangle\right),\\
\Delta E_{\mathrm{TR}}^{\mathrm{raw}}(r)
&:=\max_{|\alpha|\le \rho/\sqrt{F_{\mathrm{raw}}(r)}}
\left[g_{\mathrm{lcb}}(r)|\alpha|-\frac12\max\bigl(h_r,0\bigr)\alpha^2\right].
\end{aligned}
\] Here \(F_{\mathrm{raw}}(r)\) is the raw candidate tangent norm,
\(Q_r^{(2)}\) is the inherited-window Fubini Gram matrix, \(q_r^{(2)}\)
is the candidate-to-window tangent covariance vector,
\(\mathcal N_2(r;t)\) is the raw Phase-2 span-novelty score, \(h_r\) is
the raw local curvature, and
\(\Delta E_{\mathrm{TR}}^{\mathrm{raw}}(r)\) is the raw trust-region
gain before reduced-window
elimination.\footnote{The archaic Phase-II root used the pairwise proxy $1-\max_{j\in W(r;t)}|\langle\tau_j,\tau_r\rangle|/(\|\tau_j\|\,\|\tau_r\|)$. That proxy is retained only as a cheap diagnostic or fallback in implementations that do not materialize $Q_r^{(2)}$. It is not the canonical manuscript novelty factor, because pairwise low overlap can miss collective span redundancy.}

The authoritative Phase-II score is \[
S_2(r;t)
:=
\frac{\Delta E_{\mathrm{TR}}^{\mathrm{raw}}(r)\,\mathcal N_2(r;t)}{1+K_2(r;t)},
\qquad r\in\mathcal R_2.
\] Thus Phase II remains the raw geometric selector stage: it restores
directional discrimination beyond insertion-drop alone, uses the
lower-confidence trust-region gain to handle gradient noise, tests
candidate independence against the inherited tangent span, and filters
the Phase-1 survivors before any reduced-window refit quantity is made
authoritative. Commutation no longer modifies the Phase-2 score; it only
supplies the inherited lane label whose current usefulness is tested by
Phase-2 geometry.

For continuity with the older local-refit notation, retain the auxiliary
windowed drop \[
\mathcal C_W(r;\theta):=\{\widetilde\theta\in\Theta_{\mathcal O\oplus_p m}:\widetilde\theta_j=(\theta\oplus_p 0)_j\ \forall j\notin W(r;t)\}
\] and \[
\delta E_2(r\mid\mathcal O,\theta):=
E(\mathcal O,\theta)-\inf_{\widetilde\theta\in\mathcal C_W(r;\theta)}E(\mathcal O\oplus_p m,\widetilde\theta).
\] This \(\delta E_2\) quantity is retained only as an auxiliary
comparator for the inherited window. It may be used for notation
continuity, diagnostics, or ablation studies, but it does not define a
separate selector phase, it does not enter \(S_2(r;t)\), and it does not
modify the Phase-II shortlist rule. Any authoritative reduced-window
rerank remains deferred to Phase 3.

\subsubsection{11.3.2 Inherited lane survival and shortlist
rule}\label{inherited-lane-survival-and-shortlist-rule}

Phase 2 first ranks records inside each inherited lane, then ranks lanes
against one another through a geometry-only lane-health statistic. For
each \(\ell\in\mathfrak L\), let \[
M_{2,\ell}(t):=|\mathcal R_{2,\ell}(t)|.
\] If \(M_{2,\ell}(t)>0\), order \(\mathcal R_{2,\ell}(t)\) as \[
r_{\ell,(1)}^{(2)},\dots,r_{\ell,(M_{2,\ell})}^{(2)},
\qquad
s_{\ell,j}^{(2)}(t):=S_2\!\left(r_{\ell,(j)}^{(2)};t\right),
\] with
\(s_{\ell,1}^{(2)}(t)\ge s_{\ell,2}^{(2)}(t)\ge\cdots\ge s_{\ell,M_{2,\ell}}^{(2)}(t)\),
ties resolved by \(\prec\). The default lane-health statistic is the
best resolved raw-geometric score in the lane, \[
H_{2,\ell}(t):=
\begin{cases}
s_{\ell,1}^{(2)}(t),&M_{2,\ell}(t)>0,\\
0,&M_{2,\ell}(t)=0.
\end{cases}
\] This default encodes the rule that one strong Phase-2 candidate is
enough to keep a lane alive for the current selector step. A top-few
mean may be used as a robustness ablation, but it is not the canonical
lane-health statistic. Define the relative lane health \[
H_{2,\star}(t):=\max_{\ell'\in\mathfrak L}H_{2,\ell'}(t),
\qquad
\widehat H_{2,\ell}(t):=
\begin{cases}
H_{2,\ell}(t)/H_{2,\star}(t),&H_{2,\star}(t)>0,\\
0,&H_{2,\star}(t)=0.
\end{cases}
\] Let \(\tau_2^{\mathrm{abs}}(t)\ge0\) be the common absolute
lane-survival scale and \(\tau_2^{\mathrm{rel}}(t)\in[0,1]\) the common
relative lane-survival scale. Lane-specific patience factors
\(c_\ell^{\mathrm{abs}}(t)>0\) and \(c_\ell^{\mathrm{rel}}(t)>0\) define
\[
\tau_{2,\ell}^{\mathrm{abs}}(t):=c_\ell^{\mathrm{abs}}(t)\tau_2^{\mathrm{abs}}(t),
\qquad
\tau_{2,\ell}^{\mathrm{rel}}(t):=
\operatorname{clip}_{[0,1]}\!\left(c_\ell^{\mathrm{rel}}(t)\tau_2^{\mathrm{rel}}(t)\right).
\] Smaller \(c_\ell^{\mathrm{abs}}\) or \(c_\ell^{\mathrm{rel}}\) gives
lane \(\ell\) more patience; larger values require stronger geometric
evidence. These factors may be fixed, scheduled by scaffold maturity, or
learned by the outer policy, but they affect only lane survival and lane
budgets, not the intra-lane record ranking.

The live-lane gate is \[
\Gamma_{2,\ell}^{\mathrm{live}}(t)
:=
\mathbf 1[M_{2,\ell}(t)>0]\,
\mathbf 1[H_{2,\ell}(t)\ge\tau_{2,\ell}^{\mathrm{abs}}(t)]\,
\mathbf 1[\widehat H_{2,\ell}(t)\ge\tau_{2,\ell}^{\mathrm{rel}}(t)].
\] Lane failure is local to the current selector step:
\(\Gamma_{2,\ell}^{\mathrm{live}}(t)=0\) does not delete the lane from
the generator universe and does not certify algebraic redundancy.

For live lanes, define the lane-local record acceptance set by a record
floor and an intra-lane closeness rule, \[
\mathcal A_{2,\ell}(t)
:=
\left\{
r\in\mathcal R_{2,\ell}(t):
S_2(r;t)\ge\tau_2(t),\
S_2(r;t)\ge \eta_{2,\ell}^{\mathrm{in}}(t)s_{\ell,1}^{(2)}(t)
\right\},
\] where \(0\le\eta_{2,\ell}^{\mathrm{in}}(t)\le1\). Let
\(B_{2,\ell}(t)\in\mathbb N\) be the Phase-2 budget assigned to lane
\(\ell\). The lane-wise Phase-2 shortlist is \[
\mathcal S_{2,\ell}(t)
:=
\begin{cases}
\operatorname{Top}_{B_{2,\ell}(t)}\!\left(S_2;\mathcal A_{2,\ell}(t)\right),
&
\Gamma_{2,\ell}^{\mathrm{live}}(t)=1,\\[1mm]
\varnothing,
&
\Gamma_{2,\ell}^{\mathrm{live}}(t)=0.
\end{cases}
\] The retained Phase-2 shortlist is \[
\mathcal S_2(t):=\bigcup_{\ell\in\mathfrak L}\mathcal S_{2,\ell}(t).
\] Of course, \(\mathcal G^{(2)}:=\pi_m(\mathcal S_2(t))\), and
\(\mathcal P_m^{(2)}:=\{p:(m,p)\in\mathcal S_2(t)\}\). Thus Phase 1
gives algebraic lanes a chance, Phase 2 keeps or retires inherited lanes
using raw geometry, and Phase 3 receives only the surviving records.

\subsection{11.4 Phase 3 reduced-window rerank and post-shortlist
admission}\label{phase-3-reduced-window-rerank-and-post-shortlist-admission}

This section refines the surviving Phase-2 records with the more
expensive reduced-window geometry, then uses that final rerank surface
for shortlist formation, batched admission, pruning, and continuation.

\subsubsection{11.4.1 Reduced-window geometry, curvature, and
trust-region
gain}\label{reduced-window-geometry-curvature-and-trust-region-gain}

Phase 3 starts from the Phase-II raw quantities \(\tau_r\), \(g_r\),
\(h_r\), \(g_{\mathrm{lcb}}(r)\), \(F_{\mathrm{raw}}(r)\), and
\(\mathcal N_2(r;t)\), then refines them using the inherited refit
window. The raw span-novelty score \(\mathcal N_2(r;t)\) is not
recomputed here; instead the candidate direction is reduced against the
inherited window before the expensive rerank surface is formed.

\subsubsection{11.4.2 Refit window, reduced geometry, and Phase-3 record
family}\label{refit-window-reduced-geometry-and-phase-3-record-family}

In the record-inheritance mode used here, the Phase-3 input set is
\(\mathcal R_3(t):=\mathcal S_2(t)\). The inherited refit window
attached to position \(p\) at selector step \(t\) is denoted \(W(p;t)\),
with decomposition
\(W(p;t)=W_{\mathrm{new}}^{(p)}\cup W_{\mathrm{age}}^{(p)}\), where
\(W_{\mathrm{new}}^{(p)}\) keeps the newest coordinates and
\(W_{\mathrm{age}}^{(p)}\) preserves the maturity-scaled set of oldest
retained old coordinates defined later by the branch-local admission
order. For each shortlisted record \(r=(m,p)\), let \(W_r:=W(r;t)\) be
the inherited window that would actually be refit if \(m\) were admitted
at position \(p\), and let \(\{\tau_j\}_{j\in W_r}\) be the current
horizontal tangents in that inherited window. Phase 3 evaluates its
inherited reduced-window geometry under the later-phase controller shot
law \(N_{\mathrm{shot},3}(t)\).

Using the Phase-II raw tangent data, define
\((Q_r)_{jk}=\Re\langle \tau_j,\tau_k\rangle\),
\((q_r)_j=\Re\langle \tau_j,\tau_r\rangle\),
\(E_0:=E(\theta;\mathcal O)\),
\((b_r)_j:=\left.\partial_\alpha\partial_{\theta_j}E(\alpha,\delta\theta_{W_r};r)\right|_{\alpha=0,\,\delta\theta_{W_r}=0}\)
for \(j\in W_r\), and
\((H_{W_r})_{jk}:=\left.\partial_{\theta_j}\partial_{\theta_k}E(\alpha,\delta\theta_{W_r};r)\right|_{\alpha=0,\,\delta\theta_{W_r}=0}\)
for \(j,k\in W_r\). The local reduced-window energy model is \[
E(\alpha,\delta\theta_{W_r};r)
\approx
E_0+g_r\alpha+\tfrac12 h_r\alpha^2+\alpha b_r^\top\delta\theta_{W_r}+\tfrac12\delta\theta_{W_r}^\top H_{W_r}\delta\theta_{W_r}.
\] For reduction and reduced novelty,
\(\mathbf R_r=\tfrac12(H_{W_r}+H_{W_r}^{\top})+\lambda_H I\),
\(\widetilde h_r=h_r-b_r^\top \mathbf R_r^{-1}b_r\),
\(F_r^{\mathrm{red}}=F_{\mathrm{raw}}(r)-2q_r^\top \mathbf R_r^{-1}b_r+b_r^\top \mathbf R_r^{-1}Q_r\mathbf R_r^{-1}b_r\),
\(F_{r,\mathrm{safe}}^{\mathrm{red}}=\max\!\left(F_r^{\mathrm{red}},\varepsilon_{\mathrm{red}}\right)\)
with \(\varepsilon_{\mathrm{red}}>0\),
\(q_r^{\mathrm{red}}=q_r-Q_r\mathbf R_r^{-1}b_r\), and
\(z_r:=\frac{(q_r^{\mathrm{red}})^\top (Q_r+\varepsilon_{\mathrm{nov}}I)^{-1} q_r^{\mathrm{red}}}{F_{r,\mathrm{safe}}^{\mathrm{red}}}\),
\(\mathcal N_3(r;t):=\operatorname{clip}_{[0,1]}\!\left(1-z_r\right)\).
The reduced trust-region gain is \[
\Delta E_{\mathrm{TR}}(r)
:=
\max_{|\alpha|\le \rho/\sqrt{F_{r,\mathrm{safe}}^{\mathrm{red}}}}
\left[g_{\mathrm{lcb}}(r)|\alpha|-\tfrac12\max(\widetilde h_r,0)\alpha^2\right].
\] This reduced tangent is an energy-relaxed tangent, not a pure
orthogonal projection: the Hessian block \(\mathbf R_r\) encodes the
local inherited-window relaxation used to ask how the candidate would
look after nearby old coordinates are allowed to
adjust.\footnote{Equivalently, the reduction mixes tangent-space geometry with the local energetic relaxation model. That is the intended Phase-III semantics; pure span novelty is only an approximation to this energy-relaxed geometry.}

\subsubsection{11.4.3 Phase-3 rerank score and effective selector}

The authoritative Phase-3 rerank score is the reduced-window geometric
score alone: \[
S_3(r;t)
:=
\frac{\Delta E_{\mathrm{TR}}(r)\,\mathcal N_3(r;t)}{1+K_3(r;t)},
\qquad r\in\mathcal R_3(t).
\] Thus Phase 3 is geometry-authoritative: the reduced-window
lower-confidence trust-region gain, reduced novelty, and phase-local
burden determine the final rank. The inherited lane label
\(\ell_3(r;t):=\ell_2(r;t)\), commutation telemetry, operator-family
labels, and motif descriptors may be logged, may inform future
lane-budget policy, or may resolve deterministic ties after geometric
equality. They do not enter the canonical Phase-3 score as additive or
multiplicative rescue terms.

The retained Phase-3 shortlist is obtained from
\(\mathcal A_3(t):=\{r=(m,p)\in\mathcal R_3(t):S_3(r;t)\ge \tau_3(t)\}=\{r_{(1)},\dots,r_{(M)}\}\),
with \(M:=|\mathcal A_3(t)|\). Let \(s_k^{(3)}(t):=S_3(r_{(k)};t)\),
such that \(s^{(3)}_1(t)\ge s^{(3)}_2(t)\ge\cdots\ge s^{(3)}_M(t)\).
Rate the candidate's score distinction by
\(u_{\mathrm{front},k}^{(3)}(t):=\frac{s^{(3)}_{k+1}(t)+\varepsilon}{s^{(3)}_{k}(t)+\varepsilon}\),
with \(k=1,\dots,M-1\) and \(0<\eta_3<1\). Then \[
\begin{aligned}
k_3^\star(t)&:=
\begin{cases}
0, & M=0,\\[4pt]
\min\!\left\{k\in\{1,\dots,\min(M-1,N_3(t)-1)\}:u_{\mathrm{front},k}^{(3)}(t)\le \eta_3\right\},
& \text{if this set is nonempty},\\[6pt]
\min\{M,N_3(t)\}, & \text{otherwise},
\end{cases}\\
\mathcal D_3(t)&:=\{r_{(j)}:1\le j\le k_3^\star(t)\},
\qquad
\mathcal S_3(t):=\mathcal A_3(t)\cap \mathcal D_3(t)=\mathcal D_3(t).
\end{aligned}
\] Of course, \(\mathcal G^{(3)}:=\pi_m(\mathcal S_3(t))\), and
\(\mathcal P_m^{(3)}:=\{p:(m,p)\in\mathcal S_3(t)\}\).

\subsubsection{11.4.4 Reduced-plane geometric batch selection after
shortlisting}\label{reduced-plane-geometric-batch-selection-after-shortlisting}

After the Phase-3 shortlisting has produced the record family
\(\mathcal S_3(t)\subseteq \mathcal R_3(t)\), batching is resolved by
reduced-state-space geometry on a common inherited tangent plane, rather
than by heuristic pairwise compatibility penalties. Support-overlap,
commutation, and related proxy features may still be used as optional
cheap prescreens or route-level compatibility filters, but they are not
part of the principled batching law below unless the chosen route
explicitly declares them to be hard cross-record constraints. Assume
throughout that \(\mathcal S_3(t)\neq\varnothing\); otherwise no Phase-3
batch is formed.

Disjoint support plays a different role in batching than it plays in
Phase 1. The relation \(\Omega_{mn}=0\) is not a Phase-1 redundancy
signal and does not by itself justify a large exploratory shortlist
budget. However, after two records have survived Phase 2 and Phase 3
geometry, support-disjoint structure is favorable batching evidence: it
suggests lower structural conflict, weaker direct measurement overlap,
and a better chance that the two reduced directions can be co-admitted
without destructive local interference. Thus disjoint records receive
modest protected exposure before geometry, but may become preferred
low-conflict batch partners after geometry has validated them.

Near-degenerate shell and structural feasibility: let
\(s_t(r):=S_3(r;t)\) and
\(r_{\max}(t)\in \arg\max_{r\in \mathcal S_3(t)} s_t(r)\). Fix a
near-degenerate shell parameter \(\eta_{\mathrm{nd}}\in(0,1]\) and a
batch cap \(B_{\max}(t)\in\mathbb N\). Define the near-degenerate record
shell
\(\mathcal N_3(t):=\left\{ r\in \mathcal S_3(t): s_t(r)\ge \eta_{\mathrm{nd}}\,s_t\!\bigl(r_{\max}(t)\bigr) \right\}\)
and the corresponding near-degenerate batch family
\(\mathfrak B_{\mathrm{nd}}(t):=\left\{ B\subseteq \mathcal N_3(t): r_{\max}(t)\in B,\quad 1\le |B|\le B_{\max}(t) \right\}\).
Define a hard batch-structural gate
\(\Gamma_{\mathrm{batch}}^{\mathrm{struct}}(B;t)\in\{0,1\}\) to mean
that the records in \(B\) are jointly insertable on the current
branch/state, satisfy all route-declared hard cross-record constraints,
and determine a well-defined batched scaffold update
\(\mathcal O_t\oplus B\) in the fixed canonical insertion order. This is
an evaluability predicate for the route, not a theorem-level commutation
predicate: a conservative TETRIS-style route may instantiate it with
pairwise exact commutation and non-overlapping support, while the
route-neutral reduced-plane law only requires the proposed batch to be
well-defined before \(G_B(t)\), \(\widetilde H_B(t)\), and
\(\Delta E_B^{\mathrm{TR}}(t)\) are evaluated. The structurally feasible
batch family is
\(\mathfrak B_{\mathrm{struct}}(t):=\left\{ B\subseteq \mathcal S_3(t): \Gamma_{\mathrm{batch}}^{\mathrm{struct}}(B;t)=1 \right\}\).
A controller-side transient beam fallback may also be defined on the
same rerank
shell.\footnote{One convenient rule is to set $\mathcal T_{\mathrm{tie}}(t):=\{\,r\in\mathcal S_3(t): s_t(r)\ge \max(\eta_{\mathrm{tie}}\,s_t(r_{\max}(t)),\,s_t(r_{\max}(t))-\delta_{\mathrm{tie}})\,\}$. If $|\mathcal T_{\mathrm{tie}}(t)|>1$ and the controller elects not to resolve the shell by an immediate one-scaffold batch, it may retain a highest-scoring subset $\mathcal A_{\mathrm{tie}}(t)\subseteq \mathcal T_{\mathrm{tie}}(t)$ with $1\le |\mathcal A_{\mathrm{tie}}(t)|\le B_{\mathrm{tie}}(t)$ and pass one retained alternative per record into the scaffold-beam machinery of \S 11.5.2.}

Common-window reduced geometry: write the current state as
\(|\psi_t\rangle:=|\psi(\mathcal O_t,\theta_t)\rangle, \qquad Q_{\psi_t}:=I-|\psi_t\rangle\langle \psi_t|\),
where \(Q_{\psi_t}\) is the horizontal projector. For any batch
\(B\in \mathfrak B_{\mathrm{struct}}(t)\), define the common inherited
refit window \(W_B(t):=\bigcup_{r\in B} W_r(t)\) and the corresponding
inherited horizontal tangent matrix
\(T_B(t):=[\,t_j(t)\,]_{j\in W_B(t)}, \qquad t_j(t):=Q_{\psi_t}\,\partial_{\theta_j}|\psi_t\rangle\).
Let \(E_t(\alpha,\delta\theta_{W_B};B)\) denote the local energy
obtained by inserting the records in \(B\) with batch amplitudes
\(\alpha=(\alpha_r)_{r\in B}\) and varying the inherited coordinates on
\(W_B(t)\) by \(\delta\theta_{W_B}\). For each \(r\in B\), let
\(|\psi_{t,r}(\alpha_r)\rangle\) denote the state obtained by inserting
only record \(r\) at amplitude \(\alpha_r\) into the current scaffold
while keeping inherited coordinates fixed. Its zero-initialized
horizontal candidate tangent is
\(d_r(t):=Q_{\psi_t}\,\partial_{\alpha_r}|\psi_{t,r}(\alpha_r)\rangle\Big|_{\alpha_r=0}\).
Let \(\mathsf H_{W_B}(t)\) denote the inherited-window Hessian block of
\(E_t(\alpha,\delta\theta_{W_B};B)\) with respect to the coordinates on
\(W_B(t)\) at \((\alpha,\delta\theta_{W_B})=(0,0)\), and set
\(M_B(t):=\frac12\Bigl(\mathsf H_{W_B}(t)+\mathsf H_{W_B}(t)^\top\Bigr)+\lambda_H I, \qquad \lambda_H>0.\)
For each \(r\in B\), define the mixed candidate-window curvature vector
by
\(\bigl(b_r^{(B)}(t)\bigr)_j:=\left.\partial_{\alpha_r}\partial_{\theta_j} E_t(\alpha,\delta\theta_{W_B};B)\right|_{(\alpha,\delta\theta_{W_B})=(0,0)}, \qquad j\in W_B(t)\).
The reduced candidate tangent after linear relaxation of the common
inherited window is then
\(u_r^{(B)}(t):=d_r(t)-T_B(t)\,M_B(t)^{-1}b_r^{(B)}(t)\), hence the
effective batch plane is
\(\Pi_B(t):=\operatorname{span}\{u_r^{(B)}(t):r\in B\}\subseteq T_{\psi_t}\mathbb{P}(\mathcal H)\).
Its reduced Gram matrix is
\(\bigl(G_B(t)\bigr)_{rs}:=\Re\langle u_r^{(B)}(t),u_s^{(B)}(t)\rangle, \qquad r,s\in B\),
and in particular \(\bigl(G_B(t)\bigr)_{rr}=F_r^{\mathrm{red},(B)}(t)\)
is the reduced norm of record \(r\) in the common batch reduction. Next
define the symmetrized raw candidate-candidate curvature on the same
common reduction by
\(h_{rs}^{(B)}(t):=\frac12 \left. \bigl( \partial_{\alpha_r}\partial_{\alpha_s} + \partial_{\alpha_s}\partial_{\alpha_r} \bigr) E_t(\alpha,\delta\theta_{W_B};B) \right|_{(\alpha,\delta\theta_{W_B})=(0,0)}, \qquad r,s\in B,\)
and then the reduced curvature matrix
\(\bigl(\widetilde H_B(t)\bigr)_{rs}:=h_{rs}^{(B)}(t)-\bigl(b_r^{(B)}(t)\bigr)^\top M_B(t)^{-1} b_s^{(B)}(t)\).
For \(r=s\), the entry \(\bigl(\widetilde H_B(t)\bigr)_{rr}\) is the
usual single-record reduced curvature in the common batch window; for
\(r\neq s\), the entry \(\bigl(\widetilde H_B(t)\bigr)_{rs}\) is the
mixed second-order coupling after inherited-window relaxation. The
diagnostic coherences are
\(\mu_{\mathrm{tan}}(B;t):=\max_{r\neq s\in B} \frac{|(G_B(t))_{rs}|} {\sqrt{(G_B(t))_{rr}(G_B(t))_{ss}}+\varepsilon}\)
and
\(\mu_{\mathrm{curv}}(B;t):=\max_{r\neq s\in B} \frac{|(\widetilde H_B(t))_{rs}|} {\sqrt{\bigl((\widetilde H_B(t))_{rr}\bigr)^{(+)} \bigl((\widetilde H_B(t))_{ss}\bigr)^{(+)} }+\varepsilon}, \qquad x^{(+)}:=\max\{x,0\}\);
these may be used as optional early pruning diagnostics, but the actual
batch objective is the joint reduced trust-region gain below.

Joint reduced gain and additivity: let
\(g_{B,\mathrm{lcb}}(t):=\bigl(g_{\mathrm{lcb}}(r;t)\bigr)_{r\in B}\),
and let \(\bigl(\widetilde H_B(t)\bigr)^{(+)}\) denote the
positive-semidefinite part of \(\widetilde H_B(t)\). The batch
trust-region objective below is therefore a conservative convexified
local gain model on the reduced batch plane: \[
\Delta E_B^{\mathrm{TR}}(t)
:=
\max_{\alpha\in\mathbb R^{|B|},\ \alpha^\top G_B(t)\alpha\le \rho_B(t)^2}
\left[
g_{B,\mathrm{lcb}}(t)^\top \alpha
-
\frac12 \alpha^\top \bigl(\widetilde H_B(t)\bigr)^{(+)}\alpha
\right].
\] Here \(\rho_B(t)>0\) is the Phase-3 batch trust-region radius. If
record identity already includes the favorable orientation, one may
additionally impose \(\alpha\ge 0\) componentwise. For each \(r\in B\),
define the contextual single-record gain in the same common reduction by
\(\Delta E_{r\mid B}^{\mathrm{TR}}(t):=\max_{\beta\in\mathbb R,\ \beta^2 (G_B(t))_{rr}\le \rho_B(t)^2} \left[ g_{\mathrm{lcb}}(r;t)\,\beta - \frac12 \bigl((\widetilde H_B(t))_{rr}\bigr)^{(+)}\beta^2 \right].\)
Then the additivity defect is
\(\delta_{\mathrm{add}}(B;t):=\left[ 1- \frac{\Delta E_B^{\mathrm{TR}}(t)} {\sum_{r\in B}\Delta E_{r\mid B}^{\mathrm{TR}}(t)+\varepsilon} \right]_{+}.\)
Small \(\delta_{\mathrm{add}}(B;t)\) means that the joint reduced model
is approximately additive on the common batch plane.

Batch choice and greedy implementation: the geometric admissible batch
family is
\(\mathfrak B_{\mathrm{geom}}(t):=\left\{ B\in \mathfrak B_{\mathrm{nd}}(t)\cap \mathfrak B_{\mathrm{struct}}(t): \lambda_{\min}\!\bigl(G_B(t)\bigr) \ge \tau_{\mathrm{rank}}\, \frac{\operatorname{tr}G_B(t)}{|B|}, \quad \delta_{\mathrm{add}}(B;t)\le \tau_{\mathrm{add}} \right\}\).
When a more aggressive pruning rule is desired, one may additionally
require
\(\mu_{\mathrm{tan}}(B;t)\le \tau_{\mathrm{tan}}, \qquad \mu_{\mathrm{curv}}(B;t)\le \tau_{\mathrm{curv}}.\)
The principled post-shortlist batch choice is \[
B_t^\star
\in
\arg\max_{B\in \mathfrak B_{\mathrm{geom}}(t)}
\Delta E_B^{\mathrm{TR}}(t).
\] If \(\mathfrak B_{\mathrm{geom}}(t)=\varnothing\), no Phase-3 batch
is admitted from the current shortlist. A greedy implementation avoids
exhaustive search: initialize \(B_t^{(1)}:=\{r_{\max}(t)\}\). Given a
current batch \(B_t^{(k)}\), define the admissible add-on family
\(\mathcal A_k(t):=\left\{ r\in \mathcal N_3(t)\setminus B_t^{(k)}: B_t^{(k)}\cup\{r\}\in \mathfrak B_{\mathrm{nd}}(t)\cap \mathfrak B_{\mathrm{struct}}(t) \right\}\).
For each \(r\in \mathcal A_k(t)\), define the marginal joint gain
\(\Delta_{\mathrm{marg}}(r\mid B_t^{(k)};t):=\Delta E_{B_t^{(k)}\cup\{r\}}^{\mathrm{TR}}(t)-\Delta E_{B_t^{(k)}}^{\mathrm{TR}}(t).\)
Choose
\(r_k^{\mathrm{best}}(t)\in \arg\max_{r\in \mathcal A_k(t)} \Delta_{\mathrm{marg}}(r\mid B_t^{(k)};t)\)
subject to
\(\lambda_{\min}\!\bigl(G_{B_t^{(k)}\cup\{r\}}(t)\bigr) \ge \tau_{\mathrm{rank}}\, \frac{\operatorname{tr}G_{B_t^{(k)}\cup\{r\}}(t)}{|B_t^{(k)}|+1}, \qquad \delta_{\mathrm{add}}(B_t^{(k)}\cup\{r\};t)\le \tau_{\mathrm{add}},\)
and, when enabled,
\(\mu_{\mathrm{tan}}(B_t^{(k)}\cup\{r\};t)\le \tau_{\mathrm{tan}}, \qquad \mu_{\mathrm{curv}}(B_t^{(k)}\cup\{r\};t)\le \tau_{\mathrm{curv}}.\)
If no such \(r\) exists, or if the best feasible marginal gain is
nonpositive, stop and set \(B_t^\star:=B_t^{(k)}\). Otherwise append the
best add-on, \(B_t^{(k+1)}:=B_t^{(k)}\cup\{r_k^{\mathrm{best}}(t)\}\),
and continue until no further feasible positive-gain append exists or
until \(|B_t^{(k)}|=B_{\max}(t)\). The deployed Phase-3 action is the
final batch \(B_t^\star\). If \(|B_t^\star|=1\), this reduces to the
usual single-record admission. If \(|B_t^\star|>1\), the controller
inserts the batch on the current branch in the canonical batch order and
performs the common-window refit on \(W_{B_t^\star}(t)\).

\subsubsection{11.4.5 Planned size-aware Hamiltonian-agnostic Phase-3
policy
optimization}\label{planned-size-aware-hamiltonian-agnostic-phase-3-policy-optimization}

The planned policy-optimization layer treats the Phase-3 selector above
as a controlled decision system, not as a Hamiltonian-label lookup
table. Let \(h\in\mathfrak H\) denote a Hamiltonian instance, \(L\) its
size, \(\mathcal G_h(L)\) the resolved candidate-generator pool, and
\(\xi_h\) a transferable descriptor vector containing scale-normalized
couplings, graph/locality summaries, sector data, term-family counts,
operator-family features, and algebraic graph descriptors. The algebraic
portion of \(\xi_h\) includes local commutation degree histograms, local
noncommutation degree histograms, commuting-cluster size distribution,
noncommuting-bridge distribution, flat-cluster saturation statistics,
support-overlap graph statistics, and per-family commutation profiles.
For a fixed pair \((h,L)\), the instance oracle policy is \[
\phi^\star_{h,L}
\in
\arg\min_{\phi\in\Phi_L} J_{h,L}(\phi),
\qquad
J_{h,L}(\phi):=
\mathbb E_{\omega\sim\Omega_{h,L}}\!\left[\ell_{h,L,\omega}(\phi)\right]
+\lambda_{\mathrm{fail}}\,\mathbb P_{\omega}\!\left(F_{h,L,\omega}(\phi)\right),
\] where \(\omega\) ranges over seeds, shot/noise draws, and admissible
initialization choices; \(\ell\) is the terminal cost such as final
energy error plus normalized measurement and compile burden; and \(F\)
is a run failure event such as nonconvergence, invalid scaffold growth,
or violation of a hard controller contract. This oracle is an analysis
object: it is allowed to specialize to the held pair \((h,L)\), so it
gives the best per-instance policy one could hope to match.

For heavily bosonic benchmark families, the small-system training suite
should not use a single-excitation oscillator cutoff as the only
evidence surface. Let \(n_{b}^{\max}\) denote the local boson or phonon
cutoff. The default bosonic stress lane uses \(n_{b}^{\max}=2\) for the
pure boson-chain and spin-boson families, while retaining cheaper
\(n_{b}^{\max}=1\) lanes where they are intentionally used as smoke
tests or HH-specific historical baselines. This prevents the outer
policy from fitting selector weights to an artificially two-level
oscillator and exposes whether shortlisting, shot-cost penalties,
insertion, beam, and rollback still behave when each bosonic site has
three local Fock levels.

The same suite should include the closed-shell molecular route as an
electronic anchor, even though its Hamiltonian construction uses a
molecular-integral JSON rather than a lattice-coupling tuple. The
molecular lane uses explicit roles: H2/STO-3G remains the lightweight
control, LiH/STO-3G is the primary molecular train anchor, and
H2O/STO-3G remains the transfer/stress molecule. Each molecular
benchmark samples the same selector policy vector \(\phi\), so this lane
tests whether the learned shortlist, cost, beam, insertion, and rollback
weights transfer beyond lattice-generated operator families while
keeping the legacy H2/H2O benchmark identities stable.

The deployable target is instead a single Hamiltonian-agnostic,
size-aware map \[
\phi_{\mathrm{agn}}:\ (L,\xi_h)\longmapsto
\Bigl(\tau_k,\eta_k,N_k,B_{\max},\rho_B,\lambda_H,
\{B_{1,\ell},B_{2,\ell},c_\ell^{\mathrm{abs}},c_\ell^{\mathrm{rel}},
\eta_{1,\ell},\eta_{2,\ell}^{\mathrm{in}}\}_{\ell\in\mathfrak L},
\beta_{\mathrm{comm\_recover}},\beta_{\mathrm{curv\_recover}},
\beta_{\mathrm{bridge\_risk}},
\text{family lane quotas}\Bigr)_{k=1}^{3},
\] which emits Phase-1/2/3 thresholds, lane budgets, lane-survival
patience factors, batch caps, trust radii, regularizers, and
prune-recovery risk weights before the online selector is evaluated.
Algebraic descriptors in \(\xi_h\) may tune lane exposure, lane
patience, and prune-recovery caution; they do not create score bonuses
for \(S_1\), \(S_2\), or \(S_3\). Its population objective should prefer
good average behavior while protecting the worst practical cases: \[
\begin{aligned}
\mathcal J_{\mathrm{agn}}(\phi_{\mathrm{agn}})
&:=
\mathbb E_{(h,L)\sim\mathcal D}\left[
\mathbb E_{\omega}\,\ell_{h,L,\omega}\!\left(\phi_{\mathrm{agn}}(L,\xi_h)\right)
\right] \\
&\quad+
\lambda_{\mathrm{cvar}}\,
\operatorname{CVaR}_{\alpha}\!\left(
\ell_{h,L,\omega}\!\left(\phi_{\mathrm{agn}}(L,\xi_h)\right)
\right)
+\lambda_{\mathrm{rob}}\,
\sup_{(h,L)\in\mathcal U}\left[
J_{h,L}\!\left(\phi_{\mathrm{agn}}(L,\xi_h)\right)-J^{\mathrm{base}}_{h,L}
\right]_{+} \\
&\quad+
\lambda_{\mathrm{fail}}\,
\mathbb P_{(h,L),\omega}\!\left(
F_{h,L,\omega}\!\left(\phi_{\mathrm{agn}}(L,\xi_h)\right)
\right).
\end{aligned}
\] The first term optimizes mean performance; the CVaR term penalizes
the high-loss tail; the robust term prevents a small adversarial or
withheld family \(\mathcal U\) from being sacrificed; and the failure
term keeps apparently cheap policies from winning by producing brittle
or invalid runs.

The natural measure of success is the oracle gap \[
\Delta_{\mathrm{oracle}}(h,L;\phi_{\mathrm{agn}})
:=
J_{h,L}\!\left(\phi_{\mathrm{agn}}(L,\xi_h)\right)
-
J_{h,L}(\phi^\star_{h,L})
\ge 0,
\] reported both as a mean over \(\mathcal D\) and as a tail statistic
over withheld Hamiltonians and larger sizes. A small gap means the
agnostic policy has learned size and family regularities rather than
memorizing a particular Hamiltonian.

Shortlist and pool budgets should scale explicitly with both size and
available generator count. A compact family of caps is \[
N_k(L,|\mathcal G_h(L)|;\zeta_k)
:=
\min\!\left\{
|\mathcal G_h(L)|,
\left\lceil a_k L^{p_k}\bigl[\log(1+|\mathcal G_h(L)|)\bigr]^{q_k}\right\rceil
\right\},
\qquad k\in\{1,2,3\},
\] with learned or scheduled \(\zeta_k=(a_k,p_k,q_k)\), constrained so
that \(N_3\le N_2\le N_1\) after rounding. Thus larger systems may
receive larger absolute candidate budgets, while the fraction of the
full pool inspected can still shrink when \(|\mathcal G_h(L)|\) grows
faster than the useful front.

These shortlists are intended to be active control surfaces, not
disabled compatibility knobs. A concrete tuning layer should therefore
tune \(N_1,N_2\), the Phase-2 shortlist fraction, the insertion-position
probe cap \(M_{\mathrm{probe}}\), and the cost weights that trade energy
proxy against compile and measurement burden. The goal is to reduce shot
and compile cost without collapsing expressivity; broad-pool
expressivity is preserved by scaling the budgets with
\((L,|\mathcal G_h(L)|)\) and by retaining beam, insertion, rollback,
batch, and prune alternatives at the controller level.

Operator-family controls enter only as transferable lane-budget priors
and quotas. If \(f(r)\) is the abstract family of record \(r\), the
policy may emit smooth family-lane quotas \(q_{\ell,f}(L,\xi_h,t)\) and
family-lane exposure weights \(\pi_{\ell,f}(L,\xi_h,t)\). These
quantities can change how many Phase-1 or Phase-2 seats a family
receives inside a lane, for example by constraining \[
\left|\{r\in\mathcal S_{1,\ell}(t):f(r)=f\}\right|
\le q_{\ell,f}(L,\xi_h,t),
\] or by adjusting \(B_{1,\ell}(t)\) and \(B_{2,\ell}(t)\) before
records are ranked. They may say that displacement-like, hopping-like,
squeeze-like, or Hubbard-like families tend to deserve more or less
exposure in a size/coupling regime. They do not alter the geometric
score of any record and must not become exact label masks such as
\(\mathbf 1\{m\in\mathcal M^{\mathrm{best}}_{h,L}\}\), because such
masks would encode the answer for one Hamiltonian instead of a
transferable policy.

The final validation should freeze \(\phi_{\mathrm{agn}}\) after
training or tuning on smaller and mixed-family surfaces, then evaluate
it on held-out \(L=4\) and larger Hamiltonians. The report should
compare mean loss, CVaR loss, failure rate, budget use
\(N_k(L,|\mathcal G|)\), and oracle gap against the per-\((h,L)\) oracle
\(\phi^\star_{h,L}\), a size-blind baseline, and any label-specific
baseline. Passing this validation means the policy scales through
\(L=4+\) by using size-aware and family-transferable structure, not by
recognizing exact operator labels from the training set.

\subsection{\texorpdfstring{11.5 Post \(m_p\) Application Entails
Windowed-Refitting}{11.5 Post m\_p Application Entails Windowed-Refitting}}\label{post-m_p-application-entails-windowed-refitting}

In ADAPT-VQE the variational problem is
\(\min_{\boldsymbol\theta} E_L(\boldsymbol\theta)\) with
\(E_L(\boldsymbol\theta)=\langle\phi_0|(\prod_{\ell=L}^{1}e^{i\theta_\ell G_\ell})H(\prod_{\ell=1}^{L}e^{-i\theta_\ell G_\ell})|\phi_0\rangle\),
and the measurement bottleneck arises because one iteration must
estimate not only \(E_L\) but also pool-selection scores
\(s_\mu(\boldsymbol\theta)=|\langle\phi_0|(\prod_{\ell=L}^{1}e^{i\theta_\ell G_\ell}),i[H,A_\mu],(\prod_{\ell=1}^{L}e^{-i\theta_\ell G_\ell})|\phi_0\rangle|\)
for \(\mu\in\mathcal P\), reoptimization gradients
\(\partial_{\theta_j}E_L(\boldsymbol\theta)\), and, for geometry-aware
updates, metric entries
\(M_{jk}(\boldsymbol\theta)=\Re\langle \partial_{\theta_j}[(\prod_{\ell=1}^{L}e^{-i\theta_\ell G_\ell})|\phi_0\rangle]|\partial_{\theta_k}[(\prod_{\ell=1}^{L}e^{-i\theta_\ell G_\ell})|\phi_0\rangle]\rangle\),
so at precisions \(\epsilon_s,\epsilon_g,\epsilon_M\) the shot cost
obeys
\(N_{\mathrm{iter}}\gtrsim \sum_{\mu\in\mathcal P}\mathrm{Var}(\hat s_\mu)/\epsilon_s^2+\sum_{j=1}^{L}\mathrm{Var}(\widehat{\partial_{\theta_j}E_L})/\epsilon_g^2+\sum_{j,k=1}^{L}\mathrm{Var}(\hat M_{jk})/\epsilon_M^2\);
the proposed strategy is to choose disjoint sets
\(F,A\subset{1,\dots,L}\) with \(F\cup A={1,\dots,L}\), impose
\(\dot\theta_j=0\) for \(j\in F\), and optimize only \(j\in A\), which
reduces the reparameterization sector to
\(\sum_{j\in A}\mathrm{Var}(\widehat{\partial_{\theta_j}E_L})/\epsilon_g^2+\sum_{j,k\in A}\mathrm{Var}(\hat M_{jk})/\epsilon_M^2\)
and therefore removes the growth coming from repeatedly updating old
coordinates, but it leaves the pool-screening term
\(\sum_{\mu\in\mathcal P}\mathrm{Var}(\hat s_\mu)/\epsilon_s^2\)
unchanged and it does not shorten the physical circuit, since every shot
still executes the full depth-\(L\) scaffold with entangling-gate count
\(N_{2q}(L)\) and runtime \(t_{\mathrm{circ}}(L)\), so hardware fidelity
still degrades roughly as
\(e^{-t_{\mathrm{circ}}(L)/T_1}e^{-t_{\mathrm{circ}}(L)/T_2}(1-\varepsilon_{2q})^{N_{2q}(L)}\);
consequently the freezing strategy cures only the reoptimization
component of the ADAPT bottleneck, and the limiting cost after that is
either the unchanged operator-pool screening overhead or, if screening
is separately controlled, the conventional VQE bottleneck of resolving
\(E_L\) on a deep noisy ansatz, with shots scaling as
\(N_E\gtrsim \mathrm{Var}(H_{\mathrm{meas}}\mid \prod_{\ell=1}^{L}e^{-i\theta_\ell G_\ell},\phi_0,\text{noise})/(\Delta E)^2\)
to distinguish an energy change \(\Delta E\).

Let the length-\((m)\) scaffold be
\(U^{(m)}(\boldsymbol{\theta})=\prod_{j=1}^{m} e^{-i\theta_j G_j}\),
where \((G_j)\) is the \((j)\)-th appended generator and \((\theta_j)\)
its coefficient. The reparameterization motif is to optimize only an
active subset \((\mathcal A_m\subseteq {1,\dots,m})\), chosen by an
activity functional
\(a_j^{(m)}:=\lambda_{\mathrm{rec}}\,r_j^{(m)}+\lambda_{\mathrm{mob}}\,v_j^{(m)}-\lambda_{\mathrm{age}}\,\phi(m-j)-\lambda_{\mathrm{stag}}\,s_j^{(m)}\),
and then set \((\mathcal A_m=\{j: a_j^{(m)}>\tau_m\})\). Here
\((r_j^{(m)})\) is a recency weight, large for recently appended
generators;
\((v_j^{(m)}:=\sum_{k=j}^{m-1}\lvert \Delta\theta_j^{(k)}\rvert)\) is a
lifetime mobility measure, large for parameters that have historically
moved substantially; \((\phi(m-j))\) is an age penalty, biasing against
early-added generators; and \((s_j^{(m)})\) is a stagnation indicator,
large when a parameter has remained nearly inert over recent refits.
Thus the inclusion bias is toward recent and historically labile
generators, while the exclusion bias is toward old and typically
stagnant generators.

This is the same scaffold language used above:
\(U^{(m)}(\boldsymbol\theta)\) is the scaffold \(\mathcal O\), the
active set \(\mathcal A_m\) plays the role of the inherited refit window
\(W(r;t)\), the frozen set is \(\{1,\ldots,n+1\}\setminus W(r;t)\), and
the activity functional \(a_j^{(m)}\) is the same maturity-facing
bookkeeping later summarized by the eligible set \(\mathcal M_t\).

Concomitantly, as scaffold maturity increases, the controller shot laws
\(N_{\mathrm{shot},k}^{\mathrm{eff}}(t)\) should preserve decision-level
signal-to-noise ratio on the active block while respecting the live
phase mask. Formally, if \(\delta_{t,k}\) denotes a typical useful
signal scale among the phase-\(k\) unfrozen directions and
\(\sigma_{t,k}\) the corresponding estimator spread, then the capped
schedule-plus-SNR law above targets \[
\mathrm{SNR}_{t,k}\sim
\frac{\delta_{t,k}\sqrt{N_{\mathrm{shot},k}^{\mathrm{eff}}(t)}}{\sigma_{t,k}}
\ge \kappa_k
\] for every live phase \(k\), while a nulled phase has
\(N_{\mathrm{shot},k}^{\mathrm{eff}}(t)=0\). Hence the control law is:
as scaffold maturity grows, shrink the reoptimized subspace by freezing
archaic or stagnant directions, increase shots on the surviving live
active directions, and eventually null expensive geometric phases when
the runway forecast says the scaffold is mature enough.

\subsubsection{11.5.1 Branch-local ADAPT prune-recoverability ladder:
nomination, refit, and
rollback}\label{branch-local-adapt-prune-recoverability-ladder-nomination-refit-and-rollback}

All quantities in this subsection are branch-local. On a single active
branch we suppress the branch label. Let
\((\mathcal O_t^-,\theta_t^-,E_t^-)\) be the committed branch state
immediately before the accepted Phase-3 payload \(\mathcal A_t^\star\)
is applied, and let \((\mathcal O_t^+,\theta_t^+,E_t^+)\) be the
committed post-admission state after the ordinary local refit on that
same branch has completed. The admitted-step gain is \[
\Delta_{\mathrm{adm}}(t):=E_t^- - E_t^+,
\qquad
\mathcal J_t:=\{1,\dots,n_t^+\}.
\] Each committed coordinate \(j\in\mathcal J_t\) carries the
operational metadata \[
(m_j,\iota_j,r_j,\mathfrak c_j,\kappa_j,c_j(t),\{a_j(q)\}_{q\le t},
\operatorname{noncomm\_bridge\_score}_j,
\operatorname{recovery\_class}_j),
\] where \(m_j\) is the generator label, \(\iota_j\) is the branch-local
admission step, \(r_j\) is the admission record, \(\mathfrak c_j\) is
the generator-family or motif class, \(\kappa_j:=K_3(r_j;\iota_j)\) is
the stored selector burden, \(c_j(t)\) is the prune-cooldown counter,
\(a_j(q)\) is the transported normalized amplitude history,
\(\operatorname{noncomm\_bridge\_score}_j\) is the static bridge-risk
telemetry, and \(\operatorname{recovery\_class}_j\) is the auditable
prune outcome. Optional algebraic telemetry such as commuting-cluster
labels, local commutation degrees, and flat-cluster saturation may be
stored with the coordinate, but it is not part of the main deletion
authority below.

The normalized amplitude history is \[
a_j(q):=
\frac{\left|\operatorname{wrap}_{\mathfrak c_j}\theta_j(q)\right|}
{s_{\mathfrak c_j}+\varepsilon_s},
\] using the principal-angle wrap for periodic rotation coordinates and
the identity for nonperiodic coordinates. The recovery class is
initially unset, \[
\operatorname{recovery\_class}_j\in
\{\varnothing,\mathrm{flat\_redundant},\mathrm{curvature\_compensated},\mathrm{failed}\},
\] and is set only after a prune trial.

The target object is not whether coordinate \(j\) is small; it is
whether the scaffold with \(j\) deleted can recover the same variational
energy after the surviving coordinates compensate. With
\(\mathcal O_{t,-j}^+\) denoting the scaffold obtained by deleting
\(j\), define \[
\boxed{
\delta E_j^\star(t)
:=
\inf_{\xi\in\Theta_{\mathcal O_{t,-j}^+}}
E(\mathcal O_{t,-j}^+,\xi)
-
E(\mathcal O_t^+,\theta_t^+).
}
\] Small amplitude is only one sufficient-looking nomination signal, \[
|\theta_{t,j}|\approx0
\Longrightarrow
\delta E_j^\star(t)\approx0,
\qquad
\delta E_j^\star(t)\approx0
\not\Longrightarrow
|\theta_{t,j}|\approx0.
\] Thus prune pressure may increase opportunity, candidate breadth, and
remove-refit breadth, but it must not weaken final safety. The
branch-local prune chain is \[
\boxed{
\mathcal J_t
\to
\mathcal M_t
\to
\mathcal N_1(t)
\to
\{W_j^{[s]}(t)\}_{s=0}^{4}
\to
\Gamma_j^{\mathrm{proj},[s]}(t)
\to
\mathcal C_{\mathrm{trial}}(t)
\to
\mathsf D_j^{[s,N]}
\to
\operatorname{commit}/\operatorname{rollback}.
}
\] Here \(\mathcal M_t\) is the mature deletion-target universe,
\(\mathcal N_1(t)\) is the cheap deletion-target nomination set,
\(W_j^{[s]}(t)\) are surviving-coordinate refit windows for a fixed
target \(j\), \(\Gamma_j^{\mathrm{proj},[s]}\) is an optional cached
geometry screen, \(\mathcal C_{\mathrm{trial}}(t)\) is the capped set
sent to exact remove-refit, and \(\mathsf D_j^{[s,N]}\) is the
confidence-controlled decision. No amplitude score, commutation
relation, cached projection, or Schur surrogate deletes a coordinate by
itself.

\paragraph{Prune 0: permission, maturity, and protected
coordinates}\label{prune-0-permission-maturity-and-protected-coordinates}

Prune may open after an accepted admission, after a no-admit plateau, at
scheduled compression checkpoints, or at the terminal scaffold. Let
\(\mathcal T_{\mathrm{chk}}\subseteq\mathbb N\) be the scheduled
checkpoint set, let \(\operatorname{Plateau}_t\) mean that no useful
admission has survived the current controller patience window, and let
\(\operatorname{StableRefit}_t\) mean that the latest local refit
completed without optimizer or numerical failure. The permission gate is
\[
\begin{aligned}
\Gamma_{\mathrm{prune}}^{\mathrm{perm}}(t)
&:=
\mathbf 1[|\mathcal O_t^+|\ge 2]\,
\mathbf 1[\operatorname{StableRefit}_t]\,
\mathbf 1[u_{\mathrm{sat}}(t)\ge\tau_{\mathrm{med}}]\\
&\quad\times
\mathbf 1[
\text{accepted admission at }t
\ \vee\
\operatorname{Plateau}_t
\ \vee\
t\in\mathcal T_{\mathrm{chk}}
\ \vee\
t=t_{\mathrm{terminal}}
].
\end{aligned}
\] The scale used later for prune tolerance is \[
\Delta_{\mathrm{scr}}(t)
:=
\max\!\bigl(
\tau_{\mathrm{use},t},
\operatorname{EWMA}_{\mathrm{recent}}(\Delta_{\mathrm{adm}}),
\Delta_{\mathrm{num}}
\bigr).
\] A scalar prune-pressure signal may schedule candidate caps, window
fractions, and cached-screen thresholds. For example, with normalized
pressure \(\bar\lambda_{\mathrm{pr}}(t)\in[0,1]\), one may set \[
Q_{\mathrm{pr}}(t)
:=
\left\lceil
Q_{\mathrm{pr}}^{\min}+
(Q_{\mathrm{pr}}^{\max}-Q_{\mathrm{pr}}^{\min})\bar\lambda_{\mathrm{pr}}(t)
\right\rceil,
\] \[
\rho_s(t)
:=
\operatorname{clip}_{[0,1]}\!\left(
\rho_s^{\min}+(\rho_s^{\max}-\rho_s^{\min})\bar\lambda_{\mathrm{pr}}(t)
\right),
\qquad s=1,\dots,4,
\] and may tighten or relax \(\tau_F^{\mathrm{pr}}(t)\) and
\(\tau_{\mathrm{Schur}}^{\mathrm{pr}}(t)\) according to the scheduled
compression policy. This pressure changes what is worth testing; it does
not change the final deletion authority.

Newly admitted coordinates are protected for \(\ell_{\mathrm{protect}}\)
admitted steps, \[
\mathcal P_{\mathrm{protect}}(t)
:=
\{j\in\mathcal J_t:t-\iota_j<\ell_{\mathrm{protect}}\}.
\] Let \(\operatorname{age}_{j,t}:=t-\iota_j\), and let
\(s_j^{\mathrm{stag}}(t)\) be the branch-local coordinate stagnation
indicator, with larger values meaning less recent parameter motion. The
mature deletion-target universe is \[
\boxed{
\mathcal M_t
:=
\Bigl\{
 j\in\mathcal J_t:
\operatorname{age}_{j,t}\ge a_{\min},
\ j\notin\mathcal P_{\mathrm{protect}}(t),
\ c_j(t)=0
\Bigr\}.
}
\] Prune eligibility is therefore branch-local and conservative: old
enough, not recently admitted, and not under temporary cooldown.
Staleness affects cheap nomination through a normalized feature, not
through hard eligibility.

\paragraph{Prune 1: cheap deletion-target nomination and pairwise
compensator
evidence}\label{prune-1-cheap-deletion-target-nomination-and-pairwise-compensator-evidence}

Prune 1 has two outputs. First, it nominates deletion targets that are
cheap enough to test. Second, it records pairwise evidence that Prune 2
uses to build compensator windows. Let \(\mathcal H_t\) be the recent
time-history region of accepted or refit states used for empirical
co-motion estimates. This is a time-index region, not a coordinate refit
window.

Define activity, motion, and stationarity summaries \[
\widetilde A_j(t)
:=
\frac{A_j^{\mathrm{late}}(t)}{A_{\mathrm{ref}}+\varepsilon_A},
\qquad
\widetilde V_j(t)
:=
\frac{V_j^{\mathrm{late}}(t)}{V_{\mathrm{ref}}+\varepsilon_V},
\qquad
\widetilde G_j(t)
:=
\frac{|g_j^{\mathrm{late}}(t)|}{G_{\mathrm{ref}}+\varepsilon_G}.
\] Here \(A_j^{\mathrm{late}}\) is an EWMA or windowed mean of the
normalized amplitude history, \[
A_j^{\mathrm{late}}(t)
:=
\operatorname{EWMA}_{q\in\mathcal H_t} a_j(q),
\] \(V_j^{\mathrm{late}}\) is a recent parameter-motion scale, and
\(g_j^{\mathrm{late}}\) is a cached gradient or stationarity signal.
Define normalized burden, age, and staleness features by \[
\widetilde\kappa_j(t):=\frac{\kappa_j}{\kappa_{\mathrm{ref}}+\varepsilon_\kappa},
\qquad
z_j^{\mathrm{age}}(t):=\operatorname{norm}(\operatorname{age}_{j,t}),
\qquad
z_j^{\mathrm{stale}}(t):=\operatorname{norm}(s_j^{\mathrm{stag}}(t)).
\] For \(i\ne j\), define the pairwise historical co-motion score \[
\boxed{
C_{ij}^{\mathrm{hist}}(t)
:=
\max\!\left\{
|\operatorname{corr}_{\mathcal H_t}(\theta_j,\theta_i)|,
|\operatorname{corr}_{\mathcal H_t}(g_j,g_i)|,
|\operatorname{corr}_{\mathcal H_t}(\Delta\theta_j,\Delta\theta_i)|
\right\}.
}
\] If tangent data are current and cached, define the pairwise
tangent-overlap score \[
\boxed{
C_{ij}^{\mathrm{tan}}(t)
:=
\frac{
\left|\operatorname{Re}\langle\bar\tau_j,\bar\tau_i\rangle\right|
}{
\|\bar\tau_j\|_2\,\|\bar\tau_i\|_2+\varepsilon_T
},
}
\] where \(\bar\tau_i:=Q_{\psi_t}\partial_{\theta_i}|\psi_t\rangle\) is
the horizontal tangent at the current post-admission state. If tangent
data are unavailable or stale, set \(C_{ij}^{\mathrm{tan}}(t)=0\). The
combined pairwise compensability evidence is \[
\boxed{
E_{ij}^{\mathrm{corr}}(t)
:=
\max\!\left(C_{ij}^{\mathrm{hist}}(t),C_{ij}^{\mathrm{tan}}(t)\right).
}
\] The static algebraic metadata are the support-overlap bit
\(\Omega_{m_i m_j}\), the bare-commutation bit \(\mathsf C_{m_i m_j}\),
and the exactness tag \(\mathsf q_{m_i m_j}\). Define the pairwise
algebraic evidence bits \[
\boxed{
E_{ij}^{\mathrm{flat}}
:=
\mathbf 1[\Omega_{m_i m_j}=1]\,
\mathbf 1[\mathsf C_{m_i m_j}=1]\,
\mathbf 1[\mathsf q_{m_i m_j}=\mathrm{exact}],
}
\] \[
\boxed{
E_{ij}^{\mathrm{curv}}
:=
\mathbf 1[\Omega_{m_i m_j}=1]\,
\mathbf 1[\mathsf C_{m_i m_j}=0]\,
\mathbf 1[\mathsf q_{m_i m_j}=\mathrm{exact}],
}
\] and \[
\boxed{
E_{ij}^{\mathrm{disj}}
:=
\mathbf 1[\Omega_{m_i m_j}=0]\,
\mathbf 1[\mathsf C_{m_i m_j}=1]\,
\mathbf 1[\mathsf q_{m_i m_j}=\mathrm{exact}].
}
\] The disjoint commuting bit is telemetry or batching context only. It
is not flat-sector redundancy evidence.

Aggregate features for cheap deletion-target ranking are \[
C_j^{\mathrm{comp}}(t)
:=
\max_{i\in\mathcal J_t\setminus\{j\}}E_{ij}^{\mathrm{corr}}(t),
\] \[
C_j^{\mathrm{tan}}(t)
:=
\max_{i\in\mathcal J_t\setminus\{j\}}C_{ij}^{\mathrm{tan}}(t),
\qquad
C_j^{\mathrm{comm}}(t)
:=
\max_{i\in\mathcal J_t\setminus\{j\}}E_{ij}^{\mathrm{flat}}E_{ij}^{\mathrm{corr}}(t).
\] Let
\(\mathcal J_{\mathrm{down}}(j;t)\subseteq\mathcal J_t\setminus\{j\}\)
be the recent downstream or later-in-order scaffold coordinates used to
estimate whether deleting \(j\) would remove a noncommuting bridge.
Define \[
\boxed{
B_j^{\mathrm{bridge}}(t)
:=
\operatorname{clip}_{[0,1]}\!\left(
\operatorname{noncomm\_bridge\_score}_j
+
\frac{
\sum_{i\in\mathcal J_{\mathrm{down}}(j;t)}E_{ij}^{\mathrm{curv}}
}{
|\mathcal J_{\mathrm{down}}(j;t)|+\varepsilon
}
\right).
}
\] Here \(C_j^{\mathrm{comp}}\), \(C_j^{\mathrm{tan}}\), and
\(C_j^{\mathrm{comm}}\) make deletion easier to nominate, while
\(B_j^{\mathrm{bridge}}\) protects bridge-like targets by raising their
cheap deletion score. The bridge feature is a risk feature for the
target \(j\), not a compensator-window definition.

The cheap recoverability-prior score is \[
\boxed{
\begin{aligned}
R_j^{\mathrm{cheap}}(t)
:=
&\ w_A\widetilde A_j
+w_V\widetilde V_j
+w_G\widetilde G_j
-w_C C_j^{\mathrm{comp}}
-w_T C_j^{\mathrm{tan}}
-w_{\mathrm{comm}}C_j^{\mathrm{comm}}\\
&+w_{\mathrm{bridge}}B_j^{\mathrm{bridge}}
-w_{\kappa}\widetilde\kappa_j
-w_{\mathrm{age}}z_j^{\mathrm{age}}
-w_{\mathrm{stale}}z_j^{\mathrm{stale}}.
\end{aligned}
}
\] Low \(R_j^{\mathrm{cheap}}\) means
\texttt{try\ testing\ this\ target,\textquotesingle{}\textquotesingle{}\ not}delete
this target.'\,' Define the cheap-score order by \[
j_1\preceq_{\mathrm{cheap}} j_2
\Longleftrightarrow
R_{j_1}^{\mathrm{cheap}}(t)\le R_{j_2}^{\mathrm{cheap}}(t),
\] with deterministic ties resolved by the fixed scaffold order.

The pure amplitude/activity-collapse lane is \[
\boxed{
\mathcal P_{\mathrm{amp}}(t)
:=
\left\{
 j\in\mathcal M_t:
\widetilde A_j(t)\le\tau_A^{\mathrm{pr}}(t),
\ \widetilde V_j(t)\le\tau_V^{\mathrm{pr}}(t),
\ \widetilde G_j(t)\le\tau_G^{\mathrm{pr}}(t)
\right\}.
}
\] The broader cheap recoverability-prior lane is \[
\boxed{
\mathcal P_{\mathrm{cheap}}(t)
:=
\operatorname{Bottom}_{Q_{\mathrm{pr}}(t)}
\left(
R_j^{\mathrm{cheap}}(t);
\ j\in\mathcal M_t
\right).
}
\] Prune 1 ends with the cheap deletion-target nomination set \[
\boxed{
\mathcal N_1(t)
:=
\mathcal P_{\mathrm{amp}}(t)
\cup
\mathcal P_{\mathrm{cheap}}(t).
}
\] The same pairwise evidence used to rank cheap nominees now constructs
the surviving-coordinate compensator windows.

\paragraph{Prune 2: typed compensator windows and cached Fubini--Schur
screens}\label{prune-2-typed-compensator-windows-and-cached-fubinischur-screens}

For each mature coordinate \(j\in\mathcal M_t\), define the
surviving-coordinate universe \[
I_{-j}:=\mathcal J_t\setminus\{j\},
\] identified with the coordinates of \(\mathcal O_{t,-j}^+\) after the
obvious index shift. A compensator window is a subset of \(I_{-j}\); it
contains surviving amplitudes allowed to refit after \(j\) is removed.
These windows are not nomination lanes.

The typed compensator pools are constructed from the Prune-1 pairwise
evidence: \[
\boxed{
W_{\mathrm{comm}}(j;t)
:=
\{i\in I_{-j}:E_{ij}^{\mathrm{flat}}=1\},
}
\] \[
\boxed{
W_{\mathrm{corr}}(j;t)
:=
\{i\in I_{-j}:E_{ij}^{\mathrm{corr}}(t)\ge\tau_{\mathrm{corr}}^{\mathrm{pr}}(t)\},
}
\] \[
\boxed{
W_{\mathrm{nc}}(j;t)
:=
\{i\in I_{-j}:E_{ij}^{\mathrm{curv}}=1\},
}
\] and \[
\boxed{
W_{\mathrm{age}}(j;t)
:=
\{i\in I_{-j}:\operatorname{age}_{i,t}\ge a_{\mathrm{comp}},\ i\notin\mathcal P_{\mathrm{protect}}(t)\}.
}
\] Thus \(W_{\mathrm{comm}}\) is the exact overlapping commuting pool,
\(W_{\mathrm{corr}}\) is the empirical or tangent-correlated pool,
\(W_{\mathrm{nc}}\) is the overlapping noncommuting pool, and
\(W_{\mathrm{age}}\) is the broad old-coordinate pool. Support-disjoint
commuting coordinates do not enter \(W_{\mathrm{comm}}\).

For a pool \(P(j;t)\subseteq I_{-j}\), write \(T_\rho(P,p)\) for the top
\(\lceil \rho |P|\rceil\) coordinates under priority score
\(p_{ij}(t)\), with deterministic ties. A conservative choice is \[
p_{ij}^{\mathrm{comm}}(t):=E_{ij}^{\mathrm{flat}}\bigl(\varepsilon_{\mathrm{comm}}+E_{ij}^{\mathrm{corr}}(t)\bigr),
\qquad
p_{ij}^{\mathrm{corr}}(t):=E_{ij}^{\mathrm{corr}}(t),
\] \[
p_{ij}^{\mathrm{nc}}(t):=E_{ij}^{\mathrm{curv}}E_{ij}^{\mathrm{corr}}(t),
\qquad
p_i^{\mathrm{age}}(t):=z_i^{\mathrm{age}}(t)+z_i^{\mathrm{stale}}(t).
\] Define truncated pools, for example, \[
W_{\mathrm{comm}}^{\rho}(j;t):=T_\rho(W_{\mathrm{comm}},p^{\mathrm{comm}}),
\] and analogously for \(W_{\mathrm{corr}}^\rho\),
\(W_{\mathrm{nc}}^\rho\), and \(W_{\mathrm{age}}^\rho\).

The nested refit ladder is recursive: \[
\boxed{W_j^{[0]}(t):=\varnothing,}
\] \[
\boxed{W_j^{[1]}(t):=W_{\mathrm{comm}}^{\rho_1}(j;t),}
\] \[
\boxed{W_j^{[2]}(t):=W_j^{[1]}(t)\cup W_{\mathrm{comm}}^{\rho_2}(j;t)\cup W_{\mathrm{corr}}^{\rho_2}(j;t),}
\] \[
\boxed{W_j^{[3]}(t):=W_j^{[2]}(t)\cup W_{\mathrm{nc}}^{\rho_3}(j;t),}
\] \[
\boxed{W_j^{[4]}(t):=W_j^{[3]}(t)\cup W_{\mathrm{age}}^{\rho_4}(j;t)\cup W_{\mathrm{term}}(j;t).}
\] Here \(0\le\rho_1\le\rho_2\le\rho_3\le\rho_4\le1\), and
\(W_{\mathrm{term}}(j;t)\subseteq I_{-j}\) is the optional broad
terminal or scheduled-compression window. The recursion guarantees \[
\varnothing=W_j^{[0]}(t)
\subseteq
W_j^{[1]}(t)
\subseteq
W_j^{[2]}(t)
\subseteq
W_j^{[3]}(t)
\subseteq
W_j^{[4]}(t)
\subseteq I_{-j}.
\] Rungs \(s\le2\) are frozen, flat-sector, or correlated recovery
attempts. Rungs \(s\ge3\) have opened noncommuting or broad compensators
and therefore are curvature-compensated unless the first accepted rung
occurred earlier.

Before launching an expensive exact remove-refit trial, a branch may
apply a cached Fubini--Schur screen. Let \[
\operatorname{valid}_{j}^{\mathrm{cache},[s]}(t)
:=
\operatorname{valid}_{j}^{\mathrm{cache}}(W_j^{[s]}(t);t)
\] mean that the cached tangent-overlap and Hessian/Schur blocks for
coordinate \(j\) were evaluated at the current branch state
\((\mathcal O_t^+,\theta_t^+)\) and for the same rung window
\(W_j^{[s]}(t)\). If this validity flag is zero, the projected screen
for that rung abstains; stale geometry cannot authorize a prune trial.

For nonempty \(W\subseteq I_{-j}\), define \[
\bigl(Q_j(W;t)\bigr)_{ik}
:=
\operatorname{Re}\langle \bar\tau_i,\bar\tau_k\rangle,
\qquad
\bigl(q_j(W;t)\bigr)_i
:=
\operatorname{Re}\langle \bar\tau_i,\bar\tau_j\rangle,
\] for \(i,k\in W\), and \[
F_j(t):=\operatorname{Re}\langle\bar\tau_j,\bar\tau_j\rangle+\varepsilon_F.
\] With \(\lambda_F>0\), the scale-free tangent residual is \[
\boxed{
\rho_j^{\mathrm{tan}}(W;t)
:=
\operatorname{clip}_{[0,1]}\!\left(
1-
\frac{
q_j(W;t)^\top\bigl(Q_j(W;t)+\lambda_F I\bigr)^{-1}q_j(W;t)
}{F_j(t)}
\right).
}
\] The amplitude-scaled Fubini deletion residual is obtained from \[
\delta_{W,F}^{\star}(j;t)
:=
\theta_{t,j}\bigl(Q_j(W;t)+\lambda_F I\bigr)^{-1}q_j(W;t),
\] \[
\boxed{
D_{j\mid W,F}^{2}(t)
:=
\theta_{t,j}^2F_j(t)
-2\theta_{t,j}q_j(W;t)^\top\delta_{W,F}^{\star}(j;t)
+
\bigl(\delta_{W,F}^{\star}(j;t)\bigr)^\top Q_j(W;t)\delta_{W,F}^{\star}(j;t).
}
\] In the unregularized ideal case, \[
D_{j\mid W,F}^{2}(t)
=
\theta_{t,j}^2
\left(F_j(t)-q_j(W;t)^\top Q_j(W;t)^+q_j(W;t)\right).
\] Thus \(\rho_j^{\mathrm{tan}}\) is a scale-free tangent-rank
diagnostic, while \(D_{j\mid W,F}^{2}\) is the state-space prune-screen
quantity.

For the energy-side screen, let \(H_W(t)\) be the Hessian block on
\(W\), let \(h_{Wj}(t)\) be the mixed Hessian vector between \(W\) and
coordinate \(j\), let \(h_{jW}(t):=h_{Wj}(t)^\top\), let
\(g_W(t):=(g_i(t))_{i\in W}\), and let \(h_{jj}(t)\) be the
one-coordinate curvature of \(j\). Define \[
R_j(W;t):=\frac12\left(H_W(t)+H_W(t)^\top\right)+\lambda_H I.
\] The gradient-aware quadratic deletion-refit surrogate is \[
\boxed{
\widehat L_j^{\mathrm{Schur}}(W;t)
:=
-g_j(t)\theta_{t,j}
+\frac12h_{jj}(t)\theta_{t,j}^2
-\frac12
\bigl(g_W(t)-h_{Wj}(t)\theta_{t,j}\bigr)^\top
R_j(W;t)^{-1}
\bigl(g_W(t)-h_{Wj}(t)\theta_{t,j}\bigr).
}
\] Near a locally refit state, the curvature-only version is \[
\widetilde h_j^{(-)}(W;t)
:=
h_{jj}(t)-h_{jW}(t)R_j(W;t)^{-1}h_{Wj}(t),
\] \[
\boxed{
L_j^{\mathrm{Schur}}(W;t)
:=
\frac12\theta_{t,j}^2\max\{\widetilde h_j^{(-)}(W;t),0\}.
}
\] If the branch is not near a stable local refit,
\(\widehat L_j^{\mathrm{Schur}}\) is the safer nomination proxy; if the
gradient/Hessian cache is missing or stale, the cached screen abstains
rather than launching new measurements only for nomination.

For \(s\in\{1,2,3,4\}\), define the rungwise projected screen \[
\boxed{
\Gamma_j^{\mathrm{proj},[s]}(t)
:=
\mathbf 1[\operatorname{valid}_{j}^{\mathrm{cache},[s]}(t)=1]\,
\mathbf 1[D_{j\mid W_j^{[s]},F}^{2}(t)\le(\tau_F^{\mathrm{pr}}(t))^2]\,
\mathbf 1[L_j^{\mathrm{Schur}}(W_j^{[s]};t)\le\tau_{\mathrm{Schur}}^{\mathrm{pr}}(t)].
}
\] Set \(\Gamma_j^{\mathrm{proj},[s]}(t)=0\) when
\(W_j^{[s]}(t)=\varnothing\). The first geometrically passing rung is \[
\boxed{
s_{\mathrm{geo}}(j;t)
:=
\min\{s\in\{1,2,3,4\}:\Gamma_j^{\mathrm{proj},[s]}(t)=1\},
}
\] with \(s_{\mathrm{geo}}(j;t)=\infty\) if no cached rung passes. The
projected flat and curvature-compensated sets are \[
\boxed{
\mathcal C_{\mathrm{flat}}^{\mathrm{proj}}(t)
:=
\{j\in\mathcal M_t:s_{\mathrm{geo}}(j;t)\le2\},
}
\] \[
\boxed{
\mathcal C_{\mathrm{curv}}^{\mathrm{proj}}(t)
:=
\{j\in\mathcal M_t:3\le s_{\mathrm{geo}}(j;t)\le4\}.
}
\] If only one side of the cache is available, weak geometry lanes may
also be recorded: \[
\mathcal P_F(t)
:=
\{j\in\mathcal M_t:\exists s,\ \operatorname{valid}_{j}^{\mathrm{cache},[s]}(t)=1,
\ D_{j\mid W_j^{[s]},F}^{2}(t)\le(\tau_F^{\mathrm{pr}}(t))^2\},
\] \[
\mathcal P_{\mathrm{Schur}}(t)
:=
\{j\in\mathcal M_t:\exists s,\ \operatorname{valid}_{j}^{\mathrm{cache},[s]}(t)=1,
\ L_j^{\mathrm{Schur}}(W_j^{[s]};t)\le\tau_{\mathrm{Schur}}^{\mathrm{pr}}(t)\}.
\] The uncapped trial universe is \[
\boxed{
\mathcal U_{\mathrm{trial}}(t)
:=
\mathcal N_1(t)
\cup
\mathcal C_{\mathrm{flat}}^{\mathrm{proj}}(t)
\cup
\mathcal C_{\mathrm{curv}}^{\mathrm{proj}}(t)
\cup
\mathcal P_F(t)
\cup
\mathcal P_{\mathrm{Schur}}(t).
}
\] The capped exact-trial set is \[
\boxed{
\mathcal C_{\mathrm{trial}}(t)
:=
\operatorname{Top}_{Q_{\mathrm{trial}}(t)}
\left(\preceq_{\mathrm{pr}};\ j\in\mathcal U_{\mathrm{trial}}(t)\right).
}
\] A conservative ordering is lexicographic in \[
\left(
\chi_{\mathrm{curv}}(j),
\ s_{\mathrm{geo}}(j;t),
\ L_j^{\mathrm{Schur}}(W_j^{[s_{\mathrm{geo}}]};t),
\ D_{j\mid W_j^{[s_{\mathrm{geo}}]},F}^{2}(t),
\ R_j^{\mathrm{cheap}}(t)
\right),
\] where
\(\chi_{\mathrm{curv}}(j):=\mathbf 1[s_{\mathrm{geo}}(j;t)\ge3]\), and
candidates with \(s_{\mathrm{geo}}=\infty\) are ranked by the cheap
order unless a scheduled compression pass overrides this. In all cases,
\(\mathcal C_{\mathrm{trial}}\) only selects what to test. It is not a
prune decision.

Optional cache-maintenance details such as block-prune Schur formulae,
monotonicity diagnostics, damped BFGS updates, and metric-coordinate
cache transport are auxiliary. They may improve ranking and shot
allocation, but they are not part of the deletion certificate.

\paragraph{Prune 3: exact remove--refit
ladder}\label{prune-3-exact-removerefit-ladder}

Prune 3 is the first stage that actually changes the trial scaffold. For
a candidate \(j\in\mathcal C_{\mathrm{trial}}(t)\), the controller
deletes \(j\), chooses a typed compensator window \(W_j^{[s]}\), and
refits only that surviving window while all other coordinates remain
fixed. The rung \(s\) controls how much compensating freedom is opened:
\(s=0\) is frozen deletion, \(s=1,2\) are flat-sector or correlated
recovery attempts, \(s=3\) opens noncommuting compensators, and \(s=4\)
is the broad terminal or scheduled-compression attempt.

The exact remove-refit ladder is \[
\boxed{
\delta E_j^{[s]}(t)
:=
\min_{\xi\in\Theta_{\mathcal O_{t,-j}^+}}
\left\{
E(\mathcal O_{t,-j}^+,\xi):
\xi_i=((\theta_t^+)_{-j})_i
\ \forall i\notin W_j^{[s]}(t)
\right\}
-
E(\mathcal O_t^+,\theta_t^+),
\qquad s\in\{0,1,2,3,4\}.
}
\] The first rung is frozen ablation, \[
\delta E_j^{[0]}(t)
=
E(\mathcal O_{t,-j}^+,(\theta_t^+)_{-j})
-
E(\mathcal O_t^+,\theta_t^+),
\] and, assuming comparable optimization at each rung, \[
\boxed{
\delta E_j^{[0]}(t)
\ge
\delta E_j^{[1]}(t)
\ge
\delta E_j^{[2]}(t)
\ge
\delta E_j^{[3]}(t)
\ge
\delta E_j^{[4]}(t)
\ge
\delta E_j^\star(t).
}
\] Append and prune use separate refit-breadth schedules. Append may
shrink its inherited old-coordinate fraction with maturity, while prune
broadens typed compensating windows with maturity: \[
\rho_W^{\mathrm{app}}(t)\downarrow,
\qquad
\rho_W^{\mathrm{pr}}(t)\uparrow.
\] In a shot-estimated route, typed rung \(s\) with shot budget \(N\)
uses the paired remove-refit estimator \[
\widehat\theta_{-j}^{[s,N]}
:=
\arg\min_{\xi\in\Theta_{\mathcal O_{t,-j}^+}}
\left\{
\widehat E_N(\mathcal O_{t,-j}^+,\xi):
\xi_i=((\theta_t^+)_{-j})_i
\ \forall i\notin W_j^{[s]}(t)
\right\},
\] \[
\boxed{
\widehat{\delta E}_{j}^{[s,N]}
:=
\widehat E_N(\mathcal O_{t,-j}^+,\widehat\theta_{-j}^{[s,N]})
-
\widehat E_N(\mathcal O_t^+,\theta_t^+).
}
\] Use paired Pauli-group measurement, shared random seeds, or
common-shot bookkeeping where possible, so that the variance of the
difference is estimated directly: \[
\widehat\sigma_{j}^{2}(s,N)
:=
\widehat{\operatorname{Var}}
\left[
\widehat E_N(\mathcal O_{t,-j}^+,\widehat\theta_{-j}^{[s,N]})
-
\widehat E_N(\mathcal O_t^+,\theta_t^+)
\right].
\] The upper and lower confidence bounds are \[
U_j^{[s,N]}
:=
\widehat{\delta E}_{j}^{[s,N]}
+
z_\alpha\widehat\sigma_j(s,N),
\qquad
L_j^{[s,N]}
:=
\widehat{\delta E}_{j}^{[s,N]}
-
z_\alpha\widehat\sigma_j(s,N).
\] These confidence bounds now pass to the final confirmation stage.

\paragraph{Prune 4: shot confirmation, recovery class, and
rollback}\label{prune-4-shot-confirmation-recovery-class-and-rollback}

Prune 4 converts the remove-refit trial into a committed deletion or an
exact rollback. It uses confidence bounds on the paired energy
regression, a retained admitted-gain guard, and an extra curvature guard
for any rung that opened noncommuting compensators. Cached projection,
amplitude collapse, and Schur proxies do not appear as deletion
certificates here; they have already done their job by choosing which
trial to spend budget on.

The recovery tolerance is absolute, shot-aware, and scale-aware: \[
\boxed{
\Delta_{\mathrm{pr}}^{\mathrm{tol}}(t)
:=
\max\!\left(
\Delta_{\mathrm{num}},
 c_{\mathrm{shot}}\sigma_E(t),
 c_{\mathrm{scr}}\Delta_{\mathrm{scr}}(t),
 \Delta_{\mathrm{chem}},
 c_{\mathrm{rel}}|E_t^+-E_{\mathrm{target}}|
\right),
}
\] dropping unavailable terms from the maximum. The retained-gain guard
is \[
\Gamma_{\mathrm{ret}}(j,t;s,N)
:=
\mathbf 1[\Delta_{\mathrm{adm}}(t)<\Delta_{\mathrm{ret,on}}]
\vee
\mathbf 1\!\left[
E_t^-
-
\widehat E_N(\mathcal O_{t,-j}^+,\widehat\theta_{-j}^{[s,N]})
\ge
\eta_{\mathrm{ret}}\Delta_{\mathrm{adm}}(t)
+
z_\alpha\widehat\sigma_j(s,N)
\right].
\] Define the stronger retained-gain guard
\(\Gamma_{\mathrm{ret}}^{\mathrm{strong}}(j,t;s,N)\) by replacing
\(\eta_{\mathrm{ret}}\) with
\(\eta_{\mathrm{ret}}^{\mathrm{strong}}\ge\eta_{\mathrm{ret}}\). Let
\(\operatorname{CompressionMode}_t\in\{0,1\}\) denote a scheduled
high-confidence compression pass. The curvature-compensated guard is \[
\boxed{
\begin{aligned}
\Gamma_{\mathrm{curv}}(j,t;s,N)
:=
&\ \mathbf 1[s\le2]
\vee
\mathbf 1[N\ge N_{\mathrm{hi}}]
\vee
\mathbf 1[s=4]
\vee
\mathbf 1\!\left[\frac{|W_j^{[s]}(t)|}{|I_{-j}|+\varepsilon}\ge\rho_{\mathrm{near}}\right]\\
&\vee
\mathbf 1\!\left[U_j^{[s,N]}\le \gamma_{\mathrm{curv}}\Delta_{\mathrm{pr}}^{\mathrm{tol}}(t)\right]
\vee
\mathbf 1\!\left[\Gamma_{\mathrm{ret}}^{\mathrm{strong}}(j,t;s,N)=1\right]
\vee
\mathbf 1[\operatorname{CompressionMode}_t=1],
\end{aligned}
}
\] where \(0<\gamma_{\mathrm{curv}}<1\). Any accepted rung that has
opened noncommuting compensators must be backed by at least one stronger
confirmation condition. Recovery through noncommuting compensators is
interpreted as curvature-mediated reparameterization, not proof of flat
redundancy.

The rung decision is \[
\boxed{
\mathsf D_j^{[s,N]}
:=
\begin{cases}
\mathrm{accept},
&
U_j^{[s,N]}
\le
\Delta_{\mathrm{pr}}^{\mathrm{tol}}(t)
\ \wedge\
\Gamma_{\mathrm{ret}}(j,t;s,N)=1
\ \wedge\
\Gamma_{\mathrm{curv}}(j,t;s,N)=1,
\\[1mm]
\mathrm{reject},
&
s=4,\ N=N_{\mathrm{hi}},\
L_j^{[s,N]}
>
\Delta_{\mathrm{pr}}^{\mathrm{tol}}(t),
\\[1mm]
\mathrm{escalate},
&
\text{otherwise.}
\end{cases}
}
\] If the first accepted rung has \(s\le2\), set
\(\operatorname{recovery\_class}_j=\mathrm{flat\_redundant}\). If the
first accepted rung has \(s\ge3\), set
\(\operatorname{recovery\_class}_j=\mathrm{curvature\_compensated}\). If
all scheduled rungs reject or defer without acceptance, set
\(\operatorname{recovery\_class}_j=\mathrm{failed}\). Escalation may
increase \(N\), move from \(W_j^{[s]}\) to \(W_j^{[s+1]}\), or defer to
a scheduled compression pass.

The final prune rule is \[
\boxed{
\operatorname{Prune}_j(t)=1
\Longleftrightarrow
\Gamma_{\mathrm{prune}}^{\mathrm{perm}}(t)=1
\ \wedge\
j\in\mathcal C_{\mathrm{trial}}(t)
\ \wedge\
\exists(s,N):\mathsf D_j^{[s,N]}=\mathrm{accept}.
}
\] Before remove-refit trials, store the rollback snapshot \[
\Xi_t^{\mathrm{pr}}
:=
(\mathcal O_t^+,\theta_t^+,E_t^+,\{c_i(t)\}_i,\text{current local optimizer state}).
\] If \(\Gamma_{\mathrm{prune}}^{\mathrm{perm}}(t)=0\) or
\(\mathcal C_{\mathrm{trial}}(t)=\varnothing\), no prune trial is
launched and the committed branch state remains
\((\mathcal O_t^+,\theta_t^+,E_t^+)\). Otherwise test candidates in
\(\preceq_{\mathrm{pr}}\) order up to the configured probe cap. If a
prune is accepted, commit \[
(\mathcal O_{t,-j}^+,\widehat\theta_{-j}^{[s,N]},\widehat E_N(\mathcal O_{t,-j}^+,\widehat\theta_{-j}^{[s,N]}))
\] as the new branch state, write
\((j,s,N,\operatorname{recovery\_class}_j)\) to the branch prune-event
log, and then delete the live metadata carried by coordinate \(j\). If
it is rejected, restore the branch state exactly to
\(\Xi_t^{\mathrm{pr}}\), set \(c_j(t)\leftarrow c_{\mathrm{cool}}\), and
continue according to the configured candidate budget. Hence prune
failure never invalidates an already committed admission, and rollback
is exact.

\subsubsection{11.5.2 Beam-adapt over scaffold alternatives: branch
state, pruning, and effective
selector}\label{beam-adapt-over-scaffold-alternatives-branch-state-pruning-and-effective-selector}

In this section beam-adapt is \textbf{purely structural}: each branch
carries an alternative ADAPT scaffold together with its locally refit
amplitudes. There is no time checkpoint, probe rollout, or
projected-dynamics integral in this section; those belong to the
separate time-dynamics beam surface developed later.

The cleanest way to read a live scaffold branch is \[
\mathfrak b=
\begin{aligned}[t]
\bigl(&\text{identity/history};\ \text{current scaffold and amplitudes};\\
&\text{current energy/depth};\ \text{cumulative selector totals/controller telemetry/status}\bigr).
\end{aligned}
\] More explicitly, write \[
\mathfrak b=\bigl(\operatorname{id}_b,\pi_b,\mathcal H_b,\mathcal O_b,\theta_b,E_b,d_b,S_b^{\mathrm{cum}},K_b^{\mathrm{cum}},\mathcal T_b^{\mathrm{ctrl}},\mathrm{status}_b,\mathrm{term}_b\bigr),
\] with \(\mathcal H_b=(r_{b,1},\dots,r_{b,d_b})\) and
\(r_{b,j}=(m_{b,j}^{\mathrm{adm}},p_{b,j}^{\mathrm{adm}})\). Here
\(\operatorname{id}_b\) and \(\pi_b\) are the branch and parent
identifiers, \(\mathcal H_b\) is the ordered history of admitted
records, \(\mathcal O_b\) is the scaffold currently carried by branch
\(b\), \(\theta_b\) is the refit parameter vector on that scaffold,
\(E_b=E(\theta_b;\mathcal O_b)\) is the current variational energy,
\(d_b\) is the branch ADAPT depth,
\(S_b^{\mathrm{cum}}=\sum_{j=1}^{d_b}s_{b,j}\) is the cumulative
admitted selector value,
\(K_b^{\mathrm{cum}}=\sum_{j=1}^{d_b}K_{b,j}^{\mathrm{sel}}\) is the
cumulative admitted selector burden, \(\mathcal T_b^{\mathrm{ctrl}}\)
stores branch-local controller telemetry, and
\((\mathrm{status}_b,\mathrm{term}_b)\) record live/terminal status and
termination label. The increment \(K_{b,j}^{\mathrm{sel}}\) is the
selector-cost contribution attached to the \(j\)-th admission on branch
\(b\); for a newly materialized child this contribution is the
branch-local Phase-3 cost \(K_3(r;\mathfrak b)\). Thus
\(K_b^{\mathrm{cum}}\) is cumulative selector-side cost bookkeeping
rather than a raw hardware count.

If beam search is entered from an incumbent scaffold
\((\mathcal O^{(0)},\theta^{(0)})\), the root branch is
\(\mathfrak b_{\mathrm{root}}=\bigl(\operatorname{id}_0,\varnothing,\mathcal H^{(0)},\mathcal O^{(0)},\theta^{(0)},E(\theta^{(0)};\mathcal O^{(0)}),d_0,S_{\mathrm{root}}^{\mathrm{cum}},K_{\mathrm{root}}^{\mathrm{cum}},\mathcal T_{\mathrm{root}}^{\mathrm{ctrl}},\mathrm{frontier},\varnothing\bigr)\),
with \(d_0=|\mathcal H^{(0)}|\),
\(S_{\mathrm{root}}^{\mathrm{cum}}=\sum_{j=1}^{d_0}s_j^{(0)}\), and
\(K_{\mathrm{root}}^{\mathrm{cum}}=\sum_{j=1}^{d_0}K_j^{\mathrm{sel},(0)}\),
where \(K_j^{\mathrm{sel},(0)}\) denotes the selector-cost contribution
attached to the \(j\)-th historical admission on the incumbent scaffold.
If the beam is started from a fresh scaffold, then
\(\mathcal H^{(0)}=\varnothing\), \(d_0=0\), and both cumulative sums
vanish.

For clock and budget, set \(t_b:=d_b\) and
\(D_{\mathrm{left}}(\mathfrak b):=d_{\max}-d_b\). Now apply the branch
lift rule: for any live-state quantity \(X\) from Phase 3 whose
definition is written in terms of
\((\mathcal H,\mathcal O,\theta,t,\mathcal T^{\mathrm{ctrl}})\), define
\(X(\mathfrak b):=X\big|_{(\mathcal H,\mathcal O,\theta,t,\mathcal T^{\mathrm{ctrl}})\mapsto(\mathcal H_b,\mathcal O_b,\theta_b,t_b,\mathcal T_b^{\mathrm{ctrl}})}\);
for any record-local quantity \(X(r;t)\), define
\(X(r;\mathfrak b):=X(r;t)\big|_{(\mathcal H,\mathcal O,\theta,t,\mathcal T^{\mathrm{ctrl}})\mapsto(\mathcal H_b,\mathcal O_b,\theta_b,t_b,\mathcal T_b^{\mathrm{ctrl}})}\).
In particular, \(R(\mathfrak b)\), \(u_{\mathrm{sat}}(\mathfrak b)\),
\(\widehat N_{\mathrm{rem}}(\mathfrak b)\), \(K_3(r;\mathfrak b)\), and
\(g_{\mathrm{lcb}}(r;\mathfrak b)\) are branch-lifted evaluations of
their previously defined live-state counterparts.

The Phase-3 chain is branch-local by construction: all of its
dependencies factor through the state tuple
\((\mathcal O_b,\theta_b,t_b,\mathcal T_b^{\mathrm{ctrl}})\). Hence each
live branch induces its own raw candidate-position surface
\(\mathcal R_b^{\mathrm{raw}}\), its own branch-local shortlist
\(\mathcal S_3(\mathfrak b)\), and its own retained admission set
\(\mathcal A_b\).

For stage and enabled family, let
\(\eta_b\in\{\mathrm{core/seed},\mathrm{residual}\}\) and \[
\mathcal G_{\eta_b}^{(3)}(\mathfrak b)=
\begin{cases}
\mathcal G_{\mathrm{core/seed}}^{(3)}(\mathfrak b),&\eta_b=\mathrm{core/seed},\\
\mathcal G_{\mathrm{residual}}^{(3)}(\mathfrak b),&\eta_b=\mathrm{residual}.
\end{cases}
\] Let
\(\mathcal P_{\mathrm{residual}}(\mathfrak b)\subseteq\mathcal G^{(3)}\)
denote the residual-only generator family on branch \(b\). Then the
controller-stage gate and branch-local available Phase-3 generator
family are \[
\Gamma_{\mathrm{stage},b}(m)=
\begin{cases}
1,&\eta_b=\mathrm{residual},\\
1,&\eta_b=\mathrm{core/seed}\ \text{and}\ m\notin\mathcal P_{\mathrm{residual}}(\mathfrak b),\\
0,&\eta_b=\mathrm{core/seed}\ \text{and}\ m\in\mathcal P_{\mathrm{residual}}(\mathfrak b),
\end{cases}
\qquad
\mathcal G_{\mathrm{avail}}^{(3)}(\mathfrak b)=\{m\in\mathcal G^{(3)}:\Gamma_{\mathrm{stage},b}(m)=1\}.
\] For scaffold length, append slot, and active probe window, let
\(n_b=|\mathcal O_b|\), \(p_{\mathrm{app}}(\mathfrak b)=n_b\), and
\(W_{\mathrm{act}}(\mathfrak b)\subseteq\{0,\dots,n_b-1\}\). The ordered
probe list is
\(\mathcal L_{\mathrm{probe}}(\mathfrak b):=\bigl(p_{\mathrm{app}}(\mathfrak b),0,W_{\mathrm{act}}(\mathfrak b)\bigr)\),
and the distinct retained probe positions are
\(\mathcal P_{\mathrm{probe}}(\mathfrak b)=\operatorname{Head}_{M_{\mathrm{probe}}}\!\Bigl(\operatorname{Dedup}_{\mathrm{ord}}\bigl(\mathcal L_{\mathrm{probe}}(\mathfrak b)\bigr)\Bigr)\subseteq\{0,\dots,n_b\}\).
For symmetry and admissible positions, let \[
\Gamma_{\mathrm{sym},b}(m,p)=
\begin{cases}
0,&\text{the symmetry audit for }(m,p;\mathcal O_b)\text{ returns a hard violation},\\
1,&\text{otherwise},
\end{cases}
\] so the branch-local raw candidate-position surface is \[
\mathcal R_b^{\mathrm{raw}}
:=
\{\,r=(m,p):m\in\mathcal G_{\mathrm{avail}}^{(3)}(\mathfrak b),\ p\in\mathcal P_{\mathrm{probe}}(\mathfrak b),\ \Gamma_{\mathrm{sym},b}(m,p)=1\,\}.
\] The resulting local shortlist and retained admission set are \[
\mathcal S_3(\mathfrak b)
:=
\operatorname{Top}_{N_3(t_b)}\!\left(S_3(\cdot;\mathfrak b);\mathcal R_b^{\mathrm{raw}}\right),
\qquad
\mathcal A_b
:=
\operatorname{Top}_{K_b}\!\left(S_3(\cdot;\mathfrak b);\{r\in\mathcal S_3(\mathfrak b):S_3(r;\mathfrak b)>0\}\right).
\] When a branch-local rerank shell is judged tie-degenerate, the
retained admission set may be replaced by a transient tie-shell
selector.\footnote{For example, if $r_{b,\max}\in\arg\max_{r\in\mathcal S_3(\mathfrak b)}S_3(r;\mathfrak b)$, one may define $\mathcal T_{\mathrm{tie}}(\mathfrak b):=\{\,r\in\mathcal S_3(\mathfrak b): S_3(r;\mathfrak b)\ge \max(\eta_{\mathrm{tie}}\,S_3(r_{b,\max};\mathfrak b),\,S_3(r_{b,\max};\mathfrak b)-\delta_{\mathrm{tie}})\,\}$ and replace $\mathcal A_b$ by a retained subset $\mathcal A_b^{\mathrm{tie}}\subseteq \mathcal T_{\mathrm{tie}}(\mathfrak b)$ with $1\le |\mathcal A_b^{\mathrm{tie}}|\le B_{\mathrm{tie}}(t_b)$, chosen by descending $S_3(\cdot;\mathfrak b)$.}

\textbf{(ii) Materialize either stop or one new admission.} The proposal
family is \(\mathcal Q(b)=\{\mathrm{stop}\}\cup\mathcal A_b\). The
symbol \(\mathrm{stop}\) terminates the scaffold as-is, while each
\(r=(m,p)\in\mathcal A_b\) inserts one new generator into the branch and
refits. If \(\mathcal O_b=(m_{b,1},\dots,m_{b,n_b})\) and
\(0\le p\le n_b\), then the scaffold insertion map is
\(\operatorname{Insert}(\mathcal O_b,m,p)=(m_{b,1},\dots,m_{b,p},m,m_{b,p+1},\dots,m_{b,n_b})\),
the insertion index transport is
\(I_p(j)=\begin{cases} j,&1\le j\le p,\\ j+1,&p<j\le n_b, \end{cases}\),
and the zero-lifted parameter vector is
\(\theta_b^{\uparrow(r)}=(\theta_{b,1},\dots,\theta_{b,p},0,\theta_{b,p+1},\dots,\theta_{b,n_b})\).

For contiguous inserted-window cap \(\omega_{\mathrm{new}}\), let
\(\omega_{\mathrm{eff}}^{(p)}:=\max\{1,\min(\omega_{\mathrm{new}},n_b+1)\}\)
and
\(\ell_p:=\max\!\left\{1,\min\!\left(p+1-\left\lfloor\frac{\omega_{\mathrm{eff}}^{(p)}-1}{2}\right\rfloor,\;n_b+2-\omega_{\mathrm{eff}}^{(p)}\right)\right\}\).
Then
\(W_{\mathrm{new}}^{(p)}=\{\ell_p,\;\ell_p+1,\;\ldots,\;\ell_p+\omega_{\mathrm{eff}}^{(p)}-1\}\),
so \(|W_{\mathrm{new}}^{(p)}|=\omega_{\mathrm{eff}}^{(p)}\) and
\(p+1\in W_{\mathrm{new}}^{(p)}\). Let the retained old-coordinate
fraction shrink with maturity according to
\(\rho_W(t):=\rho_W^{\min}+(\rho_W^{\max}-\rho_W^{\min})e_t\), with
\(0\le \rho_W^{\min}\le \rho_W^{\max}\le 1\), and define the
old-coordinate complement by
\(C_{\mathrm{old}}^{(p)}:=\{1,\dots,n_b+1\}\setminus W_{\mathrm{new}}^{(p)}\),
so \(|C_{\mathrm{old}}^{(p)}|=n_b+1-\omega_{\mathrm{eff}}^{(p)}\). Then
the retained old-coordinate count is
\(n_{\mathrm{old}}(t,p):=\left\lceil \rho_W(t)\,|C_{\mathrm{old}}^{(p)}|\right\rceil=\left\lceil \rho_W(t)\,\bigl(n_b+1-\omega_{\mathrm{eff}}^{(p)}\bigr)\right\rceil\).
For each current scaffold coordinate \(j\in\{1,\dots,n_b\}\), let
\(a_{b,j}^{\mathrm{adm}}\) denote the admission step at which that
coordinate entered the current scaffold on branch \(b\), with smaller
values meaning older coordinates, and let the insertion-lifted
admission-order vector be
\(a_b^{\uparrow(r)}:=(a_{b,1}^{\mathrm{adm}},\dots,a_{b,p}^{\mathrm{adm}},+\infty,a_{b,p+1}^{\mathrm{adm}},\dots,a_{b,n_b}^{\mathrm{adm}})\),
so the newly inserted slot is never selected as an old coordinate. The
retained old-coordinate window is then
\(W_{\mathrm{age}}^{(p)}:=\{\,j\in C_{\mathrm{old}}^{(p)}:\#\{\,i\in C_{\mathrm{old}}^{(p)}:(a_b^{\uparrow(r)})_i<(a_b^{\uparrow(r)})_j\,\}<n_{\mathrm{old}}(t,p)\,\}\),
namely the set of the \(n_{\mathrm{old}}(t,p)\) oldest old coordinates
outside \(W_{\mathrm{new}}^{(p)}\), and hence the inherited refit window
on the inserted scaffold is
\(W_r=W(r;t):=W_{\mathrm{new}}^{(p)}\cup W_{\mathrm{age}}^{(p)}\).

Then, for \(r=(m,p)\) and
\(\mathcal O_b\oplus_p m:=\operatorname{Insert}(\mathcal O_b,m,p)\), the
branch-local refit map is
\(\operatorname{Refit}_{W_r}(\theta_b;r):=\arg\min_{\widetilde\theta\in\Theta_{\mathcal O_b\oplus_p m}}\left\{E(\mathcal O_b\oplus_p m,\widetilde\theta):\ \widetilde\theta_j=(\theta_b^{\uparrow(r)})_j\ \forall j\notin W_r\right\}\).

The transfer map is
\(\mathcal T(\mathfrak b,\mathrm{stop})=\mathfrak b'\) with
\(\mathcal H_{b'}=\mathcal H_b\), \(\mathcal O_{b'}=\mathcal O_b\),
\(\theta_{b'}=\theta_b\), \(E_{b'}=E_b\), \(d_{b'}=d_b\),
\(S_{b'}^{\mathrm{cum}}=S_b^{\mathrm{cum}}\),
\(K_{b'}^{\mathrm{cum}}=K_b^{\mathrm{cum}}\),
\(\mathrm{status}_{b'}=\mathrm{terminal}\), and
\(\mathrm{term}_{b'}=\begin{cases} \mathrm{empty},&\mathcal A_b=\varnothing,\\ \mathrm{stop},&\mathcal A_b\neq\varnothing. \end{cases}\)
For \(r=(m,p)\), \(\mathcal T(\mathfrak b,r)=\mathfrak b'\) with
\(\mathcal O_{b'}=\operatorname{Insert}(\mathcal O_b,m,p)\),
\(\theta_{b'}=\operatorname{Refit}_{W_r}(\theta_b;r)\),
\(\mathcal H_{b'}=(\mathcal H_b,r)\),
\(E_{b'}=E(\theta_{b'};\mathcal O_{b'})\), \(d_{b'}=d_b+1\),
\(S_{b'}^{\mathrm{cum}}=S_b^{\mathrm{cum}}+S_3(r;\mathfrak b)\),
\(K_{b'}^{\mathrm{cum}}=K_b^{\mathrm{cum}}+K_3(r;\mathfrak b)\),
\(\mathrm{status}_{b'}=\begin{cases} \mathrm{terminal},&d_{b'}\ge d_{\max}\ \text{or}\ \mathcal A_{b'}=\varnothing,\\ \mathrm{frontier},&d_{b'}<d_{\max}\ \text{and}\ \mathcal A_{b'}\neq\varnothing, \end{cases}\),
and
\(\mathrm{term}_{b'}=\begin{cases} \mathrm{depth\_cap},&d_{b'}\ge d_{\max},\\ \mathrm{empty},&d_{b'}<d_{\max}\ \text{and}\ \mathcal A_{b'}=\varnothing,\\ \varnothing,&d_{b'}<d_{\max}\ \text{and}\ \mathcal A_{b'}\neq\varnothing. \end{cases}\)
where \(\mathcal A_{b'}\) is recomputed from step (i) on the child
branch itself. So a non-stop child is exactly one more ADAPT admission
applied to the parent scaffold, followed by the inherited-window refit
already defined above, and its cumulative \(K\)-coordinate is updated by
the admitted branch-local Phase-3 selector burden
\(K_3(r;\mathfrak b)\). If \(\mathcal A_b=\varnothing\), then
\(\mathcal Q(b)=\{\mathrm{stop}\}\) and the branch contributes only a
terminal child.

Branch-local rollback law: before materializing \(r=(m,p)\), store the
branch snapshot
\(\Xi_b=(\mathcal O_b,\theta_b,E_b,\mathcal H_b,S_b^{\mathrm{cum}},K_b^{\mathrm{cum}},\text{optimizer memory})\).
If the post-refit energy violates \(E_{b'}>E_b+\epsilon_{\mathrm{rb}}\),
the parameter rollback mode restores only the zero-lifted pre-refit
parameters on the enlarged scaffold, while the structural rollback mode
rejects the admission and restores \(\Xi_b\) exactly. Thus insertion,
beam branching, batching, and pruning remain problem-generic, while an
outer controller-tuning layer may tune the insertion trigger, rollback
mode, and \(\epsilon_{\mathrm{rb}}\) as policy parameters rather than
Hamiltonian-specific exceptions.

\textbf{(iii) Classify, deduplicate, and prune the scaffold children.}
The frontier and terminal descendants of branch \(b\) are
\(\mathfrak C_b^{\mathrm{front}}=\{\mathfrak b' : \mathfrak b'=\mathcal T(\mathfrak b,r),\ r\in\mathcal A_b,\ \mathrm{status}_{b'}=\mathrm{frontier}\}\)
and
\(\mathfrak C_b^{\mathrm{term}}=\{\mathfrak b' : \mathfrak b'=\mathcal T(\mathfrak b,q),\ q\in\mathcal Q(b),\ \mathrm{status}_{b'}=\mathrm{terminal}\}\).
Two branches are treated as duplicates when they represent the same
scaffold and essentially the same refit point, for example under the
fingerprint
\(\operatorname{rnd}_{10}(x):=10^{-10}\operatorname{round}(10^{10}x)\),
\(\operatorname{round}_{10}(\theta_b):=\bigl(\operatorname{rnd}_{10}(\theta_{b,1}),\dots,\operatorname{rnd}_{10}(\theta_{b,n_b})\bigr)\),
and
\(\mathrm{fp}(\mathfrak b)=\bigl(d_b,\operatorname{labels}(\mathcal O_b),\operatorname{round}_{10}(\theta_b)\bigr)\).
Among near-equivalent branches, keep the one with best lexicographic
prune key
\(\pi_{\mathrm{prune}}(\mathfrak b)=\bigl(E_b,-S_b^{\mathrm{cum}},K_b^{\mathrm{cum}},|\mathcal O_b|,\operatorname{labels}(\mathcal O_b),\operatorname{round}_{10}(\theta_b),\operatorname{id}_b\bigr)\).

The beam operators are therefore
\(\operatorname{Dedup}(\mathcal X)=\{\text{best branch per fingerprint under }\pi_{\mathrm{prune}}\}\)
and
\(\operatorname{Prune}(\mathcal X;B)=\operatorname{Top}^{\min}_B\!\left(\pi_{\mathrm{prune}};\operatorname{Dedup}(\mathcal X)\right)\).
Hence the live and terminal pools update by
\(\mathfrak F_{c+1}=\operatorname{Prune}\!\left(\bigcup_{b\in\mathfrak F_c}\mathfrak C_b^{\mathrm{front}};B_{\mathrm{live}}\right)\)
and
\(\mathfrak T_{c+1}=\operatorname{Prune}\!\left(\mathfrak T_c\cup\bigcup_{b\in\mathfrak F_c}\mathfrak C_b^{\mathrm{term}};B_{\mathrm{term}}\right)\).

\textbf{(iv) Choose the final scaffold branch.} After \(C\)
scaffold-beam rounds, the surviving candidate set is
\(\mathfrak W_{\mathrm{final}}=\mathfrak F_C\cup\mathfrak T_C\), and the
effective beam selector is
\(\mathfrak b_*=\arg\min_{\mathfrak b\in\mathfrak W_{\mathrm{final}}}\pi_{\mathrm{prune}}(\mathfrak b)\).

A compact scaffold-beam algorithm block is therefore \[
\begin{aligned}
\text{A: }&
\mathfrak F_0=\{\mathfrak b_{\mathrm{root}}\},
\qquad
\mathfrak T_0=\varnothing,\\
\text{B: }&
\mathcal R_b^{\mathrm{raw}}
\to
\mathcal S_3(\mathfrak b)
\to
\mathcal A_b
\qquad(\mathfrak b\in\mathfrak F_c),\\
\text{C: }&
\mathcal Q(b)=\{\mathrm{stop}\}\cup\mathcal A_b,
\qquad
\mathfrak C_b=\{\mathcal T(\mathfrak b,q):q\in\mathcal Q(b)\},\\
\text{D: }&
\mathfrak C_b^{\mathrm{front}}
=
\{\mathfrak b'\in\mathfrak C_b:\mathrm{status}_{b'}=\mathrm{frontier}\},
\qquad
\mathfrak C_b^{\mathrm{term}}
=
\{\mathfrak b'\in\mathfrak C_b:\mathrm{status}_{b'}=\mathrm{terminal}\},\\
\text{E: }&
\mathfrak F_{c+1}
=
\operatorname{Prune}\!\left(\bigcup_{b\in\mathfrak F_c}\mathfrak C_b^{\mathrm{front}};B_{\mathrm{live}}\right),
\qquad
\mathfrak T_{c+1}
=
\operatorname{Prune}\!\left(\mathfrak T_c\cup\bigcup_{b\in\mathfrak F_c}\mathfrak C_b^{\mathrm{term}};B_{\mathrm{term}}\right),\\
\text{F: }&
\mathfrak W_{\mathrm{final}}=\mathfrak F_C\cup\mathfrak T_C,
\qquad
\mathfrak b_*=\arg\min_{\mathfrak b\in\mathfrak W_{\mathrm{final}}}\pi_{\mathrm{prune}}(\mathfrak b).
\end{aligned}
\]

\subsection{11.6 Active HH pool surface and selected scaffold}

Identify the nested active HH pool with the retained generator images by
\(\Omega_{\mathrm{HH}}^{(k)}:=\mathcal G^{(k)}\) for \(k\in\{1,2,3\}\).
The active HH adaptive surface is therefore carried by
\(\Omega_{\mathrm{HH}}^{(3)}\subseteq\Omega_{\mathrm{HH}}^{(2)}\subseteq\Omega_{\mathrm{HH}}^{(1)}\),
with
\(\mathcal G_{\mathrm{adapt}}^{(k)}=\{G_m\}_{m\in\Omega_{\mathrm{HH}}^{(k)}}\).

If beam continuation is used, set
\(\mathcal O_*:=\mathcal O_{\mathfrak b_*}\) and
\(\theta_*^{\mathrm{adapt}}:=\theta_{\mathfrak b_*}\); otherwise
\((\mathcal O_*,\theta_*^{\mathrm{adapt}})\) denotes the final admitted
scaffold and its refit amplitudes on the main branch. Write
\(n_*:=|\mathcal O_*|\) for the final scaffold length, and write the
selected ordered scaffold as \(\mathcal O_*=(G_1,\dots,G_{n_*})\). For
each \(j\in\{1,\dots,n_*\}\), write
\(G_j=\sum_{a=1}^{q_j} c_{ja}P_{ja}\), where \(q_j\) is the number of
active nonidentity Pauli terms in the \(j\)-th selected generator,
\(P_{ja}\) is the \(a\)-th such Pauli term, and \(c_{ja}\in\mathbb R\)
its coefficient. Decompose the selected amplitudes as
\(\theta_*^{\mathrm{adapt}}=(\theta_{*,1}^{\mathrm{adapt}},\dots,\theta_{*,n_*}^{\mathrm{adapt}})\)
with \(\theta_{*,j}^{\mathrm{adapt}}\in\mathbb R^{q_j}\). Thus each
selected generator carries a parameter block rather than a singleton
scalar.\footnote{When block-valued selected generators are used, the clean index convention is $j$ for the scaffold block, $a$ for the coordinate inside block $j$, and $\mu=(j,a)$ for the flattened global coordinate. Metric, force, regularized normal, and tangent matrices are indexed by global coordinates $\mu,\nu$, while append and prune decisions may still operate on block labels $j$.}
The resulting state is
\(|\psi_*\rangle=U(\theta_*^{\mathrm{adapt}};\mathcal O_*)|\phi_0\rangle\),
so the active HH scaffold manifold selected by ADAPT may be written as
\(\mathcal M_{\mathrm{scaf}}(\mathcal O_*)=\{\,U(\vartheta;\mathcal O_*)|\phi_0\rangle:\vartheta\in\prod_{j=1}^{n_*}\mathbb R^{q_j}\,\}\).

\subsection{Appendix 11-S. Supplemental algebraic
telemetry}\label{appendix-11-s.-supplemental-algebraic-telemetry}

The canonical selector uses exact local relation counts to assign
algebraic lanes. The continuous quantities below are supplemental
telemetry only. They may be logged, used in lane-budget diagnostics,
used in prune-risk bookkeeping, or activated in ablation studies, but
they do not multiply \(S_{1,\mathrm{rank}}\), \(S_2\), or \(S_3\).

Let \(\varepsilon_{\mathrm{alg}}>0\) be the small denominator stabilizer
used in algebraic ratios. For a candidate-position record \(r=(m,p)\),
define \[
\begin{aligned}
\operatorname{local\_comm\_degree}(r)
&:=
\frac{
\sum_{n\in\mathcal N_{\mathrm{loc}}(m)}
\Omega_{mn}\mathsf C_{mn}
}{
\sum_{n\in\mathcal N_{\mathrm{loc}}(m)}
\Omega_{mn}+\varepsilon_{\mathrm{alg}}
},\\[1mm]
\operatorname{local\_noncomm\_degree}(r)
&:=
\frac{
\sum_{n\in\mathcal N_{\mathrm{loc}}(m)}
\Omega_{mn}(1-\mathsf C_{mn})
}{
\sum_{n\in\mathcal N_{\mathrm{loc}}(m)}
\Omega_{mn}+\varepsilon_{\mathrm{alg}}
},\\[1mm]
\operatorname{noncomm\_bridge\_score}(r)
&:=
\operatorname{clip}_{[0,1]}\!\left(
\frac{
\#\{(a,b):a\neq b,\ a,b\in\mathcal N_{\mathrm{loc}}(m),\
\mathsf C_{ab}=1,\ \mathsf C_{ma}=0,\ \mathsf C_{mb}=0\}
}{
\#\{(a,b):a\neq b,\ a,b\in\mathcal N_{\mathrm{loc}}(m),\ \mathsf C_{ab}=1\}
+\varepsilon_{\mathrm{alg}}
}
\right),\\[1mm]
\operatorname{flat\_cluster\_saturation}(r;t)
&:=
\sum_{j\in\mathcal O_t}
\mathbf 1[c(m_j)=c(m)]\,\mathbf 1[\Omega_{m_jm}=1].
\end{aligned}
\] The corresponding supplemental record fields are
\(\operatorname{comm\_cluster\_id}\),
\(\operatorname{local\_comm\_degree}\),
\(\operatorname{local\_noncomm\_degree}\),
\(\operatorname{noncomm\_bridge\_score}\), and
\(\operatorname{flat\_cluster\_saturation}\). These fields describe
algebraic context after the exact lane assignment has already been made.

For continuity with the retired continuous-prior route, let
\(\mathcal J_{\mathrm{rec}}(t)\) be the newest protected or active
scaffold coordinates considered by the selector, with recency weights
\(\omega_{j,t}\ge0\), and let
\(q_{0.9}^{\mathrm{cluster}}(t):=\operatorname{Quantile}_{0.9}\bigl(\{\operatorname{flat\_cluster\_saturation}(u;t):u\in\mathcal R_k(t)\}\bigr)\)
on the current phase-input record family. The retired telemetry
quantities are \[
\begin{aligned}
\rho_{\mathrm{nc}}^{\mathrm{rec}}(r;t)
&:=
\frac{
\sum_{j\in\mathcal J_{\mathrm{rec}}(t)}
\omega_{j,t}\,\Omega_{m_jm}\,(1-\mathsf C_{m_jm})
}{
\sum_{j\in\mathcal J_{\mathrm{rec}}(t)}
\omega_{j,t}\,\Omega_{m_jm}+\varepsilon_{\mathrm{alg}}
},\\[1mm]
\pi_{\mathrm{flat}}(r;t)
&:=
\operatorname{clip}_{[0,1]}\!\left(
\frac{\operatorname{flat\_cluster\_saturation}(r;t)}{
q_{0.9}^{\mathrm{cluster}}(t)+\varepsilon_{\mathrm{alg}}}
\right).
\end{aligned}
\] These ratios summarize recent noncommutation and flat-cluster
saturation as continuous diagnostics. They are not authoritative
selection terms.

\section{12. SPSA and Optimizer
Semantics}\label{spsa-and-optimizer-semantics}

\subsection{12.1 SPSA schedules}\label{spsa-schedules}

The standard SPSA schedules are \[
c_k=\frac{c}{(k+1)^\gamma},
\qquad
a_k=\frac{a}{(A+k+1)^\alpha}.
\]

\subsection{12.2 Two-point stochastic
gradient}\label{two-point-stochastic-gradient}

At iteration \(k\), with Rademacher direction \(\Delta_k\), evaluate \[
y_+=f(x_k+c_k\Delta_k),
\qquad
y_-=f(x_k-c_k\Delta_k),
\] and form the gradient estimate \[
\hat g_k = \frac{y_+-y_-}{2c_k}\,\Delta_k.
\]

\subsection{12.3 Update and projection}\label{update-and-projection}

The update is \[
x_{k+1}=x_k-a_k\hat g_k.
\] If projection or clipping is imposed, it is applied to both the
perturbed points and the updated iterate.

\subsection{12.4 Repeat aggregation}\label{repeat-aggregation}

If the same point is sampled multiple times, the repeated evaluations
may be aggregated by \[
\operatorname{mean}
\qquad\text{or}\qquad
\operatorname{median}.
\]

\subsection{12.5 Return policy}\label{return-policy}

Two standard return policies are:

\begin{itemize}
\tightlist
\item
  return an averaged late-trajectory iterate,
\item
  return the best sampled point observed during the run.
\end{itemize}

\subsection{12.6 Surface note}\label{surface-note}

Stochastic inner optimizers are especially natural for large, noisy, or
nondifferentiable effective surfaces, while deterministic optimizers are
often effective on smaller and smoother fixed-structure manifolds.

\subsection{12.7 Opt-in QN-SPSA
extension}\label{opt-in-qn-spsa-extension}

Quantum natural SPSA keeps the same stochastic gradient sample
\(\hat g_k\) but also samples a local metric from a fidelity surface
\(F(x,y)\). With independent Rademacher directions \(\Delta_k,\Xi_k\),
define \[
\begin{aligned}
F_+&=F(x_k,x_k+c_k\Delta_k),\qquad
F_-=F(x_k,x_k-c_k\Delta_k),\\
F_{++}&=F(x_k,x_k+c_k(\Delta_k+\Xi_k)),\qquad
F_{-+}=F(x_k,x_k+c_k(-\Delta_k+\Xi_k)).
\end{aligned}
\] The scalar curvature probe is \[
d_k=\frac{(F_{++}-F_+)-(F_{-+}-F_-)}{2c_k^2},
\] and the symmetric metric sample is \[
\widehat G_k=-\frac12 d_k\,\frac{\Delta_k\Xi_k^{\top}+\Xi_k\Delta_k^{\top}}{2}.
\] A smoothed metric \(\bar G_k\) is symmetrized, regularized, clipped
to a positive semidefinite floor, and used to solve \[
\bar G_k v_k=\hat g_k,
\qquad
x_{k+1}=x_k-a_k v_k.
\] If the metric solve is singular, ill-conditioned, or non-finite, the
safe fallback is the vanilla SPSA direction \(v_k=\hat g_k\). In this
repository QN-SPSA is explicit opt-in only; the default stochastic inner
optimizer remains SPSA.

\section{17A. Append--Prune McLachlan Dynamics for Driven Mixed
Fermion--Boson
Evolution}\label{a.-appendprune-mclachlan-dynamics-for-driven-mixed-fermionboson-evolution}

This chapter is the active mathematical source for the dynamics
manuscript. The method is Append--Prune McLachlan dynamics
(AP-McLachlan): a variational real-time evolution in which the ansatz
support is fixed on open time intervals and may be patched at discrete
time points by no edit, zero-initialized append, verified prune, or a
simultaneous append--prune exchange. Exact/reference trajectories are
diagnostic outputs only. Controller decisions use measurement-compatible
tangent geometry and observable estimates of the prepared variational
state.

The older hybrid-checkpoint, secant-forecast, turbulent-regime,
adaptive-integrator, and large-amplitude theta heuristics are not the
active route identity and should not guide Paper-II implementation. They
are archival derivation material only. The main law below intentionally
keeps time integration as a declared numerical setting; the paper-facing
mechanism is support-patch control through grouped McLachlan gain and
loss.

\subsection{17A.1 Time grid, ansatz support, and projected McLachlan
geometry}\label{a.1-time-grid-ansatz-support-and-projected-mclachlan-geometry}

Let the control time points be \[
0=  au_0<   au_1<\cdots<    au_T=T,
\qquad
 t_k=t_0+\tau_k,
\qquad
\Delta\tau_k=\tau_{k+1}-\tau_k .
\] At time point \(k\), the ordered variational support is
\(\mathcal S_k\). Its persistent block labels are \(b\in\mathcal B_k\).
A block may be a Pauli-polynomial generator \[
A_b=\sum_{\ell=1}^{m_b}c_{b\ell}P_{b\ell}.
\] The support is block-valued, while the McLachlan solve is
coordinate-valued. Chapter 17A therefore treats the coordinate
parameterization as part of the benchmark protocol. In the per
Pauli/polynomial-term convention, \[
I_{b,k}^{\mathrm{term}}:=\{(b,1),\ldots,(b,m_b)\},
\qquad
U_{k,\mathrm{term}}(\theta)
:=\prod_{b\in\mathcal B_k}^{\leftarrow}
\prod_{\ell=m_b}^{1}
\exp\!\left(-i\theta_{b\ell}c_{b\ell}P_{b\ell}\right).
\] In the per logical/macro-generator convention, \[
I_{b,k}^{\mathrm{log}}:=\{b\},
\qquad
U_{k,\mathrm{log}}(\alpha)
:=\prod_{b\in\mathcal B_k}^{\leftarrow}
\prod_{\ell=m_b}^{1}
\exp\!\left(-i\alpha_{b}c_{b\ell}P_{b\ell}\right),
\] with \(U_{k,\mathrm{log}}(\alpha)\) written as
\(\prod_b^{\leftarrow}\exp(-i\alpha_bA_b)\) only when the block is
represented as a grouped exact generator. The per-term tangent family is
indexed by \((b,\ell)\), whereas the logical tangent for block \(b\) is
the diagonal derivative through all component factors, \[
\bar T_{b,k}^{\mathrm{log}}
:=Q_{\psi_k}\partial_{\alpha_b}|\psi(\alpha;\mathcal S_k)\rangle
=Q_{\psi_k}\sum_{\ell=1}^{m_b}
\left.\partial_{\theta_{b\ell}}|\psi(\theta;\mathcal S_k)\rangle
\right|_{\theta_{b1}=\cdots=\theta_{bm_b}=\alpha_b}.
\] Thus the same block support can produce different tangent matrices,
different normal equations, and different active parameter counts.
Below, \(I_{b,k}\) and \(\theta_k\) denote the declared coordinate
convention, and the active coordinate support is \[
J_k:=\{1,\ldots,N_k\},
\qquad
N_k:=\sum_{b\in\mathcal B_k}|I_{b,k}| .
\] The prepared state is \[
|\psi_k\rangle
:=|\psi(\theta_k;\mathcal S_k)\rangle
=U_k(\theta_k)|\psi_{\mathrm{ref}}\rangle,
\qquad
U_k(\theta):=U(\theta;\mathcal S_k).
\] If AP-McLachlan is initialized from a completed static ansatz, the
handoff state is denoted \(|\psi_{\mathrm{initial}}\rangle\). A
same-seed dynamics row requires the declared coordinate vector to
reproduce that state up to ray tolerance: \[
\epsilon_{\mathrm{handoff}}
:=\min_{\varphi\in\mathbb R}
\left\|
U_k(\theta_0)|\psi_{\mathrm{ref}}\rangle
-e^{i\varphi}|\psi_{\mathrm{initial}}\rangle
\right\|_2
\le \tau_{\mathrm{handoff}}.
\] For the per logical/macro-generator convention this check uses the
logical seed coordinates \(\alpha_0\) themselves. It is not valid to
silently replace a per-term seed by averaged logical coordinates unless
the ray check above passes. During propagation the support labels are
fixed and only the driven Hamiltonian coefficients change: \[
H(t)=\sum_{\alpha\in\mathcal I_{\mathrm{stat}}}h_\alpha P_\alpha
+\sum_{\beta\in\mathcal I_{\mathrm{drv}}}c_\beta(t)P_\beta .
\] The drive profile is the scalar coefficient \(c_\beta(t)\); the drive
operator is the Pauli-polynomial channel multiplying it. If a driven
benchmark initializes the ansatz with a zero-angle drive-aligned block
\(D_\beta=\sum_\ell d_{\beta\ell}P_{\beta\ell}\), the Hamiltonian is
unchanged while the tangent plane gains \[
\left.\partial_{\phi_\beta}
e^{-i\phi_\beta D_\beta}|\psi_k\rangle
\right|_{\phi_\beta=0}
=-iD_\beta|\psi_k\rangle,
\qquad
\bar T_{\beta,k}^{\mathrm{drv}}
:=Q_{\psi_k}(-iD_\beta|\psi_k\rangle)
\] as an available McLachlan direction. This initialization changes
\(\mathcal S_0\) and the measured tangent geometry; it does not change
\(H(t)\). For a sampled time \(t_k\) or declared quadrature stage,
define the horizontal tangent block and the projected Schrodinger drift
\[
\bar T_k
:=Q_{\psi_k}\left[\partial_{\theta_{k,1}}|\psi_k\rangle,\ldots,
\partial_{\theta_{k,N_k}}|\psi_k\rangle\right],
\qquad
Q_{\psi_k}:=I-|\psi_k\rangle\langle\psi_k|,
\] \[
\bar b_k
:=-i\bigl(H(t_k)-\langle H(t_k)\rangle_{\psi_k}\bigr)|\psi_k\rangle .
\] The McLachlan metric, force, and regularized normal matrix are \[
G_k:=\Re(\bar T_k^\dagger\bar T_k),
\qquad
f_k:=\Re(\bar T_k^\dagger\bar b_k),
\qquad
K_k:=G_k+\Lambda_k .
\] The regularizer \(\Lambda_k\succeq0\) represents declared numerical
stabilization. It is not a new physical force and it is not a
route-selection score.

For any coordinate support \(S\) inside the active or zero-initialized
augmented tangent support, define the restricted metric, regularized
metric, and force \[
G_{k,SS}:=\Re(\bar T_{k,S}^{\dagger}\bar T_{k,S}),
\qquad
K_{k,SS}:=G_{k,SS}+\Lambda_{k,SS},
\qquad
f_{k,S}:=\Re(\bar T_{k,S}^{\dagger}\bar b_k),
\] and let the structural and regularized support velocities be \[
\dot\theta_{k,S}^{G}:=(G_{k,SS})_\tau^\dagger f_{k,S},
\qquad
\dot\theta_{k,S}^{K}:=(K_{k,SS})_{\tau,\varepsilon}^{\oplus} f_{k,S}.
\] Here \((\cdot)_\tau^\dagger\) is the supported geometric
pseudoinverse and \((\cdot)_{\tau,\varepsilon}^{\oplus}\) is the damped
supported inverse used by the implemented velocity. The scalar \[
f_{k,S}^{\top}(G_{k,SS})_\tau^\dagger f_{k,S}
=f_{k,S}^{\top}\dot\theta_{k,S}^{G}
=
(\dot\theta_{k,S}^{G})^\top G_{k,SS}\dot\theta_{k,S}^{G}
=\|\bar T_{k,S}\dot\theta_{k,S}^{G}\|_2^2
\] is the force vector overlapped with its metric-scaled response. The
retained normal equation is
\(G_{k,SS}\dot\theta_{k,S}^{G}=P_{k,S}^{G,\tau}f_{k,S}\), with
\(P_{k,S}^{G,\tau}:=G_{k,SS}(G_{k,SS})_\tau^\dagger\); if the force lies
inside the retained geometric range, this reduces to
\(f_{k,S}=G_{k,SS}\dot\theta_{k,S}^{G}\). The corresponding regularized
scalar \[
f_{k,S}^{\top}(K_{k,SS})_{\tau,\varepsilon}^{\oplus}f_{k,S}
=f_{k,S}^{\top}\dot\theta_{k,S}^{K}
\] is used consistently for propagation, append gain, prune loss, and
exchange scoring.

The propagated velocity on an open time interval is \[
\dot\theta_k
:=K_{k,\tau,\varepsilon}^{\oplus}f_k,
\qquad
K_k\dot\theta_k^{\mathrm{LS}}=f_k
\quad\text{for the formal least-squares solution.}
\] The parameter ODE between adjacent control time points is therefore
\[
\frac{d\theta}{dt}=F_{\mathcal S_k}(\theta,t),
\qquad
F_{\mathcal S}(\theta,t):=\bigl(G_{\mathcal S}(\theta)+\Lambda_{\mathcal S}(\theta,t)\bigr)_{\tau,\varepsilon}^{\oplus} f_{\mathcal S}(\theta,t),
\] with \(\mathcal S=\mathcal S_k\) held fixed until the next
support-patch decision. The benchmark may declare Euler, RK4, or another
fixed one-step method, but adaptive integrator switching is not part of
the active AP-McLachlan mechanism.

\subsection{17A.2 Structural residual and support
inadequacy}\label{a.2-structural-residual-and-support-inadequacy}

The structural miss is the component of the sampled horizontal
Schrodinger drift that is not represented by the current tangent support
before damping artifacts are interpreted as physics: \[
\begin{aligned}
\epsilon_{\mathrm{expr},k}^2
&:=\operatorname{clip}\!\left(
\|\bar b_k\|_2^2
-f_{k,J_k}^{\top}(G_{k,J_kJ_k})_\tau^\dagger f_{k,J_k},
0,
\|\bar b_k\|_2^2
\right),\\
\rho_{\mathrm{expr},k}
&:=\operatorname{clip}\!\left(
\frac{\epsilon_{\mathrm{expr},k}^2}{\|\bar b_k\|_2^2+\varepsilon},
0,1
\right).
\end{aligned}
\] The realized regularized residual is \[
\begin{aligned}
\epsilon_{\mathrm{real},k}^2
&:=\|\bar T_k K_{k,\tau,\varepsilon}^{\oplus}f_k-\bar b_k\|_2^2,\\
\rho_{\mathrm{num},k}
&:=\left[
\frac{\epsilon_{\mathrm{real},k}^2}{\|\bar b_k\|_2^2+\varepsilon}
-\rho_{\mathrm{expr},k}
\right]_+ .
\end{aligned}
\] Large \(\rho_{\mathrm{expr},k}\) means the ansatz tangent span lacks
physical drift directions and may justify append. Large
\(\rho_{\mathrm{num},k}\) means the solve is failing to realize
available tangent directions and should trigger solve repair, damping
review, support threshold review, or time-step reduction; it is not by
itself evidence for a new ansatz block.

A typical high-miss trigger is \[
\operatorname{HighMiss}_k
:=\mathbf 1\!\left[
\rho_{\mathrm{expr},k}>\tau_{\mathrm{miss}}
\ \wedge\
\epsilon_{\mathrm{expr},k}^2>\epsilon_{\mathrm{abs}}^{\mathrm{miss}}
\right]
\] with optional persistence over recent time points. Persistence may
reduce noise sensitivity, but it cannot import exact target-state
information into the decision rule.

\subsection{17A.3 Zero-initialized append and grouped
gain}\label{a.3-zero-initialized-append-and-grouped-gain}

An append candidate is a support record \(r=(m,p)\): generator block
\(m\) inserted at ordered position \(p\). If the block contains \(m_r\)
Pauli-polynomial components, then \(|I_{r,k}^+|=m_r\) in the per
Pauli/polynomial-term convention and \(|I_{r,k}^+|=1\) in the per
logical/macro-generator convention. Denote the new coordinate set by
\(I_{r,k}^+\). The append is zero-initialized, so \[
|\psi(\theta_k,0;\mathcal S_k^+(r))\rangle=|\psi(\theta_k;\mathcal S_k)\rangle,
\] and the state is continuous at the time point. The new block changes
the tangent plane available for subsequent McLachlan propagation.

The canonical normalized append gain is the grouped forward
extra-sum-of-squares gain in the regularized McLachlan objective: \[
\boxed{
\rho_{r,k}^{\mathrm{gain}}
:=
\frac{
\left[
f_{k,J_k\cup I_{r,k}^+}^{\top}
(K_{k,(J_k\cup I_{r,k}^+)(J_k\cup I_{r,k}^+)})_{\tau,\varepsilon}^{\oplus}
f_{k,J_k\cup I_{r,k}^+}
-
f_{k,J_k}^{\top}
(K_{k,J_kJ_k})_{\tau,\varepsilon}^{\oplus}
f_{k,J_k}
\right]_+
}{
\|\bar b_k\|_2^2+\varepsilon
} .
}
\] For a batched append decision \(R_k^+\), define \[
I_{R,k}^+:=\bigcup_{r\in R_k^+}I_{r,k}^+,
\qquad
J_{R,k}^+:=J_k\cup I_{R,k}^+ .
\] The joint batched gain is \[
\boxed{
\rho_{R,k}^{\mathrm{gain}}
:=
\frac{
\left[
f_{k,J_{R,k}^+}^{\top}
(K_{k,J_{R,k}^+J_{R,k}^+})_{\tau,\varepsilon}^{\oplus}
f_{k,J_{R,k}^+}
-
f_{k,J_k}^{\top}
(K_{k,J_kJ_k})_{\tau,\varepsilon}^{\oplus}
f_{k,J_k}
\right]_+
}{
\|\bar b_k\|_2^2+\varepsilon
} .
}
\] The singleton formula is recovered by \(R_k^+=\{r\}\). Both
quantities measure the amount of regularized horizontal drift newly
explained after the old coordinates and candidate coordinates are
considered together under the same retained-support, ridge, and damping
convention. Static hardware cost may rank candidate evaluation, \[
\Xi_{r,k}^{\mathrm{app}}
:=\frac{\rho_{r,k}^{\mathrm{gain}}}{(\kappa_{r,k}^{\mathrm{hw}})^{\alpha_{\mathrm{app}}}},
\] but the cost-pressure score is not the admission quantity. A
sub-threshold tangent gain is not promoted by being cheap.

For a candidate tangent block \(\bar U_{r,k}\), set \[
B_{r,k}:=\Re(\bar T_k^\dagger\bar U_{r,k}),
\qquad
C_{r,k}:=\Re(\bar U_{r,k}^\dagger\bar U_{r,k}),
\qquad
q_{r,k}:=\Re(\bar U_{r,k}^\dagger\bar b_k).
\] Under the same old-support inverse, ridge, damping, and retained
spectral support as the propagated velocity, the grouped gain can be
evaluated through the Schur block \[
w_{r,k}:=q_{r,k}-B_{r,k}^\top K_{k,\tau,\varepsilon}^{\oplus}f_k,
\qquad
S_{r,k}:=C_{r,k}+\Lambda_r-B_{r,k}^\top K_{k,\tau,\varepsilon}^{\oplus}B_{r,k},
\] \[
\Delta_{r,k}^{\mathrm{Schur}}
:=w_{r,k}^\top(S_{r,k})_\tau^\dagger w_{r,k}.
\] The Schur expression is a compact evaluation or scout/confirm
diagnostic for the grouped gain. The direct grouped before--after
quantity \(\rho_{r,k}^{\mathrm{gain}}\) remains the definition of append
acceptance when whitening, thresholding, or compression makes the
shortcut non-identical.

An append is admitted only from present-time local data: \[
\operatorname{Admit}_{\mathrm{app}}(r,k)=1
\iff
\operatorname{HighMiss}_k=1
\ \wedge\
\rho_{r,k}^{\mathrm{gain}}\ge\tau_{\mathrm{gain}}
\ \wedge\
S_{r,k}^{\mathrm{conf}}\ge\tau_{\mathrm{conf}} .
\] The confirmation score may include Schur consistency, candidate
independence, direction-change control, and measurement budget, but it
may not use exact future fidelity, exact evolved states, or benchmark
reference trajectories.

\subsection{17A.4 Verified prune and grouped deletion
loss}\label{a.4-verified-prune-and-grouped-deletion-loss}

A prune proposal deletes a persistent active block \(b\) with coordinate
set \(I_{b,k}\). Deletion is the backward grouped companion of append.
The canonical normalized full-support deletion loss is \[
\boxed{
L_{b,k}^{\mathrm{full}}
:=
\frac{
\left[
f_{k,J_k}^{\top}
(K_{k,J_kJ_k})_{\tau,\varepsilon}^{\oplus}
f_{k,J_k}
-
f_{k,J_k\setminus I_{b,k}}^{\top}
(K_{k,(J_k\setminus I_{b,k})(J_k\setminus I_{b,k})})_{\tau,\varepsilon}^{\oplus}
f_{k,J_k\setminus I_{b,k}}
\right]_+
}{
\|\bar b_k\|_2^2+\varepsilon
} .
}
\] For a prune batch \(B_k^-\), define \[
I_{B,k}^-:=\bigcup_{b\in B_k^-}I_{b,k},
\qquad
J_{B,k}^-:=J_k\setminus I_{B,k}^- .
\] The batched full-support deletion loss is \[
\boxed{
L_{B,k}^{\mathrm{full}}
:=
\frac{
\left[
f_{k,J_k}^{\top}
(K_{k,J_kJ_k})_{\tau,\varepsilon}^{\oplus}
f_{k,J_k}
-
f_{k,J_{B,k}^-}^{\top}
(K_{k,J_{B,k}^-J_{B,k}^-})_{\tau,\varepsilon}^{\oplus}
f_{k,J_{B,k}^-}
\right]_+
}{
\|\bar b_k\|_2^2+\varepsilon
} .
}
\] The singleton formula is recovered by \(B_k^-=\{b\}\). The loss is
computed using the same support convention, damping, ridge, and
measurement-compatible tangent data as propagation and append. Small
coordinate amplitude, old age, pairwise commutation, or a cheap local
proxy can nominate a candidate, but none of those quantities certifies
deletion.

Prune acceptance requires a reduced-state projection rather than raw
coordinate dropping. Let \(\mathcal S_k^{(-b)}\) be the support with
block \(b\) removed. The reduced parameter point is \[
\theta_k^{(-b)}
:=\arg\min_{\xi\in\Theta_{\mathcal S_k^{(-b)}}}
\left[
 d_{\mathrm{FS}}^2\!\left(\psi(\xi;\mathcal S_k^{(-b)}),\psi(\theta_k;\mathcal S_k)\right)
 +\lambda_Y\|Y(\xi;\mathcal S_k^{(-b)})-Y_k\|_{\Sigma_Y^{-1}}^2
 +\lambda_v\|v(\xi;\mathcal S_k^{(-b)})-v_k\|_2^2
\right],
\] where \(Y_k\) denotes the tracked observable bundle and
\(v_k=\bar T_k\dot\theta_k\) is the represented horizontal velocity. In
strict measurement-compatible routes, every component of \(Y_k\) used
here is estimated from the incumbent or candidate prepared circuit
state; exact/reference observables remain diagnostic side channels.

The prune acceptance predicate has the form \[
\operatorname{Accept}_{\mathrm{pr}}(b,k)=1
\iff
\Gamma_{\mathrm{perm}}^{\mathrm{pr}}(b,k)=1
\ \wedge\
L_{b,k}^{\mathrm{full}}\le\tau_{\mathrm{loss}}^{\mathrm{pr}}
\ \wedge\
\Delta_{\mathrm{ray},b,k}^{\mathrm{pr}}\le\tau_{\mathrm{ray}}^{\mathrm{pr}}
\ \wedge\
\Delta\rho_{b,k}^{\mathrm{pr}}\le\tau_{\rho}^{\mathrm{pr}}
\ \wedge\
\mathfrak C_{b,k}^{\mathrm{shadow}}=1
\ \wedge\
\operatorname{Persist}_{b,k}^{\mathrm{pr}}=1 .
\] Here \(\Gamma_{\mathrm{perm}}^{\mathrm{pr}}\) enforces age, cooldown,
protected-block, solve-health, and candidate-budget rules. The ray and
differential-miss terms prevent state or tangent-span damage. The shadow
certificate propagates the incumbent and reduced candidate for a short
horizon under the same declared integrator and Hamiltonian sampling,
then checks observable no-harm: \[
\mathfrak C_{b,k}^{\mathrm{shadow}}
:=\mathbf 1\!\left[
\max_{1\le h\le H_{\mathrm{pr}}}
\|Y_{k+h}^{(-b)}-Y_{k+h}^{\mathrm{inc}}\|_{\Sigma_Y^{-1}}
+\lambda_{\Delta Y}
\max_{1\le h\le H_{\mathrm{pr}}}
\|\Delta Y_{k+h}^{(-b)}-\Delta Y_{k+h}^{\mathrm{inc}}\|_{\Sigma_{\Delta Y}^{-1}}
\le\tau_{\mathrm{shadow}}^{\mathrm{pr}}
\right].
\] Rejected prune trials roll back exactly and place the tested block on
cooldown. Accepted prune events commit the projected reduced state
\(\bigl(\mathcal S_k^{(-b)},\theta_k^{(-b)}\bigr)\).

\subsection{17A.5 Support-exchange patch and event
map}\label{a.5-support-exchange-patch-and-event-map}

The structural AP-McLachlan object at time point \(k\) is the support
patch \[
\nu_k=(B_k^-,R_k^+),
\] where \(B_k^-\) is a set of active blocks proposed for deletion and
\(R_k^+\) is a set of candidate blocks proposed for zero-initialized
insertion. Empty choices recover stay, pure append, and pure prune.
Nonempty choices on both sides give an exchange patch that replaces
stale tangent coordinates with new tangent coordinates in one event.

For a patch, define \[
I_{\nu,k}^-:=\bigcup_{b\in B_k^-}I_{b,k},
\qquad
I_{\nu,k}^+:=\bigcup_{r\in R_k^+}I_{r,k}^+,
\qquad
J_k^\nu:=(J_k\setminus I_{\nu,k}^-)\cup I_{\nu,k}^+,
\] \[
\Delta_{\nu,k}^{\mathrm{patch}}
:=
\frac{
f_{k,J_k^\nu}^{\top}
(K_{k,J_k^\nu J_k^\nu})_{\tau,\varepsilon}^{\oplus}
f_{k,J_k^\nu}
-
f_{k,J_k}^{\top}
(K_{k,J_kJ_k})_{\tau,\varepsilon}^{\oplus}
f_{k,J_k}
}{
\|\bar b_k\|_2^2+\varepsilon
} .
\] The separate append and prune criteria remain hard side constraints:
every appended block must pass grouped gain and confirmation, every
deleted block must pass full-support loss and nonlinear verification,
and the combined patch score is recomputed on \(J_k^\nu\). Exchange can
therefore be implemented conservatively after pure append and pure prune
are tested; if the combined support check disagrees with the side
scouts, the controller falls back to pure append, pure prune, or stay.

The structural map is \[
\mathsf S_{\nu_k}(\mathcal S_k,\theta_k)=
\begin{cases}
(\mathcal S_k,\theta_k),&B_k^-=\varnothing,\ R_k^+=\varnothing,\\
(\mathcal S_k^+(R_k^+),\theta_k\oplus 0),&B_k^-=\varnothing,
\ R_k^+\ne\varnothing,\\
(\mathcal S_k^{(-B_k^-)},\theta_k^{(-B_k^-)}),&B_k^-\ne\varnothing,
\ R_k^+=\varnothing,\\
((\mathcal S_k^{(-B_k^-)})^+(R_k^+),\theta_k^{(-B_k^-)}\oplus0),&B_k^-\ne\varnothing,
\ R_k^+\ne\varnothing .
\end{cases}
\] After the patch is committed, the next interval is propagated by the
declared fixed integrator: \[
(\mathcal S_k,\theta_k)
\xrightarrow{\mathrm{diagnose/select}}
\nu_k^\star
\xrightarrow{\mathsf S_{\nu_k^\star}}
(\mathcal S_k^\nu,\theta_k^{\nu,+})
\xrightarrow{\operatorname{Step}_{\mathrm{declared}}}
(\mathcal S_{k+1},\theta_{k+1}).
\] Solve repair has priority over structural edits when the tangent
solve is unsupported, badly conditioned, or dominated by numerical
damping. Exact benchmarking may be computed during or after the run for
plots, errors, spectra, and audits, but it has no code path into
\(\nu_k^\star\) in a strict AP-McLachlan row.

\subsection{17A.6 Paper-facing evidence
contract}\label{a.6-paper-facing-evidence-contract}

The dynamics manuscript should describe the active route as AP-McLachlan
support-patch control over time points or time iterations. Legacy
implementation fields may still contain \texttt{checkpoint} for
compatibility, but the mathematical object is the time-indexed support
patch \(\nu_k=(B_k^-,R_k^+)\) and the projective McLachlan velocity on
the active variational manifold.

Every same-seed AP-McLachlan or fixed-McLachlan row must record its
coordinate parameterization. The two supported labels are per
Pauli/polynomial term and per logical/macro generator; implementation
records may abbreviate these as \texttt{per\_pauli\_term} and
\texttt{logical\_shared}. The record should expose the active coordinate
count, the logical block count, and the underlying Pauli-component
count, because appending a block changes \(N_k\) by \(m_r\) in the
per-term convention and by one in the logical convention. Rows using
different coordinate parameterizations are different variational tangent
geometries unless the benchmark explicitly presents them as a
coordinate-mode ablation.

Main-text claims should stay inside the demonstrated evidence surface:

\begin{itemize}
\tightlist
\item
  append is supported when selected rows show accepted zero-initialized
  tangent-block additions and improved grouped McLachlan support score;
\item
  prune is supported only when selected rows show committed verified
  deletions under the full-support loss, projection, shadow,
  persistence, cooldown, and rollback checks;
\item
  drive-aligned ansatz initialization is reported as a zero-angle
  support augmentation and not as a Hamiltonian modification;
\item
  logical/macro-generator rows are supported only when the handoff
  parity check against \(|\psi_{\mathrm{initial}}\rangle\) passes;
  otherwise the benchmark must fail loudly or use the per-term
  coordinate convention;
\item
  integrator switching, secant forecasts, turbulent-regime labels, and
  large-amplitude theta heuristics are archival diagnostics, not active
  Paper-II implementation guidance;
\item
  exact/reference trajectories are valid diagnostic comparators for
  energy, observables, fidelity, spectra, and plotting, but are
  forbidden as controller-decision inputs in strict AP-McLachlan runs.
\end{itemize}

For the present same-seed dynamics surface, the paper-facing language
should reflect the measured support-edit pattern: append can be
discussed as an active AP-McLachlan mechanism; pruning should be
presented as a verified controller capability and ablation surface
unless the selected evidence rows actually commit prune events.

\section{18. Geometry-Selected QSE, Spectral Functions, and
Excited-State
Response}\label{geometry-selected-qse-spectral-functions-and-excited-state-response}

This chapter is the source-of-truth mathematical layer for the
excited-state and spectral-response route. The purpose is to define the
objects before any numerical table is trusted: the nonorthogonal QSE
matrix pencil, the selected excitation record set, the spectral weights,
the time-domain correlations, the Green-function sectors, and the
frozen/live dynamics handoff. Results belong later; this chapter defines
what those results are supposed to measure.

The high-level construction is a Ritz reduction over physically
generated excitation kets. The ground-state circuit prepares a seed
state, the excitation registry proposes operators that act on that seed,
and QSE diagonalizes \(H_0\) inside the span of the resulting kets. The
central mathematical issue is metric redundancy: the selected kets
generally satisfy \(S\ne I\) and can satisfy
\({\rm rank}(S)<|\mathcal A|\). Consequently every coefficient vector is
interpreted through the overlap metric \(S\), every reported root is
checked for conditioning and transition visibility, and every spectral
claim specifies the physical probe that makes the root observable.

The optimized ground-state ansatz supplies one prepared seed state, \[
|\Phi_0\rangle:=|\mathbf z_{\rm ref}\rangle,\qquad
V_0:=\prod_{j=1}^{m_0}\exp[-i\theta_j^\star A_j],
\qquad
|\psi_0\rangle:=V_0|\Phi_0\rangle .
\] The excitation-record source is the same problem-local
\texttt{full\_meta} registry used by the adaptive ground-state program:
a candidate record \(r\) carries an encoded operator \(\Omega_r\),
sector/channel metadata, and measurement or compiled-cost metadata. QSE
is defined only after a retained record set \(\mathcal A\) has been
chosen from that registry.

The section has four layers. First, QSE turns selected excitation
records into a nonorthogonal matrix pencil. Second, spectral functions
and correlations turn the QSE roots into observable response. Third,
Green functions use the same subspace logic in particle-addition and
particle-removal sectors. Fourth, frozen/adaptive/live dynamics test
whether the selected spectral manifold remains closed under a drive.

The same tradeoff appears in all four layers. Adding records can improve
the physical span, reveal missing spectral weight, and reduce residual
error; it also enlarges the matrix pencil, worsens near-linear
dependence, increases measurement burden, and can make finite-shot root
tracking unstable. The method is a controlled construction of a
measurable, well-conditioned spectral manifold.

Interpretive compression: Section 18 is an operator-geometric reduction
of excited-state physics. The general object is a response-relevant
finite subspace \(\mathcal K_{\mathcal A}\subset\mathcal H\) generated
by records acting on a seed; the specific algebraic machinery is a
Hermitian generalized Ritz pencil \((M,S)\), its supported metric
normalization, and probe-dependent channel functionals applied to the
resulting Ritz roots. The physics is encoded in which operators generate
the subspace, which probes interrogate it, and which metric directions
survive measurement-compatible conditioning.

\subsection{18.1 Excitation kets and the nonorthogonal QSE
pencil}\label{excitation-kets-and-the-nonorthogonal-qse-pencil}

The primitive object is the generated excitation ket
\(Q_{\mathfrak s,0}\Omega_a|\psi_0\rangle\). Two different records can
generate nearly the same ket after acting on \(|\psi_0\rangle\), while a
record with a complicated operator expression can be physically dark for
the probes of interest. The QSE pencil is determined by the
Hilbert-space vectors \(|\phi_a\rangle\), their overlaps, their
Hamiltonian matrix elements, and their transition shadows against the
seed.

Fix the time-independent Hamiltonian \(H_0\) and a target symmetry
sector \(\mathfrak s\). Let \(\Pi_{\mathfrak s}\) be the sector
projector. For neutral same-sector excitations the seed component is
removed by
\(Q_{\mathfrak s,0}:=\Pi_{\mathfrak s}(I-|\psi_0\rangle\langle\psi_0|)\Pi_{\mathfrak s}\).
In particle-changing Green-function sectors,
\(|\psi_0\rangle\notin\mathcal H_{N\pm1}\), and the sector projector
replaces \(Q_{\mathfrak s,0}\).

Each retained record \(a\in\mathcal A\) supplies one unnormalized
excitation ket
\(|\phi_a\rangle:=Q_{\mathfrak s,0}\Omega_a|\psi_0\rangle\). With
\(N:=|\mathcal A|\), collect the kets and their projected matrices as \[
\begin{aligned}
\Phi_{\mathcal A}
&:=\begin{bmatrix}|\phi_{a_1}\rangle&\cdots&|\phi_{a_M}\rangle\end{bmatrix},\\
S_{ab}
&:=\langle\phi_a|\phi_b\rangle
=(\Phi_{\mathcal A}^\dagger\Phi_{\mathcal A})_{ab},\\
M_{ab}
&:=\langle\phi_a|H_0|\phi_b\rangle
=(\Phi_{\mathcal A}^\dagger H_0\Phi_{\mathcal A})_{ab}.
\end{aligned}
\] Thus \(S\) is the metric of the nonorthogonal excitation coordinates
and \(M\) is the Hamiltonian restricted to the same coordinate system.

This is the finite-dimensional version of a variational subspace
problem. The column matrix \(\Phi_{\mathcal A}\) maps coefficient
coordinates into the physical Hilbert space; \(S\) measures lengths and
angles after that map, while \(M\) measures the Hamiltonian quadratic
form after that same map. If two coefficient vectors differ by a
direction in the null space of \(S\), they synthesize the same physical
ket, so the coefficient space has more coordinates than physical
directions.

For \(c\in\mathbb C^N\), the synthesized QSE state, norm, energy
numerator, and Ritz quotient are \[
\begin{aligned}
|\Psi(c)\rangle
&:=\Phi_{\mathcal A}c
=\sum_{a\in\mathcal A}c_a|\phi_a\rangle,\\
\langle\Psi(c)|\Psi(c)\rangle
&=c^\dagger S c,\qquad
\langle\Psi(c)|H_0|\Psi(c)\rangle
=c^\dagger M c,\\
\mathcal R(c)
&:=\frac{c^\dagger M c}{c^\dagger S c}.
\end{aligned}
\] The coefficient vector is therefore a synthesis coordinate: it tells
how the selected nonorthogonal excitation kets are mixed before the
physical norm is evaluated.

The quotient \(\mathcal R(c)\) asks the usual variational question:
among all nonzero states that can be synthesized from the selected
records, which directions make the Hamiltonian expectation stationary?
Since \(S\ne I\) in a generic nonorthogonal record set, the stationary
equation is a generalized eigenproblem in \(c\). The metric \(S\) turns
coefficient algebra back into Hilbert-space geometry.

Impose \(c^\dagger S c=1\) and vary the constrained functional \[
\begin{aligned}
\mathcal L(c,\varepsilon)
&:=c^\dagger M c-\varepsilon(c^\dagger S c-1),\\
0
&=\partial_{c^\dagger}\mathcal L
=M c-\varepsilon S c .
\end{aligned}
\] The QSE roots are the generalized Ritz pairs \[
\begin{aligned}
M c^{(\nu)}
&=\varepsilon_\nu S c^{(\nu)},&
(c^{(\mu)})^\dagger S c^{(\nu)}
&=\delta_{\mu\nu},\\
|\Psi_\nu^{\rm QSE}\rangle
&:=\Phi_{\mathcal A}c^{(\nu)}
=\sum_{a\in\mathcal A}c_a^{(\nu)}|\phi_a\rangle,&
\omega_\nu^{\rm QSE}
&:=\varepsilon_\nu-E_0,
\end{aligned}
\] where \(E_0:=\langle\psi_0|H_0|\psi_0\rangle\). The generalized
eigenproblem is forced by the nonorthogonal metric: physical
normalization is \(c^\dagger S c=1\), and Euclidean normalization
matches it exactly in the special case \(S=I\).

\subsection{18.2 Supported metric whitening and finite-shot
stability}\label{supported-metric-whitening-and-finite-shot-stability}

\subsubsection{§18.2 lesson: supported metric whitening in
QSE}\label{lesson-supported-metric-whitening-in-qse}

Interpretive compression: metric whitening is a coordinate gauge on a
finitely generated, nonorthogonal Ritz space. QSE takes a finite set of
excitation kets \(\{|\phi_a\rangle\}_{a\in\mathcal A}\), forms their
span, and solves the Hamiltonian only through the pullback geometry
induced by the synthesis map. The synthesized ket is a state of the
original physical Hilbert space restricted to the finite QSE subspace \[
\mathcal K_{\mathcal A}:={\rm range}(\Phi_{\mathcal A})
={\rm span}\{|\phi_a\rangle:a\in\mathcal A\}\subset\mathcal H .
\] The coordinate and physical objects are \[
c\in\mathbb C^N,\qquad
|\Psi(c)\rangle=\Phi_{\mathcal A}c\in\mathcal H.
\] The pair \((\mathbb C^N,\langle c,d\rangle_S:=c^\dagger S d)\) is the
finite coordinate Hilbert model of \(\mathcal K_{\mathcal A}\). In that
coordinate model, \(S\) supplies the inner product and \(M\) supplies
the Hamiltonian quadratic form.

The index expansion has one job: it proves which matrix implements
physical length and energy after the coefficient vector is embedded into
the physical Hilbert space. Let \[
\mathcal A=\{a_1,\ldots,a_N\},\qquad
|\phi_i\rangle:=|\phi_{a_i}\rangle,\qquad i=1,\ldots,N,
\] and collect the kets as the synthesis map \[
\Phi_{\mathcal A}:=
\begin{bmatrix}
|\phi_1\rangle&|\phi_2\rangle&\cdots&|\phi_N\rangle
\end{bmatrix}.
\] For \(c\in\mathbb C^N\), \[
|\Psi(c)\rangle
:=\Phi_{\mathcal A}c
=\sum_{j=1}^N c_j|\phi_j\rangle .
\] Here \(|\Psi(c)\rangle\) is the physical QSE trial state described
inside the selected quantum subspace. The vector \(c\) is its coordinate
representation in the excitation-ket coordinate system.

Define the overlap and projected Hamiltonian matrices by \[
S_{ij}:=\langle\phi_i|\phi_j\rangle,\qquad
M_{ij}:=\langle\phi_i|H_0|\phi_j\rangle .
\] Then \[
\begin{aligned}
\langle\Psi(c)|\Psi(c)\rangle
&=
\left(\sum_{i=1}^N c_i|\phi_i\rangle\right)^\dagger
\left(\sum_{j=1}^N c_j|\phi_j\rangle\right)\\
&=
\left(\sum_{i=1}^N c_i^*\langle\phi_i|\right)
\left(\sum_{j=1}^N c_j|\phi_j\rangle\right)\\
&=
\sum_{i=1}^N\sum_{j=1}^N c_i^*c_j\langle\phi_i|\phi_j\rangle\\
&=
\sum_{i=1}^N\sum_{j=1}^N c_i^*S_{ij}c_j
=c^\dagger S c .
\end{aligned}
\] Thus \(c^\dagger S c=1\) means the synthesized ket
\(|\Psi(c)\rangle\) has physical norm one. The Euclidean coordinate norm
\(c^\dagger c\) measures coefficient size; the Gram norm
\(c^\dagger S c\) measures physical ket size.

Similarly, \[
\begin{aligned}
\langle\Psi(c)|H_0|\Psi(c)\rangle
&=
\left(\sum_{i=1}^N c_i^*\langle\phi_i|\right)
H_0
\left(\sum_{j=1}^N c_j|\phi_j\rangle\right)\\
&=
\sum_{i=1}^N\sum_{j=1}^N c_i^*c_j
\langle\phi_i|H_0|\phi_j\rangle\\
&=
\sum_{i=1}^N\sum_{j=1}^N c_i^*M_{ij}c_j
=c^\dagger M c .
\end{aligned}
\] Thus \(M\) is the matrix representation of \(H_0\) on the finite QSE
subspace in the raw excitation coordinates, and \(c^\dagger M c\) is the
energy numerator of the physical ket \(\Phi_{\mathcal A}c\).

The QSE Rayleigh quotient is therefore \[
\mathcal R(c)
:=
\frac{\langle\Psi(c)|H_0|\Psi(c)\rangle}
{\langle\Psi(c)|\Psi(c)\rangle}
=\frac{c^\dagger M c}{c^\dagger S c}.
\] The variational problem is the ordinary Ritz problem on
\(\mathcal K_{\mathcal A}\) written in a nonorthogonal coordinate chart.

The matrix types follow from the same definitions: \[
\begin{aligned}
S_{ji}
&=\langle\phi_j|\phi_i\rangle
=\langle\phi_i|\phi_j\rangle^*
=S_{ij}^*,\\
c^\dagger S c
&=\langle\Psi(c)|\Psi(c)\rangle\ge0,
\end{aligned}
\] so \(S=S^\dagger\succeq0\). It is a Gram matrix. If
\(H_0=H_0^\dagger\), then \[
\begin{aligned}
M_{ji}
&=\langle\phi_j|H_0|\phi_i\rangle\\
&=\langle\phi_i|H_0^\dagger|\phi_j\rangle^*
=\langle\phi_i|H_0|\phi_j\rangle^*
=M_{ij}^*,
\end{aligned}
\] so \(M=M^\dagger\). The positive-semidefinite structure belongs to
\(S\). The projected Hamiltonian \(M\) is Hermitian, and its signs
depend on the chosen energy zero and on the spectrum of \(H_0\) over the
selected subspace.

Calling \((M,S)\) a Hermitian generalized eigenproblem combines two
facts. First, \(M\) and \(S\) have Hermitian symmetry, with \(S\) acting
as the coordinate metric. Second, stationarity of the \(S\)-normalized
Rayleigh quotient supplies the generalized eigen-equation. Impose \[
c^\dagger S c=1
\] and vary \[
\mathcal L(c,c^\dagger,\varepsilon)
:=
c^\dagger M c-\varepsilon(c^\dagger S c-1).
\] In indices, \[
\mathcal L
=
\sum_{i,j=1}^N c_i^*M_{ij}c_j
-\varepsilon\left(
\sum_{i,j=1}^N c_i^*S_{ij}c_j-1
\right).
\] Treat \(c_j\) and \(c_i^*\) as independent variables in the complex
variational derivative. For fixed \(k\), \[
\begin{aligned}
\frac{\partial}{\partial c_k^*}
\sum_{i,j=1}^N c_i^*M_{ij}c_j
&=
\sum_{i,j=1}^N\delta_{ik}M_{ij}c_j
=\sum_{j=1}^N M_{kj}c_j
=(Mc)_k,\\
\frac{\partial}{\partial c_k^*}
\sum_{i,j=1}^N c_i^*S_{ij}c_j
&=
\sum_{i,j=1}^N\delta_{ik}S_{ij}c_j
=\sum_{j=1}^N S_{kj}c_j
=(Sc)_k .
\end{aligned}
\] Stationarity gives \[
0=\frac{\partial\mathcal L}{\partial c_k^*}
=(Mc)_k-\varepsilon(Sc)_k
\] for every \(k=1,\ldots,N\), hence \[
Mc=\varepsilon Sc .
\] This is the nonorthogonal QSE pencil. Positive semidefiniteness of
\(S\) supplies the Gram diagonalization and supported inverse; the
constrained Ritz stationarity supplies the generalized eigen-equation.

The generalized normalization convention says that the synthesized
physical QSE roots are orthonormal after the metric \(S\) evaluates
their physical inner product: \[
\begin{aligned}
\langle\Psi_\mu^{\rm QSE}|\Psi_\nu^{\rm QSE}\rangle
&=
\langle\Phi_{\mathcal A}c^{(\mu)}|\Phi_{\mathcal A}c^{(\nu)}\rangle\\
&=(c^{(\mu)})^\dagger\Phi_{\mathcal A}^\dagger\Phi_{\mathcal A}c^{(\nu)}
=(c^{(\mu)})^\dagger S c^{(\nu)}
=\delta_{\mu\nu}.
\end{aligned}
\] In the raw coordinate dot product the same coefficient vectors can
have different Euclidean overlaps. In the physical QSE subspace, \(S\)
supplies the inner product, so \(S\)-orthonormal coefficients synthesize
orthonormal kets.

Since \(S=S^\dagger\succeq0\), diagonalize the known overlap matrix: \[
S=U\,{\rm diag}(s_1,\ldots,s_N)\,U^\dagger,\qquad
U^\dagger U=UU^\dagger=I,\qquad
s_j\ge0 .
\] This factorization represents the \(N\) coupled eigenpair equations
\[
U=
\begin{bmatrix}
u_1&\cdots&u_N
\end{bmatrix},
\qquad
S u_j=s_j u_j,\qquad
u_i^\dagger u_j=\delta_{ij}.
\] The entries of \(S\) are known matrix elements. The diagonalization
returns the eigenvalues \(s_j\) and the coefficient-space eigenvectors
\(u_j\) together, up to phases and rotations inside degenerate
eigenspaces. Each \(u_j\) is a special coefficient direction; its
associated physical mode is \[
\Phi_{\mathcal A}u_j
=\sum_{i=1}^N (u_j)_i|\phi_i\rangle .
\] If \(c=\alpha u_j\), then \[
\begin{aligned}
c^\dagger S c
&=(\alpha u_j)^\dagger S(\alpha u_j)
=|\alpha|^2 u_j^\dagger S u_j\\
&=|\alpha|^2u_j^\dagger(s_j u_j)
=|\alpha|^2s_j .
\end{aligned}
\] Thus \(s_j\) is the physical norm squared per unit coefficient
amplitude along the metric eigen-direction \(u_j\). Tiny \(s_j\) means a
large coefficient displacement produces a tiny physical ket
displacement, so inverse factors \(s_j^{-1}\) and \(s_j^{-1/2}\) amplify
finite-shot errors.

The support cut keeps the trusted metric eigen-directions: \[
J_\tau:=\{j:s_j>\tau_S\},\qquad
r:=|J_\tau|,\qquad
J_\tau=\{j_1,\ldots,j_r\}.
\] Collect those eigenvectors and eigenvalues as \[
U_\tau:=
\begin{bmatrix}
u_{j_1}&\cdots&u_{j_r}
\end{bmatrix}
\in\mathbb C^{N\times r},
\qquad
D_\tau:={\rm diag}(s_{j_1},\ldots,s_{j_r})
\in\mathbb R^{r\times r}.
\] Then \[
\begin{aligned}
U_\tau^\dagger U_\tau&=I_r,\\
U_\tau^\dagger S U_\tau
&=
\left[
u_{j_p}^\dagger S u_{j_q}
\right]_{p,q=1}^r\\
&=
\left[
u_{j_p}^\dagger s_{j_q}u_{j_q}
\right]_{p,q=1}^r
=\left[s_{j_q}\delta_{pq}\right]_{p,q=1}^r
=D_\tau .
\end{aligned}
\] The support projector \[
P_\tau:=U_\tau U_\tau^\dagger
\] is the Euclidean projector in coefficient space onto the span of the
retained metric eigenvectors. Applying \(P_\tau\) to a raw coefficient
vector keeps its components along trusted directions: \[
P_\tau c
=\sum_{p=1}^r u_{j_p}\,u_{j_p}^\dagger c .
\] Solving on retained support means \[
c=P_\tau c=U_\tau x,\qquad x:=U_\tau^\dagger c\in\mathbb C^r .
\] The corresponding physical ket is \(\Phi_{\mathcal A}U_\tau x\), a
state in the original Hilbert space synthesized from retained metric
modes.

Substitute \(c=U_\tau x\) into the norm and Hamiltonian numerator: \[
\begin{aligned}
c^\dagger S c
&=(U_\tau x)^\dagger S(U_\tau x)
=x^\dagger U_\tau^\dagger S U_\tau x
=x^\dagger D_\tau x\\
&=\sum_{p=1}^r s_{j_p}|x_p|^2,\\
c^\dagger M c
&=(U_\tau x)^\dagger M(U_\tau x)
=x^\dagger(U_\tau^\dagger M U_\tau)x .
\end{aligned}
\] The support-cut Rayleigh quotient is \[
\mathcal R_\tau(x)
=
\frac{x^\dagger(U_\tau^\dagger M U_\tau)x}
{x^\dagger D_\tau x}.
\] This quotient is already restricted to the trusted support. Whitening
is the next coordinate choice: it rescales \(x\) so the denominator
becomes an ordinary Euclidean norm.

Define the exact supported Lowdin whitening map \[
W_\tau:=U_\tau D_\tau^{-1/2},\qquad
D_\tau^{-1/2}:={\rm diag}(s_{j_1}^{-1/2},\ldots,s_{j_r}^{-1/2}),
\qquad
c=W_\tau y .
\] The map \(W_\tau\) is built entirely from the eigenvectors and
eigenvalues of \(S\) on the retained support. It maps a whitened
coordinate vector \(y\in\mathbb C^r\) into the original \(N\)-component
record-coordinate vector \(c\). The associated physical ket is \[
|\Psi(y)\rangle=\Phi_{\mathcal A}W_\tau y .
\] Compute the whitened metric: \[
\begin{aligned}
W_\tau^\dagger S W_\tau
&=
(U_\tau D_\tau^{-1/2})^\dagger S(U_\tau D_\tau^{-1/2})\\
&=
D_\tau^{-1/2}U_\tau^\dagger S U_\tau D_\tau^{-1/2}
=D_\tau^{-1/2}D_\tau D_\tau^{-1/2}\\
&=
{\rm diag}(s_{j_p}^{-1/2}s_{j_p}s_{j_p}^{-1/2})_{p=1}^r
=I_r .
\end{aligned}
\] Therefore \[
c^\dagger S c
=(W_\tau y)^\dagger S(W_\tau y)
=y^\dagger W_\tau^\dagger S W_\tau y
=y^\dagger y .
\] This is the precise meaning of whitening: the physical overlap metric
becomes the identity in the retained \(y\)-coordinates. The QSE state
remains the embedded ket \(\Phi_{\mathcal A}W_\tau y\), while the
coefficient geometry becomes Euclidean.

The Hamiltonian numerator becomes \[
\begin{aligned}
c^\dagger M c
&=(W_\tau y)^\dagger M(W_\tau y)
=y^\dagger W_\tau^\dagger M W_\tau y,\\
\widetilde M_\tau
&:=W_\tau^\dagger M W_\tau .
\end{aligned}
\] Since \(M=M^\dagger\), \[
\widetilde M_\tau^\dagger
=(W_\tau^\dagger M W_\tau)^\dagger
=W_\tau^\dagger M^\dagger W_\tau
=W_\tau^\dagger M W_\tau
=\widetilde M_\tau .
\] The whitened Rayleigh quotient is \[
\mathcal R_\tau(y)
=
\frac{y^\dagger\widetilde M_\tau y}{y^\dagger y}.
\] Stationarity under \(y^\dagger y=1\) gives the ordinary Hermitian
eigenproblem \[
\widetilde M_\tau y^{(\nu)}=\varepsilon_\nu y^{(\nu)},\qquad
c^{(\nu)}=W_\tau y^{(\nu)},\qquad
|\Psi_\nu^{\rm QSE}\rangle=\Phi_{\mathcal A}W_\tau y^{(\nu)} .
\] The separate diagonalization of \(\widetilde M_\tau\) finds the QSE
roots after the \(S\)-metric has been flattened. Thus \(W_\tau\) comes
from \(S\), and the final root vectors \(y^{(\nu)}\) come from
\(\widetilde M_\tau\).

The supported inverse is the inverse metric corresponding to the same
retained coordinates: \[
S_\tau^+:=U_\tau D_\tau^{-1}U_\tau^\dagger .
\] The whitening map satisfies \[
\begin{aligned}
W_\tau W_\tau^\dagger
&=(U_\tau D_\tau^{-1/2})(U_\tau D_\tau^{-1/2})^\dagger\\
&=U_\tau D_\tau^{-1/2}D_\tau^{-1/2}U_\tau^\dagger
=U_\tau D_\tau^{-1}U_\tau^\dagger
=S_\tau^+ .
\end{aligned}
\] The map \(W_\tau\) transports whitened coordinates into record
coordinates; \(S_\tau^+\) acts directly as the supported inverse metric
in the original record-coordinate system.

The damping variant replaces the exact inverse square root by a softened
inverse square root: \[
B_\tau^{(\lambda)}:=U_\tau(D_\tau+\lambda_S I)^{-1/2}.
\] The retained metric in those coordinates is \[
\begin{aligned}
G_\tau^{(\lambda)}
&:=\bigl(B_\tau^{(\lambda)}\bigr)^\dagger S B_\tau^{(\lambda)}\\
&=
(D_\tau+\lambda_S I)^{-1/2}
U_\tau^\dagger S U_\tau
(D_\tau+\lambda_S I)^{-1/2}\\
&=
(D_\tau+\lambda_S I)^{-1/2}
D_\tau
(D_\tau+\lambda_S I)^{-1/2}\\
&=
D_\tau(D_\tau+\lambda_S I)^{-1}
={\rm diag}\!\left(\frac{s_j}{s_j+\lambda_S}\right)_{j\in J_\tau}.
\end{aligned}
\] With \[
M_\tau^{(\lambda)}
:=
\bigl(B_\tau^{(\lambda)}\bigr)^\dagger M B_\tau^{(\lambda)},
\] the damped retained coordinate equation is \[
M_\tau^{(\lambda)}y=\varepsilon\,G_\tau^{(\lambda)}y .
\] The exact-whitening case is recovered at \(\lambda_S=0\), where
\(B_\tau^{(\lambda)}=W_\tau\) and \(G_\tau^{(\lambda)}=I_r\). Damping
stabilizes weak retained modes by replacing \(s_j^{-1/2}\) with
\((s_j+\lambda_S)^{-1/2}\); the support threshold \(\tau_S\) controls
which directions enter the retained subspace.

The finite-shot reason for support and whitening is visible entrywise.
In the exact whitened basis, \[
\begin{aligned}
(\widetilde M_\tau)_{pq}
&=
(W_\tau^\dagger M W_\tau)_{pq}\\
&=
\frac{u_{j_p}^\dagger M u_{j_q}}
{\sqrt{s_{j_p}s_{j_q}}}.
\end{aligned}
\] A finite-shot perturbation \(\delta M\) enters as \[
(\delta\widetilde M_\tau)_{pq}
=
\frac{u_{j_p}^\dagger(\delta M)u_{j_q}}
{\sqrt{s_{j_p}s_{j_q}}}
\] before accounting for perturbations in \(U_\tau\) and \(D_\tau\).
Small retained overlap eigenvalues amplify the same matrix-element
noise.

Root separation controls how that matrix noise rotates eigenvectors. For
\[
\widetilde M_\tau y_\nu=\varepsilon_\nu y_\nu,\qquad
y_\mu^\dagger y_\nu=\delta_{\mu\nu},
\] first-order perturbation gives \[
\delta y_\nu
=
\sum_{\mu\ne\nu}
\frac{y_\mu^\dagger\delta\widetilde M_\tau y_\nu}
{\varepsilon_\nu-\varepsilon_\mu}\,y_\mu
+\text{gauge/normalization component}.
\] Thus \[
\|\delta y_\nu\|
\lesssim
\frac{\|\delta\widetilde M_\tau\|}{\delta_\nu},
\qquad
\delta_\nu:=\min_{\mu\ne\nu}|\varepsilon_\mu-\varepsilon_\nu| .
\] The retained condition and root-separation diagnostics are therefore
\[
\kappa_\tau(S)
:=\frac{\max_{j\in J_\tau}s_j}{\min_{j\in J_\tau}s_j},
\qquad
\delta_\nu:=\min_{\mu\ne\nu}|\varepsilon_\mu-\varepsilon_\nu| .
\] Small retained \(s_j\) increases \(\kappa_\tau(S)\), and small
\(\delta_\nu\) lets finite-shot perturbations rotate nearby roots. These
two quantities control root flips, unstable transition amplitudes, and
spurious spectral peaks.

\subsection{18.3 Transition amplitudes and spectral
functions}\label{transition-amplitudes-and-spectral-functions}

Section 18.2 builds the potential structure: a supported QSE subspace, a
stable coordinate system on that subspace, and approximate roots of the
projected Hamiltonian. Section 18.3 turns that structure into results. A
spectrum asks which of those roots are visible to a physical probe and
at which excitation frequencies that visibility appears.

Interpretive compression: a QSE eigensolve produces candidate poles,
while a probe vector supplies the channel functional that assigns
residues to those poles. Eigenvalues alone describe the algebraic Ritz
ladder; transition amplitudes describe the representation of a
physically prepared disturbance \(O|\psi_0\rangle\) in the dual basis of
QSE roots. A spectral function is therefore the probe-filtered spectral
measure induced by the pair \((H_0,O)\) after projection into
\(\mathcal K_{\mathcal A}\).

Physical semantics: a probe \(O\) is the operator conjugate to a weak
source amplitude \(\lambda(t)\) in a perturbed Hamiltonian
\(H_\lambda(t)=H_0-\lambda(t)O\). Linear response studies the
first-order dependence of measured observables on \(\lambda\), while the
zero-source spectral function records the latent response channels
already encoded in \((H_0,|\psi_0\rangle,O)\). A pole identifies an
energy at which the system has an available excitation; the residue or
intensity identifies how strongly the source \(O\) couples the seed to
that excitation. Thus \(O\) specifies the physical channel,
\(\omega_\nu\) specifies the energetic threshold, \(t_\nu^O\) specifies
the source-to-mode amplitude, and \(I_\nu^O\) specifies the plotted
spectral weight.

\subsubsection{§18.3 lesson: transition amplitudes and spectral
functions}\label{lesson-transition-amplitudes-and-spectral-functions}

A ket is a vector in the encoded Hilbert space. After fermion, boson,
and qubit encodings, the same ket can be represented by amplitudes over
computational bitstrings, but its physical meaning comes from the
Hamiltonian, observables, and inner products defined on that Hilbert
space. A QSE route is the workflow \[
|\psi_0\rangle
\xrightarrow{\Omega_a}
\{|\phi_a\rangle\}_{a\in\mathcal A}
\xrightarrow{S,M}
\{|\Psi_\nu^{\rm QSE}\rangle,\omega_\nu^{\rm QSE}\}_\nu
\xrightarrow{O}
A_O^{\mathcal A}(\omega),
\] where the first half constructs a reduced quantum subspace and the
second half asks what that subspace predicts for a chosen response
channel.

Physical probes matter because they define what an experiment or
diagnostic is asking. Density operators see charge rearrangement and
charge-order excitations. Phonon displacement operators see vibrational
motion and sidebands. Current operators see optical or transport
response. Particle-addition and particle-removal operators see
tunneling, photoemission, spectral gaps, and quasiparticle dressing.
Mixed electron-phonon probes test joint electronic-lattice dressing.

Let \[
\mathcal A=\{a_1,\ldots,a_N\},\qquad
|\phi_i\rangle:=|\phi_{a_i}\rangle,\qquad
\Phi_{\mathcal A}:\mathbb C^N\to\mathcal H,\qquad
\Phi_{\mathcal A}c=\sum_{i=1}^N c_i|\phi_i\rangle .
\] The synthesis map \(\Phi_{\mathcal A}\) takes a finite coefficient
vector and embeds it as a physical ket in the original Hilbert space.
The columns \(|\phi_i\rangle\) are the selected excitation kets; a
coefficient vector says how to mix those kets.

For a QSE root, \[
|\Psi_\nu^{\rm QSE}\rangle
:=\Phi_{\mathcal A}c^{(\nu)}
=\sum_{a\in\mathcal A}c_a^{(\nu)}|\phi_a\rangle .
\] The root ket is the approximate answer mode produced by the QSE
eigensolve. The coefficient vector \(c^{(\nu)}\) gives its coordinates
in the raw excitation-ket system; in supported whitened coordinates
\(c^{(\nu)}=W_\tau y^{(\nu)}\).

Choose a probe operator \(O\). The disturbed seed vector is \[
|\chi_O\rangle:=O|\psi_0\rangle .
\] This vector is the physical question asked of the ground state by the
probe. The seed-to-basis transition vector records the shadows of this
disturbed vector against the QSE excitation bras: \[
d_i^O
:=\langle\phi_i|O|\psi_0\rangle
=\langle\phi_i|\chi_O\rangle,\qquad
d^O:=\Phi_{\mathcal A}^\dagger O|\psi_0\rangle .
\] The vector \(d^O\) has one component per excitation ket. It is the
probe vector in the raw QSE basis. It is formed from overlaps with
\(O|\psi_0\rangle\). The role of \(S\) is to evaluate inner products
among excitation-coordinate vectors after they are embedded by
\(\Phi_{\mathcal A}\).

If \(\{|r\rangle\}_{r=1}^D\) is the encoded computational basis, then
the same component can be written as \[
\begin{aligned}
|\phi_i\rangle&=\sum_{r=1}^D(\phi_i)_r|r\rangle,\qquad
|\psi_0\rangle=\sum_{s=1}^D(\psi_0)_s|s\rangle,\qquad
O=\sum_{r,s=1}^D O_{rs}|r\rangle\langle s|,\\
d_i^O
&=\langle\phi_i|O|\psi_0\rangle
=\sum_{r,s=1}^D(\phi_i)_r^*O_{rs}(\psi_0)_s .
\end{aligned}
\] This line connects the qubit-bitstring representation to the QSE
coordinate representation: the full encoded vector is compressed into
one transition component by the overlap with \(\langle\phi_i|\).

The transition amplitude into root \(\nu\) is the scalar overlap between
the probed seed and that synthesized QSE root: \[
\begin{aligned}
t_\nu^O
&:=\langle\Psi_\nu^{\rm QSE}|O|\psi_0\rangle\\
&=
\left(
\sum_{a\in\mathcal A}c_a^{(\nu)}|\phi_a\rangle
\right)^\dagger O|\psi_0\rangle\\
&=
\sum_{a\in\mathcal A}(c_a^{(\nu)})^*
\langle\phi_a|O|\psi_0\rangle\\
&=
(c^{(\nu)})^\dagger d^O .
\end{aligned}
\] The transition vector \(d^O\) contains the probe's shadow against
every selected excitation ket. The transition amplitude \(t_\nu^O\)
contracts that vector with one QSE root. Thus \(d^O\) is channel-level
information, while \(t_\nu^O\) is root-level visibility for that
channel.

Using the supported-whitened coordinate \(c^{(\nu)}=W_\tau y^{(\nu)}\),
\[
t_\nu^O
=(W_\tau y^{(\nu)})^\dagger d^O
=(y^{(\nu)})^\dagger W_\tau^\dagger d^O .
\] Define the retained whitened probe vector \[
g^O:=W_\tau^\dagger d^O .
\] Then \[
t_\nu^O=(y^{(\nu)})^\dagger g^O .
\] The vector \(g^O\) is the probe represented in the supported whitened
coordinates from Section 18.2. It is the filtered probe vector that can
be compared directly with the orthonormal QSE eigenvectors
\(y^{(\nu)}\). The intensity comes from squaring the scalar overlap with
an individual root.

Expanding \(g^O\) shows the full path through the retained metric modes.
With \[
W_\tau=U_\tau D_\tau^{-1/2},\qquad
U_\tau=
\begin{bmatrix}
u_{j_1}&\cdots&u_{j_r}
\end{bmatrix},
\qquad
D_\tau={\rm diag}(s_{j_1},\ldots,s_{j_r}),
\] we have \[
\begin{aligned}
t_\nu^O
&=(y^{(\nu)})^\dagger D_\tau^{-1/2}U_\tau^\dagger d^O\\
&=
\sum_{p=1}^r
(y_p^{(\nu)})^*s_{j_p}^{-1/2}
u_{j_p}^\dagger d^O\\
&=
\sum_{p=1}^r
\sum_{a\in\mathcal A}
(y_p^{(\nu)})^*s_{j_p}^{-1/2}
(u_{j_p})_a^*
\langle\phi_a|O|\psi_0\rangle .
\end{aligned}
\] The path is \[
O|\psi_0\rangle
\longrightarrow
d^O=\Phi_{\mathcal A}^\dagger O|\psi_0\rangle
\longrightarrow
g^O=W_\tau^\dagger d^O
\longrightarrow
t_\nu^O=(y^{(\nu)})^\dagger g^O .
\] The first arrow projects the probed seed into the raw excitation
bras. The second arrow applies the retained support and whitening. The
third arrow asks how much of that retained probe vector lies along root
\(\nu\). This reduced route avoids full exact diagonalization in the
original Hilbert space by diagonalizing the \(r\times r\) supported QSE
Hamiltonian; the bottleneck becomes selecting enough records, measuring
the projected matrices and transition vectors, and keeping the overlap
metric stable.

The transition intensity is the squared amplitude \[
I_\nu^O:=|t_\nu^O|^2 .
\] Qualitatively, \(I_\nu^O\) is the brightness or spectral weight of
root \(\nu\) for probe \(O\). It is large when \(O|\psi_0\rangle\)
points strongly along the QSE root, and small when the probe has little
component along that root. The root energy determines the line location;
the intensity determines how visible that line is.

Algebraically, \[
\begin{aligned}
I_\nu^O
&=
\left((c^{(\nu)})^\dagger d^O\right)^*
\left((c^{(\nu)})^\dagger d^O\right)\\
&=
(c^{(\nu)})^\dagger d^O(d^O)^\dagger c^{(\nu)} .
\end{aligned}
\] Define the rank-one probe-response matrix \[
R_O:=d^O(d^O)^\dagger,\qquad
(R_O)_{ab}=d_a^O(d_b^O)^* .
\] Then \[
I_\nu^O=(c^{(\nu)})^\dagger R_Oc^{(\nu)} .
\] In whitened coordinates, \[
\begin{aligned}
\widetilde R_O
&:=W_\tau^\dagger R_OW_\tau
=W_\tau^\dagger d^O(d^O)^\dagger W_\tau\\
&=(W_\tau^\dagger d^O)(W_\tau^\dagger d^O)^\dagger
=g^O(g^O)^\dagger,\\
I_\nu^O
&=(y^{(\nu)})^\dagger\widetilde R_Oy^{(\nu)} .
\end{aligned}
\] The response matrix is the probe's visibility operator inside the QSE
coordinate space. The Hamiltonian eigenproblem supplies the roots;
\(R_O\) or \(\widetilde R_O\) supplies which roots the chosen physical
operator can see.

The spectral line position is \[
\omega_\nu^{\rm QSE}:=\varepsilon_\nu-E_0,
\] with \(E_0:=\langle\psi_0|H_0|\psi_0\rangle\), unless the projected
Hamiltonian has already been shifted by \(E_0\), in which case
\(\varepsilon_\nu\) itself is the excitation frequency.

In the zero-temperature Lehmann
picture,\footnote{The Lehmann representation writes a response by inserting a complete energy eigenbasis between probe operators; at zero temperature this gives peaks or poles at energy differences $E_\nu-E_0$, with weights such as $|\langle\Psi_\nu|O|\psi_0\rangle|^2$.}
a scalar probe spectrum is a weighted collection of lines: \[
A_O(\omega)
=
\sum_\nu
|\langle\Psi_\nu|O|\psi_0\rangle|^2
\delta(\omega-(E_\nu-E_0)).
\] The QSE approximation replaces exact states, exact gaps, and exact
weights by QSE roots: \[
|\Psi_\nu\rangle\mapsto|\Psi_\nu^{\rm QSE}\rangle,\qquad
E_\nu-E_0\mapsto\omega_\nu^{\rm QSE},\qquad
|\langle\Psi_\nu|O|\psi_0\rangle|^2\mapsto I_\nu^O .
\] For plotted or finite-resolution spectra, the delta line is replaced
by a broadening kernel: \[
\delta(\omega-\omega_\nu^{\rm QSE})
\mapsto
L_\eta(\omega-\omega_\nu^{\rm QSE}).
\] The broadened QSE spectral function is \[
\begin{aligned}
A_O^{\mathcal A}(\omega)
&:=
\sum_{\nu\in\mathcal W_{\rm targ}}
I_\nu^O L_\eta(\omega-\omega_\nu^{\rm QSE}),\\
I_\nu^O&:=|t_\nu^O|^2 .
\end{aligned}
\] A spectral function is therefore a frequency-resolved visible
response: it places peaks at QSE excitation frequencies and scales each
peak by the probe-dependent intensity. The spectral window
\(\mathcal W_{\rm targ}\) is the set of retained roots or frequency
interval chosen for reporting. It determines which roots contribute to a
plotted spectrum, table metric, or integrated weight.

If the broadening kernel is normalized, \[
\int_{-\infty}^{\infty}L_\eta(x)\,dx=1,
\] then \[
\begin{aligned}
\int_{-\infty}^{\infty}A_O^{\mathcal A}(\omega)\,d\omega
&=
\sum_{\nu\in\mathcal W_{\rm targ}}I_\nu^O
\int_{-\infty}^{\infty}
L_\eta(\omega-\omega_\nu^{\rm QSE})\,d\omega\\
&=
\sum_{\nu\in\mathcal W_{\rm targ}}I_\nu^O .
\end{aligned}
\] The broadening controls visual resolution and comparison smoothness;
a normalized kernel preserves total reported spectral weight over the
chosen window. Two common kernels are \[
L_\eta^{\rm Lor}(x)=\frac1\pi\frac{\eta}{x^2+\eta^2},
\qquad
L_\eta^{\rm Gauss}(x)=\frac{1}{\sqrt{2\pi}\eta}
\exp\!\left(-\frac{x^2}{2\eta^2}\right).
\]

\begin{figure}
\centering
\pandocbounded{\includegraphics[keepaspectratio,alt={Schematic spectral semantics. Left: a probe-filtered scalar spectrum A\_O\^{}\{\textbackslash mathcal A\}(\textbackslash omega)=\textbackslash sum\_\textbackslash nu I\_\textbackslash nu\^{}O L\_\textbackslash eta(\textbackslash omega-\textbackslash omega\_\textbackslash nu\^{}\{\textbackslash rm QSE\}), where line positions are QSE excitation frequencies and line heights/areas encode probe-dependent intensities. Right: a single-particle spectrum A\_\{pp\}(\textbackslash omega)=-(1/\textbackslash pi)\{\textbackslash rm Im\}\textbackslash,G\_\{pp\}\^{}R(\textbackslash omega), with removal weight at negative frequency, addition weight at positive frequency, a spectral gap around the chosen zero, a coherent quasiparticle pole, and phonon sidebands generated by dressed electron-phonon propagation.}]{figures/qse_spectral_semantics.png}}
\caption{Schematic spectral semantics. Left: a probe-filtered scalar
spectrum
\(A_O^{\mathcal A}(\omega)=\sum_\nu I_\nu^O L_\eta(\omega-\omega_\nu^{\rm QSE})\),
where line positions are QSE excitation frequencies and line
heights/areas encode probe-dependent intensities. Right: a
single-particle spectrum
\(A_{pp}(\omega)=-(1/\pi){\rm Im}\,G_{pp}^R(\omega)\), with removal
weight at negative frequency, addition weight at positive frequency, a
spectral gap around the chosen zero, a coherent quasiparticle pole, and
phonon sidebands generated by dressed electron-phonon propagation.}
\end{figure}

Bright and dark roots are defined relative to a probe. A root with
\(I_\nu^O\approx0\) is dark for \(O\) and contributes little to
\(A_O^{\mathcal A}(\omega)\). A root with large \(I_\nu^O\) is bright
for \(O\) and can dominate the visible response above the lowest
excitation. A table of gaps tests root placement; a spectral function
tests whether the selected QSE records carry the physical response
channel.

The projection interpretation is compact in whitened coordinates. If the
retained eigenvectors \(y^{(\nu)}\) span the retained support, then \[
\sum_{\nu=1}^r y^{(\nu)}(y^{(\nu)})^\dagger=I_r .
\] With \(g^O=W_\tau^\dagger d^O\), \[
\begin{aligned}
\sum_{\nu=1}^r |t_\nu^O|^2
&=
\sum_{\nu=1}^r
(g^O)^\dagger y^{(\nu)}(y^{(\nu)})^\dagger g^O\\
&=
(g^O)^\dagger g^O
=\|W_\tau^\dagger d^O\|^2 .
\end{aligned}
\] Thus the full retained spectrum distributes the retained probe-vector
norm across the retained roots. A narrow \(\mathcal W_{\rm targ}\)
reports only the portion of that weight assigned to the chosen roots or
frequency band.

Finite-shot instability enters through the same objects. With \[
t_\nu^O=(y^{(\nu)})^\dagger W_\tau^\dagger d^O,
\] first-order perturbations give \[
\delta t_\nu^O
=
(\delta y^{(\nu)})^\dagger W_\tau^\dagger d^O
+(y^{(\nu)})^\dagger(\delta W_\tau^\dagger)d^O
+(y^{(\nu)})^\dagger W_\tau^\dagger\delta d^O .
\] The intensity variation is \[
\delta I_\nu^O
\approx
2\,{\rm Re}\!\left[(t_\nu^O)^*\delta t_\nu^O\right].
\] Root rotation, unstable whitening, and noisy transition-vector
estimates can therefore produce unstable spectral intensities even when
the energy error looks moderate.

The compact summary is \[
\begin{aligned}
|\Psi_\nu^{\rm QSE}\rangle
&=\Phi_{\mathcal A}c^{(\nu)},\\
d^O
&=\Phi_{\mathcal A}^\dagger O|\psi_0\rangle,\\
g^O
&=W_\tau^\dagger d^O,\\
t_\nu^O
&=(c^{(\nu)})^\dagger d^O
=(y^{(\nu)})^\dagger g^O,\\
I_\nu^O
&=|t_\nu^O|^2,\\
A_O^{\mathcal A}(\omega)
&=
\sum_{\nu\in\mathcal W_{\rm targ}}
I_\nu^O L_\eta(\omega-\omega_\nu^{\rm QSE}) .
\end{aligned}
\] The eigenproblem tells where the possible lines are. The transition
vector tells which lines the physical probe can see. The spectral
function is the visible response: root locations weighted by
probe-dependent brightness.

\subsection{18.4 Correlation functions, response matrices, and spectral
moments}\label{correlation-functions-response-matrices-and-spectral-moments}

Section 18.3 used one probe \(O\) and asked which QSE roots are bright
under that probe. Section 18.4 uses two probes, \(A\) and \(B\). The
operator \(B\) creates a disturbance from the seed, the Hamiltonian
evolves that disturbed component, and \(A\) reads out how much of the
evolved disturbance appears in another physical channel. A response
function is the time- or frequency-dependent object that records this
create-evolve-read process. In frequency space, a sharp response appears
as a pole: a term such as \(R/(\omega-\omega_\nu+i\eta)\), whose
denominator becomes small near the excitation frequency \(\omega_\nu\).
The pole location gives the excitation energy, and the residue \(R\)
gives the channel weight carried by that excitation.

Interpretive compression: \(C_{AB}(t)\), \(S_{AB}(\omega)\), and
\(m_k^{AB}\) are three representations of the same bilinear spectral
measure. The time correlator tracks phase evolution of residues, the
frequency response resolves the same residues into broadened poles, and
the moments compress the distribution into weighted polynomial sums of
excitation frequencies. The response matrix upgrades scalar brightness
into channel covariance: diagonal entries report self-channel spectral
weight, off-diagonal entries report coherent mixing between distinct
physical probes.

Physical semantics: the pair \((A,B)\) defines an input-output channel.
The source operator \(B\) selects a disturbed seed vector
\(B|\psi_0\rangle\); the readout operator \(A\) selects the dual
functional \(\langle\psi_0|A^\dagger\) after time evolution. The object
\(C_{AB}(t)\) is therefore a causal channel correlator before Fourier
analysis, and \(\mathsf S_{AB}(\omega)\) is its frequency-resolved
transfer measure. Diagonal entries \(\mathsf S_{AA}\) quantify the
spectral weight returned to the same channel; off-diagonal entries
\(\mathsf S_{AB}\) quantify coherent conversion between channels,
including relative phase information in complex residues
\((t_\nu^A)^*t_\nu^B\).

\subsubsection{§18.4 lesson: correlation functions, response matrices,
and spectral
moments}\label{lesson-correlation-functions-response-matrices-and-spectral-moments}

The Heisenberg and Schrödinger pictures put the same dynamics in
different places. In the Schrödinger picture, states carry the time
dependence. In the Heisenberg picture, the reference state is fixed and
operators carry the time dependence: \[
A(t):=e^{iH_0t}Ae^{-iH_0t}.
\] The correlation convention used here places the disturbance at time
zero and the readout at time \(t\): \[
B(0):=B,\qquad
C_{AB}(t):=\langle\psi_0|A^\dagger(t)B(0)|\psi_0\rangle .
\] For a stationary seed of \(H_0\), only the time difference matters.
Writing \(B(0)\) makes the convention explicit: \(B\) prepares the
disturbed vector at the initial time, and \(A^\dagger(t)\) tests that
disturbed vector after Hamiltonian evolution.

Assume first an exact finite Hilbert-space problem, \[
H_0|\psi_0\rangle=E_0|\psi_0\rangle,\qquad
H_0|\Psi_\nu\rangle=E_\nu|\Psi_\nu\rangle .
\] Then \[
\begin{aligned}
C_{AB}(t)
&=
\langle\psi_0|
e^{iH_0t}A^\dagger e^{-iH_0t}B
|\psi_0\rangle .
\end{aligned}
\] Insert the exact resolution of identity between the readout operator
and the evolved disturbance: \[
I=\sum_\nu|\Psi_\nu\rangle\langle\Psi_\nu|.
\] Then \[
\begin{aligned}
C_{AB}(t)
&=
\langle\psi_0|
e^{iH_0t}A^\dagger
\left(\sum_\nu|\Psi_\nu\rangle\langle\Psi_\nu|\right)
e^{-iH_0t}B
|\psi_0\rangle\\
&=
\sum_\nu
\langle\psi_0|e^{iH_0t}A^\dagger|\Psi_\nu\rangle
\langle\Psi_\nu|e^{-iH_0t}B|\psi_0\rangle .
\end{aligned}
\] Use \[
\langle\psi_0|e^{iH_0t}=e^{iE_0t}\langle\psi_0|,
\qquad
\langle\Psi_\nu|e^{-iH_0t}=e^{-iE_\nu t}\langle\Psi_\nu|.
\] This gives the Lehmann expansion \[
\begin{aligned}
C_{AB}(t)
&=
\sum_\nu
e^{iE_0t}
\langle\psi_0|A^\dagger|\Psi_\nu\rangle
e^{-iE_\nu t}
\langle\Psi_\nu|B|\psi_0\rangle\\
&=
\sum_\nu
\langle\psi_0|A^\dagger|\Psi_\nu\rangle
\langle\Psi_\nu|B|\psi_0\rangle
e^{-i(E_\nu-E_0)t}.
\end{aligned}
\] Each term has three pieces: \(B|\psi_0\rangle\) supplies the
component of the seed disturbance along eigenmode \(\nu\),
\(e^{-i(E_\nu-E_0)t}\) rotates that component at its excitation
frequency, and \(\langle\psi_0|A^\dagger\) reads out whether that
evolved component is visible to channel \(A\).

Use the transition-amplitude convention from Section 18.3: \[
t_\nu^O:=\langle\Psi_\nu|O|\psi_0\rangle .
\] Then \[
t_\nu^A=\langle\Psi_\nu|A|\psi_0\rangle,\qquad
t_\nu^B=\langle\Psi_\nu|B|\psi_0\rangle,
\] and \[
\begin{aligned}
(t_\nu^A)^*
&=
\left(\langle\Psi_\nu|A|\psi_0\rangle\right)^*\\
&=
\langle\psi_0|A^\dagger|\Psi_\nu\rangle .
\end{aligned}
\] Thus \[
\begin{aligned}
C_{AB}(t)
&=
\sum_\nu
(t_\nu^A)^*t_\nu^B e^{-i\omega_\nu t},
\qquad
\omega_\nu:=E_\nu-E_0 .
\end{aligned}
\] The phrase ``amplitude from \(B\), read out by \(A\)'' means exactly
this product: \(t_\nu^B\) is how strongly the perturbing vector
\(B|\psi_0\rangle\) enters eigenmode \(\nu\), while \((t_\nu^A)^*\) is
how strongly the same eigenmode is visible to the \(A\)-readout bra.

Now replace the exact eigenstates and exact gaps by QSE roots: \[
|\Psi_\nu\rangle\mapsto|\Psi_\nu^{\rm QSE}\rangle,\qquad
\omega_\nu\mapsto\omega_\nu^{\rm QSE}.
\] The QSE root is synthesized as \[
|\Psi_\nu^{\rm QSE}\rangle
=
\Phi_{\mathcal A}c^{(\nu)}
=
\Phi_{\mathcal A}W_\tau y^{(\nu)}
=
\sum_{a\in\mathcal A}c_a^{(\nu)}|\phi_a\rangle .
\] For any probe \(O\), define the raw and whitened probe vectors \[
d^O:=\Phi_{\mathcal A}^\dagger O|\psi_0\rangle,
\qquad
g^O:=W_\tau^\dagger d^O .
\] Then \[
\begin{aligned}
t_\nu^O
&=
\langle\Psi_\nu^{\rm QSE}|O|\psi_0\rangle\\
&=
(c^{(\nu)})^\dagger d^O\\
&=
(y^{(\nu)})^\dagger W_\tau^\dagger d^O\\
&=
(y^{(\nu)})^\dagger g^O .
\end{aligned}
\] Therefore the QSE correlation reconstruction is \[
C_{AB}^{\mathcal A}(t)
:=
\sum_{\nu\in\mathcal W_{\rm targ}}
(t_\nu^A)^*t_\nu^B
e^{-i\omega_\nu^{\rm QSE}t}.
\]

The cross-channel residue of root \(\nu\) is \[
R_\nu^{AB,\mathcal A}:=(t_\nu^A)^*t_\nu^B .
\] It is rank one in channel space because, for one fixed root, all
channel-pair weights are built from one vector of transition amplitudes.
In raw QSE coordinates, \[
\begin{aligned}
R_\nu^{AB,\mathcal A}
&=
\left((c^{(\nu)})^\dagger d^A\right)^*
\left((c^{(\nu)})^\dagger d^B\right)\\
&=
(d^A)^\dagger c^{(\nu)}(c^{(\nu)})^\dagger d^B .
\end{aligned}
\] In supported whitened coordinates, \[
\begin{aligned}
R_\nu^{AB,\mathcal A}
&=
\left((y^{(\nu)})^\dagger g^A\right)^*
\left((y^{(\nu)})^\dagger g^B\right)\\
&=
(g^A)^\dagger y^{(\nu)}(y^{(\nu)})^\dagger g^B .
\end{aligned}
\] In raw excitation-basis indices, \[
\begin{aligned}
R_\nu^{AB,\mathcal A}
&=
\sum_{a,b\in\mathcal A}
c_a^{(\nu)}
(c_b^{(\nu)})^*
(d_a^A)^*d_b^B\\
&=
\sum_{a,b\in\mathcal A}
c_a^{(\nu)}
(c_b^{(\nu)})^*
\langle\psi_0|A^\dagger|\phi_a\rangle
\langle\phi_b|B|\psi_0\rangle .
\end{aligned}
\] The \(A\)-brightness and \(B\)-brightness of a root are the
magnitudes and phases of \(t_\nu^A\) and \(t_\nu^B\). A diagonal residue
\((t_\nu^O)^*t_\nu^O=|t_\nu^O|^2\) is a scalar brightness. An
off-diagonal residue \((t_\nu^A)^*t_\nu^B\) records how the same root
connects two different physical channels.

Frequency space is obtained by Fourier transforming the time signal, or
by directly replacing each oscillatory mode with its frequency-domain
line. Formally, \[
\int_{-\infty}^{\infty} e^{i\omega t}e^{-i\omega_\nu t}\,dt
=
2\pi\,\delta(\omega-\omega_\nu),
\] so a pure oscillation becomes a sharp line at \(\omega=\omega_\nu\).
For finite-resolution plotting and finite-shot comparison, the sharp
distribution is replaced by a normalized kernel \(L_\eta\). The QSE
response entry is \[
\mathsf S_{AB}^{\mathcal A}(\omega)
:=
\sum_{\nu\in\mathcal W_{\rm targ}}
R_\nu^{AB,\mathcal A}
L_\eta(\omega-\omega_\nu^{\rm QSE})
=
\sum_{\nu\in\mathcal W_{\rm targ}}
(t_\nu^A)^*t_\nu^B
L_\eta(\omega-\omega_\nu^{\rm QSE}) .
\] For \(A=B=O\), \[
\begin{aligned}
\mathsf S_{OO}^{\mathcal A}(\omega)
&=
\sum_{\nu\in\mathcal W_{\rm targ}}
|t_\nu^O|^2
L_\eta(\omega-\omega_\nu^{\rm QSE})\\
&=
A_O^{\mathcal A}(\omega).
\end{aligned}
\] Thus the scalar spectral function from Section 18.3 is the diagonal
same-probe response entry.

The matrix-valued response keeps all probe channels together. Choose \[
\mathcal O_{\rm resp}=\{O_1,\ldots,O_P\},
\qquad
T_{\nu\alpha}:=t_\nu^{O_\alpha}.
\] For a fixed frequency \(\omega\), define the diagonal line-shape
matrix \[
D_\omega:=
{\rm diag}\!\left(
L_\eta(\omega-\omega_\nu^{\rm QSE})
\right)_{\nu\in\mathcal W_{\rm targ}} .
\] Then \[
\mathsf S^{\mathcal A}(\omega):=T^\dagger D_\omega T
\] has entries \[
\begin{aligned}
\left(\mathsf S^{\mathcal A}(\omega)\right)_{\alpha\beta}
&=
\sum_{\mu,\nu}
(T^\dagger)_{\alpha\nu}
(D_\omega)_{\nu\mu}
T_{\mu\beta}\\
&=
\sum_{\mu,\nu}
T_{\nu\alpha}^*
\delta_{\nu\mu}
L_\eta(\omega-\omega_\nu^{\rm QSE})
T_{\mu\beta}\\
&=
\sum_\nu
(t_\nu^{O_\alpha})^*t_\nu^{O_\beta}
L_\eta(\omega-\omega_\nu^{\rm QSE})\\
&=
\mathsf S_{O_\alpha O_\beta}^{\mathcal A}(\omega).
\end{aligned}
\] This response matrix has indices in probe-channel space. Its rows and
columns are density, displacement, current, particle-addition,
particle-removal, or mixed probes. The overlap metric
\(S_{ab}=\langle\phi_a|\phi_b\rangle\) from Section 18.2 has indices in
excitation-basis space. The index sets identify the object.

If \(L_\eta(x)\) is real, then \[
\begin{aligned}
\left(\mathsf S_{AB}^{\mathcal A}(\omega)\right)^*
&=
\left[
\sum_\nu
(t_\nu^A)^*t_\nu^B
L_\eta(\omega-\omega_\nu^{\rm QSE})
\right]^*\\
&=
\sum_\nu
(t_\nu^B)^*t_\nu^A
L_\eta(\omega-\omega_\nu^{\rm QSE})\\
&=
\mathsf S_{BA}^{\mathcal A}(\omega).
\end{aligned}
\] If \(L_\eta(x)\ge0\), then for any channel vector
\(z\in\mathbb C^P\), \[
\begin{aligned}
z^\dagger\mathsf S^{\mathcal A}(\omega)z
&=
z^\dagger T^\dagger D_\omega Tz\\
&=
(Tz)^\dagger D_\omega(Tz)\\
&=
\sum_\nu
|(Tz)_\nu|^2
L_\eta(\omega-\omega_\nu^{\rm QSE})\\
&\ge0 .
\end{aligned}
\] Thus the channel-response matrix is Hermitian for real broadening and
positive semidefinite for nonnegative broadening. These are
implementation sanity checks.

Spectral moments summarize a distribution by weighted sums over
frequency powers. The zeroth moment is total channel weight in the
reported window. The first moment is the frequency-weighted channel
weight. Higher moments emphasize higher-frequency parts of the response.
For the QSE response, \[
\begin{aligned}
m_k^{AB,\mathcal A}
&:=
\sum_{\nu\in\mathcal W_{\rm targ}}
(\omega_\nu^{\rm QSE})^k
R_\nu^{AB,\mathcal A}\\
&=
\sum_{\nu\in\mathcal W_{\rm targ}}
(\omega_\nu^{\rm QSE})^k
(t_\nu^A)^*t_\nu^B\\
&=
\sum_{\nu\in\mathcal W_{\rm targ}}
(\omega_\nu^{\rm QSE})^k
(d^A)^\dagger c^{(\nu)}(c^{(\nu)})^\dagger d^B\\
&=
\sum_{\nu\in\mathcal W_{\rm targ}}
(\omega_\nu^{\rm QSE})^k
(g^A)^\dagger y^{(\nu)}(y^{(\nu)})^\dagger g^B .
\end{aligned}
\] The first two QSE moments are \[
\begin{aligned}
m_0^{AB,\mathcal A}
&=
\sum_{\nu\in\mathcal W_{\rm targ}}
(t_\nu^A)^*t_\nu^B,\\
m_1^{AB,\mathcal A}
&=
\sum_{\nu\in\mathcal W_{\rm targ}}
\omega_\nu^{\rm QSE}
(t_\nu^A)^*t_\nu^B .
\end{aligned}
\] They are the same residues used in the response function, with
different powers of frequency attached.

Moments are also short-time Taylor data for the correlation function.
Starting from \[
C_{AB}^{\mathcal A}(t)
=
\sum_\nu
R_\nu^{AB,\mathcal A}e^{-i\omega_\nu^{\rm QSE}t},
\] differentiate \(k\) times: \[
\begin{aligned}
\frac{d^k}{dt^k}C_{AB}^{\mathcal A}(t)
&=
\sum_\nu
R_\nu^{AB,\mathcal A}
(-i\omega_\nu^{\rm QSE})^k
e^{-i\omega_\nu^{\rm QSE}t}.
\end{aligned}
\] At \(t=0\), \[
\left.\frac{d^k}{dt^k}C_{AB}^{\mathcal A}(t)\right|_{t=0}
=
(-i)^k m_k^{AB,\mathcal A},
\qquad
m_k^{AB,\mathcal A}
=
i^k
\left.\frac{d^k}{dt^k}C_{AB}^{\mathcal A}(t)\right|_{t=0}.
\] So moment errors are short-time correlation errors written in
frequency language.

The exact same-cutoff sum rules give reference values for these moments.
``Same cutoff'' means the exact comparison is evaluated in the same
finite physical representation as the QSE calculation: the same particle
sector, symmetry sector, boson cutoff, lattice size, operator
convention, and Hamiltonian \(H_0\). This assigns the residual to
QSE-subspace incompleteness inside a fixed finite representation.

In that same finite representation, a complete exact eigenbasis gives \[
m_k^{AB}
=
\sum_\nu
(E_\nu-E_0)^k
\langle\psi_0|A^\dagger|\Psi_\nu\rangle
\langle\Psi_\nu|B|\psi_0\rangle .
\] For \(k=0\), \[
\begin{aligned}
m_0^{AB}
&=
\sum_\nu
\langle\psi_0|A^\dagger|\Psi_\nu\rangle
\langle\Psi_\nu|B|\psi_0\rangle\\
&=
\langle\psi_0|A^\dagger
\left(\sum_\nu|\Psi_\nu\rangle\langle\Psi_\nu|\right)
B|\psi_0\rangle\\
&=
\langle\psi_0|A^\dagger B|\psi_0\rangle .
\end{aligned}
\] For \(k=1\), use \[
(E_\nu-E_0)\langle\Psi_\nu|B|\psi_0\rangle
=
\langle\Psi_\nu|(H_0-E_0)B|\psi_0\rangle .
\] Then \[
\begin{aligned}
m_1^{AB}
&=
\sum_\nu
\langle\psi_0|A^\dagger|\Psi_\nu\rangle
\langle\Psi_\nu|(H_0-E_0)B|\psi_0\rangle\\
&=
\langle\psi_0|A^\dagger
\left(\sum_\nu|\Psi_\nu\rangle\langle\Psi_\nu|\right)
(H_0-E_0)B|\psi_0\rangle\\
&=
\langle\psi_0|A^\dagger(H_0-E_0)B|\psi_0\rangle .
\end{aligned}
\] The general same-cutoff identity is \[
m_k^{AB}
=
\langle\psi_0|A^\dagger(H_0-E_0)^kB|\psi_0\rangle .
\] The selected-QSE residual is \[
\Delta m_k^{AB,\mathcal A}
:=
m_k^{AB}-m_k^{AB,\mathcal A}.
\] For the two lowest moments, \[
\begin{aligned}
\Delta m_0^{AB,\mathcal A}
&=
\langle\psi_0|A^\dagger B|\psi_0\rangle
-
\sum_{\nu\in\mathcal W_{\rm targ}}
(t_\nu^A)^*t_\nu^B,\\
\Delta m_1^{AB,\mathcal A}
&=
\langle\psi_0|A^\dagger(H_0-E_0)B|\psi_0\rangle
-
\sum_{\nu\in\mathcal W_{\rm targ}}
\omega_\nu^{\rm QSE}
(t_\nu^A)^*t_\nu^B .
\end{aligned}
\] These residuals audit spectral closure. A few accurate visible peak
locations can coexist with a large moment residual when the selected QSE
basis misses response-relevant roots or distributes the channel weight
incorrectly.

The connected correlator removes the elastic disconnected background. If
the seed already has nonzero expectations \[
\langle A^\dagger\rangle_0:=\langle\psi_0|A^\dagger|\psi_0\rangle,
\qquad
\langle B\rangle_0:=\langle\psi_0|B|\psi_0\rangle,
\] then the ground-state term in the Lehmann sum has \(\omega_0=0\) and
contributes \[
\langle A^\dagger\rangle_0\langle B\rangle_0 e^{-i0t}
=
\langle A^\dagger\rangle_0\langle B\rangle_0 .
\] This is a static zero-frequency contribution: it records the
background expectation already present in the seed. The connected
correlator subtracts this piece: \[
C_{AB}^{\rm conn}(t)
:=
C_{AB}(t)
-
\langle\psi_0|A^\dagger|\psi_0\rangle
\langle\psi_0|B|\psi_0\rangle .
\] Equivalently, it studies the fluctuation operators \[
\delta A:=A-\langle A\rangle_0,\qquad
\delta B:=B-\langle B\rangle_0,
\] so the response focuses on dynamical fluctuations around the seed
background. The raw correlator reports static expectation plus
fluctuations; the connected correlator reports created excitations with
the elastic background line removed.

The retarded linear-response susceptibility is the causal version of the
response for Hermitian bosonic or even observable channels such as
density, displacement, and current. It uses a commutator and the step
function: \[
\chi_{AB}^{R}(t)
:=
-i\,\Theta(t)
\langle\psi_0|[A(t),B(0)]|\psi_0\rangle ,
\qquad
[A(t),B(0)]:=A(t)B(0)-B(0)A(t).
\] The two ordered pieces have Lehmann forms \[
\begin{aligned}
\langle\psi_0|A(t)B|\psi_0\rangle
&=
\sum_\nu
(t_\nu^A)^*t_\nu^B e^{-i\omega_\nu t},\\
\langle\psi_0|BA(t)|\psi_0\rangle
&=
\sum_\nu
(t_\nu^B)^*t_\nu^A e^{+i\omega_\nu t}.
\end{aligned}
\] Thus \[
\chi_{AB}^{R}(t)
=
-i\,\Theta(t)
\sum_\nu
\left[
(t_\nu^A)^*t_\nu^B e^{-i\omega_\nu t}
-
(t_\nu^B)^*t_\nu^A e^{+i\omega_\nu t}
\right].
\] Using \[
\chi_{AB}^{R}(\omega)
:=
\int_{-\infty}^{\infty}dt\,e^{i\omega t}\chi_{AB}^{R}(t)
\] with convergence factor \(e^{-\eta t}\) for \(t>0\), the first term
gives \[
\begin{aligned}
-i\int_0^\infty dt\,
e^{i\omega t}e^{-i\omega_\nu t}e^{-\eta t}
&=
-i\int_0^\infty dt\,
e^{-[\eta-i(\omega-\omega_\nu)]t}\\
&=
\frac{1}{\omega-\omega_\nu+i\eta},
\end{aligned}
\] and the second term gives \[
\begin{aligned}
i\int_0^\infty dt\,
e^{i\omega t}e^{+i\omega_\nu t}e^{-\eta t}
&=
i\int_0^\infty dt\,
e^{-[\eta-i(\omega+\omega_\nu)]t}\\
&=
-
\frac{1}{\omega+\omega_\nu+i\eta}.
\end{aligned}
\] Therefore \[
\chi_{AB}^{R}(\omega)
=
\sum_\nu
\left[
\frac{(t_\nu^A)^*t_\nu^B}{\omega-\omega_\nu+i\eta}
-
\frac{(t_\nu^B)^*t_\nu^A}{\omega+\omega_\nu+i\eta}
\right].
\] This is the pole expansion for the Hermitian response channels used
in this subsection. In the ideal limit \(\eta\to0^+\), the denominators
identify resonances at \(\omega=\omega_\nu\) and \(\omega=-\omega_\nu\).
In the QSE estimate, replace \(\omega_\nu,t_\nu^A,t_\nu^B\) by
\(\omega_\nu^{\rm QSE},t_{\nu,\mathcal A}^A,t_{\nu,\mathcal A}^B\)
computed from the selected supported QSE basis. Non-Hermitian
particle-addition and particle-removal probes use the paired adjoint
sectors described in Section 18.6.

The compact summary is \[
\begin{aligned}
B|\psi_0\rangle
&\text{ creates the disturbed seed component},\\
e^{-iH_0t}
&\text{ evolves that component},\\
A
&\text{ reads out the evolved component},\\
C_{AB}^{\mathcal A}(t)
&=
\sum_{\nu\in\mathcal W_{\rm targ}}
(t_\nu^A)^*t_\nu^B
e^{-i\omega_\nu^{\rm QSE}t},\\
\mathsf S_{AB}^{\mathcal A}(\omega)
&=
\sum_{\nu\in\mathcal W_{\rm targ}}
(t_\nu^A)^*t_\nu^B
L_\eta(\omega-\omega_\nu^{\rm QSE}),\\
m_k^{AB,\mathcal A}
&=
\sum_{\nu\in\mathcal W_{\rm targ}}
(\omega_\nu^{\rm QSE})^k
(t_\nu^A)^*t_\nu^B .
\end{aligned}
\] Section 18.3 is one probe seeing bright and dark roots. Section 18.4
is pairs of probes revealing how response channels talk to each other.
The moment identities are the audit: if the selected QSE basis captures
the response channel, its summed residues reproduce the same-cutoff
low-order sum rules; if it misses response-relevant states, the
residuals quantify the missing spectral weight.

\subsection{18.5 Hubbard-Holstein response
channels}\label{hubbard-holstein-response-channels}

In the Hubbard-Holstein model, an excitation can carry charge motion,
oscillator motion, or both at the same time. The reason is visible at
the Hamiltonian level: \[
H_{\rm HH}
=
H_{\rm Hub}
+\omega_0\sum_i\left(n_{b,i}+\frac12 I\right)
+g\sum_iX_i(n_i-\bar n I)
+\cdots,
\qquad
X_i:=b_i^\dagger+b_i .
\] The interaction term is density fluctuation times displacement. A QSE
basis that captures Hubbard-Holstein physics should therefore be
interrogated by several probes: density \(n\), phonon displacement
\(X\), phonon momentum \(P\), and mixed density-displacement dressing
\(nX\). The eigenvalues say where the roots are; the response channels
say what those roots are made of.

Interpretive compression: Hubbard-Holstein response is a
channel-resolved decomposition of polaronic spectral weight. The
eigenmode labels are secondary; the physically invariant diagnostic is
how each Ritz root distributes residue across electronic density,
oscillator quadrature, and mixed density-displacement probes.
Open-boundary form factors replace momentum quantum numbers with spatial
analysis vectors, while periodic Fourier modes are the
translation-symmetric specialization.

General-to-specific semantics: a response channel is a selected
observable subalgebra through which the many-body state is interrogated.
In a pure electronic model, density and current channels already
separate charge redistribution from charge transport. In a coupled
electron-phonon model, oscillator quadratures add independent lattice
channels, and mixed density-displacement objects test correlated
electronic-lattice motion. For Hubbard-Holstein, the interaction
\(gX_i(n_i-\bar nI)\) makes the mixed residue structure central: a root
with simultaneous \(n\)-weight and \(X\)-weight is a dressed
charge-lattice excitation, while the cutoff diagnostic
\(\ell_{\rm cut}\) checks whether the finite oscillator register
supports that interpretation.

\subsubsection{§18.5 lesson: Hubbard-Holstein response
channels}\label{lesson-hubbard-holstein-response-channels}

The base convention in this document is an open chain. Let sites be
\(j=0,\ldots,L-1\), with positions \(r_j=ja\). Open boundaries break
translation symmetry, so a label called \(k\) is a probe-shape label.
Conserved lattice momentum quantum numbers require translation symmetry.
The clean open-boundary notation is a form-factor probe: \[
f:\{0,\ldots,L-1\}\to\mathbb C .
\] It defines a spatial pattern over the chain. A local probe uses
\(f(j)=\delta_{j\ell}\). A smooth open-chain standing-wave probe can use
\[
q_m:=\frac{\pi m}{(L+1)a},
\qquad
f_m^{\rm OBC}(j):=\sqrt{\frac{2}{L+1}}\sin(q_m(j+1)a),
\qquad
m=1,\ldots,L .
\] These \(q_m\) are wavevector-like labels for standing waves on an
open interval. They diagnose spatial structure at wavelength scale
\(2\pi/q_m\), while the boundaries prevent them from being exact
conserved
momenta.\footnote{For a periodic chain, translation symmetry makes the plane wave $f_k^{\rm PBC}(j)=L^{-1/2}e^{-ikr_j}$ natural, with $k=2\pi m/(La)$, $m=0,\ldots,L-1$. Then $k$ is a reciprocal-lattice momentum modulo $2\pi/a$. The formulas below reduce to the familiar $n_k,X_k,P_k$ notation by taking $f=f_k^{\rm PBC}$. Omitting $L^{-1/2}$ is a normalization convention that rescales intensities by $L$.}

Start with the local electronic density. For site \(i\) and spin
\(\sigma\), \[
n_{i\sigma}:=c_{i\sigma}^\dagger c_{i\sigma},
\qquad
n_i:=\sum_\sigma c_{i\sigma}^\dagger c_{i\sigma}.
\] The seed-averaged density is \[
\bar n:=\frac1L\sum_i\langle\psi_0|n_i|\psi_0\rangle,
\] and the density fluctuation is \[
\delta n_i:=n_i-\bar n I .
\] The subtraction removes the static density background: \[
\begin{aligned}
\frac1L\sum_i\langle\psi_0|\delta n_i|\psi_0\rangle
&=
\frac1L\sum_i
\langle\psi_0|(n_i-\bar n I)|\psi_0\rangle\\
&=
\frac1L\sum_i\langle\psi_0|n_i|\psi_0\rangle
-
\frac1L\sum_i\bar n\\
&=
\bar n-\bar n
=0,
\end{aligned}
\] assuming \(\langle\psi_0|\psi_0\rangle=1\). In a fixed
particle-number sector with \(\bar n=N_e/L\), the uniform density
fluctuation satisfies \[
\sum_i\delta n_i
=
\sum_i n_i-L\bar n I
=
N_e I-N_e I
=0
\] inside that sector. Thus neutral density response is carried by
nonuniform spatial patterns.

The local phonon quadratures are \[
X_i:=b_i^\dagger+b_i,
\qquad
P_i:=i(b_i^\dagger-b_i).
\] Here \(X_i\) is displacement-like and \(P_i\) is momentum-like. With
\[
[b_i,b_j^\dagger]=\delta_{ij},
\qquad
[b_i,b_j]=[b_i^\dagger,b_j^\dagger]=0,
\] their commutator is \[
\begin{aligned}
\left[X_i,P_j\right]
&=
[b_i^\dagger+b_i,\;i(b_j^\dagger-b_j)]\\
&=
i\left(
[b_i^\dagger,b_j^\dagger]
-[b_i^\dagger,b_j]
+[b_i,b_j^\dagger]
-[b_i,b_j]
\right)\\
&=
i\left(0-(-\delta_{ij})+\delta_{ij}-0\right)\\
&=
2i\delta_{ij}.
\end{aligned}
\] The factor \(2\) comes from the convention \(X=b^\dagger+b\) and
\(P=i(b^\dagger-b)\), relative to the \(1/\sqrt2\)-normalized oscillator
convention.

For any spatial form factor \(f\), define the open-chain response probes
\[
\begin{aligned}
n[f]&:=\sum_j f(j)\delta n_j,\\
X[f]&:=\sum_j f(j)X_j,\\
P[f]&:=\sum_j f(j)P_j .
\end{aligned}
\] If \(f\) is real, these are Hermitian operators. If \(f\) is complex,
then \[
n[f]^\dagger=n[f^*],
\qquad
X[f]^\dagger=X[f^*],
\qquad
P[f]^\dagger=P[f^*].
\] The periodic shorthand is recovered by \(f(j)=e^{-ikr_j}\): \[
n_k:=\sum_j e^{-ikr_j}\delta n_j,
\qquad
X_k:=\sum_j e^{-ikr_j}X_j,
\qquad
P_k:=\sum_j e^{-ikr_j}P_j,
\] with \[
n_k^\dagger=n_{-k},
\qquad
X_k^\dagger=X_{-k},
\qquad
P_k^\dagger=P_{-k}.
\] Under open boundaries, the same plane-wave formula is still a valid
chosen probe shape, but the standing-wave or general \(f\)-probe
notation is the base interpretation.

The local mixed dressing observable is \[
C_{nX}(r)
:=
\sum_{i:\,0\le i,\,i+r<L}\delta n_iX_{i+r}.
\] It is a product operator. It asks whether a density fluctuation near
site \(i\) is accompanied by a lattice displacement \(r\) sites away.
This is the local polaronic question: does the charge fluctuation drag a
phonon cloud? Since density acts on the electronic factor and \(X\) acts
on the bosonic factor, \[
[\delta n_i,X_j]=0,
\] so the product has no fermion-boson ordering ambiguity.

For a periodic chain, the relative-displacement Fourier transform of
this composite operator gives the familiar product of momentum probes.
With periodic wrapping, \[
C_{nX}(r)=\sum_i\delta n_iX_{i+r},
\qquad
C_{nX}(k):=\sum_r e^{-ikr_r}C_{nX}(r).
\] Set \(\ell=i+r\), so \(r_r=r_\ell-r_i\). Then \[
\begin{aligned}
C_{nX}(k)
&=
\sum_{i,\ell}
e^{-ik(r_\ell-r_i)}
\delta n_iX_\ell\\
&=
\left(\sum_i e^{+ikr_i}\delta n_i\right)
\left(\sum_\ell e^{-ikr_\ell}X_\ell\right)\\
&=
n_{-k}X_k,
\end{aligned}
\] up to the chosen Fourier normalization. On an open chain, the
valid-pair restriction in \(C_{nX}(r)\) is the base definition, and a
form-factor version can be written as \[
C_{nX}[h]:=\sum_{i,j}h(i,j)\delta n_iX_j .
\] The cross-channel response \(\mathsf S_{nX}\) and the composite
operator \(C_{nX}\) are related diagnostics. The former uses density as
one probe and displacement as another; the latter places density and
displacement in the same operator.

Now connect these probes to the QSE response machinery from Sections
18.3 and 18.4. For any probe \(O\), \[
d_a^O:=\langle\phi_a|O|\psi_0\rangle,
\qquad
t_\nu^O
:=
\langle\Psi_\nu^{\rm QSE}|O|\psi_0\rangle
=
(c^{(\nu)})^\dagger d^O .
\] With \(c^{(\nu)}=W_\tau y^{(\nu)}\), \[
t_\nu^O
=
(y^{(\nu)})^\dagger W_\tau^\dagger d^O .
\] For the density form-factor probe, \[
\begin{aligned}
d_a^{n[f]}
&=
\langle\phi_a|n[f]|\psi_0\rangle\\
&=
\sum_j f(j)\langle\phi_a|\delta n_j|\psi_0\rangle .
\end{aligned}
\] Similarly, \[
d_a^{X[f]}=\sum_j f(j)\langle\phi_a|X_j|\psi_0\rangle,
\qquad
d_a^{P[f]}=\sum_j f(j)\langle\phi_a|P_j|\psi_0\rangle .
\] Thus \[
\begin{aligned}
t_\nu^{n[f]}
&=
\sum_{a,j}
(c_a^{(\nu)})^*f(j)
\langle\phi_a|\delta n_j|\psi_0\rangle,\\
t_\nu^{X[f]}
&=
\sum_{a,j}
(c_a^{(\nu)})^*f(j)
\langle\phi_a|X_j|\psi_0\rangle,\\
t_\nu^{P[f]}
&=
\sum_{a,j}
(c_a^{(\nu)})^*f(j)
\langle\phi_a|P_j|\psi_0\rangle .
\end{aligned}
\] The QSE coefficients select a root, the local matrix elements
determine channel visibility, and the form factor selects the spatial
pattern being interrogated.

Using the response matrix from Section 18.4, \[
\mathsf S_{AB}^{\mathcal A}(\omega)
=
\sum_{\nu\in\mathcal W_{\rm targ}}
(t_\nu^A)^*t_\nu^B
L_\eta(\omega-\omega_\nu^{\rm QSE}),
\] the main neutral Hubbard-Holstein response targets on an open chain
are \[
\begin{aligned}
\mathsf S_{nn}^{\mathcal A}[f,g](\omega)
&:=
\mathsf S_{\,n[f]\,n[g]}^{\mathcal A}(\omega)
=
\sum_{\nu\in\mathcal W_{\rm targ}}
(t_\nu^{n[f]})^*t_\nu^{n[g]}
L_\eta(\omega-\omega_\nu^{\rm QSE}),\\
\mathsf S_{XX}^{\mathcal A}[f,g](\omega)
&:=
\mathsf S_{\,X[f]\,X[g]}^{\mathcal A}(\omega)
=
\sum_{\nu\in\mathcal W_{\rm targ}}
(t_\nu^{X[f]})^*t_\nu^{X[g]}
L_\eta(\omega-\omega_\nu^{\rm QSE}),\\
\mathsf S_{nX}^{\mathcal A}[f,g](\omega)
&:=
\mathsf S_{\,n[f]\,X[g]}^{\mathcal A}(\omega)
=
\sum_{\nu\in\mathcal W_{\rm targ}}
(t_\nu^{n[f]})^*t_\nu^{X[g]}
L_\eta(\omega-\omega_\nu^{\rm QSE}) .
\end{aligned}
\] When \(f=g=f_m^{\rm OBC}\), these are the open-chain standing-wave
spectra. In compact plot labels one may write \[
\mathsf S_{nn}(q_m,\omega),
\qquad
\mathsf S_{XX}(q_m,\omega),
\qquad
\mathsf S_{nX}(q_m,\omega),
\] with \(q_m\) understood as the open-chain standing-wave label. Under
periodic boundaries, take \(f=g=f_k^{\rm PBC}\) and these become the
usual \(\mathsf S_{nn}(k,\omega)\), \(\mathsf S_{XX}(k,\omega)\), and
\(\mathsf S_{nX}(k,\omega)\).

Unfold the density channel. Let \[
N_{\nu j}:=\langle\Psi_\nu^{\rm QSE}|\delta n_j|\psi_0\rangle .
\] Then \[
t_\nu^{n[f]}
=
\sum_j f(j)N_{\nu j}.
\] Therefore \[
\begin{aligned}
|t_\nu^{n[f]}|^2
&=
\left(\sum_i f(i)N_{\nu i}\right)^*
\left(\sum_j f(j)N_{\nu j}\right)\\
&=
\sum_{i,j}
f(i)^*f(j)N_{\nu i}^*N_{\nu j}\\
&=
\sum_{i,j}
f(i)^*f(j)
\langle\psi_0|\delta n_i|\Psi_\nu^{\rm QSE}\rangle
\langle\Psi_\nu^{\rm QSE}|\delta n_j|\psi_0\rangle .
\end{aligned}
\] For the same-probe density spectrum, \[
\begin{aligned}
\mathsf S_{nn}^{\mathcal A}[f,f](\omega)
&=
\sum_{\nu\in\mathcal W_{\rm targ}}
\sum_{i,j}
f(i)^*f(j)
\langle\psi_0|\delta n_i|\Psi_\nu^{\rm QSE}\rangle
\langle\Psi_\nu^{\rm QSE}|\delta n_j|\psi_0\rangle\\
&\qquad\qquad\times
L_\eta(\omega-\omega_\nu^{\rm QSE}) .
\end{aligned}
\] This is the charge-density spectral response in the spatial pattern
\(f\). It reports which roots are visible to charge fluctuations of that
pattern.

The same derivation for displacement uses \[
X_{\nu j}:=\langle\Psi_\nu^{\rm QSE}|X_j|\psi_0\rangle,
\qquad
t_\nu^{X[f]}=\sum_j f(j)X_{\nu j}.
\] Thus \[
\begin{aligned}
\mathsf S_{XX}^{\mathcal A}[f,f](\omega)
&=
\sum_{\nu\in\mathcal W_{\rm targ}}
\sum_{i,j}
f(i)^*f(j)
\langle\psi_0|X_i|\Psi_\nu^{\rm QSE}\rangle
\langle\Psi_\nu^{\rm QSE}|X_j|\psi_0\rangle\\
&\qquad\qquad\times
L_\eta(\omega-\omega_\nu^{\rm QSE}) .
\end{aligned}
\] This channel sees oscillator/displacement spectral weight. Phonon
peaks and phonon sidebands should appear here even when density weight
is weak.

For the mixed channel, use \[
t_\nu^{n[f]}=\sum_i f(i)N_{\nu i},
\qquad
t_\nu^{X[g]}=\sum_j g(j)X_{\nu j}.
\] Then \[
\begin{aligned}
(t_\nu^{n[f]})^*t_\nu^{X[g]}
&=
\left(\sum_i f(i)N_{\nu i}\right)^*
\left(\sum_j g(j)X_{\nu j}\right)\\
&=
\sum_{i,j}
f(i)^*g(j)
\langle\psi_0|\delta n_i|\Psi_\nu^{\rm QSE}\rangle
\langle\Psi_\nu^{\rm QSE}|X_j|\psi_0\rangle .
\end{aligned}
\] Therefore \[
\begin{aligned}
\mathsf S_{nX}^{\mathcal A}[f,g](\omega)
&=
\sum_{\nu\in\mathcal W_{\rm targ}}
\sum_{i,j}
f(i)^*g(j)
\langle\psi_0|\delta n_i|\Psi_\nu^{\rm QSE}\rangle
\langle\Psi_\nu^{\rm QSE}|X_j|\psi_0\rangle\\
&\qquad\qquad\times
L_\eta(\omega-\omega_\nu^{\rm QSE}) .
\end{aligned}
\] This object can be complex. Its magnitude records joint visibility to
density and displacement probes; its phase records the relative phase
between the two channels under the chosen \(f,g\) convention. This is
the response-matrix version of electron-phonon dressing.

For periodic plane waves, substituting \(f(i)=g(i)=e^{-ikr_i}\) gives
the familiar phase factor \[
f(i)^*g(j)=e^{ik(r_i-r_j)}.
\] Thus the periodic density channel becomes \[
\mathsf S_{nn}^{\mathcal A}(k,\omega)
=
\sum_{\nu,i,j}
e^{ik(r_i-r_j)}
\langle\psi_0|\delta n_i|\Psi_\nu^{\rm QSE}\rangle
\langle\Psi_\nu^{\rm QSE}|\delta n_j|\psi_0\rangle
L_\eta(\omega-\omega_\nu^{\rm QSE}),
\] up to the chosen Fourier normalization. On the open chain, replace
\(e^{ik(r_i-r_j)}\) by \(f_m^{\rm OBC}(i)f_m^{\rm OBC}(j)\) for same
standing-wave probes, or by \(f(i)^*g(j)\) for general detector shapes.

The physical readings are:

\begin{enumerate}
\def\labelenumi{\arabic{enumi}.}
\tightlist
\item
  \(\mathsf S_{nn}^{\mathcal A}[f,f](\omega)\) is charge-visible
  spectral weight. A bright line means the root changes the electronic
  density pattern selected by \(f\).
\item
  \(\mathsf S_{XX}^{\mathcal A}[f,f](\omega)\) is phonon-visible
  spectral weight. A bright line means the root moves the oscillator
  coordinate in that spatial pattern.
\item
  \(\mathsf S_{nX}^{\mathcal A}[f,g](\omega)\) is mixed charge-lattice
  response. A large residue \((t_\nu^{n[f]})^*t_\nu^{X[g]}\) means the
  same root carries both density and displacement character.
\end{enumerate}

Now the cutoff diagnostic. Because the oscillator Hilbert space is
truncated at \(n_{\rm ph,max}\), the highest retained phonon level is
artificial: in the infinite oscillator,
\(b^\dagger|n_{\rm ph,max}\rangle\) continues to
\(|n_{\rm ph,max}+1\rangle\), while the finite register has no such
state. Track cutoff stress with \[
\ell_{\rm cut}(\Psi)
:=
\sum_\mu
\langle\Psi|\Pi_{n_\mu=n_{\rm ph,max}}|\Psi\rangle .
\] Here \(\mu\) indexes local phonon modes, and
\(\Pi_{n_\mu=n_{\rm ph,max}}\) projects onto states where mode \(\mu\)
sits at the top retained occupation.

Unfold the diagnostic by writing \[
|\Psi\rangle
=
\sum_{\xi,\mathbf n}
A_{\xi,\mathbf n}
|\xi\rangle\otimes|\mathbf n\rangle,
\qquad
\mathbf n=(n_0,\ldots,n_{L-1}),
\qquad
0\le n_\mu\le n_{\rm ph,max},
\] where \(\xi\) labels the electronic occupation configuration. The
projector acts as \[
\Pi_{n_\mu=n_{\rm ph,max}}
\left(|\xi\rangle\otimes|n_0,\ldots,n_\mu,\ldots,n_{L-1}\rangle\right)
=
\begin{cases}
|\xi\rangle\otimes|\mathbf n\rangle,&n_\mu=n_{\rm ph,max},\\
0,&n_\mu\ne n_{\rm ph,max}.
\end{cases}
\] Therefore \[
\begin{aligned}
\langle\Psi|\Pi_{n_\mu=n_{\rm ph,max}}|\Psi\rangle
&=
\sum_{\xi,\mathbf n}
|A_{\xi,\mathbf n}|^2
\mathbf 1[n_\mu=n_{\rm ph,max}],\\
\ell_{\rm cut}(\Psi)
&=
\sum_\mu
\sum_{\xi,\mathbf n}
|A_{\xi,\mathbf n}|^2
\mathbf 1[n_\mu=n_{\rm ph,max}].
\end{aligned}
\] Thus \(\ell_{\rm cut}\) counts probability mass on the phonon
boundary, summed over modes. It can exceed \(1\) because several modes
can sit at the cutoff simultaneously.

The algebraic reason for tracking this quantity is sharp. Let \[
d:=n_{\rm ph,max}+1
\] be the truncated oscillator dimension, and define \[
b_d
=
\sum_{n=1}^{d-1}\sqrt n\,|n-1\rangle\langle n|,
\qquad
b_d^\dagger
=
\sum_{n=0}^{d-2}\sqrt{n+1}\,|n+1\rangle\langle n|.
\] Then \[
\begin{aligned}
b_db_d^\dagger
&=
\sum_{m=0}^{d-2}(m+1)|m\rangle\langle m|,\\
b_d^\dagger b_d
&=
\sum_{m=1}^{d-1}m|m\rangle\langle m|.
\end{aligned}
\] Hence \[
\begin{aligned}
\left[b_d,b_d^\dagger\right]
&=
\sum_{m=0}^{d-2}(m+1)|m\rangle\langle m|
-
\sum_{m=1}^{d-1}m|m\rangle\langle m|\\
&=
\sum_{m=0}^{d-2}|m\rangle\langle m|
-
(d-1)|d-1\rangle\langle d-1|\\
&=
I-d|d-1\rangle\langle d-1|.
\end{aligned}
\] The infinite oscillator has \([b,b^\dagger]=I\). The truncated
oscillator places the commutator defect entirely on the top retained
level. Therefore a QSE root with large \(\ell_{\rm cut}\) can have
distorted \(X=b+b^\dagger\), \(P=i(b^\dagger-b)\), and phonon sideband
weights. Reporting \(\ell_{\rm cut}\) alongside response errors makes
this failure mode visible: a spectral peak near the cutoff can be a
finite-register artifact, and an accurate electronic gap can coexist
with an unreliable phonon sideband.

The channel definitions collected from this section are: \[
\begin{aligned}
\delta n_i&=n_i-\bar n I,
&
X_i&=b_i^\dagger+b_i,
&
P_i&=i(b_i^\dagger-b_i),\\
C_{nX}(r)&=\sum_{i:\,0\le i,\,i+r<L}\delta n_iX_{i+r},
&
n[f]&=\sum_j f(j)\delta n_j,
&
X[f]&=\sum_j f(j)X_j .
\end{aligned}
\] The first row asks for charge fluctuation, displacement, and the
conjugate phonon quadrature. The second row asks for local
electron-phonon dressing and for the same channels in a chosen spatial
pattern. \[
\begin{aligned}
\mathsf S_{nn}^{\mathcal A}[f,f](\omega)
&=
\sum_\nu|t_\nu^{n[f]}|^2L_\eta(\omega-\omega_\nu^{\rm QSE}),\\
\mathsf S_{XX}^{\mathcal A}[f,f](\omega)
&=
\sum_\nu|t_\nu^{X[f]}|^2L_\eta(\omega-\omega_\nu^{\rm QSE}),\\
\mathsf S_{nX}^{\mathcal A}[f,g](\omega)
&=
\sum_\nu(t_\nu^{n[f]})^*t_\nu^{X[g]}
L_\eta(\omega-\omega_\nu^{\rm QSE}),\\
\ell_{\rm cut}(\Psi)
&=
\sum_\mu\langle\Psi|\Pi_{n_\mu=n_{\rm ph,max}}|\Psi\rangle .
\end{aligned}
\] These are charge-visible spectral weight, phonon-visible spectral
weight, mixed charge-lattice dressing, and phonon cutoff stress,
respectively. Section 18.5 adds physically meaningful probes to the same
transition-amplitude machinery from Sections 18.3 and 18.4. The QSE
eigenproblem is already fixed by the selected subspace; this section
chooses probes for reading the roots. Under open boundaries, \(q_m\) or
\(f\) labels the detector shape; under periodic boundaries, \(k\)
becomes a genuine lattice momentum. In both cases, \(\mathsf S_{nn}\),
\(\mathsf S_{XX}\), and \(\mathsf S_{nX}\) separate charge-like,
phonon-like, and dressed response content.

\subsection{18.6 Single-particle Green
functions}\label{single-particle-green-functions}

A neutral response channel in Sections 18.3--18.5 applies an operator
\(O\) that preserves the particle-number sector of the seed, so
\(O|\psi_0\rangle\in\mathcal H_N\). A single-particle Green function
applies a fermion creation or annihilation operator. It therefore asks a
different question: what spectrum is seen after inserting an electron or
removing one? The relevant QSE spaces are the addition sector
\(\mathcal H_{N+1}\) and the removal sector \(\mathcal H_{N-1}\). In a
Hubbard--Holstein system these sectors can contain a coherent
quasiparticle peak together with phonon sidebands, because the added
electron or removed electron can carry oscillator displacement with it.

In compact physical language, \(c_q^\dagger|\psi_0\rangle\) means ``put
an electron into spin-orbital \(q\), then let the Hamiltonian respond,''
while \(c_p|\psi_0\rangle\) means ``remove an electron from spin-orbital
\(p\), then let the Hamiltonian respond.'' The retarded Green function
\(G_{pq}^R(\omega)\) asks whether an electron inserted at \(q\) later
appears at \(p\), plus the analogous hole propagation from removing an
electron. This is why the final frequency-domain formula has two sums:
one over \(N+1\) addition roots and one over \(N-1\) removal roots.

Interpretive compression: the single-particle Green function is a
sector-resolved resolvent matrix element of the fermionic field algebra.
Neutral QSE uses \(\mathcal K_{\mathcal A}\subset\mathcal H_N\); \(G^R\)
factors through two different projected spectral measures, one on
\(\mathcal H_{N+1}\) for electron injection and one on
\(\mathcal H_{N-1}\) for hole propagation. The residues are
quasiparticle/hole overlap factors, and in HH they also encode
phonon-cloud dressing through the addition/removal Ritz states.

Single-particle spectral terminology: a quasiparticle pole is an
addition or removal pole with large coherent residue for a bare fermion
source \(c_q^\dagger|\psi_0\rangle\) or \(c_p|\psi_0\rangle\). A band is
the graph of such pole locations as the orbital, spatial form factor, or
momentum-like label is varied. A spectral gap is a frequency interval
with suppressed \(A_{pp}(\omega)\), often measured between the lowest
visible removal edge and the lowest visible addition edge under the
chosen chemical-potential convention. A phonon sideband is a secondary
peak separated from a quasiparticle peak by an oscillator-like energy
scale, indicating that fermion insertion or removal also excites lattice
displacement. In HH language, the Green function tests propagation of
dressed electrons and dressed holes; neutral charge and phonon response
are handled by the fixed-\(N\) channels above.

\subsubsection{§18.6 lesson: single-particle Green
functions}\label{lesson-single-particle-green-functions}

Let \(c_p^\dagger\) create a fermion in spin-orbital \(p\), and let
\(c_p\) annihilate one. The canonical anticommutation relations are \[
\{c_p,c_q^\dagger\}
:=
c_pc_q^\dagger+c_q^\dagger c_p
=\delta_{pq}I,\qquad
\{c_p,c_q\}=0,\qquad
\{c_p^\dagger,c_q^\dagger\}=0 .
\] The braces \(\{\cdot,\cdot\}\) denote an anticommutator. Fermion
fields use anticommutators because exchanging two fermionic operators
carries a sign; the equal-time relation above is the algebraic statement
that a spin-orbital can be empty or occupied once. In retarded
single-fermion response, the anticommutator gives the standard causal
propagator and the spectral-weight sum rule
\(\int A_{pp}(\omega)\,d\omega=1\) for a complete
basis.\footnote{Bosonic or even observables usually use commutators in retarded susceptibilities, such as $-i\Theta(t)\langle[A(t),B(0)]\rangle$. Single fermion operators are odd under fermion parity, so the corresponding retarded Green function uses an anticommutator.}

The Heisenberg operator is \[
c_p(t):=e^{iH_0t}c_pe^{-iH_0t}.
\] The operator carries the time dependence; the seed state is held
fixed. The Heaviside step function is \[
\Theta(t)=
\begin{cases}
1,&t>0,\\
0,&t<0,
\end{cases}
\] with the value at \(t=0\) immaterial for the frequency-domain
formulas below. It enforces retardedness: the measured response begins
after the source is applied.

The zero-temperature retarded Green function is \[
G_{pq}^R(t)
:=
-i\,\Theta(t)
\langle\psi_0|\{c_p(t),c_q^\dagger(0)\}|\psi_0\rangle .
\] Expanding the anticommutator separates the two physical processes: \[
\begin{aligned}
G_{pq}^R(t)
&=
-i\Theta(t)
\left[
\langle\psi_0|c_p(t)c_q^\dagger|\psi_0\rangle
+
\langle\psi_0|c_q^\dagger c_p(t)|\psi_0\rangle
\right].
\end{aligned}
\] The first term adds a particle with \(c_q^\dagger\), evolves in
\(\mathcal H_{N+1}\), and removes it with \(c_p\). The second term
removes a particle with \(c_p\), evolves in \(\mathcal H_{N-1}\), and
fills the hole with \(c_q^\dagger\).

Assume the seed has fixed particle number: \[
H_0|\psi_0\rangle=E_0|\psi_0\rangle,\qquad
\widehat N|\psi_0\rangle=N|\psi_0\rangle .
\] Let exact eigenstates in the addition and removal sectors obey \[
\begin{aligned}
H_0|\Psi_m^{N+1}\rangle&=E_m^{N+1}|\Psi_m^{N+1}\rangle,\\
H_0|\Psi_n^{N-1}\rangle&=E_n^{N-1}|\Psi_n^{N-1}\rangle .
\end{aligned}
\] Their sector identities are \[
\begin{aligned}
I_{N+1}
&=\sum_m|\Psi_m^{N+1}\rangle\langle\Psi_m^{N+1}|,\\
I_{N-1}
&=\sum_n|\Psi_n^{N-1}\rangle\langle\Psi_n^{N-1}| .
\end{aligned}
\] These identities are inserted only in the sector reached by the
source state \(c_q^\dagger|\psi_0\rangle\) or \(c_p|\psi_0\rangle\).

For the addition term, since
\(c_q^\dagger|\psi_0\rangle\in\mathcal H_{N+1}\), \[
\begin{aligned}
\langle\psi_0|c_p(t)c_q^\dagger|\psi_0\rangle
&=
\langle\psi_0|e^{iH_0t}c_pe^{-iH_0t}c_q^\dagger|\psi_0\rangle\\
&=
e^{iE_0t}
\langle\psi_0|c_pe^{-iH_0t}c_q^\dagger|\psi_0\rangle\\
&=
e^{iE_0t}
\sum_m
\langle\psi_0|c_pe^{-iH_0t}|\Psi_m^{N+1}\rangle
\langle\Psi_m^{N+1}|c_q^\dagger|\psi_0\rangle\\
&=
\sum_m
\langle\psi_0|c_p|\Psi_m^{N+1}\rangle
\langle\Psi_m^{N+1}|c_q^\dagger|\psi_0\rangle
e^{-i(E_m^{N+1}-E_0)t}.
\end{aligned}
\] Define the addition energy \[
\Delta_m^+:=E_m^{N+1}-E_0 .
\] For the removal term, since \(c_p|\psi_0\rangle\in\mathcal H_{N-1}\),
\[
\begin{aligned}
\langle\psi_0|c_q^\dagger c_p(t)|\psi_0\rangle
&=
\langle\psi_0|c_q^\dagger e^{iH_0t}c_pe^{-iH_0t}|\psi_0\rangle\\
&=
e^{-iE_0t}
\langle\psi_0|c_q^\dagger e^{iH_0t}c_p|\psi_0\rangle\\
&=
e^{-iE_0t}
\sum_n
\langle\psi_0|c_q^\dagger e^{iH_0t}|\Psi_n^{N-1}\rangle
\langle\Psi_n^{N-1}|c_p|\psi_0\rangle\\
&=
\sum_n
\langle\psi_0|c_q^\dagger|\Psi_n^{N-1}\rangle
\langle\Psi_n^{N-1}|c_p|\psi_0\rangle
e^{+i(E_n^{N-1}-E_0)t}.
\end{aligned}
\] Define \[
\Delta_n^-:=E_n^{N-1}-E_0 .
\] Thus the time-domain Lehmann form is \[
\begin{aligned}
G_{pq}^R(t)
=
-i\Theta(t)
\Bigg[
&\sum_m
\langle\psi_0|c_p|\Psi_m^{N+1}\rangle
\langle\Psi_m^{N+1}|c_q^\dagger|\psi_0\rangle
e^{-i\Delta_m^+t}\\
&+
\sum_n
\langle\psi_0|c_q^\dagger|\Psi_n^{N-1}\rangle
\langle\Psi_n^{N-1}|c_p|\psi_0\rangle
e^{+i\Delta_n^-t}
\Bigg].
\end{aligned}
\]

Now Fourier transform with the convention \[
G_{pq}^R(\omega)
:=
\int_{-\infty}^{\infty}dt\,e^{i\omega t}G_{pq}^R(t).
\] For \(t>0\), insert the convergence factor \(e^{-\eta t}\),
\(\eta>0\), and then either take \(\eta\to0^+\) as a distributional
boundary value or keep finite \(\eta\) as the plotted
broadening.\footnote{The retarded Fourier integral is an oscillatory integral. Multiplying by $e^{-\eta t}$ makes it absolutely convergent and selects the causal boundary value $\omega+i0^+$. Physics interprets finite $\eta$ as finite resolution, finite lifetime, or deliberate line broadening. Mathematically the exact retarded function is recovered by the limit $\eta\downarrow0$ after the integral is evaluated.}
The addition exponential gives \[
\begin{aligned}
-i\int_0^\infty dt\,
e^{i\omega t}e^{-i\Delta_m^+t}e^{-\eta t}
&=
-i\int_0^\infty dt\,
e^{-[\eta-i(\omega-\Delta_m^+)]t}\\
&=
\frac{1}{\omega-\Delta_m^+ + i\eta}.
\end{aligned}
\] The removal exponential gives \[
\begin{aligned}
-i\int_0^\infty dt\,
e^{i\omega t}e^{+i\Delta_n^-t}e^{-\eta t}
&=
-i\int_0^\infty dt\,
e^{-[\eta-i(\omega+\Delta_n^-)]t}\\
&=
\frac{1}{\omega+\Delta_n^- + i\eta}.
\end{aligned}
\] Therefore \[
\begin{aligned}
G_{pq}^R(\omega)
&=
\sum_m
\frac{
\langle\psi_0|c_p|\Psi_m^{N+1}\rangle
\langle\Psi_m^{N+1}|c_q^\dagger|\psi_0\rangle}
{\omega-(E_m^{N+1}-E_0)+i\eta}\\
&\quad+
\sum_n
\frac{
\langle\psi_0|c_q^\dagger|\Psi_n^{N-1}\rangle
\langle\Psi_n^{N-1}|c_p|\psi_0\rangle}
{\omega+(E_n^{N-1}-E_0)+i\eta}.
\end{aligned}
\] This is the single-particle Lehmann representation. The denominators
locate poles. A pole is a value of complex \(\omega\) where the response
function diverges; after finite broadening, that mathematical divergence
is seen as a peak of finite width in a plotted spectrum.

The spectral function is \[
A_{pq}(\omega):=-\frac1\pi\,\operatorname{Im}G_{pq}^R(\omega).
\] The sign convention converts a retarded denominator into a positive
Lorentzian: \[
\frac{1}{x+i\eta}
=
\frac{x-i\eta}{x^2+\eta^2},
\qquad
-\frac1\pi\operatorname{Im}\frac{1}{x+i\eta}
=
\frac1\pi\frac{\eta}{x^2+\eta^2}.
\] So a visual spectral plot is a collection of broadened peaks. The
horizontal location is the addition or removal energy under the chosen
frequency zero; the peak weight is the corresponding residue. A
schematic plot can help intuition; a manuscript-facing figure is
strongest when generated from computed \(A(\omega)\) data.

The indices \(p,q\) label spin-orbitals. A local spectral density is a
trace or average over a chosen orbital set \(\mathcal P_G\): \[
\begin{aligned}
A_{\rm loc}(\omega)
&:=
\frac1{|\mathcal P_G|}
\sum_{p\in\mathcal P_G}A_{pp}(\omega)\\
&=
-\frac1{\pi|\mathcal P_G|}
\sum_{p\in\mathcal P_G}
\operatorname{Im}G_{pp}^R(\omega).
\end{aligned}
\] This contraction answers a deliberately coarse question: averaged
over these orbitals, how much electron-addition or electron-removal
spectral weight exists at frequency \(\omega\)? It is a trace over the
Green-function matrix. The same sector QSE roots supply the matrix
entries; the contraction only changes how the orbital indices are
summarized.

Mode-resolved spectra are basis changes on the orbital indices. For an
open chain, use the same detector-shape convention as Section 18.5. Let
\(f(j)\) and \(g(j)\) be spatial form factors and define \[
c_{f\sigma}:=\sum_j f(j)c_{j\sigma},
\qquad
c_{g\sigma}^\dagger:=\sum_\ell g(\ell)^*c_{\ell\sigma}^\dagger .
\] Then \[
\begin{aligned}
G_{fg,\sigma}^R(\omega)
&:=
G_{c_{f\sigma},c_{g\sigma}}^R(\omega)\\
&=
\sum_{j,\ell}
f(j)g(\ell)^*
G_{j\sigma,\ell\sigma}^R(\omega),
\end{aligned}
\] and
\(A_{f\sigma}(\omega)=-(1/\pi)\operatorname{Im}G_{ff,\sigma}^R(\omega)\).
Taking \(f=f_m^{\rm OBC}\) gives an open-chain standing-wave spectrum.
In a periodic chain, \(f_k(j)=L^{-1/2}e^{-ikr_j}\), so \[
\begin{aligned}
c_{k\sigma}
&=
\frac1{\sqrt L}\sum_j e^{-ikr_j}c_{j\sigma},\\
G_{k\sigma}^R(\omega)
&=
\frac1L
\sum_{j,\ell}
e^{-ikr_j}e^{+ikr_\ell}
G_{j\sigma,\ell\sigma}^R(\omega).
\end{aligned}
\] Thus the local, open-mode, and periodic-momentum spectra are
different contractions of the same matrix-valued Green function.

Now QSE enters. One repeats the nonorthogonal QSE construction in two
different Hilbert sectors. The neutral basis from \(\mathcal H_N\)
supplies useful notation and selected-record intuition, but the
single-particle Green function requires new addition and removal sector
kets. For source orbitals \(q,p\), define \[
|\chi_q^+\rangle:=c_q^\dagger|\psi_0\rangle,
\qquad
|\chi_p^-\rangle:=c_p|\psi_0\rangle .
\] A record \(a\) carries an operator or circuit action \(\Omega_a\). In
this section, \(\Omega_a\) means the selected excitation/action record
used to generate QSE basis kets from a source state. Applying
\(\Omega_a\) to \(c_q^\dagger|\psi_0\rangle\) or \(c_p|\psi_0\rangle\)
builds candidate states in the particle-addition or particle-removal
sector. With symmetry-sector projectors, \[
\begin{aligned}
|\phi_{a,q}^+\rangle
&:=
\Pi_{N+1,\mathfrak s_+}\Omega_a|\chi_q^+\rangle,\\
|\phi_{a,p}^-\rangle
&:=
\Pi_{N-1,\mathfrak s_-}\Omega_a|\chi_p^-\rangle .
\end{aligned}
\] Here \(\Pi_{N+1,\mathfrak s_+}\) and \(\Pi_{N-1,\mathfrak s_-}\)
enforce the intended particle-number and symmetry sectors. The basis may
be source-specific, meaning a separate basis for each source orbital, or
unioned across a set \(\mathcal P_G\) so one basis supports a full
Green-function matrix.

Write the selected addition and removal synthesis maps as \[
\Phi^+
=
\begin{bmatrix}
|\phi_1^+\rangle&\cdots&|\phi_{r_+}^+\rangle
\end{bmatrix},
\qquad
\Phi^-
=
\begin{bmatrix}
|\phi_1^-\rangle&\cdots&|\phi_{r_-}^-\rangle
\end{bmatrix}.
\] Each sector has its own overlap metric and projected Hamiltonian: \[
\begin{aligned}
S^+&:=(\Phi^+)^\dagger\Phi^+,&
M^+&:=(\Phi^+)^\dagger H_0\Phi^+,\\
S^-&:=(\Phi^-)^\dagger\Phi^-,&
M^-&:=(\Phi^-)^\dagger H_0\Phi^- .
\end{aligned}
\] The two QSE pencils are \[
M^+c_m^+=\varepsilon_m^+S^+c_m^+,\qquad
M^-c_n^-=\varepsilon_n^-S^-c_n^-,
\] with synthesized sector roots \[
|\Psi_m^+\rangle:=\Phi^+c_m^+,
\qquad
|\Psi_n^-\rangle:=\Phi^-c_n^- .
\] This answers the practical question directly: the plus and minus
sectors require separate QSE builds, separate overlap matrices, separate
eigensolves, and separate root tracking. They can share the same design
rule for selecting records, with source states satisfying
\(c_q^\dagger|\psi_0\rangle\in\mathcal H_{N+1}\) and
\(c_p|\psi_0\rangle\in\mathcal H_{N-1}\).

The addition transition vector for orbital \(p\) is \[
(d_p^+)_a:=\langle\phi_a^+|c_p^\dagger|\psi_0\rangle .
\] It records how strongly the source \(c_p^\dagger|\psi_0\rangle\)
overlaps each addition-sector QSE bra. The addition amplitude into root
\(m\) is obtained by mixing those components with the QSE root
coefficients: \[
\begin{aligned}
(d_p^+)^\dagger c_m^+
&=
\sum_a\left((d_p^+)_a\right)^*(c_m^+)_a\\
&=
\sum_a
\langle\psi_0|c_p|\phi_a^+\rangle
(c_m^+)_a\\
&=
\langle\psi_0|c_p
\left(\sum_a(c_m^+)_a|\phi_a^+\rangle\right)\\
&=
\langle\psi_0|c_p|\Psi_m^+\rangle .
\end{aligned}
\] Define \[
u_{pm}^+
:=
\langle\psi_0|c_p|\Psi_m^+\rangle
=(d_p^+)^\dagger c_m^+ .
\] This is the addition-sector residue amplitude. It says how strongly
the \(N+1\) QSE root can be annihilated back into the seed by orbital
\(p\).

The removal transition vector is \[
(d_p^-)_a:=\langle\phi_a^-|c_p|\psi_0\rangle .
\] It records how strongly the removed-particle source
\(c_p|\psi_0\rangle\) overlaps each removal-sector QSE bra. The removal
amplitude is \[
\begin{aligned}
(c_n^-)^\dagger d_p^-
&=
\sum_a(c_n^-)_a^*(d_p^-)_a\\
&=
\sum_a(c_n^-)_a^*
\langle\phi_a^-|c_p|\psi_0\rangle\\
&=
\left(\sum_a(c_n^-)_a|\phi_a^-\rangle\right)^\dagger
c_p|\psi_0\rangle\\
&=
\langle\Psi_n^-|c_p|\psi_0\rangle .
\end{aligned}
\] Define \[
v_{pn}^-
:=
\langle\Psi_n^-|c_p|\psi_0\rangle
=(c_n^-)^\dagger d_p^- .
\] Thus \(u_{pm}^+\) and \(v_{pn}^-\) are the QSE versions of the exact
Lehmann matrix elements.

Substitute these amplitudes into the Lehmann representation. The QSE
addition residue is \[
\langle\psi_0|c_p|\Psi_m^+\rangle
\langle\Psi_m^+|c_q^\dagger|\psi_0\rangle
=
u_{pm}^+(u_{qm}^+)^*,
\] and the QSE removal residue is \[
\langle\psi_0|c_q^\dagger|\Psi_n^-\rangle
\langle\Psi_n^-|c_p|\psi_0\rangle
=
(v_{qn}^-)^*v_{pn}^- .
\] Therefore \[
\begin{aligned}
G_{pq}^{R,\mathcal A}(\omega)
&=
\sum_m
\frac{
u_{pm}^+(u_{qm}^+)^*
}{
\omega-(\varepsilon_m^+-E_0)+i\eta
}\\
&\quad+
\sum_n
\frac{
(v_{qn}^-)^*v_{pn}^-
}{
\omega+(\varepsilon_n^--E_0)+i\eta
}.
\end{aligned}
\]

Metric support and whitening occur separately in the two sectors. If \[
S^\pm=U^\pm D^\pm(U^\pm)^\dagger,
\] define retained supports \[
J_\tau^\pm:=\{j:s_j^\pm>\tau_\pm\},
\qquad
W_\tau^\pm:=U_\tau^\pm(D_\tau^\pm)^{-1/2}.
\] Then \[
c_m^+=W_\tau^+y_m^+,\qquad
c_n^-=W_\tau^-y_n^- .
\] The whitened addition probe vector is \[
g_p^+:=(W_\tau^+)^\dagger d_p^+,
\] so \[
\begin{aligned}
u_{pm}^+
&=(d_p^+)^\dagger W_\tau^+y_m^+\\
&=\left((W_\tau^+)^\dagger d_p^+\right)^\dagger y_m^+\\
&=(g_p^+)^\dagger y_m^+ .
\end{aligned}
\] The whitened removal probe vector is \[
g_p^-:=(W_\tau^-)^\dagger d_p^-,
\] so \[
\begin{aligned}
v_{pn}^-
&=(W_\tau^-y_n^-)^\dagger d_p^-\\
&=(y_n^-)^\dagger(W_\tau^-)^\dagger d_p^-\\
&=(y_n^-)^\dagger g_p^- .
\end{aligned}
\] The same metric idea from Section 18.2 is applied twice. Addition and
removal have different \(S^\pm\), different retained supports, different
conditioning, and possibly different cutoff stress. Treating them as
separate sectors is part of the scientific result, because missing
weight or unstable roots can occur on one side of the spectrum while the
other side is well resolved.

The exact diagonal spectral-weight sum rule shows what complete sector
bases recover. For \(p=q\), the addition weight is \[
\begin{aligned}
W_p^+
&:=
\sum_m
\left|
\langle\Psi_m^{N+1}|c_p^\dagger|\psi_0\rangle
\right|^2\\
&=
\langle\psi_0|
c_p
\left(\sum_m|\Psi_m^{N+1}\rangle\langle\Psi_m^{N+1}|\right)
c_p^\dagger
|\psi_0\rangle\\
&=
\langle\psi_0|c_pc_p^\dagger|\psi_0\rangle\\
&=
1-\langle\psi_0|c_p^\dagger c_p|\psi_0\rangle
=1-\langle n_p\rangle .
\end{aligned}
\] The removal weight is \[
\begin{aligned}
W_p^-
&:=
\sum_n
\left|
\langle\Psi_n^{N-1}|c_p|\psi_0\rangle
\right|^2\\
&=
\langle\psi_0|
c_p^\dagger
\left(\sum_n|\Psi_n^{N-1}\rangle\langle\Psi_n^{N-1}|\right)
c_p
|\psi_0\rangle\\
&=
\langle\psi_0|c_p^\dagger c_p|\psi_0\rangle
=\langle n_p\rangle .
\end{aligned}
\] Thus \[
W_p^+ + W_p^- = 1,
\qquad
\int_{-\infty}^{\infty}A_{pp}(\omega)\,d\omega=1
\] for a complete, correctly normalized fermionic Green function with
normalized broadening. In QSE, a deficit in this sum rule means the
selected \(N+1\) and \(N-1\) subspaces miss addition or removal spectral
weight.

The frequency zero must be stated. With the seed energy convention
above, the addition poles are \[
\omega_m^+=\varepsilon_m^+-E_0,
\] and the removal poles are \[
\omega_n^-=-(\varepsilon_n^--E_0).
\] If a chemical potential is used, replace \(H_0\) by the
grand-canonical operator \[
K:=H_0-\mu\widehat N .
\] Then \[
\begin{aligned}
\omega_m^+
&=(E_m^{N+1}-E_0)-\mu,\\
\omega_n^-
&=-(E_n^{N-1}-E_0)-\mu .
\end{aligned}
\] Different Green-function papers choose different frequency zeros,
especially when comparing finite clusters, half filling, or
electron-addition/removal gaps. A benchmark must report the convention,
the chemical potential choice if any, and the broadening \(\eta\).
Literature comparison is then a matter of aligning these conventions
before comparing peak locations.

The practical contract for a Green-function benchmark is therefore: \[
\begin{gathered}
\text{state which }N+1\text{ sectors are included},\qquad
\text{state which }N-1\text{ sectors are included},\\
\text{state whether bases are source-specific or unioned over }\mathcal P_G,\\
\text{state how roots are matched and tracked},\\
\text{state the frequency zero or chemical potential},\qquad
\text{state the broadening }\eta .
\end{gathered}
\]

For Hubbard--Holstein reading, the Green function is the spectrum of
injecting or extracting a fermion. A bare electron source
\(c_q^\dagger|\psi_0\rangle\) can evolve into a dressed electron plus
lattice displacement. A bare removal source \(c_p|\psi_0\rangle\) can
leave behind a dressed hole and phonon cloud. Addition and removal peaks
therefore diagnose quasiparticle bands, spectral gaps, and polaronic
sidebands. The neutral response sections ask which fixed-\(N\) modes are
visible to density or phonon probes; this section asks which \(N\pm1\)
modes are visible to particle insertion or extraction.

The implementation contract collected from this section is \[
\begin{aligned}
\text{Neutral QSE: }&\quad O|\psi_0\rangle\in\mathcal H_N,\\
\text{Addition Green sector: }&\quad c_q^\dagger|\psi_0\rangle\in\mathcal H_{N+1},\\
\text{Removal Green sector: }&\quad c_p|\psi_0\rangle\in\mathcal H_{N-1},\\
\text{Sector pencils: }&\quad
M^\pm c^\pm=\varepsilon^\pm S^\pm c^\pm,\\
\text{Addition amplitude: }&\quad
u_{pm}^+=(d_p^+)^\dagger c_m^+,\\
\text{Removal amplitude: }&\quad
v_{pn}^-=(c_n^-)^\dagger d_p^- .
\end{aligned}
\] The output formula is \[
\begin{aligned}
G_{pq}^{R,\mathcal A}(\omega)
&=
\sum_m
\frac{u_{pm}^+(u_{qm}^+)^*}
{\omega-(\varepsilon_m^+-E_0)+i\eta}\\
&\quad+
\sum_n
\frac{(v_{qn}^-)^*v_{pn}^-}
{\omega+(\varepsilon_n^--E_0)+i\eta},
\qquad
A_{pq}(\omega)
=-\frac1\pi\operatorname{Im}G_{pq}^R(\omega).
\end{aligned}
\] Section 18.6 uses the same nonorthogonal QSE geometry as Section
18.2, applied twice. The new work is sector bookkeeping, source
construction, frequency convention, and residue tracking for
particle-addition and particle-removal spectra.

\subsection{18.7 Optical and current
response}\label{optical-and-current-response}

Optical response is the electromagnetic specialization of the response
machinery from Sections 18.3--18.6. Density response asks whether a root
carries charge fluctuation; phonon response asks whether a root carries
oscillator displacement; optical response asks how much current flows
when charged particles are driven by a weak vector potential. The
current-current part identifies roots that absorb light through current
matrix elements. The diamagnetic/contact part is the instantaneous
current response coming from the curvature of the Hamiltonian with
respect to the external vector potential.

Interpretive compression: optical conductivity is the Kubo response of
the Peierls-coupled lattice Hamiltonian. The paramagnetic sector is the
retarded \(JJ\) susceptibility and therefore fits the same pole/residue
grammar as Section 18.4; the diamagnetic sector is the curvature term
required by the electromagnetic coupling convention. The QSE-specific
operation is only the spectral replacement of exact current matrix
elements by supported QSE transition amplitudes.

Physical semantics: the vector potential \(A_{\rm ext}\) is the source
coordinate for lattice charge motion, \(J=-\partial_AH|_0\) is the
conjugate current observable, and \(D=\partial_A^2H|_0\) is the contact
curvature of the same coupling. The conductivity is the frequency-domain
coefficient relating current density to electric field,
\(j(\omega)=\sigma(\omega)E(\omega)\). Hence current response has two
inseparable parts: delayed paramagnetic transport encoded by
\(\chi_{JJ}^R\), and instantaneous diamagnetic/contact response encoded
by \(\langle D\rangle_0\). The spectral peaks in
\({\rm Re}\,\sigma(\omega)\) identify neutral excitations carrying
current matrix element \(t_\nu^J\), so optical response is the
current-channel specialization of the general response matrix.

\subsubsection{§18.7 lesson: optical and current
response}\label{lesson-optical-and-current-response}

The standard lattice contract is the explicit convention for coupling
light to the lattice Hamiltonian. Introduce a spatially uniform external
vector potential \(A_{\rm ext}\), expand the Hamiltonian at zero field,
and define the current and contact operators by the same expansion: \[
\begin{aligned}
H(A_{\rm ext})
&=
H_0-A_{\rm ext}J+\frac12A_{\rm ext}^2D+\mathcal O(A_{\rm ext}^3),\\
J
&:=-\left.\frac{\partial H}{\partial A_{\rm ext}}\right|_{A_{\rm ext}=0},
\qquad
D:=\left.\frac{\partial^2H}{\partial A_{\rm ext}^2}\right|_{A_{\rm ext}=0}.
\end{aligned}
\] This fixes the current operator, sign, normalization, and
diamagnetic/contact convention. The current retarded response is the
same retarded correlation template from Section 18.4 with the probe set
to \(J\): \[
\chi_{JJ}^R(t)
:=
-i\Theta(t)\langle\psi_0|[J(t),J(0)]|\psi_0\rangle .
\] Here \(R\) means retarded, \(JJ\) means current-current, and \(\chi\)
is the susceptibility: after a tiny vector-potential kick couples to
\(J\), \(\chi_{JJ}^R\) describes how current appears later.

Now derive the conductivity convention. Differentiate the
Peierls-expanded Hamiltonian: \[
\begin{aligned}
\frac{\partial H}{\partial A_{\rm ext}}
&=-J+A_{\rm ext}D+\mathcal O(A_{\rm ext}^2),\\
I(A_{\rm ext})
&:=-\frac{\partial H}{\partial A_{\rm ext}}
=J-A_{\rm ext}D+\mathcal O(A_{\rm ext}^2),\\
I(0)&=J .
\end{aligned}
\] Thus \(J\) is the paramagnetic current operator, while
\(-A_{\rm ext}D\) is the first-order diamagnetic/contact contribution to
the physical current.

Let the externally applied vector potential be weak and time-dependent,
\(A_{\rm ext}=A(t)\). Then \[
V(t):=H(A(t))-H_0
=-A(t)J+\frac12A(t)^2D+\cdots,
\qquad
V^{(1)}(t)=-A(t)J .
\] Here \(A(t)\) is an imposed classical drive field. In a gauge with
zero scalar potential, the electric field is \(E(t)=-\dot A(t)\). Linear
response first computes how \(J\) responds to \(A\), then converts
vector-potential response into electric-field response.

The first-order interaction-picture response is \[
\begin{aligned}
\delta\langle J(t)\rangle
&=
i\int_{-\infty}^{t}dt'\,
\langle\psi_0|[V^{(1)}(t'),J(t)]|\psi_0\rangle\\
&=
-i\int_{-\infty}^{t}dt'\,
A(t')\langle\psi_0|[J(t'),J(t)]|\psi_0\rangle\\
&=
i\int_{-\infty}^{t}dt'\,
A(t')\langle\psi_0|[J(t),J(t')]|\psi_0\rangle .
\end{aligned}
\] The two-time response kernel is \[
\chi_{JJ}^R(t-t')
:=
-i\Theta(t-t')\langle\psi_0|[J(t),J(t')]\ |\psi_0\rangle .
\] The argument \(t-t'\) is the delay between the earlier kick at \(t'\)
and the later readout at \(t\). The seed is stationary under \(H_0\), so
the kernel depends only on this time difference. For \(t'\le t\),
\(\Theta(t-t')=1\), and \[
i\langle[J(t),J(t')]\rangle_0=-\chi_{JJ}^R(t-t'),
\qquad
\delta\langle J(t)\rangle
=-\int_{-\infty}^{\infty}dt'\,
\chi_{JJ}^R(t-t')A(t').
\] Fourier transforming turns this memory convolution into
multiplication: \[
\delta\langle J(\omega)\rangle=-\chi_{JJ}^R(\omega)A(\omega).
\] The Fourier transform is useful because the time-domain response
answers ``what current appears after a kick?'', while the
frequency-domain response answers ``what current oscillates at frequency
\(\omega\) when the electric field oscillates at \(\omega\)?'' Energy
eigenstates contribute factors \(e^{-i(E_\nu-E_0)t}\), so Fourier
transformation converts those oscillations into poles or peaks at
\(\omega_\nu=E_\nu-E_0\).

The physical current includes the contact term: \[
\begin{aligned}
\delta\langle I(\omega)\rangle
&=
\delta\langle J(\omega)\rangle-\langle D\rangle_0A(\omega),\\
\langle D\rangle_0
&:=\langle\psi_0|D|\psi_0\rangle,\\
\delta\langle I(\omega)\rangle
&=
-\left(\langle D\rangle_0+\chi_{JJ}^R(\omega)\right)A(\omega).
\end{aligned}
\] Use the Fourier convention
\(A(t)=\int(d\omega/2\pi)e^{-i\omega t}A(\omega)\). Since
\(E(t)=-dA(t)/dt\), \[
\begin{aligned}
E(t)
&=
\int\frac{d\omega}{2\pi}\,i\omega e^{-i\omega t}A(\omega),\\
E(\omega)
&=i(\omega+i\eta)A(\omega),\qquad
A(\omega)=\frac{E(\omega)}{i(\omega+i\eta)} .
\end{aligned}
\] The \(\eta>0\) is the same retarded convergence/broadening parameter
used above. It can be read mathematically as adiabatic switch-on from
the remote past, which selects the causal boundary value
\(\omega+i0^+\), and physically as finite frequency resolution or
lifetime broadening when kept finite.

Substitute \(A(\omega)=E(\omega)/[i(\omega+i\eta)]\), divide current by
system size \(L\), and define
\(j(\omega):=L^{-1}\delta\langle I(\omega)\rangle=\sigma(\omega)E(\omega)\):
\[
\begin{aligned}
j(\omega)
&=
-\frac1L
\left(\langle D\rangle_0+\chi_{JJ}^R(\omega)\right)
\frac{E(\omega)}{i(\omega+i\eta)}\\
&=
\frac{i}{L(\omega+i\eta)}
\left(\langle D\rangle_0+\chi_{JJ}^R(\omega)\right)E(\omega),\\
\sigma(\omega)
&=
\frac{i}{L(\omega+i\eta)}
\left(
\langle D\rangle_0+\chi_{JJ}^R(\omega)
\right).
\end{aligned}
\] This is the Kubo/Peierls conductivity convention used here. The
QSE-specific part is the replacement of exact current-current poles and
residues by QSE roots and current transition weights.

For open boundaries, a uniform Peierls vector potential still defines a
finite-chain drive by placing the same phase convention on the hopping
bonds. Translation symmetry is absent, so periodic plane waves
\(e^{ikj}\) are probe shapes; conserved crystal-momentum labels belong
to periodic translation symmetry. Open-chain reporting can keep
real-space bond-current correlators, or use standing-wave form factors
\(f_{m}(j)\propto\sin(k_mj)\) with \(k_m=m\pi/(L+1)\), exactly as in the
density and phonon channels of Section 18.5. Periodic \(k\)-resolved
optical plots are the translation-symmetric special case.

For a concrete nearest-neighbor chain, take the oriented Peierls
convention \[
\begin{aligned}
H_t(A)
&=
-t\sum_{i,\sigma}
\left(
e^{-iA}c_{i+1,\sigma}^\dagger c_{i\sigma}
+
e^{+iA}c_{i\sigma}^\dagger c_{i+1,\sigma}
\right),\\
T_{i\sigma}
&:=c_{i+1,\sigma}^\dagger c_{i\sigma},
\qquad
T_{i\sigma}^\dagger=c_{i\sigma}^\dagger c_{i+1,\sigma},\\
H_t(A)
&=
-t\sum_{i,\sigma}
\left(e^{-iA}T_{i\sigma}+e^{+iA}T_{i\sigma}^\dagger\right).
\end{aligned}
\] Then \[
\begin{aligned}
\frac{\partial H_t}{\partial A}
&=
-t\sum_{i,\sigma}
\left(
(-i)e^{-iA}T_{i\sigma}
+
(i)e^{+iA}T_{i\sigma}^\dagger
\right)\\
&=
it\sum_{i,\sigma}
\left(
e^{-iA}T_{i\sigma}
-
e^{+iA}T_{i\sigma}^\dagger
\right),\\
J
&=
-\left.\frac{\partial H_t}{\partial A}\right|_{A=0}
=
it\sum_{i,\sigma}
\left(T_{i\sigma}^\dagger-T_{i\sigma}\right)\\
&=
it\sum_{i,\sigma}
\left(
c_{i\sigma}^\dagger c_{i+1,\sigma}
-
c_{i+1,\sigma}^\dagger c_{i\sigma}
\right).
\end{aligned}
\] The Hermiticity check is immediate from the antisymmetric hopping
combination: \[
\begin{aligned}
J^\dagger
&=
-it\sum_{i,\sigma}
\left(
c_{i+1,\sigma}^\dagger c_{i\sigma}
-
c_{i\sigma}^\dagger c_{i+1,\sigma}
\right)
=J .
\end{aligned}
\] The second derivative gives the diamagnetic/contact operator: \[
\begin{aligned}
\frac{\partial^2H_t}{\partial A^2}
&=
-t\sum_{i,\sigma}
\left(
(-i)(-i)e^{-iA}T_{i\sigma}
+
(i)(i)e^{+iA}T_{i\sigma}^\dagger
\right)\\
&=
t\sum_{i,\sigma}
\left(
e^{-iA}T_{i\sigma}
+
e^{+iA}T_{i\sigma}^\dagger
\right),\\
D
&=
\left.\frac{\partial^2H_t}{\partial A^2}\right|_{A=0}
=
t\sum_{i,\sigma}
\left(T_{i\sigma}+T_{i\sigma}^\dagger\right)\\
&=
t\sum_{i,\sigma}
\left(
c_{i+1,\sigma}^\dagger c_{i\sigma}
+
c_{i\sigma}^\dagger c_{i+1,\sigma}
\right).
\end{aligned}
\] For this hopping-only Peierls convention,
\(H_t(0)=-t\sum_{i,\sigma}(T_{i\sigma}+T_{i\sigma}^\dagger)\), hence
\(D=-H_t(0)\). This equality belongs to this coupling convention.

In Jordan--Wigner qubit form, for modes \(p<q\), \[
\begin{aligned}
c_p^\dagger c_q+c_q^\dagger c_p
&=
\frac12
\left(
X_pZ_{p+1}\cdots Z_{q-1}X_q
+
Y_pZ_{p+1}\cdots Z_{q-1}Y_q
\right),\\
i(c_p^\dagger c_q-c_q^\dagger c_p)
&=
\frac12
\left(
Y_pZ_{p+1}\cdots Z_{q-1}X_q
-
X_pZ_{p+1}\cdots Z_{q-1}Y_q
\right).
\end{aligned}
\] Thus a bond current and its matching curvature term are \[
\begin{aligned}
J_{pq}
&=
t\,i(c_p^\dagger c_q-c_q^\dagger c_p)\\
&=
\frac{t}{2}
\left(
Y_pZ_{p+1}\cdots Z_{q-1}X_q
-
X_pZ_{p+1}\cdots Z_{q-1}Y_q
\right),\\
D_{pq}
&=
t(c_p^\dagger c_q+c_q^\dagger c_p)\\
&=
\frac{t}{2}
\left(
X_pZ_{p+1}\cdots Z_{q-1}X_q
+
Y_pZ_{p+1}\cdots Z_{q-1}Y_q
\right).
\end{aligned}
\] For adjacent modes the \(Z\)-string is empty: \[
J_{p,p+1}=\frac{t}{2}(Y_pX_{p+1}-X_pY_{p+1}),
\qquad
D_{p,p+1}=\frac{t}{2}(X_pX_{p+1}+Y_pY_{p+1}).
\]

Now write the spectral current response. Let \[
H_0|\Psi_\nu\rangle=E_\nu|\Psi_\nu\rangle,\qquad
H_0|\psi_0\rangle=E_0|\psi_0\rangle,\qquad
\omega_\nu:=E_\nu-E_0,
\] and define \(t_\nu^J:=\langle\Psi_\nu|J|\psi_0\rangle\). Since
\(J=J^\dagger\), \((t_\nu^J)^*=\langle\psi_0|J|\Psi_\nu\rangle\). Then
\[
\begin{aligned}
\langle J(t)J(0)\rangle_0
&=
\sum_\nu |t_\nu^J|^2e^{-i\omega_\nu t},\\
\langle J(0)J(t)\rangle_0
&=
\sum_\nu |t_\nu^J|^2e^{+i\omega_\nu t},\\
\chi_{JJ}^R(t)
&=
-i\Theta(t)
\sum_\nu |t_\nu^J|^2
\left(e^{-i\omega_\nu t}-e^{+i\omega_\nu t}\right).
\end{aligned}
\] Fourier transforming with the same \(e^{-\eta t}\) convergence factor
gives \[
\begin{aligned}
-i\int_0^\infty dt\,
e^{i\omega t}e^{-i\omega_\nu t}e^{-\eta t}
&=
\frac{1}{\omega-\omega_\nu+i\eta},\\
+i\int_0^\infty dt\,
e^{i\omega t}e^{+i\omega_\nu t}e^{-\eta t}
&=
-\frac{1}{\omega+\omega_\nu+i\eta},
\end{aligned}
\] hence \[
\chi_{JJ}^R(\omega)
=
\sum_\nu |t_\nu^J|^2
\left[
\frac{1}{\omega-\omega_\nu+i\eta}
-
\frac{1}{\omega+\omega_\nu+i\eta}
\right].
\]

QSE now supplies approximate neutral roots and current transition
amplitudes. With \[
\Phi_{\mathcal A}
=
\begin{bmatrix}
|\phi_1\rangle&\cdots&|\phi_N\rangle
\end{bmatrix},
\qquad
|\Psi_\nu^{\rm QSE}\rangle=\Phi_{\mathcal A}c^{(\nu)},
\qquad
d_a^J:=\langle\phi_a|J|\psi_0\rangle,
\] the current transition amplitude is \[
\begin{aligned}
t_\nu^{J,\mathcal A}
&:=
\langle\Psi_\nu^{\rm QSE}|J|\psi_0\rangle\\
&=
\left(\Phi_{\mathcal A}c^{(\nu)}\right)^\dagger J|\psi_0\rangle
=(c^{(\nu)})^\dagger d^J .
\end{aligned}
\] If \(c^{(\nu)}=W_\tau y^{(\nu)}\), then \[
t_\nu^{J,\mathcal A}
=(y^{(\nu)})^\dagger W_\tau^\dagger d^J .
\] Therefore \[
\begin{aligned}
\chi_{JJ}^{R,\mathcal A}(\omega)
&=
\sum_{\nu\in\mathcal W_{\rm targ}}
|t_\nu^{J,\mathcal A}|^2
\left[
\frac{1}{\omega-\omega_\nu^{\rm QSE}+i\eta}
-
\frac{1}{\omega+\omega_\nu^{\rm QSE}+i\eta}
\right],\\
\sigma_{\mathcal A}(\omega)
&=
\frac{i}{L(\omega+i\eta)}
\left(
\langle\psi_0|D|\psi_0\rangle
+
\chi_{JJ}^{R,\mathcal A}(\omega)
\right).
\end{aligned}
\] The QSE eigensolve is the neutral-sector QSE eigensolve; the probe
vector is \(d^J=\Phi_{\mathcal A}^\dagger J|\psi_0\rangle\); the extra
optical pieces are \(D\) and the division by \(E(\omega)\).

For absorption, use the distribution identity \[
\lim_{\eta\downarrow0}\frac{1}{x+i\eta}
=
\operatorname{PV}\frac1x-i\pi\delta(x).
\] Here \(\operatorname{PV}\) is the Cauchy principal value: the
singular integral is defined by excluding a small symmetric interval
around \(x=0\) and taking the exclusion width to zero. Thus \[
\operatorname{Im}\chi_{JJ}^R(\omega)
=
-\pi\sum_\nu |t_\nu^J|^2
\left[\delta(\omega-\omega_\nu)-\delta(\omega+\omega_\nu)\right].
\] For \(\omega>0\) and \(\omega_\nu>0\), the positive-frequency
absorption comes from \(\delta(\omega-\omega_\nu)\). Define the
temporary optical kernel \[
\mathcal K_{\rm opt}(\omega)
:=
\langle D\rangle_0+\chi_{JJ}^R(\omega)
=
\mathcal K_R(\omega)+i\mathcal K_I(\omega).
\] This \(\mathcal K_{\rm opt}\) is temporary shorthand for the
bracketed conductivity numerator. Then, away from the zero-frequency
contact singularity, \[
\begin{aligned}
\sigma(\omega)
&=
\frac{i}{L\omega}\mathcal K_{\rm opt}(\omega)
=
\frac{i\mathcal K_R(\omega)-\mathcal K_I(\omega)}{L\omega},\\
\operatorname{Re}\sigma(\omega)
&=
-\frac{\operatorname{Im}\chi_{JJ}^R(\omega)}{L\omega}.
\end{aligned}
\] Therefore \[
\operatorname{Re}\sigma(\omega>0)
=
\frac{\pi}{L}
\sum_\nu
\frac{|t_\nu^J|^2}{\omega_\nu}
\delta(\omega-\omega_\nu),
\] and the broadened QSE version is \[
\operatorname{Re}\sigma_{\mathcal A}(\omega>0)
\approx
\frac{\pi}{L}
\sum_{\nu\in\mathcal W_{\rm targ}}
\frac{|t_\nu^{J,\mathcal A}|^2}{\omega_\nu^{\rm QSE}}
L_\eta(\omega-\omega_\nu^{\rm QSE}),
\] with the precise prefactor tied to the broadening and conductivity
convention.

The zero-frequency term is the Drude/contact issue. If \[
\mathcal K_{\rm opt}(0)=\langle D\rangle_0+\chi_{JJ}^R(0),
\] then \[
\begin{aligned}
\sigma(\omega)\text{ contains }
&\frac{i\mathcal K_{\rm opt}(0)}{L(\omega+i0^+)}\\
&=
\frac{i\mathcal K_{\rm opt}(0)}{L}
\operatorname{PV}\frac1\omega
+
\frac{\pi\mathcal K_{\rm opt}(0)}{L}\delta(\omega).
\end{aligned}
\] Thus \(\operatorname{Re}\sigma(\omega)\) contains a possible
\(\delta(\omega)\) contribution with weight
\(\pi\mathcal K_{\rm opt}(0)/L\). A conductivity table must state how
this Drude/contact contribution is treated.

The implementation contract collected from this section is \[
\begin{aligned}
H(A_{\rm ext})
&=
H_0-A_{\rm ext}J+\frac12A_{\rm ext}^2D+\cdots,\\
J
&=
-\left.\partial_{A_{\rm ext}}H\right|_0,
\qquad
D=
\left.\partial_{A_{\rm ext}}^2H\right|_0,\\
\chi_{JJ}^R(t)
&=
-i\Theta(t)\langle[J(t),J(0)]\rangle_0,\\
t_\nu^{J,\mathcal A}
&=
\langle\Psi_\nu^{\rm QSE}|J|\psi_0\rangle
=(c^{(\nu)})^\dagger d^J
=(y^{(\nu)})^\dagger W_\tau^\dagger d^J,\\
\chi_{JJ}^{R,\mathcal A}(\omega)
&=
\sum_\nu |t_\nu^{J,\mathcal A}|^2
\left[
\frac1{\omega-\omega_\nu^{\rm QSE}+i\eta}
-
\frac1{\omega+\omega_\nu^{\rm QSE}+i\eta}
\right],\\
\sigma_{\mathcal A}(\omega)
&=
\frac{i}{L(\omega+i\eta)}
\left(
\langle D\rangle_0+\chi_{JJ}^{R,\mathcal A}(\omega)
\right).
\end{aligned}
\] Section 18.7 is the electromagnetic specialization of Section 18.4:
set the retarded probe to \(J\), derive \(J\) and \(D\) from the
Peierls/vector-potential convention, and divide the current response by
the electric field \(E=-\dot A\). Optical conductivity is meaningful
once the vector-potential convention, current operator, broadening,
system-size normalization, boundary convention, and Drude/contact
handling are all declared.

\subsection{18.8 Geometry-selected record
admission}\label{geometry-selected-record-admission}

\subsubsection{§18.8 lesson: geometry-selected record
admission}\label{lesson-geometry-selected-record-admission}

Interpretive compression: QSE record admission is subspace-control in
the pullback metric of the excitation synthesis map. A registry contains
possible excitation records \(r\), each carrying an encoded operator
\(\Omega_r\), but admission is decided by the physical ket produced
after that operator acts on the prepared seed. The candidate is useful
when the ket \(|\phi_r\rangle=Q_{\mathfrak s,0}\Omega_r|\psi_0\rangle\)
adds a new, measurable, well-conditioned, response-relevant direction to
the selected QSE subspace. Section 18.8 therefore treats admission as a
utility rule over geometry, probe visibility, Ritz residual repair, root
tracking, and resource cost.

The old selected records span \(\operatorname{span}\Phi_\ell\), the
candidate proposes \(|\phi_r\rangle\), and the old span is allowed to
explain the candidate through a nonorthogonal least-squares projection.
Least squares here means: \(|\phi_r\rangle\) and \(\Phi_\ell\) are fixed
data, \(\beta\) is the only fitted variable, and \(\Phi_\ell\beta\) is
the best old-span reconstruction of the candidate. The optimized
coefficient \(\beta^\star\) is substituted back into the squared error,
so the reported novelty is a Schur-complement residual length in the
supported Gram geometry.

At growth step \(\ell\), write the selected record set as \[
\mathcal A_\ell=\{a_1,\ldots,a_m\},\qquad m:=|\mathcal A_\ell|,
\] and collect the selected excitation kets into the synthesis map \[
\Phi_\ell:=\Phi_{\mathcal A_\ell}
=
\begin{bmatrix}
|\phi_{a_1}\rangle&|\phi_{a_2}\rangle&\cdots&|\phi_{a_m}\rangle
\end{bmatrix}.
\] The overlap metric of the old selected span is \[
S_\ell:=\Phi_\ell^\dagger\Phi_\ell,\qquad
(S_\ell)_{ij}=\langle\phi_{a_i}|\phi_{a_j}\rangle.
\] For a candidate record \(r\), \[
|\phi_r\rangle:=Q_{\mathfrak s,0}\Omega_r|\psi_0\rangle,\qquad
F_r:=\langle\phi_r|\phi_r\rangle,\qquad
s_r:=\Phi_\ell^\dagger|\phi_r\rangle,
\] so that \[
(s_r)_i=\langle\phi_{a_i}|\phi_r\rangle.
\] Here \(F_r\) is the candidate's physical squared length and \(s_r\)
is the vector of raw shadows cast by the candidate onto the old
excitation bras.

The old-span approximation has coefficient vector
\(\beta\in\mathbb C^m\): \[
\Phi_\ell\beta=\sum_{i=1}^m\beta_i|\phi_{a_i}\rangle.
\] Choose \(\beta\) by minimizing the physical squared error \[
\mathcal E_r(\beta):=\|\phi_r-\Phi_\ell\beta\|^2
=
\langle\phi_r-\Phi_\ell\beta|\phi_r-\Phi_\ell\beta\rangle.
\] Expanding all terms gives \[
\begin{aligned}
\mathcal E_r(\beta)
&=
\langle\phi_r|\phi_r\rangle
-\langle\phi_r|\Phi_\ell\beta\rangle
-\langle\Phi_\ell\beta|\phi_r\rangle
+\langle\Phi_\ell\beta|\Phi_\ell\beta\rangle\\
&=
F_r-s_r^\dagger\beta-\beta^\dagger s_r+\beta^\dagger S_\ell\beta\\
&=
F_r-2\operatorname{Re}(\beta^\dagger s_r)+\beta^\dagger S_\ell\beta .
\end{aligned}
\] In components, \[
\mathcal E_r(\beta)
=
F_r-\sum_{i=1}^m\beta_i^\ast(s_r)_i
-\sum_{i=1}^m(s_r)_i^\ast\beta_i
+\sum_{i,j=1}^m\beta_i^\ast(S_\ell)_{ij}\beta_j .
\] Treating \(\beta_i\) and \(\beta_i^\ast\) as independent variables in
the usual complex quadratic calculus, \[
\frac{\partial\mathcal E_r}{\partial \beta_k^\ast}
=
-(s_r)_k+\sum_{j=1}^m(S_\ell)_{kj}\beta_j
=
(S_\ell\beta)_k-(s_r)_k.
\] Stationarity for every \(k\) gives the normal equation \[
S_\ell\beta^\star=s_r.
\] If \(S_\ell\) has redundant or nearly null directions, the supported
inverse from §18.2 converts raw overlaps into old-span coefficients only
on the trusted metric support: \[
\beta^\star=S_{\ell,\tau}^{+}s_r.
\] The fitted old-span ket and the residualized candidate ket are
therefore \[
|\phi_r^{\rm fit}\rangle:=\Phi_\ell\beta^\star,\qquad
|u_r\rangle:=|\phi_r\rangle-|\phi_r^{\rm fit}\rangle
=|\phi_r\rangle-\Phi_\ell S_{\ell,\tau}^{+}s_r.
\] Thus \(|u_r\rangle\) is the part of the candidate left after the old
selected records have explained the candidate as well as the retained
nonorthogonal metric allows.

Substitute \(\beta^\star\) back into the objective: \[
\begin{aligned}
\langle u_r|u_r\rangle
&=
\mathcal E_r(\beta^\star)\\
&=
F_r-(\beta^\star)^\dagger s_r-s_r^\dagger\beta^\star
+(\beta^\star)^\dagger S_\ell\beta^\star .
\end{aligned}
\] On the retained support \(S_\ell\beta^\star=s_r\), hence \[
(\beta^\star)^\dagger S_\ell\beta^\star=(\beta^\star)^\dagger s_r
\] and \[
\langle u_r|u_r\rangle
=F_r-s_r^\dagger\beta^\star
=F_r-s_r^\dagger S_{\ell,\tau}^{+}s_r.
\] The normalized novelty is the stabilized residual fraction \[
\boxed{
\mathcal N_{\rm QSE}(r;\ell)
:=
\frac{\langle u_r|u_r\rangle}{F_r+\epsilon_F}
=
\frac{F_r-s_r^\dagger S_{\ell,\tau}^{+}s_r}{F_r+\epsilon_F}.
}
\] The coefficient \(\beta^\star\) has disappeared because it has been
optimized out. Large \(\mathcal N_{\rm QSE}\) means the candidate
contributes outside-span physical length; small \(\mathcal N_{\rm QSE}\)
means the retained old span already reconstructs the candidate.

The same expression is the Schur-complement geometry of the augmented
overlap
matrix.\footnote{General Schur complement: for a block matrix $A=\left(\begin{smallmatrix}B&C\\D&E\end{smallmatrix}\right)$ with invertible $B$, the Schur complement of $B$ in $A$ is $E-DB^{-1}C$; it is the lower-right block left after the upper-left variables have been eliminated. For a Gram matrix built from old vectors and one candidate, this leftover equals the squared norm of the candidate after best projection against the old vectors.}
Appending the raw candidate gives \[
\Phi_{\ell+1}=\begin{bmatrix}\Phi_\ell&|\phi_r\rangle\end{bmatrix},
\qquad
S_{\ell+1}
=
\Phi_{\ell+1}^\dagger\Phi_{\ell+1}
=
\begin{pmatrix}
S_\ell&s_r\\
s_r^\dagger&F_r
\end{pmatrix}.
\] With an exact nonsingular \(S_\ell\), the Schur complement of the old
block is \[
F_r-s_r^\dagger S_\ell^{-1}s_r
=
\min_\beta\|\phi_r-\Phi_\ell\beta\|^2
=
\langle u_r|u_r\rangle.
\] The supported version replaces \(S_\ell^{-1}\) by
\(S_{\ell,\tau}^{+}\). The Schur complement is therefore the same
leftover length computed from the enlarged Gram matrix.

Rank counts independent physical directions. For the Hermitian Gram
matrix \[
S_\ell=U\operatorname{diag}(s_1,\ldots,s_m)U^\dagger,\qquad s_j\ge0,
\] the exact rank is \[
\operatorname{rank}(S_\ell):=\#\{j:s_j>0\}.
\] Because \(c^\dagger S_\ell c=\|\Phi_\ell c\|^2\), a zero eigenvalue
means that some nonzero coefficient direction synthesizes the zero
physical ket. The thresholded retained rank is \[
J_\tau(S_\ell):=\{j:s_j>\tau_S\},\qquad
\operatorname{rank}_\tau(S_\ell):=|J_\tau(S_\ell)|.
\] This is the number of selected QSE directions whose physical length
clears the support floor. The rank gate in the admission rule asks
whether the candidate creates at least one additional retained physical
direction: \[
\operatorname{rank}_\tau(S_{\ell+1})>\operatorname{rank}_\tau(S_\ell).
\] In the ideal positive-definite case,
\(F_r-s_r^\dagger S_\ell^{-1}s_r>0\) increases the rank by one; in the
thresholded case, the residual direction must be strong enough to
produce a new eigenvalue above \(\tau_S\).

Conditioning measures how violently the retained metric inverse can
amplify finite-shot or roundoff perturbations. With
\(J_\tau=J_\tau(S)\), \[
\kappa_\tau(S)
:=
\frac{\max_{j\in J_\tau}s_j}{\min_{j\in J_\tau}s_j},
\] with the solve marked unsupported when \(J_\tau=\varnothing\). If
\(s=s_{\rm true}+\delta s\), then \[
\delta\beta
=S_\tau^+\delta s
=
\sum_{j\in J_\tau}
u_j\frac{u_j^\dagger\delta s}{s_j},
\qquad
\|\delta\beta\|
\le
\frac{\|\delta s\|}{\min_{j\in J_\tau}s_j}.
\] Thus a tiny retained eigenvalue makes the least-squares coefficients
and whitened coordinates sensitive to small measurement errors. The
condition gate \[
\kappa_\tau(S_{\ell+1})\le\kappa_{\max}
\] keeps the enlarged QSE metric inside the declared numerical trust
region.

Novelty alone can select a direction that is new but invisible to the
spectra being reported. Let the reported response probes be \[
\mathcal O_{\rm probe}
=
\{n_k,X_k,P_k,J,c_p,c_q^\dagger,\ldots\},
\] with the actual set chosen from the density, phonon, current, and
Green-function channels in §§18.3--18.7. Candidate probe visibility is
\[
\boxed{
\mathcal T_{\rm probe}(r;\ell)
:=
\sum_{O\in\mathcal O_{\rm probe}}w_O
\frac{|\langle\phi_r|O|\psi_0\rangle|^2}{F_r+\epsilon_F}.
}
\] The numerator is the candidate's pre-root brightness for probe \(O\).
Expanding through the record operator, \[
\begin{aligned}
\langle\phi_r|O|\psi_0\rangle
&=
(Q_{\mathfrak s,0}\Omega_r|\psi_0\rangle)^\dagger O|\psi_0\rangle\\
&=
\langle\psi_0|\Omega_r^\dagger Q_{\mathfrak s,0}^\dagger O|\psi_0\rangle\\
&=
\langle\psi_0|\Omega_r^\dagger Q_{\mathfrak s,0}O|\psi_0\rangle
\qquad(Q_{\mathfrak s,0}=Q_{\mathfrak s,0}^\dagger).
\end{aligned}
\] This is a seed-state transition matrix element, so it depends on the
prepared state, the filter, the candidate record, and the chosen
physical probe.

Residual gain measures whether the candidate repairs the approximate
eigen-equation. The current QSE roots satisfy \[
|\Psi_\nu^{(\ell)}\rangle=\Phi_\ell c_\nu^{(\ell)},\qquad
M_\ell c_\nu^{(\ell)}=\varepsilon_\nu^{(\ell)}S_\ell c_\nu^{(\ell)}.
\] The Hilbert-space residual is the vector left over after the full
Hamiltonian eigen-equation is tested on the synthesized QSE root: \[
\boxed{
|R_\nu^{(\ell)}\rangle
:=
(H_0-\varepsilon_\nu^{(\ell)})|\Psi_\nu^{(\ell)}\rangle .
}
\] The Hilbert-space residual is a ket in the physical state space, and
it vanishes for an exact eigenstate. QSE enforces the projected
equation: \[
\Phi_\ell^\dagger |R_\nu^{(\ell)}\rangle
=
\Phi_\ell^\dagger H_0\Phi_\ell c_\nu^{(\ell)}
-\varepsilon_\nu^{(\ell)}\Phi_\ell^\dagger\Phi_\ell c_\nu^{(\ell)}
=
M_\ell c_\nu^{(\ell)}-\varepsilon_\nu^{(\ell)}S_\ell c_\nu^{(\ell)}
=0.
\] So the selected basis sees no residual, while directions outside the
selected span can still see the missing part of the eigen-equation. This
is why the residualized candidate \(|u_r\rangle\) is the right object to
test: \[
\boxed{
\mathcal G_{\rm res}(r,\nu;\ell)
:=
\frac{|\langle u_r|R_\nu^{(\ell)}\rangle|^2}
{\langle u_r|u_r\rangle+\epsilon_F}.
}
\] For the normalized direction
\(|\widehat u_r\rangle=|u_r\rangle/\sqrt{\langle u_r|u_r\rangle}\), \[
\mathcal G_{\rm res}(r,\nu;\ell)
\approx
|\langle\widehat u_r|R_\nu^{(\ell)}\rangle|^2,
\] so this is the squared residual component pointing along the new
candidate direction.

For a target spectral window \(\mathcal W_{\rm targ}\), the scalar
residual score used in the utility can be aggregated as \[
\mathcal G_{\rm res}(r;\ell)
:=
\sum_{\nu\in\mathcal W_{\rm targ}}w_\nu^{\rm root}
\mathcal G_{\rm res}(r,\nu;\ell),
\] or restricted to a protected root when the benchmark is
root-specific.

The two-dimensional Ritz check asks the same question by actually
allowing the current root to mix with the new residualized direction.
Define \[
\mathcal Z_{r\nu}
:=
\operatorname{span}\{|\Psi_\nu^{(\ell)}\rangle,|u_r\rangle\},
\qquad
|z_1\rangle:=|\Psi_\nu^{(\ell)}\rangle,\quad
|z_2\rangle:=|u_r\rangle.
\] The local two-vector pencil is \[
S_{ij}^{(2)}:=\langle z_i|z_j\rangle,\qquad
M_{ij}^{(2)}:=\langle z_i|H_0|z_j\rangle,\qquad i,j\in\{1,2\},
\] and the local Ritz equation is \[
M^{(2)}a_j^{(2)}=\lambda_j^{(2)}S^{(2)}a_j^{(2)}.
\] The candidate's effect enters through the off-diagonal couplings
\(\langle\Psi_\nu^{(\ell)}|H_0|u_r\rangle\),
\(\langle u_r|H_0|\Psi_\nu^{(\ell)}\rangle\), and the new diagonal
\(\langle u_r|H_0|u_r\rangle\). If those couplings and the residualized
norm vanish, the two-dimensional solve returns the old root. If they are
significant, the matched local Ritz value moves.

One explicit matching convention is \[
j^\star
:=
\arg\max_{j\in\{1,2\}}
\frac{|e_1^\dagger S^{(2)}a_j^{(2)}|^2}
{(e_1^\dagger S^{(2)}e_1)\bigl((a_j^{(2)})^\dagger S^{(2)}a_j^{(2)}\bigr)},
\qquad
e_1=(1,0)^\top,
\] then \[
\lambda_{\nu,r}^{(2)}:=\lambda_{j^\star}^{(2)},\qquad
\Delta_{\rm Ritz}(r,\nu;\ell)
:=
|\lambda_{\nu,r}^{(2)}-\varepsilon_\nu^{(\ell)}|.
\] For a lowest-root repair target, a signed lowering score such as
\([\varepsilon_\nu^{(\ell)}-\lambda_{\nu,r}^{(2)}]_+\) may be used. For
a spectral-window target, the absolute matched movement or a
residual-reduction score is cleaner. In all cases, the comparison is
between the old one-root QSE value and the local Ritz value obtained
after letting that root mix with the candidate.

The main candidate questions are complementary: \[
\mathcal N_{\rm QSE}(r;\ell)
\quad\text{asks for a new metric direction,}
\] \[
\mathcal T_{\rm probe}(r;\ell)
\quad\text{asks for visibility in reported response channels,}
\] \[
\mathcal G_{\rm res}(r;\ell)
\quad\text{asks for repair of current Ritz residuals,}
\] \[
\Delta_{\rm Ritz}(r;\ell)
\quad\text{asks for direct local spectral movement after mixing.}
\] These scores answer different questions. A new direction can be dark
for the chosen probe; a bright candidate can be nearly redundant; a
residual-repairing candidate can create an unstable metric. The
admission law keeps these roles separate.

The acquisition score is a normalized controller utility: \[
\begin{aligned}
\mathcal S_{\rm QSE}(r;\ell)
&=
w_T\widehat{\mathcal T}_{\rm probe}(r;\ell)
+w_N\widehat{\mathcal N}_{\rm QSE}(r;\ell)
+w_R\widehat{\mathcal G}_{\rm res}(r;\ell)
+w_Z\widehat{\Delta}_{\rm Ritz}(r;\ell)\\
&\quad
-w_\kappa\widehat{\mathcal K}_{\rm cond}(r;\ell)
-w_M\widehat{\mathcal C}_{\rm mat}(r;\ell)\\
&\quad
-w_{\rm 2q}\widehat{\mathcal C}_{\rm 2q}^{\rm inc}(r;\ell)
-w_d\widehat{\mathcal D}_{\rm circ}^{\rm inc}(r;\ell)
-w_\theta\widehat{\Delta K}^{\rm inc}(r;\ell).
\end{aligned}
\] The hats mean dimensionless normalized scores. The first line rewards
response visibility, novelty, residual repair, and local Ritz movement.
The second line penalizes metric conditioning damage and the new
matrix-measurement burden expanded in §18.9. The third line carries the
Paper-I hardware-cost ingredients into QSE record admission: incremental
two-qubit cost \(\mathcal C_{\rm 2q}^{\rm inc}\), incremental circuit
depth \(\mathcal D_{\rm circ}^{\rm inc}\), and incremental
parameter/runtime-term burden \(\Delta K^{\rm inc}\). A compact hardware
proxy may combine them as \[
\mathcal C_{\rm hw}^{\rm inc}
:=
\lambda_{\rm 2q}\widehat{\mathcal C}_{\rm 2q}^{\rm inc}
+\lambda_d\widehat{\mathcal D}_{\rm circ}^{\rm inc}
+\lambda_\theta\widehat{\Delta K}^{\rm inc},
\] but the expanded score shows which hardware burden is being charged.

The final admission predicate requires four simultaneous gates: \[
\boxed{
\begin{aligned}
\operatorname{Admit}_{\rm QSE}(r;\ell)
&=
\mathbf 1\!\left[
\mathcal S_{\rm QSE}(r;\ell)\ge\tau_{\rm adm}
\right]\,
\mathbf 1\!\left[
\operatorname{rank}_\tau(S_{\ell+1})>\operatorname{rank}_\tau(S_\ell)
\right]\\
&\quad\times
\mathbf 1\!\left[
\kappa_\tau(S_{\ell+1})\le\kappa_{\max}
\right]\,
\mathbf 1\!\left[
\operatorname{TrackStable}_{\mathcal W_{\rm targ}}(r;\ell)=1
\right].
\end{aligned}
}
\] The score gate requires enough net utility. The retained-rank gate
requires a new supported physical direction. The condition gate requires
a usable nonorthogonal metric. The tracking gate requires stable root
identity over the target spectral window.

Root tracking is a spectral-continuity predicate. It is evaluated after
temporarily appending \(r\), rebuilding \(S_{\ell+1},M_{\ell+1}\), and
resolving the QSE pencil. Let the old roots be
\(\{|\Psi_\nu^{(\ell)}\rangle\}_{\nu\in\mathcal W_{\rm targ}}\) and the
trial new roots be \(\{|\Psi_\mu^{(\ell+1,r)}\rangle\}\). A typical
tracking overlap is \[
\mathcal O_{\nu\mu}^{\rm root}
:=
\frac{|\langle\Psi_\nu^{(\ell)}|\Psi_\mu^{(\ell+1,r)}\rangle|^2}
{\bigl(\langle\Psi_\nu^{(\ell)}|\Psi_\nu^{(\ell)}\rangle+\epsilon_{\rm tr}\bigr)
\bigl(\langle\Psi_\mu^{(\ell+1,r)}|\Psi_\mu^{(\ell+1,r)}\rangle+\epsilon_{\rm tr}\bigr)}.
\] This is the cosine-normalized squared overlap; for normalized roots
it reduces to
\(|\langle\Psi_\nu^{(\ell)}|\Psi_\mu^{(\ell+1,r)}\rangle|^2\) up to the
small stabilizer. In coefficient form this numerator is \[
\langle\Psi_\nu^{(\ell)}|\Psi_\mu^{(\ell+1,r)}\rangle
=
(c_\nu^{(\ell)})^\dagger
\Phi_\ell^\dagger\Phi_{\ell+1,r}
c_\mu^{(\ell+1,r)}.
\] The tracking predicate can then require a one-to-one assignment
\(\pi\) on the target window such that \[
\mathcal O_{\nu,\pi(\nu)}^{\rm root}\ge\tau_{\rm ov},\qquad
|\omega_{\pi(\nu)}^{(\ell+1,r)}-\omega_\nu^{(\ell)}|\le\tau_\omega,
\] and, for the probes being reported, \[
|t_{\pi(\nu)}^{O,(\ell+1,r)}-t_\nu^{O,(\ell)}|\le\tau_t(O)
\qquad
\forall O\in\mathcal O_{\rm probe}.
\] Near a genuine degeneracy, the stable object is the subspace spanned
by the near-degenerate roots, so the same rule can track a small
invariant subspace. Commutation information belongs to structural
feasibility and batching; \(\operatorname{TrackStable}\) belongs to
eigenvalue/eigenvector continuity after the QSE pencil is solved.

Exact diagonalization may be used as a diagnostic benchmark for reported
errors, spectral overlays, and post hoc root labels. QPU-facing record
admission uses only measurement-compatible objects built from the
prepared seed, candidate records, projected matrix elements, transition
probes, conditioning estimates, root-tracking telemetry, and resource
proxies: \[
|\psi_0\rangle,\quad
\Omega_r,\quad
|\phi_r\rangle,\quad
S_\ell,\quad
M_\ell,\quad
d^O,\quad
\mathcal O_{\rm probe},\quad
\mathcal C_{\rm mat},\quad
\mathcal C_{\rm 2q},\quad
\mathcal D_{\rm circ},\quad
\Delta K.
\] The selection remains causal because exact target roots label the
error after the selected manifold has been built.

The implementation contract collected from this section is \[
\boxed{
|\phi_r\rangle=Q_{\mathfrak s,0}\Omega_r|\psi_0\rangle
}
\quad\text{candidate physical ket,}
\] \[
\boxed{
\beta^\star=S_{\ell,\tau}^{+}\Phi_\ell^\dagger|\phi_r\rangle
}
\quad\text{best old-span explanation,}
\] \[
\boxed{
|u_r\rangle=|\phi_r\rangle-\Phi_\ell\beta^\star
}
\quad\text{outside-span leftover,}
\] \[
\boxed{
\mathcal N_{\rm QSE}
=
\frac{\|u_r\|^2}{F_r+\epsilon_F}
}
\quad\text{new supported geometry,}
\] \[
\boxed{
\mathcal T_{\rm probe}
=
\sum_Ow_O
\frac{|\langle\phi_r|O|\psi_0\rangle|^2}{F_r+\epsilon_F}
}
\quad\text{pre-root response visibility,}
\] \[
\boxed{
\mathcal G_{\rm res}
=
\sum_{\nu\in\mathcal W_{\rm targ}}w_\nu^{\rm root}
\frac{|\langle u_r|R_\nu^{(\ell)}\rangle|^2}
{\|u_r\|^2+\epsilon_F}
}
\quad\text{Ritz residual repair,}
\] \[
\boxed{
\operatorname{Admit}_{\rm QSE}
=
1[\text{score}]
1[\text{rank}]
1[\text{conditioning}]
1[\text{tracking}]
}
\quad\text{Boolean record admission.}
\]

For reading: §18.8 is the point where QSE becomes a controlled
spectral-manifold construction. A candidate record is admitted when it
adds supported geometry, is relevant to the response channels being
claimed, repairs current Ritz misses, preserves trackable roots, and
stays inside the declared measurement and circuit budget.

\subsection{18.9 Matrix elements and measurement
burden}\label{matrix-elements-and-measurement-burden}

\subsubsection{§18.9 lesson: matrix elements and measurement
burden}\label{lesson-matrix-elements-and-measurement-burden}

Interpretive compression: Section 18.9 is the operational closure of the
QSE formalism. Every abstract QSE object from §§18.1--18.8 becomes a
measured scalar or a classical contraction of measured scalars. The
overlap matrix \(S\), Hamiltonian matrix \(M\), transition vector
\(d^O\), response weights, Green-function residues, and optical-current
amplitudes are built from matrix elements. After \(H_0\), the excitation
records \(\Omega_a\), the filter \(Q_{\mathfrak s,0}\), and the probes
\(O\) are encoded as Pauli sums, those matrix elements reduce to
expectation values on the same prepared seed state \(|\psi_0\rangle\).
QSE is a Ritz method, and a QPU-facing QSE calculation pays for the
matrix entries defining the Ritz space.

The shorthand formulas are \[
S_{ab}=\langle\phi_a|\phi_b\rangle,\qquad
M_{ab}=\langle\phi_a|H_0|\phi_b\rangle,\qquad
d_a^O=\langle\phi_a|O|\psi_0\rangle.
\] Now unfold \(|\phi_a\rangle\). Let \[
Q:=Q_{\mathfrak s,0},\qquad
|\phi_a\rangle:=Q\Omega_a|\psi_0\rangle,\qquad
Q^\dagger=Q,\qquad Q^2=Q.
\] Then \[
\langle\phi_a|
=
(Q\Omega_a|\psi_0\rangle)^\dagger
=
\langle\psi_0|\Omega_a^\dagger Q,
\] so the abstract QSE entries are really \[
\begin{aligned}
S_{ab}
&=\langle\psi_0|\Omega_a^\dagger Q\Omega_b|\psi_0\rangle,\\
M_{ab}
&=\langle\psi_0|\Omega_a^\dagger QH_0Q\Omega_b|\psi_0\rangle,\\
d_a^O
&=\langle\psi_0|\Omega_a^\dagger QO|\psi_0\rangle.
\end{aligned}
\] The geometric ket notation is clean; the seed-expectation notation is
operational.

Use mnemonic Pauli labels in place of generic Greek indices: \[
\begin{aligned}
H_0
&=\sum_{\lambda_H\in\mathcal I_H}h_{\lambda_H}P_{\lambda_H},\\
\Omega_a
&=\sum_{\lambda_{\Omega_a}\in\mathcal I_{\Omega_a}}
\omega_{a,\lambda_{\Omega_a}}P_{\lambda_{\Omega_a}},\\
O
&=\sum_{\lambda_{\mathcal O}\in\mathcal I_{\mathcal O}}
o_{\lambda_{\mathcal O}}P_{\lambda_{\mathcal O}},\\
Q
&=\sum_{\lambda_Q\in\mathcal I_Q}q_{\lambda_Q}P_{\lambda_Q}
\quad\text{when the filter/projector is expanded explicitly.}
\end{aligned}
\] Here \(\lambda_H\) labels Hamiltonian Pauli terms,
\(\lambda_{\Omega_a}\) labels terms in record \(a\),
\(\lambda_{\mathcal O}\) labels probe terms, and \(\lambda_Q\) labels
filter/projector terms. Pauli words are Hermitian and unitary,
\(P_\lambda^\dagger=P_\lambda\) and \(P_\lambda^2=I\), while
\(\Omega_a\) can carry complex coefficients: \[
\Omega_a^\dagger
=
\left(\sum_{\lambda_{\Omega_a}}\omega_{a,\lambda_{\Omega_a}}P_{\lambda_{\Omega_a}}\right)^\dagger
=
\sum_{\lambda_{\Omega_a}}\omega_{a,\lambda_{\Omega_a}}^\ast
P_{\lambda_{\Omega_a}}^\dagger .
\]

The overlap entry becomes \[
\begin{aligned}
S_{ab}
&=
\langle\psi_0|\Omega_a^\dagger Q\Omega_b|\psi_0\rangle\\
&=
\sum_{\lambda_{\Omega_a},\lambda_{\Omega_b}}
\omega_{a,\lambda_{\Omega_a}}^\ast
\omega_{b,\lambda_{\Omega_b}}
\langle\psi_0|
P_{\lambda_{\Omega_a}}^\dagger QP_{\lambda_{\Omega_b}}
|\psi_0\rangle\\
&=
\sum_{\lambda_{\Omega_a},\lambda_Q,\lambda_{\Omega_b}}
\omega_{a,\lambda_{\Omega_a}}^\ast
q_{\lambda_Q}
\omega_{b,\lambda_{\Omega_b}}
\langle\psi_0|
P_{\lambda_{\Omega_a}}^\dagger P_{\lambda_Q}P_{\lambda_{\Omega_b}}
|\psi_0\rangle .
\end{aligned}
\] The Hamiltonian entry becomes \[
\begin{aligned}
M_{ab}
&=
\langle\psi_0|\Omega_a^\dagger QH_0Q\Omega_b|\psi_0\rangle\\
&=
\sum_{\lambda_{\Omega_a},\lambda_H,\lambda_{\Omega_b}}
\omega_{a,\lambda_{\Omega_a}}^\ast
h_{\lambda_H}
\omega_{b,\lambda_{\Omega_b}}
\langle\psi_0|
P_{\lambda_{\Omega_a}}^\dagger QP_{\lambda_H}QP_{\lambda_{\Omega_b}}
|\psi_0\rangle\\
&=
\sum_{\lambda_{\Omega_a},\lambda_{Q,1},\lambda_H,\lambda_{Q,2},\lambda_{\Omega_b}}
\omega_{a,\lambda_{\Omega_a}}^\ast
q_{\lambda_{Q,1}}h_{\lambda_H}q_{\lambda_{Q,2}}
\omega_{b,\lambda_{\Omega_b}}
\langle\psi_0|
P_{\lambda_{\Omega_a}}^\dagger P_{\lambda_{Q,1}}P_{\lambda_H}P_{\lambda_{Q,2}}P_{\lambda_{\Omega_b}}
|\psi_0\rangle .
\end{aligned}
\] The probe transition entry becomes \[
\begin{aligned}
d_a^O
&=
\langle\psi_0|\Omega_a^\dagger QO|\psi_0\rangle\\
&=
\sum_{\lambda_{\Omega_a},\lambda_{\mathcal O}}
\omega_{a,\lambda_{\Omega_a}}^\ast
o_{\lambda_{\mathcal O}}
\langle\psi_0|
P_{\lambda_{\Omega_a}}^\dagger QP_{\lambda_{\mathcal O}}
|\psi_0\rangle\\
&=
\sum_{\lambda_{\Omega_a},\lambda_Q,\lambda_{\mathcal O}}
\omega_{a,\lambda_{\Omega_a}}^\ast
q_{\lambda_Q}o_{\lambda_{\mathcal O}}
\langle\psi_0|
P_{\lambda_{\Omega_a}}^\dagger P_{\lambda_Q}P_{\lambda_{\mathcal O}}
|\psi_0\rangle .
\end{aligned}
\] Thus \[
\boxed{
S_{ab},\ M_{ab},\ d_a^O
\quad\text{are seed-state Pauli-polynomial expectation values.}
}
\]

Pauli products stay inside the Pauli
basis.\footnote{Pauli-word multiplication: $\chi(A,B)$ is the accumulated phase from sitewise Pauli multiplication, and $A\oplus B$ is the label of the resulting Pauli word. If $P_A=X\otimes Z$ and $P_B=Y\otimes X$, then $P_AP_B=(X\otimes Z)(Y\otimes X)=(XY)\otimes(ZX)=(iZ)\otimes(iY)=-Z\otimes Y=\chi(A,B)P_{A\oplus B}$, so $\chi(A,B)=-1$ in this example. In general $\chi(A,B)\in\{\pm1,\pm i\}$.}
For two Pauli words, \[
P_AP_B=\chi(A,B)P_{A\oplus B}.
\] Therefore any product path such as \[
P_{\lambda_{\Omega_a}}^\dagger P_{\lambda_Q}P_{\lambda_{\Omega_b}}
\] collapses to \[
P_{\lambda_{\Omega_a}}^\dagger P_{\lambda_Q}P_{\lambda_{\Omega_b}}
=
\chi_{\lambda_{\Omega_a},\lambda_Q,\lambda_{\Omega_b}}
P_{\lambda_{\Omega_a}\lambda_Q\lambda_{\Omega_b}},
\] and \[
\langle\psi_0|P_{\lambda_{\Omega_a}}^\dagger P_{\lambda_Q}P_{\lambda_{\Omega_b}}|\psi_0\rangle
=
\chi_{\lambda_{\Omega_a},\lambda_Q,\lambda_{\Omega_b}}
\langle\psi_0|P_{\lambda_{\Omega_a}\lambda_Q\lambda_{\Omega_b}}|\psi_0\rangle.
\] So every QSE scalar entry has the ledger form \[
X=\sum_{\lambda}c_{X,\lambda}\langle\psi_0|P_\lambda|\psi_0\rangle.
\] If many algebraic paths produce the same final Pauli word, combine
them before measurement: \[
\sum_{\text{paths }r}c_r\langle P_{\lambda(r)}\rangle
=
\sum_\lambda
\left(\sum_{r:\lambda(r)=\lambda}c_r\right)\langle P_\lambda\rangle
=
\sum_\lambda C_\lambda\langle P_\lambda\rangle .
\] The complex phases are stored in the classical coefficients
\(C_\lambda\); the measured primitive is the expectation value of a
Hermitian Pauli word.

The seed-removal part of \(Q_{\mathfrak s,0}\) is also a concrete
expectation subtraction. If \[
Q_{\mathfrak s,0}
=
\Pi_{\mathfrak s}(I-|\psi_0\rangle\langle\psi_0|)\Pi_{\mathfrak s},
\] and the seed already lies in the sector, then
\(Q_{\mathfrak s,0}=\Pi_{\mathfrak s}-|\psi_0\rangle\langle\psi_0|\)
inside that sector, hence \[
\begin{aligned}
\langle\psi_0|AQB|\psi_0\rangle
&=
\langle\psi_0|A\Pi_{\mathfrak s}B|\psi_0\rangle
-\langle\psi_0|A|\psi_0\rangle\langle\psi_0|B|\psi_0\rangle .
\end{aligned}
\] This subtracts the seed-parallel component and assigns zero novelty
to pure ground-state overlap in \(\Omega_a|\psi_0\rangle\).

Particle-changing Green-function sectors use the same Pauli-expansion
logic with source states and sector projectors. For an addition source
\[
|\chi_q^+\rangle:=c_q^\dagger|\psi_0\rangle,\qquad
|\phi_a^+\rangle:=\Pi_+\Omega_a^+|\chi_q^+\rangle
=\Pi_+\Omega_a^+c_q^\dagger|\psi_0\rangle,
\] the addition overlap is \[
\begin{aligned}
S_{ab}^{+}
&=\langle\phi_a^+|\phi_b^+\rangle\\
&=
\langle\psi_0|
c_q(\Omega_a^+)^\dagger\Pi_+\Omega_b^+c_q^\dagger
|\psi_0\rangle .
\end{aligned}
\] If \[
c_q=\sum_{\lambda_-}\xi_{q,\lambda_-}^{(-)}P_{\lambda_-},\quad
c_q^\dagger=\sum_{\lambda_+}\xi_{q,\lambda_+}^{(+)}P_{\lambda_+},\quad
\Omega_a^+=\sum_{\lambda_{\Omega_a^+}}\omega^+_{a,\lambda_{\Omega_a^+}}P_{\lambda_{\Omega_a^+}},\quad
\Pi_+=\sum_{\lambda_{\Pi_+}}\pi^+_{\lambda_{\Pi_+}}P_{\lambda_{\Pi_+}},
\] then \[
\begin{aligned}
S_{ab}^{+}
&=
\sum_{\lambda_-,\lambda_{\Omega_a^+},\lambda_{\Pi_+},\lambda_{\Omega_b^+},\lambda_+}
\xi_{q,\lambda_-}^{(-)}
(\omega^+_{a,\lambda_{\Omega_a^+}})^\ast
\pi^+_{\lambda_{\Pi_+}}
\omega^+_{b,\lambda_{\Omega_b^+}}
\xi_{q,\lambda_+}^{(+)}
\\
&\qquad\qquad\times
\langle\psi_0|
P_{\lambda_-}P_{\lambda_{\Omega_a^+}}^\dagger
P_{\lambda_{\Pi_+}}P_{\lambda_{\Omega_b^+}}P_{\lambda_+}
|\psi_0\rangle .
\end{aligned}
\] Removal sectors replace \(c_q^\dagger|\psi_0\rangle\) by
\(c_p|\psi_0\rangle\) and \(\Pi_+\) by \(\Pi_-\). The source creation or
annihilation operator simply contributes extra Pauli words at the ends
of the seed-state expectation.

Hermitian structure halves the matrix-entry ledger. Since \[
S_{ab}=\langle\phi_a|\phi_b\rangle,\qquad
M_{ab}=\langle\phi_a|H_0|\phi_b\rangle,\qquad H_0=H_0^\dagger,
\] we have \[
S_{ba}=S_{ab}^\ast,\qquad
M_{ba}=M_{ab}^\ast.
\] An unordered matrix pair means the label \(\{a,b\}\), where \((a,b)\)
and \((b,a)\) are the same required QSE matrix element. For
\(N_{\mathcal A}:=|\mathcal A|\) retained records, only the diagonal
plus one triangle is needed: \[
(1,1),(1,2),\ldots,(1,N_{\mathcal A}),(2,2),\ldots,(N_{\mathcal A},N_{\mathcal A}).
\] The count is \[
\begin{aligned}
N_{\rm pair}(N_{\mathcal A})
&=
N_{\mathcal A}+(N_{\mathcal A}-1)+\cdots+1\\
&=
\boxed{\frac{N_{\mathcal A}(N_{\mathcal A}+1)}2}.
\end{aligned}
\] Admitting one more record changes this by \[
\Delta N_{\rm pair}
=
\frac{(N_{\mathcal A}+1)(N_{\mathcal A}+2)}2
-\frac{N_{\mathcal A}(N_{\mathcal A}+1)}2
=
N_{\mathcal A}+1.
\] Thus one neutral QSE admission adds \(N_{\mathcal A}+1\) new overlap
entries and \(N_{\mathcal A}+1\) new Hamiltonian entries. If
\(\mathcal O_{\rm probe}=\{O_1,\ldots,O_{N_O}\}\), then
\(N_{\mathcal A}N_O\) transition-vector entries exist and one new record
adds \(N_O\) more, so the scalar-entry increment is \[
\boxed{
\Delta N_{\rm scalar}
=
\underbrace{(N_{\mathcal A}+1)}_{\Delta S}
+\underbrace{(N_{\mathcal A}+1)}_{\Delta M}
+\underbrace{N_O}_{\Delta d}.
}
\]

Scalar-entry count is still an underestimate of measurement burden
because each scalar entry expands into Pauli expectations. Define
support sizes \[
L_{\Omega_a}:=|\mathcal I_{\Omega_a}|,\qquad
L_H:=|\mathcal I_H|,\qquad
L_{\mathcal O}:=|\mathcal I_{\mathcal O}|,\qquad
L_Q:=|\mathcal I_Q|.
\] The naive algebraic product counts obey \[
\#(S_{ab})\lesssim L_{\Omega_a}L_{\Omega_b}L_Q,\qquad
\#(M_{ab})\lesssim L_{\Omega_a}L_HL_{\Omega_b}L_Q^2,\qquad
\#(d_a^O)\lesssim L_{\Omega_a}L_{\mathcal O}L_Q.
\] These are upper counts before Pauli-word collisions, symmetry zeros,
reuse across entries, and commuting-group measurement.

Let \(\mathcal W(\mathcal A)\) be the distinct final Pauli words needed
after simplifying all requested \(S\), \(M\), transition vectors, and
response channels, and partition them into jointly measurable groups \[
\mathcal G_{\rm meas}(\mathcal A)=\{g_1,\ldots,g_G\}.
\] For each requested scalar \(X\in\mathcal X_{\rm req}\), \[
X=\sum_{\lambda\in\mathcal W(\mathcal A)}C_{X,\lambda}\langle P_\lambda\rangle
=
\sum_{g\in\mathcal G_{\rm meas}(\mathcal A)}
\sum_{\lambda\in g}C_{X,\lambda}\langle P_\lambda\rangle.
\] With \(N_g\) shots for group \(g\), the estimator is \[
\widehat X
=
\sum_g\sum_{\lambda\in g}C_{X,\lambda}\widehat{\langle P_\lambda\rangle}_g.
\] For one Pauli word, each shot returns \(\pm1\), so \[
\operatorname{Var}\!\left[\widehat{\langle P_\lambda\rangle}\right]
=
\frac{1-\langle P_\lambda\rangle^2}{N_g}
\le
\frac1{N_g}.
\] For a group vector, write \(C_{X,g}:=(C_{X,\lambda})_{\lambda\in g}\)
and let \(\widehat\Sigma_g\) be the empirical per-shot covariance matrix
of the jointly estimated Pauli words in group \(g\). Then \[
\operatorname{Var}(\widehat X)
\approx
\sum_g\frac{1}{N_g}C_{X,g}^\dagger\widehat\Sigma_g C_{X,g}.
\] This motivates an explicit normalized group difficulty/importance
weight: \[
\boxed{
\widehat v_g
:=
\frac{1}{v_{\rm ref}}
\sum_{X\in\mathcal X_{\rm req}}w_X\,
C_{X,g}^\dagger\widehat\Sigma_g C_{X,g},
}
\] where \(w_X\) prioritizes different scalar entries and
\(v_{\rm ref}>0\) makes the weight dimensionless. The
grouped-measurement burden proxy is then \[
\boxed{
S_{\rm mat}^{\rm proxy}(\mathcal A)
:=
\sum_{g\in\mathcal G_{\rm meas}(\mathcal A)}
N_g\,\widehat v_g.
}
\] This proxy records weighted measurement burden for comparing
selected-basis policies. A high proxy means the current encoding,
selected records, probes, grouping, and precision target are costly to
certify; a low proxy means the required entries are cheaper under the
current grouping.

The connection to §18.8 is direct. The QSE acquisition score penalizes
resource terms such as \[
-w_M\widehat{\mathcal C}_{\rm mat}(r;\ell)
-w_{\rm 2q}\widehat{\mathcal C}_{\rm 2q}^{\rm inc}(r;\ell)
-w_d\widehat{\mathcal D}_{\rm circ}^{\rm inc}(r;\ell),
\] because each admitted record adds new matrix entries, new Pauli
products, and possibly new circuit depth or two-qubit burden. ``Add
every excitation'' is therefore both a measurement policy and a
linear-algebra choice.

Finite-shot matrix noise explains why the ledger matters. Suppose the
ideal matrices are \(M,S\), while measurement gives \[
M\mapsto M+\delta M,\qquad
S\mapsto S+\delta S.
\] For an \(S\)-normalized root, \[
Mc_\nu=\varepsilon_\nu Sc_\nu,\qquad c_\nu^\dagger Sc_\nu=1.
\] The perturbed equation is \[
(M+\delta M)(c_\nu+\delta c_\nu)
=
(\varepsilon_\nu+\delta\varepsilon_\nu)(S+\delta S)(c_\nu+\delta c_\nu).
\] Keeping first-order terms and canceling
\(Mc_\nu=\varepsilon_\nu Sc_\nu\), \[
(M-\varepsilon_\nu S)\delta c_\nu
+(\delta M-\varepsilon_\nu\delta S)c_\nu
=
\delta\varepsilon_\nu S c_\nu.
\] The left-eigenvector identity comes from taking the adjoint of the
right generalized eigenproblem: \[
(Mc_\nu)^\dagger=(\varepsilon_\nu Sc_\nu)^\dagger
\Longrightarrow
c_\nu^\dagger M^\dagger=\varepsilon_\nu^\ast c_\nu^\dagger S^\dagger.
\] Since \(M=M^\dagger\), \(S=S^\dagger\), and
\(\varepsilon_\nu\in\mathbb R\), \[
c_\nu^\dagger M=\varepsilon_\nu c_\nu^\dagger S,\qquad
c_\nu^\dagger(M-\varepsilon_\nu S)=0.
\] Left-multiply the first-order equation by \(c_\nu^\dagger\): \[
\boxed{
\delta\varepsilon_\nu
=
c_\nu^\dagger(\delta M-\varepsilon_\nu\delta S)c_\nu.
}
\] Thus matrix-element noise directly perturbs the reported QSE root.
For a transition amplitude \[
t_\nu^O=(c^{(\nu)})^\dagger d^O,
\] finite-shot perturbations give, to first order, \[
\boxed{
\delta t_\nu^O
\approx
(\delta c^{(\nu)})^\dagger d^O
+(c^{(\nu)})^\dagger\delta d^O.
}
\] So transition amplitudes are affected both by noisy transition-vector
entries and by noisy root vectors caused by \(\delta M,\delta S\).

The practical use is: these formulas turn measurement burden into
spectral error budgets. If the target linewidth is \(\eta\), a root gap
is \(\delta_\nu\), or a table tolerance is \(\tau_\omega\), then the
estimated size of
\(c_\nu^\dagger(\delta M-\varepsilon_\nu\delta S)c_\nu\) tells whether
more shots, better grouping, stronger support cutoff, or record
rejection is needed before claiming a peak location. The transition
formula plays the same role for brightness: if \(\delta t_\nu^O\) is
comparable to \(t_\nu^O\), the bright/dark classification and response
intensity need more measurement support. This is the ``so what'': §18.9
supplies the measurement-facing uncertainty model used to set shot
budgets, penalize expensive records in §18.8, and decide whether a QSE
spectrum is trustworthy enough to report.

The implementation contract collected from this section is \[
\boxed{
|\phi_a\rangle=Q_{\mathfrak s,0}\Omega_a|\psi_0\rangle
}
\quad\text{geometric excitation ket,}
\] \[
\boxed{
S_{ab}=\langle\psi_0|\Omega_a^\dagger Q_{\mathfrak s,0}\Omega_b|\psi_0\rangle
}
\quad\text{overlap as a seed-state composite expectation,}
\] \[
\boxed{
M_{ab}=\langle\psi_0|\Omega_a^\dagger Q_{\mathfrak s,0}H_0Q_{\mathfrak s,0}\Omega_b|\psi_0\rangle
}
\quad\text{Hamiltonian matrix as a seed-state composite expectation,}
\] \[
\boxed{
d_a^O=\langle\psi_0|\Omega_a^\dagger Q_{\mathfrak s,0}O|\psi_0\rangle
}
\quad\text{probe shadow as a seed-state composite expectation,}
\] \[
\boxed{
H_0=\sum_{\lambda_H}h_{\lambda_H}P_{\lambda_H},\quad
\Omega_a=\sum_{\lambda_{\Omega_a}}\omega_{a,\lambda_{\Omega_a}}P_{\lambda_{\Omega_a}},\quad
O=\sum_{\lambda_{\mathcal O}}o_{\lambda_{\mathcal O}}P_{\lambda_{\mathcal O}}
}
\quad\text{Pauli expansions,}
\] \[
\boxed{
N_{\rm pair}(N_{\mathcal A})
=
\frac{N_{\mathcal A}(N_{\mathcal A}+1)}2,\qquad
\Delta N_{\rm pair}=N_{\mathcal A}+1
}
\quad\text{Hermitian matrix-entry growth,}
\] \[
\boxed{
S_{\rm mat}^{\rm proxy}(\mathcal A)
=
\sum_{g\in\mathcal G_{\rm meas}(\mathcal A)}
N_g\widehat v_g
}
\quad\text{grouped-measurement burden proxy.}
\]

For reading: §18.9 is where QSE becomes a measurement protocol. The
equations \(Mc=\varepsilon Sc\) and \(t_\nu^O=(c^{(\nu)})^\dagger d^O\)
are compact on paper, while hardware builds their ingredients from Pauli
expectations on \(|\psi_0\rangle\). A good spectral manifold must be
expressive, conditioned, probe-visible, and measurable within the
available Pauli-group and shot budget.

\subsection{18.10 Frozen QSE dynamics, adaptive QSE growth, and live
handoff}\label{frozen-qse-dynamics-adaptive-qse-growth-and-live-handoff}

\subsubsection{§18.10 lesson: frozen QSE dynamics, adaptive QSE growth,
and live
handoff}\label{lesson-frozen-qse-dynamics-adaptive-qse-growth-and-live-handoff}

After static QSE has selected an excitation manifold
\(\operatorname{span}\Phi_{\mathcal A}\), dynamics has three natural
regimes. Frozen QSE treats \(\Phi_{\mathcal A}\) as a fixed
mini-Hilbert-space and propagates only the coefficient vector. Adaptive
QSE growth keeps the same linear-subspace representation but admits new
records at checkpoints. Live handoff refits a selected QSE root into a
prepared variational circuit and then uses the McLachlan/checkpoint
machinery. The primitive choice is therefore: coefficient dynamics in a
measured linear subspace, basis growth inside that same QSE
representation, or refit into a nonlinear circuit manifold.

Interpretive compression: Section 18.10 is the dynamical extension of
the same subspace geometry. Frozen QSE uses the measured pullback metric
\(S\) to propagate a coefficient vector in a fixed linear manifold;
adaptive QSE updates the manifold when the projected Schrödinger drift
leaks outside the retained span; live handoff changes representation
class by refitting a Ritz vector into a nonlinear variational manifold.
The escape diagnostic is a variance decomposition: total instantaneous
\(H(t)\)-drift versus the portion representable by the current QSE
frame.

Physical semantics: static spectra summarize infinitesimal source
response through pole locations and residues, while driven dynamics asks
for the finite-time state trajectory under \(H(t)\). Frozen QSE assumes
the selected spectral frame contains the relevant driven ray motion, so
only \(z(t)\) evolves. Adaptive QSE growth treats a large escape
fraction as evidence that the present spectral frame lacks a drift
direction needed by the drive. Live handoff supplies variational-circuit
coordinates when the evolving object must be prepared as a circuit
state.

Let \[
\Phi_{\mathcal A}:=
\begin{bmatrix}|\phi_1\rangle&\cdots&|\phi_N\rangle\end{bmatrix},\qquad
S_{ab}:=\langle\phi_a|\phi_b\rangle,\qquad
S=\Phi_{\mathcal A}^\dagger\Phi_{\mathcal A}.
\] A frozen QSE coefficient vector \(z(t)\in\mathbb C^N\) synthesizes
the physical state \[
|\Psi(t)\rangle:=\Phi_{\mathcal A}z(t)=\sum_{a=1}^N z_a(t)|\phi_a\rangle.
\] Its norm is measured by the QSE Gram matrix: \[
\begin{aligned}
\langle\Psi(t)|\Psi(t)\rangle
&=\left(\sum_a z_a(t)|\phi_a\rangle\right)^\dagger
\left(\sum_b z_b(t)|\phi_b\rangle\right)\\
&=\sum_{a,b}z_a(t)^\ast S_{ab}z_b(t)=z(t)^\dagger S z(t).
\end{aligned}
\] Thus \(z^\dagger Sz=1\) is the physical normalization. The vector
\(z(t)\) is a linear QSE coefficient/amplitude vector. The live-circuit
parameter vector is \(\vartheta(t)\) in
\(|\psi(\vartheta(t))\rangle=U(\vartheta(t))|\Phi\rangle\). Frozen QSE
acts on \(z(t)\); McLachlan acts on \(\vartheta(t)\) after live handoff.

Now add a weak or spectrally targeted drive, \[
H(t):=H_0+\sum_\eta\lambda_\eta(t)O_\eta^{\rm drv}.
\] The measured projected matrices are \[
M_0:=\Phi_{\mathcal A}^\dagger H_0\Phi_{\mathcal A},\qquad
D_\eta:=\Phi_{\mathcal A}^\dagger O_\eta^{\rm drv}\Phi_{\mathcal A},
\] so the driven projected Hamiltonian is forced by linearity: \[
\begin{aligned}
M(t)
&:=\Phi_{\mathcal A}^\dagger H(t)\Phi_{\mathcal A}\\
&=\Phi_{\mathcal A}^\dagger
\left(H_0+\sum_\eta\lambda_\eta(t)O_\eta^{\rm drv}\right)
\Phi_{\mathcal A}\\
&=M_0+\sum_\eta\lambda_\eta(t)D_\eta.
\end{aligned}
\] The physical energy expectation inside the frozen state is likewise
\[
\langle\Psi(t)|H(t)|\Psi(t)\rangle
=z(t)^\dagger M(t)z(t),
\qquad z(t)^\dagger Sz(t)=1.
\] The projected time equation is a Galerkin/least-squares condition.
Given the exact Schrodinger vector field \(H(t)|\Psi(t)\rangle\), frozen
QSE chooses a coefficient velocity \(\dot z\) so that the residual \[
|q(\dot z,t)\rangle:=i\Phi_{\mathcal A}\dot z-H(t)\Phi_{\mathcal A}z
\] is orthogonal to the selected QSE bra tests. Equivalently, \(\dot z\)
minimizes \(\|\,i\Phi_{\mathcal A}v-H(t)\Phi_{\mathcal A}z\,\|^2\) over
coefficient velocities \(v\). The normal equation is \[
(i\Phi_{\mathcal A})^\dagger
\bigl(i\Phi_{\mathcal A}\dot z-H(t)\Phi_{\mathcal A}z\bigr)=0,
\] which gives \[
\begin{aligned}
i\Phi_{\mathcal A}^\dagger\Phi_{\mathcal A}\dot z
&=\Phi_{\mathcal A}^\dagger H(t)\Phi_{\mathcal A}z,\\
iS\dot z&=M(t)z.
\end{aligned}
\] The same result follows by substituting
\(|\Psi\rangle=\Phi_{\mathcal A}z\) into
\(i\partial_t|\Psi\rangle=H|\Psi\rangle\), using
\(\partial_t\Phi_{\mathcal A}=0\), and left-multiplying by
\(\Phi_{\mathcal A}^\dagger\). The headline equation is \[
\boxed{iS\dot z(t)=M(t)z(t).}
\] It is the time-dependent generalized Schrodinger equation in the
nonorthogonal QSE coordinate metric \(S\).

Norm conservation is a local check that the projected equation has the
right Hermitian structure. If \(S=S^\dagger\) is fixed and
\(M(t)=M(t)^\dagger\), then \(S\dot z=-iMz\), hence \[
\begin{aligned}
\frac{d}{dt}(z^\dagger Sz)
&=\dot z^\dagger Sz+z^\dagger S\dot z\\
&=(S\dot z)^\dagger z+z^\dagger(S\dot z)\\
&=(-iMz)^\dagger z+z^\dagger(-iMz)\\
&=iz^\dagger Mz-iz^\dagger Mz=0.
\end{aligned}
\] Thus the projected Hermitian flow preserves the physical QSE norm
\(z^\dagger Sz\) without invoking a full inverse of \(S\).

Exact supported whitening turns the same equation into ordinary
retained-coordinate propagation. Diagonalize the metric as
\(S=U\operatorname{diag}(s_1,\ldots,s_N)U^\dagger\), keep
\(J_\tau=\{j:s_j>\tau_S\}\), write \(U_\tau=[u_j]_{j\in J_\tau}\),
\(D_\tau=\operatorname{diag}(s_j)_{j\in J_\tau}\), and define
\(W_\tau:=U_\tau D_\tau^{-1/2}\). Then \(W_\tau^\dagger S W_\tau=I\).
With \(z=W_\tau y\), \[
\begin{aligned}
iS W_\tau\dot y&=M(t)W_\tau y,\\
iW_\tau^\dagger S W_\tau\dot y&=W_\tau^\dagger M(t)W_\tau y,\\
i\dot y&=\widetilde M_\tau(t)y,
\end{aligned}
\qquad
\widetilde M_\tau(t):=W_\tau^\dagger M(t)W_\tau
=\widetilde M_{0,\tau}+\sum_\eta\lambda_\eta(t)\widetilde D_{\eta,\tau},
\] where \(\widetilde M_{0,\tau}:=W_\tau^\dagger M_0W_\tau\) and
\(\widetilde D_{\eta,\tau}:=W_\tau^\dagger D_\eta W_\tau\). The damped
case carries its own metric. If \(z=B_\tau^{(\lambda)}y\), then \[
iG_\tau^{(\lambda)}\dot y=M_\tau^{(\lambda)}(t)y,\qquad
G_\tau^{(\lambda)}:=(B_\tau^{(\lambda)})^\dagger S B_\tau^{(\lambda)},\qquad
M_\tau^{(\lambda)}(t):=(B_\tau^{(\lambda)})^\dagger M(t)B_\tau^{(\lambda)}.
\] For \(B_\tau^{(\lambda)}=U_\tau(D_\tau+\lambda_S I)^{-1/2}\), \[
G_\tau^{(\lambda)}
=(D_\tau+\lambda_S I)^{-1/2}D_\tau(D_\tau+\lambda_S I)^{-1/2}
=D_\tau(D_\tau+\lambda_S I)^{-1}
=\operatorname{diag}\!\left(\frac{s_j}{s_j+\lambda_S}\right)_{j\in J_\tau}.
\] Therefore \(M_\tau^{(\lambda)}(t)\) alone is the exact generator only
for \(\lambda_S=0\). This is the §18.2 damping caveat applied to time
propagation.

The escape diagnostic asks whether the frozen manifold is dynamically
closed. The exact vector \(H(t)|\Psi(t)\rangle\) has a part parallel to
\(|\Psi(t)\rangle\), which changes phase, and a residual part
\((H(t)-E)|\Psi(t)\rangle\), which changes the ray. If that residual is
representable inside \(\operatorname{span}\Phi_{\mathcal A}\), frozen
QSE has the directions needed for the instantaneous drift. If part of it
remains outside, coefficient propagation hides missing physical motion.

For normalized \(z\), define \[
E(z,t):=z^\dagger M(t)z,\qquad
|r(z,t)\rangle:=(H(t)-E(z,t))|\Psi(t)\rangle,\qquad
p(z,t):=(M(t)-E(z,t)S)z.
\] The coordinate residual \(p\) is the selected-basis shadow of the
full Hilbert-space drift residual: \[
p=\Phi_{\mathcal A}^\dagger |r\rangle
=\Phi_{\mathcal A}^\dagger(H(t)-E)\Phi_{\mathcal A}z
=\bigl(M(t)-ES\bigr)z.
\] The full squared drift residual is the energy variance, \[
\begin{aligned}
\langle r|r\rangle
&=\langle\Psi|(H-E)^\dagger(H-E)|\Psi\rangle\\
&=\langle\Psi|(H-E)^2|\Psi\rangle\\
&=\langle\Psi|H(t)^2|\Psi\rangle
-E\langle\Psi|H(t)|\Psi\rangle
-E\langle\Psi|H(t)|\Psi\rangle
+E^2\langle\Psi|\Psi\rangle\\
&=\langle\Psi|H(t)^2|\Psi\rangle-E^2
=\operatorname{Var}_{\Psi(t)}[H(t)].
\end{aligned}
\] The represented part is obtained by allowing the selected QSE span to
explain \(|r\rangle\). Choose coefficients \(\alpha\) for the best
retained-span approximation \(\Phi_{\mathcal A}\alpha\). The
least-squares error is \[
\begin{aligned}
\mathcal E_{\rm drift}(\alpha)
&:=\|\,|r\rangle-\Phi_{\mathcal A}\alpha\,\|^2\\
&=\langle r|r\rangle-\alpha^\dagger\Phi_{\mathcal A}^\dagger|r\rangle
-(\Phi_{\mathcal A}^\dagger|r\rangle)^\dagger\alpha
+\alpha^\dagger S\alpha\\
&=\langle r|r\rangle-\alpha^\dagger p-p^\dagger\alpha+\alpha^\dagger S\alpha.
\end{aligned}
\] Stationarity with respect to \(\alpha^\ast\) gives \[
S\alpha^\star=p,\qquad \alpha^\star=S_\tau^+p
\] on the retained support. The represented vector is
\(\Phi_{\mathcal A}\alpha^\star\), and its represented squared norm is
\[
\|\Phi_{\mathcal A}\alpha^\star\|^2
=(\alpha^\star)^\dagger S\alpha^\star
=p^\dagger S_\tau^+SS_\tau^+p
=p^\dagger S_\tau^+p,
\] again on the same thresholded support. Hence \[
\boxed{
\rho_{\rm esc}(t)
:=
\frac{\bigl[\operatorname{Var}_{\Psi(t)}[H(t)]-p^\dagger S_\tau^+p\bigr]_+}
{\operatorname{Var}_{\Psi(t)}[H(t)]+\epsilon}
},\qquad [x]_+:=\max(x,0).
\] Large \(\rho_{\rm esc}\) means frozen coefficient propagation has an
unrepresented physical drift component outside the current selected
excitation span.

Adaptive QSE growth during dynamics is the middle route. At checkpoint
\(t_k\), keep \(|\Psi_k\rangle=\Phi_{\mathcal A_k}z_k\), keep the
nonorthogonal QSE representation, and admit new records using the static
§18.8 ingredients with \(z_k\), \(H(t_k)\), and the current drift
residual in place of a static root residual. For a candidate record
\(r\), first let the old selected span explain the candidate: \[
|\phi_r\rangle:=Q_{\mathfrak s,0}\Omega_r|\psi_0\rangle,\qquad
s_r:=\Phi_{\mathcal A_k}^\dagger|\phi_r\rangle,\qquad
\beta_r^\star:=S_{\mathcal A_k,\tau}^+s_r,
\] so the outside-span candidate direction is \[
|u_r\rangle:=|\phi_r\rangle-\Phi_{\mathcal A_k}\beta_r^\star
=|\phi_r\rangle-\Phi_{\mathcal A_k}S_{\mathcal A_k,\tau}^+\Phi_{\mathcal A_k}^\dagger|\phi_r\rangle.
\] With \(E_k=z_k^\dagger M(t_k)z_k\) and
\(|r_k\rangle:=(H(t_k)-E_k)|\Psi_k\rangle\), the dynamic residual-gain
proxy is \[
\mathcal G_{\rm dyn}(r;k):=
\frac{|\langle u_r|r_k\rangle|^2}{\langle u_r|u_r\rangle+\epsilon_F}.
\] The numerator is computable from old-span objects plus one
candidate-drift overlap: \[
\begin{aligned}
\langle u_r|r_k\rangle
&=\langle\phi_r|r_k\rangle-(\beta_r^\star)^\dagger\Phi_{\mathcal A_k}^\dagger|r_k\rangle\\
&=\langle\phi_r|(H(t_k)-E_k)|\Psi_k\rangle
-s_r^\dagger S_{\mathcal A_k,\tau}^+p_k,
\end{aligned}
\qquad
p_k:=\Phi_{\mathcal A_k}^\dagger|r_k\rangle.
\] Thus the score measures the component of the current missing drift
along the new direction after old-span redundancy has been removed.
Admission still uses novelty, conditioning, root/trajectory stability,
probe visibility, and measurement/circuit costs. After admitting \(r\),
the basis and matrices enlarge by block updates such as \[
S_{\mathcal A_k\cup\{r\}}=
\begin{pmatrix}S_{\mathcal A_k}&s_r\\ s_r^\dagger&F_r\end{pmatrix},
\qquad
F_r:=\langle\phi_r|\phi_r\rangle,
\] and analogously \(M_0\), \(D_\eta\), and every transition vector
\(d^O\) receive one new row/column or one new component. The propagation
then continues with
\(iS_{\mathcal A_{k+1}}\dot z=M_{\mathcal A_{k+1}}(t)z\).

The live handoff route turns a selected QSE root into a prepared circuit
state. The QSE target is fixed during the refit: \[
|\Psi_\nu^{\rm QSE}\rangle=\sum_{a\in\mathcal A}c_a^{(\nu)}|\phi_a\rangle,\qquad
|\psi_\nu(\vartheta)\rangle:=U_{\mathcal O_\nu}(\vartheta)|\Phi_\nu\rangle,
\] and the optimized variable is \(\vartheta\). A compact root-refit
objective is \[
\begin{aligned}
\mathcal L_\nu(\vartheta)
&=1-\bigl|\langle\Psi_\nu^{\rm QSE}|\psi_\nu(\vartheta)\rangle\bigr|^2\\
&\quad+\lambda_E\Bigl(\langle\psi_\nu(\vartheta)|H_0|\psi_\nu(\vartheta)\rangle-\varepsilon_\nu\Bigr)^2
+\lambda_\perp\sum_{\mu<\nu}\bigl|\langle\psi_\mu|\psi_\nu(\vartheta)\rangle\bigr|^2.
\end{aligned}
\] The first term is one minus the fidelity to the selected QSE root. It
is zero when the circuit state equals the QSE target up to phase. The
second term keeps the circuit-state energy near the QSE Ritz energy
\(\varepsilon_\nu\). The third term penalizes overlap with protected
lower roots \(\{|\psi_\mu\rangle:\mu<\nu\}\). The overlap term expands
through the selected basis: \[
\begin{aligned}
\langle\Psi_\nu^{\rm QSE}|\psi_\nu(\vartheta)\rangle
&=\left(\sum_a c_a^{(\nu)}|\phi_a\rangle\right)^\dagger|\psi_\nu(\vartheta)\rangle\\
&=\sum_{a\in\mathcal A}(c_a^{(\nu)})^\ast\langle\phi_a|\psi_\nu(\vartheta)\rangle\\
&=\sum_{a\in\mathcal A}(c_a^{(\nu)})^\ast
\langle\psi_0|\Omega_a^\dagger Q_{\mathfrak s,0}|\psi_\nu(\vartheta)\rangle .
\end{aligned}
\] Thus the handoff overlap is another QPU-facing transition-style
overlap between the live variational state and the QSE record kets. A
typical handoff acceptance record reports \[
F_\nu(\vartheta^\star):=
\bigl|\langle\Psi_\nu^{\rm QSE}|\psi_\nu(\vartheta^\star)\rangle\bigr|^2,\qquad
\Delta E_\nu(\vartheta^\star):=
\langle\psi_\nu(\vartheta^\star)|H_0|\psi_\nu(\vartheta^\star)\rangle-\varepsilon_\nu,
\] and protected-root overlaps
\(|\langle\psi_\mu|\psi_\nu(\vartheta^\star)\rangle|^2\). After a
successful refit, live McLachlan chooses parameter velocities by
minimizing the residual of the Schrodinger vector field inside the
circuit tangent span: \[
\dot\vartheta^\star
\in
\arg\min_v
\left\|
\sum_j v_j\,\partial_{\vartheta_j}|\psi_\nu(\vartheta)\rangle
+iH(t)|\psi_\nu(\vartheta)\rangle
\right\|^2.
\] The resulting live route uses the checkpoint McLachlan tangent
metric, supported inverse, append/prune/stay decision law, and
integrator; QSE supplies the selected root, escape trigger, protected
lower-root set, and excitation-record registry for live refinement.

The three regimes have different evidence boundaries. Frozen QSE reports
trajectory error for \(iS\dot z=M(t)z\), escape residual
\(\rho_{\rm esc}\), and observable errors computed from the frozen
coefficients. Adaptive QSE growth additionally reports admitted records,
rank/conditioning changes, updated measurement burden, and whether
\(\rho_{\rm esc}\) falls after growth. Live handoff additionally reports
refit loss, root-protection failures, handoff frequency, McLachlan
trajectory errors, and circuit-depth/two-qubit burden. Exact
trajectories may benchmark errors afterward; record admission,
frozen-to-adaptive triggers, and live-handoff decisions are determined
from the QSE/projected quantities and measurement-compatible
diagnostics.

The implementation contract collected from this section is \[
\boxed{
\Phi_{\mathcal A}\ \text{selected spectral coordinate frame},\quad
S=\Phi^\dagger\Phi,\quad
M(t)=\Phi^\dagger H(t)\Phi,\quad
|\Psi(t)\rangle=\Phi z(t)
}
\] \[
\boxed{
\text{frozen QSE: } iS\dot z=M(t)z,\qquad
\text{exact whitening: } i\dot y=W_\tau^\dagger M(t)W_\tau y
}
\] \[
\boxed{
\text{damped coordinates: } iG_\tau^{(\lambda)}\dot y=M_\tau^{(\lambda)}(t)y,\qquad
G_\tau^{(\lambda)}\ne I\ \text{unless }\lambda_S=0
}
\] \[
\boxed{
p=(M-ES)z,\qquad
\rho_{\rm esc}=
\frac{\bigl[\operatorname{Var}_{\Psi}(H)-p^\dagger S_\tau^+p\bigr]_+}
{\operatorname{Var}_{\Psi}(H)+\epsilon}
}
\] \[
\boxed{
\text{adaptive QSE growth: admit records that explain escaping drift},\qquad
\text{live handoff: minimize }\mathcal L_\nu(\vartheta)
}
\] so the QSE root becomes either a propagated frozen coefficient
vector, an enlarged QSE coefficient vector, or a prepared variational
circuit state.

\subsection{18.11 Response dictionary for spectral and dynamical
observables}\label{response-dictionary-for-spectral-and-dynamical-observables}

This section is a print-reference dictionary for Chapter 18. The general
object is a source-coupled Hamiltonian
\(H_\lambda(t)=H_0-\sum_\alpha\lambda_\alpha(t)O_\alpha\), a seed
\(|\psi_0\rangle\), and a selected QSE frame \(\Phi_{\mathcal A}\). A
source component \(\lambda_\alpha(t)\) specifies the physical channel;
the conjugate operator \(O_\alpha\) specifies the Hilbert-space
direction being interrogated; the QSE eigensolve supplies approximate
poles \(\omega_\nu^{\rm QSE}\); transition amplitudes supply residues.
In this language, a probe is an operator-valued functional direction in
Hamiltonian space, a response function is the first-order input-output
kernel associated with that direction, and a spectrum is the
frequency-domain measure obtained after resolving that kernel into
energy differences. The same formalism covers neutral probes
\(O|\psi_0\rangle\in\mathcal H_N\), particle-changing probes
\(c_q^\dagger|\psi_0\rangle\in\mathcal H_{N+1}\) and
\(c_p|\psi_0\rangle\in\mathcal H_{N-1}\), electromagnetic probes
generated by \(A_{\rm ext}\), and driven finite-time propagation under
\(H(t)\).

\[
\boxed{
\text{source}\ \lambda_\alpha
\longleftrightarrow
\text{probe }O_\alpha
\longleftrightarrow
\text{transition vector }d^{O_\alpha}
\longleftrightarrow
\text{residues }(t_\nu^{O_\alpha})^\ast t_\nu^{O_\beta}
\longleftrightarrow
\text{response in }\omega\text{ or }t .
}
\] \begingroup \small \renewcommand{\arraystretch}{1.18}

\begin{longtable}{p{0.18\textwidth}p{0.25\textwidth}p{0.26\textwidth}p{0.23\textwidth}}
\toprule
Object & General source/operation & Implemented QSE quantity & Physical content and Hubbard--Holstein specialization\\
\midrule
Scalar probe spectrum $A_O^{\mathcal A}(\omega)$ &
$H_\lambda(t)=H_0-\lambda(t)O$; seed disturbance $O|\psi_0\rangle$. &
$d^O=\Phi_{\mathcal A}^\dagger O|\psi_0\rangle$, $t_\nu^O=(y^{(\nu)})^\dagger W_\tau^\dagger d^O$, $A_O^{\mathcal A}(\omega)=\sum_\nu |t_\nu^O|^2L_\eta(\omega-\omega_\nu^{\rm QSE})$. &
Frequency-resolved visibility of QSE roots in one channel. For HH, $O=n_k$ gives charge visibility, $O=X_k$ gives displacement visibility, $O=P_k$ gives conjugate oscillator-quadrature visibility, and $O=J$ gives current absorption weights.\\
\midrule
Response matrix $S_{AB}^{\mathcal A}(\omega)$ &
Two-channel linear response: $B$ creates the disturbance and $A$ reads the evolved disturbance. &
$S_{AB}^{\mathcal A}(\omega)=\sum_\nu (t_\nu^A)^\ast t_\nu^B L_\eta(\omega-\omega_\nu^{\rm QSE})$. &
Channel covariance in frequency space. Diagonal entries are scalar spectra; off-diagonal entries diagnose hybridization between response channels. For HH, $S_{nX}$ records correlated charge-displacement dressing.\\
\midrule
Moments $m_k^{AB,\mathcal A}$ &
Weighted spectral sums from the same response kernel. &
$m_k^{AB,\mathcal A}=\sum_\nu(\omega_\nu^{\rm QSE})^k(t_\nu^A)^\ast t_\nu^B$. &
Integrated response audits. $m_0$ tests total retained channel weight; $m_1$ tests first energy-weighted weight. Exact sum rules such as $m_0^{AB}=\langle A^\dagger B\rangle_0$ and $m_1^{AB}=\langle A^\dagger(H_0-E_0)B\rangle_0$ identify missing spectral weight.\\
\midrule
HH neutral channels $S_{nn},S_{XX},S_{nX}$ &
Fourier or standing-wave components of local density and oscillator observables. &
$n_k=\sum_j f_k(j)(n_j-\bar nI)$, $X_k=\sum_j f_k(j)X_j$, $S_{nX}^{\mathcal A}(k,\omega)=\sum_\nu(t_\nu^{n_k})^\ast t_\nu^{X_k}L_\eta(\omega-\omega_\nu^{\rm QSE})$. &
Charge fluctuations, phonon displacement, and joint electron-phonon dressing. For open chains, $f_k(j)$ can be a sine/standing-wave mode or a real-space window; for periodic chains, $f_k(j)=e^{-ikr_j}$.\\
\midrule
Single-particle Green function $G_{pq}^R(\omega)$ &
Fermion addition $c_q^\dagger|\psi_0\rangle$ and removal $c_p|\psi_0\rangle$; propagation in $N+1$ and $N-1$ sectors. &
$G_{pq}^{R,\mathcal A}(\omega)=\sum_m u_{pm}^+(u_{qm}^+)^\ast/[\omega-(\varepsilon_m^+-E_0)+i\eta]+\sum_n(v_{qn}^-)^\ast v_{pn}^-/[\omega+(\varepsilon_n^--E_0)+i\eta]$. &
Dressed electron and hole propagation. Poles define quasiparticle additions/removals; the separation between occupied/removal and unoccupied/addition edges defines a single-particle gap under the chosen frequency zero; repeated phonon-scale satellite peaks are phonon sidebands.\\
\midrule
Spectral density $A_{pq}(\omega)$ &
Imaginary part of the retarded single-particle propagator. &
$A_{pq}(\omega)=-(1/\pi)\operatorname{Im}G_{pq}^R(\omega)$; $A_{\rm loc}=|\mathcal P_G|^{-1}\sum_{p\in\mathcal P_G}A_{pp}$; $A_{k\sigma}$ from orbital-index Fourier or mode projection. &
Distribution of electron-addition/removal weight over energy. Peaks correspond to coherent or incoherent particle excitations; in HH, coherent quasiparticle peaks can be accompanied by phonon sidebands from lattice dressing.\\
\midrule
Optical conductivity $\sigma(\omega)$ &
Uniform vector potential $A_{\rm ext}$ with $H(A_{\rm ext})=H_0-A_{\rm ext}J+\frac12A_{\rm ext}^2D+\cdots$. &
$\chi_{JJ}^{R,\mathcal A}(\omega)=\sum_\nu|t_\nu^{J,\mathcal A}|^2[(\omega-\omega_\nu^{\rm QSE}+i\eta)^{-1}-(\omega+\omega_\nu^{\rm QSE}+i\eta)^{-1}]$, $\sigma_{\mathcal A}(\omega)=i[\langle D\rangle_0+\chi_{JJ}^{R,\mathcal A}(\omega)]/[L(\omega+i\eta)]$. &
Current induced by an oscillatory electric field. The paramagnetic part is delayed current-current response; $D$ is the instantaneous contact curvature. HH optical peaks identify current-active dressed charge excitations.\\
\midrule
Frozen QSE dynamics &
Time-dependent Hamiltonian $H(t)=H_0+\sum_\eta\lambda_\eta(t)O_\eta^{\rm drv}$ projected into a fixed QSE frame. &
$|\Psi(t)\rangle=\Phi_{\mathcal A}z(t)$, $iS\dot z=M(t)z$, $M(t)=M_0+\sum_\eta\lambda_\eta(t)D_\eta$. &
Finite-time state trajectory inside a fixed measured spectral coordinate frame. Accuracy requires the driven Schrödinger drift to remain representable by the selected records.\\
\midrule
Escape diagnostic $\rho_{\rm esc}$ &
Variance decomposition of instantaneous drift $r=(H-E)|\Psi\rangle$. &
$p=(M-ES)z$, $\rho_{\rm esc}=[\operatorname{Var}_{\Psi}(H)-p^\dagger S_\tau^+p]_+/[\operatorname{Var}_{\Psi}(H)+\epsilon]$. &
Fraction of physical drift variance unexplained by the current QSE frame. Large values trigger adaptive record admission or live handoff.\\
\bottomrule
\end{longtable}
\endgroup

The most compact implementation dictionary is therefore \[
\boxed{
d^O=\Phi_{\mathcal A}^\dagger O|\psi_0\rangle,\qquad
t_\nu^O=(y^{(\nu)})^\dagger W_\tau^\dagger d^O,\qquad
I_\nu^O=|t_\nu^O|^2
}
\] \[
\boxed{
S_{AB}^{\mathcal A}(\omega)
=\sum_{\nu\in\mathcal W_{\rm targ}}(t_\nu^A)^\ast t_\nu^B
L_\eta(\omega-\omega_\nu^{\rm QSE}),\qquad
m_k^{AB,\mathcal A}
=\sum_{\nu\in\mathcal W_{\rm targ}}(\omega_\nu^{\rm QSE})^k(t_\nu^A)^\ast t_\nu^B
}
\] \[
\boxed{
G_{pq}^{R,\mathcal A}(\omega)
=
\sum_m\frac{u_{pm}^+(u_{qm}^+)^\ast}{\omega-(\varepsilon_m^+-E_0)+i\eta}
+
\sum_n\frac{(v_{qn}^-)^\ast v_{pn}^-}{\omega+(\varepsilon_n^--E_0)+i\eta}
}
\] \[
\boxed{
\sigma_{\mathcal A}(\omega)
=
\frac{i}{L(\omega+i\eta)}
\left(\langle D\rangle_0+\chi_{JJ}^{R,\mathcal A}(\omega)\right),
\qquad
iS\dot z=M(t)z,\qquad
\rho_{\rm esc}=
\frac{[\operatorname{Var}_{\Psi}(H)-p^\dagger S_\tau^+p]_+}
{\operatorname{Var}_{\Psi}(H)+\epsilon}
}
\] All entries reduce to the same measured primitives:
overlap/Hamiltonian matrices \(S,M\), transition vectors \(d^O\),
sector-specific addition/removal matrices for \(G^R\), current/contact
operators \(J,D\), and resource diagnostics controlling finite-shot
trust. The physics interpretation is fixed by the source operator:
density probes charge redistribution, displacement probes oscillator
motion, mixed response probes electron-phonon dressing, fermion
creation/removal probes quasiparticle and hole spectra, current probes
optical absorption, and driven \(H(t)\) probes finite-time state motion.

\subsection{18.12 Evidence targets and diagnostic
boundaries}\label{evidence-targets-and-diagnostic-boundaries}

The static spectral evidence should report the objects defined above in
addition to sorted energies. Energies test the Ritz locations,
transition strengths test observable visibility, spectral functions test
the plotted response, moments test integrated sum-rule closure, Green
functions test particle-addition/removal sectors, and conditioning
diagnostics test whether the nonorthogonal solve was numerically
trustworthy. A spectra claim is strongest when these objects agree
rather than when only the first few roots are close.

\begin{enumerate}
\def\labelenumi{\arabic{enumi}.}
\tightlist
\item
  excitation-gap error \(\omega_\nu^{\rm QSE}-\omega_\nu^{\rm ref}\);
\item
  transition-strength error
  \(\big||t_\nu^O|^2-|t_{\nu,{\rm ref}}^O|^2\big|\);
\item
  spectral-function error such as \[
  \|A_O^{\rm QSE}-A_O^{\rm ref}\|_{2,{\rm grid}}
  =
  \left(\sum_j w_j|A_O^{\rm QSE}(\omega_j)-A_O^{\rm ref}(\omega_j)|^2\right)^{1/2};
  \]
\item
  response-matrix and moment errors for \(\mathsf S_{AB}(\omega)\),
  \(C_{AB}(t)\), and \(m_k^{AB}\);
\item
  Green-function spectral error for the specified addition/removal
  sectors;
\item
  optional conductivity error only when the current, contact term, and
  broadening convention are fixed;
\item
  selected-basis rank, \(\kappa_\tau(S)\), root-tracking stability,
  cutoff stress, and \(S_{\rm mat}^{\rm proxy}\).
\end{enumerate}

The dynamics evidence should report frozen-QSE trajectory error, escape
residual, adaptive-QSE record growth, live-handoff frequency,
root-protection failures, observable trajectory errors, and the
additional measurement or circuit burden needed after static matrix
construction. These quantities separate three possible failure modes:
inaccurate static spectral manifold, dynamically leaky frozen manifold,
and successful root preservation with excessive circuit or measurement
overhead.

Exact diagonalization and exact trajectories are diagnostic references
for reported errors, spectral overlays, and root matching. QPU-facing
record admission, frozen-to-adaptive decisions, live-handoff triggers,
and checkpoint-controller decisions use measurement-compatible estimates
of the prepared seed state, candidate excitation states, projected
matrix elements, response probes, and resource proxies.

\section{19. Final Primitive-Closed
Summary}\label{final-primitive-closed-summary}

Section 19 introduces no new variational rule. It closes the dependency
graph. Every later object--Hamiltonians, drives, statevector actions,
variational circuits, ADAPT gradients, McLachlan metrics, QSE matrices,
response functions, continuation bundles, benchmark contracts, and noise
diagnostics--descends from the same finite primitive alphabet. The
reader contract is: primitive labels first, composite operators second,
direct substitution third, reduced master forms last.

Interpretive compression: the document has a single mathematical layer.
Fixed fermion ladder operators and a fixed finite boson encoding imply a
Pauli-polynomial Hamiltonian. A Pauli-polynomial Hamiltonian implies
statevector sums or QPU Pauli measurements for energies and matrix
elements. Those expectations generate gradients, tangent metrics, QSE
matrices, spectra, Green functions, optical response, and benchmark
residuals.

Start with the discrete physical labels. A site is
\(i\in\{0,\ldots,L-1\}\), a spin label is
\(\sigma\in\{\uparrow,\downarrow\}\equiv\{0,1\}\), and the fermionic
mode index is the qubit index \(p(i,\sigma)\). The two standard
orderings are \[
p_{\rm int}(i,\sigma)=2i+\sigma,
\qquad
p_{\rm blk}(i,\uparrow)=i,\qquad
p_{\rm blk}(i,\downarrow)=L+i.
\] Thus \[
(i,\sigma)\longmapsto p(i,\sigma)\longmapsto q=p(i,\sigma),
\] so every fermionic operator has a definite qubit location. The
register convention is rightmost \(q_0\), least-significant bit; Pauli
words are printed as \(q_{N_q-1}\cdots q_1q_0\). This fixes both printed
Pauli placement and computational-basis bit extraction.

The Jordan-Wigner primitives are \[
\widehat c_p^\dagger
=
\frac12(X_p-iY_p)\prod_{r=0}^{p-1}Z_r,
\qquad
\widehat c_p
=
\frac12(X_p+iY_p)\prod_{r=0}^{p-1}Z_r.
\] With \[
Z_{<p}:=\prod_{r=0}^{p-1}Z_r,
\] we have \(Z_{<p}^2=I\), and \(X_p,Y_p\) commute with \(Z_{<p}\)
because the string acts only on qubits \(r<p\). Therefore the number
operator reduces to one qubit: \[
\begin{aligned}
\widehat n_p
&=\widehat c_p^\dagger \widehat c_p\\
&=
\frac14(X_p-iY_p)Z_{<p}Z_{<p}(X_p+iY_p)\\
&=
\frac14(X_p-iY_p)(X_p+iY_p)\\
&=
\frac14\left(X_p^2+iX_pY_p-iY_pX_p+Y_p^2\right).
\end{aligned}
\] Using \[
X^2=Y^2=I,\qquad XY=iZ,\qquad YX=-iZ,
\] gives \[
\begin{aligned}
\widehat n_p
&=
\frac14\left(I+i(iZ_p)-i(-iZ_p)+I\right)\\
&=
\frac14\left(2I-Z_p-Z_p\right)\\
&=
\frac12(I-Z_p).
\end{aligned}
\] So occupation is already a Pauli polynomial: \[
\boxed{\widehat n_p=\frac{I-Z_p}{2}.}
\] The site density is then \[
\widehat n_i
:=
\widehat n_{p(i,\uparrow)}+\widehat n_{p(i,\downarrow)},
\] hence \[
\begin{aligned}
\widehat n_i
&=
\frac12(I-Z_{p(i,\uparrow)})
+
\frac12(I-Z_{p(i,\downarrow)})\\
&=
I-\frac12\left(Z_{p(i,\uparrow)}+Z_{p(i,\downarrow)}\right).
\end{aligned}
\] The doublon operator is \[
\widehat D_i:=\widehat n_{i\uparrow}\widehat n_{i\downarrow},
\] so \[
\begin{aligned}
\widehat D_i
&=
\left(\frac{I-Z_{p(i,\uparrow)}}2\right)
\left(\frac{I-Z_{p(i,\downarrow)}}2\right)\\
&=
\frac14\left(
I
-Z_{p(i,\uparrow)}
-Z_{p(i,\downarrow)}
+Z_{p(i,\uparrow)}Z_{p(i,\downarrow)}
\right).
\end{aligned}
\] Thus density and doublon observables are closed inside the Pauli
algebra.

For hopping, let \(p<q\). Then \[
\widehat c_p^\dagger \widehat c_q
=
\frac14(X_p-iY_p)Z_{<p}(X_q+iY_q)Z_{<q}.
\] Since \[
Z_{<q}=Z_{<p}Z_pZ_{p+1}\cdots Z_{q-1},
\qquad
Z_{<p}Z_{<q}=Z_pZ_{p+1}\cdots Z_{q-1},
\] and since the \(Z_p\) factor anticommutes through \(X_p\) and
\(Y_p\), the standard bilinear reductions are \[
\boxed{
\widehat c_p^\dagger\widehat c_q+\widehat c_q^\dagger\widehat c_p
=
\frac12\left(
X_pZ_{p+1}\cdots Z_{q-1}X_q
+
Y_pZ_{p+1}\cdots Z_{q-1}Y_q
\right)
}
\] and \[
\boxed{
i(\widehat c_p^\dagger\widehat c_q-\widehat c_q^\dagger\widehat c_p)
=
\frac12\left(
Y_pZ_{p+1}\cdots Z_{q-1}X_q
-
X_pZ_{p+1}\cdots Z_{q-1}Y_q
\right).
}
\] These two identities close the one-body hopping and current channels.

The Hubbard Hamiltonian is \[
\widehat H_{\rm Hub}
=
\widehat H_t+\widehat H_U+\widehat H_v,
\] with \[
\widehat H_t
=
-t\sum_{\langle i,j\rangle,\sigma}
\left(
\widehat c_{i\sigma}^\dagger\widehat c_{j\sigma}
+
\widehat c_{j\sigma}^\dagger\widehat c_{i\sigma}
\right),
\] \[
\widehat H_U
=
U\sum_i \widehat n_{i\uparrow}\widehat n_{i\downarrow},
\qquad
\widehat H_v
=
\sum_{i,\sigma}v_i\widehat n_{i\sigma}.
\] If \(p_{i\sigma}:=p(i,\sigma)\), \(p_{j\sigma}:=p(j,\sigma)\), and
the displayed string assumes \(p_{i\sigma}<p_{j\sigma}\), then
substitution gives \[
\begin{aligned}
\widehat H_t
&=
-\frac t2
\sum_{\langle i,j\rangle,\sigma}
\left(
X_{p_{i\sigma}}Z_{p_{i\sigma}+1}\cdots Z_{p_{j\sigma}-1}X_{p_{j\sigma}}
+
Y_{p_{i\sigma}}Z_{p_{i\sigma}+1}\cdots Z_{p_{j\sigma}-1}Y_{p_{j\sigma}}
\right),
\end{aligned}
\] with endpoints swapped when \(p_{j\sigma}<p_{i\sigma}\). Also \[
\widehat H_U
=
\frac U4\sum_i
\left(I-Z_{p(i,\uparrow)}\right)
\left(I-Z_{p(i,\downarrow)}\right),
\] and \[
\widehat H_v
=
\sum_i v_i
\left[
I-\frac12\left(Z_{p(i,\uparrow)}+Z_{p(i,\downarrow)}\right)
\right].
\] Therefore \[
\widehat H_{\rm Hub}\in
\operatorname{span}_{\mathbb C}\{P:\ P\text{ is a Pauli word}\}.
\] This is the first closure: after choosing \(p(i,\sigma)\), the
Hubbard model is a Pauli polynomial.

Now close the boson side. The local truncated oscillator space is \[
\mathcal H_{b,i}
=
\operatorname{span}\{|n\rangle_i:0\le n\le n_{\rm ph,max}\},
\qquad
d=n_{\rm ph,max}+1.
\] The finite ladder matrices are \[
\widehat b_i^\dagger
=
\sum_{n=0}^{d-2}\sqrt{n+1}\,|n+1\rangle\langle n|,
\qquad
\widehat b_i
=
\sum_{n=1}^{d-1}\sqrt n\,|n-1\rangle\langle n|.
\] Thus \[
\widehat n_{b,i}
=
\widehat b_i^\dagger\widehat b_i
=
\sum_{n=0}^{d-1}n|n\rangle\langle n|,
\] \[
\widehat x_i
:=
\widehat b_i+\widehat b_i^\dagger
=
\sum_{n=0}^{d-2}\sqrt{n+1}
\left(
|n+1\rangle\langle n|+|n\rangle\langle n+1|
\right),
\] and \[
\widehat P_i
:=
i(\widehat b_i^\dagger-\widehat b_i)
=
i\sum_{n=0}^{d-2}\sqrt{n+1}
\left(
|n+1\rangle\langle n|-|n\rangle\langle n+1|
\right).
\] These are Pauli polynomials after encoding because every matrix unit
\(|m\rangle\langle n|\) becomes a tensor product of one-qubit matrix
units. For one qubit, \[
|0\rangle\langle0|=\frac{I+Z}{2},
\qquad
|1\rangle\langle1|=\frac{I-Z}{2},
\] \[
|0\rangle\langle1|=\frac{X+iY}{2},
\qquad
|1\rangle\langle0|=\frac{X-iY}{2}.
\] If \(m=m_{q_b-1}\cdots m_0\) and \(n=n_{q_b-1}\cdots n_0\) are binary
codewords, then \[
|m\rangle\langle n|
=
\bigotimes_{\ell=0}^{q_b-1}|m_\ell\rangle\langle n_\ell|,
\] so it expands into Pauli words. Unary encoding uses one-hot valid
codewords, but each valid \(|m\rangle\langle n|\) is still a tensor
product of the same one-qubit matrix units. Therefore \[
\widehat b_i,\widehat b_i^\dagger,\widehat n_{b,i},\widehat x_i,\widehat P_i
\text{ are finite Pauli polynomials after encoding.}
\] The Hubbard-Holstein Hamiltonian is \[
\widehat H_{\rm HH}(t)
=
\widehat H_t+\widehat H_U+\widehat H_v+\widehat H_{\rm ph}+\widehat H_g+\widehat H_{\rm drive}(t),
\] where \[
\widehat H_{\rm ph}
=
\omega_0\sum_i\left(\widehat n_{b,i}+\frac12 I\right),
\qquad
\widehat H_g
=
g\sum_i\widehat x_i(\widehat n_i-\bar n I).
\] The operator \(\widehat x_i\) acts on the boson register and
\(\widehat n_i-\bar n I\) acts on the fermion register, so they commute
and their product is a tensor product of Pauli polynomials: \[
\widehat x_i(\widehat n_i-\bar n I)
=
\widehat x_i\otimes
\left[
(1-\bar n)I-\frac12\left(Z_{p(i,\uparrow)}+Z_{p(i,\downarrow)}\right)
\right].
\] A site-density drive has \[
\widehat H_{\rm drive}(t)
=
\sum_i v_i(t)\widehat n_i
=
\sum_i v_i(t)
\left[
I-\frac12\left(Z_{p(i,\uparrow)}+Z_{p(i,\downarrow)}\right)
\right].
\] Collecting coefficients by Pauli word, \[
\Delta c[I](t)=\sum_i v_i(t),
\qquad
\Delta c[Z_{p(i,\uparrow)}](t)=-\frac12v_i(t),
\qquad
\Delta c[Z_{p(i,\downarrow)}](t)=-\frac12v_i(t).
\] Thus time dependence changes scalar coefficients while preserving the
operator alphabet: \[
\widehat H_{\rm drive}(t)
=
\sum_{\lambda\in\mathcal L_{\rm drive}}c_\lambda(t)P_\lambda.
\] Hence \[
\boxed{
\widehat H_{\rm HH}(t)=\sum_\lambda h_\lambda(t)P_\lambda.
}
\] This is the second closure: static and driven HH physics are
represented in one finite Pauli algebra.

Those operators then feed the statevector primitives. Let \(|k\rangle\)
be a computational-basis state, and define \[
b_q(k):=\left\lfloor \frac{k}{2^q}\right\rfloor\bmod 2.
\] The one-qubit Pauli actions are \[
I|b\rangle=|b\rangle,
\qquad
Z|b\rangle=(-1)^b|b\rangle,
\] \[
X|b\rangle=|1-b\rangle,
\qquad
Y|b\rangle=i(-1)^b|1-b\rangle.
\] For a Pauli word \(P=\bigotimes_qP_q\), let \(f_P\) be the bit-flip
mask from its \(X\) and \(Y\) entries, and let \(\chi_P(k)\) be the
product of all local \(Z\)- and \(Y\)-phases. Then \[
P|k\rangle=\chi_P(k)|k\oplus f_P\rangle.
\] For \[
|\psi\rangle=\sum_k\psi_k|k\rangle,
\] the expectation of one Pauli word is \[
\begin{aligned}
\langle\psi|P|\psi\rangle
&=
\left(\sum_j\psi_j^\ast\langle j|\right)
\left(\sum_k\psi_k\chi_P(k)|k\oplus f_P\rangle\right)\\
&=
\sum_{j,k}\psi_j^\ast\psi_k\chi_P(k)\delta_{j,k\oplus f_P}\\
&=
\sum_k\psi_{k\oplus f_P}^\ast\chi_P(k)\psi_k.
\end{aligned}
\] For a Pauli polynomial \(A=\sum_\lambda a_\lambda P_\lambda\), \[
\langle A\rangle_\psi
=
\langle\psi|A|\psi\rangle
=
\sum_\lambda a_\lambda\langle\psi|P_\lambda|\psi\rangle.
\] Thus Hamiltonian energies, gradients, transition elements, QSE
matrices, response weights, and diagnostics reduce to Pauli-word
expectations.

Variational manifolds use the same operators. A fixed or adaptive ansatz
has \[
|\psi(\theta;\mathcal O)\rangle
=
U(\theta;\mathcal O)|\psi_{\rm ref}\rangle,
\qquad
U(\theta;\mathcal O)
=
\prod_{\mu=1}^{N_\theta}e^{-i\theta_\mu A_\mu},
\] where each Hermitian generator \(A_\mu\) is assembled from the same
Pauli primitives. If \(A_\mu=P\) is a single Pauli word, then \(P^2=I\),
so \[
\begin{aligned}
e^{-i\theta P}
&=
\sum_{n=0}^{\infty}\frac{(-i\theta P)^n}{n!}\\
&=
\sum_{m=0}^{\infty}\frac{(-i\theta)^{2m}}{(2m)!}I
+
\sum_{m=0}^{\infty}\frac{(-i\theta)^{2m+1}}{(2m+1)!}P\\
&=
\cos\theta\,I-i\sin\theta\,P.
\end{aligned}
\] If \(A_\mu=\sum_r a_{\mu r}P_r\), then the exponential is implemented
as a primitive Pauli-polynomial exponential, a product over commuting
Pauli terms, or a chosen approximation surface; the generator itself
still comes from the same alphabet.

ADAPT selection is primitive-closed as well. For a candidate Hermitian
generator \(A\), define \[
|\psi(\theta)\rangle=e^{-i\theta A}|\psi\rangle,
\qquad
E(\theta)=\langle\psi(\theta)|H|\psi(\theta)\rangle.
\] Then \[
E(\theta)=\langle\psi|e^{i\theta A}He^{-i\theta A}|\psi\rangle,
\] and \[
\begin{aligned}
\left.\frac{dE}{d\theta}\right|_{\theta=0}
&=
\left.
\langle\psi|
\left(
iAe^{i\theta A}He^{-i\theta A}
+
e^{i\theta A}H(-iA)e^{-i\theta A}
\right)
|\psi\rangle
\right|_{\theta=0}\\
&=
i\langle\psi|AH|\psi\rangle-i\langle\psi|HA|\psi\rangle\\
&=
i\langle\psi|[A,H]|\psi\rangle.
\end{aligned}
\] Equivalently, with \(z:=\langle H\psi|A\psi\rangle\), \[
\begin{aligned}
\left.\frac{dE}{d\theta}\right|_0
&=
i\left(\langle A\psi|H\psi\rangle-\langle H\psi|A\psi\rangle\right)\\
&=
i(z^\ast-z)\\
&=
2\operatorname{Im}z\\
&=
2\operatorname{Im}\langle H\psi|A\psi\rangle.
\end{aligned}
\] So the ADAPT gradient uses only \(H\), \(A\), and the prepared state,
all primitive-closed.

Projective McLachlan geometry also descends from the same closure. Let
\(|\psi(\theta)\rangle\) be normalized, define \[
|\partial_\mu\psi\rangle:=\frac{\partial}{\partial\theta_\mu}|\psi(\theta)\rangle,
\] and remove the physically redundant global-phase component: \[
|\widetilde\partial_\mu\psi\rangle
:=
(I-|\psi\rangle\langle\psi|)|\partial_\mu\psi\rangle.
\] The exact real-time equation is
\(\partial_t|\psi\rangle=-iH(t)|\psi\rangle\), while a path on the
variational manifold has \[
\partial_t|\psi(\theta(t))\rangle
=
\sum_\nu\dot\theta_\nu|\partial_\nu\psi\rangle.
\] Projective McLachlan chooses \(\dot\theta\) to minimize \[
\left\|
\sum_\nu\dot\theta_\nu|\widetilde\partial_\nu\psi\rangle
+
i(H-\langle H\rangle)|\psi\rangle
\right\|^2.
\] Set \[
|R\rangle
=
\sum_\nu\dot\theta_\nu|\widetilde\partial_\nu\psi\rangle
+
i(H-\langle H\rangle)|\psi\rangle.
\] Stationarity with respect to each real velocity \(\dot\theta_\mu\)
gives \[
0=\operatorname{Re}\langle\widetilde\partial_\mu\psi|R\rangle.
\] Thus \[
\begin{aligned}
0
&=
\operatorname{Re}
\left[
\sum_\nu\dot\theta_\nu
\langle\widetilde\partial_\mu\psi|\widetilde\partial_\nu\psi\rangle
+
i\langle\widetilde\partial_\mu\psi|(H-\langle H\rangle)|\psi\rangle
\right]\\
&=
\sum_\nu
\operatorname{Re}
\langle\widetilde\partial_\mu\psi|\widetilde\partial_\nu\psi\rangle
\dot\theta_\nu
-
\operatorname{Im}
\langle\widetilde\partial_\mu\psi|H|\psi\rangle.
\end{aligned}
\] With \[
G_{\mu\nu}:=
\operatorname{Re}
\langle\widetilde\partial_\mu\psi|\widetilde\partial_\nu\psi\rangle,
\qquad
f_\mu:=
\operatorname{Im}
\langle\widetilde\partial_\mu\psi|H|\psi\rangle,
\] the live dynamics equation is \[
\boxed{G\dot\theta=f.}
\] Each entry of \(G\) and \(f\) is an overlap or expectation involving
a circuit built from primitive generators and a Hamiltonian built from
Pauli words.

The local defect geometry used for append, prune, and branch decisions
is the same projection viewed as an error vector. Once \(G\dot\theta=f\)
is solved, the represented tangent drift is \[
|\dot\psi_{\rm var}\rangle
=
\sum_\mu\dot\theta_\mu|\widetilde\partial_\mu\psi\rangle,
\] while the exact projected Schrödinger drift is \[
|\dot\psi_{\rm exact}\rangle
=
-i(H-\langle H\rangle)|\psi\rangle.
\] The local defect is \[
|r\rangle
:=
|\dot\psi_{\rm exact}\rangle-|\dot\psi_{\rm var}\rangle.
\] For a candidate new direction \(|u_m\rangle\), a direct gain proxy is
\[
\mathcal G(m)
=
\frac{|\langle u_m|r\rangle|^2}{\langle u_m|u_m\rangle+\epsilon}.
\] This is the same geometry as QSE novelty and residual repair: project
the missing physical drift onto a candidate direction, normalize by
candidate size, and admit directions that explain a real defect rather
than directions that merely have an attractive label.

Continuation and handoff are state transport between already-defined
manifolds. A handoff bundle may contain \[
\mathcal B
=
(\text{manifest},\mathcal O,\theta,E,\text{observables},\text{sector labels},\text{continuation metadata}),
\] where \(\mathcal O\) is an ordered primitive-generator scaffold,
\(\theta\) are its coordinates,
\(E=\langle\psi(\theta)|H|\psi(\theta)\rangle\), and observables are
Pauli-polynomial expectations. If a QSE root is handed into a live
ansatz, the refit overlap is \[
\langle\Psi_\nu^{\rm QSE}|\psi_\nu(\vartheta)\rangle
=
\sum_{a\in\mathcal A}(c_a^{(\nu)})^\ast
\langle\phi_a|\psi_\nu(\vartheta)\rangle.
\] Since \[
|\phi_a\rangle=Q_{\mathfrak s,0}\Omega_a|\psi_0\rangle,
\] the overlap reduces again to seed/candidate/circuit matrix elements
made from primitive operators. Handoff changes the coordinate
representation while preserving the operator layer.

The benchmark and noise-validation contracts close the final loop. The
filtered exact target is \[
E_{\rm exact,filtered}
=
\min_{\psi\in\mathcal H_{\rm sector}}
\langle\psi|H|\psi\rangle,
\] whereas the unrestricted target is \[
E_{\rm exact,full}
=
\min\operatorname{spec}(H).
\] These coincide only when the computational and physical sectors
coincide. Noise deltas are observable-wise: \[
\Delta\mathcal O(t)
=
\mathcal O_{\rm noisy}(t)-\mathcal O_{\rm ideal}(t),
\] so, for example, \[
\Delta E(t)=E_{\rm noisy}(t)-E_{\rm ideal}(t),
\qquad
\Delta D(t)=D_{\rm noisy}(t)-D_{\rm ideal}(t).
\] A parity gate against an anchor trajectory is \[
\max_{t\in\mathcal T}
|\mathcal O_{\rm new}(t)-\mathcal O_{\rm anchor}(t)|
\le
\epsilon_{\rm parity}.
\] These contracts remain primitive-closed because \(H\), observables,
sectors, trajectories, and expectations all reduce to the same encoded
Hilbert space and Pauli-polynomial observables.

The whole manuscript can therefore be read as one dependency diagram: \[
\begin{aligned}
&L,\ p(i,\sigma),\ n_{\rm ph,max},\ \text{boson encoding}\\
&\Downarrow\\
&\widehat c_p^\dagger,\widehat c_p,\widehat n_p,\widehat b_i^\dagger,\widehat b_i,\widehat n_{b,i},\widehat x_i,\widehat P_i\\
&\Downarrow\\
&\widehat H_t,\widehat H_U,\widehat H_v,\widehat H_{\rm ph},\widehat H_g,\widehat H_{\rm drive}(t)\\
&\Downarrow\\
&H(t)=\sum_\lambda h_\lambda(t)P_\lambda\\
&\Downarrow\\
&\langle P_\lambda\rangle_\psi,\quad
E(\theta),\quad
g_A,\quad
G_{\mu\nu},\quad
f_\mu,\quad
S_{ab},\quad
M_{ab},\quad
d_a^O\\
&\Downarrow\\
&\text{VQE, ADAPT, McLachlan dynamics, QSE, response functions, Green functions, conductivity, diagnostics}\\
&\Downarrow\\
&\text{benchmark and noise-validation contracts}.
\end{aligned}
\] A practical way to read this section is as a mathematical audit. If a
later formula contains a symbol, walk it backward. Optical response
obeys \[
\sigma(\omega)
\leftarrow
\chi_{JJ}^R(\omega),D,J
\leftarrow
H(A_{\rm ext})
\leftarrow
\widehat c_p^\dagger\widehat c_q+\widehat c_q^\dagger\widehat c_p
\leftarrow
X_pZ\cdots X_q+Y_pZ\cdots Y_q.
\] Spectral functions obey \[
\begin{aligned}
A_O(\omega)
&\leftarrow
t_\nu^O,\omega_\nu\\
&\leftarrow
(c^{(\nu)})^\dagger d^O,\quad Mc=\varepsilon Sc\\
&\leftarrow
S_{ab},M_{ab},d_a^O\\
&\leftarrow
\langle\psi_0|\Omega_a^\dagger Q\Omega_b|\psi_0\rangle,\quad
\langle\psi_0|\Omega_a^\dagger QHQ\Omega_b|\psi_0\rangle,\quad
\langle\psi_0|\Omega_a^\dagger QO|\psi_0\rangle\\
&\leftarrow
\text{Pauli expectations}.
\end{aligned}
\] Dynamics obeys \[
\dot\theta
\leftarrow
G\dot\theta=f
\leftarrow
\operatorname{Re}\langle\widetilde\partial_\mu\psi|\widetilde\partial_\nu\psi\rangle,\quad
\operatorname{Im}\langle\widetilde\partial_\mu\psi|H|\psi\rangle
\leftarrow
U(\theta),H
\leftarrow
\text{primitive generators and Pauli Hamiltonian}.
\] The final invariant is this: every path through the manuscript
terminates at the same primitive basis. The symbolic duplicate removes
implementation labels so the mathematical skeleton is visible: \[
\begin{gathered}
\text{ordering map}
\to
\text{ladder primitives}
\to
\text{Pauli Hamiltonian}\\
\to
\text{state/action/expectation primitives}\\
\to
\text{variational, QSE, and dynamical geometry}\\
\to
\text{benchmark contracts}.
\end{gathered}
\] That closure is what makes the earlier sections mutually compatible
rather than a collection of disconnected recipes.

\section{20. Data}\label{data}

This final chapter records empirical evidence for the selector and
pruning claims developed in the main body. The goal is not to replace
the symbolic formulas above, but to show that the live Phase-3 HH
selector surfaces described in §§11 and 17 materially change the
observed outcomes on matched runs.

\subsection{20.1 Matched novelty
ablation}\label{matched-novelty-ablation}

The freshest clean ablation available is a matched four-run HH study
with \[
L=2,\qquad t=1,\qquad U=4,\qquad \omega_0=1,\qquad g=0.5,\qquad n_{\mathrm{ph,max}}=1,
\] using binary bosons, blocked ordering, open boundaries, the same
seed, the same compiled backend path, the same windowed reoptimization
policy, the same live pruning machinery, and the same
rescue/symmetry/lifetime settings. The only changes are:

\begin{enumerate}
\def\labelenumi{\arabic{enumi}.}
\tightlist
\item
  operator pool: \texttt{pareto\_lean\_l2} versus \texttt{full\_meta},
\item
  novelty exponent: \(\gamma_N=1\) versus \(\gamma_N=0\).
\end{enumerate}

Writing the final filtered-energy error as \[
\lvert \Delta E \rvert = \bigl|E_{\mathrm{final}}-E_{\mathrm{exact,filtered}}\bigr|,
\] the matched runs give:

{\def\LTcaptype{none} % do not increment counter
\begin{longtable}[]{@{}
  >{\raggedright\arraybackslash}p{(\linewidth - 14\tabcolsep) * \real{0.1000}}
  >{\raggedleft\arraybackslash}p{(\linewidth - 14\tabcolsep) * \real{0.1333}}
  >{\raggedleft\arraybackslash}p{(\linewidth - 14\tabcolsep) * \real{0.1333}}
  >{\raggedleft\arraybackslash}p{(\linewidth - 14\tabcolsep) * \real{0.1333}}
  >{\raggedleft\arraybackslash}p{(\linewidth - 14\tabcolsep) * \real{0.1333}}
  >{\raggedleft\arraybackslash}p{(\linewidth - 14\tabcolsep) * \real{0.1333}}
  >{\raggedright\arraybackslash}p{(\linewidth - 14\tabcolsep) * \real{0.1000}}
  >{\raggedleft\arraybackslash}p{(\linewidth - 14\tabcolsep) * \real{0.1333}}@{}}
\toprule\noalign{}
\begin{minipage}[b]{\linewidth}\raggedright
Pool
\end{minipage} & \begin{minipage}[b]{\linewidth}\raggedleft
\(\gamma_N\)
\end{minipage} & \begin{minipage}[b]{\linewidth}\raggedleft
\(E_{\mathrm{final}}\)
\end{minipage} & \begin{minipage}[b]{\linewidth}\raggedleft
\(\lvert \Delta E \rvert\)
\end{minipage} & \begin{minipage}[b]{\linewidth}\raggedleft
Depth
\end{minipage} & \begin{minipage}[b]{\linewidth}\raggedleft
Parameters
\end{minipage} & \begin{minipage}[b]{\linewidth}\raggedright
Stop reason
\end{minipage} & \begin{minipage}[b]{\linewidth}\raggedleft
Prune kept
\end{minipage} \\
\midrule\noalign{}
\endhead
\bottomrule\noalign{}
\endlastfoot
\texttt{pareto\_lean\_l2} & \texttt{1.0} & \texttt{0.15872408236037028}
& \texttt{5.6178234644\ \textbackslash{}times\ 10\^{}\{-5\}} &
\texttt{14} & \texttt{25} & \texttt{drop\_plateau} & \texttt{4\ /\ 5} \\
\texttt{pareto\_lean\_l2} & \texttt{0.0} & \texttt{0.48696021436469644}
& \texttt{3.2829231024\ \textbackslash{}times\ 10\^{}\{-1\}} &
\texttt{14} & \texttt{23} & \texttt{drop\_plateau} & \texttt{4\ /\ 5} \\
\texttt{full\_meta} & \texttt{1.0} & \texttt{0.15872408236037047} &
\texttt{5.6178234644\ \textbackslash{}times\ 10\^{}\{-5\}} & \texttt{13}
& \texttt{28} & \texttt{drop\_plateau} & \texttt{5\ /\ 5} \\
\texttt{full\_meta} & \texttt{0.0} & \texttt{0.48699090951182056} &
\texttt{3.2832300539\ \textbackslash{}times\ 10\^{}\{-1\}} & \texttt{13}
& \texttt{22} & \texttt{drop\_plateau} & \texttt{5\ /\ 5} \\
\end{longtable}
}

If \[
R_{\mathrm{nov}}(\mathcal G)
=
\frac{\lvert \Delta E \rvert_{\gamma_N=0,\mathcal G}}{\lvert \Delta E \rvert_{\gamma_N=1,\mathcal G}},
\] then the two matched pool ratios are \[
R_{\mathrm{nov}}(\texttt{pareto\_lean\_l2})\approx 5.84376\times 10^3,
\qquad
R_{\mathrm{nov}}(\texttt{full\_meta})\approx 5.84431\times 10^3.
\]

So, on this matched HH benchmark, turning novelty weighting off worsens
the final filtered-energy error by roughly a factor of
\(5.8\times 10^3\) in both a lean pool and a heavy pool.

\subsection{20.2 Interpretation of the novelty
ablation}\label{interpretation-of-the-novelty-ablation}

Several points survive this comparison cleanly.

\begin{enumerate}
\def\labelenumi{\arabic{enumi}.}
\item
  \textbf{Novelty-on recovers the low-error line.} In both pools,
  \(\gamma_N=1\) reaches the same final error scale, \[
  \lvert \Delta E \rvert \approx 5.62\times 10^{-5}.
  \]
\item
  \textbf{Novelty-off collapses to a much worse terminal scaffold.} In
  both pools, \(\gamma_N=0\) ends near \[
  E_{\mathrm{final}}\approx 0.487,
  \qquad
  \lvert \Delta E \rvert \approx 3.28\times 10^{-1},
  \] far above the filtered exact target near \(0.158668\).
\item
  \textbf{This is not merely a depth effect.} The on/off pairs stop at
  the same logical depth within each pool, so the dominant difference is
  the operator-ranking trajectory rather than a simple ``more steps''
  explanation.
\item
  \textbf{This is not merely a pruning-on versus pruning-off effect.}
  Pruning stays active in all four runs. What changes is the upstream
  ranking rule. In particular, novelty telemetry is still assembled when
  \(\gamma_N=0\); what is ablated is the multiplicative use of novelty
  inside the selector score, not the underlying geometric information
  itself.
\end{enumerate}

These runs therefore support the claim that the reduced-tangent-space
novelty factor is doing real work inside the ADAPT ranking rule rather
than serving as passive telemetry.

\subsection{20.3 Pool-family redundancy}\label{pool-family-redundancy}

The same data story also sharpens the pool-reduction claim. For the
studied HH regime \[
L=2,\qquad n_{\mathrm{ph,max}}=1,
\] the heavy \texttt{full\_meta} pool has size \(46\), while the lean
\texttt{pareto\_lean\_l2} pool has size \(11\). Hence the pool-reduction
fraction is \[
\eta_{\mathrm{pool}}
=
1-\frac{11}{46}
=
\frac{35}{46}
\approx 0.7609.
\]

Equivalently, the lean pool removes about \(76\%\) of the heavy-pool
candidates while still retaining the same practical final energy scale,
\[
E_{\mathrm{final}}\approx 0.15872408,
\qquad
\lvert \Delta E \rvert \approx 5.62\times 10^{-5},
\] on the current-code lean rerun.

The generator classes discussed in this L=2 comparison may be
represented symbolically by \[
\begin{aligned}
A^{\mathrm{sing}}_{ia,\sigma}
&=
i\!\left(\hat c^\dagger_{a\sigma}\hat c_{i\sigma}-\hat c^\dagger_{i\sigma}\hat c_{a\sigma}\right),\\
A^{\mathrm{dbl}}_{ijab}
&=
i\!\left(\hat c^\dagger_a\hat c^\dagger_b\hat c_j\hat c_i-\mathrm{h.c.}\right),\\
G_i^{(\mathrm{cloud})}
&=
\tilde n_i\sum_{r\in\mathcal N(i)}\alpha_{ir}P_r,\\
G_{ij}^{(\mathrm{hd})}
&=
T_{ij}^{(+)}(P_i-P_j),\\
G_i^{(\mathrm{dd})}
&=
D_iP_i,\\
G_{ij}^{(\mathrm{od})}
&=
J_{ij}^{(-)}(x_i-x_j),\\
G_{ij}^{(2)}
&=
T_{ij}^{(+)}(x_i-x_j)^2,\\
x_i
&=
\hat b_i+\hat b_i^\dagger,
\qquad
P_i
=
i(\hat b_i^\dagger-\hat b_i).
\end{aligned}
\] On this \(L=2\) benchmark, the retained/useful classes were
\(A^{\mathrm{sing}}\), the \(P\)-type polaronic classes
\(G_i^{(\mathrm{cloud})}\), \(G_{ij}^{(\mathrm{hd})}\), and
\(G_i^{(\mathrm{dd})}\), together with the quadrature seeds. By
contrast, \(A^{\mathrm{dbl}}\), HVA, and the x-type / higher-order
channels \(G_{ij}^{(\mathrm{od})}\) and \(G_{ij}^{(2)}\) were not needed
to preserve the achieved error scale.

Empirically, the following heavy-pool families are absent from the lean
pool yet are not needed to retain that achieved accuracy scale on this
benchmark:

\begin{itemize}
\tightlist
\item
  lifted UCCSD doubles,
\item
  HH unit terms,
\item
  HVA layers,
\item
  \texttt{paop\_cloud\_x},
\item
  \texttt{paop\_disp},
\item
  \texttt{paop\_dbl},
\item
  \texttt{paop\_dbl\_x},
\item
  \texttt{paop\_curdrag},
\item
  \texttt{paop\_hop2}.
\end{itemize}

This should be read as a scoped data claim, not as a universal proof:
for the present HH \(L=2\), \(n_{\mathrm{ph,max}}=1\) target problem,
the heavy pool carries substantial operator-family redundancy relative
to the observed energy objective.

\subsubsection{\texorpdfstring{20.3.1 Weak-coupling \(L=2\) class-prune
redo}{20.3.1 Weak-coupling L=2 class-prune redo}}\label{weak-coupling-l2-class-prune-redo}

A second \(L=2\) class-prune line at weaker coupling sharpens the same
conclusion. For \[
L=2,\qquad t=1,\qquad U=0.5,\qquad \omega_0=1,\qquad g=0.2,\qquad n_{\mathrm{ph,max}}=1,
\] using binary bosons, blocked ordering, open boundaries, the compiled
backend path, full-window reoptimization, beam width \(B=3\), and the
direct \texttt{phase3\_v1} HH selector, the heavy parent
\texttt{full\_meta} run reached \[
E_{\mathrm{parent}}=-0.7763497696348323,
\qquad
\lvert\Delta E\rvert_{\mathrm{parent}}=2.7761269927\times 10^{-5},
\qquad
F_{\mathrm{parent}}=0.9999861889450502,
\] at logical depth \(d_{\mathrm{parent}}=7\).

From that parent, a class-first pruning loop removed one operator family
at a time and reran fresh ADAPT after every proposed removal. A removal
was accepted only if the rerun satisfied \[
\lvert\Delta E\rvert \le 5\times 10^{-4},
\qquad
F \ge F_{\mathrm{parent}}-10^{-4}.
\] The accepted removals were

\begin{itemize}
\tightlist
\item
  \texttt{hh\_termwise\_unit},
\item
  \texttt{hva\_layer},
\item
  \texttt{paop\_cloud\_x},
\item
  \texttt{paop\_curdrag},
\item
  \texttt{paop\_dbl},
\item
  \texttt{paop\_dbl\_p},
\item
  \texttt{paop\_dbl\_x},
\item
  \texttt{paop\_disp},
\item
  \texttt{paop\_hop2},
\item
  \texttt{paop\_hopdrag},
\item
  \texttt{hh\_termwise\_quadrature}.
\end{itemize}

Hence the final validated keep-set was simply \[
\mathcal G_{\mathrm{lean,weak}}
=
\{\texttt{uccsd\_sing},\ \texttt{uccsd\_dbl},\ \texttt{paop\_cloud\_p}\}.
\] A fresh rerun using only this class-filtered pool then produced \[
E_{\mathrm{lean}}=-0.7763497696348325,
\qquad
\lvert\Delta E\rvert_{\mathrm{lean}}=2.7761269927\times 10^{-5},
\qquad
F_{\mathrm{lean}}=0.9999861889437017,
\] at logical depth \(d_{\mathrm{lean}}=4\). Thus, to numerical
precision, \[
E_{\mathrm{lean}}\approx E_{\mathrm{parent}},
\qquad
\lvert\Delta E\rvert_{\mathrm{lean}}\approx \lvert\Delta E\rvert_{\mathrm{parent}},
\qquad
F_{\mathrm{lean}}\approx F_{\mathrm{parent}},
\] while the logical depth drops from \(7\) to \(4\).

A separate optimizer check on this same weak-coupling direct Phase-3
beam surface asked a narrower question: whether SPSA, still run
locally/noiselessly with the pre-transpile proxy-cost surface
(\texttt{phase3\_backend\_cost\_mode=proxy}) and still using the same
reduced class-filtered pool with full-window refits, could at least
reach the practical noise-threshold regime even if it did not exactly
recover the Powell basin. The best completed local SPSA artifact on that
surface was \[
E_{\mathrm{SPSA,local}}=-0.7763064877481185,
\qquad
\lvert\Delta E\rvert_{\mathrm{SPSA,local}}=7.1043156641\times 10^{-5},
\qquad
F_{\mathrm{SPSA,local}}=0.9999653825260001,
\] from the direct beam-enabled class-filtered run
\texttt{20260328\_hh\_l2\_u05\_g02\_local\_beam\_b3\_spsa\_case4\_direct\_m3200\_a0p1\_c0p02\_A5}.
This did not match the Powell/class-pruned reference value
\(2.7761269927\times 10^{-5}\), but it did enter the sub-\(10^{-4}\)
regime and therefore satisfied the more practical gate of crossing below
the anticipated noise floor. In that sense the local SPSA question at
weak coupling was answered positively: exact-basin recovery remained
unresolved, but threshold-level local recovery was achieved strongly
enough to justify a direct noisy/backend follow-on on the same
algorithmic surface.

A direct noisy transfer follow-on then kept the same reduced pool
\(\{\texttt{uccsd\_sing},\texttt{uccsd\_dbl},\texttt{paop\_cloud\_p}\}\),
the same beam/full-window \texttt{phase3\_v1} surface, and the same SPSA
recipe \((a,c,A,\texttt{maxiter})=(0.1,0.02,5.0,3200)\), but turned on
\texttt{backend\_scheduled} FakeMarrakesh scoring with
\texttt{expectation\_v1}, \texttt{512} shots, \texttt{1} repeat, and no
mitigation; the exact command is recorded at
\texttt{artifacts/agent\_runs/20260328\_hh\_l2\_u05\_g02\_phase3\_beam\_b3\_fake\_marrakesh\_spsa\_case4\_recipe\_attempt3\_supervised/logs/command.sh}.
During that supervised run, the best observed heartbeat reached \[
E_{\mathrm{SPSA,noisy,best}}=-0.7762934955401648,
\qquad
\lvert\Delta E\rvert_{\mathrm{SPSA,noisy,best}}=8.4035364595\times 10^{-5},
\] at depth \(95\), as preserved in
\texttt{artifacts/agent\_runs/20260328\_hh\_l2\_u05\_g02\_phase3\_beam\_b3\_fake\_marrakesh\_spsa\_case4\_recipe\_attempt3\_supervised/best\_so\_far\_snapshot.json}.
So, before final max-depth termination, the same reduced-pool direct
Phase-3 beam SPSA lane already crossed the practical noisy
sub-\(10^{-4}\) regime under transpiled FakeMarrakesh selection/scoring.

\subsubsection{20.3.2 Live four-run supervised status
block}\label{live-four-run-supervised-status-block}

\begin{Shaded}
\begin{Highlighting}[]
\NormalTok{Objective\textless{}As of March 28, 2026 at 10:59:12 CDT, all 4 noisy}
\NormalTok{  reduced{-}pool backend\_scheduled SPSA runs are still genuinely}
\NormalTok{  running with live child PIDs and no stderr in artifacts/}
\NormalTok{  agent\_runs/20260328\_hh\_l2\_u05\_g02\_phase3\_beam\_b3\_fake\_marrake}
\NormalTok{  sh\_spsa\_case4\_recipe\_attempt4\_dropstop/progress.json,}
\NormalTok{  artifacts/}
\NormalTok{  agent\_runs/20260328\_hh\_l2\_u05\_g02\_phase3\_beam\_b3\_fake\_marrake}
\NormalTok{  sh\_spsa\_case4\_recipe\_attempt5\_dropstop\_lifetimeoff/}
\NormalTok{  progress.json, artifacts/}
\NormalTok{  agent\_runs/20260328\_hh\_l2\_u05\_g02\_phase3\_beam\_b4\_c3\_k4\_fake\_m}
\NormalTok{  arrakesh\_spsa\_case4\_recipe\_attempt6\_dropstop\_lifetimeoff/}
\NormalTok{  sh\_spsa\_case4\_recipe\_attempt7\_dropstop\_transpile\_single/}
\NormalTok{  progress.json. Best observed noisy gaps so far are: attempt4}
\NormalTok{  3.3871153069520155e{-}05 at depth 76 in artifacts/}
\NormalTok{  agent\_runs/20260328\_hh\_l2\_u05\_g02\_phase3\_beam\_b3\_fake\_marrake}
\NormalTok{  sh\_spsa\_case4\_recipe\_attempt4\_dropstop/logs/stdout.log,}
\NormalTok{  attempt6 3.733262477279009e{-}05 at depth 49 in artifacts/}
\NormalTok{  agent\_runs/20260328\_hh\_l2\_u05\_g02\_phase3\_beam\_b4\_c3\_k4\_fake\_m}
\NormalTok{  arrakesh\_spsa\_case4\_recipe\_attempt6\_dropstop\_lifetimeoff/}
\NormalTok{  logs/stdout.log, attempt7 3.899366517123859e{-}05 at depth 67}
\NormalTok{  in artifacts/}
\NormalTok{  agent\_runs/20260328\_hh\_l2\_u05\_g02\_phase3\_beam\_b3\_fake\_marrake}
\NormalTok{  sh\_spsa\_case4\_recipe\_attempt7\_dropstop\_transpile\_single/logs/}
\NormalTok{  stdout.log, and attempt5 1.0855648531227224e{-}04 at depth 67}
\NormalTok{  in artifacts/}
\NormalTok{  agent\_runs/20260328\_hh\_l2\_u05\_g02\_phase3\_beam\_b3\_fake\_marrake}
\NormalTok{  sh\_spsa\_case4\_recipe\_attempt5\_dropstop\_lifetimeoff/logs/}
\NormalTok{  stdout.log.\textgreater{}}
\NormalTok{  Why/Intent\textless{}So the current result ordering is: baseline clean{-}}
\NormalTok{  stop attempt4 best, then wider{-}beam+lifetime{-}off attempt6,}
\NormalTok{  then transpile\_single\_v1 attempt7, with lifetime{-}off only}
\NormalTok{  attempt5 trailing. None has written a final JSON yet because}
\NormalTok{  the clean stop is intentionally gated by {-}{-}adapt{-}drop{-}min{-}}
\NormalTok{  depth 96, so even though three runs are already well below}
\NormalTok{  the practical 1e{-}4 threshold, the normal drop\_plateau exit}
\NormalTok{  cannot fire before depth 96.\textgreater{}}
\NormalTok{  Suggested Next step/how this fits into broader picture\textless{}I}
\NormalTok{  recommend we keep these running until the first one crosses}
\NormalTok{  depth 96 and exits normally, because that will give us the}
\NormalTok{  first clean final JSON on this surface without interrupting a}
\NormalTok{  good basin. If the ranking stays similar, attempt4 should}
\NormalTok{  finish first, while attempt6 and attempt7 already look}
\NormalTok{  scientifically competitive enough to justify keeping them}
\NormalTok{  alive as the main comparison axes.\textgreater{}}
\end{Highlighting}
\end{Shaded}

The rejected removals show which families remained indispensable under
this contract:

\begin{itemize}
\tightlist
\item
  removing \texttt{paop\_cloud\_p} gave
  \(\lvert\Delta E\rvert\approx 1.0813\times 10^{-2}\) and
  \(F\approx 0.996632\),
\item
  removing \texttt{uccsd\_sing} gave
  \(\lvert\Delta E\rvert\approx 2.19934\) and \(F\approx 0.214233\),
\item
  removing \texttt{uccsd\_dbl} gave
  \(\lvert\Delta E\rvert\approx 1.31029\times 10^{-2}\) and
  \(F\approx 0.996691\).
\end{itemize}

So at this weaker-coupling \(L=2\) point the validated lean pool is even
smaller than the earlier heavy/lean \(U=4\), \(g=0.5\) line: the heavy
search may still visit quadrature seeds, but the final class-level
support needed to preserve the achieved energy/fidelity scale is just \[
\texttt{uccsd\_sing}\;\cup\;\texttt{uccsd\_dbl}\;\cup\;\texttt{paop\_cloud\_p}.
\]

For the backend-facing follow-up at this same weak-coupling point, the
Heron/Marrakesh optimization should be done as a \textbf{transpile-only}
sweep rather than through the noisy staged wrapper: run
\texttt{python\ -m\ pipelines.hardcoded.adapt\_circuit\_cost} on the
promoted locked scaffold
\texttt{hh\_l2\_u05\_g02\_full\_meta\_class\_pruned\_lean4\_locked\_scaffold\_v1.json}
over
\texttt{\textbackslash{}texttt\{optimization\textbackslash{}\_level\}\textbackslash{}in\textbackslash{}\{1,2\textbackslash{}\}}
and
\texttt{\textbackslash{}texttt\{seed\textbackslash{}\_transpiler\}\textbackslash{}in\textbackslash{}\{0,\textbackslash{}dots,9\textbackslash{}\}},
then rank by compiled two-qubit count, depth, and size. On
\texttt{\textbackslash{}texttt\{FakeMarrakesh\}} the winning setting was
\texttt{\textbackslash{}texttt\{optimization\textbackslash{}\_level\}=2},
\texttt{\textbackslash{}texttt\{seed\textbackslash{}\_transpiler\}=5},
reducing the same 4-operator, 6-runtime-term scaffold from about
\texttt{25} two-qubit gates, depth \texttt{63}, size \texttt{109} to
\texttt{18} two-qubit gates, depth \texttt{43}, size \texttt{83} without
changing the physics or operator order.

\subsection{20.4 Relation to live Phase-3 selector/pruning
implementation}\label{relation-to-live-phase-3-selectorpruning-implementation}

The canonical Phase-3 target contract is the selector chain stated in
§§11.4.1--11.4.3 together with the compression-prune correction in
§11.5.1. Current live repo code still lags that contract in two visible
ways.

First, the canonical manuscript selector contract remains the strict
cumulative three-stage chain
\(\mathcal S_1\to\mathcal S_2\to\mathcal S_3\) stated in §§11.1 and
11.4.1--11.4.3. Historical repo variants may evaluate auxiliary cheap
quantities during staging, but the authoritative manuscript target is
the three-stage selector above. The remaining live selector discrepancy
relevant to the corrected canonical contract is the lifetime term, which
should use the useful horizon \[
H_t=\min\!\bigl(D_{\mathrm{left}}(t),\widehat N_{\mathrm{rem}}(t)\bigr)
\] rather than a raw remaining-depth / remaining-evaluations proxy.

Second, the live pruning path still ranks removals primarily by small
amplitude together with weaker prior proxy-benefit tie-breaks and
accepts removals through a standalone post-refit regression threshold.
The corrected canonical contract in §11.5.1 instead separates prune
pressure from prune safety: pressure opens plateau/terminal compression
opportunities, broadens the candidate set beyond small angles, nominates
by a cheap compensability prior, and then decides through a windowed
remove-refit recoverability ladder with shot escalation. Schur is the
local quadratic idealization of this recoverability target, not the
default static-ADAPT live estimator.

So the honest status note is: current code still trails the canonical
Phase-3 mathematics, but the intended direction of alignment is \[
\text{live implementation}\longrightarrow \text{corrected canonical contract of §§11.4.1--11.5.1},
\] not the reverse.

\subsection{20.5 Data provenance note}\label{data-provenance-note}

The novelty-ablation numbers in this chapter come from the fresh direct
\texttt{adapt\_pipeline} artifacts
\texttt{pi\_ablate\_pareto\_lean\_l2\_novelty\_on.json},
\texttt{pi\_ablate\_pareto\_lean\_l2\_novelty\_off.json},
\texttt{pi\_ablate\_full\_meta\_novelty\_on.json}, and
\texttt{pi\_ablate\_full\_meta\_novelty\_off.json}, all dated
\texttt{2026-03-24}. The heavy-versus-lean pool-family comparison is
anchored by the \texttt{2026-03-21} comparison bundle summarized in
\texttt{pareto\_lean\_l2\_vs\_heavy\_L2\_ecut1\_comparison.md} together
with the corresponding lean rerun artifact. The data are therefore
matched and current enough for PI-facing evidence.

The honesty caveat is the same one stated in the reporting notes: these
matched ablations share the same physics, seed, and selector family, but
they are not claimed to be byte-for-byte reproductions of every earlier
March 21/22 runtime-resolved default. That caveat does not weaken the
on/off comparisons above, because the resolved settings were shared
inside each matched pair.

\subsection{20.6 Machine-agent redo note for the successful L=2 pruning
session}\label{machine-agent-redo-note-for-the-successful-l2-pruning-session}

For the validated HH prune line, keep the problem fixed at \[
L=2,\qquad t=1,\qquad U=4,\qquad \omega_0=1,\qquad g=0.5,\qquad n_{\mathrm{ph,max}}=1,
\] with the saved fullhorse parent artifact as the starting point. A
machine agent can reproduce the session in three steps from repo root:

\begin{Shaded}
\begin{Highlighting}[]
\ExtensionTok{python} \AttributeTok{{-}m}\NormalTok{ pipelines.hardcoded.hh\_prune\_nighthawk }\DataTypeTok{\textbackslash{}}
  \AttributeTok{{-}{-}input{-}json}\NormalTok{ artifacts/json/hh\_backend\_adapt\_fullhorse\_powell\_20260322T194155Z\_fakenighthawk.json }\DataTypeTok{\textbackslash{}}
  \AttributeTok{{-}{-}mode}\NormalTok{ both }\DataTypeTok{\textbackslash{}}
  \AttributeTok{{-}{-}prune{-}threshold}\NormalTok{ 1e{-}4 }\DataTypeTok{\textbackslash{}}
  \AttributeTok{{-}{-}output{-}json}\NormalTok{ artifacts/json/hh\_prune\_nighthawk\_}\OperatorTok{\textless{}}\NormalTok{timestamp}\OperatorTok{\textgreater{}}\NormalTok{.json}

\ExtensionTok{python} \AttributeTok{{-}m}\NormalTok{ pipelines.hardcoded.hh\_prune\_nighthawk }\DataTypeTok{\textbackslash{}}
  \AttributeTok{{-}{-}mode}\NormalTok{ export\_fixed\_scaffolds }\DataTypeTok{\textbackslash{}}
  \AttributeTok{{-}{-}source{-}artifact{-}json}\NormalTok{ artifacts/json/hh\_prune\_nighthawk\_aggressive\_5op.json}

\ExtensionTok{python} \AttributeTok{{-}m}\NormalTok{ pipelines.hardcoded.adapt\_circuit\_cost }\DataTypeTok{\textbackslash{}}
  \AttributeTok{{-}{-}artifact{-}json}\NormalTok{ artifacts/json/hh\_prune\_nighthawk\_gate\_pruned\_7term.json }\DataTypeTok{\textbackslash{}}
  \AttributeTok{{-}{-}backend{-}name}\NormalTok{ ibm\_marrakesh }\DataTypeTok{\textbackslash{}}
  \AttributeTok{{-}{-}seed{-}transpiler}\NormalTok{ 7 }\AttributeTok{{-}{-}optimization{-}level}\NormalTok{ 1 }\DataTypeTok{\textbackslash{}}
  \AttributeTok{{-}{-}sweep{-}backends}
\end{Highlighting}
\end{Shaded}

The redo is not complete unless it records \textbf{fidelity as well as
energy}. For each surviving scaffold and exported fixed scaffold, save
at least: \texttt{abs\_delta\_e}, exact-state fidelity
(\texttt{exact\_state\_fidelity},
\texttt{local\_exact\_state\_fidelity}, or
\texttt{expected\_local\_exact\_state\_fidelity}), retained exyz labels,
backend, transpiler seed, optimization level, compiled
\texttt{count\_2q}, compiled depth, and compiled size. The validated
executable anchor is the \texttt{gate\_pruned\_7term} artifact, not
\texttt{circuit\_optimized\_7term}: the trustworthy target is \[
\lvert\Delta E\rvert\approx 4.23\times 10^{-4},\qquad F\approx 0.9998,
\] with Marrakesh compile cost near \texttt{25} two-qubit gates and
depth \texttt{63-75}. A leaner 6-term executable may reach about
\texttt{14} two-qubit gates and depth \texttt{48}, but its exact
regression is materially worse, near \[
\lvert\Delta E\rvert\approx 4.61\times 10^{-3}.
\]

\subsection{\texorpdfstring{20.7 Direct noisy
\texttt{backend\_scheduled} L=2 fastprobe
note}{20.7 Direct noisy backend\_scheduled L=2 fastprobe note}}\label{direct-noisy-backend_scheduled-l2-fastprobe-note}

For the same HH point \[
L=2,\qquad t=1,\qquad U=4,\qquad \omega_0=1,\qquad g=0.5,\qquad n_{\mathrm{ph,max}}=1,
\] a useful data line is the direct local noisy \texttt{phase3\_v1}
fastprobe on \texttt{FakeNighthawk} with \texttt{backend\_scheduled}
oracle evaluation, \texttt{512} shots, \texttt{2} repeats, and optional
\texttt{mthree} readout mitigation. The March 26 command family is
recorded in
\texttt{artifacts/runstate/hh\_phase3\_v1\_local\_noisy\_backend\_scheduled\_fastprobe\_retry2\_20260326\_cmd.sh}
and the stabilized readout follow-up
\texttt{artifacts/runstate/hh\_phase3\_v1\_local\_noisy\_backend\_scheduled\_fastprobe\_retry3\_20260326\_cmd.sh},
with the no-readout comparator in
\texttt{artifacts/runstate/hh\_phase3\_v1\_local\_noisy\_backend\_scheduled\_noreadout\_retry1\_20260326/}.

These probes were intentionally weak
(\texttt{\textbackslash{}texttt\{adapt\textbackslash{}\_max\textbackslash{}\_depth\}=1})
and were stopped after the first scout evaluation timing split was
captured, so they are debug-quality runtime evidence rather than
convergence-quality VQE runs. The first measured \texttt{plus}
evaluation on the readout+mthree path gave \[
T_{\mathrm{total,readout}}\approx 37.04\,\mathrm{s},
\qquad
T_{\mathrm{term}}\approx 26.39\,\mathrm{s},
\qquad
T_{\mathrm{calib}}\approx 10.64\,\mathrm{s},
\] with readout-apply time negligible at about \texttt{0.01s}. The
no-readout comparator gave \[
T_{\mathrm{total,no\,readout}}\approx 26.84\,\mathrm{s},
\qquad
T_{\mathrm{term}}\approx 26.83\,\mathrm{s}.
\] So on this \(L=2\) HH debug line the dominant cost is
backend-scheduled term execution itself, while readout mitigation
contributes a secondary first-hit warmup rather than the main runtime
burden. The practical lesson is that \texttt{backend\_scheduled} is
honest enough for a very small late shortlist or confirmation pass, but
too expensive to act as the broad pre-shortlist scout surface.

\subsection{20.8 Fixed-scaffold compile-control scout
note}\label{fixed-scaffold-compile-control-scout-note}

For the locked Marrakesh/Heron 6-term HH scaffold
\texttt{artifacts/json/hh\_marrakesh\_fixed\_scaffold\_6term\_drop\_eyezee\_20260323T171528Z.json},
the useful local compile-control test is a transpile-only scout on
\texttt{FakeMarrakesh} with fixed circuit/parameters,
\texttt{shots=4096}, \texttt{oracle\_repeats=8}, readout/mthree
mitigation, and grid
\((\texttt{optimization\_level},\texttt{seed\_transpiler})\in\{1,2\}\times\{0,\dots,4\}\),
ranking by \(\Delta E=E_{\mathrm{noisy}}-E_{\mathrm{ideal}}\) then by
compiled cost; among the saved \texttt{8/10} completed candidates in
\texttt{artifacts/agent\_runs/20260328\_direct\_fixed\_scaffold\_compile\_scout\_fullaccess\_attempt2\_timeout7200/json/20260328\_direct\_fixed\_scaffold\_compile\_scout\_fullaccess\_attempt2\_timeout7200.json},
the best-so-far setting was \texttt{opt1\_seed4} with
\(\Delta E\approx 9.9899\times 10^{-2}\), two-qubit count \texttt{14},
and depth \texttt{48}, whereas the completed lower-depth
\texttt{opt2\_*} candidates kept the same two-qubit count \texttt{14}
but had depth \texttt{38} and worse \(\Delta E\):
\texttt{opt2\_seed0\textbackslash{}approx\ 1.1726\textbackslash{}times\ 10\^{}\{-1\}},
\texttt{opt2\_seed1\textbackslash{}approx\ 1.2245\textbackslash{}times\ 10\^{}\{-1\}},
\texttt{opt2\_seed2\textbackslash{}approx\ 1.2029\textbackslash{}times\ 10\^{}\{-1\}},
so lower compiled depth alone did not improve noisy accuracy and the
current hardware-facing compile preset is
\texttt{(optimization\_level=1,\ seed\_transpiler=4)}.

As an \(L=2\) run-planning prior rather than a theorem, the earlier
local FakeMarrakesh ablations were mixed between readout-only and twirl
while DD usually underperformed, but the completed 12-cell
saved-\(\theta^\star\) mitigation matrix in
\texttt{artifacts/json/20260328\_fixed\_scaffold\_saved\_theta\_mitigation\_matrix\_attempt5\_timeout86400\_delta\_logs\_fixed\_scaffold\_saved\_theta\_mitigation\_matrix.json}
now supersedes that partial prior: the top three cells were
\texttt{opt2\_seed0\_\_zne\_on\_\_twirl\_dd} with
\(\Delta E\approx-6.6431\times 10^{-4}\) and stderr
\(\approx1.2322\times 10^{-2}\) at 14 two-qubit gates and depth 38,
\texttt{opt1\_seed4\_\_zne\_on\_\_twirl\_dd} with
\(\Delta E\approx8.2990\times 10^{-3}\) and stderr
\(\approx2.6104\times 10^{-2}\) at 14 two-qubit gates and depth 48, and
\texttt{opt2\_seed0\_\_zne\_on\_\_twirl} with
\(\Delta E\approx1.4807\times 10^{-2}\) and stderr
\(\approx1.1157\times 10^{-2}\) at 14 two-qubit gates and depth 38. For
a real-QPU first pass on this locked 6-term scaffold, the current best
local starting lane is therefore
\texttt{(optimization\_level=2,\ seed\_transpiler=0)} with
\texttt{ZNE\ on\ +\ twirl\ +\ DD}, with
\texttt{(optimization\_level=1,\ seed\_transpiler=4)} plus the same
mitigation stack as the nearest fallback if one wants to hedge against
ZNE-fit variance while staying on the same two-qubit count.

For the weak-coupling locked scaffold
\texttt{artifacts/json/useful/L2/hh\_l2\_u05\_g02\_full\_meta\_class\_pruned\_lean4\_locked\_scaffold\_v1.json},
the completed 4-cell readout+mthree shortlist in
\texttt{artifacts/json/20260329\_weak\_fixed\_scaffold\_top4\_mitigation\_readout\_opt2seed5\_fixed\_scaffold\_saved\_theta\_mitigation\_matrix.json}
had top two cells \texttt{opt2\_seed5\_\_zne\_on\_\_twirl\_dd} with
\(\Delta E\approx1.0800\times 10^{-2}\) and stderr
\(\approx4.7014\times 10^{-3}\), and
\texttt{opt2\_seed5\_\_zne\_on\_\_dd} with
\(\Delta E\approx1.6984\times 10^{-2}\) and stderr
\(\approx1.9163\times 10^{-3}\), both at 18 two-qubit gates and depth
43.

The later readout-on versus readout-off ablations showed that base
readout mitigation was helping rather than hurting on both
fixed-scaffold winner lines. For the strong-coupling local winner, the
same \texttt{opt2\_seed0\_\_zne\_on\_\_twirl\_dd} lane moved from
\(\Delta E\approx-6.6431\times 10^{-4}\) with readout on in
\texttt{artifacts/json/20260328\_fixed\_scaffold\_saved\_theta\_mitigation\_matrix\_attempt5\_timeout86400\_delta\_logs\_fixed\_scaffold\_saved\_theta\_mitigation\_matrix.json}
to \(\Delta E\approx5.4456\times 10^{-2}\) with readout off in
\texttt{artifacts/json/20260329\_strong\_fixed\_scaffold\_winner\_readout\_off\_ablation\_opt2seed0\_twirl\_dd\_retry2\_fixed\_scaffold\_saved\_theta\_mitigation\_matrix.json},
at the same compiled cost (14 two-qubit gates, depth 38). For the
weak-coupling 18-two-qubit compile \texttt{opt2\_seed5}, the analogous
readout-off follow-up in
\texttt{artifacts/json/20260330\_weak\_fixed\_scaffold\_top2\_readout\_off\_ablation\_opt2seed5\_fixed\_scaffold\_saved\_theta\_mitigation\_matrix.json}
gave \texttt{opt2\_seed5\_\_zne\_on\_\_twirl\_dd} at
\(\Delta E\approx5.6830\times 10^{-2}\) and
\texttt{opt2\_seed5\_\_zne\_on\_\_dd} at
\(\Delta E\approx5.8131\times 10^{-2}\), versus the readout-on values
\(\approx1.0800\times 10^{-2}\) and \(\approx1.6984\times 10^{-2}\)
respectively, so turning readout off worsened the best weak lanes by
about \(4.60\times10^{-2}\) and \(4.11\times10^{-2}\).

The surprising regime comparison is therefore not in the noiseless
scaffold bias but in the noisy response. The strong locked scaffold had
\(E_{\mathrm{ideal}}-E_{\mathrm{exact}}\approx4.6131\times10^{-3}\) from
\texttt{artifacts/json/hh\_marrakesh\_fixed\_scaffold\_6term\_drop\_eyezee\_20260323T171528Z.json},
while the weak locked scaffold had
\(E_{\mathrm{ideal}}-E_{\mathrm{exact}}\approx2.7761\times10^{-5}\) from
\texttt{artifacts/json/useful/L2/hh\_l2\_u05\_g02\_full\_meta\_class\_pruned\_lean4\_locked\_scaffold\_v1.json}.
Thus the weak scaffold was much closer to exact in the noiseless sense,
yet still produced a worse noisy \(\Delta E\) after the current
suppression stack; the deficit is therefore a noise-response /
mitigation-effect issue rather than a larger ideal-state approximation
error.

\subsection{20.9 Heron transpile Pareto front and noisy replay for fixed
L=2
scaffolds}\label{heron-transpile-pareto-front-and-noisy-replay-for-fixed-l2-scaffolds}

All scaffolds share \(L=2\), \(n_{\mathrm{ph,max}}=1\), binary bosons,
blocked ordering, open boundaries. Each was transpiled to
\texttt{FakeMarrakesh} over
\((\texttt{optimization\_level},\texttt{seed\_transpiler})\in\{1,2\}\times\{0,\dots,9\}\);
the table reports the best compiled two-qubit count per scaffold.
Term-pruned variants start from the 7-term locked scaffold and drop
Pauli rotation terms, then reoptimize \(\theta\) via Powell.

{\def\LTcaptype{none} % do not increment counter
\begin{longtable}[]{@{}
  >{\raggedright\arraybackslash}p{(\linewidth - 12\tabcolsep) * \real{0.1200}}
  >{\raggedleft\arraybackslash}p{(\linewidth - 12\tabcolsep) * \real{0.1600}}
  >{\raggedleft\arraybackslash}p{(\linewidth - 12\tabcolsep) * \real{0.1600}}
  >{\raggedleft\arraybackslash}p{(\linewidth - 12\tabcolsep) * \real{0.1600}}
  >{\raggedleft\arraybackslash}p{(\linewidth - 12\tabcolsep) * \real{0.1600}}
  >{\raggedright\arraybackslash}p{(\linewidth - 12\tabcolsep) * \real{0.1200}}
  >{\raggedright\arraybackslash}p{(\linewidth - 12\tabcolsep) * \real{0.1200}}@{}}
\toprule\noalign{}
\begin{minipage}[b]{\linewidth}\raggedright
Scaffold
\end{minipage} & \begin{minipage}[b]{\linewidth}\raggedleft
\(\lvert\Delta E\rvert\)
\end{minipage} & \begin{minipage}[b]{\linewidth}\raggedleft
CZ gates
\end{minipage} & \begin{minipage}[b]{\linewidth}\raggedleft
Depth
\end{minipage} & \begin{minipage}[b]{\linewidth}\raggedleft
\(F\)
\end{minipage} & \begin{minipage}[b]{\linewidth}\raggedright
Coupling
\end{minipage} & \begin{minipage}[b]{\linewidth}\raggedright
Source
\end{minipage} \\
\midrule\noalign{}
\endhead
\bottomrule\noalign{}
\endlastfoot
\texttt{weak\_lean4\_locked} (4 ops) & \(2.78\times 10^{-5}\) & 18 & 43
& 0.9999 & \(U{=}0.5,g{=}0.2\) & class-prune \\
\texttt{aggressive\_5op} (12 terms) & \(5.62\times 10^{-5}\) & 26 & 72 &
--- & \(U{=}4,g{=}0.5\) & op-prune \\
\texttt{gate\_pruned\_7term} (7 terms) & \(4.23\times 10^{-4}\) & 21 &
51 & 0.9998 & \(U{=}4,g{=}0.5\) & term-prune \\
\texttt{6term\_drop\_eyezee} (6 terms) & \(4.61\times 10^{-3}\) & 14 &
38 & 0.9985 & \(U{=}4,g{=}0.5\) & term-prune \\
\texttt{6term\_drop\_eyeeez} (6 terms) & \(4.61\times 10^{-3}\) & 14 &
56 & 0.9985 & \(U{=}4,g{=}0.5\) & term-prune \\
\texttt{5term\_drop\_both\_disp} (5 terms) & \(6.92\times 10^{-3}\) & 12
& 48 & 0.9986 & \(U{=}4,g{=}0.5\) & term-prune \\
\end{longtable}
}

The Pareto front over \((\lvert\Delta E\rvert,\,\text{CZ count})\)
contains two points when both coupling regimes are pooled: \[
\bigl(2.78\times 10^{-5},\;18\;\text{CZ}\bigr)_{\text{weak}},
\qquad
\bigl(4.61\times 10^{-3},\;14\;\text{CZ}\bigr)_{\text{strong}}.
\] Within strong coupling alone, the three-point front is
\(\{(5.62\times 10^{-5},26),\;(4.61\times 10^{-3},14),\;(6.92\times 10^{-3},12)\}\).
The dominated scaffolds (\texttt{pruned\_scaffold\_11op} at 51 CZ,
\texttt{ultra\_lean\_6op}/\texttt{readapt\_6op} at 38 CZ,
\texttt{gate\_pruned\_7term} at 21 CZ) offer no accuracy benefit
relative to cheaper alternatives.

\textbf{Noisy replay on \texttt{FakeMarrakesh}
(\texttt{backend\_scheduled})} for the strong-coupling Pareto-optimal
6-term scaffold (14 CZ, depth 38, \texttt{optimization\_level=2},
\texttt{seed\_transpiler=0}):

{\def\LTcaptype{none} % do not increment counter
\begin{longtable}[]{@{}
  >{\raggedright\arraybackslash}p{(\linewidth - 4\tabcolsep) * \real{0.2727}}
  >{\raggedleft\arraybackslash}p{(\linewidth - 4\tabcolsep) * \real{0.3636}}
  >{\raggedleft\arraybackslash}p{(\linewidth - 4\tabcolsep) * \real{0.3636}}@{}}
\toprule\noalign{}
\begin{minipage}[b]{\linewidth}\raggedright
Configuration
\end{minipage} & \begin{minipage}[b]{\linewidth}\raggedleft
\(E_{\mathrm{noisy}}\)
\end{minipage} & \begin{minipage}[b]{\linewidth}\raggedleft
\(E_{\mathrm{noisy}}-E_{\mathrm{exact}}\)
\end{minipage} \\
\midrule\noalign{}
\endhead
\bottomrule\noalign{}
\endlastfoot
Saved \(\theta^*\) + readout/mthree & \(0.2773\) & \(+0.119\) \\
Saved \(\theta^*\) + readout + gate twirling & \(0.2208\) &
\(+0.062\) \\
SPSA-128 best + readout + gate twirling & \(0.1368\) & \(-0.022\) \\
Saved \(\theta^*\) ideal (noiseless) & \(0.1633\) & \(+0.005\) \\
Exact GS & \(0.1587\) & \(0\) \\
\end{longtable}
}

Mitigation stack: readout/mthree plus local 2Q Pauli twirling, symmetry
\texttt{verify\_only}, no ZNE, no DD. SPSA: 128 iterations, 1024 shots,
1 repeat, seed 19. The SPSA-optimized noisy energy undershoots exact by
\(0.022\), indicating noise-induced bias below the variational floor;
the ideal energy at the SPSA-best \(\theta\) is \(0.189\), confirming
parameter drift. At the noiseless-optimal \(\theta^*\), gate twirling
halves the noise bias from \(0.119\) to \(0.062\).

Reproduction:

\begin{Shaded}
\begin{Highlighting}[]
\ExtensionTok{python} \AttributeTok{{-}m}\NormalTok{ pipelines.exact\_bench.heron\_pareto\_sweep }\AttributeTok{{-}{-}backend}\NormalTok{ FakeMarrakesh }\AttributeTok{{-}{-}coupling}\NormalTok{ both}
\ExtensionTok{python} \AttributeTok{{-}m}\NormalTok{ pipelines.hardcoded.hh\_staged\_noise }\DataTypeTok{\textbackslash{}}
  \AttributeTok{{-}{-}L}\NormalTok{ 2 }\AttributeTok{{-}{-}t}\NormalTok{ 1.0 }\AttributeTok{{-}{-}u}\NormalTok{ 4.0 }\AttributeTok{{-}{-}g{-}ep}\NormalTok{ 0.5 }\AttributeTok{{-}{-}omega0}\NormalTok{ 1.0 }\AttributeTok{{-}{-}n{-}ph{-}max}\NormalTok{ 1 }\DataTypeTok{\textbackslash{}}
  \AttributeTok{{-}{-}ordering}\NormalTok{ blocked }\AttributeTok{{-}{-}boundary}\NormalTok{ open }\AttributeTok{{-}{-}boson{-}encoding}\NormalTok{ binary }\DataTypeTok{\textbackslash{}}
  \AttributeTok{{-}{-}include{-}fixed{-}scaffold{-}noisy{-}replay} \DataTypeTok{\textbackslash{}}
  \AttributeTok{{-}{-}fixed{-}final{-}state{-}json}\NormalTok{ artifacts/json/hh\_marrakesh\_fixed\_scaffold\_6term\_drop\_eyezee\_20260323T171528Z.json }\DataTypeTok{\textbackslash{}}
  \AttributeTok{{-}{-}use{-}fake{-}backend} \AttributeTok{{-}{-}backend{-}name}\NormalTok{ FakeMarrakesh }\DataTypeTok{\textbackslash{}}
  \AttributeTok{{-}{-}fixed{-}scaffold{-}runtime{-}transpile{-}optimization{-}level}\NormalTok{ 2 }\DataTypeTok{\textbackslash{}}
  \AttributeTok{{-}{-}fixed{-}scaffold{-}runtime{-}seed{-}transpiler}\NormalTok{ 0 }\DataTypeTok{\textbackslash{}}
  \AttributeTok{{-}{-}final{-}method}\NormalTok{ SPSA }\AttributeTok{{-}{-}final{-}maxiter}\NormalTok{ 128 }\AttributeTok{{-}{-}final{-}seed}\NormalTok{ 19 }\DataTypeTok{\textbackslash{}}
  \AttributeTok{{-}{-}mitigation}\NormalTok{ readout }\AttributeTok{{-}{-}local{-}readout{-}strategy}\NormalTok{ mthree }\AttributeTok{{-}{-}local{-}gate{-}twirling} \DataTypeTok{\textbackslash{}}
  \AttributeTok{{-}{-}symmetry{-}mitigation{-}mode}\NormalTok{ verify\_only }\DataTypeTok{\textbackslash{}}
  \AttributeTok{{-}{-}shots}\NormalTok{ 1024 }\AttributeTok{{-}{-}oracle{-}repeats}\NormalTok{ 1 }\AttributeTok{{-}{-}noise{-}seed}\NormalTok{ 7 }\DataTypeTok{\textbackslash{}}
  \AttributeTok{{-}{-}smoke{-}test{-}intentionally{-}weak} \AttributeTok{{-}{-}skip{-}pdf}
\end{Highlighting}
\end{Shaded}

Artifacts:
\texttt{artifacts/json/heron\_pareto\_front\_20260327T223120Z.json},
\texttt{artifacts/json/hh\_6term\_pareto\_readout\_twirl\_fast\_20260328.json}.

\section{Appendix A. Drive, Exact Propagation, and Propagator
Semantics}\label{appendix-a.-drive-exact-propagation-and-propagator-semantics}

\subsection{A.1 Static drive labels and time-dependent
coefficients}\label{a.1-static-drive-labels-and-time-dependent-coefficients}

Let \(\mathcal L_{\mathrm{drive}}\) denote the fixed set of density
operators appearing in the drive. Then the drive always has the form \[
\hat H_{\mathrm{drive}}(t)=\sum_{\lambda\in\mathcal L_{\mathrm{drive}}} c_\lambda(t) P_\lambda,
\] with fixed labels \(P_\lambda\) and time-dependent coefficients
\(c_\lambda(t)\).

\subsection{A.2 Reference coefficient
map}\label{a.2-reference-coefficient-map}

For a site-density drive, \[
\Delta c[Z_{p(i,\uparrow)}](t)=\Delta c[Z_{p(i,\downarrow)}](t)=-\frac{1}{2}v_i(t),
\] and the identity shift is \[
\Delta c[I](t)=\sum_i v_i(t)
\] if one chooses to retain the scalar offset explicitly.

\subsection{A.3 Propagator surfaces}\label{a.3-propagator-surfaces}

A common propagator menu is:

\begin{itemize}
\tightlist
\item
  second-order Suzuki,
\item
  piecewise exact propagation.
\end{itemize}

\subsection{A.4 Reference propagator versus reported trajectory
propagator}\label{a.4-reference-propagator-versus-reported-trajectory-propagator}

Write the reported propagator as \(U_{\mathrm{traj}}(t)\) and the
high-accuracy reference propagator as \(U_{\mathrm{ref}}(t)\). These
need not be the same approximation.

\subsubsection{A.4.1 Static reference
branch}\label{a.4.1-static-reference-branch}

In the static case, one may take \[
U_{\mathrm{ref}}(t)=e^{-it\hat H}.
\]

\subsubsection{A.4.2 Drive-enabled reference
branch}\label{a.4.2-drive-enabled-reference-branch}

For time-dependent \(\hat H(t)\), a refined piecewise-constant reference
is \[
U_{\mathrm{ref}}(t_{n+1},t_n) \approx e^{-i\Delta t\hat H(t_n^*)},
\] with \(t_n^*\) sampled on a refined grid and the full reference
obtained as an ordered product over the refined slices.

\subsection{A.5 Filtered exact manifold and subspace
fidelity}\label{a.5-filtered-exact-manifold-and-subspace-fidelity}

Let \(\Pi_{\mathrm{sector}}\) denote the projector onto the target
symmetry sector. Then the filtered exact energy is \[
E_{\mathrm{exact,filtered}} = \min_{\psi\in\Pi_{\mathrm{sector}}\mathcal H} \langle \psi|\hat H|\psi\rangle,
\] and a sector-aware fidelity may be defined by \[
F_{\mathrm{sector}}(t) = \frac{|\langle \psi_{\mathrm{exact}}(t)|\Pi_{\mathrm{sector}}|\psi(t)\rangle|^2}{\langle \psi(t)|\Pi_{\mathrm{sector}}|\psi(t)\rangle}.
\]

\subsection{A.7 Branch-resolved observable
contracts}\label{a.7-branch-resolved-observable-contracts}

If several trajectory branches are carried simultaneously, each
observable should be indexed by branch label \(b\): \[
\mathcal O_b(t),
\qquad
b\in\{\mathrm{var},\mathrm{traj},\mathrm{ref},\mathrm{noisy},\mathrm{ideal}\}.
\] Typical observables are energy, density, doublon density, and
fidelity.

\subsection{A.8 State import and handoff
contract}\label{a.8-state-import-and-handoff-contract}

A portable state bundle can be written abstractly as \[
\mathcal B = (\Lambda,\theta,a,E,\mathcal C),
\] where \(\Lambda\) is the physical manifest, \(\theta\) are
variational parameters when available, \(a\) is an amplitude map for the
working state, \(E\) stores energy summaries, and \(\mathcal C\) is
continuation metadata.

\subsection{A.9 Highest-tier continuation
surfaces}\label{a.9-highest-tier-continuation-surfaces}

At the highest continuation tier, \(\mathcal C\) may include split-event
summaries, motif data, symmetry diagnostics, replay policy data, and
limited branch histories. These are naturally viewed as auxiliary
mathematical metadata rather than as part of the operator algebra
itself.

\section{Appendix B. Continuation and Handoff
Contract}\label{appendix-b.-continuation-and-handoff-contract}

\subsection{B.1 Handoff state bundle}\label{b.1-handoff-state-bundle}

A handoff bundle may contain:

\begin{itemize}
\tightlist
\item
  the parameter manifest \(\Lambda\),
\item
  the current energy summary,
\item
  the normalized statevector amplitudes,
\item
  an exact comparison target,
\item
  optional continuation metadata.
\end{itemize}

One convenient symbolic amplitude map is \[
\{b_{N_q-1}\cdots b_0 \mapsto (\Re a_b,\Im a_b)\}.
\]

\subsection{Appendix 17A-S. Archaic secant forecast and live-score
diagnostic,
non-canonical}\label{appendix-17a-s.-archaic-secant-forecast-and-live-score-diagnostic-non-canonical}

This appendix archives the secant/forecast machinery as a historical
diagnostic and possible oracle-facing benchmark route. It is not the
canonical QPU-facing time-dynamics controller, is not a route-selection
law, and must not enter append admission or default runtime policy. The
construction below is oracle-free when its rollout is generated only by
the controller's own one-step map, but it remains semantically close to
the older exact-state secant route and is therefore isolated from the
main algorithm.

\subsubsection{Appendix 17A-S.1 Archived predicted-conditional secant
and recursive local-projective
forecast}\label{appendix-17a-s.1-archived-predicted-conditional-secant-and-recursive-local-projective-forecast}

\[
\begin{aligned}
\theta_{k+1\mid k}^{\mathrm{pred}}
&:=\operatorname{Step}_{\iota_k^{\mathrm{pol}}}(\theta_k,t_k,\Delta\tau_k,\sigma_k;\mathcal S_k),
\qquad
|\psi_{k+1\mid k}^{\mathrm{pred}}\rangle
:=U(\theta_{k+1\mid k}^{\mathrm{pred}};\mathcal S_k)|\psi_{\mathrm{ref}}\rangle,\\[1mm]
|\widetilde\psi_{k+1\mid k}^{\mathrm{pred}}\rangle
&:=\exp\!\bigl(-i\arg\langle\psi_k,\psi_{k+1\mid k}^{\mathrm{pred}}\rangle\bigr)
|\psi_{k+1\mid k}^{\mathrm{pred}}\rangle,\\[1mm]
\delta_k^{\mathrm{sec}}
&:=Q_{\psi_k}\bigl(|\widetilde\psi_{k+1\mid k}^{\mathrm{pred}}\rangle-|\psi_k\rangle\bigr),\\[1mm]
\xi_k^{\mathrm{sec}}
&:=\Re\!\bigl(\bar T_k^{\dagger}\delta_k^{\mathrm{sec}}\bigr),
\qquad
v_k^{\mathrm{sec}}
:=\frac{1}{\Delta\tau_k}K_{k,\tau,\varepsilon}^{\oplus}\xi_k^{\mathrm{sec}},\\[1mm]
v_k^{\mathrm{sec}}
&\leftarrow
\min\!\left\{1,\frac{R_k}{\|v_k^{\mathrm{sec}}\|_{\mathsf M_k}}\right\}v_k^{\mathrm{sec}},
\qquad
\widehat v_k^{\mathrm{sec}}
:=\begin{cases}
v_k^{\mathrm{sec}}/\|v_k^{\mathrm{sec}}\|_{\mathsf M_k},&\|v_k^{\mathrm{sec}}\|_{\mathsf M_k}>\varepsilon_v,\\
\varnothing,&\text{otherwise},
\end{cases}\\[1mm]
\zeta_k^{\mathrm{int},r}
&:=\zeta^{\mathrm{int}}(\mathcal S_k^{+}(r),\theta_k\oplus_r0,t_k,\Delta\tau_k,\sigma_k),
\qquad
\iota_k^{\mathrm{pol},r}:=\Pi_{\mathrm{int}}(\zeta_k^{\mathrm{int},r}),\\[1mm]
(\mathcal S_k^{(u)},\theta_{k\mid k}^{(u)},\iota_u)
&:=\begin{cases}
(\mathcal S_k,\theta_k,\iota_k^{\mathrm{pol}}),
&u=(\mathrm{stay},\iota_k^{\mathrm{pol}}),\\[1mm]
(\mathcal S_k^{+}(r),\theta_k\oplus_r0,\iota_k^{\mathrm{pol},r}),
&u=(\mathrm{append},r),
\end{cases}\\[1mm]
|\psi_{k+q\mid k}^{(u)}\rangle
&:=U(\theta_{k+q\mid k}^{(u)};\mathcal S_k^{(u)})|\psi_{\mathrm{ref}}\rangle,
\qquad
q=0,\dots,N_k^{\mathrm{hor}}+1,\\[1mm]
T_{k+q\mid k}^{(u)}
&:=\left[\partial_{\vartheta_\mu}U(\theta_{k+q\mid k}^{(u)};\mathcal S_k^{(u)})|\psi_{\mathrm{ref}}\rangle\right]_\mu,\\[1mm]
\bar T_{k+q\mid k}^{(u)}
&:=Q_{\psi_{k+q\mid k}^{(u)}}\,T_{k+q\mid k}^{(u)},
\qquad
\bar b_{k+q\mid k}^{(u)}
:=Q_{\psi_{k+q\mid k}^{(u)}}\bigl(-iH_{k+q}^{\sharp}|\psi_{k+q\mid k}^{(u)}\rangle\bigr),\\[1mm]
G_{k+q\mid k}^{(u)}
&:=\Re\!\bigl((\bar T_{k+q\mid k}^{(u)})^{\dagger}\bar T_{k+q\mid k}^{(u)}\bigr),
\qquad
f_{k+q\mid k}^{(u)}
:=\Re\!\bigl((\bar T_{k+q\mid k}^{(u)})^{\dagger}\bar b_{k+q\mid k}^{(u)}\bigr),\\[1mm]
K_{k+q\mid k}^{(u)}
&:=G_{k+q\mid k}^{(u)}+\Lambda_{k+q\mid k}^{(u)},
\qquad
\dot\theta_{k+q\mid k}^{(u)}
:=(K_{k+q\mid k}^{(u)})_{\tau,\varepsilon}^{\oplus}f_{k+q\mid k}^{(u)},\\[1mm]
\zeta_{k+q\mid k}^{\mathrm{int},(u)}
&:=\zeta^{\mathrm{int}}(\mathcal S_k^{(u)},\theta_{k+q\mid k}^{(u)},t_{k+q},\Delta\tau_{k+q},\sigma_{k+q}),\\[1mm]
\iota_{k+q\mid k}^{(u)}
&:=\Pi_{\mathrm{int}}\!\left(\zeta_{k+q\mid k}^{\mathrm{int},(u)}\right),\\[1mm]
\theta_{k+q+1\mid k}^{(u)}
&:=\operatorname{Step}_{\iota_{k+q\mid k}^{(u)}}\!\left(
\theta_{k+q\mid k}^{(u)},t_{k+q},\Delta\tau_{k+q},\sigma_{k+q};\mathcal S_k^{(u)}
\right).
\end{aligned}
\]

The same-support rollout starts from \((\mathcal S_k,\theta_k)\) and
evolves by the scheduled one-step map. Append forecasts start from the
zero-initialized augmented ansatz support; the new block acquires
amplitude only through the first forecasted TDVP step, exactly as it
would after a committed append. Any nonzero append seed is an ablation
or trust-radius variant, not part of the canonical live action. Forecast
rollouts use the same declared one-step method recursively on predicted
conditional states. The recursion uses \(q=0,\ldots,N_k^{\mathrm{hor}}\)
for the step
\(\theta_{k+q\mid k}^{(u)}\mapsto\theta_{k+q+1\mid k}^{(u)}\), while the
score uses \(h=1,\ldots,N_k^{\mathrm{hor}}\) for forecasted checkpoints.
The forecast quantities \(\pi_{k+q\mid k}^{\mathrm{col},(u)}\),
\(\chi_{k+q\mid k}^{\mathrm{curv},(u)}\), and
\(e_{k+q\mid k}^{\mathrm{Eul},(u)}\) are archived forecast diagnostics,
not active Paper-II route conditions. The stay proposal family is a
diagnostic and forecast-comparison device. The secant proposal is
retained only as a secant of the controller's own predicted one-step
rollout state, computed from the same one-step map the controller would
use at this checkpoint.

\subsubsection{Appendix 17A-S.2 Archived live checkpoint forecast
score}\label{appendix-17a-s.2-archived-live-checkpoint-forecast-score}

Let \(N_k^{\mathrm{hor}}\) denote the forecast horizon length and
\(W_k:=\sum_{h=1}^{N_k^{\mathrm{hor}}}\omega_h\). The live score remains
oracle-free: it scores the self-predicted ansatz rollout by local
projective miss, conditioning, Fubini--Study trust-region excess, and
Ehrenfest consistency, with the benchmark reference path excluded from
the checkpoint
selector.\footnote{Safe implementation note: $\mathcal J_k^{\mathrm{live}}$ is a local, finite-horizon, self-predicted consistency score. It can prefer a rollout that is locally self-consistent yet globally wrong beyond the horizon, especially under rapidly varying drive, delayed append benefit, or poorly calibrated weights. The benchmark exact/reference path may be logged to diagnose this failure, but it must not re-enter the live selector as an oracle target.}

\[
\begin{aligned}
W_k
&:=\sum_{h=1}^{N_k^{\mathrm{hor}}}\omega_h,\\[1mm]
\epsilon_{\mathrm{expr},k+h\mid k}^{2,(u)}
&:=\operatorname{clip}\!\left(
\|\bar b_{k+h\mid k}^{(u)}\|_2^2
-(f_{k+h\mid k}^{(u)})^{\top}(G_{k+h\mid k}^{(u)})_{\tau}^{\dagger}f_{k+h\mid k}^{(u)},
0,
\|\bar b_{k+h\mid k}^{(u)}\|_2^2
\right),\\[1mm]
\rho_{\mathrm{miss},k+h\mid k}^{(u)}
&:=\rho_{\mathrm{expr},k+h\mid k}^{(u)}
:=\operatorname{clip}\!\left(
\frac{\epsilon_{\mathrm{expr},k+h\mid k}^{2,(u)}}{\|\bar b_{k+h\mid k}^{(u)}\|_2^2+\varepsilon},
0,
1
\right),\\[1mm]
\epsilon_{\mathrm{real},k+h\mid k}^{2,(u)}
&:=\|\bar T_{k+h\mid k}^{(u)}\dot\theta_{k+h\mid k}^{(u)}-\bar b_{k+h\mid k}^{(u)}\|_2^2,\\[1mm]
\rho_{\mathrm{real},k+h\mid k}^{(u)}
&:=\operatorname{clip}\!\left(
\frac{\epsilon_{\mathrm{real},k+h\mid k}^{2,(u)}}{\|\bar b_{k+h\mid k}^{(u)}\|_2^2+\varepsilon},\,0,\,\rho_{\max}
\right),\\[1mm]
\rho_{\mathrm{num},k+h\mid k}^{(u)}
&:=\bigl[\rho_{\mathrm{real},k+h\mid k}^{(u)}-\rho_{\mathrm{expr},k+h\mid k}^{(u)}\bigr]_+,\\[1mm]
\kappa_{\mathrm{eff},k+h\mid k}^{(u)}
&:=\kappa_{\mathrm{eff}}\!\left(K_{k+h\mid k}^{(u)}\right),\\[1mm]
\end{aligned}
\]

\[
\begin{aligned}
\theta_{h,i}^{(u)}
&:=\theta_{k+h\mid k}^{(u)}+\Delta\tau_{k+h}\sum_{j<i}a_{ij}v_{h,j}^{(u)},
\qquad
t_{h,i}:=t_{k+h}+c_i\Delta\tau_{k+h},\\[1mm]
v_{h,i}^{(u)}
&:=F_{\mathcal S_k^{(u)}}(\theta_{h,i}^{(u)},t_{h,i}),
\qquad
|\psi_{h,i}^{(u)}\rangle:=U(\theta_{h,i}^{(u)};\mathcal S_k^{(u)})|\psi_{\mathrm{ref}}\rangle,\\[1mm]
d_{k+h\mid k}^{\mathrm{FS}}(u)
&:=\begin{cases}
d_{k+h\mid k}^{\mathrm{FS,Euler}}(u),&\iota_{k+h\mid k}^{(u)}=\mathrm{Euler},\\[1mm]
d_{k+h\mid k}^{\mathrm{FS,RK4}}(u),&\iota_{k+h\mid k}^{(u)}=\mathrm{RK4},
\end{cases}\\[1mm]
d_{k+h\mid k}^{\mathrm{FS,Euler}}(u)
&:=\Delta\tau_{k+h}
\sqrt{(\dot\theta_{k+h\mid k}^{(u)})^{\top}G_{k+h\mid k}^{(u)}\dot\theta_{k+h\mid k}^{(u)}},\\[1mm]
d_{k+h\mid k}^{\mathrm{FS,RK4}}(u)
&:=\Delta\tau_{k+h}\sum_{i=1}^{4}b_i
\sqrt{(v_{h,i}^{(u)})^{\top}G_{\mathcal S_k^{(u)}}(\theta_{h,i}^{(u)})\,v_{h,i}^{(u)}},
\qquad
b=\left(\frac16,\frac13,\frac13,\frac16\right),\\[1mm]
R_{k+h}^{\mathrm{FS}}
&:=\text{local Fubini--Study trust-region radius},\\[1mm]
P_{\mathrm{disp},k+h}^{(u)}
&:=\left[
\frac{d_{k+h\mid k}^{\mathrm{FS}}(u)-R_{k+h}^{\mathrm{FS}}}
{R_{k+h}^{\mathrm{FS}}+\varepsilon}
\right]_+^2,\\[1mm]
\end{aligned}
\]

\[
\begin{aligned}
\hat N&:=\hat n_0+\hat n_1,
\qquad
\hat d:=\hat n_0-\hat n_1,
\qquad
\mathcal A_{\mathrm{Eh}}:=\{\hat N,\hat d,\hat D\},\\[1mm]
A[\phi]&:=\langle\phi|A|\phi\rangle,\quad A\in\mathcal A_{\mathrm{Eh}},\quad \partial_tA=0,\\[1mm]
\Delta A_{k+h\mid k}^{\mathrm{path}}(u)
&:=A[\psi_{k+h+1\mid k}^{(u)}]-A[\psi_{k+h\mid k}^{(u)}],\\[1mm]
\Delta A_{k+h\mid k}^{\mathrm{Eh}}(u)
&:=\begin{cases}
\Delta\tau_{k+h}\operatorname{Re}\!\left(
i\langle\psi_{k+h\mid k}^{(u)}|[H_{k+h}^{\sharp},A]|\psi_{k+h\mid k}^{(u)}\rangle
\right),
&\iota_{k+h\mid k}^{(u)}=\mathrm{Euler},\\[1mm]
\Delta\tau_{k+h}\displaystyle\sum_{i=1}^{4}b_i
\operatorname{Re}\!\left(
i\langle\psi_{h,i}^{(u)}|[H(t_{h,i}),A]|\psi_{h,i}^{(u)}\rangle
\right),
&\iota_{k+h\mid k}^{(u)}=\mathrm{RK4},
\end{cases}\\[1mm]
s_{A,k}^{\mathrm{loc}}
&:=\max\!\left\{\varepsilon_A^{\mathrm{floor}},
\operatorname{median}_{u,h}\left|\frac{\Delta A_{k+h\mid k}^{\mathrm{path}}(u)}{\Delta\tau_{k+h}}\right|,
\operatorname{median}_{u,h}\left|\frac{\Delta A_{k+h\mid k}^{\mathrm{Eh}}(u)}{\Delta\tau_{k+h}}\right|\right\},\\[1mm]
R_{k+h}^{A}(u)
&:=\frac{|\Delta A_{k+h\mid k}^{\mathrm{path}}(u)-\Delta A_{k+h\mid k}^{\mathrm{Eh}}(u)|}{\Delta\tau_{k+h}(s_{A,k}^{\mathrm{loc}}+\varepsilon)},\\[1mm]
\end{aligned}
\]

\[
\begin{aligned}
E_{k+h\mid k}^{(u),0}
&:=\langle\psi_{k+h\mid k}^{(u)}|H(t_{k+h})|\psi_{k+h\mid k}^{(u)}\rangle,\\[1mm]
\Delta E_{k+h\mid k}^{\mathrm{path}}(u)
&:=E_{k+h+1\mid k}^{(u),0}-E_{k+h\mid k}^{(u),0},\\[1mm]
\Delta E_{k+h\mid k}^{\mathrm{drv}}(u)
&:=\begin{cases}
\Delta\tau_{k+h}\operatorname{Re}\!\left\langle\partial_tH(t_{k+h}+\sigma_{k+h}\Delta\tau_{k+h})\right\rangle_{\psi_{k+h\mid k}^{(u)}},
&\iota_{k+h\mid k}^{(u)}=\mathrm{Euler},\\[1mm]
\Delta\tau_{k+h}\displaystyle\sum_{i=1}^{4}b_i
\operatorname{Re}\!\left\langle\psi_{h,i}^{(u)}\middle|\partial_tH(t_{h,i})\middle|\psi_{h,i}^{(u)}\right\rangle,
&\iota_{k+h\mid k}^{(u)}=\mathrm{RK4},
\end{cases}\\[1mm]
\partial_tH(t)
&\approx\frac{H(t+\Delta\tau_{k+h})-H(t)}{\Delta\tau_{k+h}}\quad\text{when no analytic drive derivative is available},\\[1mm]
s_{E,k}^{\mathrm{loc}}
&:=\max\!\left\{\varepsilon_E^{\mathrm{floor}},
\operatorname{median}_{u,h}\left|\frac{\Delta E_{k+h\mid k}^{\mathrm{path}}(u)}{\Delta\tau_{k+h}}\right|,
\operatorname{median}_{u,h}\left|\frac{\Delta E_{k+h\mid k}^{\mathrm{drv}}(u)}{\Delta\tau_{k+h}}\right|\right\},\\[1mm]
R_{k+h}^{E}(u)
&:=\frac{|\Delta E_{k+h\mid k}^{\mathrm{path}}(u)-\Delta E_{k+h\mid k}^{\mathrm{drv}}(u)|}{\Delta\tau_{k+h}(s_{E,k}^{\mathrm{loc}}+\varepsilon)},\\[1mm]
\end{aligned}
\]

\[
\begin{aligned}
g_k^{\mathrm{imm}}(u)
&:=\begin{cases}
0,&u\in\mathcal U_k^{\mathrm{stay}},\\[1mm]
\rho_{r,k}^{\mathrm{gain}},&u=(\mathrm{append},r),
\end{cases}\\[1mm]
\mathcal J_k^{\mathrm{live}}(u)
&:=-\lambda_{\mathrm{gain}}\,g_k^{\mathrm{imm}}(u)
+\frac{1}{W_k}\sum_{h=1}^{N_k^{\mathrm{hor}}}\omega_h\,\mathcal L_{k+h\mid k}(u),\\[1mm]
\mathcal L_{k+h\mid k}(u)
&:=\lambda_{\mathrm{miss}}\,\rho_{\mathrm{miss},k+h\mid k}^{(u)}
+\lambda_{\mathrm{num}}\,\rho_{\mathrm{num},k+h\mid k}^{(u)}\
&\quad+\lambda_{\mathrm{cond}}\log\!\bigl(1+\kappa_{\mathrm{eff},k+h\mid k}^{(u)}\bigr)
+\lambda_{\mathrm{disp}}\,P_{\mathrm{disp},k+h}^{(u)}\
&\quad+\sum_{A\in\mathcal A_{\mathrm{Eh}}}\lambda_A R_{k+h}^{A}(u)
+\lambda_E R_{k+h}^{E}(u),\\[1mm]
u_{k,\mathrm{stay}}^{\star}
&:=(\mathrm{stay},\iota_k^{\mathrm{pol}}),\\[1mm]
u_{k,\mathrm{app}}^{\star}
&:=\begin{cases}
\displaystyle \arg\min_{u\in\mathcal U_k^{\mathrm{app}}}\mathcal J_k^{\mathrm{live}}(u),&\mathcal U_k^{\mathrm{app}}\neq\varnothing,\\[2mm]
\varnothing,&\mathcal U_k^{\mathrm{app}}=\varnothing,
\end{cases}\\[1mm]
\mathcal J_k^{\mathrm{gap}}
&:=\begin{cases}
\mathcal J_k^{\mathrm{live}}(u_{k,\mathrm{stay}}^{\star})-\mathcal J_k^{\mathrm{live}}(u_{k,\mathrm{app}}^{\star}),&u_{k,\mathrm{app}}^{\star}\neq\varnothing,\\[1mm]
-\infty,&u_{k,\mathrm{app}}^{\star}=\varnothing,
\end{cases}\\[1mm]
\operatorname{Admit}_k(u_{k,\mathrm{app}}^{\star})=1
&\iff
u_{k,\mathrm{app}}^{\star}=(\mathrm{append},r_k^\star)\\
&\quad\wedge\ \operatorname{HighMiss}_k=1
\wedge\ \rho_{r_k^{\star},k}^{\mathrm{gain}}\ge\tau_{\mathrm{gain}}\\
&\quad\wedge\ S_{r_k^{\star},k}^{\mathrm{conf}}\ge\tau_{\mathrm{conf}}
\wedge\ \mathcal J_k^{\mathrm{gap}}\ge\delta_{\mathrm{live}}.
\end{aligned}
\]

The RK4 displacement uses the same stage velocities and stage metrics as
the RK4 step. The checkpoint comparator therefore uses self-predicted
stay and append rollouts together with the current local Schur-gain
gate. The scored forecast miss is the expressivity miss; the realized
forecast miss and \(\rho_{\mathrm{num}}\) penalize trajectories whose
tangent plane is expressive but whose damped implementation velocity
does not realize that expressivity. When
\(\|\bar b_{k+h\mid k}^{(u)}\|_2^2\) falls below a stationary-drift
floor, the normalized miss ratio is treated as unreliable; the
controller uses absolute residual size as a diagnostic and suppresses
structural adaptation unless miss growth persists over the forecast
horizon. The displacement term is a Fubini--Study trust-region excess
penalty, not a reward for slow motion: physical ray speed is penalized
only after it exceeds \(R_{k+h}^{\mathrm{FS}}\). For the \(L=2\) HH
controller, the Ehrenfest channel uses the independent observable basis
\((N,d,D)\) rather than double-counting the density-difference mode
through both \((n_0,n_1)\) and \(\rho\). The Ehrenfest identity used
above is
\(d\langle A\rangle/dt=i\langle[H,A]\rangle+\langle\partial_tA\rangle\),
specialized to the listed time-independent observables, and the
diagnostic quadrature follows the integrator order. In finite precision
the Ehrenfest and drive-energy right-hand sides are evaluated by taking
real parts. The energy channel compares the variational path's
endpoint-Hamiltonian energy change against
\(d\langle H(t)\rangle/dt=\langle\partial_tH(t)\rangle\) for exact
Schrödinger evolution under a time-dependent Hamiltonian. For
time-independent Hamiltonians, exact Schrödinger evolution conserves
energy; along the regularized projected path, nonzero energy drift is
interpreted as a composite diagnostic of projection error, numerical
time-discretization error, and regularization-induced damping. In
benchmark mode, this drift should decrease under ansatz-support
enrichment, timestep refinement, and damping reduction.

\subsection{B.2 Continuation block}\label{b.2-continuation-block}

The continuation payload can be treated abstractly as \[
\mathcal C = (m,\mathcal S,\mathcal M_{\mathrm{opt}},\Gamma,\Sigma,\mathcal L,\mathcal U,\Xi,\mathcal R,\mathcal P),
\] where:

\begin{itemize}
\tightlist
\item
  \(m\) is the continuation mode,
\item
  \(\mathcal S\) is inherited manifold-structure data,
\item
  \(\mathcal M_{\mathrm{opt}}\) is auxiliary optimizer state,
\item
  \(\Gamma\) is selected-generator metadata,
\item
  \(\Sigma\) is split-event data,
\item
  \(\mathcal L\) is a motif library,
\item
  \(\mathcal U\) is motif usage,
\item
  \(\Xi\) is symmetry information,
\item
  \(\mathcal R\) is rescue or recovery history,
\item
  \(\mathcal P\) is optional historical replay-policy metadata when such
  provenance is still being carried.
\end{itemize}

When controller forecast telemetry is persisted, quantities such as
\(D_{\mathrm{left}}(t)\), \(\widehat N_{\mathrm{rem}}(t)\), and \(H_t\)
belong naturally to this policy/provenance layer. They remain heuristic
controller state rather than operator-level observables.

\subsection{B.3 General handoff map}\label{b.3-general-handoff-map}

Let \(\mathcal H_{A\to B}\) be a handoff map from source manifold
\(\mathcal M_A\) to target manifold \(\mathcal M_B\): \[
\mathcal H_{A\to B}: (|\psi_A\rangle,\theta_A,\mathcal C_A) \mapsto (|\psi_B^{(0)}\rangle,\theta_B^{(0)},\mathcal C_B).
\] The map is intentionally abstract here: it records state,
coordinates, and continuation payload transfer between manifolds without
committing the main body to any obsolete replay-specific burn-in policy.

\subsection{B.4 Symmetry and tier-three payload
note}\label{b.4-symmetry-and-tier-three-payload-note}

Symmetry information belongs in the handoff if later stages will use it
for validation, postselection, or projector-renormalized estimation.
Mathematically, the carried object is a projector or a list of commuting
symmetry generators.

\subsection{B.5 Historical warm/refine chain
note}\label{b.5-historical-warmrefine-chain-note}

The obsolete warm-start \(\to\) refine \(\to\) ADAPT \(\to\) replay path
is kept only in Appendix A.2 and is not part of the active main-body
symbolic contract.

\subsection{Appendix 17A-H. Archived hybrid checkpoint-controller
formalism,
non-canonical}\label{appendix-17a-h.-archived-hybrid-checkpoint-controller-formalism-non-canonical}

This archived block preserves the pre-AP-McLachlan Chapter 17A draft for
provenance. Statements inside this block that call the hybrid checkpoint
law ``canonical'' are historical language from the old draft; the active
route authority is the AP-McLachlan law in Chapter 17A above. Do not use
this appendix to set route identity, defaults, append admission, prune
acceptance, or paper-facing claims without an explicit
diagnostic/ablation label.

The archived real-time law was piecewise smooth in physical time and
piecewise discrete in ansatz support. On each open checkpoint interval
the ansatz support is fixed and the state follows the projective
McLachlan projection of the sampled Schrödinger drift. At the checkpoint
itself one chooses exactly one structural action: keep the support,
append one new ordered block, or delete one old block. In what follows,
each object is defined once, the main quantities are expanded linearly
down to the operator/expectation-value primitives, and the condensed
shorthand is given only after the fully expanded form has already
appeared.

The archival driven \(L=2\) calibration printed near the front of this
manuscript is historical evidence only; it is not the route authority
and must not define code defaults. The formulas below define the
canonical Chapter 17A realtime law for a general ansatz support with
arbitrary checkpoint index \(k\) and arbitrary internal branch index.
The physical time is denoted \(t\) and checkpoint times are
\(0=\tau_0<\tau_1<\cdots<\tau_T=T\).

All variational generators \(A_{k,\mu}\) are Hermitian, the Hamiltonian
\(H(t)\) is Hermitian for every physical time, the reference state is
normalized, and units are chosen so that \(\hbar=1\). The driven
coefficients are assumed sufficiently smooth on each checkpoint interval
for the chosen time integrator. All expectation values are computed in
normalized variational states. We use the Fubini--Study ray-distance
convention \[
d_{\mathrm{FS}}(\psi,\phi):=\arccos\!\left(\operatorname{clip}(|\langle\psi|\phi\rangle|,0,1)\right),
\qquad
ds_{\mathrm{FS}}^2=d\theta^\top G(\theta)d\theta .
\] This convention fixes the scale of displacement penalties, prune
state-jump tests, append seed penalties when used, Euler embedded error
checks, and benchmark ray-distance diagnostics.

\subsubsection{Appendix 17A-H.1 Fixed-support projective McLachlan
geometry}\label{appendix-17a-h.1-fixed-support-projective-mclachlan-geometry}

The ansatz support is block-valued, while McLachlan geometry is
coordinate-valued. Let \(B_k:=|\mathcal S_k|\) be the number of logical
support blocks, let block \(b\) carry \(m_{k,b}\) real coordinates, let
\[
N_k:=\sum_{b=1}^{B_k}m_{k,b},
\qquad
I_{b,k}\subseteq\{1,\dots,N_k\}
\] be the active coordinate count and the coordinate set owned by block
\(b\), and let \(\mu,\nu\in\{1,\dots,N_k\}\) denote flattened coordinate
indices. Append and prune decisions act on persistent block labels;
\(G_k,f_k,K_k,\dot\theta_k\) act on flattened coordinate indices. The
Hilbert-space row index is \(a\in\{1,\dots,d\}\), where
\(d:=\dim\mathcal H\).

\[
\begin{aligned}
|\psi_k\rangle
&:=|\psi(\theta_k;\mathcal S_k)\rangle
 =U_k(\theta_k)|\psi_{\mathrm{ref}}\rangle,
\qquad
U_k(\theta):=U(\theta;\mathcal S_k),
\qquad
U(\theta;\mathcal S_k):=\prod_{\mu=N_k}^{1}e^{-i\theta_{\mu}A_{k,\mu}},\\[1mm]
U_{k,>\mu}
&:=\prod_{\ell=N_k}^{\mu+1}e^{-i\theta_{k,\ell}A_{k,\ell}},
\qquad
U_{k,<\mu}:=\prod_{\ell=\mu-1}^{1}e^{-i\theta_{k,\ell}A_{k,\ell}},
\qquad
\widetilde A_{k,\mu}:=U_{k,>\mu}A_{k,\mu}U_{k,>\mu}^{\dagger},\\[1mm]
\tau_{k,\mu}
&:=\partial_{\theta_{k,\mu}}|\psi_k\rangle\\
&=-iU_{k,>\mu}A_{k,\mu}e^{-i\theta_{k,\mu}A_{k,\mu}}U_{k,<\mu}|\psi_{\mathrm{ref}}\rangle
 =-i\widetilde A_{k,\mu}|\psi_k\rangle,\\[1mm]
\bar\tau_{k,\mu}
&:=\tau_{k,\mu}-\langle\psi_k,\tau_{k,\mu}\rangle|\psi_k\rangle
 =Q_{\psi_k}\tau_{k,\mu}
 =-i\bigl(\widetilde A_{k,\mu}-\langle\widetilde A_{k,\mu}\rangle_{\psi_k}\bigr)|\psi_k\rangle,
\qquad
Q_{\psi_k}:=I-|\psi_k\rangle\langle\psi_k|,\\[1mm]
\bar T_k
&:=\bigl[(\bar T_k)_{a\mu}\bigr]_{a=1,\dots,d}^{\,\mu=1,\dots,N_k}
 =\bigl[(\bar\tau_{k,\mu})_a\bigr]_{a=1,\dots,d}^{\,\mu=1,\dots,N_k}\\
&=\left[
\begin{array}{ccc}
-i\bigl((\widetilde A_{k,1}-\langle\widetilde A_{k,1}\rangle_{\psi_k})\psi_k\bigr)_1 & \cdots & -i\bigl((\widetilde A_{k,N_k}-\langle\widetilde A_{k,N_k}\rangle_{\psi_k})\psi_k\bigr)_1\\
\vdots & \ddots & \vdots\\
-i\bigl((\widetilde A_{k,1}-\langle\widetilde A_{k,1}\rangle_{\psi_k})\psi_k\bigr)_d & \cdots & -i\bigl((\widetilde A_{k,N_k}-\langle\widetilde A_{k,N_k}\rangle_{\psi_k})\psi_k\bigr)_d
\end{array}
\right],\\[1mm]
\bar b_k
&:=Q_{\psi_k}\bigl(-iH_k^{\sharp}|\psi_k\rangle\bigr)
 =-i\bigl(H_k^{\sharp}-E_{\psi_k,k}^{\sharp}\bigr)|\psi_k\rangle,
\qquad
E_{\psi_k,k}^{\sharp}:=\langle\psi_k|H_k^{\sharp}|\psi_k\rangle,\\[1mm]
&=\left[
\begin{array}{c}
\bigl(-i(H_k^{\sharp}-E_{\psi_k,k}^{\sharp})\psi_k\bigr)_1\\
\vdots\\
\bigl(-i(H_k^{\sharp}-E_{\psi_k,k}^{\sharp})\psi_k\bigr)_d
\end{array}
\right].
\end{aligned}
\]

The metric, force, regularized system, quadratic fit, and projected
velocity are then defined in one
chain:\footnote{Core intuition: the variational state $|\psi(\theta)\rangle$, with $\theta=(\theta_1,\ldots,\theta_p)\in\mathbb R^p$, lives on a variational manifold. The coordinates $\theta_i$ are not themselves physical; what is physical is the Hilbert-space motion produced when $\theta$ changes. The tangent map sends a real parameter direction $u$ to the state tangent $\bar T_k u=\sum_\mu u_\mu\bar\tau_{k,\mu}$. Therefore $u^\top G_k u=u^\top\Re(\bar T_k^\dagger\bar T_k)u=\|\bar T_k u\|_2^2$. Thus $G_k$ is not merely a numerical array: it is the local metric tensor of the ansatz manifold, measuring how strongly a coordinate direction moves the physical ray. A direction with very small $u^\top G_k u$ changes parameters while barely moving the state; such directions are redundant, gauge-like, flat, or numerically unresolved. Naively solving $K\dot\theta=f$ can divide by these tiny metric strengths and manufacture huge coordinate velocities from almost no physical tangent signal.}

\[
\begin{aligned}
G_k
&:=\bigl[(G_k)_{\mu\nu}\bigr]_{\mu,\nu=1}^{N_k}
 =\Re(\bar T_k^{\dagger}\bar T_k)
 =\left[\Re\!\left(\sum_{a=1}^{d}(\bar T_k)_{a\mu}^{*}(\bar T_k)_{a\nu}\right)\right]_{\mu,\nu=1}^{N_k}\\
&=\left[\Re\!\Bigl(\langle\psi_k|\widetilde A_{k,\mu}\widetilde A_{k,\nu}|\psi_k\rangle-\langle\widetilde A_{k,\mu}\rangle_{\psi_k}\langle\widetilde A_{k,\nu}\rangle_{\psi_k}\Bigr)\right]_{\mu,\nu=1}^{N_k},\\[1mm]
f_k
&:=\bigl[(f_k)_\mu\bigr]_{\mu=1}^{N_k}
 =\Re(\bar T_k^{\dagger}\bar b_k)
 =\left[\Re\!\left(\sum_{a=1}^{d}(\bar T_k)_{a\mu}^{*}(\bar b_k)_a\right)\right]_{\mu=1}^{N_k}\\
&=\left[\Re\!\Bigl(\langle\psi_k|\widetilde A_{k,\mu}H_k^{\sharp}|\psi_k\rangle-\langle\widetilde A_{k,\mu}\rangle_{\psi_k}E_{\psi_k,k}^{\sharp}\Bigr)\right]_{\mu=1}^{N_k},\\[1mm]
K_k
&:=\bigl[(K_k)_{\mu\nu}\bigr]_{\mu,\nu=1}^{N_k}
 =G_k+\Lambda_k
 =\left[(G_k)_{\mu\nu}+(\Lambda_k)_{\mu\nu}\right]_{\mu,\nu=1}^{N_k}\\
&=\left[\Re\!\Bigl(\langle\psi_k|\widetilde A_{k,\mu}\widetilde A_{k,\nu}|\psi_k\rangle-\langle\widetilde A_{k,\mu}\rangle_{\psi_k}\langle\widetilde A_{k,\nu}\rangle_{\psi_k}\Bigr)+(\Lambda_k)_{\mu\nu}\right]_{\mu,\nu=1}^{N_k},\\[1mm]
\mathfrak m_k(v)
&:=\|\bar T_kv-\bar b_k\|_2^2+v^{\top}\Lambda_kv\\
&=v^{\top}K_kv-2v^{\top}f_k+\operatorname{Var}_{\psi_k}(H_k^{\sharp}),\\[1mm]
\dot\theta_k^{\mathrm{LS}}
&\in \operatorname*{arg\,min}_{v\in\mathbb R^{N_k}}\mathfrak m_k(v),\\
K_k\dot\theta_k^{\mathrm{LS}}
&=f_k,\\[1mm]
\dot\theta_k
&:=K_{k,\tau,\varepsilon}^{\oplus}f_k\\
&=(G_k+\Lambda_k)_{\tau,\varepsilon}^{\oplus}f_k,\\[1mm]
\operatorname{Var}_{\psi_k}(H_k^{\sharp})
&:=\|\bar b_k\|_2^2
 =\langle\psi_k|(H_k^{\sharp}-E_{\psi_k,k}^{\sharp})^2|\psi_k\rangle.
\end{aligned}
\]

Equivalently, the damped implementation velocity is the stabilized
least-squares solve \(K_{k,\tau,\varepsilon}^{\oplus}f_k\), while
\(\dot\theta_k^{\mathrm{LS}}\) denotes the formal normal-equation
solution
\(K_k\dot\theta_k^{\mathrm{LS}}=f_k\).\footnote{McLachlan derivation, in condensed inline form: for real parameters, $\frac{d}{dt}|\psi(\theta(t))\rangle=\sum_j\dot\theta_j|\partial_j\psi\rangle$, while the exact Schrödinger velocity is $|b\rangle=-iH|\psi\rangle$, or horizontally $\bar b=Q_\psi|b\rangle$. McLachlan chooses $\dot\theta$ to minimize $\mathcal R(\dot\theta)=\|\sum_j\dot\theta_j|\partial_j\psi\rangle-|b\rangle\|^2$. If $|a(\dot\theta)\rangle=\sum_j\dot\theta_j|\partial_j\psi\rangle$, then $\mathcal R=\langle a-b|a-b\rangle=\langle a|a\rangle-\langle a|b\rangle-\langle b|a\rangle+\langle b|b\rangle$. Since $\dot\theta_j\in\mathbb R$, $\langle a|a\rangle=\sum_{j,k}\dot\theta_j\dot\theta_k\langle\partial_j\psi|\partial_k\psi\rangle$, and $\langle a|b\rangle+\langle b|a\rangle=2\sum_j\dot\theta_j\Re\langle\partial_j\psi|b\rangle$. Defining $K_{jk}:=\Re\langle\partial_j\psi|\partial_k\psi\rangle$ and $f_j:=\Re\langle\partial_j\psi|b\rangle=\Re\langle\partial_j\psi|-iH|\psi\rangle$, one obtains $\mathcal R(\dot\theta)=\sum_{j,k}\dot\theta_jK_{jk}\dot\theta_k-2\sum_j\dot\theta_jf_j+\langle b|b\rangle$. Differentiating with respect to $\dot\theta_i$ gives $\partial_{\dot\theta_i}\sum_{j,k}\dot\theta_jK_{jk}\dot\theta_k=\sum_{j,k}(\delta_{ij}K_{jk}\dot\theta_k+\dot\theta_jK_{jk}\delta_{ik})=\sum_kK_{ik}\dot\theta_k+\sum_j\dot\theta_jK_{ji}=2\sum_kK_{ik}\dot\theta_k$, using $K_{ji}=K_{ij}$, while $\partial_{\dot\theta_i}[-2\sum_j\dot\theta_jf_j]=-2f_i$ and $\partial_{\dot\theta_i}\langle b|b\rangle=0$. Setting the derivative to zero gives $2\sum_kK_{ik}\dot\theta_k-2f_i=0$, hence $\sum_kK_{ik}\dot\theta_k=f_i$, i.e. $K\dot\theta=f$.}

The spectral solve convention separates the algebraic inverse used in
exact residual identities from the damped implementation inverse used
for propagation. Here \(A\) denotes a measured real matrix intended to
represent a symmetric geometric object such as \(K_k\), \(G_k\), or a
Schur block. Exact arithmetic gives symmetry, but finite-shot,
finite-difference, and roundoff estimates can leave small antisymmetric
contamination. The controller therefore first replaces \(A\) by its
symmetric part, \[
A_s:=\frac12(A+A^\top),
\qquad
(A_s)_{ij}=\frac12(A_{ij}+A_{ji}).
\] This averages the two measured entries that should coincide. The
discarded skew part \(A_a:=\frac12(A-A^\top)\) cannot affect any real
quadratic least-squares term, since \(x^\top A_a x=0\) for real \(x\).
After this replacement, \(A_s=A_s^\top\), so the real spectral theorem
applies: there is an orthogonal matrix \(V\) and real eigenvalues
\(s_i\) such that \[
A_s=V\operatorname{diag}(s_1,\ldots,s_N)V^\top,
\qquad
V^\top V=VV^\top=I.
\]
\footnote{Symmetry and spectrum: in exact arithmetic, $K_{ij}=\Re\langle\partial_i\psi|\partial_j\psi\rangle$ and $K_{ji}=\Re\langle\partial_j\psi|\partial_i\psi\rangle=\Re(\langle\partial_i\psi|\partial_j\psi\rangle^*)=K_{ij}$, so $K=K^\top$. Measured, finite-shot, finite-difference, or roundoff estimates may give $K\ne K^\top$, so the real geometric object is restored by $K_s=\frac12(K+K^\top)$. The antisymmetric part $A=\frac12(K-K^\top)$ does not contribute to the real least-squares quadratic: for real $x$, $x^\top A x=(x^\top A x)^\top=x^\top A^\top x=x^\top(-A)x=-x^\top A x$, hence $x^\top A x=0$. Keeping the antisymmetric part only corrupts the spectral analysis. Since $K_s=K_s^\top$, the real spectral theorem gives $K_s=V\operatorname{diag}(s_1,\ldots,s_p)V^\top$, with $V^\top V=VV^\top=I$. Writing $v_a$ for the $a$-th column of $V$, $K_s=\sum_a s_a v_av_a^\top$. The entrywise identity is $(K_s)_{ij}=(VDV^\top)_{ij}=\sum_{a,b}V_{ia}D_{ab}V_{jb}=\sum_{a,b}V_{ia}s_a\delta_{ab}V_{jb}=\sum_a s_a(v_a)_i(v_a)_j$. Each $v_a$ is a parameter-space mode, and each $s_a$ is its tangent strength.}
We then overwrite \(A\leftarrow A_s\) and project measurement-induced
negative eigenvalues to zero unless the negative curvature exceeds a
rejection tolerance: \[
\begin{aligned}
A&\leftarrow \frac12(A+A^\top),\qquad
A=V\operatorname{diag}(s_1,\ldots,s_N)V^\top,\qquad s_1\ge\cdots\ge s_N\ge 0,\\[1mm]
A^{\dagger}
&:=V\operatorname{diag}\!\left(
\frac{\mathbf 1[s_i>0]}{s_i}
\right)_{i=1}^{N}V^\top,\\[1mm]
A_{\tau}^{\dagger}
&:=V\operatorname{diag}\!\left(
\frac{\mathbf 1[s_i\ge \tau_{\mathrm{pinv}}s_1]}{s_i}
\right)_{i=1}^{N}V^\top,\\[1mm]
A_{\tau,\varepsilon}^{\oplus}
&:=V\operatorname{diag}\!\left(
\frac{\mathbf 1[s_i\ge \tau_{\mathrm{pinv}}s_1]}{s_i+\varepsilon_{\mathrm{solve}}}
\right)_{i=1}^{N}V^\top,\\[1mm]
\kappa_{\mathrm{eff}}(A)
&:=\frac{s_{\max}(A)}{\max\{s_{\min}^{\mathrm{supp}}(A),\varepsilon_\kappa\}}.
\end{aligned}
\] The Moore--Penrose inverse \(A^\dagger\) is reserved for exact
algebraic statements; \(A_\tau^\dagger\) is the undamped supported
inverse used in structural expressivity identities; and
\(A_{\tau,\varepsilon}^{\oplus}\) is the damped supported implementation
inverse used for propagated velocities. The dense classical eigensolve
scales as \(\mathcal O(N_k^3)\) in the number of active variational
parameters, not in
\(\dim\mathcal H\).\footnote{Eigenbasis solve and cutoff: solving $K_s\dot\theta=f$ with $K_s=VDV^\top$, set $\dot\theta=Vy$ and $g=V^\top f$, so $VDV^\top Vy=f$, $VDy=f$, $VDy=Vg$, and therefore $Dy=g$. Componentwise, $s_ay_a=g_a$, with $g_a=(V^\top f)_a=\sum_iV_{ia}f_i=v_a^\top f$. Hence $y_a=(v_a^\top f)/s_a$ when $s_a\ne0$, and $\dot\theta=\sum_a((v_a^\top f)/s_a)v_a$. This is the exact place instability enters: division by tiny $s_a$. Geometrically, for the ansatz map $\Psi:\theta\mapsto|\psi(\theta)\rangle$, the tangent map is its differential $T_\theta\Psi$; for any candidate parameter velocity $u\in T_\theta\mathbb R^{N_k}$, with coordinates along the $\theta_j$ axes, write $Tu:=(T_\theta\Psi)u=\sum_j u_j|\partial_j\psi\rangle$, so $u^\top Ku=\Re\langle Tu|Tu\rangle$ (or with $G$ before regularization). Thus this quadratic is the state-space size induced by candidate parameter motion, not an extra controller score. For an eigenvector $v_a$ with $v_a^\top v_a=1$, $s_a=v_a^\top Kv_a=\Re\langle Tv_a|Tv_a\rangle$. Equivalently, if $|\partial_{v_a}\psi\rangle:=\sum_j(v_a)_j|\partial_j\psi\rangle$, then $s_a=\Re\langle\partial_{v_a}\psi|\partial_{v_a}\psi\rangle$. Thus $s_a\approx0$ means that parameter direction $v_a$ barely moves the quantum state, so inverting along it is physically meaningless. Algorithmically, the eigenmode strengths $s_a$ determine support retention, damping, conditioning diagnostics, and the Schur geometry used for append/prune decisions. The thresholded inverse keeps $I_\tau=\{a:s_a\ge\tau_{\mathrm{pinv}}s_{\max}\}$, sets $\chi_a=\mathbf 1[a\in I_\tau]=\mathbf 1[s_a\ge\tau_{\mathrm{pinv}}s_{\max}]$, and gives $K_\tau^\dagger=V\operatorname{diag}(\chi_a/s_a)V^\top=\sum_a(\chi_a/s_a)v_av_a^\top=\sum_{a\in I_\tau}(1/s_a)v_av_a^\top$. The implemented damped inverse $K_{\tau,\varepsilon}^{\oplus}$ uses denominator $s_a+\varepsilon_{\mathrm{solve}}$ on retained modes. Therefore $\dot\theta_\tau=K_\tau^\dagger f=\sum_a\chi_a((v_a^\top f)/s_a)v_a=\sum_{a\in I_\tau}((v_a^\top f)/s_a)v_a$ for exact supported algebra, while propagation uses $K_{\tau,\varepsilon}^{\oplus}f$. The retained-mode projector is $P_\tau=V\operatorname{diag}(\chi_a)V^\top=\sum_{a\in I_\tau}v_av_a^\top$, and $K_\tau^\dagger K_s=V\operatorname{diag}(\chi_a/s_a)V^\top V\operatorname{diag}(s_a)V^\top=V\operatorname{diag}(\chi_a)V^\top=P_\tau$. Thus the cutoff means: project onto trusted modes, then invert trusted modes only. The slogan is: do not invert coordinates; invert only physically resolved tangent modes.}
If no eigenvalue is retained, the solve returns zero velocity and flags
the local solve as
unsupported.\footnote{Redundancy example: suppose $|\psi(\theta_1,\theta_2)\rangle=|\varphi(\theta_1+\theta_2)\rangle$ and $\phi=\theta_1+\theta_2$. Then by the chain rule $|\partial_1\psi\rangle=|\partial_\phi\varphi\rangle$ and $|\partial_2\psi\rangle=|\partial_\phi\varphi\rangle$. With $c:=\Re\langle\partial_\phi\varphi|\partial_\phi\varphi\rangle$, the metric entries are $K_{11}=c$, $K_{12}=c$, $K_{21}=c$, and $K_{22}=c$, hence $K=c\left(\begin{smallmatrix}1&1\\1&1\end{smallmatrix}\right)$. The mode $v_+=(1,1)^\top/\sqrt2$ obeys $Kv_+=c\left(\begin{smallmatrix}1&1\\1&1\end{smallmatrix}\right)(1,1)^\top/\sqrt2=(c/\sqrt2)(2,2)^\top=2c\,v_+$, so $s_+=2c$. This is the real physical direction: increase both parameters and change $\phi$. The mode $v_-=(1,-1)^\top/\sqrt2$ obeys $Kv_-=c\left(\begin{smallmatrix}1&1\\1&1\end{smallmatrix}\right)(1,-1)^\top/\sqrt2=(c/\sqrt2)(0,0)^\top=0\,v_-$, so $s_-=0$. This is a null coordinate direction: increase $\theta_1$, decrease $\theta_2$, keep $\theta_1+\theta_2$ fixed, and move nowhere in Hilbert space. The singularity is therefore not a failure of physics; it reveals redundant coordinates. Spectral cutoff keeps $v_+$ and discards $v_-$.}\footnote{Condition-number meaning: if $s_{\min}^{\mathrm{supp}}:=\min_{a\in I_\tau}s_a$, then $\kappa_{\mathrm{eff}}(K)=s_{\max}(K)/\max\{s_{\min}^{\mathrm{supp}}(K),\varepsilon_\kappa\}$ measures how uneven the retained tangent geometry remains. Perturbing $f$ to $f+\delta f$ gives $\delta\dot\theta_\tau=K_\tau^\dagger\delta f=\sum_{a\in I_\tau}((v_a^\top\delta f)/s_a)v_a$ for the exact supported solve. Since the $v_a$ are orthonormal, $\|\delta\dot\theta_\tau\|^2=\sum_{a\in I_\tau}((v_a^\top\delta f)/s_a)^2\le(1/(s_{\min}^{\mathrm{supp}})^2)\sum_{a\in I_\tau}(v_a^\top\delta f)^2=(1/(s_{\min}^{\mathrm{supp}})^2)\|P_\tau\delta f\|^2\le(1/(s_{\min}^{\mathrm{supp}})^2)\|\delta f\|^2$. Thus $\|\delta\dot\theta_\tau\|\le\|\delta f\|/s_{\min}^{\mathrm{supp}}$. The smallest retained eigenvalue controls worst-case amplification of force-estimation error, while $s_{\max}$ gives the strongest retained tangent direction. If $\kappa_{\mathrm{eff}}\approx1$, retained modes have comparable strength; if $\kappa_{\mathrm{eff}}\gg1$, even the trusted tangent geometry is anisotropic and inversion remains delicate.}

\subsubsection{Appendix 17A-H.2 Driven Hamiltonian sampling and live
TDVP right-hand
side}\label{appendix-17a-h.2-driven-hamiltonian-sampling-and-live-tdvp-right-hand-side}

Here \(H_{\mathrm{stat}}\) is the static Pauli polynomial,
\(H_{\mathrm{drv}}(t)\) is the driven Pauli polynomial,
\(I_k:=[\tau_k,\tau_{k+1}]\) is the \(k\)th checkpoint interval,
\(\Delta\tau_k:=\tau_{k+1}-\tau_k\), \(M_k^{\mathrm{ref}}\ge 1\) is the
refined benchmark micro-step count, and \(\sigma_k\in\{0,\tfrac12,1\}\)
selects left, midpoint, or right checkpoint sampling. The absolute-time
convention is \[
t_k:=t_0+\tau_k,
\qquad
H_k^\sharp:=H(t_k+\sigma_k\Delta\tau_k).
\] For forecast depth \(h\), \[
t_{k+h}:=t_0+\tau_{k+h},
\qquad
H_{k+h}^\sharp:=H(t_{k+h}+\sigma_{k+h}\Delta\tau_{k+h}).
\] The driven Hamiltonian is

\(H(t):=H_{\mathrm{stat}}+H_{\mathrm{drv}}(t)=\sum_{\alpha\in\mathcal I_{\mathrm{stat}}}h_{\alpha}P_{\alpha}+\sum_{\beta\in\mathcal I_{\mathrm{drv}}}c_{\beta}(t)P_{\beta}\),
so
\(H_k^{\sharp}=\sum_{\alpha\in\mathcal I_{\mathrm{stat}}}h_{\alpha}P_{\alpha}+\sum_{\beta\in\mathcal I_{\mathrm{drv}}}c_{\beta}(t_k+\sigma_k\Delta\tau_k)P_{\beta}\).

With this sampled Hamiltonian, the force entries of §17A.1 become
\((f_k)_\mu=\sum_{\alpha\in\mathcal I_{\mathrm{stat}}}h_{\alpha}\Re\langle\bar\tau_{k,\mu},-iP_{\alpha}\psi_k\rangle+\sum_{\beta\in\mathcal I_{\mathrm{drv}}}c_{\beta}(t_k+\sigma_k\Delta\tau_k)\Re\langle\bar\tau_{k,\mu},-iP_{\beta}\psi_k\rangle\),
while \(G_k\) is unchanged because it depends only on tangent overlaps.

For any fixed ansatz support \(\mathcal S\), the physical projective
McLachlan velocity is the minimum-residual tangent velocity computed
from the horizontal tangent metric. The propagated controller velocity
may replace this by a regularized solve for numerical stability: \[
\begin{aligned}
|\psi(\theta;\mathcal S)\rangle
&:=U(\theta;\mathcal S)|\psi_{\mathrm{ref}}\rangle,\qquad
\bar T_{\mathcal S}(\theta):=Q_{\psi(\theta;\mathcal S)}T_{\mathcal S}(\theta),\\[1mm]
\bar b_{\mathcal S}(\theta,t)
&:=Q_{\psi(\theta;\mathcal S)}\bigl(-iH(t)|\psi(\theta;\mathcal S)\rangle\bigr),\\[1mm]
G_{\mathcal S}(\theta)
&:=\Re\!\bigl(\bar T_{\mathcal S}(\theta)^\dagger\bar T_{\mathcal S}(\theta)\bigr),\\[1mm]
f_{\mathcal S}(\theta,t)
&:=\Re\!\bigl(\bar T_{\mathcal S}(\theta)^\dagger\bar b_{\mathcal S}(\theta,t)\bigr),\\[1mm]
K_{\mathcal S}(\theta,t)
&:=G_{\mathcal S}(\theta)+\Lambda_{\mathcal S}(\theta,t),\\[1mm]
F_{\mathcal S}^{\mathrm{phys}}(\theta,t)
&:=G_{\mathcal S}(\theta)_{\tau}^{\dagger}f_{\mathcal S}(\theta,t),\\[1mm]
F_{\mathcal S}^{\mathrm{damp}}(\theta,t)
&:=K_{\mathcal S}(\theta,t)_{\tau,\varepsilon}^{\oplus}f_{\mathcal S}(\theta,t),\\[1mm]
F_{\mathcal S}(\theta,t)
&:=F_{\mathcal S}^{\mathrm{damp}}(\theta,t).
\end{aligned}
\] Thus the regularized velocity is the damped implementation of the
projective law, while the unregularized residual remains the structural
expressivity diagnostic used by append and prune. The regularizer
\(\Lambda_{\mathcal S}\) is symmetric positive semidefinite, has the
same units as \(G_{\mathcal S}\), and is chosen smoothly enough across
checkpoint intervals that it does not inject artificial high-frequency
coordinate motion. The sampled objects \(H_k^\sharp,G_k,f_k,K_k\) are
checkpoint diagnostics and cached score inputs. Open-interval
propagation uses \(F_{\mathcal S}(\theta,t)\), with stage-time
Hamiltonians when a multi-stage integrator is used.

The refined benchmark reference path is the high-resolution
piecewise-constant surrogate \[
\delta\tau_k:=\Delta\tau_k/M_k^{\mathrm{ref}},
\qquad
t_{k,m}^{\sharp}:=t_k+(m+\sigma_k)\delta\tau_k,
\qquad
H_{k,m}^{\sharp}:=H(t_{k,m}^{\sharp}),
\] \[
|\psi_{k,m+1}^{\mathrm{ref}}\rangle
:=
\exp\bigl(-i\delta\tau_kH_{k,m}^{\sharp}\bigr)
|\psi_{k,m}^{\mathrm{ref}}\rangle .
\] It is a benchmark/debug path for the sampled piecewise-constant
surrogate; the live controller uses local drift quantities rather than
future reference states.

\subsubsection{Appendix 17A-H.3 Residual geometry and local inadequacy
of the ansatz
support}\label{appendix-17a-h.3-residual-geometry-and-local-inadequacy-of-the-ansatz-support}

This subsection separates intrinsic ansatz-support inadequacy from
numerical damping. The intrinsic expressivity miss asks how much of the
sampled horizontal Schrödinger drift is outside the current tangent
plane before regularization. The realized miss asks how much mismatch
remains after the regularized live step. Structural append and prune
gates use the expressivity miss.

\[
\begin{aligned}
\epsilon_{\mathrm{expr},k}^{2}
&:=\operatorname{clip}\!\left(
\min_{v\in\mathbb R^{N_k}}\|\bar T_kv-\bar b_k\|_2^2,\,
0,\,
\operatorname{Var}_{\psi_k}(H_k^{\sharp})
\right)\\
&=\operatorname{clip}\!\left(
\|\bar b_k\|_2^2-f_k^{\top}G_{k,\tau}^{\dagger}f_k,\,
0,\,
\operatorname{Var}_{\psi_k}(H_k^{\sharp})
\right)\\
&=\operatorname{clip}\!\left(
\operatorname{Var}_{\psi_k}(H_k^{\sharp})-\sum_{\mu,\nu=1}^{N_k}(f_k)_\mu(G_{k,\tau}^{\dagger})_{\mu\nu}(f_k)_\nu,\,
0,\,
\operatorname{Var}_{\psi_k}(H_k^{\sharp})
\right),\\[1mm]
\rho_{\mathrm{expr},k}
&:=\operatorname{clip}\!\left(
\frac{\epsilon_{\mathrm{expr},k}^{2}}{\operatorname{Var}_{\psi_k}(H_k^{\sharp})+\varepsilon},\,0,\,1
\right),\\[1mm]
\bar r_{\mathrm{real},k}
&:=\bar T_kK_{k,\tau,\varepsilon}^{\oplus}f_k-\bar b_k\\
&=\sum_{\mu=1}^{N_k}(K_{k,\tau,\varepsilon}^{\oplus}f_k)_\mu\bigl[-i(\widetilde A_{k,\mu}-\langle\widetilde A_{k,\mu}\rangle_{\psi_k})|\psi_k\rangle\bigr]+i(H_k^{\sharp}-E_{\psi_k,k}^{\sharp})|\psi_k\rangle,\\[1mm]
\epsilon_{\mathrm{real},k}^{2}
&:=\|\bar r_{\mathrm{real},k}\|_2^2\\
&=\left\|\sum_{\mu=1}^{N_k}(K_{k,\tau,\varepsilon}^{\oplus}f_k)_\mu\bigl[-i(\widetilde A_{k,\mu}-\langle\widetilde A_{k,\mu}\rangle_{\psi_k})|\psi_k\rangle\bigr]+i(H_k^{\sharp}-E_{\psi_k,k}^{\sharp})|\psi_k\rangle\right\|_2^2,\\[1mm]
\rho_{\mathrm{real},k}
&:=\operatorname{clip}\!\left(
\frac{\epsilon_{\mathrm{real},k}^{2}}{\operatorname{Var}_{\psi_k}(H_k^{\sharp})+\varepsilon},\,0,\,\rho_{\max}
\right),\\[1mm]
\rho_{\mathrm{miss},k}
&:=\rho_{\mathrm{expr},k},
\qquad
\rho_{\mathrm{num},k}
:=\bigl[\rho_{\mathrm{real},k}-\rho_{\mathrm{expr},k}\bigr]_+,\\[1mm]
\operatorname{HighMiss}_k
&:=\mathbf 1\!\left[
\left(
\rho_{\mathrm{expr},k}>\tau_{\mathrm{miss}}
\;\wedge\;
\epsilon_{\mathrm{expr},k}^{2}>\epsilon_{\mathrm{abs}}^{\mathrm{miss}}
\right)
\vee
\sum_{j=k-L_{\mathrm{miss}}+1}^{k}
\mathbf 1[\rho_{\mathrm{expr},j}>\tau_{\mathrm{miss}}]
\ge N_{\mathrm{persist}}
\right].
\end{aligned}
\]

After this split, \(\rho_{\mathrm{miss},k}\) means the structural miss
\(\rho_{\mathrm{expr},k}\) unless explicitly labeled otherwise. The
diagnostic \(\rho_{\mathrm{real},k}\) tracks the actual regularized
step, and \(\rho_{\mathrm{num},k}\) isolates the extra miss attributable
to regularization, damping, conditioning, or other numerical
stabilization rather than missing ansatz directions. The current
checkpoint append lane uses \(\operatorname{HighMiss}_k\), not the
relative miss alone: near stationary drift, structural adaptation
requires absolute miss evidence or persistence across recent
checkpoints.\footnote{Safe implementation note: do not collapse $\rho_{\mathrm{expr}}$, $\rho_{\mathrm{real}}$, and $\rho_{\mathrm{num}}$ into one residual channel. Large $\rho_{\mathrm{expr}}$ means the current tangent plane is structurally missing Schrödinger drift and may justify append. Large $\rho_{\mathrm{num}}=[\rho_{\mathrm{real}}-\rho_{\mathrm{expr}}]_+$ means the damped or cutoff solve is failing to realize otherwise available tangent-plane expressivity and should trigger solve repair, not append by itself. Appending to fix only $\rho_{\mathrm{num}}$ can create a feedback loop in which ill-conditioning causes false append, the larger ansatz support worsens conditioning, and the realized residual grows again.}

\subsubsection{Appendix 17A-H.4 Checkpoint-local decision surfaces and
calmness
gate}\label{appendix-17a-h.4-checkpoint-local-decision-surfaces-and-calmness-gate}

The checkpoint law uses block-identifier transport before comparing
consecutive velocities. The transport map \[
\Pi_{k-1\to k}:\mathbb R^{N_{k-1}}\to\mathbb R^{N_k}
\] copies every surviving block's coordinates into its current slots and
fills newly created slots with zero. Let \(\mathsf M_k\succeq0\) be the
checkpoint comparison metric, typically \(G_k\) on its supported
subspace or the identity fallback. Use its induced norm and inner
product in the formulas below. The previous accepted velocity is
transported before calmness and rate-change tests are evaluated.

\[
\begin{aligned}
\dot\theta_{k-1\to k}^{\mathrm{acc}}
&:=\Pi_{k-1\to k}\dot\theta_{k-1}^{\mathrm{acc}},\\[1mm]
a_k^{\mathrm{vel}}
&:=\|\dot\theta_k\|_{\mathsf M_k},
\qquad
a_{k,-}^{\mathrm{vel}}
:=\|\dot\theta_{k-1\to k}^{\mathrm{acc}}\|_{\mathsf M_k},\\[1mm]
c_k^{\mathrm{dir}}
&:=\begin{cases}
1,&a_k^{\mathrm{vel}}\le\varepsilon_v\ \wedge\
a_{k,-}^{\mathrm{vel}}\le\varepsilon_v,\\[1mm]
\dfrac{\langle \dot\theta_k,\dot\theta_{k-1\to k}^{\mathrm{acc}}\rangle_{\mathsf M_k}}
{a_k^{\mathrm{vel}}a_{k,-}^{\mathrm{vel}}},
&a_k^{\mathrm{vel}}>\varepsilon_v\ \wedge\
a_{k,-}^{\mathrm{vel}}>\varepsilon_v,\\[2mm]
0,&\text{otherwise,}
\end{cases}\\[1mm]
\eta_k^{\mathrm{rate}}
&:=\frac{
\|\dot\theta_k-\dot\theta_{k-1\to k}^{\mathrm{acc}}\|_{\mathsf M_k}
}{
\max\{
a_k^{\mathrm{vel}},
a_{k,-}^{\mathrm{vel}},
\varepsilon
\}
},\\[1mm]
\mathfrak C_k
&:=\mathbf 1\bigl[c_k^{\mathrm{dir}}\ge \tau_{\mathrm{calm}}^{\cos}\bigr]\,\mathbf 1\bigl[\eta_k^{\mathrm{rate}}\le \tau_{\mathrm{calm}}^{\mathrm{rate}}\bigr],\\[1mm]
\kappa_{\mathrm{eff},k}
&:=\kappa_{\mathrm{eff}}(K_k),\qquad
F_k:=F_{\mathcal S_k}(\theta_k,t_k+\sigma_k\Delta\tau_k),
\qquad
\dot\theta_k=F_k,
\qquad
v_k^{\mathrm{state}}:=\bar T_k\dot\theta_k,\\[1mm]
\bar T_{k,b}
&:=\bar T_k[:,I_{b,k}],\qquad
\alpha_{b,k}^{\star}\in
\arg\min_{\alpha\in\mathbb R^{|I_{b,k}|}}
\|\bar T_{k,b}\alpha-v_k^{\mathrm{state}}\|_2^2,\\[1mm]
\pi_k^{\mathrm{col}}
&:=\frac{
\max_{b=1,\dots,B_k}
\|\bar T_{k,b}\alpha_{b,k}^{\star}\|_2^2
}{
\|v_k^{\mathrm{state}}\|_2^2+\varepsilon
}.
\end{aligned}
\]

\[
\begin{aligned}
\theta_k^{1/2}
&:=\theta_k+\frac12\Delta\tau_k\dot\theta_k,\\[1mm]
\dot\theta_k^{1/2}
&:=F_{\mathcal S_k}\!\left(\theta_k^{1/2},t_k+\frac12\Delta\tau_k\right),\\[1mm]
\chi_k^{\mathrm{curv}}
&:=\frac{
\left\|
\bar T_{\mathcal S_k}(\theta_k^{1/2})\dot\theta_k^{1/2}-\bar T_k\dot\theta_k
\right\|_2
}{
\|\bar T_k\dot\theta_k\|_2+\varepsilon
},\\[1mm]
\theta_{k,\mathrm{1E}}
&:=\operatorname{Step}_{\mathrm{Euler}}(\theta_k,t_k,\Delta\tau_k,\sigma_k;\mathcal S_k),\\[1mm]
\theta_{k,\mathrm{1/2E}}
&:=\operatorname{Step}_{\mathrm{Euler}}\!\left(\theta_k,t_k,\frac12\Delta\tau_k,\sigma_k;\mathcal S_k\right),\\[1mm]
\theta_{k,\mathrm{2E}}
&:=\operatorname{Step}_{\mathrm{Euler}}\!\left(\theta_{k,\mathrm{1/2E}},t_k+\frac12\Delta\tau_k,\frac12\Delta\tau_k,\sigma_k;\mathcal S_k\right),\\[1mm]
e_k^{\mathrm{Eul}}
&:=\arccos\!\left(\operatorname{clip}\!\left(\left|
\left\langle
\psi(\theta_{k,\mathrm{1E}};\mathcal S_k)
\middle|
\psi(\theta_{k,\mathrm{2E}};\mathcal S_k)
\right\rangle
\right|,0,1\right)\right),\\[1mm]
\operatorname{Col}_k
&:=\mathbf 1[\pi_k^{\mathrm{col}}\ge p_{\mathrm{col}}]\,
\mathbf 1[\chi_k^{\mathrm{curv}}\le\tau_{\mathrm{curv}}]\,
\mathbf 1[e_k^{\mathrm{Eul}}\le\tau_{\mathrm{Eul}}]\\
&\quad\times
\mathbf 1[\rho_{\mathrm{miss},k}\le\tau_{\mathrm{miss}}^{\mathrm{col}}]\,
\mathbf 1[\kappa_{\mathrm{eff},k}\le\kappa_{\mathrm{col}}].
\end{aligned}
\]

\[
\begin{aligned}
\zeta^{\mathrm{int}}(\mathcal S,\theta,t,\Delta\tau,\sigma)
&:=\bigl(\operatorname{Col},\pi^{\mathrm{col}},\chi^{\mathrm{curv}},e^{\mathrm{Eul}},\rho_{\mathrm{miss}},\kappa_{\mathrm{eff}}\bigr)[\mathcal S,\theta,t,\Delta\tau,\sigma],\\[1mm]
\zeta_k^{\mathrm{int}}
&:=\zeta^{\mathrm{int}}(\mathcal S_k,\theta_k,t_k,\Delta\tau_k,\sigma_k),\\[1mm]
\iota_k^{\mathrm{pol}}
&:=\Pi_{\mathrm{int}}(\zeta_k^{\mathrm{int}})
:=\begin{cases}
\mathrm{Euler},&\operatorname{Col}_k=1,\\[1mm]
\mathrm{RK4},&\operatorname{Col}_k=0,
\end{cases}\\[1mm]
\mathcal I_k^{\mathrm{int}}
&:=\{\iota_k^{\mathrm{pol}}\},\\[1mm]
\mathcal U_k^{\mathrm{stay}}
&:=\{(\mathrm{stay},\iota_k^{\mathrm{pol}})\},\\[1mm]
\mathcal U_k^{\mathrm{app}}:=\begin{cases}
\{(\mathrm{append},r):r\in\mathcal R_k^{\mathrm{conf}}\},&\mathcal R_k^{\mathrm{conf}}\neq\varnothing,\\[1mm]
\varnothing,&\mathcal R_k^{\mathrm{conf}}=\varnothing.
\end{cases}
\end{aligned}
\]

The callable surface
\(\zeta^{\mathrm{int}}(\mathcal S,\theta,t,\Delta\tau,\sigma)\) is
evaluated by replacing
\((\mathcal S_k,\theta_k,t_k,\Delta\tau_k,\sigma_k)\) in the definitions
above by the supplied ansatz-support, state, time, and sampling tuple.
In the paper-facing AP-McLachlan benchmark, time integration is a
declared numerical setting rather than a route-defining novelty. The
Euler/RK4 switch material in this archival subsection is not an active
Paper-II route; fixed Euler/RK4 comparisons may be reported only as
numerical sensitivity diagnostics.

\subsubsection{Appendix 17A-H.5 Position-aware append geometry and
confirm-grade
gain}\label{appendix-17a-h.5-position-aware-append-geometry-and-confirm-grade-gain}

An append move is a record \(r=(m,p)\) rather than a bare generator
because the new block is inserted at a definite ordered position. Let
\(m_r\) denote the number of new coordinates carried by the appended
block, let \(\eta=(\eta_1,\dots,\eta_{m_r})\) denote those new
coordinates, let \(\vartheta=(\theta,\eta)\) denote the augmented
coordinate vector, and let \(\eta=0\) at the append instant. The
physical state therefore does not jump at the checkpoint, but the
tangent plane acquires new columns.

The controller keeps two least-squares functionals separate. The
geometric functional answers the expressivity question used by
\(\operatorname{HighMiss}_k\); the regularized controller functional
uses the same supported inverse convention as the propagated McLachlan
velocity and is the authority for append gain and prune deletion loss.
Let \[
V_k:=\|\bar b_k\|_2^2=\operatorname{Var}_{\psi_k}(H_k^\sharp),
\qquad
J_k:=\{1,\ldots,N_k\},
\] and for any coordinate support \(S\) in the current tangent support
or in a zero-initialized augmented tangent support, \[
G_{k,SS}:=\Re(\bar T_{k,S}^{\dagger}\bar T_{k,S}),
\qquad
K_{k,SS}:=G_{k,SS}+\Lambda_{k,SS},
\qquad
f_{k,S}:=\Re(\bar T_{k,S}^{\dagger}\bar b_k).
\] The undamped geometric explained drift is \[
\Gamma_k^{G}(S):=
f_{k,S}^{\top}(G_{k,SS})_{\tau}^{\dagger}f_{k,S},
\qquad
\epsilon_{\mathrm{expr},k}^{2}(S):=
\operatorname{clip}\!\left(V_k-\Gamma_k^{G}(S),0,V_k\right),
\] while the regularized controller explained drift is \[
\Gamma_k^{K}(S):=
f_{k,S}^{\top}(K_{k,SS})_{\tau,\varepsilon}^{\oplus}f_{k,S}.
\] Hard support thresholding can break exact monotonicity across
different supports, so monotone Schur-ladder claims below assume a
common retained support or a common positive ridge. Static resource
costs inherited from the ground-state ansatz-construction stage are
denoted \(c_{\mathrm{hw}}\) or \(\kappa_{\mathrm{hw}}\), not \(K\), so
they cannot be confused with the McLachlan normal matrix \(K_k\).

At zero initialization, \[
|\psi(\theta_k,0;\mathcal S_k^{+}(r))\rangle=|\psi_k\rangle,
\qquad
Q_{\psi(\theta_k,0;\mathcal S_k^{+}(r))}=Q_{\psi_k}.
\] Hence the old tangent columns are unchanged on the augmented ansatz
support, the old-old blocks of the local McLachlan geometry are
inherited exactly, and only the old-new and new-new blocks require
candidate-local data. Thus zero-initialized append does not jump the
state and does not require remeasurement of old-old geometry.

For internal append primitive \(a\), let \(A_{r,a}\) be inserted at
position \(p\) and let \(U_{k,>p}^{+}\) denote the product of old
ansatz-support blocks lying after that insertion position on the
enlarged ordering. The appended transported generator and its horizontal
tangent column are \[
\widetilde A_{r,a}^{(p)}
:=
U_{k,>p}^{+}A_{r,a}(U_{k,>p}^{+})^\dagger,
\qquad
\bar u_{r,a}
:=
Q_{\psi_k}
\left.\partial_{\eta_a}|\psi(\theta_k,\eta;\mathcal S_k^+(r))\rangle\right|_{\eta=0}
=-iQ_{\psi_k}\widetilde A_{r,a}^{(p)}|\psi_k\rangle .
\] \[
\begin{aligned}
\bar U_{r,k}
&:=\bigl[(\bar U_{r,k})_{sa}\bigr]_{s=1,\dots,d}^{\,a=1,\dots,m_r}
 =\bigl[(Q_{\psi_k}\partial_{\eta_a}|\psi(\theta_k,\eta;\mathcal S_k^{+}(r))\rangle\vert_{\eta=0})_s\bigr]_{s=1,\dots,d}^{\,a=1,\dots,m_r}\\
&=\left[
\begin{array}{ccc}
(\bar u_{r,1})_1 & \cdots & (\bar u_{r,m_r})_1\\
\vdots & \ddots & \vdots\\
(\bar u_{r,1})_d & \cdots & (\bar u_{r,m_r})_d
\end{array}
\right],\\[1mm]
B_{r,k}
&:=\bigl[(B_{r,k})_{\mu a}\bigr]_{\mu=1,\dots,N_k}^{\,a=1,\dots,m_r}
 =\Re(\bar T_k^{\dagger}\bar U_{r,k})
 =\left[\Re\!\left(\sum_{s=1}^{d}(\bar T_k)_{s\mu}^{*}(\bar U_{r,k})_{sa}\right)\right]_{\mu=1,\dots,N_k}^{\,a=1,\dots,m_r},\\[1mm]
C_{r,k}
&:=\bigl[(C_{r,k})_{ab}\bigr]_{a,b=1}^{m_r}
 =\Re(\bar U_{r,k}^{\dagger}\bar U_{r,k})
 =\left[\Re\!\left(\sum_{s=1}^{d}(\bar U_{r,k})_{sa}^{*}(\bar U_{r,k})_{sb}\right)\right]_{a,b=1}^{m_r},\\[1mm]
q_{r,k}
&:=\bigl[(q_{r,k})_{a}\bigr]_{a=1}^{m_r}
 =\Re(\bar U_{r,k}^{\dagger}\bar b_k)
 =\left[\Re\!\left(\sum_{s=1}^{d}(\bar U_{r,k})_{sa}^{*}(\bar b_k)_s\right)\right]_{a=1}^{m_r},\\[1mm]
\mathfrak m_{r,k}(v,\xi)
&:=\|\bar T_kv+\bar U_{r,k}\xi-\bar b_k\|_2^2+v^{\top}\Lambda_kv+\xi^{\top}\Lambda_r\xi\\
&=\begin{bmatrix}v\\ \xi\end{bmatrix}^{\!\top}
\begin{bmatrix}K_k & B_{r,k}\\ B_{r,k}^{\top} & C_{r,k}+\Lambda_r\end{bmatrix}
\begin{bmatrix}v\\ \xi\end{bmatrix}
-2\begin{bmatrix}v\\ \xi\end{bmatrix}^{\!\top}\begin{bmatrix}f_k\\ q_{r,k}\end{bmatrix}
+\langle\psi_k|(H_k^{\sharp}-E_{\psi_k,k}^{\sharp})^2|\psi_k\rangle,\\[1mm]
\begin{bmatrix}\dot\theta_{r,k}^{\mathrm{old}}\\ \dot\eta_{r,k}\end{bmatrix}
&:=\begin{bmatrix}K_k & B_{r,k}\\ B_{r,k}^{\top} & C_{r,k}+\Lambda_r\end{bmatrix}_{\tau,\varepsilon}^{\oplus}\begin{bmatrix}f_k\\ q_{r,k}\end{bmatrix},
\qquad
\dot\vartheta_{r,k}^{\mathrm{aug}}
:=
\begin{bmatrix}\dot\theta_{r,k}^{\mathrm{old}}\\ \dot\eta_{r,k}\end{bmatrix},\\[1mm]
\end{aligned}
\]

\[
\begin{aligned}
\Pi_{k-1\to k}^{(r)}
&:\mathbb R^{N_{k-1}}\to\mathbb R^{N_k+m_r}
\\
&\quad\text{copies surviving old block coordinates},\\[-1mm]
&\quad\text{and fills appended coordinates with zero},\\[1mm]
G_{r,k}^{\mathrm{aug}}
&:=\begin{bmatrix}G_k & B_{r,k}\\ B_{r,k}^{\top} & C_{r,k}\end{bmatrix},\\[1mm]
\mathsf M_{r,k}
&:=\begin{cases}
G_{r,k}^{\mathrm{aug}},&\text{on the supported augmented tangent subspace},\\
I,&\text{otherwise},
\end{cases}
\qquad
v^{\star}_{K}(\xi):=K_{k,\tau,\varepsilon}^{\oplus}(f_k-B_{r,k}\xi),\\[1mm]
\mathfrak m_{r,k}(v^{\star}_{K}(\xi),\xi)
&=\left\|\bar T_kv^{\star}_{K}(\xi)+\bar U_{r,k}\xi-\bar b_k\right\|_2^2+(v_{K}^{\star}(\xi))^\top\Lambda_kv_{K}^{\star}(\xi)+\xi^\top\Lambda_r\xi,\\[1mm]
w_{r,k}^{G}
&:=q_{r,k}-B_{r,k}^{\top}G_{k,\tau}^{\dagger}f_k,\\[1mm]
S_{r,k}^{G}
&:=C_{r,k}-B_{r,k}^{\top}G_{k,\tau}^{\dagger}B_{r,k},\\[1mm]
\Delta_{r,k}^{G}
&:=(w_{r,k}^{G})^{\top}(S_{r,k}^{G})_{\tau}^{\dagger}w_{r,k}^{G},\\[1mm]
\Delta_{r,k}^{G,\mathrm{clip}}
&:=\operatorname{clip}\!\left(
\Delta_{r,k}^{G},\,0,\,\epsilon_{\mathrm{expr},k}^{2}
\right),\\[1mm]
\epsilon_{\mathrm{expr},r,k}^{2,+}
&:=\operatorname{clip}\!\left(
\min_{v,\xi}\|\bar T_kv+\bar U_{r,k}\xi-\bar b_k\|_2^2,\,
0,\,
\|\bar b_k\|_2^2
\right),\\[1mm]
\Delta_{r,k}^{\mathrm{auth}}
&:=\bigl[\epsilon_{\mathrm{expr},k}^{2}-\epsilon_{\mathrm{expr},r,k}^{2,+}\bigr]_+,\\[1mm]
I_{r,k}^{+}
&:=\{N_k+1,\ldots,N_k+m_r\},\\[1mm]
\Delta_{r,k}^{\mathrm{app}}
&:=\left[
\Gamma_k^{K}(J_k\cup I_{r,k}^{+})
-
\Gamma_k^{K}(J_k)
\right]_+,\\[1mm]
\rho_{r,k}^{\mathrm{gain,Schur}}
&:=\frac{\Delta_{r,k}^{G,\mathrm{clip}}}{\operatorname{Var}_{\psi_k}(H_k^\sharp)+\varepsilon},\\[1mm]
\rho_{r,k}^{\mathrm{gain,expr}}
&:=\frac{\Delta_{r,k}^{\mathrm{auth}}}{\operatorname{Var}_{\psi_k}(H_k^\sharp)+\varepsilon},\\[1mm]
\rho_{r,k}^{\mathrm{gain}}
&:=\frac{\Delta_{r,k}^{\mathrm{app}}}{V_k+\varepsilon},\\[1mm]
\end{aligned}
\]

\[
\begin{aligned}
w_{r,k}
&:=q_{r,k}-B_{r,k}^{\top}K_{k,\tau,\varepsilon}^{\oplus}f_k\\
&=\left[\Re\!\left(\sum_{s=1}^{d}(\bar U_{r,k})_{sa}^{*}(\bar b_k)_s\right)-\sum_{\mu=1}^{N_k}(B_{r,k})_{\mu a}(K_{k,\tau,\varepsilon}^{\oplus}f_k)_\mu\right]_{a=1}^{m_r},\\[1mm]
S_{r,k}
&:=C_{r,k}+\Lambda_r-B_{r,k}^{\top}K_{k,\tau,\varepsilon}^{\oplus}B_{r,k}\\
&=\left[(C_{r,k})_{ab}+(\Lambda_r)_{ab}-\sum_{\mu,\nu=1}^{N_k}(B_{r,k})_{\mu a}(K_{k,\tau,\varepsilon}^{\oplus})_{\mu\nu}(B_{r,k})_{\nu b}\right]_{a,b=1}^{m_r},\\[1mm]
\Delta_{r,k}^{\mathrm{Schur}}
&:=w_{r,k}^{\top}(S_{r,k})_{\tau}^{\dagger}w_{r,k}\\
&=(q_{r,k}-B_{r,k}^{\top}K_{k,\tau,\varepsilon}^{\oplus}f_k)^{\top}(C_{r,k}+\Lambda_r-B_{r,k}^{\top}K_{k,\tau,\varepsilon}^{\oplus}B_{r,k})_{\tau}^{\dagger}(q_{r,k}-B_{r,k}^{\top}K_{k,\tau,\varepsilon}^{\oplus}f_k),\\[1mm]
\rho_{r,k}^{\mathrm{gain,num}}
&:=\frac{\Delta_{r,k}^{\mathrm{Schur}}}{V_k+\varepsilon}.
\end{aligned}
\]

The grouped quantity \(\Delta_{r,k}^{\mathrm{app}}\) is the canonical
current-time append gain for AP-McLachlan. The Schur expression
\(\Delta_{r,k}^{\mathrm{Schur}}\) is the block evaluation of the same
grouped gain when the old support, candidate support, damping, and
retained spectral support are shared consistently; otherwise it is a
scout/confirm diagnostic inside the same tangent geometry. The geometric
quantity \(\Delta_{r,k}^{\mathrm{auth}}\) remains a useful structural
expressivity audit: it compares the current expressivity defect against
the directly recomputed augmented expressivity defect. The clipped
geometric Schur gain cannot exceed the currently available expressivity
defect.

The confirm-grade compressed path keeps the same append object and
resolves only the old-coordinate contribution on a retained spectral
support. If \(K_k=Z_k\Sigma_k^2Z_k^{\top}\) is the supported
factorization, \(z_k:=\Sigma_k^{-1}Z_k^{\top}f_k\), and
\(\Gamma_{r,k}:=\Sigma_k^{-1}Z_k^{\top}B_{r,k}\), then
\(w_{r,k}=q_{r,k}-\Gamma_{r,k}^{\top}z_k\), the compressed Schur block
on a retained mode set \(J_r\) is
\(\widetilde S_{r,k}(J_r):=C_{r,k}+\Lambda_r-\Gamma_{r,k;J_r}^{\top}\Gamma_{r,k;J_r}\),
the compressed gain is
\(\widetilde\Delta_{r,k}(J_r):=(q_{r,k}-\Gamma_{r,k}^{\top}z_k)^{\top}\widetilde S_{r,k}(J_r)_{\tau}^{\dagger}(q_{r,k}-\Gamma_{r,k}^{\top}z_k)\),
and the dimensionless transported direction penalty is \[
\widehat\delta_{r,k}^{\mathrm{dir}}
:=
\frac{
\|\dot\vartheta_{r,k}^{\mathrm{aug}}-\Pi_{k-1\to k}^{(r)}\dot\theta_{k-1}^{\mathrm{acc}}\|_{\mathsf M_{r,k}}
}{
\max\{\|\dot\vartheta_{r,k}^{\mathrm{aug}}\|_{\mathsf M_{r,k}},\|\Pi_{k-1\to k}^{(r)}\dot\theta_{k-1}^{\mathrm{acc}}\|_{\mathsf M_{r,k}},\varepsilon\}
}.
\] The confirm score is as follows. Here \(g_{r,k}^{\mathrm{new}}\) is
the estimated new measurement burden of record \(r\), such as new
expectation values, Pauli groups, shot cost, or acquisition time. \[
\begin{aligned}
S_{r,k}^{\mathrm{conf}}
:=\begin{cases}
\rho_{r,k}^{\mathrm{gain,num}}-\lambda_{\mathrm{dir}}\,\widehat\delta_{r,k}^{\mathrm{dir}}-\lambda_{\mathrm{meas}}g_{r,k}^{\mathrm{new}},&\text{full local confirm},\\[2mm]
\dfrac{(q_{r,k}-\Gamma_{r,k}^{\top}z_k)^{\top}\widetilde S_{r,k}(J_r)_{\tau}^{\dagger}(q_{r,k}-\Gamma_{r,k}^{\top}z_k)}{\operatorname{Var}_{\psi_k}(H_k^\sharp)+\varepsilon}-\lambda_{\mathrm{dir}}\,\widehat\delta_{r,k}^{\mathrm{dir}}-\lambda_{\mathrm{meas}}g_{r,k}^{\mathrm{new}},&\text{compressed confirm}.
\end{cases}
\end{aligned}
\] Thus the confirmed append shortlist is
\(\mathcal R_k^{\mathrm{conf}}:=\{r:\rho_{r,k}^{\mathrm{gain}}\ge \tau_{\mathrm{gain}},\ S_{r,k}^{\mathrm{conf}}\ge \tau_{\mathrm{conf}}\}\).

The compressed confirm surface is a certified monotone Schur ladder only
under a fixed numerator, nested retained mode sets, and a common
positive-definite retained Schur ridge. If \[
K_k=Z_k\Sigma_k^{2}Z_k^{\top},
\qquad
z_k:=\Sigma_k^{-1}Z_k^{\top}f_k,
\qquad
\Gamma_{r,k}:=\Sigma_k^{-1}Z_k^{\top}B_{r,k},
\] and if \[
J_r^{(0)}=\varnothing\subseteq J_r^{(1)}\subseteq \cdots \subseteq J_r^{(m_{\max})},
\] then with \[
\underline\Delta_{r,k}
:=
w_{r,k}^{\top}(C_{r,k}+\Lambda_r)_{\tau}^{\dagger}w_{r,k},
\] \[
\widetilde S_{r,k}(J_r^{(m)})
:=
C_{r,k}+\Lambda_r-\Gamma_{r,k;J_r^{(m)}}^{\top}\Gamma_{r,k;J_r^{(m)}},
\] \[
\widetilde\Delta_{r,k}^{(\lambda_S)}(J_r^{(m)})
:=
(q_{r,k}-\Gamma_{r,k}^{\top}z_k)^{\top}
\bigl(\widetilde S_{r,k}(J_r^{(m)})+\lambda_SI\bigr)^{-1}
(q_{r,k}-\Gamma_{r,k}^{\top}z_k),
\] \[
\Delta_{r,k}^{\mathrm{Schur},(\lambda_S)}
:=
w_{r,k}^{\top}(S_{r,k}+\lambda_SI)^{-1}w_{r,k},
\] with \(\lambda_S>0\) and
\(\widetilde S_{r,k}(J_r^{(m)})+\lambda_SI\succ0\) for every retained
support, then one has \[
\widetilde\Delta_{r,k}^{(\lambda_S)}(J_r^{(0)})
\le
\widetilde\Delta_{r,k}^{(\lambda_S)}(J_r^{(1)})
\le
\cdots
\le
\widetilde\Delta_{r,k}^{(\lambda_S)}(J_r^{(m_{\max})})
\le
\Delta_{r,k}^{\mathrm{Schur},(\lambda_S)}.
\] Thus scout, confirm, and commit remain inside the same local Schur
geometry only when the ridge-supported certificate assumptions hold. If
noisy estimates make a retained Schur block indefinite, the
implementation either uses the common ridge above or rejects the
compressed certificate. If thresholded pseudoinverse support changes,
the numerator is truncated, or the positive-definite ridge is not used,
the compressed sequence is a heuristic refinement rather than a monotone
certificate.\footnote{Safe implementation note: full local append confirmation and compressed confirmation do not have the same authority. The commit-grade structural object is the directly recomputed augmented expressivity improvement $\Delta_{r,k}^{\mathrm{auth}}=[\epsilon_{\mathrm{expr},k}^{2}-\epsilon_{\mathrm{expr},r,k}^{2,+}]_+$. Compressed Schur scores are safe as scout/confirm aids, and they are monotone certificates only under fixed numerator, nested retained supports, and a common positive-definite ridge or positive semidefinite retained Schur blocks. If those assumptions are not enforced, a compressed gain should not be the sole reason to commit append.}

\subsubsection{Appendix 17A-H.6 Diagnostic proposal family for
within-support
motion}\label{appendix-17a-h.6-diagnostic-proposal-family-for-within-support-motion}

A stay action keeps the ordered ansatz support fixed and uses the
declared checkpoint integrator. The optional diagnostic policy
\(\iota_k^{\mathrm{pol}}\) from §17A.4 may be used for ablation
bookkeeping. The same-support directions below are diagnostic and
forecast-proposal directions, not independent physical stay controls.
Here \(I\) denotes the identity operator or identity matrix from
context, \(I_k:=[\tau_k,\tau_{k+1}]\) is the current checkpoint
interval, \(I_k^{\mathrm{drv}}\subseteq\{1,\dots,N_k\}\) is the runtime
coordinate block occupied by the injected drive-aligned density
direction, and \(\mathcal B_k\) is the finite blend-weight set. A trial
vector \(v\in\mathbb R^{N_k}\) is a diagnostic candidate parameter
velocity and \(p\) is a normalized proposal direction.

\[
\begin{aligned}
\mathsf M_k
&:=\begin{cases}
G_k,&\text{when the tangent metric is used on its supported subspace},\\
I,&\text{otherwise},
\end{cases}
\qquad
\|v\|_{\mathsf M_k}:=\sqrt{v^{\top}\mathsf M_kv},
\qquad
\langle x,y\rangle_{\mathsf M_k}:=x^{\top}\mathsf M_ky,\\[1mm]
v_k^{\mathrm{McL}}
&:=\dot\theta_k,\\[1mm]
v_{k,\mathrm{drv}}^{\mathrm{cur}}[I_k^{\mathrm{drv}}]
&:=K_k[I_k^{\mathrm{drv}},I_k^{\mathrm{drv}}]_{\tau,\varepsilon}^{\oplus}f_k[I_k^{\mathrm{drv}}],
\qquad
v_{k,\mathrm{drv}}^{\mathrm{cur}}[(I_k^{\mathrm{drv}})^{c}]:=0,\\[1mm]
H_{k+1}^{\sharp}
&:=H(t_{k+1}+\sigma_{k+1}\Delta\tau_{k+1}),
\qquad
f_{k+1,\mu}^{\sharp}:=\Re\langle\bar\tau_{k,\mu},Q_{\psi_k}(-iH_{k+1}^{\sharp}|\psi_k\rangle)\rangle,\\[1mm]
K_{k+1}^{\sharp}
&:=G_k+\Lambda_{k+1}^{\sharp},\\[1mm]
v_{k,\mathrm{drv}}^{\mathrm{next}}[I_k^{\mathrm{drv}}]
&:=K_{k+1}^{\sharp}[I_k^{\mathrm{drv}},I_k^{\mathrm{drv}}]_{\tau,\varepsilon}^{\oplus}f_{k+1}^{\sharp}[I_k^{\mathrm{drv}}],
\qquad
v_{k,\mathrm{drv}}^{\mathrm{next}}[(I_k^{\mathrm{drv}})^{c}]:=0,\\[1mm]
d_k^{\perp}
&:=v_{k,\mathrm{drv}}^{\mathrm{next}}
-\frac{\langle v_k^{\mathrm{McL}},v_{k,\mathrm{drv}}^{\mathrm{next}}\rangle_{\mathsf M_k}}{\|v_k^{\mathrm{McL}}\|_{\mathsf M_k}^{2}+\varepsilon_v^2}\,v_k^{\mathrm{McL}},\\[1mm]
\widetilde d_k^{\perp}
&:=\sqrt{\frac{\|v_k^{\mathrm{McL}}\|_{\mathsf M_k}^{2}}{\|d_k^{\perp}\|_{\mathsf M_k}^{2}}}\,d_k^{\perp}
\qquad \text{when }\|d_k^{\perp}\|_{\mathsf M_k}>\varepsilon_v
\text{ and }\|v_k^{\mathrm{McL}}\|_{\mathsf M_k}>\varepsilon_v,\\[1mm]
v_{k,\beta}^{\mathrm{blend}}
&:=\sqrt{\frac{\|v_k^{\mathrm{McL}}\|_{\mathsf M_k}^{2}}{\|v_k^{\mathrm{McL}}+\beta\widetilde d_k^{\perp}\|_{\mathsf M_k}^{2}+\varepsilon_v^2}}\,
\bigl(v_k^{\mathrm{McL}}+\beta\widetilde d_k^{\perp}\bigr),
\qquad \beta\in\mathcal B_k,\\[1mm]
\widehat v
&:=\begin{cases}
v/\|v\|_{\mathsf M_k},&\|v\|_{\mathsf M_k}>\varepsilon_v,\\
\varnothing,&\text{otherwise},
\end{cases}\\[1mm]
\mathcal P_k^{\mathrm{stay}}
&\subseteq\{\widehat v_k^{\mathrm{McL}},\widehat v_{k,\mathrm{drv}}^{\mathrm{cur}},\widehat v_{k,\mathrm{drv}}^{\mathrm{next}},\widehat v_{k,\beta}^{\mathrm{blend}}\},\\[1mm]
\chi_k^{\mathrm{dir}}
&:=\frac{1+c_k^{\mathrm{dir}}}{2},
\qquad
\chi_k^{\mathrm{rate}}
:=\frac{\tau_{\mathrm{calm}}^{\mathrm{rate}}}{\tau_{\mathrm{calm}}^{\mathrm{rate}}+\eta_k^{\mathrm{rate}}},
\qquad
\chi_k^{\mathrm{miss}}
:=\frac{\tau_{\mathrm{miss}}}{\tau_{\mathrm{miss}}+\rho_{\mathrm{miss},k}},\\[1mm]
s_k
&:=\chi_k^{\mathrm{dir}}\chi_k^{\mathrm{rate}}\chi_k^{\mathrm{miss}},
\qquad
0\le s_k\le 1,\\[1mm]
\mathcal U_k^{\mathrm{stay}}
&:=\{(\mathrm{stay},\iota_k^{\mathrm{pol}})\}.
\end{aligned}
\]

The committed stay trajectory is generated by integrating
\(F_{\mathcal S_k}(\theta,t)=K_{\mathcal S_k}(\theta,t)_{\tau,\varepsilon}^{\oplus}f_{\mathcal S_k}(\theta,t)\)
with the declared checkpoint integrator. Empty normalized directions are
omitted from \(\mathcal P_k^{\mathrm{stay}}\). The checkpoint-local
scalar \(s_k\) is the joint direction, rate, and structural-miss
confidence of the current same-support regime: it increases when the
accepted motion remains directionally coherent and changes slowly, and
it decreases when the expressivity miss is large. The proposal family
\(\mathcal P_k^{\mathrm{stay}}\) diagnoses curvature and drive response.
Archived secant proposals are not part of the canonical stay family; see
Appendix 17A-S.

\subsubsection{Appendix 17A-H.7 Canonical append admission law without
secant
forecast}\label{appendix-17a-h.7-canonical-append-admission-law-without-secant-forecast}

The canonical QPU-facing checkpoint controller admits append using
current-time local McLachlan diagnostics only. The append candidate has
already passed the scout-confirm ladder of §17A.5; no self-predicted
secant, finite-horizon rollout score, exact reference state, or
benchmark path enters the route-defining admission predicate.

\[
\operatorname{Admit}_k(r)=1
\iff
\operatorname{HighMiss}_k=1
\wedge
\rho_{r,k}^{\mathrm{gain}}\ge\tau_{\mathrm{gain}}
\wedge
S_{r,k}^{\mathrm{conf}}\ge\tau_{\mathrm{conf}}.
\]

Here \(\operatorname{HighMiss}_k\) is the current structural
expressivity trigger, \(\rho_{r,k}^{\mathrm{gain}}\) is the normalized
grouped regularized append gain, and \(S_{r,k}^{\mathrm{conf}}\) is the
confirmation score produced by the staged Schur/authoritative
recomputation pipeline. Thus append is controlled by present local
residual geometry and confirmed structural gain, not by a forecasted
trajectory comparison. Historical candidate data may be reused only when
it has already been measured for another controller purpose; the
canonical forward append law does not acquire extra quantum measurements
merely to build a history for unadmitted candidates. If a static
ansatz-construction resource cost \(c_{r}^{\mathrm{hw}}\) is available
for a candidate block, it ranks or budgets candidate evaluation after
the absolute gain checks above; it does not turn a sub-threshold tangent
gain into an accepted append.

\subsubsection{Appendix 17A-H.8 Retired secant-forecast
slot}\label{appendix-17a-h.8-retired-secant-forecast-slot}

The former predicted-conditional secant and live forecast score are not
part of the canonical 17A route. They are retained only as the archived
diagnostic/oracle-facing construction in Appendix 17A-S. Euler and RK4
are numerical propagators for fixed-integrator runs or appendix-level
switching diagnostics. The archived secant construction instead builds a
chord-like finite-difference proposal and a forecast score, so it must
not set append admission, route identity, or QPU-facing defaults.

\subsubsection{Appendix 17A-H.9 Compression-prune law with projection,
shadow verification, and
rollback}\label{appendix-17a-h.9-compression-prune-law-with-projection-shadow-verification-and-rollback}

The prune law is a compression controller, not a hard test that one
block has zero physical content. A block may carry nonzero coordinates
and still be redundant after the surviving coordinates move. Therefore
prune pressure is increased by opening more opportunities, testing more
mature candidates, using broader compensating projections, and verifying
a short shadow horizon. The final acceptance threshold remains
recoverability and observable safety. Here \(b\) is a persistent logical
block identifier, \(I_{b,k}\) is the runtime coordinate set owned by
block \(b\) at checkpoint \(k\), and \(I_{b,k}^{c}\) is its complement.
Append and prune act on blocks, while the McLachlan arrays act on
flattened coordinates.

This time-dynamics prune law deliberately does \textbf{not} import the
static-ADAPT quasi-Newton Hessian surrogate as its controller object.
The live local object is the McLachlan normal geometry already measured
or simulated through the controller interface. For a block \(B\) with
runtime coordinates \(I_{B,k}\subseteq J_k\), set
\(J_k^{-B}:=J_k\setminus I_{B,k}\). The canonical tangent-space deletion
loss is the full active-support grouped extra-sum-of-squares loss \[
\boxed{
L_{B,k}^{\mathrm{full}}
:=
\frac{
\left[
\Gamma_k^{K}(J_k)-\Gamma_k^{K}(J_k^{-B})
\right]_+
}{
V_k+\varepsilon
}.
}
\] This is the backward companion of the append gain. It uses the same
damping, retained support, and measurement-compatible McLachlan data as
the propagated velocity. For an optional compensator window
\(W\subseteq J_k^{-B}\), define \[
L_{B,k}(W)
:=
\frac{
\left[
\Gamma_k^{K}(J_k)-\Gamma_k^{K}(W)
\right]_+
}{
V_k+\varepsilon
}.
\] Nested window losses may stage, explain, or label a prune attempt,
and under a common positive ridge/support convention they decrease as
the compensator window grows. The final canonical window is
\(W=J_k^{-B}\). Legacy raw-amplitude, pairwise-correlation, or proxy
losses are not interchangeable with \(L_{B,k}^{\mathrm{full}}\). The
linear grouped loss ranks and screens prune attempts, but hard deletion
still requires projected reduced-state initialization, differential miss
safety, persistence, and the short observable shadow certificate below.
Exact target/reference trajectories may be plotted or logged for
diagnostics, but they are not inputs to this controller loss or to the
shadow accept decision.

Exact static commutation metadata may be shared with the checkpoint
controller as prune anatomy. It may annotate dynamic prune events,
define bounded ordering under a candidate cap, construct optional typed
live-Schur diagnostic windows, or label a successful deletion. It is not
a deletion certificate and it does not replace the live McLachlan loss,
projected reduced-state safety, differential miss guard, shadow/no-harm
check, or persistence. Thus commutation classifies the prune anatomy,
while live McLachlan geometry decides the prune.

Because the current active-support McLachlan quantities are already
measured for the propagated velocity, prune may also use a weighted
history of full deletion losses without adding candidate measurements.
Let \[
\mathcal H_{b,k}^{\mathrm{pr}}
:=\{\ell\le k:\ b\in\mathcal B_\ell,\ L_{b,\ell}^{\mathrm{full}}\ \text{was recorded}\},
\qquad
w_{b,k,\ell}\ge0,\qquad
\sum_{\ell\in\mathcal H_{b,k}^{\mathrm{pr}}}w_{b,k,\ell}=1.
\] Define \[
\overline L_{b,k}^{\mathrm{full}}
:=
\sum_{\ell\in\mathcal H_{b,k}^{\mathrm{pr}}}
w_{b,k,\ell}L_{b,\ell}^{\mathrm{full}},
\qquad
\pi_{b,k}^{\mathrm{safe}}
:=
\sum_{\ell\in\mathcal H_{b,k}^{\mathrm{pr}}}
w_{b,k,\ell}
\mathbf 1[L_{b,\ell}^{\mathrm{full}}\le\tau_{\mathrm{loss}}^{\mathrm{hist}}].
\] The history predicate \[
\operatorname{HistSafe}_{b,k}^{\mathrm{pr}}
:=
\mathbf 1[
\overline L_{b,k}^{\mathrm{full}}\le\tau_{\mathrm{loss}}^{\mathrm{hist}}
]\,
\mathbf 1[
\pi_{b,k}^{\mathrm{safe}}\ge p_{\mathrm{safe}}^{\mathrm{hist}}
]
\] summarizes persistence of low full-support deletion loss. It can
rank, prioritize, or gate a prune only together with the current full
loss \(L_{b,k}^{\mathrm{full}}\); history does not replace the current
checkpoint test. No analogous forward append history is part of the
canonical law unless candidate tangent data were already acquired for
another reason.

Let \(\mathcal B_k\) be the active block set at checkpoint \(k\). Each
block \(b\in\mathcal B_k\) owns runtime coordinates \(I_{b,k}\) and,
when available, a multiset of inherited static generator labels \[
\mathcal G_b:=\{m_{b,1},\ldots,m_{b,r_b}\}.
\] For generator labels \(m,n\), reuse the exact static metadata \[
\Omega_{mn}:=\mathbf 1[\operatorname{supp}(m)\cap\operatorname{supp}(n)\neq\varnothing],
\qquad
\mathsf C_{mn}:=\mathbf 1[[A_m,A_n]=0],
\qquad
\mathsf q_{mn}\in\{\mathrm{exact},\mathrm{approx}\}.
\] For any generator pair, set \[
\begin{aligned}
e_{mn}^{\mathrm{flat}}
&:=
\mathbf 1[\Omega_{mn}=1]\,
\mathbf 1[\mathsf C_{mn}=1]\,
\mathbf 1[\mathsf q_{mn}=\mathrm{exact}],\\
e_{mn}^{\mathrm{curv}}
&:=
\mathbf 1[\Omega_{mn}=1]\,
\mathbf 1[\mathsf C_{mn}=0]\,
\mathbf 1[\mathsf q_{mn}=\mathrm{exact}],\\
e_{mn}^{\mathrm{disj}}
&:=
\mathbf 1[\Omega_{mn}=0]\,
\mathbf 1[\mathsf C_{mn}=1]\,
\mathbf 1[\mathsf q_{mn}=\mathrm{exact}].
\end{aligned}
\] For single-generator dynamic blocks \(i,b\), define the block
relation directly: \[
\begin{aligned}
E_{ib}^{\mathrm{flat,dyn}}
&:=e_{m_i m_b}^{\mathrm{flat}},\\
E_{ib}^{\mathrm{curv,dyn}}
&:=e_{m_i m_b}^{\mathrm{curv}},\\
E_{ib}^{\mathrm{disj,dyn}}
&:=e_{m_i m_b}^{\mathrm{disj}}.
\end{aligned}
\] The first relation is exact support-overlapping commutation; the
second is exact support-overlapping noncommutation; the third is exact
support-disjoint commutation. Support-disjoint commutation may be logged
or used as batching context, but it is not dynamic redundancy evidence:
\[
E_{ib}^{\mathrm{disj,dyn}}=1
\quad\not\Longrightarrow\quad
E_{ib}^{\mathrm{flat,dyn}}=1.
\] For composite blocks, the default first-pass controller does not lift
a single favorable generator pair into a block-level flat relation. If a
block has \(|\mathcal G_b|>1\) or \(|\mathcal G_i|>1\), and no explicit
typed diagnostic mode is enabled, set the algebraic block-pair class to
unclassified and use the live-only pruning path. In an optional typed
diagnostic mode, choose normalized internal weights \[
\alpha_{b,n}\ge0,\qquad
\sum_{n\in\mathcal G_b}\alpha_{b,n}=1,
\] and summarize pair types by weighted histograms \[
\phi_{ib}^{x}
:=
\sum_{m\in\mathcal G_i}
\sum_{n\in\mathcal G_b}
\alpha_{i,m}\alpha_{b,n}e_{mn}^{x},
\qquad
x\in\{\mathrm{flat},\mathrm{curv},\mathrm{disj}\}.
\] These histograms are diagnostic or ordering metadata unless the route
explicitly enables a typed-prune diagnostic; they do not alter the
default acceptance predicate.

When typed diagnostics are enabled, construct live compensator windows
from active blocks, not from bare generator metadata alone: \[
\begin{aligned}
W_{b,\mathrm{flat}}^{\mathrm{dyn}}(k)
&:=
\bigcup_{\substack{i\in\mathcal B_k\setminus\{b\}\\
E_{ib}^{\mathrm{flat,dyn}}=1}}
I_{i,k},\\
W_{b,\mathrm{curv}}^{\mathrm{dyn}}(k)
&:=
W_{b,\mathrm{flat}}^{\mathrm{dyn}}(k)
\cup
\bigcup_{\substack{i\in\mathcal B_k\setminus\{b\}\\
E_{ib}^{\mathrm{curv,dyn}}=1}}
I_{i,k},\\
W_{b,\mathrm{full}}^{\mathrm{dyn}}(k)
&:=J_k\setminus I_{b,k}.
\end{aligned}
\] For \(W\subseteq J_k\setminus I_{b,k}\), reuse the grouped
deletion-loss functional \(L_{b,k}(W)\) defined above. Then \[
L_{b,k}^{\mathrm{flat}}
:=L_{b,k}(W_{b,\mathrm{flat}}^{\mathrm{dyn}}),
\qquad
L_{b,k}^{\mathrm{curv}}
:=L_{b,k}(W_{b,\mathrm{curv}}^{\mathrm{dyn}}),
\qquad
L_{b,k}^{\mathrm{full}}
:=L_{b,k}(W_{b,\mathrm{full}}^{\mathrm{dyn}}).
\] Let \(\operatorname{TypedLoss}_{b,k}=1\) when
\(L_{b,k}^{\mathrm{flat}},L_{b,k}^{\mathrm{curv}},L_{b,k}^{\mathrm{full}}\)
are all recorded, and set it to zero otherwise. If typed losses are
recorded, an accepted prune may be labeled by the first live window that
explains the deletion: \[
\operatorname{dyn\_recovery\_class}_{b,k}
:=
\begin{cases}
\mathrm{flat\_dynamic\_redundant},
&\operatorname{accept}_{\mathrm{prune}}(b)=1,\ 
L_{b,k}^{\mathrm{flat}}\le\tau_{\mathrm{loss}}^{\mathrm{pr}},\\
\mathrm{curvature\_dynamic\_compensated},
&\operatorname{accept}_{\mathrm{prune}}(b)=1,\ 
L_{b,k}^{\mathrm{flat}}>\tau_{\mathrm{loss}}^{\mathrm{pr}},\
L_{b,k}^{\mathrm{curv}}\le\tau_{\mathrm{loss}}^{\mathrm{pr}},\\
\mathrm{broad\_dynamic\_compensated},
&\operatorname{accept}_{\mathrm{prune}}(b)=1,\ 
L_{b,k}^{\mathrm{curv}}>\tau_{\mathrm{loss}}^{\mathrm{pr}},\
L_{b,k}^{\mathrm{full}}\le\tau_{\mathrm{loss}}^{\mathrm{pr}},\\
\mathrm{unclassified},
&\operatorname{accept}_{\mathrm{prune}}(b)=1,\
\operatorname{TypedLoss}_{b,k}=0,\\
\mathrm{failed},
&\operatorname{accept}_{\mathrm{prune}}(b)=0.
\end{cases}
\] The label is explanatory. It cannot accept a deletion whose live
loss, ray displacement, differential miss, shadow/no-harm, or
persistence guard fails.

\[
\begin{aligned}
\mathcal P_{\mathrm{protect}}(k)
&:=\{\,b:\operatorname{age}_{b,k}<N_{\mathrm{prot}}\,\},
\qquad
\mathcal I_b^{\mathrm{stag}}(k)
:=
\{\,\ell:\max(k-\operatorname{age}_{b,k}+1,k-L_{\mathrm{stag}}+1)\le \ell\le k\,\},\\[1mm]
m_{b,k}^{\mathrm{mov}}
&:=
\frac{1}{|\mathcal I_b^{\mathrm{stag}}(k)|}
\sum_{\ell\in\mathcal I_b^{\mathrm{stag}}(k)}
\frac{
\sqrt{\Delta\theta_{I_{b,\ell},\ell}^{\top}
G_{\ell}[I_{b,\ell},I_{b,\ell}]
\Delta\theta_{I_{b,\ell},\ell}}
}{
\sqrt{\Delta\theta_{\ell}^{\top}G_{\ell}\Delta\theta_{\ell}}+\varepsilon_{\theta}
},\\[1mm]
m_{b,k}^{\mathrm{fit}}
&:=
\frac{1}{|\mathcal I_b^{\mathrm{stag}}(k)|}
\sum_{\ell\in\mathcal I_b^{\mathrm{stag}}(k)}
\frac{
\left|\Delta\theta_{I_{b,\ell},\ell}^{\top}f_{I_{b,\ell},\ell}\right|
}{
\sum_{c=1}^{B_\ell}
\left|\Delta\theta_{I_{c,\ell},\ell}^{\top}f_{I_{c,\ell},\ell}\right|
+\varepsilon_f
},\\[1mm]
\sigma_{b,k}^{\mathrm{stag}}
&:=\alpha_{\mathrm{pr}}
\bigl(1-\operatorname{clip}(m_{b,k}^{\mathrm{mov}},0,1)\bigr)
+
(1-\alpha_{\mathrm{pr}})
\bigl(1-\operatorname{clip}(m_{b,k}^{\mathrm{fit}},0,1)\bigr),\\[1mm]
\mathfrak C_k^{\mathrm{calm}}
&:=
\mathbf 1[\rho_{\mathrm{miss},k}\le\tau_{\mathrm{prune}}^{\mathrm{miss}}]\,
\mathbf 1[c_k^{\mathrm{dir}}\ge\tau_{\mathrm{calm}}^{\cos}]\,
\mathbf 1[\eta_k^{\mathrm{rate}}\le\tau_{\mathrm{calm}}^{\mathrm{rate}}].
\end{aligned}
\]

The full deletion loss \(L_{b,k}^{\mathrm{full}}\), the history
predicate, the block norm, and the stagnation score are nomination and
trust-region quantities. They do not by themselves certify deletion. The
committed reduced state must be produced by a reduced-manifold
projection, because the raw coordinate projection \(\Pi_{-b}\theta_k\)
is generally discontinuous in state, velocity, and observables when
\(\theta_{I_{b,k}}\neq0\). Let \(Y_k\) be the tracked observable bundle
used by the route, for example energy, occupations/density, doublon,
staggered density, and local miss summaries. In strict QPU-facing
routes, every component of \(Y_k\) used here is an observable estimate
of the prepared incumbent or candidate state; exact/reference
trajectories may appear only in diagnostic exact-assisted routes. Define
\[
\begin{aligned}
\mathcal P_{-b}(\theta_k)
&:=
\arg\min_{\xi\in\Theta_{\mathcal S_k^{(-b)}}}
\Bigl[
d_{\mathrm{FS}}^2\!\left(
\psi(\xi;\mathcal S_k^{(-b)}),
\psi(\theta_k;\mathcal S_k)
\right)\\
&\qquad\qquad
+\lambda_v
\left\|
\bar T_{\mathcal S_k^{(-b)}}(\xi)F_{\mathcal S_k^{(-b)}}(\xi,t_k)
-
\bar T_{\mathcal S_k}(\theta_k)F_{\mathcal S_k}(\theta_k,t_k)
\right\|_2^2\\
&\qquad\qquad
+\lambda_Y
\left\|
Y_{\mathcal S_k^{(-b)}}(\xi)-Y_k
\right\|_{\Sigma_Y^{-1}}^2
\Bigr],\\[1mm]
\theta_k^{(-b)}
&:=\mathcal P_{-b}(\theta_k),\qquad
|\psi_k^{(-b)}\rangle
&:=U(\theta_k^{(-b)};\mathcal S_k^{(-b)})|\psi_{\mathrm{ref}}\rangle,
\qquad
\Theta_{b,k}^{\mathrm{blk}}:=\|\theta_{k,I_{b,k}}\|_2,\qquad
\|x\|_{\Sigma^{-1}}:=\sqrt{x^\top\Sigma^{-1}x},\\[1mm]
\Delta_{\mathrm{ray},b,k}^{\mathrm{pr}}
&:=\arccos\!\left(\operatorname{clip}\!\left(\left|
\langle\psi_k^{(-b)}|\psi_k\rangle
\right|,0,1\right)\right),\\[1mm]
T_k^{(-b)}
&:=\left[\partial_{\vartheta_\mu}U(\theta_k^{(-b)};\mathcal S_k^{(-b)})|\psi_{\mathrm{ref}}\rangle\right]_\mu,
\qquad
\bar T_k^{(-b)}:=Q_{\psi_k^{(-b)}}T_k^{(-b)},\\[1mm]
\bar b_k^{(-b)}
&:=Q_{\psi_k^{(-b)}}\bigl(-iH_k^{\sharp}|\psi_k^{(-b)}\rangle\bigr),\\[1mm]
G_k^{(-b)}
&:=\Re\!\bigl((\bar T_k^{(-b)})^\dagger\bar T_k^{(-b)}\bigr),
\qquad
f_k^{(-b)}
:=\Re\!\bigl((\bar T_k^{(-b)})^\dagger\bar b_k^{(-b)}\bigr),\\[1mm]
\epsilon_{\mathrm{expr},k}^{2,(-b)}
&:=\operatorname{clip}\!\left(
\|\bar b_k^{(-b)}\|_2^2-(f_k^{(-b)})^{\top}(G_k^{(-b)})_{\tau}^{\dagger}f_k^{(-b)},\,0,\,\|\bar b_k^{(-b)}\|_2^2
\right),\\[1mm]
\rho_{\mathrm{expr},k}^{(-b)}
&:=\operatorname{clip}\!\left(\frac{\epsilon_{\mathrm{expr},k}^{2,(-b)}}{\|\bar b_k^{(-b)}\|_2^2+\varepsilon},\,0,\,1\right),
\qquad
\Delta\rho_{b,k}^{\mathrm{pr}}:=\rho_{\mathrm{expr},k}^{(-b)}-\rho_{\mathrm{expr},k}.
\end{aligned}
\]

The permissibility gate is candidate-specific. A low-miss calm regime is
the cheap ordinary path, but high-miss episodes may still prune a block
when the block is differentially irrelevant to the miss: \[
\begin{aligned}
\Gamma_k^{\mathrm{prune,perm,new}}(b)
&:=
\mathbf 1[B_k\ge2]\,
\mathbf 1[b\notin\mathcal P_{\mathrm{protect}}(k)]\,
\mathbf 1[\operatorname{cooldown}_{b,k}=0]\,
\mathbf 1[\operatorname{NeedsSolveRepair}_k=0]\,
\mathbf 1[\sigma_{b,k}^{\mathrm{stag}}\ge\tau_{\mathrm{stale}}]\\
&\quad\times
\mathbf 1\!\left[
\mathfrak C_k^{\mathrm{calm}}=1
\ \vee\
\left(
L_{b,k}^{\mathrm{full}}\le\tau_{\mathrm{loss}}^{\mathrm{pr}}
\ \wedge\
\Delta\rho_{b,k}^{\mathrm{pr}}\le\tau_{\rho}^{\mathrm{pr}}
\right)
\right].
\end{aligned}
\] This differential alternative is the formal pre-append compression
route: under high miss, test an old block only if projected deletion
shows the block is not the source of the current missing tangent
direction. After a committed compression, recompute the miss before
admitting another append.

Pressure expands the tested candidate set rather than lowering the
accept threshold. With maturity pressure \(g_k\), historical low-loss
persistence, and optional saved static hardware cost
\(\kappa_b^{\mathrm{hw}}\), \[
\begin{aligned}
Q_{\mathrm{pr},k}
&:=
\left\lceil
Q_{\mathrm{pr}}^{\min}
+(Q_{\mathrm{pr}}^{\max}-Q_{\mathrm{pr}}^{\min})g_k
\right\rceil,\\[1mm]
\lambda_{\mathrm{pr}}^{\mathrm{td}}(k)
&:=
\lambda_{\mathrm{pr}}^{\min}
+(\lambda_{\mathrm{pr}}^{\max}-\lambda_{\mathrm{pr}}^{\min})g_k,\\[1mm]
P_{b,k}^{\mathrm{pr,td}}
&:=
L_{b,k}^{\mathrm{full}}
+w_{\rho}[\Delta\rho_{b,k}^{\mathrm{pr}}]_+
+w_{\mathrm{ray}}\Delta_{\mathrm{ray},b,k}^{\mathrm{pr}}
-\lambda_{\mathrm{pr}}^{\mathrm{td}}(k)
\left(
w_{\mathrm{stag}}\sigma_{b,k}^{\mathrm{stag}}
+w_{\mathrm{age}}z_{\mathrm{age},b,k}
+w_{\mathrm{hist}}\operatorname{HistSafe}_{b,k}^{\mathrm{pr}}
+w_{\mathrm{hw}}\kappa_{b}^{\mathrm{hw}}
\right),\\[1mm]
\mathcal D_k^{\mathrm{pr}}
&:=
\operatorname{Top}^{\min}_{Q_{\mathrm{pr},k}}
\left(
P_{b,k}^{\mathrm{pr,td}};
\Gamma_k^{\mathrm{prune,perm,new}}(b)=1
\right).
\end{aligned}
\] If the algebraic prune anatomy above is available, it may modify this
capped ordering only through a bounded perturbation. A simple
live-algebraic context pair is \[
\begin{aligned}
C_{ib,k}^{\mathrm{live}}
&:=
\frac{
\|G_k[I_{i,k},I_{b,k}]\|_{\mathrm F}
}{
\sqrt{
\|G_k[I_{i,k},I_{i,k}]\|_{\mathrm F}
\|G_k[I_{b,k},I_{b,k}]\|_{\mathrm F}
}
+\varepsilon
},\\
F_{b,k}^{\mathrm{flat}}
&:=
\max_{i\in\mathcal B_k\setminus\{b\}}
E_{ib}^{\mathrm{flat,dyn}}\,C_{ib,k}^{\mathrm{live}},\\
B_{b,k}^{\mathrm{bridge}}
&:=
\operatorname{clip}_{[0,1]}
\left(
B_b^{\mathrm{static}}
+
\max_{i\in\mathcal B_k\setminus\{b\}}
E_{ib}^{\mathrm{curv,dyn}}\,C_{ib,k}^{\mathrm{live}}
\right),
\end{aligned}
\] with \(B_b^{\mathrm{static}}:=0\) when no inherited bridge score is
available. Then, for example, \[
\begin{aligned}
\delta P_{b,k}^{\mathrm{alg}}
&:=
-w_{\mathrm{flat}}F_{b,k}^{\mathrm{flat}}
+w_{\mathrm{bridge}}B_{b,k}^{\mathrm{bridge}},\\
\left|\delta P_{b,k}^{\mathrm{alg}}\right|
&\le
\mu_{\mathrm{alg}}
\left(
L_{b,k}^{\mathrm{full}}
+\tau_{\mathrm{loss}}^{\mathrm{pr}}
+\varepsilon
\right),
\qquad
0<\mu_{\mathrm{alg}}\ll1,\\
P_{b,k}^{\mathrm{pr,td,alg}}
&:=
P_{b,k}^{\mathrm{pr,td}}
+\delta P_{b,k}^{\mathrm{alg}}.
\end{aligned}
\] The controller may use \(P_{b,k}^{\mathrm{pr,td,alg}}\) in the
candidate cap when this option is enabled, but the acceptance predicate
below is unchanged. The block norm \(\Theta_{b,k}^{\mathrm{blk}}\) may
rank numerical trust or trigger a wider projection solve, but it is not
a hard mathematical proof of nonredundancy.

For each \(b\in\mathcal D_k^{\mathrm{pr}}\), run a shadow-prune branch
over \(H_{\mathrm{pr}}\) checkpoints from the incumbent state and from
the projected reduced state, using the same declared one-step method and
the same sampled Hamiltonian schedule. Let \(Y_{k+h}^{\mathrm{inc}}\)
and \(Y_{k+h}^{(-b)}\) denote the tracked observable bundles along these
two branches, and let \(\Delta Y_{k+h}:=Y_{k+h}-Y_{k+h-1}\). Define \[
\begin{aligned}
\mathcal J_{b,k}^{\mathrm{shadow}}
&:=
\max_{1\le h\le H_{\mathrm{pr}}}
\left\|
Y_{k+h}^{(-b)}-Y_{k+h}^{\mathrm{inc}}
\right\|_{\Sigma_Y^{-1}}\\
&\quad+
\lambda_{\Delta Y}
\max_{1\le h\le H_{\mathrm{pr}}}
\left\|
\Delta Y_{k+h}^{(-b)}-\Delta Y_{k+h}^{\mathrm{inc}}
\right\|_{\Sigma_{\Delta Y}^{-1}}\\
&\quad+
\lambda_\rho
\max_{1\le h\le H_{\mathrm{pr}}}
\left[
\rho_{\mathrm{expr},k+h}^{(-b)}
-
\rho_{\mathrm{expr},k+h}^{\mathrm{inc}}
\right]_+,\\[1mm]
\mathfrak C_{b,k}^{\mathrm{shadow}}
&:=
\mathbf 1[
\mathcal J_{b,k}^{\mathrm{shadow}}
\le
\tau_{\mathrm{shadow}}^{\mathrm{pr}}
].
\end{aligned}
\] In an exact-assisted diagnostic route only, the incumbent comparison
may be replaced by the reference-relative guard
\(\mathcal J_{b,k}^{\mathrm{shadow,ref}}\le\mathcal J_{\mathrm{stay},k}^{\mathrm{ref}}+\tau_{\mathrm{pr}}\).
This diagnostic reference guard is forbidden as a strict QPU-facing
controller input.

Do not commit a hard deletion at the first checkpoint where it barely
passes. Mark a block dormant-prunable and require persistence: \[
\begin{aligned}
\operatorname{Dormant}_{b,k}
&:=\mathbf 1[
L_{b,k}^{\mathrm{full}}\le\tau_{\mathrm{loss}}^{\mathrm{pr}}
]\,
\mathbf 1[
\operatorname{HistSafe}_{b,k}^{\mathrm{pr}}=1
]\,
\mathbf 1[
\Delta_{\mathrm{ray},b,k}^{\mathrm{pr}}\le\tau_{\mathrm{ray}}^{\mathrm{pr}}
]\,
\mathbf 1[
\Delta\rho_{b,k}^{\mathrm{pr}}\le\tau_{\rho}^{\mathrm{pr}}
]\,
\mathfrak C_{b,k}^{\mathrm{shadow}},\\[1mm]
\operatorname{Persist}_{b,k}^{\mathrm{pr}}
&:=
\mathbf 1\!\left[
\sum_{q=k-q_{\mathrm{pers}}+1}^{k}
\operatorname{Dormant}_{b,q}
\ge q_{\mathrm{req}}
\right],\\[1mm]
\operatorname{accept}_{\mathrm{prune}}(b)
&\Longleftrightarrow
\Gamma_k^{\mathrm{prune,perm,new}}(b)=1
\;\wedge\;
\operatorname{Dormant}_{b,k}=1
\;\wedge\;
\operatorname{Persist}_{b,k}^{\mathrm{pr}}=1,\\[1mm]
\operatorname{PruneReady}_k
&:=\mathbf 1[
\exists b\in\mathcal D_k^{\mathrm{pr}}:
\operatorname{accept}_{\mathrm{prune}}(b)=1
],\\[1mm]
b_k^{\star}
&:=\begin{cases}
\displaystyle\arg\min_{\substack{b\in\mathcal D_k^{\mathrm{pr}}\\ \operatorname{accept}_{\mathrm{prune}}(b)=1}}
P_{b,k}^{\mathrm{pr,td}},&\operatorname{PruneReady}_k=1,\\[2mm]
\varnothing,&\operatorname{PruneReady}_k=0.
\end{cases}
\end{aligned}
\]

Equivalently, an implementation may use a homotopy deletion variable
\(s_b\in[0,1]\), replace the block factor by
\(e^{-is_b(\tau)\theta_bG_b}\), project/refit the surviving coordinates
while \(s_b\downarrow0\), and commit only when \(s_b=0\), the
persistence condition holds, and the shadow observable score remains
safe. This is the canonical smooth-prune principle: prune a little after
possibility, not at the first barely feasible raw deletion. If a test
fails, the ansatz support rolls back exactly and block \(b\) enters
cooldown. If \(\mathcal D_k^{\mathrm{pr}}=\varnothing\) or
\(\operatorname{PruneReady}_k=0\), the prune lane falls back to the next
admissible lane.

\subsubsection{Appendix 17A-H.10 Canonical hybrid event
map}\label{appendix-17a-h.10-canonical-hybrid-event-map}

We now assemble the live local-projective controller as a hybrid event
map. At checkpoint \(k\), the controller diagnoses the current pair
\((\mathcal S_k,\theta_k)\), either commits one accepted structural/stay
event or opens a no-advance repair event, and then integrates only an
accepted ansatz support across the next interval by the declared
checkpoint integrator. The Euler/RK4 policy below is a diagnostic
option, not a main Paper-II mechanism.

\[
\begin{aligned}
\operatorname{Step}_{\mathrm{Euler}}(\theta,t,\Delta\tau,\sigma;\mathcal S)
&:=\theta+\Delta\tau\,F_{\mathcal S}(\theta,t+\sigma\Delta\tau),\\[1mm]
v_i
&:=F_{\mathcal S}\!\left(\theta+\Delta\tau\sum_{j<i}a_{ij}v_j,\ t+c_i\Delta\tau\right),
\qquad i=1,\ldots,4,\\[1mm]
\operatorname{Step}_{\mathrm{RK4}}(\theta,t,\Delta\tau,\sigma;\mathcal S)
&:=\theta+\Delta\tau\left(\frac16v_1+\frac13v_2+\frac13v_3+\frac16v_4\right),\\[1mm]
\mathsf S_{\mathrm{stay}}(\mathcal S_k,\theta_k)
&:=(\mathcal S_k,\theta_k),\\[1mm]
\mathsf S_{\mathrm{append}(r)}(\mathcal S_k,\theta_k)
&:=(\mathcal S_k^{+}(r),\theta_k\oplus_r0),\\[1mm]
\mathsf S_{\mathrm{prune}(b)}(\mathcal S_k,\theta_k)
&:=(\mathcal S_k^{(-b)},\mathcal P_{-b}(\theta_k)).
\end{aligned}
\]

\[
\begin{aligned}
\operatorname{NeedsSolveRepair}_k
&:=\mathbf 1\!\left[
\operatorname{unsupported}(K_k)
\;\vee\;
\kappa_{\mathrm{eff},k}>\kappa_{\max}
\;\vee\;
\rho_{\mathrm{num},k}>\tau_{\mathrm{num}}
\right],\\[1mm]
\operatorname{NeedsNoAdmitResolution}_k
&:=\operatorname{HighMiss}_k\,
\mathbf 1\!\left[
\mathcal U_k^{\mathrm{app}}=\varnothing
\;\vee\;
\nu_{k,\mathrm{app}}^\star=\varnothing
\;\vee\;
\operatorname{Admit}_k(u_{k,\mathrm{app}}^\star)=0
\right],\\[1mm]
\pi_k^{\mathrm{HNA}}
&\in\left\{
\mathrm{bounded\_stay\_advance},\,\mathrm{repair\_stop},\,\mathrm{repair\_retry}
\right\},\qquad
\pi_{\mathrm{default}}^{\mathrm{HNA}}=\mathrm{bounded\_stay\_advance},\\[1mm]
\operatorname{lane}_k
&:=\begin{cases}
\mathrm{repair}_{\mathrm{solve}},&\operatorname{NeedsSolveRepair}_k=1,\\[1mm]
\mathrm{prune},&\operatorname{PruneReady}_k=1,\\[1mm]
\mathrm{append},&\operatorname{HighMiss}_k=1\ \wedge\ \mathcal U_k^{\mathrm{app}}\neq\varnothing\ \wedge\ \operatorname{Admit}_k(u_{k,\mathrm{app}}^\star)=1,\\[1mm]
\mathrm{bounded}_{\mathrm{stay}},&\operatorname{NeedsNoAdmitResolution}_k=1\ \wedge\ \pi_k^{\mathrm{HNA}}=\mathrm{bounded\_stay\_advance},\\[1mm]
\mathrm{repair}_{\mathrm{miss}},&\operatorname{NeedsNoAdmitResolution}_k=1\ \wedge\ \pi_k^{\mathrm{HNA}}\in\{\mathrm{repair\_stop},\mathrm{repair\_retry}\},\\[1mm]
\mathrm{stay},&\text{otherwise,}
\end{cases}.
\end{aligned}
\]

\[
\begin{aligned}
\mathfrak r_k^{\mathrm{solve}}
&\in
\left\{
\Lambda_k\uparrow,\,
\tau_{\mathrm{pinv}}\uparrow,\,
\Delta\tau_k\downarrow,\,
\operatorname{flag}_{\mathrm{unsupported}}
\right\},\\[1mm]
\mathfrak r_k^{\mathrm{miss}}
&\in
\left\{
\Delta\tau_k\downarrow,\,
N_k^{\mathrm{scout}}\uparrow,\,
N_k^{\mathrm{conf}}\uparrow,\,
\Lambda_k\uparrow,\,
\tau_{\mathrm{pinv}}\uparrow,\,
\operatorname{flag}_{\mathrm{underresolved}}
\right\},\\[1mm]
\nu_k^{\star}
&:=\begin{cases}
\mathrm{append}(r_k^\star),&\operatorname{lane}_k=\mathrm{append},\ \operatorname{Admit}_k(u_{k,\mathrm{app}}^\star)=1,\\[1mm]
\mathrm{stay}_{\mathrm{soft}}(u_{k,\mathrm{stay}}^\star),&\operatorname{lane}_k=\mathrm{bounded}_{\mathrm{stay}},\\[1mm]
\mathrm{repair}_{\mathrm{solve}}(\mathfrak r_k^{\mathrm{solve}}),&\operatorname{lane}_k=\mathrm{repair}_{\mathrm{solve}},\\[1mm]
\mathrm{repair}_{\mathrm{miss}}(\mathfrak r_k^{\mathrm{miss}}),&\operatorname{lane}_k=\mathrm{repair}_{\mathrm{miss}},\\[1mm]
\mathrm{prune}(b_k^\star),&\operatorname{lane}_k=\mathrm{prune},\ b_k^\star\neq\varnothing,\ \operatorname{accept}_{\mathrm{prune}}(b_k^\star)=1,\\[1mm]
\mathrm{stay}(u_{k,\mathrm{stay}}^\star),&\text{otherwise,}
\end{cases}\\[1mm]
\mathcal V_k^{\mathrm{repair}}
&:=\left\{
\mathrm{repair}_{\mathrm{solve}}(\mathfrak r_k^{\mathrm{solve}}),\,
\mathrm{repair}_{\mathrm{miss}}(\mathfrak r_k^{\mathrm{miss}})
\right\}.
\end{aligned}
\]

\[
\begin{aligned}
(\mathcal S_k^\nu,\theta_k^{\nu,+})
&:=\begin{cases}
\mathsf S_{\nu_k^\star}(\mathcal S_k,\theta_k),&\nu_k^\star\notin\mathcal V_k^{\mathrm{repair}},\\[1mm]
(\mathcal S_k,\theta_k),&\nu_k^\star\in\mathcal V_k^{\mathrm{repair}},
\end{cases}\\[1mm]
\zeta_k^{\mathrm{int},\nu}
&:=\zeta^{\mathrm{int}}(\mathcal S_k^\nu,\theta_k^{\nu,+},t_k,\Delta\tau_k,\sigma_k),\\[1mm]
\iota_k^{\mathrm{pol},\nu}
&:=\Pi_{\mathrm{int}}(\zeta_k^{\mathrm{int},\nu}),\\[1mm]
\theta_{k+1}
&:=\begin{cases}
\operatorname{Step}_{\iota_k^{\mathrm{pol},\nu}}\!\left(\theta_k^{\nu,+},t_k,\Delta\tau_k,\sigma_k;\mathcal S_k^\nu\right),&\nu_k^\star\notin\mathcal V_k^{\mathrm{repair}},\\[1mm]
\theta_k,&\nu_k^\star\in\mathcal V_k^{\mathrm{repair}},
\end{cases}\\[1mm]
\mathcal S_{k+1}
&:=\begin{cases}
\mathcal S_k^\nu,&\nu_k^\star\notin\mathcal V_k^{\mathrm{repair}},\\[1mm]
\mathcal S_k,&\nu_k^\star\in\mathcal V_k^{\mathrm{repair}}.
\end{cases}
\end{aligned}
\]

The RK4 coefficients are \(c=(0,\frac12,\frac12,1)\),
\(a_{21}=a_{32}=\frac12\), \(a_{43}=1\), and all other \(a_{ij}=0\).
Append is state-continuous because \(\theta_k\oplus_r0\) inserts zero
coordinates; the enlarged tangent plane becomes available immediately,
and the new coordinates acquire amplitude through the subsequent TDVP
integration. Nonzero append seeds may be studied as ablations, but they
are not canonical live committed state jumps.

If \(\operatorname{NeedsSolveRepair}_k=1\), the controller repairs the
linear solve before any structural stay, append, or prune event; RK4 is
not treated as a cure for an unsupported or badly conditioned McLachlan
inverse. If \(\operatorname{HighMiss}_k=1\) but no append action is
admitted, the production implementation uses the explicit policy
parameter \(\pi_k^{\mathrm{HNA}}\). The default policy is
\texttt{bounded\_stay\_advance}: it does \textbf{not} relax append
thresholds and does \textbf{not} classify the event as repair; instead
it emits loud underresolved-fallback telemetry, preserves the strict
no-admit reason, and advances one physical \texttt{state\_sample} on the
current ansatz support. Explicit \texttt{repair\_stop} remains the
strict terminal diagnostic mode, and explicit \texttt{repair\_retry}
remains a no-advance experimental retry/rescue ladder. A solve-repair or
miss-repair event leaves \((\mathcal S_k,\theta_k,t_k)\) unchanged,
updates the selected controller knob, and recomputes the same physical
checkpoint before advancing. Append rejection under high miss under
\texttt{repair\_stop}/\texttt{repair\_retry} therefore means the
controller failed to repair a detected tangent-plane inadequacy, not
that same-support propagation has become
reliable.\footnote{Safe implementation note: the repair sets are finite escalation ladders, not free choices. For solve repair, first classify the trigger. If the solve is unsupported, accept a zero or near-zero step only when the sampled drift is below the stationary floor; otherwise emit $\operatorname{flag}_{\mathrm{unsupported}}$. If $\kappa_{\mathrm{eff},k}$ is too large, try $\Lambda_k\uparrow$, then $\tau_{\mathrm{pinv}}\uparrow$, then $\Delta\tau_k\downarrow$. If $\rho_{\mathrm{num},k}$ is too large while $\kappa_{\mathrm{eff},k}$ is acceptable, the failure is over-damping or cutoff-induced under-realization; do not repair it by adding ansatz-support blocks. Reduce unnecessary damping or cutoff if safe, otherwise shrink $\Delta\tau_k$ or flag unsupported. For explicit miss repair, first enlarge append scout/confirm resources and retry append admission; if no append is admitted after the finite budget, shrink $\Delta\tau_k$ only as damage control and then emit $\operatorname{flag}_{\mathrm{underresolved}}$. A failed explicit repair never falls through to ordinary stay, and repeated repair must stop rather than loop forever. The default $\texttt{bounded\_stay\_advance}$ branch is deliberately separate: it is a bounded physical advance with warning telemetry, not a repair event and not evidence that the ansatz support is safe.}

In implementation-loop form, \[
(\mathcal S_k,\theta_k)
\xrightarrow{\mathrm{diagnose/select}}
\nu_k^\star
\xrightarrow{\mathrm{structural\;map\;or\;repair}}
(\mathcal S_k^\nu,\theta_k^{\nu,+})
\xrightarrow{\operatorname{Step}_{\iota_k^{\mathrm{pol},\nu}}}
(\mathcal S_{k+1},\theta_{k+1}).
\] For \(\nu_k^\star\in\mathcal V_k^{\mathrm{repair}}\), the final
integration arrow is suppressed and the same checkpoint is retried after
the knob update. For \(\nu_k^\star=\mathrm{stay}_{\mathrm{soft}}\) under
\texttt{bounded\_stay\_advance}, the integration arrow is not
suppressed: the row is physical, carries
\texttt{high\_miss\_no\_admit\_soft\_fallback=true}, and records the
no-admit cause separately from the fallback resolution. Thus the
controller propagates by moving-state regularized projective McLachlan
dynamics, repairs unsupported or numerically damped solves before
propagation, measures local expressivity miss at checkpoints, opens
append only when the current ansatz support is locally inadequate,
admits append only when local Schur gain and confirmation score both
win, treats default high-miss/no-admit fallback as an underresolved
physical advance rather than repair, and allows prune through projected
compression only when differential expressivity, persistence, and
shadow-observable safety all pass.

\section{Appendix C. Cross-Checks and Exact-Benchmark
Contracts}\label{appendix-c.-cross-checks-and-exact-benchmark-contracts}

\subsection{C.1 Trial matrix by problem
family}\label{c.1-trial-matrix-by-problem-family}

For the Hubbard model, a natural benchmark trial set is \[
\mathcal T_{\mathrm{Hub}} = \{\text{layerwise},\;\text{excitation-based},\;\text{adaptive(exc)},\;\text{adaptive(full-H)}\}.
\] For the HH model, the canonical sector-preserving core set is \[
\mathcal T_{\mathrm{HH}} = \{\text{HH-layerwise},\;\text{HH-physical-termwise},\;\text{adaptive(full-H)}\}.
\] A seed-surface extension can enlarge \(\mathcal T_{\mathrm{HH}}\) by
adding polaronic refinements built on the same sector-preserving
fixed-family surface.

\subsection{C.2 Auto-scaled parameter
resolution}\label{c.2-auto-scaled-parameter-resolution}

Let \[
S_{\mathrm{trot}}(L,n_{\mathrm{ph,max}}),
\quad
R_{\mathrm{restart}}(L,n_{\mathrm{ph,max}}),
\quad
I_{\max}(L,n_{\mathrm{ph,max}})
\] be monotone nondecreasing control maps for Trotter steps, restart
count, and optimizer effort. The key contract is monotonicity with
problem size, not a single universal constant setting.

\subsection{C.3 Exact target and shared reference-state
construction}\label{c.3-exact-target-and-shared-reference-state-construction}

Within a benchmark matrix, all trial families should be compared against
the same exact target and the same conserved-sector reference
construction. Thus every trial is judged against the same
\(E_{\mathrm{exact,sector}}\) and the same physical Hamiltonian
instance.

\subsection{C.4 Per-trial energy and trajectory
semantics}\label{c.4-per-trial-energy-and-trajectory-semantics}

Each trial should at minimum record:

\begin{itemize}
\tightlist
\item
  the variational ground-state energy,
\item
  the exact comparison target,
\item
  trajectory observables on a shared time grid when dynamics are
  studied,
\item
  fidelity or overlap relative to a common reference branch.
\end{itemize}

\subsection{C.5 Artifact contracts}\label{c.5-artifact-contracts}

Every report artifact should begin with a concise parameter manifest
listing:

\begin{itemize}
\tightlist
\item
  model family,
\item
  ansatz family,
\item
  drive enabled or disabled,
\item
  \(t\), \(U\), \(\omega_0\), \(g\), and \(n_{\mathrm{ph,max}}\),
\item
  any propagation parameters needed to reproduce the result.
\end{itemize}

For dynamics benchmark tables, the manifest must also separate
static-scaffold provenance from propagation-controller tuning
provenance. The static scaffold may be the best benchmark-point seed
available for that Hamiltonian instance, provided every comparator uses
the same prepared scaffold for that point. McLachlan/controller settings
used to support class-level claims should instead carry an explicit
coarse class-tuning contract:

\begin{itemize}
\tightlist
\item
  \texttt{tuning\_granularity\ =\ coarse\_hamiltonian\_class},
\item
  \texttt{tuning\_class} one of
  \texttt{\{fermionic,\ bosonic,\ mixed\_fermion\_boson\}},
\item
  \texttt{source\_table\_class} carrying the finer executable
  benchmark-table class used for reporting, fixture selection, and audit
  trails,
\item
  \texttt{settings\_id} derived from controller/comparator knobs and
  coarse class only, not from \texttt{case\_id}, artifact paths, time
  grids, output paths, command lines, drive parameters, Hamiltonian
  family, or fine table class,
\item
  \texttt{static\_scaffold\_scope\ =\ benchmark\_point},
\item
  \texttt{class\_tuned\_result\_locked} indicating whether the row comes
  from a relocked class-tuned sweep or only a provisional fixture fill.
\end{itemize}

This contract prevents a time-dynamics table from silently presenting
Hamiltonian-specific controller retuning as a coarse Hamiltonian-class
result while still allowing the static ADAPT seed to remain
problem-specific and shared fairly across all dynamics methods. The
intended paper-II tuning classes are fermionic, bosonic, and mixed
fermion--boson; finer table classes remain useful metadata but must not
define controller settings identifiers.

\subsection{\texorpdfstring{C.6 Fixed-accuracy measurement ledger
\(S_{\rm norm}\)}{C.6 Fixed-accuracy measurement ledger S\_\{\textbackslash rm norm\}}}\label{c.6-fixed-accuracy-measurement-ledger-s_rm-norm}

Fix one displayed comparison row: one qubit Hamiltonian
\(H=\sum_{\alpha\in\Lambda_H}h_\alpha P_\alpha\), one target tolerance
\(\tau\), one reference energy \(E_{\rm ref}\), and one reference
computational state \[
\mathcal H=(\mathbb C^2)^{\otimes n_q},\qquad
x_0\in\{0,1\}^{n_q},\qquad
|\phi_0\rangle=|x_0\rangle=\prod_{q:x_{0q}=1}X_q|0\rangle^{\otimes n_q}.
\] For checkpoint \(q\), \[
|\psi_q\rangle=U_q(\theta_q)|\phi_0\rangle,\qquad
E_q=\langle\psi_q|H|\psi_q\rangle,\qquad
q_\star=\min\{q:\ |E_q-E_{\rm ref}|\le\tau\}.
\] The table reads \(N_{2q},D_{2q},D_{\rm circ}\), and \(S_{\rm norm}\)
at \(q_\star\).

The primitive ledger row consumed before \(q_\star\) is \[
j=\bigl(|\phi_0\rangle,\ U_j(\theta_j),\ O_j,\ f_j\bigr),
\qquad
|\psi_j\rangle=U_j(\theta_j)|\phi_0\rangle,\qquad
f_j\in\{0,1\}.
\] Here \(f_j=1\) marks new measurement work and \(f_j=0\) marks a
cached expectation. The circuit and observable entries are \[
U_j(\theta_j)=V_{j,d_j}(\theta_{j,d_j})\cdots V_{j,1}(\theta_{j,1}),\qquad
V_{j\ell}(\theta_{j\ell})=\exp[-i\theta_{j\ell}G_{j\ell}],
\] \[
G_{j\ell}=\sum_{\beta\in\Lambda_{j\ell}}g_{j\ell\beta}P_{j\ell\beta},\qquad
O_j=\sum_{\alpha\in\Lambda_j}h_{j\alpha}P_{j\alpha}.
\] Thus the expectation primitive under every ledger row is \[
\mu_j
=\langle\psi_j|O_j|\psi_j\rangle
=\langle\phi_0|U_j^\dagger O_jU_j|\phi_0\rangle
=\sum_{\alpha\in\Lambda_j}h_{j\alpha}
\langle\phi_0|U_j^\dagger P_{j\alpha}U_j|\phi_0\rangle .
\] For a single Pauli word \(P\), a hardware repetition measures bits
\(z\in\{0,1\}^{n_q}\) after the basis rotation for \(P\) and returns \[
Y_P(z)=\prod_{q:P_q\ne I}(-1)^{z_q},\qquad
\mathbb E(Y_P)=\langle\psi_j|P|\psi_j\rangle,\qquad
\operatorname{Var}(Y_P)=1-\langle\psi_j|P|\psi_j\rangle^2.
\] For a commuting Pauli group \(b\subseteq\Lambda_j\), \[
O_{j,b}=\sum_{\alpha\in b}h_{j\alpha}P_{j\alpha},\qquad
v_{j,b}=\langle O_{j,b}^2\rangle_j-\langle O_{j,b}\rangle_j^2
=\sum_{\alpha,\beta\in b}h_{j\alpha}h_{j\beta}
\left(\langle P_{j\alpha}P_{j\beta}\rangle_j-\langle P_{j\alpha}\rangle_j\langle P_{j\beta}\rangle_j\right),
\] where \(\langle X\rangle_j:=\langle\psi_j|X|\psi_j\rangle\). If the
event were expanded into calibrated repetitions with event precision
\(\epsilon_j\), the grouped Born-variance scaling would be \[
n_j^{\rm phys}(\epsilon_j)
\approx
\frac{1}{\epsilon_j^2}
\left(\sum_{b\in\mathcal B_j}\sqrt{\max(v_{j,b},0)}\right)^2,\qquad
S_{\rm phys}(q_\star)=\sum_{j=1}^{J(q_\star)} f_j\,n_j^{\rm phys}(\epsilon_j).
\]

The reported fixed-accuracy ledger is the unit count over four primitive
row predicates. With \(J_\star:=J(q_\star)\) and \[
\mathbf 1[\mathcal C]=
\begin{cases}
1,&\mathcal C\ {\rm is\ true},\\
0,&\mathcal C\ {\rm is\ false},
\end{cases}
\] the four count components are:

\begin{itemize}
\item
  \textbf{Outer Hamiltonian rows.} \[
  \begin{aligned}
  \mathcal C_{\rm out}(j):\quad
  &f_j=1,\qquad |\psi_j\rangle=U_j(\theta_j)|\phi_0\rangle,\qquad
  O_j=H=\sum_{\alpha\in\Lambda_H}h_\alpha P_\alpha,\\
  \mu_j^{\rm out}
  &:=\langle\psi_j|H|\psi_j\rangle
  =\sum_{\alpha\in\Lambda_H}h_\alpha
  \langle\phi_0|U_j^\dagger P_\alpha U_j|\phi_0\rangle,\\
  N_{\rm out}
  &:=\sum_{j=1}^{J_\star}\mathbf 1[\mathcal C_{\rm out}(j)].
  \end{aligned}
  \] The short label for this predicate is \(H_{\rm out}\).
\item
  \textbf{Candidate-gradient rows.} \[
  \begin{aligned}
  A_j&=\sum_{\gamma\in\Lambda_{A,j}}a_{j\gamma}P_{j\gamma},\\
  O_j&=i[H,A_j]
  =i\sum_{\alpha\in\Lambda_H}\sum_{\gamma\in\Lambda_{A,j}}
  h_\alpha a_{j\gamma}(P_\alpha P_{j\gamma}-P_{j\gamma}P_\alpha)
  =\sum_{\eta\in\Lambda_j^{\rm grad}}h_{j\eta}^{\rm grad}P_{j\eta},\\
  \mathcal C_{\rm grad}(j):\quad
  &f_j=1,\qquad |\psi_j\rangle=U_j(\theta_j)|\phi_0\rangle,\qquad O_j=i[H,A_j],\\
  \mu_j^{\rm grad}
  &:=\langle\psi_j|i[H,A_j]|\psi_j\rangle
  =\sum_{\eta\in\Lambda_j^{\rm grad}}h_{j\eta}^{\rm grad}
  \langle\phi_0|U_j^\dagger P_{j\eta}U_j|\phi_0\rangle,\\
  N_{\rm grad}
  &:=\sum_{j=1}^{J_\star}\mathbf 1[\mathcal C_{\rm grad}(j)].
  \end{aligned}
  \] The short label for this predicate is \(G\).
\item
  \textbf{Geometry rows.} \[
  \begin{aligned}
  \operatorname{Herm}(X)&:=\frac12(X+X^\dagger),\\
  B_{j,u}&=\sum_{\gamma\in\Lambda_{j,u}}b_{j,u\gamma}P_{j,u\gamma},
  \qquad
  B_{j,v}=\sum_{\delta\in\Lambda_{j,v}}b_{j,v\delta}P_{j,v\delta},\\
  O_{j}^{\rm FS}(u,v)
  &:=\operatorname{Herm}(B_{j,u}B_{j,v})
  =\operatorname{Herm}\!\left(
  \sum_{\gamma,\delta}b_{j,u\gamma}b_{j,v\delta}P_{j,u\gamma}P_{j,v\delta}
  \right),\\
  O_{j}^{H}(u,v)
  &:=\operatorname{Herm}(B_{j,u}HB_{j,v})
  =\operatorname{Herm}\!\left(
  \sum_{\gamma,\alpha,\delta}b_{j,u\gamma}h_\alpha b_{j,v\delta}
  P_{j,u\gamma}P_{\alpha}P_{j,v\delta}
  \right),\\
  O_{j}^{\rm mix}(u)
  &:=\operatorname{Herm}(B_{j,u}H)
  =\operatorname{Herm}\!\left(
  \sum_{\gamma,\alpha}b_{j,u\gamma}h_\alpha P_{j,u\gamma}P_\alpha
  \right).
  \end{aligned}
  \] \[
  \begin{aligned}
  \mathcal C_{\rm geom}(j):\quad
  &f_j=1,\qquad |\psi_j\rangle=U_j(\theta_j)|\phi_0\rangle,\\
  &O_j\in\{O_j^{\rm FS}(u,v),O_j^{H}(u,v),O_j^{\rm mix}(u)\},\qquad
  O_j=\sum_{\eta\in\Lambda_j^{\rm geom}}h_{j\eta}^{\rm geom}P_{j\eta},\\
  \mu_j^{\rm geom}
  &:=\langle\psi_j|O_j|\psi_j\rangle
  =\sum_{\eta\in\Lambda_j^{\rm geom}}h_{j\eta}^{\rm geom}
  \langle\phi_0|U_j^\dagger P_{j\eta}U_j|\phi_0\rangle,\\
  N_{\rm geom}
  &:=\sum_{j=1}^{J_\star}\mathbf 1[\mathcal C_{\rm geom}(j)].
  \end{aligned}
  \] The short label for this predicate is \(M\). The symbols \(u,v\)
  index tangent generators used for Fubini--Study metric, Gram, Hessian,
  Schur, rank, or additivity entries.
\item
  \textbf{Refit Hamiltonian rows.} \[
  \begin{aligned}
  U_j^{\rm fit}(\theta_j^{\rm fit})
  &=V_{j,d_j^{\rm fit}}^{\rm fit}(\theta_{j,d_j^{\rm fit}}^{\rm fit})
  \cdots V_{j,1}^{\rm fit}(\theta_{j,1}^{\rm fit}),\qquad
  V_{j\ell}^{\rm fit}(\theta)=\exp[-i\theta G_{j\ell}^{\rm fit}],\\
  G_{j\ell}^{\rm fit}
  &=\sum_{\beta\in\Lambda_{j\ell}^{\rm fit}}g_{j\ell\beta}^{\rm fit}P_{j\ell\beta}^{\rm fit},\\
  |\psi_j^{\rm fit}\rangle
  &=U_j^{\rm fit}(\theta_j^{\rm fit})|\phi_0\rangle,\qquad
  O_j=H=\sum_{\alpha\in\Lambda_H}h_\alpha P_\alpha,\\
  \mathcal C_{\rm fit}(j):\quad
  &f_j=1,\qquad |\psi_j\rangle=|\psi_j^{\rm fit}\rangle,\qquad O_j=H,\\
  \mu_j^{\rm fit}
  &:=\langle\psi_j^{\rm fit}|H|\psi_j^{\rm fit}\rangle
  =\sum_{\alpha\in\Lambda_H}h_\alpha
  \langle\phi_0|(U_j^{\rm fit})^\dagger P_\alpha U_j^{\rm fit}|\phi_0\rangle,\\
  N_{\rm fit}
  &:=\sum_{j=1}^{J_\star}\mathbf 1[\mathcal C_{\rm fit}(j)].
  \end{aligned}
  \] The short label for this predicate is \(H_{\rm fit}\).
\end{itemize}

The four predicates are disjoint and exhaustive for the displayed ledger
rows, so \[
f_j=\mathbf 1[\mathcal C_{\rm out}(j)]+\mathbf 1[\mathcal C_{\rm grad}(j)]
+\mathbf 1[\mathcal C_{\rm geom}(j)]+\mathbf 1[\mathcal C_{\rm fit}(j)].
\] Thus \[
\boxed{
S_{\rm norm}
=N_{\rm out}+N_{\rm grad}+N_{\rm geom}+N_{\rm fit}
=\sum_{j=1}^{J_\star}f_j
}
\]

Algorithmic specializations are zero-patterns of the same ledger:

\begin{itemize}
\tightlist
\item
  \textbf{HEA VQE and family-informed VQE.} \[
  N_{\rm grad}=N_{\rm geom}=N_{\rm fit}=0,\qquad
  S_{\rm norm}=N_{\rm out}.
  \]
\item
  \textbf{Append-only ADAPT and TETRIS-ADAPT.} \[
  N_{\rm out}=N_{\rm geom}=0,\qquad
  S_{\rm norm}=N_{\rm grad}+N_{\rm fit}.
  \]
\item
  \textbf{Qubit/QEB-ADAPT-VQE.} \[
  A_j\in\mathcal A_{\rm qeb},\qquad
  N_{\rm out}=N_{\rm geom}=0,\qquad
  S_{\rm norm}=N_{\rm grad}+N_{\rm fit}.
  \]
\item
  \textbf{Pos-Geo-ADAPT.} \[
  N_{\rm out}=0,\qquad
  S_{\rm norm}=N_{\rm grad}+N_{\rm geom}+N_{\rm fit}.
  \]
\item
  \textbf{SNAKE.} \[
  S_{\rm norm}=N_{\rm out}+N_{\rm grad}+N_{\rm geom}+N_{\rm fit}.
  \]
\end{itemize}

\subsection{C.7 HH seed-surface preset and resource
proxies}\label{c.7-hh-seed-surface-preset-and-resource-proxies}

A seed-surface comparison is naturally evaluated not only by energy but
also by simple resource proxies such as \[
C_{\mathrm{2q}},
\qquad
D_{\mathrm{proxy}},
\qquad
C_{\mathrm{1q}}.
\] One can then compare improvement per marginal resource, \[
\frac{\Delta E}{C_{\mathrm{2q}}},
\qquad
\frac{\Delta E}{D_{\mathrm{proxy}}}.
\]

\section{Appendix D. Noise Validation, Symmetry-Mitigation, and Parity
Contracts}\label{appendix-d.-noise-validation-symmetry-mitigation-and-parity-contracts}

\subsection{D.1 Filtered versus full exact
targets}\label{d.1-filtered-versus-full-exact-targets}

The filtered target is \[
E_{\mathrm{exact,filtered}} = \min_{\psi\in\mathcal H_{\mathrm{sector}}} \langle \psi|\hat H|\psi\rangle,
\] while the unrestricted target is \[
E_{\mathrm{exact,full}} = \min \operatorname{spec}(\hat H).
\] The two coincide only if the physical and computational sectors match
exactly.

\subsection{D.2 Noisy-minus-ideal observable
deltas}\label{d.2-noisy-minus-ideal-observable-deltas}

For any observable \(\mathcal O(t)\), define \[
\Delta \mathcal O(t) = \mathcal O_{\mathrm{noisy}}(t)-\mathcal O_{\mathrm{ideal}}(t).
\] In particular, \[
\Delta E(t)=E_{\mathrm{noisy}}(t)-E_{\mathrm{ideal}}(t),
\qquad
\Delta D(t)=D_{\mathrm{noisy}}(t)-D_{\mathrm{ideal}}(t).
\]

\subsection{D.3 Anchor parity gate}\label{d.3-anchor-parity-gate}

If a new idealized path is compared against a fixed anchor trajectory, a
strict parity gate takes the form \[
\max_{t\in\mathcal T}\,|\mathcal O_{\mathrm{new}}(t)-\mathcal O_{\mathrm{anchor}}(t)| \le \varepsilon_{\mathrm{parity}}.
\] This is an anchor-based equality test, not merely a qualitative
agreement statement.

\subsection{D.4 Same-anchor paired comparison
semantics}\label{d.4-same-anchor-paired-comparison-semantics}

Two noisy methods should be compared against the same ideal anchor, on
the same time grid, with the same initial state, so that differences can
be attributed to the noise or mitigation model rather than to different
starting conditions.

\subsection{D.5 Symmetry-mitigation mode
surface}\label{d.5-symmetry-mitigation-mode-surface}

Three increasingly strong symmetry surfaces are:

\begin{enumerate}
\def\labelenumi{\arabic{enumi}.}
\tightlist
\item
  verify-only sector diagnostics,
\item
  bitstring postselection,
\item
  projector-renormalized estimators.
\end{enumerate}

If \(\Pi\) is the projector onto the target symmetry sector, the
renormalized estimator is \[
\langle O\rangle_{\Pi} = \frac{\langle \psi|\Pi O \Pi|\psi\rangle}{\langle \psi|\Pi|\psi\rangle}.
\]

\subsection{D.6 Parameter manifest and provenance
digest}\label{d.6-parameter-manifest-and-provenance-digest}

Even in a symbolic setting, it is useful to view each artifact as
carrying a manifest \(\mathfrak M\) and a provenance digest
\(h(\mathfrak M)\). The purpose is not algebraic; it is to guarantee
that the reported mathematics can be traced back to a unique physical
parameter set.

\section{Appendix E. Spec-only or not-yet-formalized surfaces kept out
of the main
body}\label{appendix-e.-spec-only-or-not-yet-formalized-surfaces-kept-out-of-the-main-body}

\subsection{E.1 Broader continuation-theory
items}\label{e.1-broader-continuation-theory-items}

Several topics belong naturally in a later symbolic pass rather than in
the present core manuscript:

\begin{itemize}
\tightlist
\item
  multi-source motif transfer and motif tiling,
\item
  richer split-event formalism,
\item
  more detailed rescue dynamics,
\item
  exact handoff-policy closure beyond the present abstract maps,
\item
  measured-geometry and backend-specific closure beyond the exact
  symbolic statevector surface.
\end{itemize}

\subsection{\texorpdfstring{E.2 Obsolete historical warm-start \(\to\)
refine \(\to\) ADAPT \(\to\) replay
path}{E.2 Obsolete historical warm-start \textbackslash to refine \textbackslash to ADAPT \textbackslash to replay path}}\label{e.2-obsolete-historical-warm-start-to-refine-to-adapt-to-replay-path}

The historical staged path can be written symbolically as \[
\begin{aligned}
\mathcal M_{\mathrm{warm}}
&=\{\,U_{\mathrm{warm}}(\theta_{\mathrm{w}})|\phi_0\rangle:\theta_{\mathrm{w}}\in\mathbb R^{m_{\mathrm{w}}}\,\},\\
\mathcal M_{\mathrm{refine}}
&=\{\,U_{\mathrm{refine}}(\theta_{\mathrm{r}})U_{\mathrm{warm}}(\theta_{\mathrm{w}}^*)|\phi_0\rangle:\theta_{\mathrm{r}}\in\mathbb R^{m_{\mathrm{r}}}\,\},\\
\mathcal M_{\mathrm{replay}}(\mathcal O_*)
&=\{\,U_{\mathrm{replay}}(\varphi;\mathcal O_*)|\phi_0\rangle:\varphi\in\mathbb R^{m_{\mathrm{rep}}}\,\}.
\end{aligned}
\] The corresponding historical handoff chain is \[
\begin{aligned}
(\theta_{\mathrm{w}}^*,|\psi_{\mathrm{warm}}\rangle)
&\in\mathcal M_{\mathrm{warm}}
\mapsto
(\theta_{\mathrm{r}}^*,|\psi_{\mathrm{refine}}\rangle)
\in\mathcal M_{\mathrm{refine}}\\
&\mapsto
(\mathcal O_*,\theta_*^{\mathrm{adapt}},|\psi_*\rangle)\\
&\mapsto
(\varphi^{(0)},|\psi_{\mathrm{replay}}^{(0)}\rangle)
=
\mathcal H_{\mathrm{replay}}(|\psi_*\rangle,\theta_*^{\mathrm{adapt}},\mathcal O_*,\mathcal C_*)
\in\mathcal M_{\mathrm{replay}}(\mathcal O_*).
\end{aligned}
\] A typical replay-style historical policy is:

\begin{enumerate}
\def\labelenumi{\arabic{enumi}.}
\tightlist
\item
  preserve inherited scaffold parameters,
\item
  initialize residual parameters at zero,
\item
  optionally freeze inherited coordinates for a burn-in stage,
\item
  then unfreeze and refit in a controlled window.
\end{enumerate}

This path is kept for archival provenance only and is not part of the
active main-body symbolic contract.

\subsection{E.3 Drive-amplitude comparison
path}\label{e.3-drive-amplitude-comparison-path}

A drive-amplitude comparison path can be expressed symbolically by
comparing several amplitudes \(A_0,A_1,\dots\) at fixed Hamiltonian
parameters and examining \[
\Delta E_A(t),\qquad \Delta D_A(t),\qquad F_A(t)
\] across the amplitude family. The formal structure is clear even if
the exact report layout is deferred.

\subsection{E.4 Appendix rule}\label{e.4-appendix-rule}

The rule of this symbolic manuscript is therefore:

\begin{enumerate}
\def\labelenumi{\arabic{enumi}.}
\tightlist
\item
  main body = closed mathematical definitions and comparison contracts,
\item
  appendix = promising but not yet fully formalized extensions.
\end{enumerate}

\section{References}\label{references}

\begin{enumerate}
\def\labelenumi{\arabic{enumi}.}
\item
  Grimsley, Economou, Barnes, and Mayhall, An adaptive variational
  algorithm for exact molecular simulations on a quantum computer,
  Nature Communications 10, 3007 (2019), DOI:
  10.1038/s41467-019-10988-2.
\item
  Provost and Vallée, Riemannian structure on manifolds of quantum
  states, Communications in Mathematical Physics 76, 289-301 (1980),
  DOI: 10.1007/BF02193559.
\item
  Stokes, Izaac, Killoran, and Carleo, Quantum Natural Gradient, Quantum
  4, 269 (2020), DOI: 10.22331/q-2020-05-25-269.
\item
  Nocedal and Wright, Trust-Region Methods, in Numerical Optimization,
  Springer, DOI: 10.1007/0-387-22742-3\_4.
\item
  Ramôa, Santos, Mayhall, Barnes, and Economou, Reducing measurement
  costs by recycling the Hessian in adaptive variational quantum
  algorithms, Quantum Science and Technology 10, 015031 (2025), DOI:
  10.1088/2058-9565/ad904e.
\item
  Ramôa, Anastasiou, Santos, Mayhall, Barnes, and Economou, Reducing the
  resources required by ADAPT-VQE using coupled exchange operators and
  improved subroutines, npj Quantum Information (2025), DOI:
  10.1038/s41534-025-01039-4.
\item
  Verteletskyi, Yen, and Izmaylov, Measurement optimization in the
  variational quantum eigensolver using a minimum clique cover, The
  Journal of Chemical Physics 152, 124114 (2020), DOI:
  10.1063/1.5141458.
\item
  Yen, Verteletskyi, and Izmaylov, Measuring All Compatible Operators in
  One Series of Single-Qubit Measurements Using Unitary Transformations,
  Journal of Chemical Theory and Computation (2020), DOI:
  10.1021/acs.jctc.0c00008.
\item
  Ikhtiarudin, Sunnardianto, Fathurrahman, Agusta, and Dipojono,
  Shot-Efficient ADAPT-VQE via Reused Pauli Measurements and
  Variance-Based Shot Allocation, arXiv:2507.16879 (2025).
\item
  Stadelmann, Übelher, Ramôa, Sambasivam, Barnes, and Economou,
  Strategies for Overcoming Gradient Troughs in the ADAPT-VQE Algorithm,
  arXiv:2512.25004 (2025).
\end{enumerate}

\subsubsection{\texorpdfstring{Addendum (2026-04-04): manuscript-aligned
direct HH ADAPT / \texttt{phase3\_v1} sweep
notes}{Addendum (2026-04-04): manuscript-aligned direct HH ADAPT / phase3\_v1 sweep notes}}\label{addendum-2026-04-04-manuscript-aligned-direct-hh-adapt-phase3_v1-sweep-notes}

This addendum records the new manuscript-aligned direct HH ADAPT
behavior on the public \texttt{phase3\_v1} surface, with
\texttt{runtime\_split=off} treated as the canonical manuscript/public
setting.

Physics point for the main recovery study. - \texttt{L=2} -
\texttt{t=1.0} - \texttt{U=4.0} - \texttt{g\_ep=0.5} -
\texttt{omega0=1.0} - \texttt{n\_ph\_max=1} - \texttt{problem=hh} -
binary bosons - blocked ordering - open boundary

Authority targets for interpretation. - Manuscript matched-ADAPT target
from §19.1: \texttt{E\ =\ 0.1587240823603703},
\texttt{\textbar{}ΔE\textbar{}\ =\ 5.6178234644e-05},
\texttt{stop\_reason\ =\ drop\_plateau}. - Fixed-scaffold anchor from
§19.6: \texttt{\textbar{}ΔE\textbar{}\ ≈\ 4.23e-04},
\texttt{F\ ≈\ 0.9998}.

\paragraph{Current status on the new manuscript-aligned
surface}\label{current-status-on-the-new-manuscript-aligned-surface}

The best live strong-coupling beam line on the direct public surface is
\texttt{beam\_runtime\_split-off\_\_batching-off\_\_lifetime-off\_\_profile-tight\_\_gamma-1.0}
with - \texttt{E\ =\ 0.15872511715315915} -
\texttt{\textbar{}ΔE\textbar{}\ =\ 5.721302743347256e-05} -
\texttt{ansatz\_depth\ =\ 20} - \texttt{logical\_num\_parameters\ =\ 20}
- \texttt{num\_parameters\ =\ 48} -
\texttt{stop\_reason\ =\ drop\_plateau} -
\texttt{exact\_state\_fidelity\ =\ 0.9999715655261739}

This line clearly beats the fixed-scaffold anchor, but it does not yet
numerically match the manuscript matched-ADAPT target. The residual gap
to the manuscript matched-ADAPT line is \texttt{1.0347927889e-06} in
\texttt{\textbar{}ΔE\textbar{}}.

A historical diagnostic reoptimization run with windowed \texttt{COBYLA}
reached - \texttt{E\ =\ 0.15872468271899687} -
\texttt{\textbar{}ΔE\textbar{}\ =\ 5.6778593271189504e-05} - residual
gap to the manuscript matched-ADAPT line: \texttt{6.0035862640e-07}

This is useful as a numerical diagnostic, but it is not the forward
recipe for new runs. Going forward, the intended optimizer family is
\texttt{SPSA}, not \texttt{COBYLA}.

\paragraph{Executed run families and
outcomes}\label{executed-run-families-and-outcomes}

\begin{enumerate}
\def\labelenumi{\arabic{enumi}.}
\tightlist
\item
  Strong-coupling beam sweep.
\end{enumerate}

\begin{itemize}
\tightlist
\item
  Artifact root:
  \texttt{artifacts/agent\_runs/20260403\_hh\_l2\_u4\_g05\_beam\_sweep\_stage1/}
\item
  Best case:
  \texttt{beam\_runtime\_split-off\_\_batching-off\_\_lifetime-off\_\_profile-tight\_\_gamma-1.0}
\item
  Result: \texttt{\textbar{}ΔE\textbar{}\ =\ 5.721302743347256e-05}
\item
  Interpretation: beam search nearly recovers the manuscript
  matched-ADAPT quality.
\end{itemize}

\begin{enumerate}
\def\labelenumi{\arabic{enumi}.}
\setcounter{enumi}{1}
\tightlist
\item
  Strong-coupling runtime sweep.
\end{enumerate}

\begin{itemize}
\tightlist
\item
  Artifact root:
  \texttt{artifacts/agent\_runs/20260403\_hh\_l2\_u4\_g05\_runtime\_sweep\_stage1/}
\item
  Best case still bad:
  \texttt{\textbar{}ΔE\textbar{}\ =\ 0.3282923042519191}
\item
  Interpretation: the single-trajectory runtime line collapses into the
  same bad basin seen in the manuscript novelty-off failure scale.
\end{itemize}

\begin{enumerate}
\def\labelenumi{\arabic{enumi}.}
\setcounter{enumi}{2}
\tightlist
\item
  Phase-3 beam-vs-batching diagnostic.
\end{enumerate}

\begin{itemize}
\tightlist
\item
  Artifact root:
  \texttt{artifacts/agent\_runs/20260403\_hh\_l2\_u4\_g05\_phase3\_diag\_beam\_vs\_batch\_v1/}
\item
  Single-trajectory cases stayed at
  \texttt{\textbar{}ΔE\textbar{}\ ≈\ 0.3282923}.
\item
  Beam \texttt{b3,c2,k3} immediately recovered
  \texttt{\textbar{}ΔE\textbar{}\ =\ 5.721302743347256e-05}.
\item
  Interpretation: the decisive lever is frontier diversity from beam
  search, not batching.
\end{itemize}

\begin{enumerate}
\def\labelenumi{\arabic{enumi}.}
\setcounter{enumi}{3}
\tightlist
\item
  Beam-capacity recovery slice.
\end{enumerate}

\begin{itemize}
\tightlist
\item
  Artifact root:
  \texttt{artifacts/agent\_runs/20260403\_hh\_l2\_u4\_g05\_phase3\_beam\_capacity\_recovery\_v1/}
\item
  \texttt{b3,c2,k3} and \texttt{b4,c2,k4} tied at
  \texttt{\textbar{}ΔE\textbar{}\ =\ 5.721302743347256e-05}.
\item
  Interpretation: larger beam width by itself does not improve the
  recovered solution once the beam is already large enough.
\end{itemize}

\begin{enumerate}
\def\labelenumi{\arabic{enumi}.}
\setcounter{enumi}{4}
\tightlist
\item
  Prune/shortlist micro-sweep.
\end{enumerate}

\begin{itemize}
\tightlist
\item
  Artifact root:
  \texttt{artifacts/agent\_runs/20260404\_hh\_l2\_u4\_g05\_phase3\_prune\_shortlist\_micro\_v1/}
\item
  All cases were numerically identical to saved precision.
\item
  Interpretation: \texttt{phase1\_prune\_fraction},
  \texttt{phase1\_prune\_max\_candidates}, and shortlist width are not
  the active driver of the remaining manuscript gap on this beam-fixed
  line.
\end{itemize}

\begin{enumerate}
\def\labelenumi{\arabic{enumi}.}
\setcounter{enumi}{5}
\tightlist
\item
  Reoptimization slice.
\end{enumerate}

\begin{itemize}
\tightlist
\item
  Artifact root:
  \texttt{artifacts/agent\_runs/20260404\_hh\_l2\_u4\_g05\_phase3\_reopt\_slice\_v1/}
\item
  Windowed \texttt{COBYLA} improved the gap to
  \texttt{\textbar{}ΔE\textbar{}\ =\ 5.6778593271189504e-05}.
\item
  \texttt{POWELL} windowed stayed at the original beam winner.
\item
  \texttt{POWELL} periodic/full refits got worse.
\item
  \texttt{append\_only} failed badly at
  \texttt{\textbar{}ΔE\textbar{}\ =\ 0.6244148638396572}.
\item
  Interpretation: optimizer/refit policy can move the line, but the
  \texttt{COBYLA} improvement should be treated as diagnostic only, not
  as the new manuscript-forward recipe.
\end{itemize}

\begin{enumerate}
\def\labelenumi{\arabic{enumi}.}
\setcounter{enumi}{6}
\tightlist
\item
  Score-weight slice.
\end{enumerate}

\begin{itemize}
\tightlist
\item
  Artifact root:
  \texttt{artifacts/agent\_runs/20260404\_hh\_l2\_u4\_g05\_phase3\_score\_weight\_slice\_v1/}
\item
  \texttt{near=0.90} gave a small improvement to
  \texttt{\textbar{}ΔE\textbar{}\ =\ 5.721150236287498e-05} and also
  reduced the scaffold to \texttt{logical\_num\_parameters\ =\ 19},
  \texttt{num\_parameters\ =\ 46}, \texttt{ansatz\_depth\ =\ 19}.
\item
  Other \texttt{rho} and \texttt{lambda\_H} variants were essentially
  flat at the saved precision.
\end{itemize}

\begin{enumerate}
\def\labelenumi{\arabic{enumi}.}
\setcounter{enumi}{7}
\tightlist
\item
  Seed robustness slice.
\end{enumerate}

\begin{itemize}
\tightlist
\item
  Artifact root:
  \texttt{artifacts/agent\_runs/20260404\_hh\_l2\_u4\_g05\_phase3\_seed\_robustness\_slice\_v1/}
\item
  Seeds \texttt{5,\ 7,\ 11,\ 13,\ 17} all reproduced the same
  \texttt{POWELL} baseline result exactly to saved precision.
\item
  Interpretation: the current beam-fixed \texttt{POWELL} plateau is
  deterministic across this seed set.
\end{itemize}

\begin{enumerate}
\def\labelenumi{\arabic{enumi}.}
\setcounter{enumi}{8}
\tightlist
\item
  Weak-coupling diagnostic lanes.
\end{enumerate}

\begin{itemize}
\tightlist
\item
  Canonical weak sweep root:
  \texttt{artifacts/agent\_runs/20260403\_hh\_l2\_u05\_g02\_weak\_logical\_scaffold\_canonical\_sweep\_v1/}
\item
  Weak beam sweep root:
  \texttt{artifacts/agent\_runs/20260403\_hh\_l2\_u05\_g02\_beam\_sweep\_stage1\_parent\_full\_meta/}
\item
  Best weak result in both lanes stayed at
  \texttt{\textbar{}ΔE\textbar{}\ ≈\ 0.0131029358626}.
\item
  Parent-quality target from §19.3.1 is much smaller, so weak coupling
  remains a diagnostic regime rather than a recovered manuscript-quality
  regime.
\end{itemize}

\paragraph{Good parameters versus bad
parameters}\label{good-parameters-versus-bad-parameters}

Parameters/choices that consistently helped or were at least required
for recovery. - Direct HH ADAPT on the public \texttt{phase3\_v1}
surface. - \texttt{runtime\_split=off}. - Beam search on. - Beam at or
above \texttt{live\_branches=3}, \texttt{children\_per\_parent=2},
\texttt{terminated\_keep=3}. - \texttt{gamma\_N=1.0} on the public
strong-coupling line. - \texttt{batching=off} was slightly better than
\texttt{batching=on} in the best recovered beam lane. -
\texttt{lifetime\_cost\_mode=off} was at least as good as
\texttt{phase3\_v1} within the explored resolution. - \texttt{near=0.90}
was the only score-weight variation that gave a measurable improvement
in the explored \texttt{POWELL} lane.

Parameters/choices that were neutral or nearly inert in the explored
region. - Beam width beyond the \texttt{b3,c2,k3} threshold. -
\texttt{profile=tight} versus \texttt{profile=wide}. -
\texttt{phase1\_prune\_fraction} in the explored \texttt{0.15} to
\texttt{0.35} range. - \texttt{phase1\_prune\_max\_candidates} in the
explored \texttt{4} to \texttt{8} range. - shortlist width in the
explored micro-sweep. - Most \texttt{rho} and \texttt{lambda\_H}
score-weight perturbations.

Parameters/choices that consistently hurt. - Single-trajectory runtime
mode with no beam. - Runtime line behavior that reproduces
\texttt{\textbar{}ΔE\textbar{}\ ≈\ 0.3282923}. -
\texttt{runtime\_split=shortlist\_pauli\_children\_v1} relative to
\texttt{off} on this manuscript/public lane. - \texttt{POWELL}
periodic/full refits relative to the windowed baseline. -
\texttt{append\_only} refit policy, which failed catastrophically.

\paragraph{Practical interpretation}\label{practical-interpretation}

For the manuscript-aligned direct HH ADAPT path at strong coupling, the
main qualitative lesson is now clear: beam-induced frontier diversity is
the key ingredient that recovers the near-manuscript solution, while
pruning pressure, shortlist width, profile choice, and modest
score-weight changes are second-order. The live public-surface recipe
already beats the fixed-scaffold anchor by a large margin, but it still
does not numerically match the manuscript matched-ADAPT target.

\paragraph{2026-04-04 update: public-CLI runtime weights and
reduced-pool
sweeps}\label{update-public-cli-runtime-weights-and-reduced-pool-sweeps}

After exposing the manuscript-facing selector/runtime knobs on the
public direct \texttt{phase3\_v1} CLI, two additional conclusions
emerged for the same strong-coupling \(L=2\) point.

\begin{enumerate}
\def\labelenumi{\arabic{enumi}.}
\tightlist
\item
  Public-CLI \texttt{full\_meta\ +\ SPSA} improved once frontier ratios
  were decoupled from the near-degenerate batch shell.
\end{enumerate}

\begin{itemize}
\tightlist
\item
  Artifact root:
  \texttt{artifacts/agent\_runs/20260404\_hh\_l2\_u4\_g05\_phase3\_public\_weights\_frontier\_spsa\_v1/}
\item
  Best public-CLI \texttt{full\_meta\ +\ SPSA} line: \[
  E=0.15878350593338042,\qquad |\Delta E|=1.1560180765\times 10^{-4},
  \] with \texttt{logical\_num\_parameters=21},
  \texttt{num\_parameters=53}, \texttt{ansatz\_depth=21}, and
  \texttt{fidelity=0.9999480673}.
\item
  The corresponding \texttt{weights\_zero} line was only slightly worse,
  \[
  |\Delta E|=1.1722877711\times 10^{-4},
  \] while the tighter coupled-frontier setting \[
  (\text{phase2 frontier},\text{phase3 frontier})=(0.98,0.98)
  \] reproduced the older worse basin \[
  |\Delta E|=1.4383993770\times 10^{-4}.
  \]
\item
  So the main gain from that patch was not a delicate burden-weight
  choice; it was fixing the implementation coupling between shortlist
  frontiering and the batch-shell threshold so that
  \texttt{phase2\_frontier\_ratio}, \texttt{phase3\_frontier\_ratio},
  and \texttt{eta\_\{nd\}} became genuinely independent controls.
\end{itemize}

\begin{enumerate}
\def\labelenumi{\arabic{enumi}.}
\setcounter{enumi}{1}
\tightlist
\item
  The reduced-pool \texttt{pareto\_lean\ +\ SPSA} lane remained rigid.
\end{enumerate}

\begin{itemize}
\tightlist
\item
  Artifact roots:

  \begin{itemize}
  \tightlist
  \item
    \texttt{artifacts/agent\_runs/20260404\_hh\_l2\_u4\_g05\_phase3\_pool\_family\_spsa\_fixed\_v1/}
  \item
    \texttt{artifacts/agent\_runs/20260404\_hh\_l2\_u4\_g05\_phase3\_pareto\_lean\_gamma\_spsa\_fixed\_v1/}
  \item
    \texttt{artifacts/agent\_runs/20260404\_hh\_l2\_u4\_g05\_phase3\_pareto\_lean\_drop\_policy\_spsa\_fixed\_v1/}
  \item
    \texttt{artifacts/agent\_runs/20260404\_hh\_l2\_u4\_g05\_pareto\_lean\_new\_cli\_weights\_spsa\_v1/}
  \item
    \texttt{artifacts/agent\_runs/20260404\_hh\_l2\_u4\_g05\_pareto\_lean\_compat\_proxy\_spsa\_v1/}
  \item
    \texttt{artifacts/agent\_runs/20260404\_hh\_l2\_u4\_g05\_pareto\_lean\_frontier\_ratio\_spsa\_v1/}
  \end{itemize}
\item
  Baseline reduced-pool result: \[
  E=0.1587895750333393,\qquad |\Delta E|=1.2167090761\times 10^{-4},
  \] with \texttt{logical\_num\_parameters=21},
  \texttt{num\_parameters=52}, \texttt{ansatz\_depth=21}, and
  \texttt{fidelity=0.9999443566}.
\item
  This value persisted across all tested variations of:

  \begin{itemize}
  \tightlist
  \item
    \texttt{gamma\_N} in the explored \texttt{0.75} to \texttt{1.5}
    range,
  \item
    drop-policy parameters,
  \item
    motif bonus,
  \item
    duplicate penalty,
  \item
    leakage penalty \texttt{eta\_L},
  \item
    compatibility-prescreen weights,
  \item
    remaining-evaluations proxy mode,
  \item
    frontier ratios from \texttt{(0.85,0.85)} up through
    \texttt{(0.95,0.95)} and the asymmetric \texttt{(0.90,0.75)},
    \texttt{(0.90,0.95)}, \texttt{(0.95,0.90)} cases.
  \end{itemize}
\item
  Loosening the phase-2 frontier too far did hurt: \texttt{(0.75,0.75)}
  gave \[
  |\Delta E|=1.5234589586\times 10^{-4},
  \] and \texttt{(0.75,0.90)} gave \[
  |\Delta E|=1.5457582580\times 10^{-4}.
  \]
\item
  Hence the reduced \texttt{pareto\_lean} line appears structurally
  plateaued under the currently exposed runtime-law knobs: it is not
  rescued by more local tuning of \texttt{gamma\_N}, batching pressure,
  motif/repetition heuristics, or compatibility/proxy terms.
\end{itemize}

\begin{enumerate}
\def\labelenumi{\arabic{enumi}.}
\setcounter{enumi}{2}
\tightlist
\item
  Cost-aware ranking versus raw-energy ranking now separate.
\end{enumerate}

\begin{itemize}
\tightlist
\item
  On raw strong-coupling energy error alone, the best historical rerank
  observed in these sweeps was the diagnostic \texttt{COBYLA} line \[
  E=0.15872468271899687,\qquad |\Delta E|=5.6778593271\times 10^{-5},
  \] with \texttt{logical\_num\_parameters=20},
  \texttt{num\_parameters=50}, \texttt{ansatz\_depth=20}.
\item
  However, the best currently observed \(|\Delta E|/K\)-style point
  among runs that reported logical scaffold size was the
  \texttt{near=0.90} strong-coupling line from the score-weight slice:
  \[
  E=0.15872511562808855,\qquad |\Delta E|=5.7211502363\times 10^{-5},
  \] with \[
  K_{\text{logical}}=19,\qquad K_{\text{runtime}}=46,\qquad d=19.
  \]
\item
  So the present picture is:

  \begin{itemize}
  \tightlist
  \item
    best raw-\(|\Delta E|\) line: slightly better than the old beam
    winner, but achieved on a non-preferred optimizer branch,
  \item
    best cost-aware line: slightly worse in \(|\Delta E|\) but leaner in
    both logical scaffold size and runtime parameter count,
  \item
    best reduced-pool SPSA line: still stuck near \(1.22\times 10^{-4}\)
    and therefore materially behind the broad-beam public lane.
  \end{itemize}
\end{itemize}

\paragraph{\texorpdfstring{2026-04-04 update: completed reduced-pool
\texttt{L=2/L=3}
matrix}{2026-04-04 update: completed reduced-pool L=2/L=3 matrix}}\label{update-completed-reduced-pool-l2l3-matrix}

The overnight reduced-pool matrix
\texttt{artifacts/agent\_runs/20260404\_hh\_phase3\_reduced\_pool\_matrix\_v1/}
completed through the strong confirmation and weak-transfer waves. Its
campaign summary and final strong templates are recorded in

\begin{itemize}
\tightlist
\item
  \texttt{artifacts/agent\_runs/20260404\_hh\_phase3\_reduced\_pool\_matrix\_v1/campaign\_summary.json},
\item
  \texttt{artifacts/agent\_runs/20260404\_hh\_phase3\_reduced\_pool\_matrix\_v1/final\_strong\_templates\_by\_L.json},
\item
  \texttt{artifacts/agent\_runs/20260404\_hh\_phase3\_reduced\_pool\_matrix\_v1/matrix\_leaderboard.json}.
\end{itemize}

Data field: the completed matrix contained \texttt{212} successful
cases, \texttt{36} failed cases, and \texttt{184} unique successful
configurations.

For the strong-coupling \(L=2\) target point, the best matrix result was
\[
E=0.15877612622031317,\qquad |\Delta E|=1.0822209459\times 10^{-4},
\] with \[
K_{\mathrm{logical}}=20,\qquad K_{\mathrm{runtime}}=52,\qquad d=20,\qquad F=0.9999519211,
\] from the confirmed \texttt{pareto\_lean} template with \[
(\text{beam},\text{frontier})=(2,2,2;\,0.85,0.85).
\] This improved over the earlier reduced-pool SPSA plateau near
\(1.2167\times10^{-4}\), but it still did not recover the
manuscript-quality strong-coupling target at
\(5.6178234644\times10^{-5}\) and did not beat the earlier broad-lane
live recovery line near \(5.7213027433\times10^{-5}\).

For the strong-coupling \(L=3\) point, the reduced-pool matrix failed
badly. Its best confirmed case was \[
E=-0.2564741176143438,\qquad |\Delta E|=5.0141481774\times 10^{-1},
\] with fidelity only \[
F\approx 5.81\times 10^{-8}.
\] The best matrix lane there was \texttt{uccsd\_paop\_lf\_full}, not
the more aggressively reduced families. Hence the completed reduced-pool
matrix should be read as negative evidence for the current public
reduced execution surface at \(L=3\), not as evidence that the
historical \(L=3\) HH target itself is unattainable.

For weak coupling, the transferred \(L=2\) reduced-pool line remained
poor: \[
E=-0.7632521982953032,\qquad |\Delta E|=1.3125332609\times10^{-2},
\] so the matrix did not close the weak-coupling parent-quality gap
there. By contrast, the transferred \(L=3\) weak baseline was
numerically much better, \[
E=-1.0439438365685594,\qquad |\Delta E|=4.1313925567\times10^{-4},
\] with fidelity \[
F=0.9998116001.
\] That \(L=3\) weak result is still not a manuscript anchor, but it
does show that the reduced families are not uniformly broken across all
\(L=3\) HH regimes; the main failure is the strong-coupling \(L=3\)
execution surface.

To make that contrast explicit, the surviving historical broad-pool
\(L=3\) reference from
\texttt{artifacts/reports/hh\_heavy\_scaffold\_best\_yet\_20260321.md}
was \[
E=0.24502564220194084,\qquad |\Delta E|=8.4942074026\times10^{-5},
\] while the matched lean comparison report
\texttt{artifacts/reports/pareto\_lean\_vs\_heavy\_L3\_comparison.md}
recorded the lean/class-pruned line at about \[
|\Delta E|\approx 1.15\times10^{-4}.
\] Both are therefore qualitatively incompatible with the present
reconstructed direct-surface reruns, which all collapsed to the same bad
basin.

In particular, the explicit three-way reproduction run
\texttt{artifacts/agent\_runs/20260404\_hh\_l3\_u4\_g05\_historical\_surface\_pool\_compare\_v1/}
reconstructed the documented historical settings \[
\texttt{phase3\_v1},\ \texttt{POWELL},\ \texttt{phase1 shortlist}=256,\ \texttt{phase2 shortlist}=128,\ \texttt{split}=\texttt{shortlist\_pauli\_children\_v1},
\] and compared

\begin{itemize}
\tightlist
\item
  \texttt{full\_meta},
\item
  class-filtered \texttt{full\_meta},
\item
  \texttt{pareto\_lean\_l3}.
\end{itemize}

All three produced the same failed result, \[
E=-0.25838912148429044,\qquad |\Delta E|=5.0332982161\times10^{-1},
\] with negligible fidelity. Hence the current issue is not simply which
reduced family is chosen. Rather, some additional execution-surface
ingredient from the older successful \(L=3\) lane is still missing from
the present direct CLI reconstruction.

\paragraph{Appendix: raw HF-start direct ADAPT burdens versus historical
post-processed / fixed-scaffold
submissions}\label{appendix-raw-hf-start-direct-adapt-burdens-versus-historical-post-processed-fixed-scaffold-submissions}

To avoid conflating the heavy raw recent ADAPT replay artifacts with the
older locked fixed-manifold scaffold family, the following tables record
the exact transpile-style cost snapshots discussed during the April 7
accounting pass.

These counts are all on the same local hardware proxy surface, namely
\texttt{FakeMarrakesh}, and they should be interpreted as scaffold
burden rather than full campaign burden. In particular, the recent raw
ADAPT ladders below are \emph{not} the old locked 6/7-term artifacts.

Throughout this appendix, \textbf{raw HF-start direct ADAPT} means
\texttt{adapt\_ref\_json\ =\ null} with no later JSON prune/recovery
substitution. The post-processed \texttt{31/37/57}-2Q readapt/prune
artifacts listed later are comparison-only and are not the object a
fresh real-QPU ADAPT run from HF would execute.

\subparagraph{A. Raw HF-start direct ADAPT / replay scaffold burden
snapshots}\label{a.-raw-hf-start-direct-adapt-replay-scaffold-burden-snapshots}

{\def\LTcaptype{none} % do not increment counter
\begin{longtable}[]{@{}
  >{\raggedright\arraybackslash}p{(\linewidth - 12\tabcolsep) * \real{0.1111}}
  >{\raggedleft\arraybackslash}p{(\linewidth - 12\tabcolsep) * \real{0.1481}}
  >{\raggedleft\arraybackslash}p{(\linewidth - 12\tabcolsep) * \real{0.1481}}
  >{\raggedleft\arraybackslash}p{(\linewidth - 12\tabcolsep) * \real{0.1481}}
  >{\raggedleft\arraybackslash}p{(\linewidth - 12\tabcolsep) * \real{0.1481}}
  >{\raggedleft\arraybackslash}p{(\linewidth - 12\tabcolsep) * \real{0.1481}}
  >{\raggedleft\arraybackslash}p{(\linewidth - 12\tabcolsep) * \real{0.1481}}@{}}
\toprule\noalign{}
\begin{minipage}[b]{\linewidth}\raggedright
Case
\end{minipage} & \begin{minipage}[b]{\linewidth}\raggedleft
Logical blocks
\end{minipage} & \begin{minipage}[b]{\linewidth}\raggedleft
Runtime rotations
\end{minipage} & \begin{minipage}[b]{\linewidth}\raggedleft
Raw \texttt{cx}
\end{minipage} & \begin{minipage}[b]{\linewidth}\raggedleft
Transpiled 2Q
\end{minipage} & \begin{minipage}[b]{\linewidth}\raggedleft
Depth
\end{minipage} & \begin{minipage}[b]{\linewidth}\raggedleft
Size
\end{minipage} \\
\midrule\noalign{}
\endhead
\bottomrule\noalign{}
\endlastfoot
Noisy weak raw run (\texttt{20260407} ladder) & 19 & 47 & 142 & 221 &
631 & 922 \\
Noisy mid raw run (\texttt{20260407} ladder) & 21 & 53 & 154 & 256 & 697
& 1084 \\
Noisy high raw run (\texttt{20260407} ladder) & 23 & 56 & 160 & 262 &
764 & 1113 \\
Noiseless Math front-page anchor
(\texttt{w5\_l2\_strong\_pareto\_lean\_b2c2k2\_f0p85\_0p85\_f1\_seed5})
& 20 & 52 & 152 & 218 & 633 & 939 \\
\end{longtable}
}

Interpretation: the validated front-page
\texttt{pareto\_lean\ +\ phase3\_v1} deployment anchor is still a heavy
raw HF-start ADAPT replay scaffold on Marrakesh (\texttt{218} two-qubit
gates), i.e.~it sits near the recent weak raw-noise lane and far above
the post-processed prune/readapt and locked fixed-scaffold families. The
lighter \texttt{81}-2Q direct HF-start winner is the separate scientific
winner listed on the front page.

\subparagraph{B. Historical post-processed or fixed-scaffold runtime
submissions}\label{b.-historical-post-processed-or-fixed-scaffold-runtime-submissions}

The next table isolates the \(L=2\) runtime submissions that were
actually sent to the real \texttt{ibm\_marrakesh} backend. The listed
two-qubit counts are obtained by transpiling the exact submitted
scaffold artifacts on \texttt{FakeMarrakesh} with the same transpiler
settings recorded in the runtime-result JSON.

{\def\LTcaptype{none} % do not increment counter
\begin{longtable}[]{@{}
  >{\raggedright\arraybackslash}p{(\linewidth - 10\tabcolsep) * \real{0.1500}}
  >{\raggedright\arraybackslash}p{(\linewidth - 10\tabcolsep) * \real{0.1500}}
  >{\raggedright\arraybackslash}p{(\linewidth - 10\tabcolsep) * \real{0.1500}}
  >{\raggedleft\arraybackslash}p{(\linewidth - 10\tabcolsep) * \real{0.2000}}
  >{\raggedleft\arraybackslash}p{(\linewidth - 10\tabcolsep) * \real{0.2000}}
  >{\raggedright\arraybackslash}p{(\linewidth - 10\tabcolsep) * \real{0.1500}}@{}}
\toprule\noalign{}
\begin{minipage}[b]{\linewidth}\raggedright
Real-QPU submission
\end{minipage} & \begin{minipage}[b]{\linewidth}\raggedright
Runtime source artifact
\end{minipage} & \begin{minipage}[b]{\linewidth}\raggedright
Runtime transpile settings
\end{minipage} & \begin{minipage}[b]{\linewidth}\raggedleft
Transpiled 2Q
\end{minipage} & \begin{minipage}[b]{\linewidth}\raggedleft
Depth
\end{minipage} & \begin{minipage}[b]{\linewidth}\raggedright
Notes
\end{minipage} \\
\midrule\noalign{}
\endhead
\bottomrule\noalign{}
\endlastfoot
\texttt{hh\_gatepruned7\_fixedtheta\_runtime\_ibm\_marrakesh\_20260323T145219Z\_direct}
& \texttt{artifacts/json/hh\_prune\_nighthawk\_gate\_pruned\_7term.json}
& \texttt{opt=1}, \texttt{seed=42} & 25 & 63 & Locked 7-term fixed-theta
Marrakesh submission \\
\texttt{hh\_readapt5\_energyonly\_runtime\_ibm\_marrakesh\_20260323T153244Z}
& \texttt{artifacts/json/hh\_prune\_nighthawk\_readapt\_5op.json} &
\texttt{opt=1}, \texttt{seed=7} & 37 & 126 & Post-prune/recovery runtime
candidate; not fresh HF-start ADAPT \\
\texttt{hh\_gatepruned6\_fixedtheta\_runtime\_ibm\_marrakesh\_20260323T171528Z}
&
\texttt{artifacts/json/hh\_marrakesh\_fixed\_scaffold\_6term\_drop\_eyezee\_20260323T171528Z.json}
& \texttt{opt=1}, \texttt{seed=7} & 14 & 48 & Locked 6-term fixed-theta
Marrakesh submission \\
\end{longtable}
}

So the historical real-QPU gate counts for the three recorded Marrakesh
submissions were \[
25,\qquad 37,\qquad 14
\] two-qubit gates respectively for the locked 7-term, readapt-5, and
locked 6-term submissions.

\section{Appendix F. Archived L=2 Noisy Surface Run
Handoff}\label{appendix-f.-archived-l2-noisy-surface-run-handoff}

For the \textbf{local/offline L=2 HH noisy surface} (\texttt{phase3\_v1}
+ \texttt{pareto\_lean\_l2} + \texttt{FakeNighthawk}), see
\href{notes/l2_noisy_surface_run.md}{MATH/notes/l2\_noisy\_surface\_run.md}.

This is the terminal-agent handoff note for the three mitigation lanes:

\begin{itemize}
\tightlist
\item
  \texttt{none}
\item
  \texttt{readout}
\item
  \texttt{full}
\end{itemize}

This recipe is now archived here rather than pinned to the front page,
because it is a historical run handoff rather than the active
manuscript-level formalism.

\end{document}
